\documentclass[10pt,a4paper]{article}

\usepackage{iftex}
\ifPDFTeX
  \usepackage[utf8]{inputenc}
  \usepackage[T1]{fontenc}
\else
  \usepackage{fontspec}
  \defaultfontfeatures{Ligatures=TeX}
  \IfFontExistsTF{Consolas}{%
    \setmonofont{Consolas}[Scale=0.84]
  }{%
    \IfFileExists{C:/Windows/Fonts/consola.ttf}{%
      \setmonofont{C:/Windows/Fonts/consola.ttf}[Scale=0.84]
    }{}
  }
\fi

\usepackage[catalan]{babel}
\usepackage[a4paper,margin=1.8cm,headheight=14pt]{geometry}
\usepackage{array}
\usepackage{booktabs}
\usepackage{caption}
\usepackage{enumitem}
\usepackage{float}
\usepackage{hyperref}
\usepackage{longtable}
\usepackage{microtype}
\usepackage{needspace}
\usepackage{textcomp}
\usepackage{xcolor}
\usepackage{listings}
\usepackage{tikz}
\usetikzlibrary{arrows.meta,calc,fit,matrix,positioning,shapes.geometric,shapes.multipart}

\definecolor{StudioInk}{HTML}{1F2A2E}
\definecolor{StudioMuted}{HTML}{5C6F6B}
\definecolor{StudioLine}{HTML}{60746F}
\definecolor{StudioBorder}{HTML}{8DA7A1}
\definecolor{StudioFill}{HTML}{F8FBFA}
\definecolor{StudioAccent}{HTML}{EAF4F1}
\definecolor{StudioBlue}{HTML}{2E6F95}
\definecolor{StudioBlueFill}{HTML}{EEF6FF}
\definecolor{StudioAmber}{HTML}{A06D2C}
\definecolor{StudioAmberFill}{HTML}{FFF8EF}
\definecolor{CodeBg}{HTML}{FBFCFC}
\definecolor{CodeFrame}{HTML}{CBD9D6}
\definecolor{CodeKeyword}{HTML}{205D7A}
\definecolor{CodeString}{HTML}{7A4B20}
\definecolor{CodeComment}{HTML}{5F7A72}

\hypersetup{
  colorlinks=true,
  linkcolor=StudioBlue,
  urlcolor=StudioBlue,
  citecolor=StudioBlue,
  pdftitle={Annex de codi de BEMS-RL Studio},
  pdfauthor={Guillem Torm}
}

\captionsetup{font=small,labelfont=bf}
\setlist[itemize]{leftmargin=1.4em,itemsep=0.25em}
\setlist[enumerate]{leftmargin=1.6em,itemsep=0.25em}
\setlength{\parskip}{0.55em}
\setlength{\parindent}{0pt}
\setlength{\emergencystretch}{2em}
\newcommand{\diagramnote}[1]{\par\smallskip\noindent\textbf{Lectura del diagrama.} #1\par\medskip}

\tikzset{
  studioBox/.style={
    draw=StudioBorder,
    fill=StudioFill,
    rounded corners=3pt,
    align=center,
    minimum width=2.65cm,
    minimum height=0.9cm,
    inner sep=5pt,
    font=\small
  },
  studioWide/.style={studioBox, minimum width=3.6cm},
  studioBlue/.style={studioBox, draw=StudioBlue, fill=StudioBlueFill},
  studioAmber/.style={studioBox, draw=StudioAmber, fill=StudioAmberFill},
  studioCluster/.style={draw=StudioBorder, fill=StudioAccent, rounded corners=4pt, inner sep=8pt},
  studioEdge/.style={-{Latex[length=2.1mm]}, thick, draw=StudioLine},
  studioDashed/.style={studioEdge, dashed},
  classBox/.style={
    draw=StudioBorder,
    fill=StudioFill,
    rectangle split,
    rectangle split parts=3,
    rounded corners=2pt,
    align=left,
    text width=4.1cm,
    font=\scriptsize,
    inner sep=4pt
  }
}

\lstdefinestyle{studioPython}{
  language=Python,
  basicstyle=\ttfamily\scriptsize,
  keywordstyle=\bfseries\color{CodeKeyword},
  stringstyle=\color{CodeString},
  commentstyle=\itshape\color{CodeComment},
  backgroundcolor=\color{CodeBg},
  frame=single,
  rulecolor=\color{CodeFrame},
  framerule=0.4pt,
  numbers=left,
  numberstyle=\tiny\color{StudioMuted},
  stepnumber=1,
  numbersep=7pt,
  showstringspaces=false,
  keepspaces=true,
  columns=fullflexible,
  breaklines=true,
  breakatwhitespace=false,
  tabsize=4,
  upquote=true,
  captionpos=t,
  inputencoding=utf8,
  extendedchars=true,
  literate=*
    {°}{{$^\circ$}}1
    {²}{{$^2$}}1
    {·}{{$\cdot$}}1
    {À}{{\`A}}1
    {Â}{{\^A}}1
    {Ã}{{\~A}}1
    {Ó}{{\'O}}1
    {Ú}{{\'U}}1
    {à}{{\`a}}1
    {â}{{\^a}}1
    {ç}{{\c{c}}}1
    {è}{{\`e}}1
    {é}{{\'e}}1
    {í}{{\'i}}1
    {ï}{{\"i}}1
    {ò}{{\`o}}1
    {ó}{{\'o}}1
    {ú}{{\'u}}1
    {ā}{{\=a}}1
    {Δ}{{$\Delta$}}1
    {×}{{$\times$}}1
    {–}{{--}}1
    {—}{{---}}1
    {’}{{'}}1
    {•}{{$\bullet$}}1
    {…}{{\ldots}}1
    {→}{{$\to$}}1
    {−}{{$-$}}1
    {≈}{{$\approx$}}1
    {≥}{{$\ge$}}1
    {─}{{-}}1
    {▲}{{$\triangle$}}1
    {▸}{{$\triangleright$}}1
    {▼}{{$\triangledown$}}1
    {◫}{{$\square$}}1
    {☀}{{sun}}1
    {✅}{{ok}}1
    {❄}{{snow}}1
    {️}{{}}0
    {🌸}{{spring}}1
    {🍂}{{autumn}}1
    {📁}{{folder}}1
    {📄}{{file}}1
    {🔥}{{heat}}1
}


\title{\textbf{Annex de codi de BEMS-RL Studio}\\
\large Arquitectura, diagrames UML i codi funcional complet}
\author{Guillem Torm}
\date{© 2025, Guillem Torm}

\begin{document}
\maketitle
\tableofcontents
\clearpage

\section{Objectiu i abast}

Aquest annex acompanya la documentació professional de BEMS-RL Studio. El seu
objectiu és donar una visió tècnica suficient per entendre com està estructurat
el projecte i, després, conservar una còpia compilable de tot el codi Python
funcional de l'aplicació.

L'annex inclou 70 fitxers Python i 21,618 línies de codi.
Abans de les llistes completes s'hi incorporen els diagrames principals de la
documentació Sphinx, adaptats a LaTeX/TikZ, i un UML conceptual dels tipus de
dades que vertebren els workflows.

El codi reproduït en aquest annex correspon exclusivament a BEMS-RL Studio.
No s'hi copia ni s'hi reprodueix codi font del paquet Sinergym: les aparicions
del nom Sinergym dins del document indiquen integracions, importacions,
artefactes, rutes o dependències utilitzades per l'aplicació.

\subsection{Criteri d'inclusió}

\begin{itemize}
  \item S'inclou el codi funcional de pàgines, components, backend, generació de
        gràfics, informes, eines de manteniment i punts d'entrada.
  \item No s'inclouen els mòduls d'estil CSS pur: \texttt{page\_styles/*.py} i
        \texttt{ui\_theme.py}. El seu contingut és visual i ja queda cobert per
        la documentació d'interfície.
  \item El codi queda incrustat en el fitxer LaTeX generat. Les icones Unicode
        del codi es mostren com a escapes per evitar glifs buits en el PDF,
        però els fitxers originals no es modifiquen.
\end{itemize}

\subsection{Inventari resumit}

\begin{longtable}{>{\raggedright\arraybackslash}p{0.28\linewidth}rr}
\toprule
Grup & Fitxers & Línies \\
\midrule
\endfirsthead
\toprule
Grup & Fitxers & Línies \\
\midrule
\endhead
Entrypoints i navegacio & 3 & 309 \\
Pagines Streamlit & 11 & 4,753 \\
Components reutilitzables & 7 & 2,503 \\
Backend funcional & 30 & 9,954 \\
Grafics i metriques & 19 & 4,099 \\
\bottomrule
\end{longtable}

\clearpage

\section{Guia ràpida de fitxers}

Aquesta taula resumeix cada fitxer inclòs a l'annex: què fa, quina part del projecte afecta i com treballa per dins. Les descripcions surten del docstring del mòdul i d'una lectura estàtica de funcions, classes i imports.

\subsection{Entrypoints i navegacio}

\begingroup
\scriptsize
\setlength{\tabcolsep}{3pt}
\sloppy
\begin{longtable}{>{\raggedright\arraybackslash}p{0.20\linewidth}>{\raggedright\arraybackslash}p{0.29\linewidth}>{\raggedright\arraybackslash}p{0.20\linewidth}>{\raggedright\arraybackslash}p{0.23\linewidth}}
\toprule
Fitxer & Què fa & A què afecta & Com funciona \\
\midrule
\endfirsthead
\toprule
Fitxer & Què fa & A què afecta & Com funciona \\
\midrule
\endhead
\path{Inici.py} & Pàgina inicial de BEMS-RL STUDIO. & Entrypoints i navegacio. & organitza funcions com get\_energyplus\_logo\_data\_uri, detect\_energyplus\_status, render\_hero…; connecta amb sidebar\_nav. \\
\path{__init__.py} & Metadades de l'aplicació per a l'espai de treball Streamlit de BEMS-RL Studio. & Entrypoints i navegacio. & agrupa constants o inicialització de mòdul. \\
\path{sidebar_nav.py} & Navegació lateral compartida de BEMS-RL STUDIO. & Entrypoints i navegacio. & organitza funcions com configure\_studio\_page, render\_studio\_sidebar. \\
\bottomrule
\end{longtable}
\endgroup
\normalsize

\subsection{Pagines Streamlit}

\begingroup
\scriptsize
\setlength{\tabcolsep}{3pt}
\sloppy
\begin{longtable}{>{\raggedright\arraybackslash}p{0.20\linewidth}>{\raggedright\arraybackslash}p{0.29\linewidth}>{\raggedright\arraybackslash}p{0.20\linewidth}>{\raggedright\arraybackslash}p{0.23\linewidth}}
\toprule
Fitxer & Què fa & A què afecta & Com funciona \\
\midrule
\endfirsthead
\toprule
Fitxer & Què fa & A què afecta & Com funciona \\
\midrule
\endhead
\path{pages/Afegir_Entorn.py} & Sinergym Studio: vista per afegir entorns. & Creació i registre d'entorns. & organitza funcions com render\_add\_environment\_hero, render\_weather\_source\_selector, format\_current\_schedule\_range…; connecta amb backend, ui\_fragments, sidebar\_nav. \\
\path{pages/Ajuda.py} & Mòdul funcional pages/Ajuda.py. & Pagines Streamlit. & organitza funcions com inject\_help\_styles, render\_resources\_card, render\_support\_card…; connecta amb ui\_fragments, sidebar\_nav. \\
\path{pages/Avaluar_Agent.py} & Pàgina per avaluar un agent SB3 entrenat en un entorn Sinergym. & Avaluació de models SB3. & organitza funcions com render\_metadata\_grid, render\_evaluation\_runtime, render\_evaluation\_runtime\_frame; connecta amb avaluar\_agent\_backend, ui\_fragments, sidebar\_nav. \\
\path{pages/Entrenar_Agent.py} & Pàgina d'entrenament de l'agent per a l'aplicació BEMS-RL-STUDIO. & Flux d'entrenament i estat de Streamlit. & organitza funcions com train\_tab; connecta amb entrenar\_agent\_charts, entrenar\_agent\_constants, entrenar\_agent\_env…. \\
\path{pages/Gestionar_Arxius.py} & Mòdul funcional pages/Gestionar\_Arxius.py. & Gestió de fitxers del projecte. & organitza funcions com render\_weather\_view\_selector, build\_selection\_archive, build\_download\_name…; connecta amb gestionar\_arxius\_backend, mapa\_backend, ui\_fragments…. \\
\path{pages/Interaccionar_amb_l'Agent.py} & Pàgina de control en viu d'un entorn Sinergym. & Control en viu i prova manual de polítiques. & organitza funcions com render\_interaction\_hero, render\_interaction\_section, init\_session…; connecta amb interaccionar\_agent\_backend, ui\_fragments, sidebar\_nav. \\
\path{pages/Mostrar_Entorn.py} & Pàgina de visualització d'entorns Sinergym. & Visualització d'entorns i geometria 3D. & organitza funcions com render\_environment\_hero, render\_environment\_section, describe\_environment\_tags…; connecta amb common, mapa\_backend, mostrar\_entorn\_backend…. \\
\path{pages/Resultats.py} & Streamlit frontend per a la pàgina de resultats de BEMS-RL STUDIO. & Selecció de runs, dashboard i descàrregues. & organitza funcions com build\_download\_filename, render\_results\_page; connecta amb resultats\_backend, resultats\_dashboard, ui\_fragments…. \\
\path{pages/Simular_Entorn.py} & Mòdul funcional pages/Simular\_Entorn.py. & Baseline sense agent i comparació posterior. & organitza funcions com empty\_baseline\_runtime, ensure\_baseline\_runtime, reset\_baseline\_runtime…; connecta amb common, simular\_entorn\_backend, ui\_fragments…. \\
\path{pages/Visor_Climes_EPW.py} & Mòdul funcional pages/Visor\_Climes\_EPW.py. & Inspecció de fitxers EPW i clima. & Defineix EpwDashboardState; organitza funcions com format\_number, data\_frame\_has\_plot\_values, wind\_rose\_tables\_have\_plot\_values…; connecta amb backend, epw\_figures, mapa\_backend…. \\
\path{pages/__init__.py} & Streamlit Marcador de paquet de pàgina per a BEMS-RL Studio. & Pagines Streamlit. & agrupa constants o inicialització de mòdul. \\
\bottomrule
\end{longtable}
\endgroup
\normalsize

\subsection{Components reutilitzables}

\begingroup
\scriptsize
\setlength{\tabcolsep}{3pt}
\sloppy
\begin{longtable}{>{\raggedright\arraybackslash}p{0.20\linewidth}>{\raggedright\arraybackslash}p{0.29\linewidth}>{\raggedright\arraybackslash}p{0.20\linewidth}>{\raggedright\arraybackslash}p{0.23\linewidth}}
\toprule
Fitxer & Què fa & A què afecta & Com funciona \\
\midrule
\endfirsthead
\toprule
Fitxer & Què fa & A què afecta & Com funciona \\
\midrule
\endhead
\path{page_components/__init__.py} & Components Streamlit reutilitzables per a pàgines BEMS-RL Studio. & Components reutilitzables. & agrupa constants o inicialització de mòdul. \\
\path{page_components/resultats_dashboard.py} & Component del panell en línia per a la pàgina de resultats. & Dashboard incrustat de resultats. & Defineix DashboardTabSpec; organitza funcions com render\_inline\_dashboard; connecta amb grafics.battery, grafics.comfort, grafics.control…. \\
\path{page_components/training_agent.py} & Controls d'hiperparàmetres de l'agent per a la pàgina d'entrenament. & Formularis compartits d'entrenament. & organitza funcions com render\_agent\_params\_section; connecta amb entrenar\_agent\_session. \\
\path{page_components/training_rewards.py} & Controls de configuració de recompenses per a la pàgina de entrenament. & Formularis compartits d'entrenament. & organitza funcions com render\_reward\_kwargs\_section; connecta amb entrenar\_agent\_constants, entrenar\_agent\_session. \\
\path{page_components/training_shared.py} & Constants de pàgina d'entrenament compartides i components de capçalera reutilitzables. & Formularis compartits d'entrenament. & organitza funcions com render\_training\_hero, render\_training\_section, format\_saved\_training\_label…; connecta amb entrenar\_agent\_artifacts, entrenar\_agent\_constants, entrenar\_agent\_session. \\
\path{page_components/training_wrappers.py} & Controls de wrappers i registre per a la pàgina d'entrenament. & Formularis compartits d'entrenament. & organitza funcions com render\_observation\_wrappers, render\_action\_wrappers, render\_energy\_logging\_wrappers…; connecta amb entrenar\_agent\_constants, entrenar\_agent\_env, entrenar\_agent\_session. \\
\path{page_components/ui_fragments.py} & Fragments UI compartits per pàgines Streamlit de BEMS-RL Studio. & Components reutilitzables. & organitza funcions com render\_hero, render\_section\_title, render\_section\_card…. \\
\bottomrule
\end{longtable}
\endgroup
\normalsize

\subsection{Backend funcional}

\begingroup
\scriptsize
\setlength{\tabcolsep}{3pt}
\sloppy
\begin{longtable}{>{\raggedright\arraybackslash}p{0.20\linewidth}>{\raggedright\arraybackslash}p{0.29\linewidth}>{\raggedright\arraybackslash}p{0.20\linewidth}>{\raggedright\arraybackslash}p{0.23\linewidth}}
\toprule
Fitxer & Què fa & A què afecta & Com funciona \\
\midrule
\endfirsthead
\toprule
Fitxer & Què fa & A què afecta & Com funciona \\
\midrule
\endhead
\path{backend/__init__.py} & Lògica d'aplicació de BEMS-RL Studio segura per importar. & Backend funcional. & agrupa constants o inicialització de mòdul. \\
\path{backend/afegir_entorn_analysis.py} & Lectura i anàlisi d'edificis per preparar la configuració dels entorns. & Anàlisi, conversió i YAML d'entorns. & Defineix ProbeReward; organitza funcions com load\_building\_json, extract\_schedule\_type\_index, resolve\_schedule\_name…; connecta amb afegir\_entorn\_common. \\
\path{backend/afegir_entorn_backend.py} & Façana de compatibilitat per al backend afegir entorns. & Anàlisi, conversió i YAML d'entorns. & connecta amb afegir\_entorn\_analysis, afegir\_entorn\_common, afegir\_entorn\_config…. \\
\path{backend/afegir_entorn_common.py} & Constants, tipus i helpers compartits pel flux d'afegir entorns. & Anàlisi, conversió i YAML d'entorns. & Defineix UpgradeResult, EnvironmentValidationResult, ActuatorOption…; organitza funcions com save\_uploaded\_bytes, sanitize\_identifier, normalize\_token…. \\
\path{backend/afegir_entorn_config.py} & Crea i valideu les configuracions de Sinergym des de l'entrada del formulari Studio. & Anàlisi, conversió i YAML d'entorns. & organitza funcions com build\_variable\_output\_names, add\_variable\_spec, build\_training\_observation\_config…; connecta amb afegir\_entorn\_analysis, afegir\_entorn\_common. \\
\path{backend/afegir_entorn_conversion.py} & Conversió i actualització d'IDF/epJSON dins del flux per afegir entorns. & Anàlisi, conversió i YAML d'entorns. & organitza funcions com detect\_idf\_version, needs\_idf\_upgrade, get\_transition\_updater\_dir…; connecta amb afegir\_entorn\_common. \\
\path{backend/avaluar_agent_backend.py} & Temps d'execució d'avaluació asíncrona per a agents Studio entrenats. & Execució de models entrenats sobre entorns. & Defineix EvaluationResult; organitza funcions com push\_evaluation\_progress, build\_evaluation\_env, empty\_evaluation\_runtime…; connecta amb common, model\_metadata, sb3\_utils. \\
\path{backend/common.py} & Utilitats compartides de projecte per al backend de BEMS-RL Studio. & Backend funcional. & organitza funcions com registered\_env\_ids, list\_registered\_env\_ids, is\_registered\_env\_id…. \\
\path{backend/entrenar_agent_artifacts.py} & Rutes d'artefacte d'entrenament, metadades i postprocessament CSV. & Preparació, execució i persistència d'entrenaments. & organitza funcions com sanitize\_training\_name\_component, get\_training\_artifacts\_root, build\_training\_artifact\_paths…; connecta amb entrenar\_agent\_constants. \\
\path{backend/entrenar_agent_charts.py} & Gràfics d'entrenament en directe per a la pàgina d'entrenament de l'agent. & Preparació, execució i persistència d'entrenaments. & organitza funcions com render\_live\_training\_charts; connecta amb entrenar\_agent\_session, grafics.data\_loader, grafics.episode. \\
\path{backend/entrenar_agent_constants.py} & Constants i registres per al backend de l'agent de entrenament. & Preparació, execució i persistència d'entrenaments. & Defineix RewardClassRegistry; connecta amb common. \\
\path{backend/entrenar_agent_env.py} & Descobriment d'entorns i metadades per configurar entrenaments. & Preparació, execució i persistència d'entrenaments. & organitza funcions com list\_registered\_envs, list\_files, with\_detailed\_hvac\_meters…; connecta amb common, entrenar\_agent\_constants. \\
\path{backend/entrenar_agent_rewards.py} & Recompenses, algorismes i creadors de càrrega útil de entrenament completa. & Preparació, execució i persistència d'entrenaments. & organitza funcions com build\_multizone\_temp\_mapping, build\_multizone\_occupancy\_mapping, build\_reward\_kwargs…; connecta amb common, entrenar\_agent\_artifacts, entrenar\_agent\_constants…. \\
\path{backend/entrenar_agent_runtime.py} & Cicle de vida d'execució per a sessions d'entrenament en directe. & Preparació, execució i persistència d'entrenaments. & organitza funcions com clear\_training\_runtime, learn\_training\_chunk, save\_training\_artifacts…; connecta amb entrenar\_agent\_artifacts, entrenar\_agent\_constants, entrenar\_agent\_env…. \\
\path{backend/entrenar_agent_session.py} & Gestió de l'estat de sessió per a la pàgina d'entrenament de l'agent. & Preparació, execució i persistència d'entrenaments. & organitza funcions com training\_form\_key, reset\_widget\_state\_if\_disabled, ensure\_select\_state…; connecta amb entrenar\_agent\_constants, entrenar\_agent\_runtime. \\
\path{backend/entrenar_agent_wrappers.py} & Configuració i aplicació de wrapper per a entorns d'entrenament. & Preparació, execució i persistència d'entrenaments. & Defineix EnergyCostFileLogger; organitza funcions com wrapper\_details, build\_wrapper\_summary\_rows, apply\_training\_wrappers…; connecta amb common, entrenar\_agent\_constants. \\
\path{backend/epw_figures.py} & Plotly creadors de figures per al visor de clima EPW. & Lectura i gràfics de clima EPW. & organitza funcions com ui\_text, build\_active\_month\_axis, apply\_figure\_style…; connecta amb visor\_epw\_backend. \\
\path{backend/gestionar_arxius_backend.py} & Gestió de fitxers per a actius, models i dades meteorològiques de Studio. & Exploració, pujada, descàrrega i eliminació de fitxers. & Defineix ExplorerItem; organitza funcions com get\_size\_format, parse\_epw\_overview, load\_weather\_map\_rows…. \\
\path{backend/interaccionar_agent_backend.py} & Execució en directe per provar agents entrenats dins d'un entorn Sinergym. & Execució de models entrenats sobre entorns. & organitza funcions com mode\_requires\_model, load\_model\_training\_metadata, build\_env\_factory…; connecta amb common, model\_metadata, sb3\_utils. \\
\path{backend/mapa_backend.py} & Utilitats de representació de mapes compartides construïdes a pydeck i Streamlit. & Backend funcional. & organitza funcions com render\_location\_map, render\_multipoint\_map. \\
\path{backend/model_metadata.py} & Utilitats compartides per reconstruir models entrenats de l'Studio. & Backend funcional. & organitza funcions com load\_model\_metadata, build\_gym\_kwargs\_from\_metadata, resolve\_reward\_class…; connecta amb sb3\_utils. \\
\path{backend/mostrar_entorn_backend.py} & Backend de metadades de l'entorn: carrega actius, PV, emmagatzematge i dades per pintar la UI. & Lectura d'actius i representació 3D. & organitza funcions com load\_epjson, parse\_pv\_and\_storage, load\_environment\_assets…; connecta amb viewer\_3d\_backend. \\
\path{backend/resultats_backend.py} & Càrrega de dades, preparació de KPI i descobriment d'execució per a la pàgina de resultats. & Càrrega de runs, mètriques i exportació d'informes. & Defineix RunOption, ResultsPageState, RunArtifacts…; organitza funcions com build\_results\_page\_state, get\_run\_artifacts, load\_action\_data…; connecta amb grafics.time\_utils. \\
\path{backend/resultats_report_backend.py} & Generació d'informes HTML i PDF per a les execucions de Studio completades. & Càrrega de runs, mètriques i exportació d'informes. & Defineix ReportFigure; organitza funcions com generate\_report\_bytes; connecta amb grafics.battery, grafics.comfort, grafics.comfort\_scope…. \\
\path{backend/sb3_utils.py} & Descoberta, càrrega i format d'accions per a models Stable-Baselines3. & Backend funcional. & organitza funcions com scan\_model\_zips, candidate\_vecnorm, load\_sb3\_model\_bytes…; connecta amb common. \\
\path{backend/simular_entorn_backend.py} & Simulacions de referència per a execucions no controlades o basades en regles. & Simulació baseline i artefactes comparables. & Defineix EnvironmentSummary, BaselineSimulationResult, ScheduleBaselineReward; organitza funcions com list\_environment\_ids, describe\_environment, run\_baseline\_simulation; connecta amb common, entrenar\_agent\_co… \\
\path{backend/training_scripts.py} & Scripts d'ajuda reproduïbles per a execucions BEMS-RL desades. & Backend funcional. & organitza funcions com build\_training\_eval\_script, build\_training\_repro\_script. \\
\path{backend/viewer_3d_backend.py} & Backend del visor d'edificis 3D: processament de geometria i generació de figures Plotly. & Lectura d'actius i representació 3D. & organitza funcions com extract\_vertices\_general, build\_surface\_records, plot\_3d\_model. \\
\path{backend/visor_epw_backend.py} & Utilitats de preparació de dades per al visor climàtic EPW. & Lectura i gràfics de clima EPW. & organitza funcions com variable\_label, list\_epw\_catalog, load\_epw\_bundle…. \\
\path{backend/weather_profiles.py} & Anàlisi de perfils meteorològics i vista prèvia per crear entorns. & Backend funcional. & Defineix WeatherProfileSuggestion; organitza funcions com summarize\_weather\_file, load\_epw\_dry\_bulb\_series, build\_weather\_temperature\_preview…. \\
\bottomrule
\end{longtable}
\endgroup
\normalsize

\subsection{Grafics i metriques}

\begingroup
\scriptsize
\setlength{\tabcolsep}{3pt}
\sloppy
\begin{longtable}{>{\raggedright\arraybackslash}p{0.20\linewidth}>{\raggedright\arraybackslash}p{0.29\linewidth}>{\raggedright\arraybackslash}p{0.20\linewidth}>{\raggedright\arraybackslash}p{0.23\linewidth}}
\toprule
Fitxer & Què fa & A què afecta & Com funciona \\
\midrule
\endfirsthead
\toprule
Fitxer & Què fa & A què afecta & Com funciona \\
\midrule
\endhead
\path{backend/grafics/__init__.py} & Gràfics Plotly i mètriques compartides per a resultats i informes. & KPIs, normalització temporal i figures Plotly. & agrupa constants o inicialització de mòdul. \\
\path{backend/grafics/battery.py} & Dades del quadre de comandament de la bateria, la xarxa i les tarifes. & KPIs, normalització temporal i figures Plotly. & organitza funcions com make\_battery\_power\_plot, make\_battery\_soc\_plot, make\_energy\_price\_plot…; connecta amb grafics.column\_utils, grafics.energy\_units, grafics.plot\_helpers…. \\
\path{backend/grafics/column_utils.py} & Detecció de columnes disponibles per construir les figures del panell. & KPIs, normalització temporal i figures Plotly. & agrupa constants o inicialització de mòdul. \\
\path{backend/grafics/comfort.py} & Compliment de la comoditat i xifres d'incompliment de la temperatura. & KPIs, normalització temporal i figures Plotly. & organitza funcions com make\_comfort\_compliance, make\_violation\_bars, make\_violation\_single\_line; connecta amb grafics.column\_utils, grafics.comfort\_scope, grafics.style…. \\
\path{backend/grafics/comfort_scope.py} & Filtres d'abast de confort compartits per KPI, gràfics i informes. & KPIs, normalització temporal i figures Plotly. & organitza funcions com derive\_occupied\_mask, has\_occupied\_data, filter\_by\_comfort\_scope…. \\
\path{backend/grafics/control.py} & Figures de control i acció per a HVAC i terra radiant. & KPIs, normalització temporal i figures Plotly. & organitza funcions com make\_radiant\_control\_plot, make\_agent\_actions\_plot; connecta amb grafics.style, grafics.time\_utils. \\
\path{backend/grafics/data_loader.py} & Carrega i normalitza els fitxers CSV de resultats per a panells i informes. & KPIs, normalització temporal i figures Plotly. & organitza funcions com load\_data; connecta amb grafics.observation\_columns, grafics.time\_utils. \\
\path{backend/grafics/energy_units.py} & Conversions d'energia, preu i unitats de bateria per als gràfics. & KPIs, normalització temporal i figures Plotly. & connecta amb grafics.time\_utils. \\
\path{backend/grafics/episode.py} & Xifres del quadre de comandament a nivell d'episodi. & KPIs, normalització temporal i figures Plotly. & organitza funcions com make\_episode\_metrics; connecta amb grafics.style. \\
\path{backend/grafics/figures_zones.py} & Detecció i filtratge de zones per als panells de resultats. & KPIs, normalització temporal i figures Plotly. & organitza funcions com detect\_zones, filter\_obs\_by\_zone, get\_zone\_options. \\
\path{backend/grafics/heatmaps.py} & Xifres de mapes de calor per als resums del panell global i de zona. & KPIs, normalització temporal i figures Plotly. & organitza funcions com make\_heatmap, make\_violation\_heatmap\_percent, compute\_zone\_temperature\_pivots…; connecta amb grafics.column\_utils, grafics.comfort, grafics.time\_utils. \\
\path{backend/grafics/hvac.py} & HVAC xifres de consum i avaria del comptador. & KPIs, normalització temporal i figures Plotly. & organitza funcions com make\_hvac\_consumption\_plot, make\_hvac\_meter\_breakdown\_plot; connecta amb grafics.plot\_helpers, grafics.style, grafics.time\_utils. \\
\path{backend/grafics/kpis.py} & Càlculs de KPI per als resums de resultats de BEMS-RL Studio. & KPIs, normalització temporal i figures Plotly. & organitza funcions com compute\_kpis; connecta amb grafics.comfort\_scope, grafics.observation\_columns. \\
\path{backend/grafics/metrics.py} & Agrega les dades d'observació i progrés en mètriques preparades per a gràfics. & KPIs, normalització temporal i figures Plotly. & organitza funcions com compute\_metrics; connecta amb grafics.observation\_columns, grafics.time\_utils. \\
\path{backend/grafics/observation_columns.py} & Utilitats de normalització de columnes per a fitxers d'observació Sinergym. & KPIs, normalització temporal i figures Plotly. & organitza funcions com find\_first\_available\_column, normalize\_observation\_columns, find\_hvac\_consumption\_source…. \\
\path{backend/grafics/plot_helpers.py} & Helpers Plotly comuns per als gràfics agregats del dashboard. & KPIs, normalització temporal i figures Plotly. & organitza funcions com build\_mode\_scatter\_trace, add\_mode\_scatter\_trace, energy\_series\_for\_mode…; connecta amb grafics.time\_utils. \\
\path{backend/grafics/style.py} & Estil Plotly compartit per a les figures del panell BEMS-RL Studio. & KPIs, normalització temporal i figures Plotly. & organitza funcions com pick\_semantic\_trace\_color, style\_figure\_semantics. \\
\path{backend/grafics/thermal.py} & Dades d'entrada de confort tèrmic: temperatura, humitat i consignes. & KPIs, normalització temporal i figures Plotly. & organitza funcions com make\_indoor\_temperature\_plot, make\_indoor\_humidity\_plot, make\_setpoints\_plot…; connecta amb grafics.column\_utils, grafics.plot\_helpers, grafics.style…. \\
\path{backend/grafics/time_utils.py} & Filtrat temporal i ajuda d'eix per a figures del panell. & KPIs, normalització temporal i figures Plotly. & organitza funcions com sanitize\_observation\_time\_columns, filter\_by\_season; connecta amb grafics.style. \\
\bottomrule
\end{longtable}
\endgroup
\normalsize

\clearpage

\section{Guia de classes}
Aquesta secció explica totes les classes i dataclasses del codi funcional inclòs a l'annex. Serveix com a mapa ràpid abans d'entrar al codi complet.
\begingroup
\scriptsize
\setlength{\tabcolsep}{3pt}
\sloppy
\begin{longtable}{>{\raggedright\arraybackslash}p{0.18\linewidth}>{\raggedright\arraybackslash}p{0.13\linewidth}>{\raggedright\arraybackslash}p{0.25\linewidth}>{\raggedright\arraybackslash}p{0.18\linewidth}>{\raggedright\arraybackslash}p{0.18\linewidth}}
\toprule
Classe & Tipus & Què fa & Fitxer & Com funciona \\
\midrule
\endfirsthead
\toprule
Classe & Tipus & Què fa & Fitxer & Com funciona \\
\midrule
\endhead
\texttt{ProbeReward} & reward Sinergym & Reward nul usat només per inspeccionar handlers disponibles. & \path{backend/afegir_entorn_analysis.py} & Hereta de BaseReward. \\
\texttt{UpgradeResult} & dataclass & Resultat d'un intent d'actualitzar una IDF a la versio suportada. & \path{backend/afegir_entorn_common.py} & guarda dades tipades sense lògica complexa. \\
\texttt{EnvironmentValidationResult} & dataclass & Espais i observacio inicial retornats en validar un entorn registrat. & \path{backend/afegir_entorn_common.py} & guarda dades tipades sense lògica complexa. \\
\texttt{ActuatorOption} & dataclass & Opcio d'actuador candidata detectada en un model EnergyPlus. & \path{backend/afegir_entorn_common.py} & guarda dades tipades sense lògica complexa. \\
\texttt{BuildingTrainingAnalysis} & dataclass & Resultat agregat de l'analisi d'un edifici per entrenament RL. & \path{backend/afegir_entorn_common.py} & guarda dades tipades sense lògica complexa. \\
\texttt{EvaluationResult} & dataclass & Resum de resultats d'una execució d'un treballador d'avaluació. & \path{backend/avaluar_agent_backend.py} & guarda dades tipades sense lògica complexa. \\
\texttt{RewardClassRegistry} & classe & Registre lazy de classes de recompensa exposades per Sinergym. & \path{backend/entrenar_agent_constants.py} & guarda dades tipades sense lògica complexa. \\
\texttt{EnergyCostFileLogger} & wrapper Gym & Logger addicional (opcional) que escriu cost per pas en un CSV independent. & \path{backend/entrenar_agent_wrappers.py} & Hereta de Wrapper; mètodes clau: step, reset, close. \\
\texttt{ExplorerItem} & dataclass & Fila del navegador de fitxers serialitzable que es mostra a la pàgina de gestió d'actius. & \path{backend/gestionar_arxius_backend.py} & guarda dades tipades sense lògica complexa. \\
\texttt{RunOption} & dataclass & Directori d'execució seleccionable que es mostra a la pàgina de resultats. & \path{backend/resultats_backend.py} & guarda dades tipades sense lògica complexa. \\
\texttt{ResultsPageState} & dataclass & S'han resolt les entrades necessàries abans que es pugui representar la pàgina de resultats. & \path{backend/resultats_backend.py} & guarda dades tipades sense lògica complexa. \\
\texttt{RunArtifacts} & dataclass & Fitxers d'artefactes derivats d'una execució seleccionada. & \path{backend/resultats_backend.py} & guarda dades tipades sense lògica complexa. \\
\texttt{DashboardData} & dataclass & Es requereix un paquet de dades carregat per representar el panell de resultats en línia. & \path{backend/resultats_backend.py} & guarda dades tipades sense lògica complexa. \\
\texttt{ReportFigure} & dataclass & Plotly gràfic més la secció d'informe on hauria d'aparèixer. & \path{backend/resultats_report_backend.py} & guarda dades tipades sense lògica complexa. \\
\texttt{EnvironmentSummary} & dataclass & Descripció compacta d'un entorn Sinergym registrat. & \path{backend/simular_entorn_backend.py} & guarda dades tipades sense lògica complexa. \\
\texttt{BaselineSimulationResult} & dataclass & Sortides persistents i metadades d'estat per a una simulació de referència. & \path{backend/simular_entorn_backend.py} & guarda dades tipades sense lògica complexa. \\
\texttt{ScheduleBaselineReward} & reward Sinergym & Reward zero que conserva les mètriques físiques per comparar contra agents entrenats. & \path{backend/simular_entorn_backend.py} & Hereta de BaseReward; mètodes clau: \_compute\_power\_demand, \_compute\_temperature\_violation. \\
\texttt{WeatherProfileSuggestion} & dataclass & Descriu un suggeriment de fitxer meteorològic per a la creació de l'entorn UI. & \path{backend/weather_profiles.py} & guarda dades tipades sense lògica complexa. \\
\texttt{DashboardTabSpec} & classe & Declaració d'una pestanya del dashboard i les figures que conté. & \path{page_components/resultats_dashboard.py} & Hereta de NamedTuple. \\
\texttt{EpwDashboardState} & dataclass & Dades preparades necessàries per representar les pestanyes del panell EPW. & \path{pages/Visor_Climes_EPW.py} & guarda dades tipades sense lògica complexa. \\
\bottomrule
\end{longtable}
\endgroup
\normalsize

\clearpage


\section{Visió general del codi}

BEMS-RL Studio està organitzat com una capa d'aplicació sobre Sinergym. Les
pàgines Streamlit gestionen la interacció amb l'usuari, els components
reutilitzables concentren formularis i panells compartits, i el backend manté
la lògica de fitxers, validació, entrenament, simulació, avaluació, gràfics i
informes. Aquesta separació és la que permet documentar el codi amb autodoc i,
alhora, mantenir les pàgines com a punts d'entrada lleugers.

\subsection{Arquitectura d'execució}

\begin{figure}[H]
\centering
\resizebox{\linewidth}{!}{%
\begin{tikzpicture}[node distance=1.4cm and 1.4cm]
  \node[studioBox] (user) {Usuari};
  \node[studioBlue, right=of user] (pages) {Pàgines\\Streamlit};
  \node[studioBox, right=of pages] (components) {Components\\reutilitzables};
  \node[studioWide, right=of components] (backend) {Backends\\d'aplicació};
  \node[studioBox, above right=0.6cm and 1.2cm of backend] (charts) {Capa de gràfics\\\texttt{backend.grafics}};
  \node[studioBox, below right=0.6cm and 1.2cm of backend] (sinergym) {Sinergym\\Gymnasium + EnergyPlus};
  \node[studioBox, right=2.1cm of charts] (reports) {Informes\\HTML/PDF};
  \node[studioBox, right=2.1cm of sinergym] (artifacts) {Artefactes\\models, logs, CSV};
  \draw[studioEdge] (user) -- (pages);
  \draw[studioEdge] (pages) -- (components);
  \draw[studioEdge] (components) -- (backend);
  \draw[studioEdge] (pages) to[bend left=20] (backend);
  \draw[studioEdge] (backend) -- (charts);
  \draw[studioEdge] (backend) -- (sinergym);
  \draw[studioEdge] (charts) -- (reports);
  \draw[studioEdge] (sinergym) -- (artifacts);
  \draw[studioEdge] (artifacts) -- (charts);
\end{tikzpicture}}
\caption{Arquitectura lògica de BEMS-RL Studio i relació amb Sinergym.}
\end{figure}
\diagramnote{El diagrama separa l'aplicació en capes. L'usuari només toca les pàgines Streamlit; aquestes pàgines deleguen en components i backends. El backend és qui parla amb Sinergym, llegeix o escriu artefactes i passa dades a la capa de gràfics i informes.}

\subsection{Topologia Docker i documentació}

\begin{figure}[H]
\centering
\resizebox{0.92\linewidth}{!}{%
\begin{tikzpicture}[node distance=1.45cm and 1.3cm]
  \node[studioBox] (browser) {Navegador};
  \node[studioBlue, right=of browser] (app) {\texttt{bemsrlstudio}\\Streamlit\\:8501};
  \node[studioBox, above right=0.65cm and 1.35cm of app] (sinDocs) {\texttt{docs-sinergym}\\Sphinx\\:8000};
  \node[studioBox, below right=0.65cm and 1.35cm of app] (studioDocs) {\texttt{docs-studio}\\Sphinx\\:8001};
  \node[studioWide, right=2.0cm of app] (source) {Arbre de codi\\compartit};
  \node[studioWide, right=of source] (build) {Sortida Sphinx\\\texttt{docs/build}};
  \draw[studioEdge] (browser) -- node[above,font=\scriptsize] {aplicació} (app);
  \draw[studioEdge] (browser) to[bend left=24] node[above,font=\scriptsize] {docs Sinergym} (sinDocs);
  \draw[studioEdge] (browser) to[bend right=24] node[below,font=\scriptsize] {docs Studio} (studioDocs);
  \draw[studioEdge] (source) -- (app);
  \draw[studioEdge] (source) -- (sinDocs);
  \draw[studioEdge] (source) -- (studioDocs);
  \draw[studioEdge] (sinDocs) -- (build);
  \draw[studioEdge] (studioDocs) -- (build);
\end{tikzpicture}}
\caption{Serveis locals: aplicació i documentació separada en dos ports.}
\end{figure}
\diagramnote{Aquest dibuix mostra la topologia d'execució local. Streamlit serveix la interfície principal, mentre que la documentació queda separada en serveis propis. Tots tres serveis llegeixen el mateix arbre de codi, però generen sortides diferents.}

\subsection{Propietat de dades i artefactes}

\begin{figure}[H]
\centering
\resizebox{\linewidth}{!}{%
\begin{tikzpicture}[node distance=1.0cm and 1.05cm]
  \node[studioBox] (idf) {IDF / epJSON};
  \node[studioBox, below=of idf] (epw) {EPW};
  \node[studioWide, right=of idf] (env) {\texttt{afegir\_entorn\_*}\\configuració i validació};
  \node[studioBox, right=of env] (yaml) {YAML\\Sinergym};
  \node[studioWide, right=of yaml] (train) {\texttt{entrenar\_agent\_*}\\workspace reproduïble};
  \node[studioBox, above right=0.35cm and 1.0cm of train] (model) {Model\\SB3};
  \node[studioBox, below right=0.35cm and 1.0cm of train] (logs) {CSV\\progress + obs.};
  \node[studioWide, right=1.15cm of logs] (results) {\texttt{resultats\_backend}\\KPIs i dades};
  \node[studioWide, above=of results] (eval) {Avaluació\\i control en viu};
  \node[studioBox, right=of results] (report) {Informe\\HTML/PDF};
  \draw[studioEdge] (idf) -- (env);
  \draw[studioEdge] (epw) -- (env);
  \draw[studioEdge] (env) -- (yaml);
  \draw[studioEdge] (yaml) -- (train);
  \draw[studioEdge] (train) -- (model);
  \draw[studioEdge] (train) -- (logs);
  \draw[studioEdge] (model) -- (eval);
  \draw[studioEdge] (logs) -- (results);
  \draw[studioEdge] (results) -- (report);
\end{tikzpicture}}
\caption{Mòduls responsables de cada família d'artefactes.}
\end{figure}
\diagramnote{Aquí es veu el recorregut dels fitxers físics. Els IDF/epJSON i EPW entren pel flux de creació d'entorns, es converteixen en YAML de Sinergym i després alimenten entrenaments. Dels entrenaments surten models i CSV, que són la base de resultats i informes.}

\subsection{Límit d'importació i autodoc}

\begin{figure}[H]
\centering
\resizebox{0.82\linewidth}{!}{%
\begin{tikzpicture}[node distance=1.25cm and 1.45cm]
  \node[studioAmber] (pages) {\texttt{pages/*.py}\\UI executada en importar};
  \node[studioBox, right=of pages] (components) {\texttt{page\_components}\\components reutilitzables de pantalla};
  \node[studioBox, below=of components] (backend) {\texttt{backend}\\lògica importable};
  \node[studioBox, above=of components] (styles) {\texttt{page\_styles}\\CSS importable};
  \node[studioBlue, right=2.0cm of components] (autodoc) {Sphinx\\autodoc};
  \node[studioBox, below=of autodoc] (staticref) {Referència\\AST estàtica};
  \draw[studioEdge] (pages) -- (components);
  \draw[studioEdge] (pages) -- (backend);
  \draw[studioEdge] (pages) -- (styles);
  \draw[studioEdge] (autodoc) -- (components);
  \draw[studioEdge] (autodoc) -- (backend);
  \draw[studioEdge] (autodoc) -- (styles);
  \draw[studioDashed] (staticref) -- node[below,font=\scriptsize] {sense executar UI} (pages);
\end{tikzpicture}}
\caption{Frontera entre lògica importable i punts d'entrada de Streamlit.}
\end{figure}
\diagramnote{La frontera important és que les pàgines Streamlit s'executen en importar-les, mentre que backend i components han de ser importables sense efectes laterals pesats. Per això la documentació automàtica tracta les pàgines amb més prudència.}

\section{Fluxos principals}

\subsection{Creació d'entorns}

\begin{figure}[H]
\centering
\resizebox{\linewidth}{!}{%
\begin{tikzpicture}[node distance=1.0cm and 0.95cm]
  \node[studioBox] (upload) {Pujar fitxers\\edifici + clima};
  \node[studioBox, right=of upload] (inspect) {Inspeccionar\\variables i actuadors};
  \node[studioBox, right=of inspect] (configure) {Construir\\YAML};
  \node[studioBox, right=of configure] (validate) {Validar\\Gymnasium};
  \node[studioBox, right=of validate] (register) {Registrar\\entorn};
  \node[studioBox, right=of register] (use) {Entrenar\\simular\\avaluar};
  \draw[studioEdge] (upload) -- (inspect);
  \draw[studioEdge] (inspect) -- (configure);
  \draw[studioEdge] (configure) -- (validate);
  \draw[studioEdge] (validate) -- (register);
  \draw[studioEdge] (register) -- (use);
\end{tikzpicture}}
\caption{Flux d'autoria d'entorns de la documentació de Studio.}
\end{figure}
\diagramnote{El flux va d'esquerra a dreta: primer es carreguen fitxers, després s'inspeccionen variables i actuadors, es construeix el YAML, es valida amb Gymnasium i finalment l'entorn queda llest per entrenar, simular o avaluar.}

\subsection{Cicle d'entrenament}

\begin{figure}[H]
\centering
\resizebox{0.92\linewidth}{!}{%
\begin{tikzpicture}[node distance=1.1cm and 1.15cm]
  \node[studioBox] (form) {Formulari\\d'entrenament};
  \node[studioBox, right=of form] (state) {Estat UI\\validat};
  \node[studioBox, right=of state] (workspace) {Workspace\\d'entrenament};
  \node[studioBox, above right=0.35cm and 1.15cm of workspace] (scripts) {Scripts\\reproduïbles};
  \node[studioBox, right=of workspace] (model) {Model\\\texttt{.zip}};
  \node[studioBox, below right=0.35cm and 1.15cm of workspace] (logs) {Logs i\\observacions};
  \draw[studioEdge] (form) -- (state);
  \draw[studioEdge] (state) -- (workspace);
  \draw[studioEdge] (workspace) -- (scripts);
  \draw[studioEdge] (workspace) -- (model);
  \draw[studioEdge] (workspace) -- (logs);
\end{tikzpicture}}
\caption{Cicle de vida dels artefactes d'entrenament.}
\end{figure}
\diagramnote{El formulari no entrena directament. Primer es converteix en estat validat, després es crea un workspace reproduïble i, a partir d'aquí, es guarden model, logs, observacions i scripts que permeten repetir l'experiment.}

\subsection{Resultats i selecció de gràfics}

\begin{figure}[H]
\centering
\resizebox{\linewidth}{!}{%
\begin{tikzpicture}[node distance=1.05cm and 1.05cm]
  \node[studioBox] (progress) {\texttt{progress.csv}};
  \node[studioBox, below=of progress] (obs) {\texttt{observations.csv}};
  \node[studioWide, right=of progress] (loader) {\texttt{DashboardData}};
  \node[studioBox, right=of loader] (kpis) {KPIs};
  \node[studioWide, below=of kpis] (figures) {Constructors\\Plotly};
  \node[studioBox, right=of kpis] (dashboard) {Dashboard\\Resultats};
  \node[studioBox, right=of figures] (report) {Informe\\HTML/PDF};
  \node[studioBox, below=of figures] (skip) {Ometre secció\\si no hi ha dades};
  \draw[studioEdge] (progress) -- (loader);
  \draw[studioEdge] (obs) -- (loader);
  \draw[studioEdge] (loader) -- (kpis);
  \draw[studioEdge] (loader) -- (figures);
  \draw[studioEdge] (kpis) -- (dashboard);
  \draw[studioEdge] (figures) -- (dashboard);
  \draw[studioEdge] (figures) -- node[above,font=\scriptsize] {traça vàlida} (report);
  \draw[studioDashed] (figures) -- node[right,font=\scriptsize] {sense dades} (skip);
\end{tikzpicture}}
\caption{Ruta de dades de resultats i criteri d'inclusió de figures.}
\end{figure}
\diagramnote{Els resultats sempre parteixen de \texttt{progress.csv} i \texttt{observations.csv}. El carregador els converteix en \texttt{DashboardData}; les figures només passen al dashboard o a l'informe quan tenen dades reals, així s'eviten gràfics buits.}

\subsection{Control en viu}

\begin{figure}[H]
\centering
\resizebox{0.94\linewidth}{!}{%
\begin{tikzpicture}[node distance=1.05cm and 1.05cm]
  \node[studioBox] (model) {Model SB3};
  \node[studioBox, right=of model] (metadata) {Metadades\\del model};
  \node[studioBox, below=of metadata] (env) {Entorn\\seleccionat};
  \node[studioWide, right=of metadata] (validate) {Compatibilitat\\d'espai d'accions};
  \node[studioBox, right=of validate] (predict) {Predir\\acció};
  \node[studioBox, right=of predict] (step) {Pas\\EnergyPlus};
  \node[studioBox, below=of step] (status) {Progrés\\i avisos};
  \node[studioBox, left=of status] (ui) {Interfície\\Streamlit};
  \draw[studioEdge] (model) -- (metadata);
  \draw[studioEdge] (metadata) -- (validate);
  \draw[studioEdge] (env) -- (validate);
  \draw[studioEdge] (validate) -- (predict);
  \draw[studioEdge] (predict) -- (step);
  \draw[studioEdge] (step) -- (status);
  \draw[studioEdge] (status) -- (ui);
  \draw[studioEdge] (step) to[bend left=35] node[above,font=\scriptsize] {següent observació} (predict);
\end{tikzpicture}}
\caption{Bucle d'execució de control en viu.}
\end{figure}
\diagramnote{El bucle de control comprova metadades i compatibilitat abans de predir. Després aplica l'acció a EnergyPlus, rep la següent observació i actualitza la UI. La fletxa de retorn indica que el procés es repeteix a cada pas.}

\section{Dependències entre fitxers}

Els diagrames següents mostren les dependències principals entre fitxers. No
inclouen cada import petit de suport, sinó les relacions que cal entendre per
defensar el projecte: pàgina, components, backend, artefactes i classes de dades.

\subsection{Creació d'entorns}

\begin{figure}[H]
\centering
\resizebox{\linewidth}{!}{%
\begin{tikzpicture}[node distance=1.0cm and 1.0cm]
  \node[studioBlue, text width=3.8cm] (page) {Pàgina\\\texttt{pages/Afegir\_Entorn.py}};
  \node[studioWide, text width=4.4cm, right=of page] (facade) {Orquestració\\\texttt{afegir\_entorn\_backend.py}};
  \node[studioWide, text width=4.7cm, right=of facade] (analysis) {Lectura i conversió\\\texttt{afegir\_entorn\_analysis}\\\texttt{afegir\_entorn\_conversion}};
  \node[studioWide, text width=4.7cm, right=of analysis] (config) {Configuració segura\\\texttt{afegir\_entorn\_config}\\\texttt{afegir\_entorn\_types}\\\texttt{afegir\_entorn\_utils}};
  \node[studioWide, text width=4.2cm, right=of config] (support) {Context del projecte\\\texttt{weather\_profiles}\\\texttt{paths}};
  \node[studioAmber, text width=3.3cm, right=of support] (yaml) {YAML\\entorn registrat};
  \draw[studioEdge] (page) -- (facade);
  \draw[studioEdge] (facade) -- (analysis);
  \draw[studioEdge] (analysis) -- (config);
  \draw[studioEdge] (config) -- (support);
  \draw[studioEdge] (support) -- (yaml);
\end{tikzpicture}}
\caption{Dependències dels fitxers que creen i registren entorns Sinergym.}
\end{figure}
\diagramnote{Aquest flux comença a la pàgina d'alta d'entorns i passa pel backend principal, que reparteix la feina entre lectura del model, conversió de fitxers, construcció del YAML i validació. Els fitxers de tipus i utilitats eviten repetir estructures; \texttt{weather\_profiles} i \texttt{paths} situen l'entorn dins del projecte.}

\subsection{Entrenament d'agents}

\begin{figure}[H]
\centering
\resizebox{\linewidth}{!}{%
\begin{tikzpicture}[node distance=1.0cm and 0.95cm]
  \node[studioBlue, text width=3.6cm] (page) {Pàgina\\\texttt{pages/Entrenar\_Agent.py}};
  \node[studioWide, text width=4.5cm, right=of page] (ui) {Components UI\\\texttt{training\_page}\\\texttt{training\_agent}\\\texttt{training\_rewards}\\\texttt{training\_wrappers}\\\texttt{training\_shared}};
  \node[studioWide, text width=4.2cm, right=of ui] (facade) {Backend principal\\\texttt{entrenar\_agent\_backend.py}};
  \node[studioWide, text width=4.6cm, right=of facade] (prepare) {Sessió i configuració\\\texttt{entrenar\_agent\_session}\\\texttt{entrenar\_agent\_env}\\\texttt{entrenar\_agent\_rewards}\\\texttt{entrenar\_agent\_wrappers}};
  \node[studioWide, text width=4.4cm, right=of prepare] (run) {Execució RL\\\texttt{entrenar\_agent\_runtime}\\\texttt{training\_scripts}\\\texttt{sb3\_utils}};
  \node[studioAmber, text width=3.6cm, right=of run] (out) {Artefactes\\\texttt{entrenar\_agent\_artifacts}\\models\\logs\\config};
  \draw[studioEdge] (page) -- (ui);
  \draw[studioEdge] (ui) -- (facade);
  \draw[studioEdge] (facade) -- (prepare);
  \draw[studioEdge] (prepare) -- (run);
  \draw[studioEdge] (run) -- (out);
\end{tikzpicture}}
\caption{Dependències principals del flux d'entrenament.}
\end{figure}
\diagramnote{La pàgina només recull opcions i mostra estat. Els components separen pestanyes i formularis, el backend principal valida la petició i els mòduls de sessió, entorn, recompenses i wrappers preparen una configuració coherent. L'execució real queda a \texttt{entrenar\_agent\_runtime}, \texttt{training\_scripts} i \texttt{sb3\_utils}; al final es deixen models, logs i configuració reproduïble.}

\subsection{Resultats, dashboard i informes}

\begin{figure}[H]
\centering
\resizebox{\linewidth}{!}{%
\begin{tikzpicture}[node distance=1.0cm and 1.0cm]
  \node[studioBlue, text width=3.4cm] (page) {Pàgina\\\texttt{pages/Resultats.py}};
  \node[studioWide, text width=4.4cm, right=of page] (backend) {Selecció de runs\\\texttt{resultats\_backend}\\\texttt{resultats\_dashboard}};
  \node[studioWide, text width=4.4cm, right=of backend] (data) {Dades i KPIs\\\texttt{data\_loader}\\\texttt{kpis}\\\texttt{metrics}\\\texttt{figures\_zones}};
  \node[studioWide, text width=4.8cm, right=of data] (figures) {Famílies de gràfics\\\texttt{figures\_common}\\battery, comfort, hvac, thermal\\control, energy, indoor};
  \node[studioWide, text width=4.2cm, right=of figures] (report) {Exportació\\\texttt{resultats\_report\_backend}\\\texttt{style}\\\texttt{time\_utils}};
  \node[studioAmber, text width=3.2cm, right=of report] (out) {Dashboard\\HTML/PDF};
  \draw[studioEdge] (page) -- (backend);
  \draw[studioEdge] (backend) -- (data);
  \draw[studioEdge] (data) -- (figures);
  \draw[studioEdge] (figures) -- (report);
  \draw[studioEdge] (report) -- (out);
\end{tikzpicture}}
\caption{Dependències entre fitxers del flux de resultats i exportació.}
\end{figure}
\diagramnote{El flux de resultats primer localitza la run i els fitxers necessaris. Després carrega CSV, calcula mètriques i prepara dades de zona. Les famílies de gràfics comparteixen format i criteris visuals; l'exportador reutilitza aquestes figures per generar l'informe HTML/PDF sense duplicar càlculs.}

\subsection{Simulació, avaluació i control en viu}

\begin{figure}[H]
\centering
\resizebox{\linewidth}{!}{%
\begin{tikzpicture}[node distance=1.0cm and 1.0cm]
  \node[studioBlue, text width=4.0cm] (pages) {Pàgines\\\texttt{Simular\_Entorn.py}\\\texttt{Avaluar\_Agent.py}\\\texttt{Interaccionar\_amb\_l'Agent.py}};
  \node[studioWide, text width=4.5cm, right=of pages] (backends) {Backends d'execució\\\texttt{simular\_entorn\_backend}\\\texttt{avaluar\_agent\_backend}\\\texttt{interaccionar\_agent\_backend}};
  \node[studioWide, text width=4.2cm, right=of backends] (shared) {Peces compartides\\\texttt{sb3\_utils}\\\texttt{entrenar\_agent\_backend}};
  \node[studioAmber, text width=3.6cm, right=of shared] (engine) {Sinergym\\EnergyPlus\\models\\VecNormalize};
  \node[studioAmber, text width=3.3cm, right=of engine] (out) {CSV\\mètriques\\estat en viu};
  \draw[studioEdge] (pages) -- (backends);
  \draw[studioEdge] (backends) -- (shared);
  \draw[studioEdge] (shared) -- (engine);
  \draw[studioEdge] (engine) -- (out);
\end{tikzpicture}}
\caption{Dependències dels fluxos d'execució sense editar entorns.}
\end{figure}
\diagramnote{Aquests tres fluxos comparteixen la idea d'executar un entorn ja existent. La simulació genera una línia base, l'avaluació carrega un agent entrenat i el control en viu repeteix prediccions pas a pas. \texttt{sb3\_utils} centralitza càrrega i formes d'acció perquè els tres camins no tinguin regles diferents.}

\subsection{Visors i gestió de fitxers}

\begin{figure}[H]
\centering
\resizebox{\linewidth}{!}{%
\begin{tikzpicture}[node distance=1.0cm and 1.0cm]
  \node[studioBlue, text width=4.0cm] (pages) {Pàgines\\\texttt{Mostrar\_Entorn.py}\\\texttt{Visor\_Climes\_EPW.py}\\\texttt{Gestionar\_Arxius.py}};
  \node[studioWide, text width=4.4cm, right=of pages] (backends) {Backends\\\texttt{mostrar\_entorn\_backend}\\\texttt{visor\_epw\_backend}\\\texttt{gestionar\_arxius\_backend}};
  \node[studioWide, text width=4.2cm, right=of backends] (viewers) {Visualització\\\texttt{viewer\_3d\_backend}\\\texttt{mapa\_backend}\\\texttt{epw\_figures}};
  \node[studioWide, text width=3.7cm, right=of viewers] (paths) {Fitxers\\\texttt{paths}\\IDF\\EPW\\models};
  \node[studioAmber, text width=3.1cm, right=of paths] (out) {Mapa\\3D\\gràfics\\gestió};
  \draw[studioEdge] (pages) -- (backends);
  \draw[studioEdge] (backends) -- (viewers);
  \draw[studioEdge] (viewers) -- (paths);
  \draw[studioEdge] (paths) -- (out);
\end{tikzpicture}}
\caption{Dependències dels visors d'entorn/clima i del gestor d'arxius.}
\end{figure}
\diagramnote{Aquest bloc agrupa les pantalles que no entrenen agents. Els backends preparen metadades i fitxers; els mòduls de visualització converteixen aquesta informació en mapa, visor 3D o gràfics EPW. \texttt{paths} manté rutes consistents perquè el gestor d'arxius i els visors mirin el mateix lloc.}

\section{Classes i dataclasses principals}

\begin{figure}[H]
\centering
\resizebox{\linewidth}{!}{%
\begin{tikzpicture}[node distance=0.8cm and 0.9cm]
  \node[classBox] (upgrade) {
    \textbf{UpgradeResult}
    \nodepart{second}
    path\\version\\warnings
    \nodepart{third}
    resultat de conversió IDF
  };
  \node[classBox, right=of upgrade] (validation) {
    \textbf{EnvironmentValidationResult}
    \nodepart{second}
    ok\\message\\warnings
    \nodepart{third}
    prova de fum de l'entorn
  };
  \node[classBox, right=of validation] (analysis) {
    \textbf{BuildingTrainingAnalysis}
    \nodepart{second}
    variables\\meters\\actuators
    \nodepart{third}
    base per crear YAML
  };
  \node[classBox, below=of analysis] (actuator) {
    \textbf{ActuatorOption}
    \nodepart{second}
    name\\component\\control type
    \nodepart{third}
    opció d'acció candidata
  };
  \node[classBox, left=of actuator] (probe) {
    \textbf{ProbeReward}
    \nodepart{second}
    reward zero\\inspecció segura
    \nodepart{third}
    evita efectes durant l'anàlisi
  };
  \node[classBox, right=of actuator] (weather) {
    \textbf{WeatherProfileSuggestion}
    \nodepart{second}
    climate\\profile\\rationale
    \nodepart{third}
    proposta de perfil EPW
  };
  \draw[studioEdge] (upgrade) -- (validation);
  \draw[studioEdge] (validation) -- (analysis);
  \draw[studioEdge] (analysis) -- (actuator);
  \draw[studioEdge] (probe) -- (actuator);
  \draw[studioEdge] (actuator) -- (weather);
\end{tikzpicture}}
\caption{Classes i dataclasses del flux de creació d'entorns.}
\end{figure}
\diagramnote{\texttt{UpgradeResult} i \texttt{EnvironmentValidationResult} descriuen resultats de processos que poden fallar o generar avisos. \texttt{BuildingTrainingAnalysis} concentra allò que s'ha llegit de l'edifici i es connecta amb actuadors, recompenses de prova i suggeriments de clima per acabar generant una configuració entrenable.}

\begin{figure}[H]
\centering
\resizebox{\linewidth}{!}{%
\begin{tikzpicture}[node distance=0.8cm and 0.9cm]
  \node[classBox] (summary) {
    \textbf{EnvironmentSummary}
    \nodepart{second}
    env\\variables\\meters
    \nodepart{third}
    resum per simulació base
  };
  \node[classBox, right=of summary] (baseline) {
    \textbf{BaselineSimulationResult}
    \nodepart{second}
    run path\\elapsed\\cancelled
    \nodepart{third}
    retorn del baseline
  };
  \node[classBox, right=of baseline] (eval) {
    \textbf{EvaluationResult}
    \nodepart{second}
    elapsed\\cancelled\\warnings
    \nodepart{third}
    retorn de l'avaluació SB3
  };
  \node[classBox, below=of baseline] (baselineReward) {
    \textbf{ScheduleBaselineReward}
    \nodepart{second}
    reward zero\\mètriques físiques
    \nodepart{third}
    conserva schedules EnergyPlus
  };
  \node[classBox, right=of baselineReward] (costLogger) {
    \textbf{EnergyCostFileLogger}
    \nodepart{second}
    step\\cost\\energy
    \nodepart{third}
    CSV addicional de cost
  };
  \node[classBox, right=of costLogger] (epwState) {
    \textbf{EpwDashboardState}
    \nodepart{second}
    path\\metadata\\figures
    \nodepart{third}
    estat del visor EPW
  };
  \draw[studioEdge] (summary) -- (baseline);
  \draw[studioEdge] (baseline) -- (eval);
  \draw[studioEdge] (baselineReward) -- (baseline);
  \draw[studioEdge] (baselineReward) -- (costLogger);
  \draw[studioEdge] (costLogger) -- (epwState);
\end{tikzpicture}}
\caption{Classes del flux d'execució: baseline, entrenament, avaluació i clima.}
\end{figure}
\diagramnote{\texttt{EnvironmentSummary}, \texttt{BaselineSimulationResult} i \texttt{EvaluationResult} són els retorns grans que les pàgines poden mostrar sense conèixer els detalls d'EnergyPlus. A sota hi ha classes de suport: una recompensa neutra per executar baselines, un logger de cost i l'estat del visor EPW.}

\begin{figure}[H]
\centering
\resizebox{\linewidth}{!}{%
\begin{tikzpicture}[node distance=0.8cm and 0.9cm]
  \node[classBox] (runOption) {
    \textbf{RunOption}
    \nodepart{second}
    label\\path\\kind
    \nodepart{third}
    opció al selector de runs
  };
  \node[classBox, right=of runOption] (pageState) {
    \textbf{ResultsPageState}
    \nodepart{second}
    options\\selected\\warning
    \nodepart{third}
    estat inicial de Resultats
  };
  \node[classBox, right=of pageState] (artifacts) {
    \textbf{RunArtifacts}
    \nodepart{second}
    progress\\observations\\metadata
    \nodepart{third}
    fitxers necessaris d'una run
  };
  \node[classBox, right=of artifacts] (dashboard) {
    \textbf{DashboardData}
    \nodepart{second}
    progress\\obs\\metrics
    \nodepart{third}
    alimenta dashboard i informe
  };
  \node[classBox, right=of dashboard] (figure) {
    \textbf{ReportFigure}
    \nodepart{second}
    title\\figure\\description
    \nodepart{third}
    gràfic exportable
  };
  \node[classBox, below=of pageState] (explorer) {
    \textbf{ExplorerItem}
    \nodepart{second}
    path\\kind\\metadata
    \nodepart{third}
    element del gestor d'arxius
  };
  \draw[studioEdge] (runOption) -- (pageState);
  \draw[studioEdge] (pageState) -- (artifacts);
  \draw[studioEdge] (artifacts) -- (dashboard);
  \draw[studioEdge] (dashboard) -- (figure);
\end{tikzpicture}}
\caption{Classes i dataclasses de resultats, informes i gestió de fitxers.}
\end{figure}
\diagramnote{El primer tram explica el camí normal de Resultats: opció seleccionable, estat de pàgina, fitxers trobats, dades preparades i figures exportables. \texttt{ExplorerItem} queda separat perquè representa elements del gestor d'arxius, no del càlcul principal de KPIs.}


\appendix

\section{Annex de codi funcional}

Les llistes següents inclouen el codi Python funcional de BEMS-RL Studio dins d'aquest mateix document LaTeX. S'exclouen només els mòduls d'estil CSS pur (\texttt{page\_styles/*.py} i \texttt{ui\_theme.py}) perquè no aporten lògica d'aplicació, dades, entrenament, simulació o generació d'informes. Les icones Unicode es representen com a escapes \texttt{\textbackslash u...} o \texttt{\textbackslash U...} per evitar glifs buits en el PDF.

\section{Codi: Entrypoints i navegacio}

\Needspace{8\baselineskip}

\subsection{\texorpdfstring{\texttt{Inici.py}}{Inici.py}}

\begin{lstlisting}[style=studioPython,caption={Inici.py},label={lst:inici-py}]
"""Pàgina inicial de BEMS-RL STUDIO."""

from __future__ import annotations

import base64
from html import escape
from pathlib import Path
from subprocess import STDOUT, CalledProcessError, check_output

import streamlit as st

from page_styles.inici import inject_home_styles
from sidebar_nav import configure_studio_page


PAGE_TITLE = "BEMS-RL STUDIO"
PAGE_LAYOUT = "wide"
INTRODUCTION_TEXT = (
    "Benvingut a BEMS-RL STUDIO, una interfície gràfica per treballar "
    "amb Reinforcement Learning en simulacions energètiques d'edificis."
)
CAPABILITIES = (
    "Afegir nous entorns amb edificis i meteorologia.",
    "Visualitzar i inspeccionar els entorns disponibles abans de treballar-hi.",
    "Entrenar agents d'Aprenentatge per Reforç en un entorn Eplus personalitzat.",
    "Simular entorns energètics per validar comportaments i condicions de partida.",
    "Avaluar el comportament del teu agent.",
    "Interaccionar amb agents entrenats mitjançant el control en viu.",
    "Explorar resultats amb dashboards i gràfics.",
    "Gestionar fitxers, execucions, models, climes i entorns del projecte.",
    "Analitzar fitxers climàtics EPW amb resum, patrons anuals i roses dels vents.",
    "Consultar el centre d'ajuda per entendre el flux de treball de l'aplicació.",
)
NAVIGATION_HINT = "Utilitza el menú lateral per accedir a cada secció."
RESOURCES_DIR = Path(__file__).resolve().parent / "resources"
ENERGYPLUS_LOGO_PATH = RESOURCES_DIR / "energyplus-logo.png"


@st.cache_data(show_spinner=False)
def get_energyplus_logo_data_uri() -> str:
    """Carrega el logo d'EnergyPlus com a data URI per incrustar-lo a la UI."""

    if not ENERGYPLUS_LOGO_PATH.exists():
        return ""

    encoded_logo = base64.b64encode(ENERGYPLUS_LOGO_PATH.read_bytes()).decode("ascii")
    return f"data:image/png;base64,{encoded_logo}"


@st.cache_data(show_spinner=False)
def detect_energyplus_status() -> tuple[bool, str]:
    """Retorna si EnergyPlus està disponible i el missatge per mostrar a la UI."""

    try:
        version = check_output(["energyplus", "--version"], stderr=STDOUT).decode("utf-8").strip()
    except (CalledProcessError, FileNotFoundError, OSError):
        return (
            False,
            "No s'ha detectat una instal·lació d'EnergyPlus al sistema operatiu o no està disponible al PATH.",
        )

    if not version:
        return False, "EnergyPlus respon sense versió. Revisa la instal·lació del motor."

    return True, f"Motor actiu detectat: {version}"


def render_hero() -> None:
    """Mostra el bloc principal de presentació."""

    # Hero home
    st.markdown(
        f"""
        <section class="home-hero">
            <div class="hero-kicker">Plataforma de simulació energètica</div>
            <h1 class="hero-title">{escape(PAGE_TITLE)}</h1>
            <div class="hero-copy">{escape(INTRODUCTION_TEXT)}</div>
        </section>
        """,
        unsafe_allow_html=True,
    )


def render_status_panel(energyplus_available: bool, energyplus_message: str) -> None:
    """Mostra l'estat d'EnergyPlus amb una targeta visual."""

    status_title = "EnergyPlus detectat" if energyplus_available else "EnergyPlus no disponible"
    status_class = "status-ok" if energyplus_available else "status-warning"
    logo_data_uri = get_energyplus_logo_data_uri()
    logo_html = ""
    if logo_data_uri:
        logo_html = f'<img class="status-logo" src="{logo_data_uri}" alt="EnergyPlus logo">'

    # Targeta estat EnergyPlus
    st.markdown(
        f"""
        <section class="home-panel home-top-panel status-shell {status_class}">
            {logo_html}
            <div class="panel-kicker">Entorn de simulació</div>
            <div class="panel-title">{escape(status_title)}</div>
            <div class="panel-copy status-copy">{escape(energyplus_message)}</div>
        </section>
        """,
        unsafe_allow_html=True,
    )


def render_navigation_panel() -> None:
    """Dibuixa la pista de navegació lateral."""

    # Targeta navegacio home
    st.markdown(
        f"""
        <section class="home-panel home-top-panel">
            <div class="panel-kicker">Navegació</div>
            <div class="panel-title">Començar a treballar</div>
            <div class="panel-copy">{escape(NAVIGATION_HINT)}</div>
        </section>
        """,
        unsafe_allow_html=True,
    )


def render_capability_card(index: int, capability: str) -> None:
    """Crea una targeta de funcionalitat."""

    # Targeta funcionalitat home
    st.markdown(
        f"""
        <section class="feature-card">
            <div class="feature-kicker">Funció {index:02d}</div>
            <p class="feature-copy">{escape(capability)}</p>
        </section>
        """,
        unsafe_allow_html=True,
    )


def render_capability_grid() -> None:
    """Organitza les funcionalitats en files de tres targetes alineades."""

    row_size = 3
    for row_start in range(0, len(CAPABILITIES), row_size):
        row_items = CAPABILITIES[row_start : row_start + row_size]
        # Columnes funcionalitats
        row_columns = st.columns(row_size, gap="large")
        for column, capability_index, capability in zip(
            row_columns,
            range(row_start + 1, row_start + len(row_items) + 1),
            row_items,
        ):
            with column:
                render_capability_card(capability_index, capability)


def render_organization_footer() -> None:
    """Afegeix el bloc informatiu final."""

    # Targeta peu projecte
    st.markdown(
        """
        <section class="footer-card">
            <div class="panel-kicker">Projecte</div>
            <div class="footer-title">Desenvolupat per SUNO</div>
            <div class="footer-copy">
                Aquest projecte ha estat desenvolupat per
                <a href="https://www.suno.cat/" target="_blank">SUNO</a>,
                una empresa pionera en simulació energètica d'edificis,
                optimització de sistemes i integració de solucions
                d'intel·ligència artificial.
            </div>
        </section>
        """,
        unsafe_allow_html=True,
    )


def render_home_page() -> None:
    """Munta la pàgina inicial."""

    configure_studio_page(PAGE_TITLE, layout=PAGE_LAYOUT)
    inject_home_styles()
    energyplus_available, energyplus_message = detect_energyplus_status()

    render_hero()

    # Columnes estat i navegacio
    status_col, navigation_col = st.columns(2, gap="large")
    with status_col:
        render_status_panel(energyplus_available, energyplus_message)
    with navigation_col:
        render_navigation_panel()

    # Separador home
    st.markdown("<div class='studio-spacer-115'></div>", unsafe_allow_html=True)
    # Titol funcionalitats
    st.markdown('<h2 class="section-title">Què pots fer?</h2>', unsafe_allow_html=True)
    render_capability_grid()
    # Separador funcionalitats
    st.markdown('<div class="home-features-gap"></div>', unsafe_allow_html=True)
    render_organization_footer()


render_home_page()
\end{lstlisting}

\Needspace{8\baselineskip}

\subsection{\texorpdfstring{\texttt{\_\_init\_\_.py}}{\_\_init\_\_.py}}

\begin{lstlisting}[style=studioPython,caption={\_\_init\_\_.py},label={lst:init---py}]
"""Metadades de l'aplicació per a l'espai de treball Streamlit de BEMS-RL Studio.

El punt d'entrada executable és :mod:`Inici`; les pàgines individuals viuen sota
``pages`` i Streamlit les carrega automàticament. Aquest mòdul es manté sense
efectes secundaris perquè la documentació i les eines de diagnòstic puguin
importar l'arrel de Studio sense renderitzar cap pàgina.
"""

from __future__ import annotations

APP_NAME = "BEMS-RL Studio"
APP_ENTRYPOINT = "Inici.py"

__all__ = ["APP_NAME", "APP_ENTRYPOINT"]
\end{lstlisting}

\Needspace{8\baselineskip}

\subsection{\texorpdfstring{\texttt{sidebar\_nav.py}}{sidebar\_nav.py}}

\begin{lstlisting}[style=studioPython,caption={sidebar\_nav.py},label={lst:sidebar-nav-py}]
"""Navegació lateral compartida de BEMS-RL STUDIO."""

from __future__ import annotations

import streamlit as st
from page_styles.sidebar import inject_sidebar_shell_styles


APP_NAME = "BEMS-RL STUDIO"

SIDEBAR_SECTIONS: tuple[tuple[str, str, tuple[tuple[str, str, str], ...]], ...] = (
    (
        "ENTORNS I AGENTS",
        "Configura, entrena i analitza",
        (
            ("Inici.py", "Inici", ":material/space_dashboard:"),
            ("pages/Afegir_Entorn.py", "Crear Entorn", ":material/home_work:"),
            ("pages/Mostrar_Entorn.py", "Mostrar Entorn", ":material/apartment:"),
            ("pages/Entrenar_Agent.py", "Entrenar Agent", ":material/model_training:"),
            ("pages/Resultats.py", "Resultats", ":material/insights:"),
            ("pages/Simular_Entorn.py", "Simular Entorn", ":material/play_circle:"),
            ("pages/Avaluar_Agent.py", "Avaluar Agent", ":material/fact_check:"),
            ("pages/Interaccionar_amb_l'Agent.py", "Control en Viu", ":material/joystick:"),
        ),
    ),
    (
        "GESTIÓ",
        "Organitza i inspecciona els fitxers del projecte",
        (
            ("pages/Gestionar_Arxius.py", "Arxius", ":material/folder_open:"),
            ("pages/Visor_Climes_EPW.py", "Visor EPW", ":material/cloud:"),
        ),
    ),
    (
        "SUPORT",
        "Guies i ajuda",
        (
            ("pages/Ajuda.py", "Centre d'Ajuda", ":material/support_agent:"),
        ),
    ),
)


def configure_studio_page(page_title: str, *, layout: str = "wide") -> None:
    """Configura una pàgina i renderitza la barra lateral compartida."""

    browser_title = page_title if page_title.upper().startswith(APP_NAME) else f"{APP_NAME} - {page_title}"
    st.set_page_config(
        page_title=browser_title,
        page_icon=":material/neurology:",
        layout=layout,
        initial_sidebar_state="expanded",
    )
    render_studio_sidebar()


def render_studio_sidebar() -> None:
    """Mostra la navegació personalitzada de la barra lateral de Streamlit."""

    inject_sidebar_shell_styles()

    with st.sidebar:
        # Marca de la barra lateral
        st.markdown(
            """
            <div class="studio-sidebar-brand">
                <div class="studio-sidebar-brand-title">BEMS-RL</div>
                <div class="studio-sidebar-brand-subtitle">STUDIO</div>
            </div>
            <div class="studio-sidebar-divider"></div>
            """,
            unsafe_allow_html=True,
        )

        for section_title, section_copy, items in SIDEBAR_SECTIONS:
            _render_section_header(section_title, section_copy)
            for page_path, label, icon in items:
                # Enllaç de pàgina a la barra lateral
                st.page_link(page_path, label=label, icon=icon)


def _render_section_header(title: str, copy_text: str) -> None:
    """Crea la capçalera de la secció."""
    # Capçalera de secció de la barra lateral
    st.markdown(
        f"""
        <div class="studio-sidebar-section">{title}</div>
        <div class="studio-sidebar-section-copy">{copy_text}</div>
        """,
        unsafe_allow_html=True,
    )
\end{lstlisting}

\section{Codi: Pagines Streamlit}

\Needspace{8\baselineskip}

\subsection{\texorpdfstring{\texttt{pages/Afegir\_Entorn.py}}{pages/Afegir\_Entorn.py}}

\begin{lstlisting}[style=studioPython,caption={pages/Afegir\_Entorn.py},label={lst:pages-afegir-entorn-py}]
"""Sinergym Studio: vista per afegir entorns.

Aquesta vista gestiona la UI de Streamlit per carregar models d'edifici
EnergyPlus (.idf/.epJSON) i fitxers meteorològics (.epw), configurar els
actuadors HVAC i la variabilitat climàtica, i registrar un entorn Gym nou a
Sinergym mitjançant un fitxer YAML generat.
"""

import subprocess
from pathlib import Path
from typing import Any, Dict, List, Optional

import pandas as pd
import streamlit as st
import yaml

from backend import afegir_entorn_backend as add_env_backend
from page_components.ui_fragments import render_card_header, render_hero, render_section_title
from page_styles.afegir_entorn import inject_add_environment_styles
from sidebar_nav import configure_studio_page


PAGE_TITLE = "Crear Nou Entorn"
PAGE_LAYOUT = "wide"
INTRODUCTION_TEXT = (
    "Carrega un edifici i un únic clima, detecta controladors HVAC i bateria "
    "usables del model i configura l'espai d'accions abans de registrar el nou entorn."
)


def render_add_environment_hero() -> None:
    """Mostra la capçalera principal amb el mateix llenguatge visual de l'app."""

    render_hero("add-environment-hero", "CREACIÓ D'ENTORNS", PAGE_TITLE, INTRODUCTION_TEXT)


def render_weather_source_selector(target=st) -> str:
    """Mostra un control segmentat o un grup de ràdio per seleccionar la font meteorològica.

    Paràmetres:
        target: un contenidor Streamlit per vincular el widget. Per defecte a st.

    Retorna:
        str: El mode seleccionat ("Escollir existents" o "Pujar nous fitxers").
    """

    options = ["Escollir existents", "Pujar nous fitxers"]
    current_value = st.session_state.get("add_env_weather_mode", options[0])
    target.markdown('<div class="studio-segmented-pill-anchor"></div>', unsafe_allow_html=True)

    # Selector origen clima
    selected_value = target.segmented_control(
        "Origen del clima",
        options,
        default=current_value,
        key="add_env_weather_mode",
    )
    return selected_value or current_value


def format_current_schedule_range(option) -> str:
    """Formata l'interval numèric vàlid d'un actuador per mostrar-lo en funció de la seva planificació inicial.

    Paràmetres:
        option (ActuatorOption): l'opció de l'actuador detectada des del model.

    Retorna:
        str: una cadena preformatada que presenta l'interval numèric detectat.
    """
    if option.current_low is None or option.current_high is None:
        return "Sense resum automàtic"
    if option.category == "availability":
        if abs(option.current_low - option.current_high) < 1e-6:
            return f"{option.current_low:.0f}"
        return f"{option.current_low:.0f} – {option.current_high:.0f}"
    if abs(option.current_low - option.current_high) < 1e-6:
        return f"{option.current_low:.2f}"
    return f"{option.current_low:.2f} – {option.current_high:.2f}"


def render_controller_selection_table(
    title: str,
    options,
    *,
    default_selected: bool,
    key: str,
    empty_message: str,
):
    """Mostra una taula interactiva per incloure/excloure controladors HVAC.

    Paràmetres:
        title (str): l'encapçalament que es mostra a sobre de la taula.
        options (Sequence[ActuatorOption]): les opcions de l'actuador disponibles.
        default_selected (bool): Si True, les caselles de selecció estan marcades inicialment.
        key (str): clau única d'estat de sessió Streamlit per a l'editor de dades.
        empty_message (str): Text alternativa quan no hi ha opcions per mostrar.

    Retorna:
        List[str]: una llista de les cadenes `option_id` seleccionades.
    """
    if not options:
        # Text taula buida
        st.caption(empty_message)
        return []

    # Titol taula controladors
    st.markdown(f"**{title}**")
    selection_df = pd.DataFrame(
        [
            {
                "Afegir": default_selected,
                "Control": option.label,
                "Afecta": option.reference,
                "Programa actual": format_current_schedule_range(option),
            }
            for option in options
        ],
        index=[option.option_id for option in options],
    )
    # Editor seleccio controladors
    edited_selection_df = st.data_editor(
        selection_df,
        hide_index=True,
        width="stretch",
        key=key,
        disabled=["Control", "Afecta", "Programa actual"],
        column_config={
            "Afegir": st.column_config.CheckboxColumn("Incloure", width="small"),
            "Control": st.column_config.TextColumn("Controlador"),
            "Afecta": st.column_config.TextColumn("Afecta"),
            "Programa actual": st.column_config.TextColumn("Programa actual"),
        },
    )
    return [
        option_id
        for option_id, row in edited_selection_df.iterrows()
        if bool(row["Afegir"])
    ]


def render_building_upload_card() -> Any:
    """Mostra la targeta d'upload de l'edifici i retorna el fitxer penjat.

    Retorna:
        Any: l'objecte del fitxer carregat Streamlit o None si no s'ha carregat cap fitxer.
    """
    # Targeta upload edifici
    building_card = st.container(border=True)
    render_card_header(
        building_card,
        anchor_class="upload-card-shell-anchor",
        kicker="Edifici",
        title="Fitxer d'edifici",
        description="Puja un fitxer IDF o epJSON per preparar el model base.",
    )
    # Upload fitxer edifici
    return building_card.file_uploader(
        "Fitxer d'edifici",
        type=["idf", "epjson", "epJSON"],
        label_visibility="collapsed",
    )


def render_weather_upload_card() -> tuple[Path, Any]:
    """Mostra la targeta del fitxer meteorològic i torna el camí resolt al fitxer EPW seleccionat.

    Gestiona tant el selector de fitxers existents com els modes de càrrega de fitxers nous.
    Crida st.stop() si no hi ha cap fitxer meteorològic vàlid disponible o seleccionat.

    Retorna:
        tuple[Path, Any]: el camí del fitxer EPW seleccionat i el seu contenidor de la targeta Streamlit.
    """
    # Targeta fitxer clima
    climate_card = st.container(border=True)
    render_card_header(
        climate_card,
        anchor_class="upload-card-shell-anchor",
        kicker="Clima",
        title="Fitxer climàtic",
        description="Escull un .epw existent o puja'n un de nou per registrar l'entorn.",
    )
    mode_weather = render_weather_source_selector(climate_card)

    selected_weather_path = None

    if mode_weather == "Escollir existents":
        weather_files = sorted(weather_file.name for weather_file in add_env_backend.WEATHERS_DIR.glob("*.epw"))
        if not weather_files:
            # Error sense EPW
            st.error("No hi ha fitxers .epw disponibles al catàleg actual.")
            st.stop()
        with climate_card:
            # Selector fitxer clima
            selected_weather_name = st.selectbox(
                "Selecciona el fitxer de clima",
                weather_files,
            )
        selected_weather_path = add_env_backend.WEATHERS_DIR / selected_weather_name
    else:
        with climate_card:
            # upload fitxer clima
            uploaded_weather_file = st.file_uploader(
                "Puja un fitxer .epw",
                type=["epw"],
                key="weather_upload",
            )
        if uploaded_weather_file:
            weather_upload_token = (uploaded_weather_file.name, uploaded_weather_file.size)
            cached_weather_path = st.session_state.get("add_env_weather_path")
            cached_weather_file = Path(cached_weather_path) if cached_weather_path else None
            if (
                st.session_state.get("add_env_weather_token") != weather_upload_token
                or cached_weather_file is None
                or not cached_weather_file.exists()
            ):
                cached_weather_file = add_env_backend.WEATHERS_DIR / uploaded_weather_file.name
                add_env_backend.save_uploaded_bytes(cached_weather_file, uploaded_weather_file.read())
                st.session_state.add_env_weather_token = weather_upload_token
                st.session_state.add_env_weather_path = str(cached_weather_file)
            selected_weather_path = cached_weather_file
            with climate_card:
                # Missatge clima carregat
                st.success(f"Fitxer de clima carregat: {uploaded_weather_file.name}")

    if selected_weather_path is None:
        # Avís seleccionar clima
        st.info("Selecciona un fitxer climàtic per definir l'entorn.")
        st.stop()

    return selected_weather_path, climate_card


def render_weather_variability_section() -> tuple[Optional[Dict], bool]:
    """Mostra la configuració opcional de variabilitat meteorològica estocàstica.

    Retorna:
        tuple[Optional[Dict], bool]: configuració de la variabilitat del clima i si
        les variants estocàstiques estan habilitades.
    """
    render_section_title("Variabilitat climàtica opcional", "section-title", level=3)
    # Targeta variabilitat clima
    variability_card = st.container(border=True)
    variability_card.markdown('<div class="weather-variability-card-anchor"></div>', unsafe_allow_html=True)
    # Checkbox variants estocastiques
    enable_stochastic = variability_card.checkbox(
        "Generar també variants estocàstiques",
        value=True,
        help="Quan s'activa, es registren també entorns amb el sufix -continuous-stochastic-v1.",
    )

    if not enable_stochastic:
        return None, False

    default_sigma, default_mu, default_tau = add_env_backend.DEFAULT_WEATHER_VARIABILITY["Dry Bulb Temperature"]
    # Inputs variabilitat clima
    sigma_col, mu_col, tau_col = variability_card.columns(3)
    with sigma_col:
        # Input sigma clima
        sigma = st.number_input("Sigma", min_value=0.0, value=float(default_sigma), step=0.1)
    with mu_col:
        # Input mu clima
        mu = st.number_input("Mu", value=float(default_mu), step=0.1)
    with tau_col:
        # Input tau clima
        tau = st.number_input("Tau (hores)", min_value=0.1, value=float(default_tau), step=1.0)

    return {"Dry Bulb Temperature": [float(sigma), float(mu), float(tau)]}, True


@st.cache_data(show_spinner=False)
def load_weather_temperature_preview(
    weather_path: str,
    weather_mtime_ns: int,
    sigma: float,
    mu: float,
    tau: float,
) -> List[Dict[str, float]]:
    """Carrega les dades de previsualització meteorològica anual a la memòria cau per al fitxer EPW seleccionat.

    Paràmetres:
        weather_path (str): camí cap al fitxer meteorològic EPW.
        weather_mtime_ns (int): marca de temps de modificació de fitxers utilitzada com a clau d'invalidació de la memòria cau.
        sigma (float): paràmetre de desviació estàndard per al procés d'OU.
        mu (float): paràmetre mitjà per al procés d'OU.
        tau (float): paràmetre constant de temps per al procés d'OU.

    Retorna:
        List[Dict[str, float]]: registres de previsualització diaris preparats per a la traça.
    """

    del weather_mtime_ns
    return add_env_backend.build_weather_temperature_preview(Path(weather_path), sigma, mu, tau)


def render_weather_temperature_preview(
    selected_weather_path: Path,
    weather_variability: Optional[Dict],
) -> None:
    """Mostra la previsualització anual de la temperatura EPW per als canvis meteorològics estocàstics.

    Paràmetres:
        selected_weather_path (Path): fitxer EPW seleccionat o carregat actualment.
        weather_variability (Optional[Dict]): configuració activa de la variabilitat del clima.
    """

    if not weather_variability:
        return

    sigma, mu, tau = weather_variability["Dry Bulb Temperature"]
    try:
        weather_mtime_ns = selected_weather_path.stat().st_mtime_ns
    except OSError:
        weather_mtime_ns = 0

    preview_records = load_weather_temperature_preview(
        str(selected_weather_path),
        weather_mtime_ns,
        float(sigma),
        float(mu),
        float(tau),
    )
    if not preview_records:
        # Avís preview clima
        st.info("No s'ha pogut llegir la temperatura seca del fitxer EPW seleccionat.")
        return

    # Targeta previsualitzacio clima
    preview_card = st.container(border=True)
    render_card_header(
        preview_card,
        anchor_class="weather-preview-card-anchor",
        kicker="Clima modificat",
        title="Previsualització anual de temperatura",
    )
    # Grafic temperatura clima
    preview_card.plotly_chart(
        add_env_backend.build_weather_temperature_preview_figure(preview_records),
        width="stretch",
        config={"displayModeBar": False},
    )


def prepare_building_file(uploaded_file: Any) -> tuple[Path, Path]:
    """Desa, actualitza si cal i converteix el fitxer de construcció penjat a epJSON.

    Utilitza l'estat de la sessió per emmagatzemar el resultat a la memòria cau i evitar el processament redundant en les reexecucions.
    Crida st.stop() per errors irrecuperables.

    Paràmetres:
        uploaded_file: l'objecte de fitxer carregat Streamlit per al model d'edifici.

    Retorna:
        Tuple[Path, Path]: el camí desat en brut (tmp_path) i el camí final epJSON (building_path).
    """
    uploaded_token = (uploaded_file.name, uploaded_file.size)
    cached_tmp_path = st.session_state.get("add_env_tmp_path")
    cached_building_path = st.session_state.get("add_env_building_path")

    tmp_path = Path(cached_tmp_path) if cached_tmp_path else None
    building_path = Path(cached_building_path) if cached_building_path else None

    # Streamlit reexecuta la pàgina cada vegada que es toca un widget; aquest token evita
    # reconvertir el mateix fitxer una vegada i una altra.
    needs_prepare = (
        st.session_state.get("add_env_current_building_upload_token") != uploaded_token
        or tmp_path is None
        or building_path is None
        or not tmp_path.exists()
        or not building_path.exists()
    )

    if not needs_prepare:
        return Path(st.session_state["add_env_tmp_path"]), Path(st.session_state["add_env_building_path"])

    st.session_state.add_env_current_building_upload_token = uploaded_token
    tmp_path = add_env_backend.BUILDINGS_DIR / uploaded_file.name
    add_env_backend.save_uploaded_bytes(tmp_path, uploaded_file.read())

    if tmp_path.suffix.lower() == ".idf":
        with st.spinner("Escanejant i actualitzant l'arxiu IDF, si cal..."):
            detected_version = add_env_backend.detect_idf_version(tmp_path)
            if detected_version and add_env_backend.needs_idf_upgrade(detected_version):
                # EnergyPlus només converteix bé IDF de versions properes. Si el fitxer és
                # antic, el passem primer pels Transition tools oficials.
                updater_dir = add_env_backend.get_transition_updater_dir()
                if updater_dir is None:
                    with st.spinner("Descarregant les eines de transició d'EnergyPlus..."):
                        updater_dir = add_env_backend.download_transition_tools()

                if updater_dir is not None:
                    # Avís actualitzacio IDF
                    st.info(
                        "S'ha detectat un fitxer IDF antic "
                        f"({detected_version[0]}.{detected_version[1]}). "
                        "S'actualitzarà automàticament abans de convertir-lo."
                    )
                    try:
                        upgrade_result = add_env_backend.upgrade_idf_version(tmp_path, updater_dir)
                    except subprocess.TimeoutExpired:
                        # Error timeout actualitzacio
                        st.error(
                            "L'actualització automàtica de l'IDF ha superat el temps límit. "
                            "Revisa el fitxer o converteix-lo prèviament abans de continuar."
                        )
                        st.stop()
                    if upgrade_result.failed_transition_version is not None:
                        failed_major, failed_minor = upgrade_result.failed_transition_version
                        # Avís transicio fallida
                        st.warning(
                            "La transició automàtica ha fallat a la versió "
                            f"{failed_major}.{failed_minor}. Revisa l'IDF abans de continuar."
                        )
                    elif upgrade_result.upgraded and upgrade_result.final_version is not None:
                        # Missatge IDF actualitzat
                        st.success(
                            "L'arxiu s'ha actualitzat correctament a la versió "
                            f"{upgrade_result.final_version[0]}.{upgrade_result.final_version[1]}."
                        )

    with st.spinner("Convertint i preparant el model de l'edifici..."):
        try:
            # A partir d'aquí treballem sempre amb epJSON; simplifica l'anàlisi posterior
            # i evita duplicar parser per IDF i epJSON.
            building_path = (
                add_env_backend.convert_idf_to_epjson(tmp_path) if tmp_path.suffix.lower() == ".idf" else tmp_path
            )
        except FileNotFoundError:
            # Error EnergyPlus absent
            st.error("No s'ha trobat `energyplus`. Assegura't que està instal·lat i disponible al PATH.")
            st.stop()
        except subprocess.TimeoutExpired:
            # Error timeout conversio
            st.error(
                "La conversió o preparació del model ha superat el temps límit. "
                "Revisa el fitxer d'edifici o prova una conversió prèvia a epJSON."
            )
            st.stop()
        except Exception as error:
            # Error conversio edifici
            st.error(f"Error en la conversió amb EnergyPlus: {error}")
            st.stop()

    st.session_state.add_env_tmp_path = str(tmp_path)
    st.session_state.add_env_building_path = str(building_path)
    # Missatge edifici preparat
    st.success(f"Fitxer carregat i preparat correctament: {building_path.name}")
    return tmp_path, building_path


def load_or_run_training_analysis(building_path: Path, weather_path: Path):
    """Retorna la memòria cau BuildingTrainingAnalysis o l'executa nova si les entrades han canviat.

    Paràmetres:
        building_path (Path): camí cap al model d'edifici epJSON.
        weather_path (Path): camí cap al fitxer meteorològic EPW.

    Retorna:
        BuildingTrainingAnalysis: Detectats termòstats, actuadors, variables i comptadors.
    """
    analysis_cache_key = (str(building_path.resolve()), str(weather_path.resolve()))
    if st.session_state.get("add_env_training_analysis_key") != analysis_cache_key:
        with st.spinner("Detectant termòstats i actuadors del model..."):
            st.session_state.add_env_training_analysis = add_env_backend.build_training_analysis(
                building_path,
                weather_path,
                probe_handlers=False,
            )
            st.session_state.add_env_training_analysis_key = analysis_cache_key
    return st.session_state.add_env_training_analysis


def render_environment_id_section(
    building_path: Path,
    weather_profiles: List[Dict],
    enable_stochastic: bool,
) -> str:
    """Mostra l'entrada d'ID base de l'entorn i previsualitza tots els ID que es registraran.

    Paràmetres:
        building_path (Path): S'utilitza per obtenir un nom base predeterminat de la base del fitxer.
        weather_profiles (List[Dict]): perfils meteorològics per calcular tots els ID generats.
        enable_stochastic (bool): si s'han d'incloure els ID de variants estocàstiques.

    Retorna:
        str: l'identificador de l'entorn base validat i netejat.
    """
    render_section_title("Configuració del nou entorn")

    with st.container(border=True):
        default_id_base = add_env_backend.sanitize_identifier(building_path.stem, "nou_entorn")
        # Input ID entorn
        id_base_input = st.text_input("**ID base de l'entorn**", value=default_id_base)
        id_base = add_env_backend.sanitize_identifier(id_base_input, default_id_base)
        if id_base != id_base_input:
            st.caption(f"L'identificador base es registrarà com a `{id_base}`.")

        generated_env_ids = [
            add_env_backend.build_registered_env_id(id_base, profile["key"])
            for profile in weather_profiles
        ]
        if enable_stochastic:
            generated_env_ids.extend(
                add_env_backend.build_registered_env_id(id_base, profile["key"], stochastic=True)
                for profile in weather_profiles
            )
        # Etiqueta IDs generats
        st.markdown(
            "<p class='add-env-generated-ids-label'>IDs que es registraran:</p>",
            unsafe_allow_html=True,
        )
        # Llista IDs generats
        st.code("\n".join(generated_env_ids), language="text")

    return id_base


def render_controller_selection_section(analysis) -> tuple[dict, list]:
    """Mostra la detecció del controlador i retorna els ID dels actuadors seleccionats.

    Crida st.stop() si no hi ha actuadors disponibles o no n'hi ha cap seleccionat.

    Paràmetres:
        analysis (BuildingTrainingAnalysis): resultat de l'anàlisi de l'edifici detectat.

    Retorna:
        Tuple[dict, list]: les opcions completes de l'actuador en forma de diccionari i la llista d'ID d'opcions seleccionades.
    """
    render_section_title("Controladors detectats")

    detected_actuator_options = [
        *list(analysis.thermostat_actuators),
        *list(analysis.hvac_actuators),
    ]
    actuator_options = {option.option_id: option for option in detected_actuator_options}

    if not actuator_options:
        # Error sense controladors
        st.error(
            "No s'ha detectat cap controlador usable al model. "
            "Cal preparar prèviament l'edifici amb termòstats, setpoints, bateria o altres controls editables abans de registrar l'entorn."
        )
        st.stop()

    thermostat_options = list(analysis.thermostat_actuators)
    scheduled_setpoint_options = [
        option for option in analysis.hvac_actuators if option.category == "scheduled_temperature_setpoint"
    ]
    availability_options = [
        option for option in analysis.hvac_actuators if option.category == "availability"
    ]
    battery_options = [
        option for option in analysis.hvac_actuators if option.category in {"battery_charge", "battery_discharge"}
    ]

    # Separem els controladors per família perquè l'usuari pugui deixar fora actuadors
    # massa experimentals sense perdre els termòstats principals.
    # Resum controladors detectats
    st.caption(
        f"S'han detectat {len(detected_actuator_options)} controladors usables. "
        "Selecciona només els que vulguis incloure."
    )

    selected_actuator_ids = []
    selected_actuator_ids.extend(
        render_controller_selection_table(
            "Termòstats de zona",
            thermostat_options,
            default_selected=True,
            key="add_env_select_thermostats",
            empty_message="No s'han detectat termòstats de zona seleccionables.",
        )
    )
    selected_actuator_ids.extend(
        render_controller_selection_table(
            "Setpoints de sistema HVAC",
            scheduled_setpoint_options,
            default_selected=False,
            key="add_env_select_hvac_setpoints",
            empty_message="No s'han detectat setpoints HVAC programats addicionals.",
        )
    )
    selected_actuator_ids.extend(
        render_controller_selection_table(
            "Disponibilitat HVAC",
            availability_options,
            default_selected=False,
            key="add_env_select_hvac_availability",
            empty_message="No s'han detectat controls de disponibilitat HVAC.",
        )
    )
    selected_actuator_ids.extend(
        render_controller_selection_table(
            "Bateria",
            battery_options,
            default_selected=True,
            key="add_env_select_battery",
            empty_message="No s'han detectat controls de bateria.",
        )
    )

    selected_actuator_ids = list(dict.fromkeys(selected_actuator_ids))
    if not selected_actuator_ids:
        # Avís sense actuadors
        st.info("Selecciona com a mínim un controlador per poder definir l'espai d'accions.")
        st.stop()

    return actuator_options, selected_actuator_ids


def render_action_space_section(actuator_options: dict, selected_actuator_ids: List[str]) -> dict:
    """Mostra l'editor de límits de l'espai d'acció i retorna els límits validats.

    Crida a st.stop() si algun actuador té un límit no vàlid (mínim >= màxim).

    Paràmetres:
        actuator_options (dict): Dict complet dels objectes ActuatorOption disponibles per ID.
        selected_actuator_ids (List[str]): ID dels actuadors seleccionats per l'usuari.

    Retorna:
        dict: Mapeig de option_id a tuples flotants (baix, alt).
    """
    render_section_title("Espai d'accions")

    selected_actuator_options = [actuator_options[option_id] for option_id in selected_actuator_ids]
    controls_df = pd.DataFrame(
        [
            {
                "Control": option.label,
                "Afecta": option.reference,
                "Programa actual": format_current_schedule_range(option),
                "Mínim": float(option.default_low),
                "Màxim": float(option.default_high),
            }
            for option in selected_actuator_options
        ],
        index=[option.option_id for option in selected_actuator_options],
    )
    # Editor rangs accions
    edited_controls_df = st.data_editor(
        controls_df,
        hide_index=True,
        width="stretch",
        key="add_env_actuator_bounds_editor",
        disabled=["Control", "Afecta", "Programa actual"],
        column_config={
            "Control": st.column_config.TextColumn("Controlador"),
            "Afecta": st.column_config.TextColumn("Afecta"),
            "Programa actual": st.column_config.TextColumn("Programa actual"),
            "Mínim": st.column_config.NumberColumn("Mínim", step=0.5, format="%.2f"),
            "Màxim": st.column_config.NumberColumn("Màxim", step=0.5, format="%.2f"),
        },
    )

    actuator_bounds = {}
    bounds_are_valid = True
    for option_id, row in edited_controls_df.iterrows():
        low_value = float(row["Mínim"])
        high_value = float(row["Màxim"])
        if low_value >= high_value:
            bounds_are_valid = False
        actuator_bounds[option_id] = (low_value, high_value)

    if not bounds_are_valid:
        # Error rangs accio
        st.error("Cada actuador ha de tenir un límit mínim estrictament inferior al màxim.")
        st.stop()

    return actuator_bounds


def render_registration_section(
    id_base: str,
    building_path: Path,
    tmp_path: Path,
    weather_profiles: List[Dict],
    analysis,
    actuator_options: dict,
    selected_actuator_ids: List[str],
    actuator_bounds: dict,
    weather_variability: Optional[Dict],
) -> None:
    """Gestiona el botó de registre, la creació i la validació de l'entorn.

    En cas d'èxit, mostra el YAML generat i executa una simulació de validació ràpida.
    En cas d'error, reverteix els fitxers escrits parcialment i mostra els detalls de l'error.

    Paràmetres:
        id_base (str): identificador de base netejat per al entorn.
        building_path (Path): Camí al fitxer de construcció epJSON preparat.
        tmp_path (Path): camí cap al fitxer de construcció original penjat.
        weather_profiles (List[Dict]): llista de perfils meteorològics utilitzat per a la generació de la configuració.
        analysis (BuildingTrainingAnalysis): resultat de l'anàlisi de l'edifici detectat.
        actuator_options (dict): totes les opcions de l'actuador detectades teclejades per option_id.
        selected_actuator_ids (List[str]): ID dels actuadors seleccionats per l'usuari.
        actuator_bounds (dict): Mapeig de option_id a tuples de lligat (baix, alt).
        weather_variability (Optional[Dict]): configuració de la variabilitat del clima estocàstic, o None.
    """
    render_section_title("Creació i registre")

    # Botó crear entorn
    if not st.button("Crear entorn automàticament"):
        return

    cfg_path = add_env_backend.CFG_DIR / f"{id_base}.yaml"
    selected_actuator_objects = [actuator_options[option_id] for option_id in selected_actuator_ids]

    try:
        # Primer escrivim i registrem l'entorn. Si qualsevol pas falla, netegem el YAML
        # i el fitxer temporal per no deixar configuracions mig creades.
        env_cfg = add_env_backend.build_environment_config(
            id_base=id_base,
            building_file_name=building_path.name,
            weather_profiles=weather_profiles,
            analysis=analysis,
            selected_actuators=selected_actuator_objects,
            actuator_bounds=actuator_bounds,
            weather_variability=weather_variability,
        )
        add_env_backend.write_yaml_config(cfg_path, env_cfg)
        add_env_backend.register_environment_from_yaml(cfg_path)
    except Exception as error:
        # Error registre entorn
        st.error(f"No s'ha pogut construir o registrar l'entorn: {error}")
        cleanup_messages = add_env_backend.cleanup_failed_environment(cfg_path, tmp_path)
        for level, message in cleanup_messages:
            if level == "warning":
                # Avís neteja entorn
                st.warning(message)
            else:
                # Error neteja entorn
                st.error(message)
        st.stop()

    # Missatge entorn creat
    st.success(f"Entorn creat i registrat correctament: {id_base}")
    with st.expander("Veure el YAML generat"):
        # YAML generat
        st.code(yaml.dump(env_cfg, sort_keys=False, allow_unicode=True), language="yaml")

    validation_env_id = add_env_backend.build_registered_env_id(id_base, weather_profiles[0]["key"])
    try:
        # La validació curta és una prova de fum: si falla aquí, fallaria després en training.
        with st.spinner(f"Muntant l'entorn {validation_env_id} per validar la configuració..."):
            add_env_backend.validate_registered_environment(validation_env_id, cfg_path)
    except Exception as error:
        # Error validacio entorn
        st.error(f"No s'ha pogut validar l'entorn registrat. Es revertirà la configuració: {error}")
        cleanup_messages = add_env_backend.cleanup_failed_environment(cfg_path, tmp_path)
        for level, message in cleanup_messages:
            if level == "warning":
                # Avís rollback entorn
                st.warning(message)
            else:
                # Error rollback entorn
                st.error(message)
        st.stop()

    # Missatge validacio correcta
    st.success(
        "S'ha realitzat una validació ràpida de simulació amb l'entorn registrat procedimentalment de Sinergym. "
        "La configuració d'accions i observacions s'ha validat correctament."
    )


def main() -> None:
    """Punt d'entrada a la pàgina Afegir Entorn Streamlit.

    Organitza el flux complet de la pàgina: estils, heroi, càrregues de fitxers, processament d'edificis,
    detecció d'actuadors, configuració de l'espai d'acció i registre de l'entorn.
    """
    configure_studio_page(PAGE_TITLE, layout=PAGE_LAYOUT)
    inject_add_environment_styles()
    render_add_environment_hero()

    render_section_title("Pujar fitxers")
    # Columnes upload fitxers
    building_column, climate_column = st.columns([1, 1.35], gap="large")
    with building_column:
        uploaded_file = render_building_upload_card()
    with climate_column:
        selected_weather_path, climate_card = render_weather_upload_card()

    selected_weather_paths = [selected_weather_path]
    weather_profile_suggestions = add_env_backend.build_weather_profile_suggestions(selected_weather_paths)
    with climate_card:
        # Text clima suggerit
        st.caption(weather_profile_suggestions[0].climate_label)
    weather_profiles = [
        {"key": suggestion.suggested_key, "weather_file": suggestion.file_name}
        for suggestion in weather_profile_suggestions
    ]

    weather_variability, enable_stochastic = render_weather_variability_section()
    render_weather_temperature_preview(selected_weather_path, weather_variability)

    if not uploaded_file:
        # Avís upload edifici
        st.info("Puja un fitxer d'edifici per continuar amb la detecció de termòstats i actuadors.")
        st.stop()

    tmp_path, building_path = prepare_building_file(uploaded_file)
    analysis = load_or_run_training_analysis(building_path, selected_weather_paths[0])

    if not analysis.thermostat_zones:
        # Error sense termostats
        st.error(
            "No s'han detectat zones amb termòstat. Aquesta alta automàtica està pensada per edificis amb heating/cooling setpoints configurables."
        )
        st.stop()

    id_base = render_environment_id_section(building_path, weather_profiles, enable_stochastic)
    actuator_options, selected_actuator_ids = render_controller_selection_section(analysis)
    actuator_bounds = render_action_space_section(actuator_options, selected_actuator_ids)

    render_registration_section(
        id_base=id_base,
        building_path=building_path,
        tmp_path=tmp_path,
        weather_profiles=weather_profiles,
        analysis=analysis,
        actuator_options=actuator_options,
        selected_actuator_ids=selected_actuator_ids,
        actuator_bounds=actuator_bounds,
        weather_variability=weather_variability,
    )


main()
\end{lstlisting}

\Needspace{8\baselineskip}

\subsection{\texorpdfstring{\texttt{pages/Ajuda.py}}{pages/Ajuda.py}}

\begin{lstlisting}[style=studioPython,caption={pages/Ajuda.py},label={lst:pages-ajuda-py}]
from html import escape

import streamlit as st
from page_components.ui_fragments import render_hero as render_page_hero
from sidebar_nav import configure_studio_page
from page_styles.theme import inject_studio_theme


PAGE_TITLE = "Ajuda"
INTRODUCTION_TEXT = (
    "BEMS-RL STUDIO permet controlar el cicle complet de l'entrenament "
    "d'un agent RL aplicat a simulacións d'EnergyPlus."
)
INTRODUCTION_HINT = (
    "A continuació trobaràs una explicació detallada de cada funcionalitat "
    "disponible al menú de l'esquerra."
)
RESOURCES = (
    (
        "Documentació de Sinergym",
        "https://ugr-sail.github.io/sinergym/compilation/main/index.html",
    ),
    ("Stable Baselines 3", "https://stable-baselines3.readthedocs.io/"),
    ("EnergyPlus", "https://energyplus.net/"),
    ("Gymnasium", "https://gymnasium.farama.org/"),
)
SUPPORT_CONTACT = {
    "name": "Guillem Torm",
    "email": "guillem@suno.cat",
}
HELP_SECTIONS = (
    (
        "1. Inici",
        """
        - **Objectiu:** Oferir una entrada ràpida al flux complet de treball: crear o revisar entorns, entrenar agents, simular baselines, avaluar models i analitzar resultats.
        - **Quan usar-ho:** Quan vulguis situar-te abans de començar una sessió o saltar cap a una tasca concreta des de la barra lateral.
        - **Consell:** Si no tens clar per on començar, segueix l'ordre natural: **Crear Entorn** o **Mostrar Entorn**, després **Simular Entorn**, **Entrenar Agent**, **Avaluar Agent** i finalment **Resultats**.
        """,
    ),
    (
        "2. Crear Entorn",
        """
        - **Objectiu:** Registrar nous entorns Sinergym a partir d'un edifici, un fitxer climàtic i la configuració d'accions, observacions i actius energètics.
        - **Què pots fer:** Pujar o seleccionar un clima EPW, detectar controladors HVAC i bateria disponibles, ajustar l'espai d'accions i revisar el YAML generat abans de crear l'entorn.
        - **Consell:** Revisa el YAML final i comprova que les variables d'observació i actuació existeixen al model abans de llançar entrenaments llargs.
        """,
    ),
    (
        "3. Mostrar Entorn",
        """
        - **Objectiu:** Inspeccionar els entorns registrats i entendre què veu i què pot controlar l'agent.
        - **Què pots consultar:** Observacions, accions, configuració de reward, clima associat, geometria, zones, actius energètics i configuració crua.
        - **Consell:** Usa aquesta pàgina per validar compatibilitats abans d'entrenar: un entorn discret encaixa millor amb algorismes discrets, mentre que un entorn continu requereix polítiques preparades per accions contínues.
        """,
    ),
    (
        "4. Entrenar Agent",
        """
        - **Objectiu:** Configurar i llançar entrenaments amb Stable Baselines3 sobre entorns Sinergym.
        - **Què pots configurar:** Entorn, algorisme, política, reward, wrappers, paràmetres SB3 avançats, timesteps i configuració climàtica o de cost energètic quan la reward ho necessita.
        - **Consell:** Dona noms curts i descriptius a les execucions, sense espais ni caràcters especials, per facilitar la lectura dels resultats i evitar problemes amb rutes de fitxer.
        """,
    ),
    (
        "5. Resultats",
        """
        - **Objectiu:** Analitzar execucions guardades i convertir els CSV d'entrenament, simulació o avaluació en gràfics i KPIs.
        - **Què pots fer:** Seleccionar una execució, obrir el dashboard integrat, comparar agregacions horàries, diàries, mensuals o sèries reals, revisar confort per zones i descarregar dades.
        - **Consell:** No eliminis els fitxers de monitoratge d'una execució si encara vols visualitzar-ne el dashboard; la pàgina necessita els CSV generats durant el procés.
        """,
    ),
    (
        "6. Simular Entorn",
        """
        - **Objectiu:** Executar un baseline d'EnergyPlus sense control RL per tenir una referència física i energètica de l'entorn.
        - **Què pots fer:** Triar l'entorn, ajustar la configuració del baseline, iniciar o cancel·lar la simulació i comparar-ne els resultats amb execucions guardades.
        - **Consell:** Genera sempre un baseline abans de valorar si un agent RL millora el comportament; així pots comparar consum, confort i desviacions amb una referència clara.
        """,
    ),
    (
        "7. Avaluar Agent",
        """
        - **Objectiu:** Carregar un model entrenat i executar una avaluació controlada sobre un entorn Sinergym.
        - **Què pots fer:** Seleccionar models `.zip`, revisar metadades detectades, ajustar wrappers manuals si el model no inclou tota la configuració i iniciar, cancel·lar o reiniciar l'avaluació.
        - **Consell:** Mantén l'entorn, la normalització i els wrappers coherents amb l'entrenament original; canvis petits poden alterar molt les accions que produeix la política.
        """,
    ),
    (
        "8. Control en Viu",
        """
        - **Objectiu:** Interaccionar pas a pas amb el simulador, ja sigui deixant actuar l'agent o provant accions manuals.
        - **Què pots fer:** Inicialitzar un entorn, carregar un model opcional, generar observacions aleatòries, avançar passos amb l'agent o aplicar accions manuals i veure l'impacte en els KPIs.
        - **Consell:** Fes servir aquest mode per diagnosticar decisions concretes, no per substituir una avaluació completa. Quan acabis, atura o tanca la sessió abans de canviar de pàgina.
        """,
    ),
    (
        "9. Arxius",
        """
        - **Objectiu:** Gestionar fitxers del projecte des de la interfície sense obrir el terminal.
        - **Què pots organitzar:** Entrenaments, scripts, models, fitxers climàtics i configuracions d'entorns. També pots pujar nous fitxers o eliminar artefactes que ja no necessitis.
        - **Consell:** Abans d'esborrar una carpeta d'entrenament, comprova si encara la necessites a **Resultats** o com a referència d'avaluació.
        """,
    ),
    (
        "10. Visor EPW",
        """
        - **Objectiu:** Explorar fitxers climàtics EPW abans d'usar-los en entorns o simulacions.
        - **Què pots veure:** Resum climàtic, metadades, sèries temporals, patrons mensuals, distribucions i dades tabulars del fitxer.
        - **Consell:** Consulta el visor abans de crear un entorn nou; ajuda a detectar climes incoherents, períodes estranys o diferències fortes entre ciutats.
        """,
    ),
    (
        "Consells generals",
        """
        - **Ordre recomanat:** Primer valida l'entorn, després genera un baseline, entrena l'agent, avalua'l i compara els resultats amb el dashboard.
        - **Compatibilitat:** Comprova sempre que l'algorisme, l'espai d'accions i els wrappers coincideixen amb l'entorn i amb el model que vols carregar.
        - **Execucions llargues:** Evita tancar el navegador o canviar de pàgina mentre una simulació, entrenament o avaluació està en curs.
        - **Noms d'execució:** Usa noms llegibles i estables, com `PPO_Radiant_Hivern` o `SAC_OfficeGrid_Test01`.
        - **Comparació justa:** Quan comparis models, intenta mantenir el mateix clima, horitzó temporal, reward i configuració de confort.
        - **Neteja:** Usa el gestor d'arxius amb prudència: eliminar models o monitoratges pot deixar sense dades les pàgines d'avaluació i resultats.
        """,
    ),
)


def inject_help_styles() -> None:
    """Aplica el tema global d'estudi a la pàgina d'ajuda."""
    inject_studio_theme(max_width=1180)


def render_resources_card() -> None:
    """Mostra els recursos addicionals."""

    resource_items = "\n".join(
        f'<li><a href="{escape(url)}" target="_blank">{escape(label)}</a></li>'
        for label, url in RESOURCES
    )
    # Targeta recursos ajuda
    st.markdown(
        f"""
        <section class="resources-card">
            <div class="panel-kicker">Recursos</div>
            <div class="panel-title">Enllaços útils</div>
            <div class="panel-copy">
                Referències externes per aprofundir en les tecnologies que fa servir la plataforma.
            </div>
            <ul>
                {resource_items}
            </ul>
        </section>
        """,
        unsafe_allow_html=True,
    )


def render_support_card() -> None:
    """Mostra les dades de contacte del centre d'ajuda."""

    support_name = SUPPORT_CONTACT["name"]
    support_email = SUPPORT_CONTACT["email"]
    # Targeta contacte ajuda
    st.markdown(
        f"""
        <section class="section-card">
            <div class="panel-kicker">Contacte</div>
            <div class="panel-title">Suport del centre d'ajuda</div>
            <div class="panel-copy">
                Per dubtes o incidències, pots contactar amb
                <strong>{escape(support_name)}</strong> a
                <a href="mailto:{escape(support_email)}">{escape(support_email)}</a>.
            </div>
        </section>
        """,
        unsafe_allow_html=True,
    )


def render_help_page() -> None:
    """Munta la pàgina d'ajuda."""

    configure_studio_page(PAGE_TITLE, layout="wide")
    inject_help_styles()

    render_page_hero("help-hero", "Guia d'ús", PAGE_TITLE, [INTRODUCTION_TEXT, INTRODUCTION_HINT])

    # Separador ajuda
    st.markdown("<div class='studio-spacer-115'></div>", unsafe_allow_html=True)
    # Titol seccions ajuda
    st.markdown('<h2 class="section-title">Funcionalitats i consells</h2>', unsafe_allow_html=True)

    for title, content in HELP_SECTIONS:
        # Acordio seccio ajuda
        with st.expander(title):
            st.markdown(content)

    render_support_card()
    # Separador recursos ajuda
    st.markdown("<div class='studio-spacer-100'></div>", unsafe_allow_html=True)
    render_resources_card()


render_help_page()
\end{lstlisting}

\Needspace{8\baselineskip}

\subsection{\texorpdfstring{\texttt{pages/Avaluar\_Agent.py}}{pages/Avaluar\_Agent.py}}

\begin{lstlisting}[style=studioPython,caption={pages/Avaluar\_Agent.py},label={lst:pages-avaluar-agent-py}]
"""Pàgina per avaluar un agent SB3 entrenat en un entorn Sinergym.

Escaneja models, detecta l'entorn més probable i executa l'avaluació en segon pla
perquè la interfície continuï responent mentre EnergyPlus treballa.
"""

from __future__ import annotations

from html import escape
from pathlib import Path

import streamlit as st
from page_components.ui_fragments import (
    render_hero,
    render_info_panel,
    render_runtime_progress,
    render_section_title,
)
from page_styles.runtime_pages import inject_evaluation_styles
from sidebar_nav import configure_studio_page

from backend.avaluar_agent_backend import (
    ONE_YEAR_STEPS,
    candidate_vecnorm,
    env_id_from_meta_or_name,
    load_model_metadata,
    load_sb3_model_bytes,
    consume_evaluation_runtime_rerun_flag,
    request_evaluation_stop,
    reset_evaluation_runtime,
    scan_model_zips,
    start_evaluation_run,
    sync_evaluation_runtime,
)
PAGE_TITLE = "Avaluar Agent"
PAGE_LAYOUT = "wide"
INTRODUCTION_TEXT = (
    "Carrega un model entrenat, revisa la configuració detectada i executa una "
    "avaluació controlada sobre un entorn Sinergym."
)
def _key(s: str) -> str:
    """Key."""
    return f"eval_{s}"

def render_metadata_grid(cards: list[tuple[str, list[tuple[str, str]]]]) -> None:
    """Mostra un resum en graella amb targetes d'alçada uniforme."""

    cards_markup = ""
    for kicker, items in cards:
        items_markup = "".join(
            (
                '<div class="eval-meta-item">'
                f'<span class="eval-meta-label">{escape(label)}</span>'
                f'<span class="eval-meta-value">{escape(value)}</span>'
                "</div>"
            )
            for label, value in items
        )
        cards_markup += (
            '<section class="eval-meta-card">'
            f'<div class="eval-meta-kicker">{escape(kicker)}</div>'
            f"{items_markup}"
            "</section>"
        )

    # Graella configuracio detectada
    st.markdown(
        f'<div class="eval-meta-grid">{cards_markup}</div>',
        unsafe_allow_html=True,
    )


def render_evaluation_runtime(runtime: dict[str, object]) -> None:
    """Mostra el progrés i el resultat de l'avaluació activa."""

    render_section_title("Estat de l'avaluació")
    render_runtime_progress(runtime, progress_label="Progrés")

    if runtime.get("cancel_requested") and runtime.get("running"):
        # Avís cancel·lació avaluació
        st.warning("S'ha sol·licitat l'aturada. L'avaluació es tancarà al proper punt segur.")

    if runtime.get("error"):
        # Error runtime avaluació
        st.error(str(runtime.get("error")))
        if runtime.get("traceback"):
            with st.expander("Traceback intern", expanded=False):
                st.code(str(runtime.get("traceback")), language="text")
        return

    result = runtime.get("result")
    if result is None or runtime.get("running"):
        return

    for warn_msg in getattr(result, "warnings", ()):
        # Avís resultat avaluacio
        st.warning(warn_msg)

    if bool(getattr(result, "cancelled", False)):
        # Avís avaluacio cancel·lada
        st.warning(
            f"Avaluació cancel·lada per l'usuari després de {float(getattr(result, 'elapsed_seconds', 0.0)):.2f} s."
        )
    else:
        # Missatge avaluacio finalitzada
        st.success(
            f"Avaluació finalitzada en {float(getattr(result, 'elapsed_seconds', 0.0)):.2f} s."
        )


def render_evaluation_runtime_frame() -> None:
    """Manté sincronitzat el bloc d'estat amb el runtime en segon pla."""

    runtime = sync_evaluation_runtime()
    if runtime.get("running") or runtime.get("completed") or runtime.get("error") or runtime.get("result"):
        st.divider()
        render_evaluation_runtime(runtime)

    if consume_evaluation_runtime_rerun_flag(runtime):
        st.rerun()


render_evaluation_runtime_fragment = st.fragment(run_every="1s")(render_evaluation_runtime_frame)


configure_studio_page(PAGE_TITLE, layout=PAGE_LAYOUT)
inject_evaluation_styles()

st.session_state.setdefault(_key("dirs"), ",".join([str(Path.cwd() / "models"), str(Path.cwd())]))
st.session_state.setdefault(_key("selected_idx"), 0)
st.session_state.setdefault(_key("deterministic"), True)
st.session_state.setdefault(_key("use_vecnorm"), True)
st.session_state.setdefault(_key("steps"), ONE_YEAR_STEPS)

render_hero("eval-hero", "Validació del model", PAGE_TITLE, INTRODUCTION_TEXT)

# Panells resum avaluacio
summary_col, scale_col = st.columns(2, gap="large")
with summary_col:
    render_info_panel(
        "eval-panel",
        "Avaluació controlada",
        "Objectiu",
        "Carrega un model SB3, detecta l'entorn probable i executa una passada completa amb seguiment de progrés.",
    )
    with scale_col:
        render_info_panel(
        "eval-panel",
        f"{ONE_YEAR_STEPS} passos per defecte",
        "Escala",
            "La configuració inicial parteix d'una avaluació equivalent a un any aproximat de simulació.",
        )

    # Separador localitzacio models
    st.markdown("<div class='studio-spacer-115'></div>", unsafe_allow_html=True)
render_section_title("Localització de models")
dirs_text = st.session_state[_key("dirs")]
df_models = scan_model_zips(dirs_text)
runtime = sync_evaluation_runtime()
controls_disabled = bool(runtime.get("running"))

if df_models.empty:
    # Avís models no trobats
    st.info("No s'han trobat models .zip als directoris indicats.")
    st.stop()

options = list(range(len(df_models)))
# Selector model entrenat
idx = st.selectbox(
    "Selecciona el model entrenat (.zip)",
    options=options,
    format_func=lambda i: str(df_models.iloc[i]["stem"]),
    index=min(st.session_state[_key("selected_idx")], len(df_models) - 1),
    disabled=controls_disabled,
)
st.session_state[_key("selected_idx")] = idx

row = df_models.iloc[idx]
model_path = Path(row["path"])
zip_stem = row["stem"]

with open(model_path, "rb") as fh:
    model_tmp, meta = load_sb3_model_bytes(fh.read(), device="cpu")

algo_name = type(model_tmp).__name__
policy_name = type(getattr(model_tmp, "policy", object())).__name__

model_metadata = load_model_metadata(model_path)
env_id_guess = model_metadata.get("env_id") or env_id_from_meta_or_name(meta, zip_stem)
vecnorm_path = candidate_vecnorm(model_path)
wrapper_configs = model_metadata.get("wrapper_configs", [])
wrapper_names = [
    str(wrapper.get("name"))
    for wrapper in wrapper_configs
    if isinstance(wrapper, dict) and wrapper.get("name")
]
wrappers_label = ", ".join(wrapper_names) if wrapper_names else "Cap (metadades no trobades)"

render_section_title("Configuració detectada")
render_metadata_grid(
    [
        (
            "Model",
            [
                ("Agent", zip_stem),
                ("Algorisme", algo_name),
            ],
        ),
        (
            "Configuració",
            [
                ("Política", policy_name),
                ("VecNormalize", vecnorm_path.name if vecnorm_path else "No"),
            ],
        ),
        (
            "Context",
            [
                ("Entorn detectat", env_id_guess or "No detectat"),
                ("Wrappers d'entrenament", wrappers_label),
                ("Observacions esperades", str(getattr(model_tmp.observation_space, "shape", ("?",))[0])),
            ],
        ),
    ]
)

st.divider()

render_section_title("Paràmetres d'avaluació")
env_id = env_id_guess or ""
# Controls opcions avaluacio
col_b, col_c = st.columns(2)
with col_b:
    # Checkbox mode deterministic
    st.checkbox(
        "Deterministic",
        key=_key("deterministic"),
        value=st.session_state[_key("deterministic")],
        disabled=controls_disabled,
    )
with col_c:
    # Checkbox VecNormalize
    st.checkbox(
        "Usar VecNormalize si hi és",
        key=_key("use_vecnorm"),
        value=bool(vecnorm_path),
        disabled=controls_disabled,
    )

# Input limit passos
steps_limit = st.number_input(
    "Límit de passos (per defecte 35040)",
    min_value=1,
    max_value=5_000_000,
    value=st.session_state[_key("steps")],
    step=35040,
    disabled=controls_disabled,
)

if not wrapper_configs:
    # Avís wrappers no trobats
    st.info(
        "No s'han trobat metadades de wrappers per aquest model. "
        "L'avaluació inferirà automàticament l'espai d'accions a partir del model i de l'entorn seleccionat."
    )

st.divider()

render_section_title("Control")
# Botons control avaluacio
col_run, col_cancel, col_reset = st.columns(3, gap="medium")
with col_run:
    # Botó iniciar avaluacio
    start_requested = st.button(
        "Iniciar avaluació",
        type="primary",
        width="stretch",
        disabled=controls_disabled or not env_id,
    )
with col_cancel:
    # Botó cancel·lar avaluacio
    cancel_requested = st.button(
        "Cancel·lar",
        width="stretch",
        disabled=not bool(runtime.get("running")),
    )
with col_reset:
    # Botó netejar estat
    reset_requested = st.button(
        "Netejar estat",
        width="stretch",
        disabled=not bool(runtime.get("completed") or runtime.get("error")),
    )

if cancel_requested and request_evaluation_stop(runtime):
    st.rerun()

if reset_requested:
    reset_evaluation_runtime()
    st.rerun()

if start_requested:
    start_result = start_evaluation_run(
        {
            "model_path": model_path,
            "env_id": env_id,
            "use_vecnorm": bool(st.session_state[_key("use_vecnorm")]),
            "vecnorm_path": vecnorm_path,
            "steps_target": int(steps_limit),
            "deterministic": bool(st.session_state[_key("deterministic")]),
            "wrapper_configs": wrapper_configs,
        }
    )
    for error in start_result.get("errors") or []:
        # Error inici avaluacio
        st.error(str(error))
    if start_result.get("started"):
        st.rerun()

render_evaluation_runtime_fragment()
\end{lstlisting}

\Needspace{8\baselineskip}

\subsection{\texorpdfstring{\texttt{pages/Entrenar\_Agent.py}}{pages/Entrenar\_Agent.py}}

\begin{lstlisting}[style=studioPython,caption={pages/Entrenar\_Agent.py},label={lst:pages-entrenar-agent-py}]
"""Pàgina d'entrenament de l'agent per a l'aplicació BEMS-RL-STUDIO.

Orquestra el flux complet d'entrenament: selecció de l'entorn, configuració de la
reward, wrappers, hiperparàmetres de l'agent i bucle incremental amb gràfics en directe.
"""

from __future__ import annotations

import streamlit as st

from sidebar_nav import configure_studio_page

from backend.entrenar_agent_constants import (
    POLICIES,
    REWARD_CLASSES,
    TRAINING_RESULT_KEY,
    TRAINING_RUNTIME_KEY,
)
from backend.entrenar_agent_env import get_env_metadata, list_registered_envs
from backend.entrenar_agent_rewards import assemble_training_payload
from backend.entrenar_agent_runtime import (
    clear_training_runtime,
    create_training_runtime,
    learn_training_chunk,
    save_training_artifacts,
)
from backend.sb3_utils import ALGOS
from backend.entrenar_agent_session import (
    TRAINING_WORKSPACE_KEY,
    normalize_training_range_state,
    selectbox_state_kwargs,
    training_form_key,
)
from backend.entrenar_agent_charts import (
    _get_workspace_from_runtime,
    render_live_training_charts,
)
from page_components.training_agent import render_agent_params_section
from page_components.training_rewards import render_reward_kwargs_section
from page_components.training_shared import (
    PAGE_LAYOUT,
    PAGE_TITLE,
    inject_training_styles,
    render_saved_training_library,
    render_training_hero,
    render_training_section,
)
from page_components.training_wrappers import render_wrappers_section


def train_tab() -> None:
    """Punt d'entrada a la pàgina d'entrenament.

    Orquestra totes les seccions UI en ordre:
    1. Inicialització de l'estat de la sessió.
    2. Entorn, algorisme i selecció de polítiques.
    3. Recompensa kwargs (pesos, rangs estacionals, comoditat).
    4. Selecció i configuració de wrappers (observació, acció, logs).
    5. Paràmetres de l'agent (taxa d'aprenentatge, n_steps, total_timesteps).
    6. Controls d'execució (inici/aturada).
    7. Cicle d'entrenament incremental amb gràfics en directe.
    """
    configure_studio_page(PAGE_TITLE, layout=PAGE_LAYOUT)
    inject_training_styles()

    # Inicialització de l'estat de la sessió
    if "is_training" not in st.session_state:
        st.session_state.is_training = False
    if "stop_training" not in st.session_state:
        st.session_state.stop_training = False

    envs = list_registered_envs()
    normalize_training_range_state()
    render_training_hero()
    render_saved_training_library(envs)
    st.divider()

    # Secció 1: configuració de l'entorn
    render_training_section(
        "Configuracio d'entorn",
        "Seleccio base",
        "Defineix l'entorn, l'algorisme i la politica abans de construir la recompensa i els wrappers.",
    )
    # Selectors entorn i agent
    env_col1, env_col2 = st.columns(2)
    with env_col1:
        # Selector entorn
        env_id = st.selectbox(
            "Entorn",
            envs,
            **selectbox_state_kwargs("env_id", envs),
        )
        algo_options = list(ALGOS.keys())
        # Selector algorisme
        algo_name = st.selectbox(
            "Algorisme SB3",
            algo_options,
            **selectbox_state_kwargs("algo_name", algo_options),
        )
    with env_col2:
        # Selector politica
        policy_name = st.selectbox(
            "Politica SB3",
            POLICIES[algo_name],
            **selectbox_state_kwargs("policy_name", POLICIES[algo_name]),
        )
        reward_options = list(REWARD_CLASSES.keys())
        # Selector recompensa
        reward_name = st.selectbox(
            "Recompensa",
            reward_options,
            **selectbox_state_kwargs("reward_name", reward_options),
        )

    env_meta = get_env_metadata(env_id)
    spec                             = env_meta["spec"]
    variables                        = env_meta["variables"]
    observation_variables            = env_meta["observation_variables"]
    action_variables                 = env_meta["action_variables"]
    action_space                     = env_meta["action_space"]
    candidate_grid_energy_vars       = env_meta["candidate_grid_energy_vars"]
    candidate_battery_charge_vars    = env_meta["candidate_battery_charge_vars"]
    candidate_battery_discharge_vars = env_meta["candidate_battery_discharge_vars"]
    candidate_battery_loss_vars      = env_meta["candidate_battery_loss_vars"]

    # Secció 2: recompensa kwargs
    st.divider()
    render_training_section(
        "Reward kwargs",
        "Confort i energia",
        "Ajusta pesos, rangs estacionals i criteris de confort per construir la reward final.",
    )
    reward_kwargs_values = render_reward_kwargs_section(
        reward_name=reward_name,
        observation_variables=observation_variables,
        candidate_grid_energy_vars=candidate_grid_energy_vars,
        candidate_battery_charge_vars=candidate_battery_charge_vars,
        candidate_battery_discharge_vars=candidate_battery_discharge_vars,
        candidate_battery_loss_vars=candidate_battery_loss_vars,
    )

    # Secció 3: wrappers i logs
    st.divider()
    render_training_section(
        "Wrappers i logs",
        "Observacio i control",
        "Activa nomes els wrappers que necessites per enriquir l'observacio o restringir l'espai d'accions.",
    )
    wrapper_values = render_wrappers_section(
        env_meta=env_meta,
        observation_variables=observation_variables,
        action_variables=action_variables,
        action_space=action_space,
        context_variables=env_meta["context_variables"],
        candidate_temp_vars=env_meta["candidate_temp_vars"],
        candidate_setpoint_vars=env_meta["candidate_setpoint_vars"],
    )

    # Apartat 4: paràmetres de l'agent
    st.divider()
    render_training_section(
        "Parametres de l'agent",
        "Aprenentatge",
        "Configura el ritme d'aprenentatge, la mida dels blocs i el nombre total de timesteps.",
    )
    agent_params = render_agent_params_section(algo_name=algo_name)

    # Secció 5: controls d'execució
    st.divider()
    render_training_section(
        "Execucio",
        "Control",
        "Llança l'entrenament o atura'l de forma segura mantenint els artefactes guardats.",
    )
    # Botons control entrenament
    controls_col1, controls_col2 = st.columns(2)
    # Botó entrenar agent
    start_clicked = controls_col1.button(
        "Entrenar agent",
        disabled=st.session_state.is_training,
        width="stretch",
    )
    # Botó aturar entrenament
    stop_clicked = controls_col2.button(
        "Aturar entrenament",
        disabled=not st.session_state.is_training,
        width="stretch",
    )

    if start_clicked:
        clear_training_runtime()
        st.session_state.pop(TRAINING_RESULT_KEY, None)
        st.session_state.pop(TRAINING_WORKSPACE_KEY, None)
        st.session_state.is_training = True
        st.session_state.stop_training = False

    if stop_clicked and st.session_state.is_training:
        st.session_state.stop_training = True

    if not st.session_state.is_training:
        return

    # A partir d'aquí Streamlit reexecuta la pàgina moltes vegades. Si ja hi ha runtime,
    # no tornem a muntar payload ni entorn: només continuem el model que és a memòria.
    runtime = st.session_state.get(TRAINING_RUNTIME_KEY)
    runtime_ui_state = runtime.get("ui_state", {}) if runtime is not None else {}
    training_payload = (
        None
        if runtime is not None
        else assemble_training_payload({
            "spec":                              spec,
            "env_id":                            env_id,
            "algo_name":                         algo_name,
            "policy_name":                       policy_name,
            "reward_name":                       reward_name,
            "is_continuous":                     env_meta["is_continuous"],
            "is_discrete":                       env_meta["is_discrete"],
            "action_space_type":                 type(action_space).__name__,
            "building_name":                     env_meta["building_name"],
            "variables":                         variables,
            **reward_kwargs_values,
            **wrapper_values,
            **agent_params,
        })
    )

    if training_payload is not None:
        dropped_reward_kwargs = training_payload["dropped_reward_kwargs"]
        if dropped_reward_kwargs:
            # Avís reward kwargs ignorats
            st.warning(
                "La classe de reward carregada no admet alguns parametres i s'han omes. "
                f"Parametres ignorats: {', '.join(dropped_reward_kwargs)}."
            )

        dropped_algo_kwargs = training_payload.get("dropped_algo_kwargs", [])
        if dropped_algo_kwargs:
            # Avís SB3 kwargs ignorats
            st.info(
                "L'algorisme seleccionat no fa servir alguns parametres SB3 avancats: "
                f"{', '.join(dropped_algo_kwargs)}."
            )

        validation_errors = training_payload.get("validation_errors", [])
        if validation_errors:
            clear_training_runtime()
            st.session_state.is_training = False
            st.session_state.stop_training = False
            for message in validation_errors:
                # Error validacio entrenament
                st.error(message)
            return

        training_config = training_payload["training_config"]
    else:
        training_config = runtime["config"]

    if runtime is not None:
        runtime_config = runtime.get("config", {})
        linear_cost_rewards = {"LinearReward", "EnergyCostLinearReward"}
        # Si una run antiga porta EnergyCostWrapper amb una configuració que ja no quadra,
        # és més segur recrear el runtime que entrenar amb una reward recalculada a mitges.
        stale_energy_wrapper = any(
            wrapper.get("name") == "EnergyCostWrapper"
            and (
                "recalculate_reward" not in wrapper.get("params", {})
                or (
                    runtime_config.get("reward_name") not in linear_cost_rewards
                    and wrapper.get("params", {}).get("recalculate_reward", True)
                )
            )
            for wrapper in runtime_config.get("wrapper_configs", [])
        )
        # El runtime creat pel botó queda congelat entre reruns; els widgets poden canviar
        # visualment, però el model en curs ha de continuar amb la configuració inicial.
        if stale_energy_wrapper:
            clear_training_runtime()
            st.session_state.pop(TRAINING_WORKSPACE_KEY, None)
            runtime = None
    if runtime is None:
        runtime = create_training_runtime(
            training_config,
            ui_state=training_payload["ui_state"] if training_payload is not None else runtime_ui_state,
        )
        st.session_state[TRAINING_RUNTIME_KEY] = runtime

    if TRAINING_WORKSPACE_KEY not in st.session_state:
        workspace = _get_workspace_from_runtime(runtime)
        if workspace:
            st.session_state[TRAINING_WORKSPACE_KEY] = workspace

    active_config = runtime["config"]

    render_training_section(
        "Estat de l'entrenament",
        "Seguiment",
        "Visualitza el progres actual, l'entorn actiu i l'estat del proces abans de guardar el model.",
    )
    st.write(f"Entorn seleccionat: `{active_config['env_id']}`")
    st.write(f"Algorisme: `{active_config['algo_name']}`")
    st.write(f"Recompensa: `{active_config['reward_name']}`")

    # Barra progrés entrenament
    progress_bar    = st.progress(0.0)
    # Text estat entrenament
    training_status = st.empty()

    current_timesteps = int(runtime["model"].num_timesteps)
    total_timesteps   = int(active_config["timesteps_per_year"])
    frac              = min(current_timesteps / max(1, total_timesteps), 1.0)
    progress_bar.progress(frac)

    render_training_section(
        "Evolucio de l'entrenament",
        "Temps real",
        "Metriques per episodi actualitzades automaticament en completar-se cada episodi.",
    )
    render_live_training_charts()

    try:
        if st.session_state.stop_training:
            training_status.warning("Aturant l'entrenament i guardant l'estat actual...")
            was_stopped = True
        elif current_timesteps >= total_timesteps:
            training_status.success("Entrenament completat. Guardant el model...")
            was_stopped = False
        else:
            training_status.info(
                f"Entrenament en curs: {current_timesteps}/{total_timesteps} timesteps ({frac * 100:.1f}%)."
            )
            remaining_timesteps = max(1, total_timesteps - current_timesteps)
            next_chunk = min(max(1, int(runtime["chunk_timesteps"])), remaining_timesteps)
            # Entrenem en blocs petits perquè la UI respongui, actualitzi gràfics i pugui
            # aturar-se sense esperar tot un any de simulació.
            learn_training_chunk(runtime["model"], next_chunk)
            st.session_state[TRAINING_RUNTIME_KEY] = runtime
            st.rerun()

        st.session_state[TRAINING_RESULT_KEY] = save_training_artifacts(runtime)
        clear_training_runtime()
        st.session_state.is_training = False
        st.session_state.stop_training = False
        st.rerun()
    except Exception:
        clear_training_runtime()
        st.session_state.is_training = False
        st.session_state.stop_training = False
        raise


# Streamlit executa aquesta funció quan carrega la pàgina.

train_tab()
\end{lstlisting}

\Needspace{8\baselineskip}

\subsection{\texorpdfstring{\texttt{pages/Gestionar\_Arxius.py}}{pages/Gestionar\_Arxius.py}}

\begin{lstlisting}[style=studioPython,caption={pages/Gestionar\_Arxius.py},label={lst:pages-gestionar-arxius-py}]
from __future__ import annotations

import time
from datetime import datetime
from html import escape
from io import BytesIO
from pathlib import Path
from zipfile import ZIP_DEFLATED, ZipFile

import streamlit as st
from page_components.ui_fragments import render_hero
from page_styles.gestionar_arxius import inject_file_manager_styles
from sidebar_nav import configure_studio_page

from backend.mapa_backend import render_multipoint_map
from backend.gestionar_arxius_backend import (
    DATA_BTC_PATH,
    DATA_WEA_PATH,
    ROOT_PATH,
    delete_item,
    filter_envs,
    filter_models,
    filter_trainings,
    filter_weather,
    list_explorer_items,
    load_weather_map_rows,
)


PAGE_TITLE = "Gestor de Fitxers Avançat"
PAGE_LAYOUT = "wide"
INTRODUCTION_TEXT = "Gestiona els fitxers del projecte organitzats per categories."
EXPLORER_COLUMN_LAYOUT = [0.8, 1.4, 5.9, 1.4, 2.5]
TRAINING_ARTIFACTS_PATH = ROOT_PATH / "trainings"


def render_weather_view_selector() -> str:
    """Mostra el selector de vista del bloc de climes sense ràdios."""

    options = ["Explorador", "Mapa"]
    current_value = st.session_state.get("weather_submenu", "Explorador")

    # Selector vista climes
    selected_value = st.segmented_control(
        "Vista de climes",
        options,
        default=current_value,
        key="weather_submenu_segmented",
    )
    if selected_value:
        st.session_state["weather_submenu"] = selected_value
    return st.session_state.get("weather_submenu", current_value)


def build_selection_archive(selected_paths: list[str]) -> bytes:
    """Construeix un zip en memòria amb els elements seleccionats."""

    archive_buffer = BytesIO()
    normalized_paths = sorted({str(Path(path_str)) for path_str in selected_paths})

    with ZipFile(archive_buffer, "w", ZIP_DEFLATED) as archive:
        for path_str in normalized_paths:
            path = Path(path_str)
            if not path.exists():
                continue

            if path.is_dir():
                archive.writestr(f"{path.name}/", "")
                for nested_path in sorted(path.rglob("*")):
                    archive_name = nested_path.relative_to(path.parent).as_posix()
                    if nested_path.is_dir():
                        try:
                            next(nested_path.iterdir())
                        except StopIteration:
                            archive.writestr(f"{archive_name}/", "")
                        except OSError:
                            pass
                        continue

                    archive.write(nested_path, archive_name)
                continue

            archive.write(path, path.name)

    archive_buffer.seek(0)
    return archive_buffer.getvalue()


def build_download_name(key_prefix: str) -> str:
    """Construeix un nom de fitxer estable per a la descàrrega."""

    timestamp = datetime.now().strftime("%Y%m%d_%H%M")
    return f"{key_prefix}_seleccio_{timestamp}.zip"


def clear_prepared_archive_if_stale(key_prefix: str, selection_signature: tuple[str, ...]) -> None:
    """Neteja l'arxiu preparat si la selecció actual ja no coincideix."""

    archive_key = f"{key_prefix}_download_archive"
    archive_name_key = f"{key_prefix}_download_name"
    archive_signature_key = f"{key_prefix}_download_signature"

    prepared_signature = st.session_state.get(archive_signature_key)
    if prepared_signature == selection_signature:
        return

    st.session_state.pop(archive_key, None)
    st.session_state.pop(archive_name_key, None)
    st.session_state.pop(archive_signature_key, None)


def delete_items(items: list[str]) -> None:
    """Esborra els elements seleccionats mostrant progrés."""

    progress_bar = st.progress(0)
    status_text = st.empty()

    for index, path_str in enumerate(items):
        path = Path(path_str)
        status_text.text(f"Eliminant: {path.name}...")
        error_message = delete_item(path_str)
        if error_message:
            # Error eliminacio fitxer
            st.error(f"Error eliminant {path.name}: {error_message}")
        progress_bar.progress((index + 1) / len(items))

    status_text.text("Fet.")
    time.sleep(1)
    st.rerun()


def handle_uploaded_files(current_path: Path, uploaded_files) -> None:
    """Desa els fitxers pujats a la ruta actual de l'explorador."""

    for uploaded_file in uploaded_files:
        destination = current_path / Path(uploaded_file.name).name
        try:
            with open(destination, "wb") as destination_handle:
                destination_handle.write(uploaded_file.getbuffer())
            st.toast(f"Fitxer carregat correctament: {destination.name}", icon="\u2705")
        except Exception as exc:
            # Error upload fitxer
            st.error(f"Error carregant {uploaded_file.name}: {exc}")

    time.sleep(1)
    st.rerun()


def navigate_file_explorer_up(path_key: str, current_path: Path, root_path: Path) -> None:
    """Mou l'explorador de fitxers un directori cap amunt sense sortir de l'arrel."""
    parent = current_path.parent
    if parent == root_path or root_path in parent.parents:
        st.session_state[path_key] = str(parent)
    else:
        st.session_state[path_key] = str(root_path)


def navigate_file_explorer_into(path_key: str, folder_path: str) -> None:
    """Mou l'explorador de fitxers cap a una carpeta filla."""
    st.session_state[path_key] = str(Path(folder_path).resolve())


def render_explorer(
    key_prefix: str,
    root_path: Path,
    filter_func=None,
    allow_nav_up: bool = True,
    allow_upload: bool = False,
    upload_label: str | None = None,
    upload_types: list[str] | None = None,
    upload_help: str | None = None,
) -> None:
    """Mostra l'explorador de fitxers per a cada categoria."""

    root_path = root_path.resolve()
    path_key = f"{key_prefix}_path"
    if path_key not in st.session_state:
        st.session_state[path_key] = str(root_path)

    current_path = Path(st.session_state[path_key]).resolve()

    # El path surt de session_state, així que el validem sempre contra l'arrel de la categoria.
    # D'aquesta manera un valor vell o manipulat no pot fer navegar fora de la carpeta prevista.
    try:
        current_path.relative_to(root_path)
    except ValueError:
        st.session_state[path_key] = str(root_path)
        current_path = root_path

    # Navegacio explorador
    nav_columns = st.columns([1.1, 5.4], gap="small")
    with nav_columns[0]:
        disabled = (current_path == root_path) or not allow_nav_up
        # Botó upload carpeta
        if st.button("Amunt", key=f"{key_prefix}_up", disabled=disabled, type="primary", width="stretch"):
            navigate_file_explorer_up(path_key, current_path, root_path)
            # Canviar carpeta modifica l'estat; forcem rerun perquè la llista es repinti de seguida.
            st.rerun()

    try:
        relative_path = current_path.relative_to(root_path)
        display_path = f"/{relative_path.as_posix()}" if str(relative_path) != "." else "/"
    except ValueError:
        display_path = str(current_path)

    with nav_columns[1]:
        # Ruta carpeta actual
        st.markdown(f"<div class='file-explorer-path'>{escape(display_path)}</div>", unsafe_allow_html=True)

    if allow_upload:
        # Upload fitxers
        uploaded_files = st.file_uploader(
            upload_label or f"Pujar fitxers a {display_path}",
            accept_multiple_files=True,
            type=upload_types,
            help=upload_help,
            key=f"{key_prefix}_upload",
        )
        if uploaded_files:
            # Les pujades es desen a la carpeta actual, no a l'arrel global del projecte.
            handle_uploaded_files(current_path, uploaded_files)

    all_items, error_message = list_explorer_items(current_path, root_path, filter_func)
    if error_message:
        # Error lectura directori
        st.error(f"Error llegint el directori: {error_message}")
        all_items = ()

    if not all_items:
        # Avís carpeta buida
        st.info("Aquesta carpeta és buida.")
        return

    # Capçalera taula fitxers
    header_columns = st.columns(EXPLORER_COLUMN_LAYOUT, gap="small")
    headers = ("Sel.", "Tipus", "Nom", "Mida", "Data")
    for column, header in zip(header_columns, headers):
        column.markdown(f"<div class='file-explorer-header'>{header}</div>", unsafe_allow_html=True)

    # Separador taula fitxers
    st.markdown("---")

    selected_items: list[str] = []
    for index, item in enumerate(all_items):
        # Fila fitxer explorador
        row_columns = st.columns(EXPLORER_COLUMN_LAYOUT, gap="small")
        checkbox_key = f"{key_prefix}_chk_{item.path}"
        # Checkbox seleccionar fitxer
        is_selected = row_columns[0].checkbox(
            f"Seleccionar {item.name}",
            key=checkbox_key,
            label_visibility="collapsed",
        )
        if is_selected:
            selected_items.append(item.path)

        item_type = "Carpeta" if item.is_dir else "Fitxer"
        item_size = "—" if item.is_dir else item.size
        item_icon = "\U0001f4c1" if item.is_dir else "\U0001f4c4"

        # Cel·la tipus fitxer
        row_columns[1].markdown(f"<div class='file-explorer-muted'>{item_type}</div>", unsafe_allow_html=True)
        name_columns = row_columns[2].columns([0.6, 7.8], gap="small")
        # Cel·la icona fitxer
        name_columns[0].markdown(
            f"<div class='file-explorer-icon-cell'><span class='file-explorer-icon'>{item_icon}</span></div>",
            unsafe_allow_html=True,
        )

        if item.is_dir:
            # Botó obrir carpeta
            if name_columns[1].button(
                item.name,
                key=f"{key_prefix}_btn_{item.path}",
                type="secondary",
                width="stretch",
            ):
                navigate_file_explorer_into(path_key, item.path)
                st.rerun()
        else:
            safe_name = escape(item.name)
            # Cel·la nom fitxer
            name_columns[1].markdown(
                (
                    "<div class='file-explorer-cell file-explorer-name'>"
                    f"<span class='file-explorer-label'>{safe_name}</span>"
                    "</div>"
                ),
                unsafe_allow_html=True,
            )

        # Cel·la mida fitxer
        row_columns[3].markdown(f"<div class='file-explorer-muted'>{escape(item_size)}</div>", unsafe_allow_html=True)
        # Cel·la data fitxer
        row_columns[4].markdown(f"<div class='file-explorer-muted'>{escape(item.mtime)}</div>", unsafe_allow_html=True)
        if index < len(all_items) - 1:
            # Separador fila fitxer
            st.markdown("<div class='file-explorer-row-spacer'></div>", unsafe_allow_html=True)

    # Separador seleccio fitxers
    st.markdown("---")

    if not selected_items:
        return

    # Controls seleccio fitxers
    selection_columns = st.columns([2.4, 1.5, 1.2], gap="small")
    selection_columns[0].markdown(
        f"<div class='file-explorer-selection'>{len(selected_items)} elements seleccionats.</div>",
        unsafe_allow_html=True,
    )
    selection_signature = tuple(sorted({str(Path(path_str)) for path_str in selected_items}))
    # Si canvia la selecció, descartem el zip preparat anterior. Evita baixar un arxiu
    # que ja no correspon al que es veu marcat a pantalla.
    clear_prepared_archive_if_stale(key_prefix, selection_signature)

    archive_key = f"{key_prefix}_download_archive"
    archive_name_key = f"{key_prefix}_download_name"
    archive_signature_key = f"{key_prefix}_download_signature"
    archive_ready = (
        st.session_state.get(archive_signature_key) == selection_signature
        and st.session_state.get(archive_key) is not None
    )

    with selection_columns[1]:
        download_slot = st.empty()
        # Botó preparar download
        if not archive_ready and download_slot.button(
            "Descarregar selecció",
            key=f"{key_prefix}_prepare_download",
            type="primary",
            width="stretch",
        ):
            try:
                with st.spinner("Preparant la descàrrega..."):
                    st.session_state[archive_key] = build_selection_archive(selected_items)
                st.session_state[archive_name_key] = build_download_name(key_prefix)
                st.session_state[archive_signature_key] = selection_signature
                archive_ready = True
            except Exception as exc:
                # Error preparar download
                st.error(f"No s'ha pogut preparar la descàrrega: {exc}")
                archive_ready = False

        if archive_ready:
            # Botó download zip
            download_slot.download_button(
                "Descarregar selecció",
                data=st.session_state[archive_key],
                file_name=st.session_state[archive_name_key],
                mime="application/zip",
                key=f"{key_prefix}_download",
                type="primary",
                width="stretch",
            )

    with selection_columns[2]:
        # Botó esborrar seleccio
        if st.button("Esborrar", key=f"{key_prefix}_del", type="primary", width="stretch"):
            delete_items(selected_items)


configure_studio_page(PAGE_TITLE, layout=PAGE_LAYOUT)
inject_file_manager_styles()

render_hero("files-hero", "Exploració i manteniment", PAGE_TITLE, INTRODUCTION_TEXT)

# Separador gestor fitxers
st.markdown("<div class='studio-spacer-045'></div>", unsafe_allow_html=True)
# Pestanyes gestor fitxers
tab_train, tab_scripts, tab_models, tab_weather, tab_envs = st.tabs(
    ["Entrenaments", "Scripts", "Models", "Climes", "Entorns"]
)

with tab_train:
    render_explorer("train", ROOT_PATH, filter_trainings)

with tab_scripts:
    if TRAINING_ARTIFACTS_PATH.exists():
        render_explorer("training_scripts", TRAINING_ARTIFACTS_PATH)
    else:
        # Avís carpeta trainings absent
        st.info("Encara no hi ha cap carpeta `trainings/`. Es creara quan es guardi el primer entrenament.")

with tab_models:
    models_root = TRAINING_ARTIFACTS_PATH if TRAINING_ARTIFACTS_PATH.exists() else ROOT_PATH
    render_explorer("models", models_root, filter_models)

with tab_weather:
    weather_root = DATA_WEA_PATH if DATA_WEA_PATH.exists() else ROOT_PATH
    # Selector vista climes
    weather_view = render_weather_view_selector()

    if weather_view == "Explorador":
        render_explorer(
            "weather",
            weather_root,
            filter_weather,
            allow_upload=True,
            upload_label="Pujar fitxers climàtics (.epw)",
            upload_types=["epw"],
            upload_help="Afegeix nous fitxers meteorològics EPW al catàleg de climes.",
        )
    else:
        # Text mapa climes
        st.caption("Mapa amb tots els fitxers `.epw` disponibles.")
        with st.spinner("Generant mapa de climes..."):
            weather_df = load_weather_map_rows(str(weather_root))

        if weather_df.empty:
            # Avís mapa climes buit
            st.warning("No s'han trobat fitxers `.epw` per mostrar al mapa.")
        else:
            render_multipoint_map(
                weather_df,
                zoom=2,
                color=(255, 0, 0, 200),
                radius=55000,
                tooltip_html=(
                    "<b>{file_name}</b><br/>{city}, {country}"
                    "<br/>Clima: {climate}<br/>Avg tmp: {avg_temp} C"
                ),
            )
            # Taula climes mapa
            st.dataframe(
                weather_df[["file_name", "city", "country", "climate", "avg_temp"]],
                hide_index=True,
                width="stretch",
            )

with tab_envs:
    env_root = DATA_BTC_PATH if DATA_BTC_PATH.exists() else ROOT_PATH
    render_explorer("envs", env_root, filter_envs)
\end{lstlisting}

\Needspace{8\baselineskip}

\subsection{\texorpdfstring{\texttt{pages/Interaccionar\_amb\_l'Agent.py}}{pages/Interaccionar\_amb\_l'Agent.py}}

\begin{lstlisting}[style=studioPython,caption={pages/Interaccionar\_amb\_l'Agent.py},label={lst:pages-interaccionar-amb-l-agent-py}]
"""Pàgina de control en viu d'un entorn Sinergym.

Permet provar un model entrenat o aplicar accions manuals pas a pas.
"""

from html import escape
from pathlib import Path

import gymnasium as gym
import numpy as np
import pandas as pd
import streamlit as st
from page_components.ui_fragments import (
    build_metric_item,
    render_hero,
    render_kicker_section,
    render_metric_row,
)
from page_styles.interaccionar_agent import inject_interaction_styles
from sidebar_nav import configure_studio_page

from backend.interaccionar_agent_backend import (
    env_id_from_meta_or_name,
    extract_display_observation,
    extract_policy_test_defaults,
    find_matching_column,
    coerce_observation_values,
    get_model_observation_size,
    get_runtime_observation_variables,
    get_wrapper_variables,
    infer_unit,
    initialize_runtime,
    load_model_training_metadata,
    load_sb3_model_bytes,
    mode_requires_model,
    normalize_policy_observation,
    prepare_action_display,
    randomize_observation_values,
    run_environment_steps,
    scan_model_zips,
)


PAGE_TITLE = "Control en Viu"
PAGE_LAYOUT = "wide"
INTRODUCTION_TEXT = (
    "Configura un entorn Sinergym, carrega un model si cal i controla "
    "la simulació pas a pas amb l'agent o amb accions manuals."
)


def render_interaction_hero() -> None:
    """Mostra la introducció de la pàgina principal a la UI de Streamlit."""

    render_hero("interaction-hero", "Simulació i control", PAGE_TITLE, INTRODUCTION_TEXT)


def render_interaction_section(title: str, kicker: str, description: str) -> None:
    """Crea una capçalera visual compacta per a una secció de pàgina."""

    render_kicker_section("interaction-section-card", title, kicker, description)


# La configuració visual es fa al principi perquè qualsevol rerun mantingui el mateix marc.
configure_studio_page(PAGE_TITLE, layout=PAGE_LAYOUT)
inject_interaction_styles()

render_interaction_hero()

# Inicialitzem les claus que Streamlit reutilitza entre reruns.
def init_session():
    """Inicialitzeu les claus de sessió utilitzades pel flux d'interacció en directe."""
    if "ix_running" not in st.session_state:
        st.session_state.ix_running = False
    if "ix_app_mode" not in st.session_state:
        st.session_state.ix_app_mode = None
    if "ix_env" not in st.session_state:
        st.session_state.ix_env = None
    if "ix_model" not in st.session_state:
        st.session_state.ix_model = None
    if "ix_obs" not in st.session_state:
        st.session_state.ix_obs = None
    if "ix_step" not in st.session_state:
        st.session_state.ix_step = 0
    if "ix_history" not in st.session_state:
        st.session_state.ix_history = []
    if "ix_reward" not in st.session_state:
        st.session_state.ix_reward = 0.0
    if "ix_info" not in st.session_state:
        st.session_state.ix_info = {}

init_session()


def stop_simulation():
    """Tanqueu l'entorn actiu i restabliu l'estat d'interacció en directe."""
    if st.session_state.ix_env is not None:
        try:
            st.session_state.ix_env.close()
        except Exception:
            pass
    st.session_state.ix_env = None
    st.session_state.ix_app_mode = None
    st.session_state.ix_model = None
    st.session_state.ix_obs = None
    st.session_state.ix_running = False
    st.session_state.ix_step = 0
    st.session_state.ix_history = []
    st.session_state.ix_reward = 0.0
    st.session_state.ix_info = {}
    st.session_state.pop("ix_observation_variables", None)
    st.session_state.pop("ix_model_observation_size", None)
    st.session_state.pop("ix_test_obs_vals", None)


def randomize_policy_test_observations(obs_vars: list[str]) -> None:
    """Substituïu les observacions de la prova de política per valors aleatoris per a les variables actives."""

    random_obs_values = randomize_observation_values(obs_vars)
    st.session_state.ix_test_obs_vals = coerce_observation_values(random_obs_values, len(obs_vars))


def apply_live_action_step(
    vec_env: object,
    core_env: object,
    action_to_take: object,
    action_source: str = "",
    repeat_n: int = 1,
) -> None:
    """Avançar la simulació en directe i mantenir l'estat d'observació resultant."""

    try:
        with st.spinner(f"Simulant els propers {repeat_n} pas(sos) EnergyPlus ({action_source})..."):
            result = run_environment_steps(
                vec_env=vec_env,
                core_env=core_env,
                action_to_take=action_to_take,
                current_step=st.session_state.ix_step,
                repeat_n=repeat_n,
                obs_vars=st.session_state.get("ix_observation_variables"),
            )

            st.session_state.ix_obs = result["next_obs"]
            st.session_state.ix_reward = result["reward"]
            st.session_state.ix_info = result["info_dict"]
            st.session_state.ix_history.extend(result["history_entries"])
            st.session_state.ix_observation_variables = get_runtime_observation_variables(
                core_env=core_env,
                obs=result["next_obs"],
                model_obj=st.session_state.ix_model,
                vec_env=vec_env,
            )

            if result["history_entries"]:
                st.session_state.ix_step = result["history_entries"][-1]["pas"]

            if result["episode_finished"]:
                # Avís episodi reiniciat
                st.warning("Episodi acabat (Terminated / Truncated). L'entorn s'ha reiniciat.")
                st.session_state.ix_obs = result["reset_obs"]
                st.session_state.ix_step = 0
                st.session_state.ix_history = []

    except Exception as exc:
        # Error aplicar accio
        st.error(f"Error aplicant l'acció: {exc}")


def _calendar_or_reward_metric(curr: dict, prev: dict | None, core_env: object) -> dict:
    """Prepara la mètrica de calendari o recompensa quan el temps no és disponible."""
    month_col = find_matching_column(curr, ["month"])
    day_col = find_matching_column(curr, ["day_of_month", "day"])
    if month_col and day_col:
        month = int(curr[month_col])
        day = int(curr[day_col])
        month_labels = ["Gen", "Feb", "Mar", "Abr", "Mai", "Jun", "Jul", "Ago", "Set", "Oct", "Nov", "Des"]
        season = "Hivern" if month in [12, 1, 2] else "Primavera" if month in [3, 4, 5] else "Estiu" if month in [6, 7, 8] else "Tardor"
        timestep_mins = "Desconegut"
        try:
            if hasattr(core_env, "get_wrapper_attr"):
                timestep_mins = getattr(core_env.unwrapped, "timestep", 15)
        except Exception:
            pass
        return {
            "label": f"Calendari (Pas: {timestep_mins} min)",
            "value": f"{day} {month_labels[month - 1] if 1 <= month <= 12 else month}",
            "delta": season,
            "delta_color": "off",
        }

    value = curr.get("reward", 0.0)
    previous = prev.get("reward", value) if prev else value
    return {
        "label": "Recompensa Agent",
        "value": f"{value:.2f}",
        "delta": f"{value - previous:+.2f}",
        "delta_color": "normal",
    }


# Flux de configuració abans de crear l'entorn.
if not st.session_state.ix_running:
    # Selecció del mode
    # Controls mode i entorn
    col_mode, col_env = st.columns([1, 1.35], gap="small")
    col_mode.markdown('<div class="studio-radio-card-anchor"></div>', unsafe_allow_html=True)
    col_env.markdown('<div class="interaction-env-card-anchor"></div>', unsafe_allow_html=True)
    # Selector mode funcionament
    use_agent = col_mode.radio(
        "Mode de Funcionament",
        [
            "Simulació en Temps Real interactiva",
            "Avaluació Manual Estàtica (Sense Simulació)",
        ],
        index=0,
    )

    # Si cal model, primer mirem el metadata del .zip; és més fiable que deduir l'entorn pel nom.
    selected_env_id = ""
    model_path = None

    if mode_requires_model(use_agent):
        model_scan_dirs = ",".join([str(Path.cwd() / "models"), str(Path.cwd())])
        df_models = scan_model_zips(model_scan_dirs)

        if not df_models.empty:
            options = list(range(len(df_models)))
            # Selector model entrenat
            idx = st.selectbox(
                "Selecciona el model entrenat (.zip)",
                options=options,
                format_func=lambda i: str(df_models.iloc[i]["stem"])
            )
            row = df_models.iloc[idx]
            model_path = Path(row["path"])

            with open(model_path, "rb") as fh:
                _, meta = load_sb3_model_bytes(fh.read(), device="cpu")

            training_metadata = load_model_training_metadata(model_path)
            suggested_env = training_metadata.get("env_id") or env_id_from_meta_or_name(meta, row["stem"])
            # Input entorn detectat
            selected_env_id = col_env.text_input("ID Entorn Sinergym", value=suggested_env or "")
        else:
            # Avís models no trobats
            st.warning("No s'han trobat models .zip")
            # Input entorn manual
            selected_env_id = col_env.text_input("ID Entorn Sinergym", value="Eplus-5zone-hot-continuous-v1")
    else:
        # En mode manual no hi ha model, però igualment cal un env_id vàlid per crear l'entorn.
        # Input entorn manual
        selected_env_id = col_env.text_input("ID Entorn Sinergym", value="Eplus-5zone-hot-continuous-v1")

    st.divider()
    init_text = "Inicialitzar Entorn i Simulació" if "Simulació" in use_agent else "Carregar Agent (Sense Simular)"
    # Botó inicialitzar entorn
    if st.button(
        init_text,
        key="ix_initialize_button",
        type="primary",
        disabled=not selected_env_id,
        width="stretch",
    ):
        with st.spinner("Creant l'entorn EnergyPlus i inicialitzant..."):
            try:
                venv, core_env, model_obj, obs, init_warnings = initialize_runtime(
                    selected_env_id=selected_env_id,
                    mode_label=use_agent,
                    model_path=model_path,
                )
                for warn_msg in init_warnings:
                    # Avís inicialitzacio
                    st.warning(warn_msg)

                st.session_state.ix_env = venv
                st.session_state.ix_core_env = core_env
                st.session_state.ix_model = model_obj
                st.session_state.ix_obs = obs
                st.session_state.ix_observation_variables = get_runtime_observation_variables(
                    core_env=core_env,
                    obs=obs,
                    model_obj=model_obj,
                    vec_env=venv,
                )
                st.session_state.ix_model_observation_size = get_model_observation_size(model_obj)
                st.session_state.ix_running = True
                st.session_state.ix_app_mode = use_agent
                st.session_state.ix_step = 0
                st.session_state.ix_reward = 0.0
                st.session_state.ix_history = []
                st.session_state.ix_info = {}
                st.rerun()

            except Exception as e:
                # Error inicialitzacio
                st.error(f"Error inicialitzant: {e}")
                stop_simulation()


# Flux d'interacció quan l'entorn ja està creat.
else:
    venv = st.session_state.ix_env
    core_env = st.session_state.get("ix_core_env", venv)
    model = st.session_state.ix_model
    obs = st.session_state.ix_obs
    app_mode = st.session_state.get("ix_app_mode", "")

    # En aquest mode només consultem la política: no avancem EnergyPlus.
    if app_mode and "Avaluació" in app_mode:
        # Avís mode estatic
        st.info("Avaluació Manual Estàtica. Aquest mode no avança el rellotge del simulador, simplement mostra què faria l'agent donat qualsevol input de l'edifici.")
        render_interaction_section(
            "Avaluació estàtica",
            "Consulta de política",
            "Construeix una observació numèrica i comprova quina acció retornaria el model sense avancar el simulador.",
        )
        # Capçalera consulta estatica
        cols_hdr = st.columns([4, 1])
        # Botó aturar consulta
        if cols_hdr[1].button(
            "Aturar i Sortir",
            key="btn_stop_test",
            type="primary",
            width="stretch",
        ):
            stop_simulation()
            st.rerun()

        st.subheader("Simulació numèrica: observacions")
        obs_vars = st.session_state.get("ix_observation_variables") or get_runtime_observation_variables(
            core_env=core_env,
            obs=obs,
            model_obj=model,
            vec_env=venv,
        )
        expected_obs_size = st.session_state.get("ix_model_observation_size") or get_model_observation_size(model)

        # Agafem una observació real com a punt de partida perquè la consulta estàtica no
        # comenci amb zeros que no tindrien sentit físic.
        if "ix_test_obs_vals" not in st.session_state or len(st.session_state.ix_test_obs_vals) != len(obs_vars):
            default_obs_values = extract_policy_test_defaults(
                obs=obs,
                info_dict=st.session_state.get("ix_info", {}),
                obs_vars=obs_vars,
                vec_env=venv,
                core_env=core_env,
            )
            st.session_state.ix_test_obs_vals = coerce_observation_values(default_obs_values, len(obs_vars))

        # Botó observacions aleatories
        st.button(
            "Generar observacions aleatòries",
            key="ix_randomize_observations",
            on_click=randomize_policy_test_observations,
            args=(obs_vars,),
        )

        custom_obs = []

        # El formulari evita reexecutar la predicció cada vegada que es toca un input numèric.
        # Formulari consulta politica
        with st.form("form_policy_test"):
            # Graella inputs observacio
            grid = st.columns(3)
            for i, var_name in enumerate(obs_vars):
                val_d = float(st.session_state.ix_test_obs_vals[i]) if i < len(st.session_state.ix_test_obs_vals) else 0.0
                unit = infer_unit(var_name)
                label = f"{var_name} ({unit})" if unit else var_name

                with grid[i % 3]:
                    # Input valor observacio
                    val = st.number_input(label, value=val_d, format="%.3f")
                    custom_obs.append(val)

            # Botó consultar agent
            submitted = st.form_submit_button("Consultar Decisió de l'Agent", type="primary")

        if submitted:
            st.divider()
            if model is None:
                # Error sense model
                st.error("No hi ha model carregat!")
            else:
                custom_obs = coerce_observation_values(custom_obs, expected_obs_size)
                norm_obs = normalize_policy_observation(custom_obs, core_env, venv, expected_obs_size)

                # Predim amb l'observació normalitzada igual que durant l'entrenament.
                action, state = model.predict(norm_obs, deterministic=True)

                # Desnormalitzem l'acció només per ensenyar-la; al simulador li passarem
                # sempre el format que espera el VecEnv.
                action_display = prepare_action_display(action, venv, core_env)

                action_vars = get_wrapper_variables(core_env, "action_variables")

                st.subheader("Resposta immediata de la xarxa neuronal")

                # Metriques accio agent
                acts_c = st.columns(max(len(action_vars), 1))
                if len(action_vars) > 0 and len(action_vars) == getattr(action_display, 'size', -1):
                    for i, act_n in enumerate(action_vars):
                        with acts_c[i]:
                            unit_a = infer_unit(act_n)
                            val_str = f"{action_display[i]:.2f} {unit_a}" if unit_a else f"{action_display[i]:.2f}"
                            # Metrica accio agent
                            st.metric(f"Acció '{act_n}'", val_str)
                else:
                    st.write(action_display)

                # Text accio crua
                st.caption(f"L'acció crua normalitzada retornada per la xarxa era: `{action}`. Si el model conté `NormalizeAction`, es tradueix als llindars de Sinergym a dalt.")

        # Interrompre execució aquí perquè no es pinti el mode dinàmic abaix
        st.stop()

    # Missatge simulacio activa
    st.success("Simulació en curs amb el motor d'**EnergyPlus real**. Esperant interacció pas a pas.")

    render_interaction_section(
        "Simulació activa",
        "Temps real",
        "Consulta l'estat actual de l'entorn, aplica accions amb l'agent o manualment i segueix els indicadors principals sense alterar el flux existent.",
    )

    # Resum ràpid de l'estat actual.
    # Metriques estat simulacio
    col1, col2, col3, col4 = st.columns(4)
    col1.metric("Pas Actual", st.session_state.ix_step)
    col2.metric("Última Recompensa", f"{st.session_state.ix_reward:.3f}")
    # Botó aturar simulacio
    if col4.button("Aturar i Tancar", key="ix_stop_live_button", width="stretch"):
        stop_simulation()
        st.rerun()

    st.divider()

    render_interaction_section(
        "Observació actual",
        "Estat de l'entorn",
        "Mostra els valors físics disponibles a la darrera observació per facilitar la lectura abans de decidir la propera acció.",
    )

    # Bloc d'observacions actuals.
    st.subheader("Observació actual (valors físics reals)")

    obs_vars = st.session_state.get("ix_observation_variables") or get_runtime_observation_variables(
        core_env=core_env,
        obs=obs,
        model_obj=model,
        vec_env=venv,
    )
    _, display_obs = extract_display_observation(
        obs=obs,
        info_dict=st.session_state.get("ix_info", {}),
        obs_vars=obs_vars,
        vec_env=venv,
    )

    if len(obs_vars) == len(display_obs):
        # Crea un DataFrame agradable
        df_obs = pd.DataFrame({"Variable": obs_vars, "Valor Original": display_obs})
        # Arrodonim una mica per que no sigui massa brut
        df_obs["Valor Original"] = df_obs["Valor Original"].apply(lambda x: round(x, 4) if isinstance(x, (int, float)) else x)
        # Taula observacio actual
        st.dataframe(df_obs, width="stretch", hide_index=True)
    else:
        st.write(display_obs)

    st.divider()

    # Bloc d'execució d'accions.
    render_interaction_section(
        "Accions",
        "Agent o manual",
        "Escull si vols delegar la decisió a la política carregada o fixar manualment els valors aplicats al simulador.",
    )

    # Selecció de l'acció a aplicar al pas següent.
    st.subheader("Acció")

    # Pestanyes accio agent/manual
    tab_agent, tab_manual = st.tabs(["Agent", "Control Manual"])

    with tab_agent:
        if model is None:
            # Avís sense model actiu
            st.info("No hi ha cap model carregat en aquest moment. Fes servir el mode manual o atura i carrega'n un.")
        else:
            # Input passos agent
            n_steps = st.number_input("Quants passos vols avançar d'un cop?", min_value=1, max_value=1000, value=1, step=1, key="n_agent")
            # Botó avançar agent
            if st.button(
                f"Avançar {n_steps} pas(sos) amb l'Agent",
                key="ix_agent_step_button",
                type="primary",
                width="stretch",
            ):
                try:
                    with st.spinner(f"L'agent està simulant {n_steps} passos..."):
                        for _ in range(n_steps):
                            current_obs = st.session_state.ix_obs
                            act, _ = model.predict(current_obs, deterministic=True)
                            apply_live_action_step(venv, core_env, act, "Agent", repeat_n=1)
                except Exception as e:
                    # Error avançar agent
                    st.error(str(e))
                st.rerun()

    with tab_manual:
        st.write("Selecciona manualment els valors. L'espai d'accions es detecta automàticament des de l'entorn actiu.")

        # Recuperem les variables exposades pels wrappers, no només l'action_space cru.
        action_vars = get_wrapper_variables(core_env, "action_variables")
        action_space = getattr(venv, "action_space", getattr(core_env, "action_space", None))

        user_action = []

        if isinstance(action_space, gym.spaces.Box):
            # Accions contínues: un slider per dimensió.
            action_dim = int(np.prod(action_space.shape))
            lows = np.asarray(action_space.low, dtype=np.float32).reshape(-1)
            highs = np.asarray(action_space.high, dtype=np.float32).reshape(-1)
            for i in range(action_dim):
                v_name = action_vars[i] if len(action_vars) > i else f"Acció_{i}"
                v_low = float(lows[i])
                v_high = float(highs[i])

                # Alguns espais declaren límits infinits; els acotem per poder pintar sliders.
                if v_low == -np.inf: v_low = -100.0
                if v_high == np.inf: v_high = 100.0

                # Slider accio continua
                val = st.slider(v_name, min_value=v_low, max_value=v_high, value=(v_low+v_high)/2, step=0.1)
                user_action.append(val)

        elif isinstance(action_space, gym.spaces.Discrete):
            # Acció discreta simple: un enter entre 0 i n-1.
            st.write(f"L'espai d'acció és Discret ({action_space.n} possibilitats).")
            # Input accio discreta
            val = st.number_input("Selecciona Acció", min_value=0, max_value=int(action_space.n)-1, step=1)
            user_action.append(val)

        elif isinstance(action_space, gym.spaces.MultiDiscrete):
            nvec = np.asarray(action_space.nvec, dtype=np.int64).reshape(-1)
            st.write(f"L'espai d'acció és MultiDiscrete ({len(nvec)} dimensions).")
            for i, n_value in enumerate(nvec):
                v_name = action_vars[i] if len(action_vars) > i else f"Acció_{i}"
                # Input accio multidiscreta
                val = st.number_input(
                    v_name,
                    min_value=0,
                    max_value=max(int(n_value) - 1, 0),
                    value=0,
                    step=1,
                    key=f"manual_multidiscrete_{i}",
                )
                user_action.append(val)

        elif isinstance(action_space, gym.spaces.MultiBinary):
            action_dim = int(np.prod(action_space.shape))
            st.write(f"L'espai d'acció és MultiBinary ({action_dim} dimensions).")
            for i in range(action_dim):
                v_name = action_vars[i] if len(action_vars) > i else f"Acció_{i}"
                # Checkbox accio binaria
                val = st.checkbox(v_name, value=False, key=f"manual_multibinary_{i}")
                user_action.append(1 if val else 0)

        else:
            # Avís espai accio no suportat
            st.warning("Espai d'acció no suportat actualment per la inserció dinàmica manual.")

        # Input passos manuals
        n_steps_mod = st.number_input("Quants passos (repetint l'acció) vols avançar d'un cop?", min_value=1, max_value=1000, value=1, step=1, key="n_manual")

        # Botó avançar manual
        if st.button(
            f"Avançar {n_steps_mod} pas(sos) (Manual)",
            key="ix_manual_step_button",
            width="stretch",
        ):
            if isinstance(action_space, (gym.spaces.Box, gym.spaces.Discrete, gym.spaces.MultiDiscrete, gym.spaces.MultiBinary)):
                apply_live_action_step(venv, core_env, user_action, "Manual", repeat_n=n_steps_mod)
                st.rerun()

    # Resum ràpid de KPIs del pas actual.
    if len(st.session_state.ix_history) > 0:
        render_interaction_section(
            "Dashboard principal",
            "Indicadors",
            "Resumeix l'estat recent de la simulació amb setpoints, temperatures, consum i confort a partir de l'històric disponible.",
        )
        st.divider()
        st.subheader("Dashboard principal (estat actual)")

        # Agafem només els dos últims passos per poder calcular deltes sense carregar
        # tot l'històric a cada rerun.
        history_len = len(st.session_state.ix_history)
        curr = st.session_state.ix_history[-1]
        prev = st.session_state.ix_history[-2] if history_len > 1 else None

        heat_col = find_matching_column(curr, ['heating', 'htg_setpoint'])
        cool_col = find_matching_column(curr, ['cooling', 'clg_setpoint'])
        in_temp_col = find_matching_column(curr, ['air_temperature', 'indoor_temperature'])
        out_temp_col = find_matching_column(curr, ['outdoor_temperature', 'site_outdoor_air_drybulb_temperature'])
        power_col = find_matching_column(curr, ['demand_rate', 'power', 'hvac_elèctricity_demand'])

        # Primera fila: què ha decidit l'actuador i com canvia respecte al pas anterior.
        # Metriques accions actuador
        st.markdown("##### Accions de l'actuador")
        render_metric_row(
            [
                build_metric_item(
                    curr,
                    prev,
                    heat_col,
                    label="Setpt. Calor (°C)",
                    value_format=lambda value: f"{value:.1f}°C",
                    delta_format=lambda delta: f"{delta:+.1f}°C",
                    delta_color="normal",
                    empty_message="Sense Setpoint de Calor",
                ),
                build_metric_item(
                    curr,
                    prev,
                    cool_col,
                    label="Setpt. Fred (°C)",
                    value_format=lambda value: f"{value:.1f}°C",
                    delta_format=lambda delta: f"{delta:+.1f}°C",
                    delta_color="inverse",
                    empty_message="Sense Setpoint de Fred",
                ),
                _calendar_or_reward_metric(curr, prev, core_env),
            ]
        )


        # Separador dashboard
        st.markdown("<div class='studio-spacer-100'></div>", unsafe_allow_html=True)

        # Segona fila: estat sensorial mínim per entendre si l'acció tenia sentit.
        # Metriques sensors
        st.markdown("##### Sensors (observacions)")
        render_metric_row(
            [
                build_metric_item(
                    curr,
                    prev,
                    in_temp_col,
                    label="Temp. Interior T1",
                    value_format=lambda value: f"{value:.1f}°C",
                    delta_format=lambda delta: f"{delta:+.2f}°C",
                    empty_message="Sense T. Interior",
                ),
                build_metric_item(
                    curr,
                    prev,
                    out_temp_col,
                    label="Meteo. exterior",
                    value_format=lambda value: f"{value:.1f}°C",
                    delta_format=lambda delta: f"{delta:+.1f}°C",
                    empty_message="Sense T. Exterior",
                ),
                build_metric_item(
                    curr,
                    prev,
                    power_col,
                    label="Consum HVAC (Act/Rate)",
                    value_format=lambda value: f"{value / 1000:.2f} kW",
                    delta_format=lambda delta: f"{delta / 1000:+.2f} kW",
                    delta_color="inverse",
                    empty_message="Sense Consum Electric",
                ),
            ]
        )


        # Separador confort
        st.markdown("<div class='studio-spacer-100'></div>", unsafe_allow_html=True)

        # Tercera fila: confort calculat amb el rang actiu de la reward.
        # Bloc estat confort
        st.markdown("##### Estat de Confort")

        # Columnes confort PMV
        comfort_col1, comfort_col2 = st.columns(2)

        with comfort_col1:
            if in_temp_col:
                current_temp = curr[in_temp_col]
                # Busquem els rangs a la reward activa. Si no hi són, fem servir valors típics
                # per poder donar feedback sense rebentar la simulació.
                if hasattr(core_env, 'get_wrapper_attr'):
                    # Sinergym acostuma a guardar aquests rangs dins reward_kwargs.
                    is_summer = False
                    try:
                        # Inferència simple de temporada a partir del mes actual.
                        month_col = find_matching_column(curr, ['month'])
                        if month_col and curr[month_col] >= 6 and curr[month_col] <= 9:
                            is_summer = True
                    except: pass

                    try:
                        # Depèn de la reward concreta, per això el fallback és intencionat.
                        cfg = core_env.get_wrapper_attr('reward_kwargs') if hasattr(core_env, 'get_wrapper_attr') else {}
                        r_comfort_winter = cfg.get('range_comfort_winter', (20.0, 23.5))
                        r_comfort_summer = cfg.get('range_comfort_summer', (23.0, 26.0))
                    except:
                        r_comfort_winter = (20.0, 23.5)
                        r_comfort_summer = (23.0, 26.0)

                    active_range = r_comfort_summer if is_summer else r_comfort_winter

                    if active_range[0] <= current_temp <= active_range[1]:
                        # Missatge confort correcte
                        st.success(f"\u2705 Dins el Rang de Confort ({active_range[0]}°C a {active_range[1]}°C)")
                    elif current_temp < active_range[0]:
                        # Error confort fred
                        st.error(f"\u2744 Massa Fred (Objectiu: {active_range[0]}°C a {active_range[1]}°C)")
                    else:
                        # Error confort calor
                        st.error(f"\U0001f525 Massa Calor (Objectiu: {active_range[0]}°C a {active_range[1]}°C)")
            else:
                # Avís confort no calculable
                st.info("No es pot calcular el confort (Sense Temperatura Interior)")

        with comfort_col2:
            pmv_col = find_matching_column(curr, ['pmv'])
            ppd_col = find_matching_column(curr, ['ppd'])

            if pmv_col and ppd_col:
                val_pmv = curr[pmv_col]
                val_ppd = curr[ppd_col]

                # PMV Ideal: 0 (Rang -0.5 a +0.5). PPD Ideal: < 10%
                pmv_class = (
                    "comfort-value-good"
                    if -0.5 <= val_pmv <= 0.5
                    else "comfort-value-warn"
                    if -1.0 <= val_pmv <= 1.0
                    else "comfort-value-bad"
                )
                ppd_class = (
                    "comfort-value-good"
                    if val_ppd <= 10.0
                    else "comfort-value-warn"
                    if val_ppd <= 20.0
                    else "comfort-value-bad"
                )


                # Text PMV
                st.markdown(
                    f"**Vot Mitjà (PMV):** <span class='{pmv_class}'>{val_pmv:.2f}</span>",
                    unsafe_allow_html=True,
                )
                # Text PPD
                st.markdown(
                    f"**Insatisfacció (PPD):** <span class='{ppd_class}'>{val_ppd:.1f}%</span>",
                    unsafe_allow_html=True,
                )
            else:
                st.write("Variables PMV/PPD no disponibles en aquest entorn.")
\end{lstlisting}

\Needspace{8\baselineskip}

\subsection{\texorpdfstring{\texttt{pages/Mostrar\_Entorn.py}}{pages/Mostrar\_Entorn.py}}

\begin{lstlisting}[style=studioPython,caption={pages/Mostrar\_Entorn.py},label={lst:pages-mostrar-entorn-py}]
"""Pàgina de visualització d'entorns Sinergym.

Mostra la geometria 3D, el clima, la configuració i els actius energètics
associats a l'entorn seleccionat.
"""

from html import escape
from typing import Any, Dict, List, Tuple

import gymnasium as gym
import numpy as np
import pandas as pd
import streamlit as st
from page_components.ui_fragments import render_hero, render_section_title
from sidebar_nav import configure_studio_page

from backend.mapa_backend import render_location_map
from backend.mostrar_entorn_backend import (
    _format_ui_value,
    _reward_summary_frames,
    _sequence_to_df,
    _tuple_mapping_to_df,
    load_environment_assets,
    parse_epw_metadata_and_temperature,
    summarize_pv,
)
from backend.common import BUILDINGS_DIR, WEATHERS_DIR
from backend.common import list_registered_env_ids
from backend.viewer_3d_backend import plot_3d_model
from page_styles.mostrar_entorn import inject_environment_styles


PAGE_TITLE = "Mostrar Entorn"
PAGE_LAYOUT = "wide"
INTRODUCTION_TEXT = (
    "Consulta els entorns registrats, explora la geometria de l'edifici "
    "i revisa configuració, clima i actius energètics des d'una sola pàgina."
)


def render_environment_hero() -> None:
    """Munta el bloc principal de la pàgina."""

    render_hero("environment-hero", "Visualització d'entorns", PAGE_TITLE, INTRODUCTION_TEXT)


def render_environment_section(title: str, kicker: str, description: str) -> None:
    """Mostra la capçalera visual d'una secció."""

    render_section_title(title)


def describe_environment_tags(env_name: str) -> List[str]:
    """Resumeix propietats de l'entorn en etiquetes curtes."""

    tags: List[str] = []
    lower_name = env_name.lower()
    if "continuous" in lower_name:
        tags.append("Accions continues")
    if "discrete" in lower_name:
        tags.append("Accions discretes")
    if "stochastic" in lower_name:
        tags.append("Entorn aleatori")
    if "mixed" in lower_name:
        tags.append("Accions mixtes")
    if "hot" in lower_name:
        tags.append("Clima calid")
    if "cool" in lower_name:
        tags.append("Clima moderat")
    if "cold" in lower_name:
        tags.append("Clima fred")
    return tags


def render_environment_tags(tags: List[str]) -> None:
    """Mostra propietats detectades en format d'etiqueta."""

    if not tags:
        return
    tags_html = "".join(f'<span class="environment-tag">{escape(tag)}</span>' for tag in tags)
    # Etiquetes entorn
    st.markdown(f'<div class="environment-tags">{tags_html}</div>', unsafe_allow_html=True)


def render_metric_card(label: str, value: str, variant: str = "") -> None:
    """Crea una targeta compacta amb una metrica o text curt."""

    extra_class = f" environment-metric-card--{variant}" if variant else ""
    # Targeta metrica entorn
    st.markdown(
        (
            f'<div class="environment-metric-card{extra_class}">'
            f'<div class="environment-metric-label">{escape(label)}</div>'
            f'<div class="environment-metric-value">{escape(value)}</div>'
            "</div>"
        ),
        unsafe_allow_html=True,
    )


def render_metric_cards(items: List[Tuple[str, str]], columns: int = 2, wide_first: bool = False) -> None:
    """Distribueix targetes compactes en files estables."""

    if not items:
        return

    remaining_items = items
    if wide_first:
        first_label, first_value = remaining_items[0]
        render_metric_card(first_label, first_value, variant="wide")
        remaining_items = remaining_items[1:]

    column_count = max(1, columns)
    for start in range(0, len(remaining_items), column_count):
        row_items = remaining_items[start : start + column_count]
        # Fila targetes metriques
        row_columns = st.columns(len(row_items), gap="small")
        for column, (label, value) in zip(row_columns, row_items):
            with column:
                render_metric_card(label, value)


def render_detail_card(title: str, rows: List[Tuple[str, str]]) -> None:
    """Crea una targeta compacta de detall."""

    rows_html = "".join(
        (
            '<div class="environment-detail-row">'
            f'<div class="environment-detail-label">{escape(str(label))}</div>'
            f'<div class="environment-detail-value">{escape(str(value))}</div>'
            '</div>'
        )
        for label, value in rows
    )
    card_parts = ['<div class="environment-detail-card">']
    if title:
        card_parts.append(f'<div class="environment-detail-title">{escape(title)}</div>')
    card_parts.append(f'<div class="environment-detail-list">{rows_html}</div>')
    card_parts.append('</div>')
    # Targeta detall entorn
    st.markdown("".join(card_parts), unsafe_allow_html=True)


def render_battery_cards(storage_list: List[dict]) -> None:
    """Mostra les bateries detectades com a targetes visuals."""

    for start in range(0, len(storage_list), 2):
        row_items = storage_list[start : start + 2]
        # Fila targetes bateria
        row_columns = st.columns(len(row_items), gap="small")
        for item_index, (column, storage) in enumerate(zip(row_columns, row_items), start=start + 1):
            rows = [
                ("Capacitat util", f"{_format_ui_value(storage.get('energy_kwh'))} kWh"),
                ("Carrega maxima", f"{_format_ui_value(storage.get('p_charge_kw'))} kW"),
                ("Descarrega maxima", f"{_format_ui_value(storage.get('p_discharge_kw'))} kW"),
                ("Rendiment", _format_ui_value(storage.get('eff_roundtrip'))),
            ]
            display_name = "Bateria" if len(storage_list) == 1 else f"Bateria {item_index}"
            rows_html = "".join(
                (
                    '<div class="environment-detail-row">'
                    f'<div class="environment-detail-label">{escape(label)}</div>'
                    f'<div class="environment-detail-value">{escape(value)}</div>'
                    "</div>"
                )
                for label, value in rows
            )
            with column:
                # Targeta bateria
                st.markdown(
                    (
                        '<div class="environment-battery-card">'
                        f'<div class="environment-battery-name">{escape(display_name)}</div>'
                        f'<div class="environment-battery-type">{escape(_format_ui_value(storage.get("type")))}</div>'
                        f'<div class="environment-battery-list">{rows_html}</div>'
                        "</div>"
                    ),
                    unsafe_allow_html=True,
                )


def _reward_label(reward_value: Any) -> str:
    """Retorna una etiqueta llegible per al reward configurat."""

    return getattr(reward_value, '__name__', _format_ui_value(reward_value))


def render_basic_config(spec: gym.envs.registration.EnvSpec, kwargs: Dict[str, Any]) -> None:
    """Omple la secció de configuració bàsica a la UI de Streamlit."""
    # Seccio configuracio general
    st.markdown("#### Configuració general")
    render_detail_card(
        "",
        [
            ("ID", _format_ui_value(spec.id)),
            ("Entry point", _format_ui_value(spec.entry_point)),
            ("Fitxer edifici", _format_ui_value(kwargs.get('building_file'))),
            ("Reward", _reward_label(kwargs.get('reward'))),
        ],
    )


def render_action_space_summary(action_space: Any) -> None:
    """Resumeix l'espai d'acció a la UI de Streamlit."""
    if action_space is None:
        st.code('-', language=None)
        return

    if isinstance(action_space, gym.spaces.Box):
        low = np.asarray(action_space.low, dtype=float).reshape(-1)
        high = np.asarray(action_space.high, dtype=float).reshape(-1)
        bounds_df = pd.DataFrame([
            {
                'Dimensió': index + 1,
                'Min': round(float(low_value), 3),
                'Max': round(float(high_value), 3),
            }
            for index, (low_value, high_value) in enumerate(zip(low, high))
        ])
        # Taula rangs accio
        st.dataframe(bounds_df, hide_index=True, width="stretch", height=min(220, 35 * len(bounds_df) + 38))
        return

    if isinstance(action_space, gym.spaces.Discrete):
        render_detail_card(
            "",
            [
                ("Tipus", "Discret"),
                ("Accions possibles", str(int(action_space.n))),
            ],
        )
        return

    if isinstance(action_space, gym.spaces.MultiDiscrete):
        nvec = np.asarray(action_space.nvec).reshape(-1)
        dims_df = pd.DataFrame([
            {
                'Dimensió': index + 1,
                'Tipus': 'MultiDiscrete',
                'Valors possibles': int(values),
            }
            for index, values in enumerate(nvec)
        ])
        # Taula accions multidiscretes
        st.dataframe(dims_df, hide_index=True, width="stretch", height=min(220, 35 * len(dims_df) + 38))
        return

    if isinstance(action_space, gym.spaces.MultiBinary):
        render_detail_card(
            "",
            [
                ("Tipus", "MultiBinary"),
                ("Bits", str(int(np.prod(action_space.shape)))),
            ],
        )
        return

    st.code(_format_ui_value(action_space), language=None)


def _actuators_mapping_df(actuators: Dict[str, Any]) -> pd.DataFrame:
    """Construeix un resum llegible del mapatge entre dimensions i actuadors."""

    rows = []
    for index, (actuator_key, metadata) in enumerate((actuators or {}).items(), start=1):
        if isinstance(metadata, dict):
            rows.append(
                {
                    'Dimensió': index,
                    'Actuador': actuator_key,
                    'Tipus objecte': _format_ui_value(metadata.get('element_type')),
                    'Camp de control': _format_ui_value(metadata.get('value_type')),
                    'Variable RL': _format_ui_value(metadata.get('variable_name')),
                }
            )
        else:
            rows.append(
                {
                    'Dimensió': index,
                    'Actuador': actuator_key,
                    'Tipus objecte': _format_ui_value(metadata),
                    'Camp de control': '-',
                    'Variable RL': '-',
                }
            )
    return pd.DataFrame(rows)


def render_kwargs_overview(kwargs: Dict[str, Any]) -> None:
    """Secció de visió general de renderització de kwargs a la UI de Streamlit."""
    # Els kwargs de Sinergym barregen variables, actuadors, reward i camps interns.
    # Els separem en pestanyes perquè sigui revisable sense llegir el JSON sencer.
    # Pestanyes kwargs entorn
    tab_obs, tab_control, tab_reward, tab_misc, tab_raw = st.tabs([
        'Variables', 'Control', 'Reward', 'Extra', 'JSON brut'
    ])

    with tab_obs:
        time_df = _sequence_to_df(kwargs.get('time_variables', []), 'Variable temporal')
        if not time_df.empty:
            # Taula variables temporals
            st.markdown('**Variables temporals**')
            st.dataframe(time_df, hide_index=True, width="stretch")

        variables_df = _tuple_mapping_to_df(kwargs.get('variables', {}), ["Variable d'EnergyPlus", "Ambit"])
        if not variables_df.empty:
            # Taula variables observades
            st.markdown('**Variables observades**')
            st.dataframe(variables_df, hide_index=True, width="stretch", height=320)
        else:
            st.info('No hi ha variables definides per aquest entorn.')

        meters_df = pd.DataFrame(
            [{'Alias': alias, 'Meter': meter} for alias, meter in (kwargs.get('meters') or {}).items()]
        )
        if not meters_df.empty:
            # Taula meters
            st.markdown('**Meters**')
            st.dataframe(meters_df, hide_index=True, width="stretch")

    with tab_control:
        # Seccio espai accio
        st.markdown("**Espai d'acció**")
        render_action_space_summary(kwargs.get('action_space'))
        st.caption("Cada dimensió correspon, en aquest ordre, als actuadors definits a continuació.")

        actuators_df = _actuators_mapping_df(kwargs.get('actuators', {}))
        if not actuators_df.empty:
            # Taula actuadors
            st.markdown('**Actuadors**')
            st.dataframe(actuators_df, hide_index=True, width="stretch", height=260)
        else:
            st.info('No hi ha actuadors configurats.')

    with tab_reward:
        scalar_reward_df, grouped_reward_df = _reward_summary_frames(kwargs.get('reward_kwargs') or {})
        if not scalar_reward_df.empty:
            # Taula reward simple
            st.markdown('**Paràmetres simples**')
            st.dataframe(scalar_reward_df, hide_index=True, width="stretch")
        if not grouped_reward_df.empty:
            # Taula reward agrupada
            st.markdown('**Paràmetres agrupats**')
            st.dataframe(grouped_reward_df, hide_index=True, width="stretch")
        if scalar_reward_df.empty and grouped_reward_df.empty:
            st.info('No hi ha reward kwargs definits.')

    with tab_misc:
        weather_rows = []
        for variable_name, values in (kwargs.get('weather_variability') or {}).items():
            # La variabilitat pot arribar com a escalar o com a tuple sigma/mu/període.
            # Normalitzem la forma només per mostrar-la en taula.
            if isinstance(values, (list, tuple)):
                sigma = values[0] if len(values) > 0 else '-'
                mu = values[1] if len(values) > 1 else '-'
                period = values[2] if len(values) > 2 else '-'
            else:
                sigma = values
                mu = '-'
                period = '-'
            weather_rows.append({
                'Variable': variable_name,
                'Sigma': _format_ui_value(sigma),
                'Mu': _format_ui_value(mu),
                'Periode': _format_ui_value(period),
            })
        weather_df = pd.DataFrame(weather_rows)
        if not weather_df.empty:
            # Taula variabilitat clima
            st.markdown('**Variabilitat meteorològica**')
            st.dataframe(weather_df, hide_index=True, width="stretch")

        misc_df = pd.DataFrame([
            {'Clau': 'building_config', 'Valor': _format_ui_value(kwargs.get('building_config'))},
            {'Clau': 'context', 'Valor': _format_ui_value(kwargs.get('context'))},
            {'Clau': 'initial_context', 'Valor': _format_ui_value(kwargs.get('initial_context'))},
        ])
        # Taula camps extra
        st.markdown('**Altres camps**')
        st.dataframe(misc_df, hide_index=True, width="stretch")

    with tab_raw:
        # JSON brut kwargs
        st.json(kwargs)


def render_3d_viewer(
    records: List[dict],
    zones: List[str],
    types: List[str],
    z_clip_range: Tuple[float, float],
    pv_list: List[dict],
    name_index: Dict[str, int],
) -> None:
    """Mostra el visor de geometria de l'edifici en 3D.

    Mostra la figura 3D interactiva Plotly. Totes les superfícies es mostren amb
    per defecte amb tots els tipus i zones visibles.
    """
    if records:
        vis_types = {surface_type: True for surface_type in types}
        if pv_list:
            vis_types["PV"] = True
        vis_zones = {zone: True for zone in zones}
        fig3d = plot_3d_model(
            records=records,
            visible_types=vis_types,
            visible_zones=vis_zones,
            z_clip=z_clip_range,
            pv_list=pv_list,
            name_index=name_index,
        )
        # Visor 3D edifici
        st.plotly_chart(fig3d, width="stretch")
        st.caption("Passeu el cursor per veure zona, area i orientacio. Aspecte data i model recentrat.")
    else:
        st.warning("No s'han trobat superficies amb coordenades.")


def render_environment_page() -> None:
    """Munta la pàgina amb la presentació visual actualitzada."""

    configure_studio_page(PAGE_TITLE, layout=PAGE_LAYOUT)
    inject_environment_styles()

    render_environment_hero()
    all_envs = list_registered_env_ids()
    render_environment_section(
        "Selecció de l'entorn",
        "Catàleg",
        "Filtra els entorns registrats i tria quin vols obrir abans de veure la representació geomètrica i les dades de configuració.",
    )

    # Controls cerca entorn
    search_col, select_col = st.columns(2, gap="large")
    with search_col:
        # Input cerca entorn
        search_query = st.text_input("Cerca un entorn per nom", placeholder="Ex: warehouse, hot, continuous...")

    filtered_envs = [env_name for env_name in all_envs if search_query.lower() in env_name.lower()] if search_query else all_envs
    with select_col:
        # Selector entorn
        env_selected = st.selectbox(
            "Selecciona un entorn",
            filtered_envs if filtered_envs else ["Cap coincidència"],
            disabled=not filtered_envs,
        )

    # Text entorns filtrats
    st.caption(f"{len(filtered_envs)} entorns disponibles amb el filtre actual.")
    if not filtered_envs:
        # Avís entorn no trobat
        st.warning("No s'ha trobat cap entorn amb aquest filtre.")
        return

    render_environment_tags(describe_environment_tags(env_selected))

    try:
        spec = gym.spec(env_selected)
        kwargs = spec.kwargs
        bfile = kwargs.get("building_file")

        weather_file = None
        if "weather_files" in kwargs:
            raw = kwargs["weather_files"]
            weather_file = raw if isinstance(raw, str) else raw[0]
        elif "weather" in kwargs:
            wf = kwargs["weather"].get("weather_files", [])
            weather_file = wf if isinstance(wf, str) else (wf[0] if wf else None)

        render_environment_section(
            "Visualització de l'edifici",
            "Geometria",
            "La pàgina mostra la geometria 3D de l'edifici amb totes les superfícies visibles.",
        )

        records = []
        name_index = {}
        pv_list = []
        storage_list = []
        z_clip_range = (0.0, 0.0)
        zones: List[str] = []
        types: List[str] = []

        if bfile and (BUILDINGS_DIR / bfile).exists():
            with st.spinner("Carregant dades de l'entorn..."):
                # load_environment_assets centralitza el parseig pesat de l'IDF/epJSON.
                # La pàgina només decideix què es mostra i amb quins filtres.
                ep_path = str(BUILDINGS_DIR / bfile)
                assets = load_environment_assets(ep_path)
                records = assets["records"]
                zones = assets["zones"]
                types = assets["types"]
                z_clip_range = assets["z_clip_range"]
                name_index = assets["name_index"]
                pv_list = assets["pv"]
                storage_list = assets["storage"]

            render_3d_viewer(records, zones, types, z_clip_range, pv_list, name_index)
        else:
            # Avís visor no disponible
            st.info("No disponible per a aquest format o falta el fitxer d'edifici.")

        render_environment_section(
            "Configuració i dades associades",
            "Context tècnic",
            "Reuneix la configuració base de l'entorn, els paràmetres exposats per Sinergym i la informació del fitxer meteorològic associat.",
        )
        # Columnes configuracio entorn
        col1, col2 = st.columns(2)

        with col1:
            # Targeta configuracio base
            render_basic_config(spec, kwargs)

            weather_loaded = False
            if weather_file:
                wpath = WEATHERS_DIR / weather_file
                if wpath.exists():
                    weather_loaded = True
                    loc, avg_temp, cz = parse_epw_metadata_and_temperature(str(wpath))
                    # Seccio clima localitzacio
                    st.markdown("### Clima i localització")
                    st.caption("Resum del fitxer meteorològic actiu i ubicació aproximada de l'entorn.")
                    render_detail_card(
                        "Informació general",
                        [
                            ("Fitxer clima", weather_file),
                            ("Temp. mitjana", f"{avg_temp:.2f} C"),
                            ("Zona climàtica", str(cz)),
                            ("Ciutat", f"{loc['city']}, {loc['country']}"),
                            ("Coordenades", f"{loc['latitude']:.3f}, {loc['longitude']:.3f}"),
                        ],
                    )
                    # Text resum clima
                    st.caption("Ubicació aproximada")
                    # Mapa ubicacio clima
                    render_location_map(
                        loc["latitude"],
                        loc["longitude"],
                        zoom=4,
                        color=(255, 0, 0, 220),
                        radius=15000,
                    )

            if not weather_loaded:
                # Avís clima no disponible
                st.info("No hi ha cap fitxer de clima disponible per mostrar en aquesta configuració.")

            # Separador clima energia
            st.markdown("<div class='studio-spacer-070'></div>", unsafe_allow_html=True)
            # Seccio solar bateria
            st.markdown("### Solar i bateria")
            st.caption("Actius energètics detectats a partir de la geometria i dels objectes associats.")

            if pv_list:
                pv_summary = summarize_pv(records, name_index, pv_list)
                surface_names = pv_summary["surfaces_with_pv"]
                # Mostrem només unes quantes superfícies perquè alguns models tenen molts
                # panells i la targeta quedaria il·legible.
                surface_preview = ", ".join(surface_names[:4]) if surface_names else "Sense coincidencies"
                if len(surface_names) > 4:
                    surface_preview += f" +{len(surface_names) - 4} mes"
                # Targeta panells fotovoltaics
                render_detail_card(
                    "Panells fotovoltaics",
                    [
                        ("Generadors PV", str(pv_summary["count"])),
                        ("Cares actives", str(len(surface_names))),
                        ("Area aprox.", "0 m2" if pv_summary["area_m2"] is None else f"{pv_summary['area_m2']:.2f} m2"),
                        ("Tilt mig", "0 deg" if pv_summary["avg_tilt"] is None else f"{pv_summary['avg_tilt']:.0f} deg"),
                        ("Azimut mig", "0 deg" if pv_summary["avg_azimuth"] is None else f"{pv_summary['avg_azimuth']:.0f} deg"),
                        ("Superficies", surface_preview),
                    ],
                )
            else:
                # Avís sense panells PV
                st.info("No s'han detectat panells fotovoltaics associats a superficies.")

            if storage_list:
                # Targetes bateria
                render_battery_cards(storage_list)
            else:
                # Avís sense bateria
                st.info("No s'han detectat objectes de bateria.")

        with col2:
            # Taules kwargs entorn
            render_kwargs_overview(kwargs)

    except Exception as error:
        st.error(f"Error carregant l'entorn: {error}")


render_environment_page()
st.stop()
\end{lstlisting}

\Needspace{8\baselineskip}

\subsection{\texorpdfstring{\texttt{pages/Resultats.py}}{pages/Resultats.py}}

\begin{lstlisting}[style=studioPython,caption={pages/Resultats.py},label={lst:pages-resultats-py}]
"""Streamlit frontend per a la pàgina de resultats de BEMS-RL STUDIO."""

from __future__ import annotations

from pathlib import Path

import streamlit as st
from sidebar_nav import configure_studio_page

from backend.resultats_backend import (
    DashboardData,
    build_results_page_state,
    get_run_artifacts,
    load_dashboard_data,
)
from page_components.resultats_dashboard import render_inline_dashboard
from page_components.ui_fragments import render_hero
from page_styles.resultats import inject_results_page_styles


PAGE_TITLE = "Resultats"
PAGE_LAYOUT = "wide"
INTRODUCTION_TEXT = (
    "Explora els resultats d'una execució RL, descarrega els CSV principals "
    "i obre el dashboard analític."
)

def build_download_filename(prefix: str, run_name: str) -> str:
    """Crea un nom de fitxer de descàrrega estable CSV per a una execució seleccionada."""

    sanitized_run_name = "".join(
        c if c.isalnum() or c in {"-", "_"} else "-" for c in run_name.strip()
    ).strip("-")
    if not sanitized_run_name:
        sanitized_run_name = "entrenament"
    return f"{prefix}-{sanitized_run_name}.csv"


def render_results_page() -> None:
    """Prepara la pàgina de resultats: selector d'execució, dashboard i downloads CSV."""

    configure_studio_page(PAGE_TITLE, layout=PAGE_LAYOUT)
    inject_results_page_styles()

    base_dir = str(Path.cwd())
    render_hero("results-hero", "Analítica i exportació", PAGE_TITLE, INTRODUCTION_TEXT)

    # La cerca de runs i la validació dels CSV queden fora de la UI principal perquè
    # aquesta funció només decideixi què es mostra i quan es carrega el dashboard.
    page_state = build_results_page_state(base_dir)

    # Separador resultats
    st.markdown("<div class='studio-spacer-115'></div>", unsafe_allow_html=True)
    # Titol seleccio resultats
    st.markdown('<h2 class="section-title">Selecció</h2>', unsafe_allow_html=True)

    if page_state.warning_message:
        # Avís resultats
        st.warning(page_state.warning_message)
        st.stop()

    try:
        # Fila seleccio i dashboard
        selection_cols = st.columns([10, 1.15], gap="small", vertical_alignment="center")
    except TypeError:
        # Fila seleccio i dashboard
        selection_cols = st.columns([10, 1.15], gap="small")
    with selection_cols[0]:
        # Selector execucio
        selected_run = st.selectbox("Selecciona l'execució", page_state.run_dirs)

    run_artifacts = get_run_artifacts(selected_run.path)
    if run_artifacts.error_message:
        # Error artefactes resultats
        st.error(run_artifacts.error_message)
        st.stop()

    with selection_cols[1]:
        try:
            # Botó generar dashboard
            generate_requested = st.button(
                " ",
                icon=":material/insert_chart:",
                help="Generar Dashboard",
                key="generate_dashboard_button",
                width="stretch",
            )
        except TypeError:
            # Botó generar dashboard
            generate_requested = st.button(
                "\u25eb",
                help="Generar Dashboard",
                key="generate_dashboard_button",
                width="stretch",
            )

    if generate_requested:
        with st.spinner("Carregant dades del dashboard..."):
            try:
                # Guardem el dashboard a session_state perquè no es torni a carregar en cada
                # interacció menor, com fer download d'un CSV o obrir una pestanya.
                dashboard_data = load_dashboard_data(
                    run_artifacts.progress_path,
                    run_artifacts.observations_path,
                )
                st.session_state["_inline_dashboard_data"] = dashboard_data
                st.session_state["_inline_dashboard_run_path"] = selected_run.path
            except Exception as exc:
                # Error carregar dashboard
                st.error(f"Error carregant el dashboard: {exc}")

    inline_data: DashboardData | None = st.session_state.get("_inline_dashboard_data")
    active_run_path: str | None = st.session_state.get("_inline_dashboard_run_path")

    if inline_data is not None and active_run_path == selected_run.path:
        # Dashboard resultats
        render_inline_dashboard(inline_data, page_state.run_dirs, selected_run.path)

    # Botons download CSV
    download_cols = st.columns(2, gap="large")
    with download_cols[0]:
        # Botó download entrenament
        st.download_button(
            "Descarregar dades d'entrenament",
            data=Path(run_artifacts.progress_path).read_bytes(),
            file_name=build_download_filename("entrenament", selected_run.name),
            mime="text/csv",
            width="stretch",
        )
    with download_cols[1]:
        # Botó download monitor
        st.download_button(
            "Descarregar dades monitor",
            data=Path(run_artifacts.observations_path).read_bytes(),
            file_name=build_download_filename("monitor", selected_run.name),
            mime="text/csv",
            width="stretch",
        )


render_results_page()
\end{lstlisting}

\Needspace{8\baselineskip}

\subsection{\texorpdfstring{\texttt{pages/Simular\_Entorn.py}}{pages/Simular\_Entorn.py}}

\begin{lstlisting}[style=studioPython,caption={pages/Simular\_Entorn.py},label={lst:pages-simular-entorn-py}]
from __future__ import annotations

from functools import partial
from html import escape
import queue
import threading
import time
import traceback

import streamlit as st
from page_components.ui_fragments import render_hero, render_runtime_progress, render_section_card
from page_styles.runtime_pages import inject_baseline_styles
from sidebar_nav import configure_studio_page

from backend.common import ONE_YEAR_STEPS
from backend.simular_entorn_backend import (
    describe_environment,
    list_environment_ids,
    run_baseline_simulation,
)


PAGE_TITLE = "Simular Entorn"
PAGE_LAYOUT = "wide"
INTRODUCTION_TEXT = (
    "Executa un baseline dels entorns amb els schedules normals d'EnergyPlus, "
    "sense control RL, i compara'l amb execucions que ja teniu guardades."
)
def _key(name: str) -> str:
    """Key."""
    return f"baseline_{name}"

def empty_baseline_runtime() -> dict[str, object]:
    """Temps d'execució de base buit."""
    return {
        "running": False,
        "completed": False,
        "cancel_requested": False,
        "needs_full_rerun": False,
        "thread": None,
        "queue": None,
        "stop_event": None,
        "result": None,
        "error": None,
        "traceback": None,
        "env_id": None,
        "steps_target": 1,
        "latest_step": 0,
        "status": "",
        "started_at": None,
        "stopped_at": None,
        "frozen_elapsed_seconds": None,
    }


def ensure_baseline_runtime() -> dict[str, object]:
    """Assegureu-vos el temps d'execució de referència."""
    runtime_key = _key("runtime")
    if runtime_key not in st.session_state:
        st.session_state[runtime_key] = empty_baseline_runtime()
    return st.session_state[runtime_key]


def reset_baseline_runtime() -> dict[str, object]:
    """Restableix el temps d'execució de referència."""
    runtime = ensure_baseline_runtime()
    previous_thread = runtime.get("thread")
    if previous_thread is not None and getattr(previous_thread, "is_alive", lambda: False)():
        return runtime

    st.session_state[_key("runtime")] = empty_baseline_runtime()
    return st.session_state[_key("runtime")]


def drain_baseline_runtime(runtime: dict[str, object]) -> None:
    """Esgota el temps d'execució de la línia base."""
    event_queue = runtime.get("queue")
    if event_queue is None:
        return

    # El worker escriu esdeveniments en una cua i Streamlit els buida a cada rerun.
    # Així la simulació pot continuar en segon pla sense bloquejar la interfície.
    while True:
        try:
            event = event_queue.get_nowait()
        except queue.Empty:
            break

        event_type = str(event.get("type") or "")
        if event_type == "progress":
            runtime["latest_step"] = max(
                int(runtime.get("latest_step") or 0),
                int(event.get("step_number") or 0),
            )
            runtime["steps_target"] = max(
                int(runtime.get("steps_target") or 1),
                int(event.get("total_steps") or 1),
            )
            runtime["status"] = str(event.get("status") or runtime.get("status") or "")
        elif event_type == "result":
            runtime["result"] = event.get("result")
            result = runtime["result"]
            if result is not None and runtime.get("frozen_elapsed_seconds") is None:
                runtime["frozen_elapsed_seconds"] = float(getattr(result, "elapsed_seconds", 0.0))
            runtime["completed"] = True
            runtime["running"] = False
            runtime["needs_full_rerun"] = True
        elif event_type == "error":
            runtime["error"] = str(event.get("error") or "Error desconegut durant la simulacio.")
            runtime["traceback"] = event.get("traceback")
            runtime["completed"] = True
            runtime["running"] = False
            runtime["needs_full_rerun"] = True
        elif event_type == "finished":
            runtime["completed"] = True
            runtime["running"] = False
            runtime["thread"] = None
            runtime["needs_full_rerun"] = True

    thread = runtime.get("thread")
    if thread is not None and not thread.is_alive() and runtime.get("running"):
        runtime["running"] = False
        runtime["completed"] = True
        runtime["thread"] = None
        runtime["needs_full_rerun"] = True


def sync_baseline_runtime() -> dict[str, object]:
    """Sincronitza el temps d'execució de referència."""
    runtime = ensure_baseline_runtime()
    drain_baseline_runtime(runtime)
    return runtime


def consume_baseline_rerun_flag(runtime: dict[str, object] | None = None) -> bool:
    """Consumeix la bandera de repetició de la línia de base."""
    active_runtime = runtime if runtime is not None else ensure_baseline_runtime()
    if active_runtime.get("needs_full_rerun") and not active_runtime.get("running"):
        active_runtime["needs_full_rerun"] = False
        return True
    return False


def request_baseline_stop(runtime: dict[str, object] | None = None) -> bool:
    """Sol·licita l'aturada de la línia de base."""
    active_runtime = runtime if runtime is not None else ensure_baseline_runtime()
    if not active_runtime.get("running") or active_runtime.get("stop_event") is None:
        return False

    active_runtime["cancel_requested"] = True
    active_runtime["status"] = "Aturant simulacio..."
    active_runtime["stopped_at"] = time.time()
    started_at = active_runtime.get("started_at")
    if started_at:
        active_runtime["frozen_elapsed_seconds"] = max(0.0, float(active_runtime["stopped_at"]) - float(started_at))
    active_runtime["stop_event"].set()
    return True


def push_baseline_progress_event(
    event_queue: "queue.Queue[dict[str, object]]",
    job_config: dict[str, object],
    step_number: int,
    total_steps: int,
) -> None:
    """Envieu un esdeveniment de progrés de referència a la cua de temps d'execució de la pàgina."""
    event_queue.put(
        {
            "type": "progress",
            "step_number": int(step_number or 0),
            "total_steps": int(total_steps or job_config["steps_target"]),
            "status": f"Passos: {int(step_number or 0)}/{int(total_steps or job_config['steps_target'])}",
        }
    )


def baseline_worker(
    job_config: dict[str, object],
    event_queue: "queue.Queue[dict[str, object]]",
    stop_event: threading.Event,
) -> None:
    """Executa la simulació baseline en segon pla."""
    try:
        result = run_baseline_simulation(
            env_id=str(job_config["env_id"]),
            steps_target=int(job_config["steps_target"]),
            seed=int(job_config["seed"]) if job_config.get("seed") is not None else None,
            should_stop=stop_event.is_set,
            on_progress=partial(push_baseline_progress_event, event_queue, job_config),
        )
        event_queue.put({"type": "result", "result": result})
    except Exception as exc:
        event_queue.put(
            {
                "type": "error",
                "error": str(exc),
                "traceback": traceback.format_exc(),
            }
        )
    finally:
        event_queue.put({"type": "finished"})


def start_baseline_run(request: dict[str, object]) -> dict[str, object]:
    """Inicia l'execució de la línia de base."""
    runtime = reset_baseline_runtime()
    response: dict[str, object] = {
        "started": False,
        "errors": [],
        "runtime": runtime,
    }

    if runtime.get("running"):
        response["errors"].append("Ja hi ha una simulacio en curs.")
        return response

    env_id = str(request.get("env_id") or "").strip()
    if not env_id:
        response["errors"].append("Selecciona un entorn valid abans d'iniciar la simulacio.")
        return response

    if env_id not in list_environment_ids():
        response["errors"].append("L'entorn seleccionat no esta registrat a Sinergym.")
        return response

    event_queue: "queue.Queue[dict[str, object]]" = queue.Queue()
    stop_event = threading.Event()
    job_config = {
        "env_id": env_id,
        "steps_target": int(request.get("steps_target") or ONE_YEAR_STEPS),
        "seed": int(request.get("seed")) if request.get("seed") is not None else None,
    }
    worker_thread = threading.Thread(
        target=baseline_worker,
        args=(job_config, event_queue, stop_event),
        daemon=True,
    )

    runtime.update(
        {
            "running": True,
            "completed": False,
            "cancel_requested": False,
            "needs_full_rerun": False,
            "thread": worker_thread,
            "queue": event_queue,
            "stop_event": stop_event,
            "result": None,
            "error": None,
            "traceback": None,
            "env_id": env_id,
            "steps_target": int(job_config["steps_target"]),
            "latest_step": 0,
            "status": "Inicialitzant simulacio baseline...",
            "started_at": time.time(),
            "stopped_at": None,
            "frozen_elapsed_seconds": None,
        }
    )
    worker_thread.start()

    response["started"] = True
    response["runtime"] = runtime
    return response


def render_environment_summary(env_id: str) -> None:
    """Mostra el resum de l'entorn a la UI de Streamlit."""
    summary = describe_environment(env_id)
    step_text = f"{summary.step_minutes:g} min" if summary.step_minutes else "No disponible"

    # Metriques resum entorn
    info_cols = st.columns(4)
    with info_cols[0]:
        # Metrica variables observades
        st.metric("Variables observades", summary.variables_count)
    with info_cols[1]:
        # Metrica meters
        st.metric("Meters", summary.meters_count)
    with info_cols[2]:
        # Metrica actuadors RL
        st.metric("Actuadors RL", summary.actuators_count)
    with info_cols[3]:
        # Metrica pas simulacio
        st.metric("Pas de simulacio", step_text)

    # Etiquetes fitxers baseline
    st.markdown(
        (
            '<div class="baseline-chip-row">'
            f'<span class="baseline-chip"><strong>Building</strong> {escape(summary.building_file or "No disponible")}</span>'
            f'<span class="baseline-chip"><strong>Weathers</strong> {summary.weather_count}</span>'
            f'<span class="baseline-chip"><strong>Mode</strong> schedules EnergyPlus</span>'
            f'<span class="baseline-chip"><strong>Resolucio</strong> pas fix de {step_text}</span>'
            "</div>"
        ),
        unsafe_allow_html=True,
    )


def render_completed_baseline(runtime: dict[str, object]) -> None:
    """Mostra la secció de referència completada a la UI de Streamlit."""
    result = runtime.get("result")
    if result is None or runtime.get("running"):
        return

    if bool(getattr(result, "cancelled", False)):
        # Avís baseline cancel·lat
        st.warning(
            f"Simulacio aturada manualment despres de {float(getattr(result, 'elapsed_seconds', 0.0)):.2f} s "
            f"i {int(getattr(result, 'completed_steps', 0) or 0)} passos."
        )
    else:
        # Missatge baseline finalitzat
        st.success(
            f"Simulacio finalitzada en {float(getattr(result, 'elapsed_seconds', 0.0)):.2f} s "
            f"i {int(getattr(result, 'completed_steps', 0) or 0)} passos."
        )

    run_path = getattr(result, "run_path", "")
    if run_path:
        # Text ruta baseline
        st.caption(f"Run guardada a: {run_path}")
    # Enllaç resultats baseline
    st.page_link("pages/Resultats.py", label="Veure resultats a Resultats", icon=":material/insights:")


def render_baseline_runtime(runtime: dict[str, object]) -> None:
    """Mostra la secció de temps d'execució de la línia de base a la UI de Streamlit."""
    render_section_card(
        "Estat de la simulacio",
        "Segueix el progres del baseline i obre les visualitzacions quan la run estigui llesta.",
        title_class="section-title",
        card_class="section-card",
    )
    render_runtime_progress(runtime, progress_label="Progres", freeze_from_result=True)

    if runtime.get("cancel_requested") and runtime.get("running"):
        # Avís aturada baseline
        st.warning("S'ha demanat l'aturada. La simulacio es tancara al proper punt segur.")

    if runtime.get("error"):
        # Error runtime baseline
        st.error(str(runtime.get("error")))
        if runtime.get("traceback"):
            # Acordio traceback baseline
            with st.expander("Traceback intern", expanded=False):
                st.code(str(runtime.get("traceback")), language="text")
        return

    render_completed_baseline(runtime)


def render_baseline_runtime_frame() -> None:
    """Mostra el marc de runtime del baseline a la UI de Streamlit."""
    runtime = sync_baseline_runtime()
    if runtime.get("running") or runtime.get("completed") or runtime.get("error") or runtime.get("result"):
        st.divider()
        render_baseline_runtime(runtime)

    if consume_baseline_rerun_flag(runtime):
        st.rerun()


render_baseline_runtime_fragment = st.fragment(run_every="1s")(render_baseline_runtime_frame)


configure_studio_page(PAGE_TITLE, layout=PAGE_LAYOUT)
inject_baseline_styles()

environment_ids = list_environment_ids()
if not environment_ids:
    # Error sense entorns baseline
    st.error("No hi ha entorns Sinergym registrats disponibles per executar un baseline.")
    st.stop()

st.session_state.setdefault(_key("selected_env"), environment_ids[0])
st.session_state.setdefault(_key("steps"), ONE_YEAR_STEPS)
st.session_state.setdefault(_key("seed"), 0)

render_hero("results-hero", "Simulacio baseline", PAGE_TITLE, INTRODUCTION_TEXT)

runtime = sync_baseline_runtime()
controls_disabled = bool(runtime.get("running"))

# Separador baseline
st.markdown("<div class='studio-spacer-115'></div>", unsafe_allow_html=True)
# Titol configuracio baseline
st.markdown('<h2 class="section-title">Configuracio del baseline</h2>', unsafe_allow_html=True)

# Selector entorn baseline
selected_env = st.selectbox(
    "Entorn Sinergym",
    options=environment_ids,
    index=environment_ids.index(st.session_state[_key("selected_env")]) if st.session_state[_key("selected_env")] in environment_ids else 0,
    disabled=controls_disabled,
)
st.session_state[_key("selected_env")] = selected_env

render_environment_summary(selected_env)

# Inputs configuracio baseline
config_cols = st.columns(2, gap="medium")
with config_cols[0]:
    # Input limit passos
    steps_limit = st.number_input(
        "Limit de passos",
        min_value=1,
        max_value=5_000_000,
        value=int(st.session_state[_key("steps")]),
        step=35040,
        disabled=controls_disabled,
    )
with config_cols[1]:
    # Input seed baseline
    seed_value = st.number_input(
        "Seed inicial",
        min_value=0,
        max_value=999_999,
        value=int(st.session_state[_key("seed")]),
        step=1,
        disabled=controls_disabled,
    )

# Nota recomanacio passos
st.markdown(
    f'<div class="baseline-inline-note">Recomanacio: usa {ONE_YEAR_STEPS} passos per generar una passada anual comparable amb una run d\'agent.</div>',
    unsafe_allow_html=True,
)

st.divider()
# Titol control baseline
st.markdown('<h2 class="section-title">Control</h2>', unsafe_allow_html=True)

# Botons control baseline
run_cols = st.columns(2, gap="medium")
with run_cols[0]:
    # Botó iniciar baseline
    start_requested = st.button(
        "Iniciar simulacio baseline",
        type="primary",
        width="stretch",
        disabled=controls_disabled,
    )
with run_cols[1]:
    # Botó cancel·lar baseline
    cancel_requested = st.button(
        "Cancelar",
        width="stretch",
        disabled=not bool(runtime.get("running")),
    )

if cancel_requested and request_baseline_stop(runtime):
    st.rerun()

if start_requested:
    start_result = start_baseline_run(
        {
            "env_id": selected_env,
            "steps_target": int(steps_limit),
            "seed": int(seed_value),
        }
    )
    for error in start_result.get("errors") or []:
        # Error inici baseline
        st.error(str(error))
    if start_result.get("started"):
        st.rerun()

render_baseline_runtime_fragment()
\end{lstlisting}

\Needspace{8\baselineskip}

\subsection{\texorpdfstring{\texttt{pages/Visor\_Climes\_EPW.py}}{pages/Visor\_Climes\_EPW.py}}

\begin{lstlisting}[style=studioPython,caption={pages/Visor\_Climes\_EPW.py},label={lst:pages-visor-climes-epw-py}]
from __future__ import annotations

from dataclasses import dataclass

import pandas as pd
import streamlit as st

from page_components.ui_fragments import (
    render_copy_block,
    render_detail_list_card,
    render_hero,
    render_metric_card_grid,
)
from page_styles.visor_climes_epw import inject_epw_viewer_styles
from sidebar_nav import configure_studio_page

from backend import visor_epw_backend as epw_backend
from backend.epw_figures import (
    build_active_month_axis,
    build_annual_heatmap_figure,
    build_comfort_radiation_figure,
    build_daily_temperature_band_figure,
    build_focus_timeseries_figure,
    build_heatmap_figure,
    build_hourly_profile_figure,
    build_monthly_climate_figure,
    build_monthly_solar_figure,
    build_monthly_wind_rose_grid_figure,
    ui_text,
)
from backend.mapa_backend import render_location_map


PAGE_TITLE = "Visor de Climes EPW"
PAGE_LAYOUT = "wide"
INTRODUCTION_TEXT = (
    "Selecciona un fitxer climàtic EPW i explora'n les sèries temporals, "
    "els patrons mensuals i els indicadors ambientals des d'un dashboard integrat."
)


PLOTLY_CHART_CONFIG = {"displayModeBar": False}


@dataclass(frozen=True)
class EpwDashboardState:
    """Dades preparades necessàries per representar les pestanyes del panell EPW."""

    selected_entry: dict[str, str]
    metadata: dict[str, object]
    full_metrics: dict[str, float]
    filtered_metrics: dict[str, float]
    month_tick_values: list[int]
    month_tick_labels: list[str]
    series_variable: str
    series_aggregation: str
    heatmap_variable: str
    secondary_heatmap_variable: str
    records_message: str
    monthly_frame: pd.DataFrame
    focus_series: pd.DataFrame
    daily_temperature: pd.DataFrame
    hourly_profile: pd.DataFrame
    heatmap_frame: pd.DataFrame
    radiation_heatmap_frame: pd.DataFrame
    annual_temperature_heatmap: pd.DataFrame
    annual_direct_radiation_heatmap: pd.DataFrame
    annual_wind_heatmap: pd.DataFrame
    annual_humidity_heatmap: pd.DataFrame
    comfort_radiation_frame: pd.DataFrame
    monthly_wind_rose_tables: dict[str, pd.DataFrame]
    download_frame: pd.DataFrame
    metric_cards: list[tuple[str, str]]


def format_number(value: float | int | None, *, decimals: int = 1, suffix: str = "") -> str:
    """Formata un valor numèric amb separadors catalans i un sufix opcional."""

    if value is None or pd.isna(value):
        return "-"
    if decimals == 0:
        return f"{int(round(float(value))):,}{suffix}".replace(",", ".")
    return (
        f"{float(value):,.{decimals}f}{suffix}"
        .replace(",", "X")
        .replace(".", ",")
        .replace("X", ".")
    )


def data_frame_has_plot_values(
    data_frame: pd.DataFrame,
    columns: list[str] | None = None,
    *,
    require_non_zero: bool = False,
) -> bool:
    """Retorna si un DataFrame conté valors numèrics que val la pena dibuixar."""

    if data_frame.empty:
        return False

    if columns is None:
        values = data_frame
    else:
        available_columns = [column for column in columns if column in data_frame.columns]
        if not available_columns:
            return False
        values = data_frame[available_columns]

    numeric_values = values.select_dtypes(include="number")
    if numeric_values.empty:
        numeric_values = values.apply(pd.to_numeric, errors="coerce")

    valid_values = numeric_values.stack().dropna()
    if valid_values.empty:
        return False

    if require_non_zero:
        return bool((valid_values.abs() > 0).any())

    return True


def wind_rose_tables_have_plot_values(monthly_rose_tables: dict[str, pd.DataFrame]) -> bool:
    """Indica si almenys una taula mensual de roses dels vents conté freqüències de vent."""

    for rose_frame in monthly_rose_tables.values():
        if data_frame_has_plot_values(rose_frame, ["frequency_pct"], require_non_zero=True):
            return True
    return False


def render_detail_card(
    metadata: dict[str, object],
    selected_entry: dict[str, str],
    total_records: int,
    filtered_records: int,
) -> None:
    """Mostra la targeta de metadades del fitxer EPW seleccionada a la UI de Streamlit."""

    detail_rows = [
        ("Fitxer", ui_text(metadata.get("file_name", "-"))),
        ("Ruta relativa", ui_text(selected_entry.get("relative_path", "-"))),
        ("Ciutat", ui_text(metadata.get("city", "-"))),
        ("País", ui_text(metadata.get("country", "-"))),
        ("Font", ui_text(metadata.get("source", "-"))),
        ("WMO", ui_text(metadata.get("wmo", "-"))),
        ("Latitud", format_number(metadata.get("latitude"), decimals=3)),
        ("Longitud", format_number(metadata.get("longitude"), decimals=3)),
        ("Fus horari", f"UTC {format_number(metadata.get('timezone'), decimals=1)}"),
        ("Elevació", f"{format_number(metadata.get('elevation_m'), decimals=0)} m"),
        ("Registres", f"{filtered_records:,} / {total_records:,}".replace(",", ".")),
        ("Classificació", ui_text(metadata.get("climate_family", "-"))),
    ]
    render_detail_list_card(
        title="Context del fitxer",
        rows=detail_rows,
        card_class="epw-detail-card",
        title_class="epw-card-title",
        list_class="epw-detail-list",
        row_class="epw-detail-row",
        label_class="epw-detail-label",
        value_class="epw-detail-value",
    )


def build_map_chart(metadata: dict[str, object]) -> None:
    """Dibuixa un mapa centrat a les coordenades de l'estació meteorològica EPW."""

    latitude = metadata.get("latitude")
    longitude = metadata.get("longitude")
    try:
        latitude = float(latitude)
        longitude = float(longitude)
    except (TypeError, ValueError):
        # Avís mapa EPW
        st.info("Aquest fitxer EPW no té coordenades vàlides per mostrar al mapa.")
        return

    city = ui_text(metadata.get("city", "-"))
    country = ui_text(metadata.get("country", "-"))
    file_name = ui_text(metadata.get("file_name", "-"))
    render_location_map(
        latitude,
        longitude,
        zoom=4.5,
        color=(95, 84, 249, 220),
        radius=22000,
        tooltip_html=f"<b>{file_name}</b><br/>{city}, {country}",
    )


def configure_epw_dashboard_page() -> None:
    """Configura Chrome de la pàgina global i renderitza l'heroi EPW."""

    configure_studio_page(PAGE_TITLE, layout=PAGE_LAYOUT)
    inject_epw_viewer_styles()
    render_hero("epw-hero", "Visualització climàtica", PAGE_TITLE, ui_text(INTRODUCTION_TEXT))


def select_epw_entry() -> tuple[dict[str, str], dict[str, object], pd.DataFrame]:
    """Mostra els controls de cerca i carrega el paquet EPW seleccionat."""

    weather_root = epw_backend.WEATHER_ROOT
    catalog = epw_backend.list_epw_catalog(str(weather_root))
    if not catalog:
        # Avís sense EPW
        st.warning(f"No s'han trobat fitxers EPW a `{weather_root}`.")
        st.stop()

    # Controls cerca EPW
    search_cols = st.columns([1.15, 2.1], gap="medium")
    with search_cols[0]:
        # Input cerca EPW
        search_query = st.text_input(
            "Cerca un fitxer EPW",
            placeholder="Ex: Granada, Lisboa, IWEC...",
        ).strip().lower()

    filtered_catalog = tuple(
        entry
        for entry in catalog
        if not search_query or search_query in entry["search_blob"]
    )
    if not filtered_catalog:
        # Avís filtre EPW buit
        st.warning("Cap fitxer EPW coincideix amb el filtre de cerca.")
        st.stop()

    entry_by_path = {entry["path"]: entry for entry in filtered_catalog}
    with search_cols[1]:
        # Selector fitxer EPW
        selected_path = st.selectbox(
            "Fitxer climàtic",
            list(entry_by_path),
            format_func=lambda option: ui_text(entry_by_path[option]["display_name"]),
        )

    bundle = epw_backend.load_epw_bundle(selected_path)
    return entry_by_path[selected_path], bundle["metadata"], bundle["data"]


def build_metric_card_values(filtered_metrics: dict[str, float]) -> list[tuple[str, str]]:
    """Crea valors de visualització per a les targetes EPW de resum KPI."""

    return [
        (
            "Temperatura mitjana",
            f"{format_number(filtered_metrics['mean_temp'], decimals=1)} {epw_backend.DEG_C}",
        ),
        (
            "Interval tèrmic",
            f"{format_number(filtered_metrics['min_temp'], decimals=1)} - "
            f"{format_number(filtered_metrics['max_temp'], decimals=1)} {epw_backend.DEG_C}",
        ),
        ("Humitat mitjana", f"{format_number(filtered_metrics['mean_rh'], decimals=1)} %"),
        ("Vent mitjà", f"{format_number(filtered_metrics['mean_wind'], decimals=2)} m/s"),
        (
            "GHI acumulada",
            f"{format_number(filtered_metrics['ghi_total_kwh_m2'], decimals=1)} {epw_backend.KWH_PER_M2}",
        ),
        ("Vent màxim", f"{format_number(filtered_metrics['max_wind'], decimals=2)} m/s"),
        (
            "DNI acumulada",
            f"{format_number(filtered_metrics['dni_total_kwh_m2'], decimals=1)} {epw_backend.KWH_PER_M2}",
        ),
        (
            "DHI acumulada",
            f"{format_number(filtered_metrics['dhi_total_kwh_m2'], decimals=1)} {epw_backend.KWH_PER_M2}",
        ),
    ]


def prepare_epw_dashboard_state(
    selected_entry: dict[str, str],
    metadata: dict[str, object],
    epw_data: pd.DataFrame,
) -> EpwDashboardState:
    """Crea totes les taules i etiquetes derivades necessàries per al panell EPW."""

    filtered_data = epw_data.copy()
    full_metrics = epw_backend.summarize_epw_metrics(epw_data)
    filtered_metrics = epw_backend.summarize_epw_metrics(filtered_data)
    month_tick_values, month_tick_labels = build_active_month_axis(filtered_data)

    series_variable = "dry_bulb_c"
    series_aggregation = "Dia"
    heatmap_variable = "dry_bulb_c"
    secondary_heatmap_variable = "direct_normal_radiation_wh_m2"
    records_message = (
        f"Mostrant {int(filtered_metrics['records']):,} registres de "
        f"{int(full_metrics['records']):,} sobre l'any EPW complet."
    ).replace(",", ".")

    return EpwDashboardState(
        selected_entry=selected_entry,
        metadata=metadata,
        full_metrics=full_metrics,
        filtered_metrics=filtered_metrics,
        month_tick_values=month_tick_values,
        month_tick_labels=month_tick_labels,
        series_variable=series_variable,
        series_aggregation=series_aggregation,
        heatmap_variable=heatmap_variable,
        secondary_heatmap_variable=secondary_heatmap_variable,
        records_message=records_message,
        monthly_frame=epw_backend.build_monthly_summary(filtered_data),
        focus_series=epw_backend.aggregate_epw_timeseries(
            filtered_data,
            series_variable,
            series_aggregation,
        ),
        daily_temperature=epw_backend.build_daily_temperature_profile(filtered_data),
        hourly_profile=epw_backend.build_hourly_profile(filtered_data),
        heatmap_frame=epw_backend.build_heatmap_table(filtered_data, heatmap_variable),
        radiation_heatmap_frame=epw_backend.build_heatmap_table(
            filtered_data,
            secondary_heatmap_variable,
        ),
        annual_temperature_heatmap=epw_backend.build_annual_hourly_heatmap_table(
            filtered_data,
            "dry_bulb_c",
        ),
        annual_direct_radiation_heatmap=epw_backend.build_annual_hourly_heatmap_table(
            filtered_data,
            "direct_normal_radiation_wh_m2",
        ),
        annual_wind_heatmap=epw_backend.build_annual_hourly_heatmap_table(
            filtered_data,
            "wind_speed_m_s",
        ),
        annual_humidity_heatmap=epw_backend.build_annual_hourly_heatmap_table(
            filtered_data,
            "relative_humidity_pct",
        ),
        comfort_radiation_frame=epw_backend.build_comfort_radiation_table(filtered_data, metadata),
        monthly_wind_rose_tables=epw_backend.build_monthly_wind_rose_tables(filtered_data),
        download_frame=epw_backend.build_download_frame(filtered_data),
        metric_cards=build_metric_card_values(filtered_metrics),
    )


def render_tab_intro(text: str) -> None:
    """Mostra un text breu al començament de cada pestanya EPW."""

    render_copy_block("epw-tab-copy", text, formatter=ui_text)


def render_annual_heatmap(
    frame: pd.DataFrame,
    variable_key: str,
    state: EpwDashboardState,
    empty_message: str,
    *,
    require_non_zero: bool = False,
) -> None:
    """Mostra un mapa de calor anual quan la taula d'origen té valors útils."""
    render_epw_chart(
        frame,
        None,
        lambda: build_annual_heatmap_figure(
            frame,
            variable_key,
            tick_values=state.month_tick_values,
            tick_labels=state.month_tick_labels,
        ),
        empty_message,
        require_non_zero=require_non_zero,
    )


def render_epw_chart(
    frame: pd.DataFrame,
    columns: list[str] | None,
    figure_builder,
    empty_message: str,
    *,
    require_non_zero: bool = False,
) -> None:
    """Mostra un gràfic EPW quan les dades tenen valors, o un avís si no."""
    if data_frame_has_plot_values(frame, columns, require_non_zero=require_non_zero):
        # Grafic EPW
        st.plotly_chart(
            figure_builder(),
            width="stretch",
            config=PLOTLY_CHART_CONFIG,
        )
    else:
        # Avís grafic EPW buit
        st.info(empty_message)


def render_overview_tab(state: EpwDashboardState) -> None:
    """Mostra la pestanya de visió general de l'EPW a la UI de Streamlit."""

    render_tab_intro(
        "Resum executiu del fitxer seleccionat, amb indicadors clau, context i ubicació."
    )
    render_metric_card_grid(
        state.metric_cards,
        shell_class="section-card epw-metric-shell",
        card_class="epw-metric-card",
        label_class="epw-metric-label",
        value_class="epw-metric-value",
        formatter=ui_text,
        label_formatter=ui_text,
    )
    # Columnes resum EPW
    overview_cols = st.columns([1.18, 0.82], gap="large")
    with overview_cols[0]:
        render_detail_card(
            state.metadata,
            state.selected_entry,
            total_records=int(state.full_metrics["records"]),
            filtered_records=int(state.filtered_metrics["records"]),
        )
    with overview_cols[1]:
        # Titol mapa EPW
        st.markdown("<div class='epw-card-title'>Ubicació</div>", unsafe_allow_html=True)
        build_map_chart(state.metadata)


def render_climate_dashboard_tab(state: EpwDashboardState) -> None:
    """Mostra mapes de calor climàtics anuals i gràfics d'orientació/vents."""

    render_tab_intro(
        "Tauler visual anual per detectar patrons ràpids de temperatura, radiació, vent i humitat."
    )
    # Graella heatmaps EPW
    dashboard_heatmaps = st.columns(2, gap="medium")
    with dashboard_heatmaps[0]:
        render_annual_heatmap(
            state.annual_temperature_heatmap,
            "dry_bulb_c",
            state,
            "No hi ha prou dades de temperatura per construir el mapa anual.",
        )
        render_annual_heatmap(
            state.annual_direct_radiation_heatmap,
            "direct_normal_radiation_wh_m2",
            state,
            "No hi ha prou dades de radiació directa per construir el mapa anual.",
            require_non_zero=True,
        )

    with dashboard_heatmaps[1]:
        render_annual_heatmap(
            state.annual_wind_heatmap,
            "wind_speed_m_s",
            state,
            "No hi ha prou dades de vent per construir el mapa anual.",
            require_non_zero=True,
        )
        render_annual_heatmap(
            state.annual_humidity_heatmap,
            "relative_humidity_pct",
            state,
            "No hi ha prou dades d'humitat per construir el mapa anual.",
        )

    render_comfort_radiation_chart(state)
    render_monthly_wind_rose_chart(state)


def render_comfort_radiation_chart(state: EpwDashboardState) -> None:
    """Genera un gràfic de radiació orientat a la comoditat quan hi hagi prou dades a la UI de Streamlit."""

    if not data_frame_has_plot_values(
        state.comfort_radiation_frame,
        ["comfort_radiation_kwh_m2", "incident_radiation_kwh_m2"],
        require_non_zero=True,
    ):
        # Avís radiacio orientada
        st.info("No hi ha prou dades per estimar la radiació orientada d'aquest fitxer.")
        return

    # Grafic radiacio orientada
    st.plotly_chart(
        build_comfort_radiation_figure(state.comfort_radiation_frame),
        width="stretch",
        config=PLOTLY_CHART_CONFIG,
    )
    # Text radiacio orientada
    st.caption(
        "Estimació sobre façanes verticals a partir de la radiació EPW "
        f"i una banda de confort tèrmic de 18-24 {epw_backend.DEG_C}."
    )


def render_monthly_wind_rose_chart(state: EpwDashboardState) -> None:
    """Mostra les roses dels vents mensuals quan hi ha dades de vent suficients."""

    if wind_rose_tables_have_plot_values(state.monthly_wind_rose_tables):
        # Grafic rosa vents mensual
        st.plotly_chart(
            build_monthly_wind_rose_grid_figure(state.monthly_wind_rose_tables),
            width="stretch",
            config=PLOTLY_CHART_CONFIG,
        )
    else:
        # Avís rosa vents
        st.info("No hi ha prou dades de vent per construir les roses mensuals.")


def render_series_tab(state: EpwDashboardState) -> None:
    """Mostra la pestanya de sèries temporals EPW a la UI de Streamlit."""

    render_tab_intro(
        "Evolució temporal anual de la temperatura, amb suport de rangs diaris i perfil horari mitjà."
    )
    render_epw_chart(
        state.focus_series,
        [state.series_variable],
        lambda: build_focus_timeseries_figure(
            state.focus_series,
            state.series_variable,
            state.series_aggregation,
        ),
        "No hi ha prou dades de temperatura per construir la sèrie temporal.",
    )

    # Columnes series EPW
    series_cols = st.columns(2, gap="medium")
    with series_cols[0]:
        render_epw_chart(
            state.daily_temperature,
            ["dry_bulb_min", "dry_bulb_mean", "dry_bulb_max"],
            lambda: build_daily_temperature_band_figure(state.daily_temperature),
            "No hi ha prou dades per construir la banda diària de temperatura.",
        )

    with series_cols[1]:
        render_epw_chart(
            state.hourly_profile,
            ["dry_bulb_mean", "dew_point_mean", "relative_humidity_mean"],
            lambda: build_hourly_profile_figure(state.hourly_profile),
            "No hi ha prou dades per construir el perfil horari mitjà.",
        )


def render_patterns_tab(state: EpwDashboardState) -> None:
    """Mostra els patrons climàtics mensuals i intradia."""

    render_tab_intro("Comparatives mensuals i patrons intradiaris del fitxer climàtic.")
    # Columnes patrons mensuals
    pattern_cols_monthly = st.columns(2, gap="medium")
    with pattern_cols_monthly[0]:
        render_epw_chart(
            state.monthly_frame,
            ["dry_bulb_mean", "dry_bulb_min", "dry_bulb_max", "relative_humidity_mean"],
            lambda: build_monthly_climate_figure(state.monthly_frame),
            "No hi ha prou dades climàtiques mensuals per construir aquest gràfic.",
        )

    with pattern_cols_monthly[1]:
        render_epw_chart(
            state.monthly_frame,
            [
                "global_horizontal_radiation_kwh_m2",
                "direct_normal_radiation_kwh_m2",
                "diffuse_horizontal_radiation_kwh_m2",
            ],
            lambda: build_monthly_solar_figure(state.monthly_frame),
            "No hi ha prou dades de radiació mensual per construir aquest gràfic.",
            require_non_zero=True,
        )

    render_pattern_heatmaps(state)


def render_pattern_heatmaps(state: EpwDashboardState) -> None:
    """Mostra els heatmaps de la pestanya de patrons."""

    # Columnes heatmaps patrons
    pattern_cols_heatmaps = st.columns(2, gap="medium")
    with pattern_cols_heatmaps[0]:
        render_epw_chart(
            state.heatmap_frame,
            None,
            lambda: build_heatmap_figure(state.heatmap_frame, state.heatmap_variable),
            "No hi ha prou dades de temperatura per construir el mapa de calor.",
        )

    with pattern_cols_heatmaps[1]:
        render_epw_chart(
            state.radiation_heatmap_frame,
            None,
            lambda: build_heatmap_figure(
                state.radiation_heatmap_frame,
                state.secondary_heatmap_variable,
            ),
            "No hi ha prou dades de radiació directa per construir el mapa de calor.",
            require_non_zero=True,
        )


def render_epw_tabs(state: EpwDashboardState) -> None:
    """Munta totes les pestanyes del panell EPW."""

    # Pestanyes visor EPW
    tab_overview, tab_dashboard, tab_series, tab_patterns = st.tabs(
        ["Resum", "Tauler climàtic", "Sèries", "Patrons"]
    )
    with tab_overview:
        render_overview_tab(state)
    with tab_dashboard:
        render_climate_dashboard_tab(state)
    with tab_series:
        render_series_tab(state)
    with tab_patterns:
        render_patterns_tab(state)


def render_epw_download(state: EpwDashboardState) -> None:
    """Mostra els controls de baixada CSV per al fitxer EPW seleccionat."""

    # Separador download EPW
    st.markdown("<div class='studio-spacer-125'></div>", unsafe_allow_html=True)
    # Botó download EPW
    st.download_button(
        "Descarregar dades del fitxer (CSV)",
        data=state.download_frame.to_csv(index=False).encode("utf-8"),
        file_name=f"{state.metadata.get('stem', 'clima')}_complet.csv",
        mime="text/csv",
        width="stretch",
        type="primary",
    )
    # Text registres EPW
    st.caption(state.records_message)


def render_epw_dashboard() -> None:
    """Mostra el flux complet del visor EPW a la UI de Streamlit."""

    configure_epw_dashboard_page()
    selected_entry, metadata, epw_data = select_epw_entry()
    # Titol lectura global
    st.markdown(
        "<div class='epw-card-title epw-global-title'>Lectura global del fitxer</div>",
        unsafe_allow_html=True,
    )
    state = prepare_epw_dashboard_state(selected_entry, metadata, epw_data)
    render_epw_tabs(state)
    render_epw_download(state)


render_epw_dashboard()
\end{lstlisting}

\Needspace{8\baselineskip}

\subsection{\texorpdfstring{\texttt{pages/\_\_init\_\_.py}}{pages/\_\_init\_\_.py}}

\begin{lstlisting}[style=studioPython,caption={pages/\_\_init\_\_.py},label={lst:pages---init---py}]
"""Streamlit Marcador de paquet de pàgina per a BEMS-RL Studio."""
\end{lstlisting}

\section{Codi: Components reutilitzables}

\Needspace{8\baselineskip}

\subsection{\texorpdfstring{\texttt{page\_components/\_\_init\_\_.py}}{page\_components/\_\_init\_\_.py}}

\begin{lstlisting}[style=studioPython,caption={page\_components/\_\_init\_\_.py},label={lst:page-components---init---py}]
"""Components Streamlit reutilitzables per a pàgines BEMS-RL Studio.

Grup de mòduls de components repetits UI seccions com ara formularis de entrenament, recompensa
controls de configuració, editors de wrappers i panells de resultats en línia. Es conserven
separat dels punts d'entrada de la pàgina de manera que es puguin documentar fragments complexos UI i
reutilitzat sense importar pàgines completes de Streamlit.
"""
\end{lstlisting}

\Needspace{8\baselineskip}

\subsection{\texorpdfstring{\texttt{page\_components/resultats\_dashboard.py}}{page\_components/resultats\_dashboard.py}}

\begin{lstlisting}[style=studioPython,caption={page\_components/resultats\_dashboard.py},label={lst:page-components-resultats-dashboard-py}]
"""Component del panell en línia per a la pàgina de resultats."""

from __future__ import annotations

from html import escape
from pathlib import Path
from typing import NamedTuple

import pandas as pd
import plotly.graph_objects as go
import streamlit as st

from backend.grafics.battery import (
    make_battery_charge_with_price_plot,
    make_battery_power_plot,
    make_battery_soc_plot,
    make_battery_vs_grid_plot,
    make_energy_price_plot,
)
from backend.grafics.comfort import (
    _fix_colors_for_visibility,
    make_comfort_compliance,
    make_violation_bars,
    make_violation_single_line,
)
from backend.grafics.control import make_agent_actions_plot, make_radiant_control_plot
from backend.grafics.episode import make_episode_metrics
from backend.grafics.heatmaps import make_heatmap, make_zone_temperature_heatmaps
from backend.grafics.hvac import make_hvac_consumption_plot, make_hvac_meter_breakdown_plot
from backend.grafics.style import style_figure_semantics
from backend.grafics.thermal import (
    make_indoor_humidity_plot,
    make_indoor_temperature_plot,
    make_setpoints_plot,
    make_setpoints_vs_indoor_plot,
)
from backend.resultats_backend import (
    DashboardData,
    add_comfort_percentage_kpi,
    apply_real_period_axis,
    build_real_period_options,
    has_action_figure_data,
    has_figure_data,
    is_radiant_run,
    load_comparison_data,
    overlay_comparison_traces,
    select_actions_for_obs,
    select_radiant_action_data,
    slice_obs_for_real_period,
)
from backend.resultats_report_backend import generate_report_bytes
_AGG_LABELS = {
    "hour": "Horaria (0..23)",
    "day": "Diaria (1..365)",
    "month": "Mensual (1..12)",
    "raw": "Sense Agregació (Crua)",
}
_SEASON_LABELS = {
    "All": "Totes",
    "Winter": "Hivern (Des, Gen, Feb)",
    "Spring": "Primavera (Mar, Abr, Mai)",
    "Summer": "Estiu (Jun, Jul, Ago)",
    "Autumn": "Tardor (Set, Oct, Nov)",
}
_VIEW_MODE_LABELS = {
    "aggregate": "Mitjana / Agregat",
    "real": "Serie real",
}
_REAL_PERIOD_LABELS = {
    "day": "Un dia",
    "week": "Una setmana",
    "month": "Un mes",
}


class DashboardTabSpec(NamedTuple):
    """Declaració d'una pestanya del dashboard i les figures que conté."""

    label: str
    visibility_group: str
    copy_text: str
    figure_items: tuple[tuple[str, str], ...]


CONTROL_FIGURE_KEYS = ("setpoints", "setpoints_vs_indoor", "radiant", "actions")
BATTERY_FIGURE_KEYS = (
    "battery_power",
    "battery_soc",
    "battery_vs_grid",
    "battery_charge_price",
)
REAL_AXIS_FIGURE_KEYS = (
    "indoor",
    "humidity",
    "setpoints",
    "setpoints_vs_indoor",
    "radiant",
    "actions",
    "hvac",
    "hvac_breakdown",
    "energy_price",
    "battery_power",
    "battery_soc",
    "battery_vs_grid",
    "battery_charge_price",
    "violation",
)
STYLE_FIGURE_KEYS = (
    "indoor",
    "humidity",
    "setpoints",
    "setpoints_vs_indoor",
    "radiant",
    "actions",
    "hvac",
    "energy_price",
    "battery_power",
    "battery_soc",
    "battery_vs_grid",
    "battery_charge_price",
    "violation",
    "zone_heatmap",
    "heatmap",
)
OVERLAY_FIGURE_KEYS = (
    "indoor",
    "humidity",
    "setpoints",
    "setpoints_vs_indoor",
    "radiant",
    "hvac",
    "hvac_breakdown",
    "energy_price",
    "battery_power",
    "battery_soc",
    "battery_vs_grid",
    "battery_charge_price",
    "violation",
)
AGGREGATE_OVERLAY_KEYS = ("comfort", "episode")
DASHBOARD_TAB_SPECS = (
    DashboardTabSpec(
        "Temperatura",
        "always",
        "Temperatura interior i exterior, percentatge d'hores en confort, infraccions tèrmiques "
        "i distribució de temperatura per zona i hora.",
        (
            ("comfort", "aggregate"),
            ("indoor", "always"),
            ("humidity", "always"),
            ("violation", "always"),
            ("zone_heatmap", "aggregate"),
        ),
    ),
    DashboardTabSpec(
        "Consum i Potència",
        "always",
        "Consum HVAC agregat, desglosament per meter, preu de l'energia i mapa de calor global "
        "del consum al llarg de l'any.",
        (
            ("hvac", "always"),
            ("hvac_breakdown", "always"),
            ("energy_price", "always"),
            ("heatmap", "aggregate"),
        ),
    ),
    DashboardTabSpec(
        "Control HVAC",
        "control",
        "Setpoints tèrmics, actuació del control radiant i accions que l'agent RL ha aplicat "
        "al sistema HVAC.",
        tuple((key, "always") for key in CONTROL_FIGURE_KEYS),
    ),
    DashboardTabSpec(
        "Bateria",
        "battery",
        "Cicles de càrrega i descàrrega, estat de càrrega (SOC), interacció amb la xarxa "
        "elèctrica i correlació amb el preu de l'energia.",
        tuple((key, "always") for key in BATTERY_FIGURE_KEYS),
    ),
    DashboardTabSpec(
        "Episodi",
        "always",
        "Mètriques de rendiment acumulades per episodi d'entrenament o avaluació.",
        (("episode", "aggregate"),),
    ),
)


def _file_signature(path: str) -> tuple[str, int, int]:
    """Retorna una signatura de memòria cau estable per a la ruta d'un fitxer."""

    file_path = Path(path)
    try:
        stat = file_path.stat()
    except OSError:
        return str(file_path), -1, -1
    return str(file_path), int(stat.st_mtime_ns), int(stat.st_size)


def _run_data_signature(progress_path: str, observations_path: str) -> tuple[object, ...]:
    """Retorna les entrades de memòria cau que canvien cada vegada que canvien les dades d'execució seleccionades."""

    run_dir = Path(progress_path).parent
    yaml_signatures = []
    try:
        yaml_signatures = [
            _file_signature(str(path))
            for path in sorted(run_dir.glob("*.yaml"))
        ]
    except OSError:
        yaml_signatures = []
    return (
        _file_signature(progress_path),
        _file_signature(observations_path),
        tuple(yaml_signatures),
    )


def _load_comparison_data_cached(run_path: str) -> DashboardData | None:
    """Carrega una comparació i la desa a la memòria cau de sessió."""

    cache_key = f"_comp_data_{hash(run_path)}"
    if cache_key in st.session_state:
        return st.session_state[cache_key]

    data = load_comparison_data(run_path)
    if data is None:
        return None

    st.session_state[cache_key] = data
    return data

@st.cache_data(show_spinner="Generant informe...", max_entries=10)
def _cached_report_bytes(
    progress_path: str,
    observations_path: str,
    agg_mode: str,
    season: str,
    comfort_scope: str,
    zone_str: str,
    data_signature: tuple[object, ...],
) -> tuple[bytes, str]:
    """Retorna bytes d'informes a la memòria cau per als filtres del panell seleccionats."""

    return generate_report_bytes(
        progress_path=progress_path,
        observations_path=observations_path,
        agg_mode=agg_mode,
        season=season,
        comfort_scope=comfort_scope,
        zone_str=zone_str,
    )


def _plot_if_available(fig: go.Figure | None) -> None:
    """Mostra un gràfic Plotly només quan conté traces visibles."""

    if has_figure_data(fig):
        # Grafic dashboard
        st.plotly_chart(fig, width="stretch")


def _dashboard_tab_is_visible(spec: DashboardTabSpec, figures: dict[str, go.Figure]) -> bool:
    """Decideix si una pestanya té dades suficients per aparèixer al dashboard."""
    if spec.visibility_group == "control":
        return any(has_figure_data(figures[key]) for key in CONTROL_FIGURE_KEYS)
    if spec.visibility_group == "battery":
        return any(has_figure_data(figures[key]) for key in BATTERY_FIGURE_KEYS)
    return True


def _render_dashboard_tab_content(
    spec: DashboardTabSpec,
    figures: dict[str, go.Figure],
    *,
    view_mode: str,
) -> None:
    """Pinta el contingut d'una pestanya de resultats segons la seva especificació."""
    # Text pestanya dashboard
    st.markdown(
        f"<div class='results-tab-copy'>{escape(spec.copy_text)}</div>",
        unsafe_allow_html=True,
    )
    if spec.label == "Episodi":
        if view_mode == "aggregate" and has_figure_data(figures["episode"]):
            # Grafic episodi
            st.plotly_chart(figures["episode"], width="stretch")
            return
        # Avís episodi agregat
        st.info("Les mètriques d'episodi només estan disponibles en mode agregat.")
        return

    for key, mode_rule in spec.figure_items:
        if mode_rule == "aggregate" and view_mode != "aggregate":
            continue
        _plot_if_available(figures[key])


def _render_report_export(
    report_action_slot,
    data: DashboardData,
    *,
    plot_mode: str,
    plot_season: str,
    comfort_scope: str,
    zone_val: str | None,
) -> None:
    """Mostra els controls de generació i baixada d'informes."""

    with report_action_slot.container():
        zone_str = zone_val if zone_val else "ALL"
        data_signature = _run_data_signature(data.progress_path, data.observations_path)
        report_key = (
            data.progress_path,
            data.observations_path,
            plot_mode,
            plot_season,
            comfort_scope,
            zone_str,
            data_signature,
        )
        if st.session_state.get("_report_cache_key") != report_key:
            st.session_state.pop("_report_cache", None)
            st.session_state["_report_cache_key"] = report_key

        if "_report_cache" not in st.session_state:
            # Botó generar informe
            if st.button(
                "Generar informe",
                icon=":material/summarize:",
                key="btn_gen_report",
                width="stretch",
            ):
                with st.spinner("Generant informe..."):
                    st.session_state["_report_cache"] = _cached_report_bytes(
                        data.progress_path,
                        data.observations_path,
                        plot_mode,
                        plot_season,
                        comfort_scope,
                        zone_str,
                        data_signature,
                    )
                st.rerun()
            return

        report_bytes, report_ext = st.session_state["_report_cache"]
        mime = "application/pdf" if report_ext == "pdf" else "text/html"
        # Botó download informe
        st.download_button(
            f"Descarregar informe (.{report_ext})",
            data=report_bytes,
            file_name=f"Sinergym_Report_{data.env_name}_{plot_season}.{report_ext}",
            mime=mime,
            key="btn_download_report",
            width="stretch",
        )


def _render_dashboard_tabs(figures: dict[str, go.Figure], *, view_mode: str) -> None:
    """Mostra els gràfics de resultats amb la mateixa estructura de pestanyes del panell."""

    active_specs = [
        spec for spec in DASHBOARD_TAB_SPECS if _dashboard_tab_is_visible(spec, figures)
    ]
    # Pestanyes dashboard
    active_tabs = st.tabs([spec.label for spec in active_specs])
    for tab, spec in zip(active_tabs, active_specs):
        with tab:
            _render_dashboard_tab_content(
                spec,
                figures,
                view_mode=view_mode,
            )


def _build_dashboard_figures(
    data: DashboardData,
    zobs: pd.DataFrame,
    action_data: pd.DataFrame,
    *,
    plot_mode: str,
    plot_season: str,
    comfort_scope: str,
    view_mode: str,
    real_period_kind: str,
) -> dict[str, go.Figure]:
    """Crea totes les figures Plotly utilitzades pel panell de resultats en línia."""

    comfort_mode = "raw" if plot_mode in ("day", "raw") else plot_mode
    hvac_raw_as_line = view_mode == "real" and real_period_kind == "month"

    # Generem totes les figures en un sol lloc perquè la comparació pugui aplicar el mateix
    # tractament a la run principal i a la run de referència.
    figures: dict[str, go.Figure] = {
        "comfort": (
            make_comfort_compliance(
                zobs,
                comfort_mode,
                plot_season,
                comfort_scope=comfort_scope,
                comfort_config=data.yaml_cfg,
            )
            if view_mode == "aggregate"
            else go.Figure()
        ),
        "indoor": make_indoor_temperature_plot(zobs, plot_mode, plot_season),
        "humidity": make_indoor_humidity_plot(zobs, plot_mode, plot_season),
        "setpoints": make_setpoints_plot(zobs, plot_mode, plot_season),
        "setpoints_vs_indoor": make_setpoints_vs_indoor_plot(zobs, plot_mode, plot_season),
        "radiant": make_radiant_control_plot(zobs, plot_mode, plot_season),
        "actions": make_agent_actions_plot(action_data, plot_mode, plot_season),
        "hvac": _fix_colors_for_visibility(
            make_hvac_consumption_plot(
                zobs,
                plot_mode,
                plot_season,
                raw_as_line=hvac_raw_as_line,
            ),
            kind="hvac",
            mode=plot_mode,
        ),
        "hvac_breakdown": make_hvac_meter_breakdown_plot(zobs, plot_mode, plot_season),
        "energy_price": make_energy_price_plot(zobs, plot_mode, plot_season),
        "battery_power": make_battery_power_plot(zobs, plot_mode, plot_season),
        "battery_soc": make_battery_soc_plot(zobs, plot_mode, plot_season),
        "battery_vs_grid": make_battery_vs_grid_plot(zobs, plot_mode, plot_season),
        "battery_charge_price": make_battery_charge_with_price_plot(zobs, plot_mode, plot_season),
        "zone_heatmap": go.Figure(),
        "heatmap": go.Figure(),
        "episode": make_episode_metrics(data.progress)
        if view_mode == "aggregate"
        else go.Figure(),
    }

    if not has_action_figure_data(figures["actions"]):
        figures["actions"] = go.Figure()

    if view_mode == "aggregate":
        figures["violation"] = _fix_colors_for_visibility(
            make_violation_bars(
                zobs,
                plot_mode,
                plot_season,
                comfort_scope=comfort_scope,
                comfort_config=data.yaml_cfg,
            ),
            kind="violation",
            mode=plot_mode,
        )
        figures["zone_heatmap"] = make_zone_temperature_heatmaps(
            zobs, season=plot_season, agg="mean", max_cols=3
        )
        pivot = data.metrics_dict.get("pivot_consumption")
        if pivot is not None and not pivot.empty:
            figures["heatmap"] = make_heatmap(pivot, "Total HVAC Consumption (kWh)")
    else:
        figures["violation"] = make_violation_single_line(
            zobs,
            plot_mode,
            plot_season,
            comfort_config=data.yaml_cfg,
        )
        for key in REAL_AXIS_FIGURE_KEYS:
            axis_data = action_data if key == "actions" else zobs
            figures[key] = apply_real_period_axis(figures[key], axis_data, real_period_kind)

    return figures


def _overlay_dashboard_figures(
    figures: dict[str, go.Figure],
    comparison_figures: dict[str, go.Figure],
    *,
    data: DashboardData,
    view_mode: str,
) -> dict[str, go.Figure]:
    """Superposeu traces de comparació a les xifres actives del panell."""

    overlay_keys = list(OVERLAY_FIGURE_KEYS)
    if view_mode == "aggregate":
        overlay_keys.extend(AGGREGATE_OVERLAY_KEYS)
    if is_radiant_run(data):
        overlay_keys.append("actions")

    for key in overlay_keys:
        figures[key] = overlay_comparison_traces(figures[key], comparison_figures[key])
    return figures


def _style_dashboard_figures(
    figures: dict[str, go.Figure],
    *,
    view_mode: str,
) -> dict[str, go.Figure]:
    """Aplica la semàntica comuna de Plotly a les figures del panell abans de representar-les."""

    styled = dict(figures)
    for key in STYLE_FIGURE_KEYS:
        styled[key] = style_figure_semantics(styled[key])
    if view_mode == "aggregate":
        styled["comfort"] = style_figure_semantics(styled["comfort"])
        styled["episode"] = style_figure_semantics(styled["episode"])
    return styled


_KPI_HIGHER_IS_BETTER = frozenset({"% Hores en Confort"})


def _kpi_delta_color(key: str) -> str:
    """Retorna el mode de color delta utilitzat per les targetes de comparació KPI."""
    if key in _KPI_HIGHER_IS_BETTER:
        return "normal"
    key_lower = key.lower()
    if "hvac" in key_lower or "violation" in key_lower:
        return "inverse"
    return "off"


def _render_kpis(kpis: dict, comp_kpis: dict | None) -> None:
    """Pinta les targetes KPI amb deltes de comparació opcionals."""
    kpi_items = [(k, v) for k, v in kpis.items()]
    if not kpi_items:
        return

    cards: list[str] = []
    for k, v in kpi_items:
        if v is None:
            val_str = "N/D"
        elif isinstance(v, float) and k.startswith("%"):
            val_str = f"{v:.1f}%"
        elif isinstance(v, (int, float)):
            val_str = f"{float(v):.2f}"
        else:
            val_str = str(v)

        delta_html = ""
        if comp_kpis is not None and k in comp_kpis:
            # El signe bo del delta depèn del KPI: menys cost o violació és millor,
            # però més confort pot ser positiu. _kpi_delta_color encapsula aquesta regla.
            cv = comp_kpis[k]
            if (
                isinstance(v, (int, float)) and isinstance(cv, (int, float))
                and v is not None and cv is not None
            ):
                delta = round(float(v) - float(cv), 2)
                color_hint = _kpi_delta_color(k)
                if color_hint == "inverse":
                    is_good = delta < 0
                elif color_hint == "normal":
                    is_good = delta > 0
                else:
                    is_good = None
                if delta == 0:
                    delta_class = "results-kpi-delta-neutral"
                    arrow = ""
                elif is_good is None:
                    delta_class = "results-kpi-delta-neutral"
                    arrow = "▲" if delta > 0 else "▼"
                elif is_good:
                    delta_class = "results-kpi-delta-good"
                    arrow = "▲" if delta > 0 else "▼"
                else:
                    delta_class = "results-kpi-delta-bad"
                    arrow = "▲" if delta > 0 else "▼"
                sign = "+" if delta > 0 else ""
                delta_html = (
                    f'<div class="results-kpi-delta {delta_class}">'
                    f'{arrow} {sign}{delta:.2f}</div>'
                )

        k_esc = escape(k)
        v_esc = escape(val_str)
        cards.append(
            f'<div class="results-kpi-card">'
            f'<div class="results-kpi-label" title="{k_esc}">{k_esc}</div>'
            f'<div class="results-kpi-value">{v_esc}</div>'
            f'{delta_html}'
            f'</div>'
        )

    # Targetes KPI
    st.markdown(
        '<section class="section-card results-kpi-shell">' + "".join(cards) + "</section>",
        unsafe_allow_html=True,
    )

def render_inline_dashboard(
    data: DashboardData,
    all_run_dirs: tuple,
    current_run_path: str,
) -> None:
    """Mostra el panell en línia complet per a una execució carregada a la UI de Streamlit."""

    from backend.grafics.figures_zones import filter_obs_by_zone
    from backend.grafics.kpis import compute_kpis

    # Separador dashboard
    st.markdown("---")

    # --- Capçalera del panell ---
    try:
        # Capçalera dashboard
        head_cols = st.columns([8, 2], gap="small", vertical_alignment="center")
    except TypeError:
        # Capçalera dashboard
        head_cols = st.columns([8, 2], gap="small")
    with head_cols[0]:
        # Titol dashboard
        st.markdown('<h2 class="section-title">Dashboard</h2>', unsafe_allow_html=True)
    with head_cols[1]:
        report_action_slot = st.empty()

    # Etiquetes resum dashboard
    st.markdown(
        f"""
        <div class="results-dashboard-header">
          <span class="results-dashboard-tag">
            <strong>Entorn</strong>&nbsp; {escape(data.env_name)}
          </span>
          <span class="results-dashboard-tag">
            <strong>Timesteps</strong>&nbsp; {data.timesteps_num:,}
          </span>
          <span class="results-dashboard-tag">
            <strong>Model</strong>&nbsp; {escape(data.model_name)}
          </span>
        </div>
        """,
        unsafe_allow_html=True,
    )

    # Separador selector comparacio
    st.markdown("<div class='studio-spacer-075'></div>", unsafe_allow_html=True)

    # --- Selector de comparació ---
    comp_options = [("", "Cap (sense comparació)")] + [
        (r.path, str(r)) for r in all_run_dirs if r.path != current_run_path
    ]
    # Selector execucio comparacio
    comp_path = st.selectbox(
        "Execució de referència per a comparació",
        options=[p for p, _ in comp_options],
        format_func=lambda x: dict(comp_options).get(x, x),
        key="dash_comp_run",
    )

    comp_data: DashboardData | None = None
    if comp_path:
        with st.spinner("Carregant run de comparació..."):
            comp_data = _load_comparison_data_cached(comp_path)
        if comp_data is None:
            # Avís comparacio no carregada
            st.warning("No s'han pogut carregar les dades de la run de referència.")

    # Separador KPIs dashboard
    st.markdown("<div class='studio-spacer-075'></div>", unsafe_allow_html=True)

    # Llegim els filtres des de session_state abans de pintar els widgets. Així els KPI
    # surten a dalt de tot, però sempre calculats amb l'última selecció real de l'usuari.
    view_mode = st.session_state.get("dash_view_mode", "aggregate")
    agg_mode = "hour"
    season = "All"
    real_period_kind = st.session_state.get("dash_real_period_kind", "day")
    if real_period_kind not in _REAL_PERIOD_LABELS:
        real_period_kind = "month"
        st.session_state["dash_real_period_kind"] = real_period_kind
    real_period_value = ""
    real_period_label = None

    if view_mode == "aggregate":
        agg_mode = st.session_state.get("dash_agg_mode", "hour")
        season = st.session_state.get("dash_season", "All")
    else:
        agg_mode = "raw"
        real_period_value = st.session_state.get(
            f"dash_real_period_value_{real_period_kind}", ""
        )
        if not real_period_value:
            _period_opts = build_real_period_options(data.obs, real_period_kind)
            real_period_value = _period_opts[0][0] if _period_opts else ""

    _default_comfort_scope = (
        data.comfort_scope_options[0]["value"] if data.comfort_scope_options else "all"
    )
    comfort_scope = st.session_state.get("dash_comfort_scope", _default_comfort_scope)
    zone_val = st.session_state.get("dash_zone", "ALL")

    sel_zone = None if (not data.has_zones or zone_val == "ALL") else zone_val
    zobs = filter_obs_by_zone(data.obs, sel_zone, data.yaml_cfg)
    if view_mode == "real":
        zobs, real_period_label = slice_obs_for_real_period(
            zobs,
            real_period_kind,
            real_period_value,
        )
        if zobs.empty:
            # Avís periode real buit
            st.warning("No hi ha dades reals disponibles per al període seleccionat.")
            return
        # Text periode real
        st.caption(f"Comparant dades reals del període: {real_period_label}")
    action_data = select_actions_for_obs(data.actions, zobs)
    action_data = select_radiant_action_data(action_data, data)

    # KPI principals del període i la zona actius.
    kpis = compute_kpis(
        obs=zobs,
        cost_hourly=data.metrics_dict.get("cost_hourly"),
        selected_zone=sel_zone,
        comfort_scope=comfort_scope,
        comfort_config=data.yaml_cfg,
    )
    kpis.pop("Avg Energy Cost (EUR/h)", None)

    add_comfort_percentage_kpi(kpis, zobs, data.yaml_cfg)

    comp_kpis: dict | None = None
    comp_sel_zone = None
    comp_zobs = None
    comp_action_data = pd.DataFrame()
    if comp_data is not None:
        # La comparació intenta aplicar exactament el mateix filtre temporal i de zona.
        # Si la run de referència no té aquell any concret, fem fallback per mes/dia equivalent.
        comp_sel_zone = sel_zone if comp_data.has_zones else None
        comp_zobs = filter_obs_by_zone(comp_data.obs, comp_sel_zone, comp_data.yaml_cfg)
        if view_mode == "real":
            comp_zobs, comp_period_label = slice_obs_for_real_period(
                comp_zobs,
                real_period_kind,
                real_period_value,
                allow_fallback=True,
            )
            if comp_zobs.empty:
                # Avís referencia buida
                st.warning("La run de referència no té dades per al període real seleccionat.")
                comp_data = None
            elif comp_period_label and comp_period_label != real_period_label:
                # Text periode referencia
                st.caption(f"La referència s'ha alineat amb el període: {comp_period_label}")
        if comp_data is None:
            comp_zobs = None
        else:
            comp_action_data = select_actions_for_obs(comp_data.actions, comp_zobs)
            comp_action_data = select_radiant_action_data(comp_action_data, comp_data)
            comp_kpis = compute_kpis(
                obs=comp_zobs,
                cost_hourly=comp_data.metrics_dict.get("cost_hourly"),
                selected_zone=comp_sel_zone,
                comfort_scope=comfort_scope,
                comfort_config=comp_data.yaml_cfg,
            )
            comp_kpis.pop("Avg Energy Cost (EUR/h)", None)
            add_comfort_percentage_kpi(comp_kpis, comp_zobs, comp_data.yaml_cfg)
            # Text delta comparacio
            st.caption(
                f"Delta = execucio actual - referencia ({dict(comp_options).get(comp_path, comp_path)})"
            )

    # Pintem els KPI abans dels filtres: és una decisió de lectura, no de càlcul.
    _render_kpis(kpis, comp_kpis)

    # Separador filtres dashboard
    st.markdown("<div class='studio-spacer-075'></div>", unsafe_allow_html=True)

    # Els filtres es pinten sota els KPI perquè el panell funcioni com un resum primer
    # i com una eina d'exploració just després.
    # Filtres dashboard
    filter_cols = st.columns(5, gap="small")
    with filter_cols[0]:
        # Selector vista temporal
        st.selectbox(
            "Vista Temporal",
            options=list(_VIEW_MODE_LABELS.keys()),
            format_func=lambda x: _VIEW_MODE_LABELS[x],
            key="dash_view_mode",
        )
    with filter_cols[1]:
        if view_mode == "aggregate":
            # Selector agregacio temporal
            st.selectbox(
                "Agregació Temporal",
                options=["hour", "day", "month"],
                format_func=lambda x: _AGG_LABELS[x],
                key="dash_agg_mode",
            )
        else:
            # Selector periode real
            st.selectbox(
                "Període Real",
                options=list(_REAL_PERIOD_LABELS.keys()),
                format_func=lambda x: _REAL_PERIOD_LABELS[x],
                key="dash_real_period_kind",
            )
    with filter_cols[2]:
        if view_mode == "aggregate":
            # Selector estacio
            st.selectbox(
                "Filtrar per Estació",
                options=list(_SEASON_LABELS.keys()),
                format_func=lambda x: _SEASON_LABELS[x],
                key="dash_season",
            )
        else:
            real_period_options = build_real_period_options(data.obs, real_period_kind)
            if real_period_options:
                option_map = dict(real_period_options)
                # Selector valor periode
                st.selectbox(
                    "Selecciona el Període",
                    options=[value for value, _ in real_period_options],
                    format_func=lambda x: option_map.get(x, x),
                    key=f"dash_real_period_value_{real_period_kind}",
                )
            else:
                # Selector periode buit
                st.selectbox(
                    "Selecciona el Període",
                    options=["Sense dades disponibles"],
                    disabled=True,
                    key=f"dash_real_period_empty_{real_period_kind}",
                )
    with filter_cols[3]:
        scope_labels = {opt["value"]: opt["label"] for opt in data.comfort_scope_options}
        # Selector ambit confort
        st.selectbox(
            "Ambit de Confort",
            options=[opt["value"] for opt in data.comfort_scope_options],
            format_func=lambda x: scope_labels[x],
            key="dash_comfort_scope",
        )
    with filter_cols[4]:
        if data.has_zones:
            zone_labels = {opt["value"]: opt["label"] for opt in data.all_zone_options}
            # Selector zona
            st.selectbox(
                "Zona",
                options=[opt["value"] for opt in data.all_zone_options],
                format_func=lambda x: zone_labels[x],
                key="dash_zone",
            )
        else:
            # Selector zona global
            st.selectbox(
                "Zona",
                options=["Global (Totes)"],
                key="dash_zone",
                disabled=True,
            )

    # Separador grafics dashboard
    st.markdown("<div class='studio-spacer-100'></div>", unsafe_allow_html=True)

    # --- Construeix totes les figures del panell ---
    plot_mode = agg_mode
    plot_season = season
    figures = _build_dashboard_figures(
        data,
        zobs,
        action_data,
        plot_mode=plot_mode,
        plot_season=plot_season,
        comfort_scope=comfort_scope,
        view_mode=view_mode,
        real_period_kind=real_period_kind,
    )

    if comp_data is not None and comp_zobs is not None:
        comp_action_data = select_radiant_action_data(comp_action_data, comp_data)
        comparison_figures = _build_dashboard_figures(
            comp_data,
            comp_zobs,
            comp_action_data,
            plot_mode=plot_mode,
            plot_season=plot_season,
            comfort_scope=comfort_scope,
            view_mode=view_mode,
            real_period_kind=real_period_kind,
        )
        figures = _overlay_dashboard_figures(
            figures,
            comparison_figures,
            data=data,
            view_mode=view_mode,
        )

    figures = _style_dashboard_figures(figures, view_mode=view_mode)

    _render_report_export(
        report_action_slot,
        data,
        plot_mode=plot_mode,
        plot_season=plot_season,
        comfort_scope=comfort_scope,
        zone_val=zone_val,
    )

    _render_dashboard_tabs(figures, view_mode=view_mode)
\end{lstlisting}

\Needspace{8\baselineskip}

\subsection{\texorpdfstring{\texttt{page\_components/training\_agent.py}}{page\_components/training\_agent.py}}

\begin{lstlisting}[style=studioPython,caption={page\_components/training\_agent.py},label={lst:page-components-training-agent-py}]
"""Controls d'hiperparàmetres de l'agent per a la pàgina d'entrenament."""

from __future__ import annotations

import streamlit as st
from backend.entrenar_agent_session import training_form_key


def render_agent_params_section(algo_name: str) -> dict:
    """Mostra els widgets d'hiperparàmetres de l'agent i retorna els seus valors.

    Cobreix la taxa d'aprenentatge, n_steps, els passos de temps totals i tots els SB3 avançats
    paràmetres (batch_size, gamma, n_epochs, gae_lambda, clip_range,
    ent_coef, vf_coef, max_grad_norm).

    Retorna:
        Un dict amb tots els valors dels paràmetres d'agent recollits dels widgets.
    """
    # Els tres primers camps són els que més es toquen en proves ràpides; els avançats
    # queden agrupats perquè la pantalla no sembli una paret de paràmetres.
    # Inputs principals agent
    agent_col1, agent_col2, agent_col3 = st.columns(3)
    with agent_col1:
        # Input learning rate
        learning_rate = st.number_input(
            "learning_rate",
            1e-5,
            1e-2,
            3e-4,
            format="%.5f",
            key=training_form_key("learning_rate"),
        )
    with agent_col2:
        # Input passos rollout
        n_steps = st.number_input(
            "n_steps (PPO/A2C)",
            256,
            4096,
            2048,
            step=256,
            key=training_form_key("n_steps"),
        )
    with agent_col3:
        # Input timesteps totals
        timesteps_per_year = st.number_input(
            "total_timesteps",
            35040,
            10_000_000,
            35040,
            step=35040,
            key=training_form_key("timesteps_per_year"),
        )

    with st.container():
        # Seccio parametres avançats
        st.markdown("**Parametres SB3 avancats**")
        # Aquests valors tenen defaults raonables per PPO/A2C. Si un algorisme no els usa,
        # assemble_training_payload els filtra abans de crear el model.
        # Inputs SB3 avançats
        sb3_col1, sb3_col2, sb3_col3, sb3_col4 = st.columns(4)
        with sb3_col1:
            # Input batch size
            batch_size = st.number_input(
                "batch_size",
                8,
                8192,
                256,
                step=8,
                key=training_form_key("batch_size"),
            )
            # Input gamma
            gamma = st.number_input(
                "gamma",
                0.0,
                0.9999,
                0.995,
                step=0.0005,
                format="%.4f",
                key=training_form_key("gamma"),
            )
        with sb3_col2:
            # Input epochs PPO
            n_epochs = st.number_input(
                "n_epochs",
                1,
                50,
                10,
                step=1,
                key=training_form_key("n_epochs"),
            )
            # Input GAE lambda
            gae_lambda = st.number_input(
                "gae_lambda",
                0.0,
                1.0,
                0.95,
                step=0.01,
                format="%.2f",
                key=training_form_key("gae_lambda"),
            )
        with sb3_col3:
            # Input clip range
            clip_range = st.number_input(
                "clip_range",
                0.01,
                1.0,
                0.2,
                step=0.01,
                format="%.2f",
                key=training_form_key("clip_range"),
            )
            # Input entropy coef
            ent_coef = st.number_input(
                "ent_coef",
                0.0,
                1.0,
                0.005,
                step=0.001,
                format="%.4f",
                key=training_form_key("ent_coef"),
            )
        with sb3_col4:
            # Input value loss coef
            vf_coef = st.number_input(
                "vf_coef",
                0.0,
                2.0,
                0.5,
                step=0.05,
                format="%.2f",
                key=training_form_key("vf_coef"),
            )
            # Input gradient clipping
            max_grad_norm = st.number_input(
                "max_grad_norm",
                0.0,
                10.0,
                0.5,
                step=0.1,
                format="%.1f",
                key=training_form_key("max_grad_norm"),
            )

    return {
        "learning_rate":    learning_rate,
        "n_steps":          n_steps,
        "timesteps_per_year": timesteps_per_year,
        "batch_size":       batch_size,
        "gamma":            gamma,
        "n_epochs":         n_epochs,
        "gae_lambda":       gae_lambda,
        "clip_range":       clip_range,
        "ent_coef":         ent_coef,
        "vf_coef":          vf_coef,
        "max_grad_norm":    max_grad_norm,
    }


# ── Pàgina principal ──────────────────────────────── ────────────────────────────────
\end{lstlisting}

\Needspace{8\baselineskip}

\subsection{\texorpdfstring{\texttt{page\_components/training\_rewards.py}}{page\_components/training\_rewards.py}}

\begin{lstlisting}[style=studioPython,caption={page\_components/training\_rewards.py},label={lst:page-components-training-rewards-py}]
"""Controls de configuració de recompenses per a la pàgina de entrenament."""

from __future__ import annotations

import streamlit as st

from backend.entrenar_agent_constants import COST_REWARDS, HOURLY_REWARDS, MULTIZONE_REWARDS
from backend.entrenar_agent_session import (
    BATTERY_REWARDS,
    sanitize_multiselect_state,
    selectbox_state_kwargs,
    training_form_key,
)


def render_reward_kwargs_section(
    reward_name: str,
    observation_variables: list[str],
    candidate_grid_energy_vars: list[str],
    candidate_battery_charge_vars: list[str],
    candidate_battery_discharge_vars: list[str],
    candidate_battery_loss_vars: list[str],
) -> dict:
    """Mostra tots els widgets de recompensa i retorna els valors recollits.

    Cobreix pesos d'energia/confort, rangs de confort estacionals, específics de la bateria
    paràmetres, paràmetres de recompensa de cost i configuració d'hores ocupades.

    Retorna:
        Un dict amb tots els kwargs de recompensa recollits dels widgets.
    """
    # Defaults de partida. Els ajustem després segons la reward escollida perquè una
    # bateria, un cost horari i una reward simple no tenen la mateixa escala.
    temperature_weight              = 0.4
    lambda_energy_cost              = 1.0
    comfort_threshold               = 0.5
    range_comfort_hours             = (9, 19)
    occupied_hours                  = (8, 18)
    occupied_discomfort_multiplier  = 20.0
    off_hours_energy_multiplier     = 5.0
    occupied_comfort_multiplier     = 2.0
    unoccupied_comfort_multiplier   = 0.3
    energy_cost_weight              = 0.20
    grid_energy_weight              = 0.2
    battery_cycle_weight            = 0.04
    battery_loss_weight             = 0.005
    simultaneous_battery_weight     = 0.005
    lambda_grid                     = 0.0001
    lambda_battery                  = 0.0001
    grid_energy_variables           = [v for v in candidate_grid_energy_vars if v in observation_variables]
    battery_charge_variables        = [v for v in candidate_battery_charge_vars if v in observation_variables]
    battery_discharge_variables     = [v for v in candidate_battery_discharge_vars if v in observation_variables]
    battery_loss_variables          = [v for v in candidate_battery_loss_vars if v in observation_variables]

    # Aquestes banderes fan que la funció sigui llarga, però mantenen la UI en una sola
    # passada: només es pinten els camps que realment necessita la reward seleccionada.
    occupied_reward_selected      = reward_name == "OccupiedHoursDiscomfortReward"
    battery_reward_selected       = reward_name in BATTERY_REWARDS
    new_battery_cost_selected     = reward_name == "MultizoneEnergyCostBatteryHVACReward"
    energy_weight_default         = 0.2 if (occupied_reward_selected or battery_reward_selected) else 0.5
    lambda_energy_default         = 0.001 if occupied_reward_selected else 0.0001
    temperature_weight_default    = 0.40 if new_battery_cost_selected else (0.55 if battery_reward_selected else 0.4)
    sanitize_multiselect_state("grid_energy_variables", observation_variables, grid_energy_variables)
    sanitize_multiselect_state("battery_charge_variables", observation_variables, battery_charge_variables)
    sanitize_multiselect_state("battery_discharge_variables", observation_variables, battery_discharge_variables)
    sanitize_multiselect_state("battery_loss_variables", observation_variables, battery_loss_variables)

    # Inputs pesos reward base
    _p1, _p2, _p3 = st.columns(3)
    with _p1:
        # Input pes energia
        energy_weight = st.number_input(
            "energy_weight",
            0.0,
            1.0,
            energy_weight_default,
            format="%.2f",
            key=training_form_key("energy_weight"),
        )
    with _p2:
        # Input lambda energia
        lambda_energy = st.number_input(
            "lambda_energy",
            0.0,
            10.0,
            lambda_energy_default,
            format="%.5f",
            key=training_form_key("lambda_energy"),
        )
    with _p3:
        # Input lambda temperatura
        lambda_temperature = st.number_input(
            "lambda_temperature",
            0.0,
            10.0,
            1.0,
            format="%.2f",
            key=training_form_key("lambda_temperature"),
        )

    if reward_name in COST_REWARDS:
        # Inputs reward cost
        _c1, _c2 = st.columns(2)
        with _c1:
            # Input pes temperatura
            temperature_weight = st.number_input(
                "temperature_weight",
                0.0,
                1.0,
                0.4,
                format="%.2f",
                key=training_form_key("temperature_weight"),
            )
        with _c2:
            # Input lambda cost energia
            lambda_energy_cost = st.number_input(
                "lambda_energy_cost",
                0.0,
                10.0,
                0.0001,
                format="%.5f",
                key=training_form_key("lambda_energy_cost"),
            )

    if reward_name in BATTERY_REWARDS:
        # En bateria és fàcil confondre potència de xarxa, càrrega i descàrrega. Per això
        # deixem els multiselect visibles i amb candidats ja filtrats per l'entorn actiu.
        # Inputs pesos bateria
        _b1, _b2, _b3 = st.columns(3)
        with _b1:
            # Input pes temperatura bateria
            temperature_weight = st.number_input(
                "temperature_weight",
                0.0,
                1.0,
                temperature_weight_default,
                format="%.2f",
                key=training_form_key("temperature_weight"),
            )
        with _b2:
            # Input pes xarxa
            grid_energy_weight = st.number_input(
                "grid_energy_weight",
                0.0,
                1.0,
                0.2,
                format="%.2f",
                key=training_form_key("grid_energy_weight"),
            )
        with _b3:
            # Input pes cicle bateria
            battery_cycle_weight = st.number_input(
                "battery_cycle_weight",
                0.0,
                1.0,
                0.04,
                format="%.3f",
                key=training_form_key("battery_cycle_weight"),
            )

        # Inputs penalitzacions bateria
        _b4, _b5, _b6 = st.columns(3)
        with _b4:
            # Input perdues bateria
            battery_loss_weight = st.number_input(
                "battery_loss_weight",
                0.0,
                1.0,
                0.005,
                format="%.3f",
                key=training_form_key("battery_loss_weight"),
            )
        with _b5:
            # Input simultaneitat bateria
            simultaneous_battery_weight = st.number_input(
                "simultaneous_battery_weight",
                0.0,
                1.0,
                0.005,
                format="%.3f",
                key=training_form_key("simultaneous_battery_weight"),
            )
        with _b6:
            # Input lambda xarxa
            lambda_grid = st.number_input(
                "lambda_grid",
                0.0,
                10.0,
                0.0001,
                format="%.5f",
                key=training_form_key("lambda_grid"),
            )

        # Input lambda bateria
        lambda_battery = st.number_input(
            "lambda_battery",
            0.0,
            10.0,
            0.0001,
            format="%.5f",
            key=training_form_key("lambda_battery"),
        )
        # Selector variables xarxa
        grid_energy_variables = st.multiselect(
            "Variables xarxa",
            observation_variables,
            default=grid_energy_variables,
            key=training_form_key("grid_energy_variables"),
        )
        # Selector variables carrega bateria
        battery_charge_variables = st.multiselect(
            "Variables carrega bateria",
            observation_variables,
            default=battery_charge_variables,
            key=training_form_key("battery_charge_variables"),
        )
        # Selector variables battery discharge
        battery_discharge_variables = st.multiselect(
            "Variables descarrega bateria",
            observation_variables,
            default=battery_discharge_variables,
            key=training_form_key("battery_discharge_variables"),
        )
        # Selector variables perdues bateria
        battery_loss_variables = st.multiselect(
            "Variables perdues bateria",
            observation_variables,
            default=battery_loss_variables,
            key=training_form_key("battery_loss_variables"),
        )

    if new_battery_cost_selected:
        # Aquesta reward barreja confort, cost econòmic i ocupació; separar el bloc evita
        # que aquests multiplicadors es vegin com a paràmetres genèrics de totes les rewards.
        # Seccio cost i ocupacio
        st.markdown("**Cost energetic i hores ocupades**")
        # Inputs cost i ocupacio
        _nc1, _nc2, _nc3 = st.columns(3)
        with _nc1:
            # Input pes cost energia
            energy_cost_weight = st.number_input(
                "energy_cost_weight", 0.0, 1.0, 0.20, format="%.2f",
                help="Pes del terme de cost economic (grid_kW * EUR/kWh).",
                key=training_form_key("energy_cost_weight"),
            )
        with _nc2:
            # Input escala cost energia
            lambda_energy_cost = st.number_input(
                "lambda_energy_cost", 0.0, 10.0, 1.0, format="%.4f",
                help="Escala del terme de cost economic.",
                key=training_form_key("lambda_energy_cost"),
            )
        with _nc3:
            # Slider hores ocupades
            occupied_hours = st.slider(
                "occupied_hours",
                0,
                23,
                (8, 18),
                key=training_form_key("occupied_hours"),
            )

        # Inputs multiplicadors confort
        _nm1, _nm2 = st.columns(2)
        with _nm1:
            # Input confort ocupat
            occupied_comfort_multiplier = st.number_input(
                "occupied_comfort_multiplier", 0.0, 20.0, 2.0, step=0.5, format="%.1f",
                help="Multiplicador del terme de confort durant hores ocupades.",
                key=training_form_key("occupied_comfort_multiplier"),
            )
        with _nm2:
            # Input confort no ocupat
            unoccupied_comfort_multiplier = st.number_input(
                "unoccupied_comfort_multiplier", 0.0, 5.0, 0.3, step=0.1, format="%.1f",
                help="Multiplicador del terme de confort fora d'hores ocupades.",
                key=training_form_key("unoccupied_comfort_multiplier"),
            )

    if reward_name in MULTIZONE_REWARDS:
        # Input llindar confort
        comfort_threshold = st.number_input(
            "comfort_threshold (C)",
            0.0,
            5.0,
            0.5,
            step=0.1,
            format="%.1f",
            key=training_form_key("comfort_threshold"),
        )

    if reward_name in HOURLY_REWARDS:
        # Slider hores confort
        range_comfort_hours = st.slider(
            "range_comfort_hours",
            0,
            23,
            (9, 19),
            key=training_form_key("range_comfort_hours"),
        )

    if reward_name == "OccupiedHoursDiscomfortReward":
        # Inputs reward ocupacio
        _o1, _o2, _o3 = st.columns(3)
        with _o1:
            # Slider hores ocupades
            occupied_hours = st.slider(
                "occupied_hours",
                0,
                23,
                (8, 18),
                key=training_form_key("occupied_hours"),
            )
        with _o2:
            # Input penalitzacio ocupada
            occupied_discomfort_multiplier = st.number_input(
                "occupied_discomfort_multiplier",
                min_value=1.0,
                max_value=100.0,
                value=20.0,
                step=1.0,
                key=training_form_key("occupied_discomfort_multiplier"),
            )
        with _o3:
            # Input energia fora hores
            off_hours_energy_multiplier = st.number_input(
                "off_hours_energy_multiplier",
                min_value=1.0,
                max_value=100.0,
                value=5.0,
                step=0.5,
                key=training_form_key("off_hours_energy_multiplier"),
            )

    st.divider()
    _MONTHS = ["Gen", "Feb", "Mar", "Abr", "Mai", "Jun", "Jul", "Ago", "Set", "Oct", "Nov", "Des"]
    # Les rewards de confort separen hivern i estiu; deixem la frontera editable perquè
    # el mateix edifici pot treballar amb climes molt diferents.
    # Slider rang confort hivern
    range_winter = st.slider(
        "range_comfort_winter",
        0.0,
        50.0,
        (20.0, 23.5),
        step=0.5,
        key=training_form_key("range_winter"),
    )
    # Slider rang confort estiu
    range_summer = st.slider(
        "range_comfort_summer",
        0.0,
        50.0,
        (23.0, 26.0),
        step=0.5,
        key=training_form_key("range_summer"),
    )

    st.divider()
    # Selectors mesos estiu
    _mc1, _mc2 = st.columns(2)
    with _mc1:
        month_options = list(range(1, 13))
        # Selector mes inici estiu
        summer_start_m = st.selectbox(
            "Inici estiu (mes)",
            month_options,
            format_func=lambda x: _MONTHS[x - 1],
            **selectbox_state_kwargs("summer_start_m", month_options, fallback=6),
        )
    with _mc2:
        # Selector mes final estiu
        summer_final_m = st.selectbox(
            "Final estiu (mes)",
            month_options,
            format_func=lambda x: _MONTHS[x - 1],
            **selectbox_state_kwargs("summer_final_m", month_options, fallback=9),
        )
    # Inputs dies estiu
    _dc1, _dc2 = st.columns(2)
    with _dc1:
        # Input dia inici estiu
        summer_start_d = st.number_input(
            "Dia inici",
            1,
            31,
            1,
            key=training_form_key("summer_start_d"),
        )
    with _dc2:
        # Input dia final estiu
        summer_final_d = st.number_input(
            "Dia final",
            1,
            31,
            30,
            key=training_form_key("summer_final_d"),
        )

    return {
        "energy_weight":                    energy_weight,
        "lambda_energy":                    lambda_energy,
        "lambda_temperature":               lambda_temperature,
        "temperature_weight":               temperature_weight,
        "lambda_energy_cost":               lambda_energy_cost,
        "comfort_threshold":                comfort_threshold,
        "range_comfort_hours":              range_comfort_hours,
        "occupied_hours":                   occupied_hours,
        "occupied_discomfort_multiplier":   occupied_discomfort_multiplier,
        "off_hours_energy_multiplier":      off_hours_energy_multiplier,
        "occupied_comfort_multiplier":      occupied_comfort_multiplier,
        "unoccupied_comfort_multiplier":    unoccupied_comfort_multiplier,
        "energy_cost_weight":               energy_cost_weight,
        "grid_energy_weight":               grid_energy_weight,
        "battery_cycle_weight":             battery_cycle_weight,
        "battery_loss_weight":              battery_loss_weight,
        "simultaneous_battery_weight":      simultaneous_battery_weight,
        "lambda_grid":                      lambda_grid,
        "lambda_battery":                   lambda_battery,
        "grid_energy_variables":            grid_energy_variables,
        "battery_charge_variables":         battery_charge_variables,
        "battery_discharge_variables":      battery_discharge_variables,
        "battery_loss_variables":           battery_loss_variables,
        "range_winter":                     range_winter,
        "range_summer":                     range_summer,
        "summer_start_m":                   summer_start_m,
        "summer_start_d":                   summer_start_d,
        "summer_final_m":                   summer_final_m,
        "summer_final_d":                   summer_final_d,
    }
\end{lstlisting}

\Needspace{8\baselineskip}

\subsection{\texorpdfstring{\texttt{page\_components/training\_shared.py}}{page\_components/training\_shared.py}}

\begin{lstlisting}[style=studioPython,caption={page\_components/training\_shared.py},label={lst:page-components-training-shared-py}]
"""Constants de pàgina d'entrenament compartides i components de capçalera reutilitzables."""

from __future__ import annotations

from html import escape
from pathlib import Path

import streamlit as st

from backend.entrenar_agent_artifacts import list_saved_training_runs
from backend.entrenar_agent_constants import TRAINING_RESULT_KEY
from backend.entrenar_agent_session import (
    TRAINING_LOADED_ARTIFACT_KEY,
    TRAINING_SAVED_RUN_KEY,
    apply_saved_training_ui_state,
)
from page_styles.training import inject_training_styles

PAGE_TITLE = "Entrenar model RL"
PAGE_LAYOUT = "wide"
INTRODUCTION_TEXT = (
    "Configura l'entorn, la reward, els wrappers i els parametres de l'agent "
    "per llançar entrenaments SB3 amb una presentacio mes clara i ordenada."
)

# ── Components de renderització ───────────────────────────────────────────────


def render_training_hero() -> None:
    """Mostra el bloc de capçalera principal de la pàgina a la UI de Streamlit."""
    # Hero entrenament
    st.markdown(
        f"""
        <section class="training-hero">
            <div class="hero-kicker">Entrenament i configuracio</div>
            <h1 class="hero-title">{escape(PAGE_TITLE)}</h1>
            <div class="hero-copy">{escape(INTRODUCTION_TEXT)}</div>
        </section>
        """,
        unsafe_allow_html=True,
    )


def render_training_section(title: str, kicker: str, description: str) -> None:
    """Crea la capçalera visual d'una secció de pàgina a la UI de Streamlit."""
    # Titol seccio entrenament
    st.markdown(f'<h2 class="page-section-title">{escape(title)}</h2>', unsafe_allow_html=True)


def format_saved_training_label(run: dict) -> str:
    """Retorna una etiqueta llegible per persones per a una cursa d'entrenament desada."""
    created_at = run.get("created_at") or "-"
    env_id = run.get("env_id") or "-"
    reward_name = run.get("reward_name") or "-"
    return f"{run.get('artifact_name', 'training')} | {env_id} | {reward_name} | {created_at}"


def render_saved_training_field(label: str, value: object) -> None:
    """Mostra un únic camp etiquetat dins de la biblioteca de entrenament desada."""
    # Camp entrenament guardat
    st.markdown(
        f"""
        <div class="training-saved-field">
            <div class="training-saved-label">{escape(label)}</div>
            <div class="training-saved-value">{escape(str(value or "-"))}</div>
        </div>
        """,
        unsafe_allow_html=True,
    )


def render_saved_training_library(envs: list[str]) -> None:
    """Mostra la biblioteca d'entrenaments desats amb controls d'upload i download."""
    render_training_section(
        "Scripts guardats",
        "Historial",
        "Recarrega configuracions anteriors i descarrega l'script literal de cada experiment.",
    )

    saved_runs = list_saved_training_runs()
    last_saved_run = st.session_state.get(TRAINING_RESULT_KEY)
    if last_saved_run:
        # Missatge entrenament guardat
        st.success(f"S'ha guardat l'entrenament `{last_saved_run.get('artifact_name', 'training')}`.")

    if not saved_runs:
        # Avís biblioteca buida
        st.info("Encara no hi ha scripts d'entrenament guardats.")
        return

    # Fem servir artifact_name com a clau estable: és el que també apareix al nom del model
    # i evita confondre runs amb el mateix entorn però configuracions diferents.
    run_map = {run["artifact_name"]: run for run in saved_runs if run.get("artifact_name")}
    run_names = list(run_map.keys())
    selected_run_name = st.session_state.get(TRAINING_SAVED_RUN_KEY)
    if selected_run_name not in run_map and run_names:
        st.session_state[TRAINING_SAVED_RUN_KEY] = run_names[0]

    # Selector entrenament guardat
    selected_run_name = st.selectbox(
        "Entrenament guardat",
        run_names,
        key=TRAINING_SAVED_RUN_KEY,
        format_func=lambda name: format_saved_training_label(run_map[name]),
    )
    selected_run = run_map[selected_run_name]

    # Camps resum entrenament
    info_col1, info_col2 = st.columns(2)
    with info_col1:
        render_saved_training_field("Entorn", selected_run.get("env_id", "-"))
    with info_col2:
        render_saved_training_field("Reward", selected_run.get("reward_name", "-"))

    script_path = Path(selected_run.get("training_script_path", ""))
    # Botons entrenament guardat
    controls_col1, controls_col2 = st.columns(2)
    # Botó carregar configuracio
    if controls_col1.button("Carregar configuracio", width="stretch"):
        ui_state = selected_run.get("ui_state") or {}
        if not ui_state:
            # Avís configuracio absent
            st.warning("Aquest entrenament no te una configuracio recarregable.")
        else:
            # Recarregar només toca l'estat dels widgets; l'entrenament encara no comença
            # fins que l'usuari prem el botó principal.
            apply_saved_training_ui_state(ui_state, envs)
            st.session_state[TRAINING_LOADED_ARTIFACT_KEY] = selected_run.get("artifact_name")
            st.rerun()

    if script_path.is_file():
        # Botó download script
        controls_col2.download_button(
            "Descarregar script",
            data=script_path.read_bytes(),
            file_name=script_path.name,
            mime="text/x-python",
            width="stretch",
        )
    else:
        # Avís script absent
        controls_col2.warning("No s'ha trobat el fitxer .py d'aquest entrenament.")

    loaded_artifact = st.session_state.get(TRAINING_LOADED_ARTIFACT_KEY)
    if loaded_artifact:
        # Missatge configuracio carregada
        st.info(f"Configuracio carregada a la UI: `{loaded_artifact}`")
\end{lstlisting}

\Needspace{8\baselineskip}

\subsection{\texorpdfstring{\texttt{page\_components/training\_wrappers.py}}{page\_components/training\_wrappers.py}}

\begin{lstlisting}[style=studioPython,caption={page\_components/training\_wrappers.py},label={lst:page-components-training-wrappers-py}]
"""Controls de wrappers i registre per a la pàgina d'entrenament."""

from __future__ import annotations

import streamlit as st

from backend.entrenar_agent_constants import DEFAULT_FORECAST_COLUMNS, ENERGY_COST_DIR
from backend.entrenar_agent_env import list_files
from backend.entrenar_agent_session import (
    reset_widget_state_if_disabled,
    sanitize_multiselect_state,
    seed_discrete_incremental_widget_defaults,
    seed_incremental_widget_defaults,
    training_form_key,
)


def _checkbox_kwargs(field: str, default: bool = False, *, disabled: bool = False) -> dict:
    """Prepara els arguments comuns dels checkboxes del formulari d'entrenament."""
    key = training_form_key(field)
    kwargs = {"disabled": disabled, "key": key}
    if key not in st.session_state:
        kwargs["value"] = default
    return kwargs


def render_observation_wrappers(
    observation_variables: list[str],
    context_variables: list[str],
    candidate_temp_vars: list[str],
    candidate_setpoint_vars: list[str],
) -> dict:
    """Pinta els controls que modifiquen les observacions abans d'arribar a l'agent."""
    # Seccio wrappers observacio
    st.caption("OBSERVACIO")

    # Streamlit conserva valors antics encara que canviï l'entorn; sanejar-los aquí
    # evita que un multiselect apunti a variables que ja no existeixen.
    sanitize_multiselect_state("previous_variables", observation_variables, candidate_temp_vars[:1])
    sanitize_multiselect_state("weather_forecast_columns", DEFAULT_FORECAST_COLUMNS, list(DEFAULT_FORECAST_COLUMNS))
    sanitize_multiselect_state("delta_temp_variables", observation_variables, candidate_temp_vars)
    sanitize_multiselect_state("delta_setpoint_variables", observation_variables, candidate_setpoint_vars[:1])
    sanitize_multiselect_state("reduced_observations", observation_variables, [])

    # Checkbox datetime wrapper
    enable_datetime_wrapper = st.checkbox(
        "DatetimeWrapper",
        False,
        key=training_form_key("enable_datetime_wrapper"),
    )

    # Checkbox previous observation
    enable_previous_wrapper = st.checkbox(
        "PreviousObservationWrapper",
        False,
        key=training_form_key("enable_previous_wrapper"),
    )
    previous_variables = []
    if enable_previous_wrapper:
        # Selector variables pas anterior
        previous_variables = st.multiselect(
            "Variables pas anterior",
            observation_variables,
            default=candidate_temp_vars[:1],
            key=training_form_key("previous_variables"),
        )

    # Checkbox multi observacio
    enable_multi_obs_wrapper = st.checkbox(
        "MultiObsWrapper",
        False,
        key=training_form_key("enable_multi_obs_wrapper"),
    )
    multi_obs_n = 5
    multi_obs_flatten = True
    if enable_multi_obs_wrapper:
        # Slider observacions apilades
        multi_obs_n = st.slider(
            "Observacions apilades",
            2,
            12,
            5,
            key=training_form_key("multi_obs_n"),
        )
        # Checkbox aplanar observacio
        multi_obs_flatten = st.checkbox(
            "Aplanar observacio",
            True,
            key=training_form_key("multi_obs_flatten"),
        )

    # Checkbox normalitzar observacio
    enable_normalize_obs_wrapper = st.checkbox(
        "NormalizeObservation",
        False,
        key=training_form_key("enable_normalize_obs_wrapper"),
    )
    normalize_obs_auto = True
    normalize_obs_epsilon = 1e-8
    normalize_obs_mean = ""
    normalize_obs_var = ""
    if enable_normalize_obs_wrapper:
        # Checkbox actualitzar normalitzacio
        normalize_obs_auto = st.checkbox(
            "Actualitzacio automatica mean/var",
            True,
            key=training_form_key("normalize_obs_auto"),
        )
        # Input epsilon normalitzacio
        normalize_obs_epsilon = st.number_input(
            "epsilon",
            1e-10,
            1e-2,
            1e-8,
            format="%.10f",
            key=training_form_key("normalize_obs_epsilon"),
        )
        # Input ruta mean
        normalize_obs_mean = st.text_input(
            "Ruta mean.txt (opcional)",
            "",
            key=training_form_key("normalize_obs_mean"),
        )
        # Input ruta var
        normalize_obs_var = st.text_input(
            "Ruta var.txt (opcional)",
            "",
            key=training_form_key("normalize_obs_var"),
        )

    # Checkbox previsio clima
    enable_weather_forecast_wrapper = st.checkbox(
        "WeatherForecastingWrapper",
        False,
        key=training_form_key("enable_weather_forecast_wrapper"),
    )
    weather_forecast_n = 5
    weather_forecast_delta = 1
    weather_forecast_columns = list(DEFAULT_FORECAST_COLUMNS)
    if enable_weather_forecast_wrapper:
        # Slider mostres previsio
        weather_forecast_n = st.slider(
            "Mostres de previsi",
            1,
            24,
            5,
            key=training_form_key("weather_forecast_n"),
        )
        # Slider separacio previsio
        weather_forecast_delta = st.slider(
            "Separacio entre mostres",
            1,
            24,
            1,
            key=training_form_key("weather_forecast_delta"),
        )
        # Selector variables previsio
        weather_forecast_columns = st.multiselect(
            "Variables de previsi",
            DEFAULT_FORECAST_COLUMNS,
            default=DEFAULT_FORECAST_COLUMNS,
            key=training_form_key("weather_forecast_columns"),
        )

    # Checkbox delta temperatura
    enable_delta_temp_wrapper = st.checkbox(
        "DeltaTempWrapper",
        False,
        key=training_form_key("enable_delta_temp_wrapper"),
    )
    delta_temp_variables = []
    delta_setpoint_variables = []
    if enable_delta_temp_wrapper:
        # Selector variables temperatura
        delta_temp_variables = st.multiselect(
            "Variables temperatura",
            observation_variables,
            default=candidate_temp_vars,
            key=training_form_key("delta_temp_variables"),
        )
        # Selector variables setpoint
        delta_setpoint_variables = st.multiselect(
            "Variables setpoint",
            observation_variables,
            default=candidate_setpoint_vars[:1],
            key=training_form_key("delta_setpoint_variables"),
        )

    # Checkbox reduir observacions
    enable_reduce_obs_wrapper = st.checkbox(
        "ReduceObservationWrapper",
        False,
        key=training_form_key("enable_reduce_obs_wrapper"),
    )
    reduced_observations = []
    if enable_reduce_obs_wrapper:
        # Selector observacions excloses
        reduced_observations = st.multiselect(
            "Variables a excloure",
            observation_variables,
            default=[],
            key=training_form_key("reduced_observations"),
        )

    variability_disabled = not bool(context_variables)
    # Si l'entorn no exposa variables de context, el wrapper es desactiva i també netegem
    # l'estat antic perquè no torni a aparèixer marcat en una reexecució.
    reset_widget_state_if_disabled("enable_variability_context_wrapper", variability_disabled)
    # Checkbox variabilitat context
    enable_variability_context_wrapper = st.checkbox(
        "VariabilityContextWrapper",
        **_checkbox_kwargs("enable_variability_context_wrapper", disabled=variability_disabled),
    )
    variability_context_low = []
    variability_context_high = []
    variability_delta_value = 1.0
    variability_step_frequency_min = 96
    variability_step_frequency_max = 96 * 7
    if enable_variability_context_wrapper:
        # Seccio rang context
        st.caption("Rang de variabilitat del context")
        for context_variable in context_variables:
            # Inputs rang context
            low_col, high_col = st.columns(2)
            with low_col:
                variability_context_low.append(
                    # Input context minim
                    st.number_input(
                        f"{context_variable} min",
                        value=0.0,
                        key=training_form_key(f"variability_context_low_{context_variable}"),
                    )
                )
            with high_col:
                variability_context_high.append(
                    # Input context maxim
                    st.number_input(
                        f"{context_variable} max",
                        value=2.0,
                        key=training_form_key(f"variability_context_high_{context_variable}"),
                    )
                )
        # Input delta context
        variability_delta_value = st.number_input(
            "Delta context",
            min_value=0.01,
            value=1.0,
            step=0.1,
            key=training_form_key("variability_delta_value"),
        )
        # Inputs frequencia context
        freq_col1, freq_col2 = st.columns(2)
        with freq_col1:
            # Input frequencia minima
            variability_step_frequency_min = st.number_input(
                "Freq. min passos",
                min_value=1,
                value=96,
                step=1,
                key=training_form_key("variability_step_frequency_min"),
            )
        with freq_col2:
            # Input frequencia maxima
            variability_step_frequency_max = st.number_input(
                "Freq. max passos",
                min_value=2,
                value=96 * 7,
                step=1,
                key=training_form_key("variability_step_frequency_max"),
            )

    return {
        "enable_datetime_wrapper": enable_datetime_wrapper,
        "enable_previous_wrapper": enable_previous_wrapper,
        "previous_variables": previous_variables,
        "enable_multi_obs_wrapper": enable_multi_obs_wrapper,
        "multi_obs_n": multi_obs_n,
        "multi_obs_flatten": multi_obs_flatten,
        "enable_normalize_obs_wrapper": enable_normalize_obs_wrapper,
        "normalize_obs_auto": normalize_obs_auto,
        "normalize_obs_epsilon": normalize_obs_epsilon,
        "normalize_obs_mean": normalize_obs_mean,
        "normalize_obs_var": normalize_obs_var,
        "enable_weather_forecast_wrapper": enable_weather_forecast_wrapper,
        "weather_forecast_n": weather_forecast_n,
        "weather_forecast_delta": weather_forecast_delta,
        "weather_forecast_columns": weather_forecast_columns,
        "enable_delta_temp_wrapper": enable_delta_temp_wrapper,
        "delta_temp_variables": delta_temp_variables,
        "delta_setpoint_variables": delta_setpoint_variables,
        "enable_reduce_obs_wrapper": enable_reduce_obs_wrapper,
        "reduced_observations": reduced_observations,
        "enable_variability_context_wrapper": enable_variability_context_wrapper,
        "variability_context_low": variability_context_low,
        "variability_context_high": variability_context_high,
        "variability_delta_value": variability_delta_value,
        "variability_step_frequency_min": variability_step_frequency_min,
        "variability_step_frequency_max": variability_step_frequency_max,
    }


def render_action_wrappers(
    env_meta: dict,
    action_variables: list[str],
    action_space,
) -> dict:
    """Pinta els controls que canvien com l'agent envia accions a l'entorn."""
    # Seccio wrappers accio
    st.caption("ACCIO")

    # Els wrappers incrementals només tenen sentit quan l'acció és contínua: converteixen
    # una ordre absoluta en canvis petits, més semblants a moure un termòstat real.
    sanitize_multiselect_state("incremental_variables", action_variables, action_variables[:1])
    # Checkbox incremental wrapper
    enable_incremental_wrapper = st.checkbox(
        "IncrementalWrapper",
        False,
        disabled=not env_meta["is_continuous"],
        key=training_form_key("enable_incremental_wrapper"),
    )
    incremental_variables = []
    incremental_definition = {}
    incremental_initial_values = []
    if enable_incremental_wrapper:
        seed_incremental_widget_defaults(action_variables)
        # Selector accions incrementals
        incremental_variables = st.multiselect(
            "Variables d'accio incrementals",
            action_variables,
            default=action_variables[:1],
            key=training_form_key("incremental_variables"),
        )
        for variable_name in incremental_variables:
            action_index = action_variables.index(variable_name)
            low_value = float(action_space.low[action_index])
            high_value = float(action_space.high[action_index])
            default_initial = float((low_value + high_value) / 2.0)
            # Titol variable incremental
            st.markdown(f"**{variable_name}**")
            # Inputs accio incremental
            init_col, delta_col, step_col = st.columns(3)
            with init_col:
                # Input valor inicial
                initial_value = st.number_input(
                    "Valor inicial",
                    value=default_initial,
                    key=training_form_key(f"incr_init_{variable_name}"),
                )
            with delta_col:
                # Input delta max
                delta_value = st.number_input(
                    "Delta max",
                    min_value=0.1,
                    value=max((high_value - low_value) / 4.0, 0.5),
                    step=0.1,
                    key=training_form_key(f"incr_delta_{variable_name}"),
                )
            with step_col:
                # Input pas incremental
                step_value = st.number_input(
                    "Pas",
                    min_value=0.1,
                    value=max((high_value - low_value) / 20.0, 0.1),
                    step=0.1,
                    key=training_form_key(f"incr_step_{variable_name}"),
                )
            incremental_definition[variable_name] = (delta_value, step_value)
            incremental_initial_values.append(initial_value)
        st.session_state[training_form_key("incremental_definition")] = incremental_definition
        st.session_state[training_form_key("incremental_initial_values")] = incremental_initial_values

    # Checkbox incremental discret
    enable_discrete_incremental_wrapper = st.checkbox(
        "DiscreteIncrementalWrapper",
        False,
        disabled=not env_meta["is_continuous"],
        key=training_form_key("enable_discrete_incremental_wrapper"),
    )
    discrete_incremental_initial_values = []
    discrete_incremental_delta = 2.0
    discrete_incremental_step = 0.5
    if enable_discrete_incremental_wrapper:
        seed_discrete_incremental_widget_defaults(action_variables)
        for action_index, variable_name in enumerate(action_variables):
            low_value = float(action_space.low[action_index])
            high_value = float(action_space.high[action_index])
            default_initial = float((low_value + high_value) / 2.0)
            discrete_incremental_initial_values.append(
                # Input valor inicial discret
                st.number_input(
                    f"Valor inicial {variable_name}",
                    value=default_initial,
                    key=training_form_key(f"disc_incr_init_{variable_name}"),
                )
            )
        # Input delta discret
        discrete_incremental_delta = st.number_input(
            "Delta discret max",
            min_value=0.1,
            value=2.0,
            step=0.1,
            key=training_form_key("discrete_incremental_delta"),
        )
        # Input pas discret
        discrete_incremental_step = st.number_input(
            "Pas discret",
            min_value=0.1,
            value=0.5,
            step=0.1,
            key=training_form_key("discrete_incremental_step"),
        )
        st.session_state[training_form_key("discrete_incremental_initial_values")] = (
            discrete_incremental_initial_values
        )

    normalize_action_disabled = (
        not env_meta["is_continuous"]
        or enable_incremental_wrapper
        or enable_discrete_incremental_wrapper
    )
    # NormalizeAction és còmode com a default en Box, però no el deixem conviure amb
    # wrappers incrementals perquè llavors l'escala de control queda doblement transformada.
    normalize_action_key = training_form_key("enable_normalize_action")
    normalize_action_seed_key = training_form_key("enable_normalize_action_default_seeded")
    reset_widget_state_if_disabled("enable_normalize_action", normalize_action_disabled)
    if normalize_action_disabled:
        st.session_state.pop(normalize_action_seed_key, None)
    elif (
        not st.session_state.get(normalize_action_seed_key)
        and normalize_action_key not in st.session_state
    ):
        st.session_state[normalize_action_key] = True
        st.session_state[normalize_action_seed_key] = True
    # Checkbox normalitzar accio
    enable_normalize_action = st.checkbox(
        "NormalizeAction",
        **_checkbox_kwargs(
            "enable_normalize_action",
            default=not normalize_action_disabled,
            disabled=normalize_action_disabled,
        ),
    )
    if normalize_action_disabled:
        enable_normalize_action = False

    storage_smoothing_disabled = env_meta.get("building_name") != "OfficeGridStorageSmoothing.epJSON"
    # Aquest wrapper està escrit per a un edifici concret amb bateria i restriccions de xarxa;
    # en la resta d'entorns el mostrem desactivat per evitar errors difícils d'entendre.
    reset_widget_state_if_disabled("enable_storage_smoothing_wrapper", storage_smoothing_disabled)
    # Checkbox storage smoothing
    enable_storage_smoothing_wrapper = st.checkbox(
        "OfficeGridStorageSmoothingActionConstraintsWrapper",
        **_checkbox_kwargs(
            "enable_storage_smoothing_wrapper",
            disabled=storage_smoothing_disabled,
        ),
    )

    return {
        "enable_incremental_wrapper": enable_incremental_wrapper,
        "incremental_variables": incremental_variables,
        "incremental_definition": incremental_definition,
        "incremental_initial_values": incremental_initial_values,
        "enable_discrete_incremental_wrapper": enable_discrete_incremental_wrapper,
        "discrete_incremental_initial_values": discrete_incremental_initial_values,
        "discrete_incremental_delta": discrete_incremental_delta,
        "discrete_incremental_step": discrete_incremental_step,
        "enable_normalize_action": enable_normalize_action,
        "enable_storage_smoothing_wrapper": enable_storage_smoothing_wrapper,
    }


def render_energy_logging_wrappers(
    energy_cost_files: list[str],
    initial_energy_cost_path: str,
    initial_file_logger_name: str,
) -> dict:
    """Pinta els controls de cost energètic i de fitxer de registre."""
    # Seccio energia i logs
    st.caption("ENERGIA I LOGS")

    energy_cost_path = initial_energy_cost_path
    file_logger_name = initial_file_logger_name

    # Checkbox energy cost
    use_energy_cost = st.checkbox(
        "EnergyCostWrapper",
        False,
        key=training_form_key("use_energy_cost"),
    )
    if use_energy_cost:
        if energy_cost_files:
            # Selector fitxer cost
            selected_energy_cost = st.selectbox(
                "Fitxer cost",
                energy_cost_files,
                index=energy_cost_files.index(energy_cost_path) if energy_cost_path in energy_cost_files else 0,
            )
            energy_cost_path = selected_energy_cost
        # Input ruta cost energia
        energy_cost_path = st.text_input(
            "Ruta fitxer cost",
            energy_cost_path,
            key=training_form_key("energy_cost_path"),
        )

    # Checkbox file logger
    use_file_logger = st.checkbox(
        "EnergyCostFileLogger",
        False,
        key=training_form_key("use_file_logger"),
    )
    if use_file_logger:
        # Input nom log energia
        file_logger_name = st.text_input(
            "Nom fitxer registre",
            "energy_cost_log.csv",
            key=training_form_key("file_logger_name"),
        )

    return {
        "energy_cost_path": energy_cost_path,
        "file_logger_name": file_logger_name,
        "use_energy_cost": use_energy_cost,
        "use_file_logger": use_file_logger,
    }


def render_wrappers_section(
    env_meta: dict,
    observation_variables: list[str],
    action_variables: list[str],
    action_space,
    context_variables: list[str],
    candidate_temp_vars: list[str],
    candidate_setpoint_vars: list[str],
) -> dict:
    """Agrupa els controls de wrappers i retorna una configuració plana per al backend."""
    energy_cost_files = list_files(ENERGY_COST_DIR, {".csv"})
    default_energy_cost_path = (
        energy_cost_files[0] if energy_cost_files
        else str(ENERGY_COST_DIR / "PVPC_active_energy_billing_Iberian_Peninsula_2023.csv")
    )
    default_file_logger_name = "energy_cost_log.csv"
    initial_energy_cost_path = st.session_state.get(
        training_form_key("energy_cost_path"),
        default_energy_cost_path,
    )
    initial_file_logger_name = st.session_state.get(
        training_form_key("file_logger_name"),
        default_file_logger_name,
    )

    # Columnes wrappers observacio accio
    observation_column, action_column = st.columns(2)
    with observation_column:
        observation_wrapper_values = render_observation_wrappers(
            observation_variables=observation_variables,
            context_variables=context_variables,
            candidate_temp_vars=candidate_temp_vars,
            candidate_setpoint_vars=candidate_setpoint_vars,
        )
    with action_column:
        action_wrapper_values = render_action_wrappers(
            env_meta=env_meta,
            action_variables=action_variables,
            action_space=action_space,
        )
        energy_logging_values = render_energy_logging_wrappers(
            energy_cost_files=energy_cost_files,
            initial_energy_cost_path=initial_energy_cost_path,
            initial_file_logger_name=initial_file_logger_name,
        )

    # Es retorna un diccionari pla perquè el muntatge final de training pugui barrejar
    # reward, wrappers i paràmetres de l'agent sense dependre de l'ordre visual de la pàgina.
    return {
        **energy_logging_values,
        **observation_wrapper_values,
        **action_wrapper_values,
    }
\end{lstlisting}

\Needspace{8\baselineskip}

\subsection{\texorpdfstring{\texttt{page\_components/ui\_fragments.py}}{page\_components/ui\_fragments.py}}

\begin{lstlisting}[style=studioPython,caption={page\_components/ui\_fragments.py},label={lst:page-components-ui-fragments-py}]
"""Fragments UI compartits per pàgines Streamlit de BEMS-RL Studio."""

from __future__ import annotations

from collections.abc import Callable
from html import escape
import time
from typing import Any, Iterable

import streamlit as st


def render_hero(class_name: str, kicker: str, title: str, copy_text: str | Iterable[str]) -> None:
    """Crea una capçalera hero amb les classes CSS de la pàgina."""
    copy_html = _copy_html(copy_text)
    # Hero pagina
    st.markdown(
        f"""
        <section class="{escape(class_name)}">
            <div class="hero-kicker">{escape(kicker)}</div>
            <h1 class="hero-title">{escape(title)}</h1>
            <div class="hero-copy">{copy_html}</div>
        </section>
        """,
        unsafe_allow_html=True,
    )


def render_section_title(title: str, class_name: str = "page-section-title", level: int = 2) -> None:
    """Mostra un títol de secció escapant el text."""
    tag = f"h{level}"
    # Titol seccio
    st.markdown(f'<{tag} class="{escape(class_name)}">{escape(title)}</{tag}>', unsafe_allow_html=True)


def render_section_card(title: str, description: str, *, title_class: str, card_class: str) -> None:
    """Combina títol de secció i una targeta descriptiva simple."""
    render_section_title(title, title_class)
    # Targeta seccio
    st.markdown(
        f"""
        <section class="{escape(card_class)} section-card-spaced">
            <div class="section-copy">{escape(description)}</div>
        </section>
        """,
        unsafe_allow_html=True,
    )


def render_card_header(
    target: Any,
    *,
    anchor_class: str,
    kicker: str,
    title: str,
    description: str = "",
    kicker_class: str = "upload-card-kicker",
    title_class: str = "upload-card-title",
    description_class: str = "upload-card-copy",
) -> None:
    """Pinta la capçalera HTML comuna d'una targeta Streamlit."""
    description_html = (
        f'<div class="{escape(description_class)}">{escape(description)}</div>'
        if description
        else ""
    )
    # Capcalera targeta
    target.markdown(
        f"""
        <div class="{escape(anchor_class)}"></div>
        <div class="{escape(kicker_class)}">{escape(kicker)}</div>
        <div class="{escape(title_class)}">{escape(title)}</div>
        {description_html}
        """,
        unsafe_allow_html=True,
    )


def render_info_panel(class_name: str, title: str, kicker: str, copy_text: str) -> None:
    """Mostra un panell informatiu compacte."""
    # Panell informatiu
    st.markdown(
        f"""
        <section class="{escape(class_name)}">
            <div class="panel-kicker">{escape(kicker)}</div>
            <div class="panel-title">{escape(title)}</div>
            <div class="panel-copy">{escape(copy_text)}</div>
        </section>
        """,
        unsafe_allow_html=True,
    )


def render_metric_card_grid(
    metrics: Iterable[tuple[Any, Any]],
    *,
    shell_class: str,
    card_class: str,
    label_class: str,
    value_class: str,
    formatter: Callable[[Any], str] | None = None,
    label_formatter: Callable[[Any], str] | None = None,
) -> None:
    """Mostra una graella HTML de targetes KPI compactes."""
    metric_items = list(metrics)
    if not metric_items:
        return

    format_value = formatter or str
    format_label = label_formatter or str
    cards = "".join(
        (
            f'<section class="{escape(card_class)}">'
            f'<div class="{escape(label_class)}">{escape(format_label(label))}</div>'
            f'<div class="{escape(value_class)}">{escape(format_value(value))}</div>'
            "</section>"
        )
        for label, value in metric_items
    )
    # Graella targetes KPI
    st.markdown(
        f'<section class="{escape(shell_class)}">{cards}</section>',
        unsafe_allow_html=True,
    )


def render_detail_list_card(
    *,
    title: str,
    rows: Iterable[tuple[Any, Any]],
    card_class: str,
    title_class: str,
    list_class: str,
    row_class: str,
    label_class: str,
    value_class: str,
    formatter: Callable[[Any], str] | None = None,
    label_formatter: Callable[[Any], str] | None = None,
) -> None:
    """Mostra una targeta de detalls amb parelles etiqueta/valor."""
    format_value = formatter or str
    format_label = label_formatter or str
    row_html = "".join(
        f"""
        <div class="{escape(row_class)}">
            <div class="{escape(label_class)}">{escape(format_label(label))}</div>
            <div class="{escape(value_class)}">{escape(format_value(value))}</div>
        </div>
        """
        for label, value in rows
    )
    # Targeta llista detall
    st.markdown(
        f"""
        <section class="{escape(card_class)}">
            <div class="{escape(title_class)}">{escape(title)}</div>
            <div class="{escape(list_class)}">{row_html}</div>
        </section>
        """,
        unsafe_allow_html=True,
    )


def render_copy_block(class_name: str, text: str, *, formatter: Callable[[str], str] | None = None) -> None:
    """Mostra un bloc de text HTML escapant-ne el contingut."""
    display_text = formatter(text) if formatter is not None else text
    # Text informatiu
    st.markdown(
        f"<div class='{escape(class_name)}'>{escape(display_text)}</div>",
        unsafe_allow_html=True,
    )


def render_metric_row(
    items: Iterable[dict[str, Any]],
    *,
    columns: int = 3,
) -> None:
    """Mostra mètriques Streamlit en columnes amb fallback informatiu."""
    item_list = list(items)
    if not item_list:
        return
    cols = st.columns(columns)
    for index, item in enumerate(item_list):
        with cols[index % columns]:
            empty_message = item.get("empty_message")
            if item.get("empty") and empty_message:
                # Avís mètrica buida
                st.info(str(empty_message))
                continue
            # Metrica fila
            st.metric(
                str(item.get("label", "")),
                str(item.get("value", "")),
                delta=item.get("delta"),
                delta_color=item.get("delta_color", "normal"),
            )


def build_metric_item(
    current: dict,
    previous: dict | None,
    field: str | None,
    *,
    label: str,
    value_format: Callable[[Any], str],
    delta_format: Callable[[Any], str],
    empty_message: str,
    delta_color: str = "off",
) -> dict[str, Any]:
    """Prepara una mètrica declarativa a partir d'un camp opcional."""
    if not field or field not in current:
        return {"empty": True, "empty_message": empty_message}
    value = current[field]
    old_value = previous[field] if previous and field in previous else value
    return {
        "label": label,
        "value": value_format(value),
        "delta": delta_format(value - old_value),
        "delta_color": delta_color,
    }


def render_kicker_section(class_name: str, title: str, kicker: str, description: str) -> None:
    """Munta una secció amb títol, kicker i text descriptiu."""
    render_section_title(title)
    # Bloc seccio amb kicker
    st.markdown(
        f"""
        <section class="{escape(class_name)}">
            <div class="section-kicker">{escape(kicker)}</div>
            <div class="section-copy">{escape(description)}</div>
        </section>
        """,
        unsafe_allow_html=True,
    )


def render_runtime_progress(
    runtime: dict[str, Any],
    *,
    progress_label: str = "Progrés",
    freeze_from_result: bool = False,
) -> None:
    """Mostra barra i mètriques de progrés per a un runtime en segon pla."""
    steps_target = max(int(runtime.get("steps_target") or 1), 1)
    latest_step = min(int(runtime.get("latest_step") or 0), steps_target)
    progress_value = min(latest_step / steps_target, 1.0)
    elapsed_seconds = _elapsed_seconds(runtime, freeze_from_result=freeze_from_result)

    # Barra progres runtime
    st.progress(progress_value)
    # Metriques runtime
    metric_cols = st.columns(3)
    metric_cols[0].metric("Passos", f"{latest_step}/{steps_target}")
    metric_cols[1].metric(progress_label, f"{progress_value * 100:.0f}%")
    metric_cols[2].metric("Temps", f"{elapsed_seconds // 60}m {elapsed_seconds % 60:02d}s")

    status_text = str(runtime.get("status") or "").strip()
    if status_text:
        # Text estat runtime
        st.caption(status_text)


def _copy_html(copy_text: str | Iterable[str]) -> str:
    """Escapa text o paràgrafs per inserir-los dins d'un bloc HTML segur."""
    if isinstance(copy_text, str):
        return escape(copy_text)
    return "<br><br>".join(escape(str(part)) for part in copy_text)


def _elapsed_seconds(runtime: dict[str, Any], *, freeze_from_result: bool) -> int:
    """Calcula segons transcorreguts a partir de l'estat runtime o del resultat final."""
    result = runtime.get("result")
    if freeze_from_result and result is not None:
        return int(float(getattr(result, "elapsed_seconds", 0.0)))
    if freeze_from_result and runtime.get("frozen_elapsed_seconds") is not None:
        return int(float(runtime.get("frozen_elapsed_seconds") or 0.0))
    started_at = runtime.get("started_at")
    if started_at:
        return max(0, int(time.time() - float(started_at)))
    return 0
\end{lstlisting}

\section{Codi: Backend funcional}

\Needspace{8\baselineskip}

\subsection{\texorpdfstring{\texttt{backend/\_\_init\_\_.py}}{backend/\_\_init\_\_.py}}

\begin{lstlisting}[style=studioPython,caption={backend/\_\_init\_\_.py},label={lst:backend---init---py}]
"""Lògica d'aplicació de BEMS-RL Studio segura per importar.

El paquet backend conté el codi que dona suport a la interfície Streamlit:
creació d'entorns, anàlisi climàtica, orquestració de l'entrenament, avaluació,
control en directe, agregació de resultats i generació d'informes. Els mòduls
d'aquest paquet eviten renderitzar la UI directament perquè es puguin importar
des de proves, scripts i Sphinx autodoc.
"""
\end{lstlisting}

\Needspace{8\baselineskip}

\subsection{\texorpdfstring{\texttt{backend/afegir\_entorn\_analysis.py}}{backend/afegir\_entorn\_analysis.py}}

\begin{lstlisting}[style=studioPython,caption={backend/afegir\_entorn\_analysis.py},label={lst:backend-afegir-entorn-analysis-py}]
"""Lectura i anàlisi d'edificis per preparar la configuració dels entorns."""

from __future__ import annotations

import csv
import io
import json
from pathlib import Path
from typing import Any, Dict, List, Optional, Sequence, Tuple

import gymnasium as gym
import numpy as np
from sinergym.envs.eplus_env import EplusEnv
from sinergym.utils.rewards import BaseReward

from backend.afegir_entorn_common import DEFAULT_TIME_VARIABLES
from backend.afegir_entorn_common import ActuatorOption, BuildingTrainingAnalysis
from backend.afegir_entorn_common import normalize_token, sanitize_identifier

class ProbeReward(BaseReward):
    """Reward nul usat només per inspeccionar handlers disponibles."""

    def __call__(self, obs_dict: Dict[str, Any]) -> Tuple[float, Dict[str, Any]]:
        """Executa l'objecte cridable."""
        return 0.0, {}


def load_building_json(building_path: Path) -> Dict[str, Any]:
    """Carrega i analitza un fitxer d'edifici de format epJSON en un diccionari.

    Paràmetres:
        building_path (Path): camí cap al fitxer epJSON.

    Retorna:
        Dict[str, Any]: el diccionari carregat que renderitza les dades de l'edifici.
    """
    with open(building_path, "r", encoding="utf-8") as file_handle:
        return json.load(file_handle)


def extract_schedule_type_index(building_data: Dict[str, Any]) -> Dict[str, str]:
    """Extreu una assignació de noms d'objectes de planificació als seus respectius tipus de blocs de planificació.

    Paràmetres:
        building_data (Dict[str, Any]): les dades de l'edifici epJSON carregades.

    Retorna:
        Dict[str, str]: un diccionari que assigna noms de planificació als seus tipus de bloc d'objectes
        (e.g., 'Programació:Compact').
    """
    schedule_types: Dict[str, str] = {}
    for object_type, objects in building_data.items():
        if not object_type.startswith("Schedule:") or not isinstance(objects, dict):
            continue
        for schedule_name in objects.keys():
            schedule_types[schedule_name] = object_type
    return schedule_types


def resolve_schedule_name(
    schedule_types: Dict[str, str],
    schedule_name: str,
) -> Optional[str]:
    """Resol les referències de planificació amb noms d'estil EnergyPlus que no distingeixen entre majúscules i minúscules."""
    if not schedule_name:
        return None
    if schedule_name in schedule_types:
        return schedule_name

    target_casefold = schedule_name.casefold()
    for actual_name in schedule_types.keys():
        if actual_name.casefold() == target_casefold:
            return actual_name

    target_normalized = normalize_token(schedule_name)
    for actual_name in schedule_types.keys():
        if normalize_token(actual_name) == target_normalized:
            return actual_name

    return None


def _register_thermostat_candidate(
    detected: Dict[Tuple[str, str], Dict[str, Any]],
    schedule_types: Dict[str, str],
    schedule_name: str,
    category: str,
    zone_name: Optional[str],
) -> None:
    """Registra una planificació de termòstat com a entrada de controlador candidat si encara no s'ha fet un seguiment.

    Crea una nova entrada a `detected` per a la clau donada (categoria, schedule_name) si no hi ha,
    aleshores hi associa el nom de la zona.

    Paràmetres:
        detected (Dict): Registre de candidats acumulat clausurat per (categoria, schedule_name).
        schedule_types (Dict[str, str]): Mapeig de noms de planificació amb les seves cadenes del tipus de bloc.
        schedule_name (str): el nom del programa EnergyPlus per registrar-se.
        category (str): categoria de consigna, e.g. 'heating_setpoint' o 'cooling_setpoint'.
        zone_name (Optional[str]): Nom de la zona a associar amb aquest candidat.
    """
    resolved_schedule_name = resolve_schedule_name(schedule_types, schedule_name)
    if resolved_schedule_name is None:
        return
    key = (category, resolved_schedule_name)
    if key not in detected:
        detected[key] = {
            "category": category,
            "schedule_name": resolved_schedule_name,
            "zones": set(),
            "references": set(),
            "source_kind": "thermostat",
            "element_type": schedule_types[resolved_schedule_name],
        }
    if zone_name:
        detected[key]["zones"].add(zone_name)
        detected[key]["references"].add(zone_name)


def extract_thermostat_schedule_candidates(
    building_data: Dict[str, Any],
    schedule_types: Dict[str, str],
) -> Tuple[List[str], List[Dict[str, Any]]]:
    """Detecta programes de termòstat implícits basats en els punts de consigna Dual i Single del model.

    Paràmetres:
        building_data (Dict[str, Any]): les dades epJSON carregades.
        schedule_types (Dict[str, str]): una assignació de noms de planificació a tipus de bloc de planificació.

    Retorna:
        Tuple[List[str], List[Dict[str, Any]]]: una tupla que conté una llista de detectats
        zones del termòstat i una llista de diccionaris estructurats amb detalls dels actuadors candidats.
    """
    thermostat_controls = building_data.get("ZoneControl:Thermostat", {})
    dual_setpoints = building_data.get("ThermostatSetpoint:DualSetpoint", {})
    single_heating = building_data.get("ThermostatSetpoint:SingleHeating", {})
    single_cooling = building_data.get("ThermostatSetpoint:SingleCooling", {})

    thermostat_zones: List[str] = []
    detected: Dict[Tuple[str, str], Dict[str, Any]] = {}

    for _, thermostat_spec in thermostat_controls.items():
        zone_name = thermostat_spec.get("zone_or_zonelist_name") or thermostat_spec.get("zone_name")
        if zone_name and zone_name not in thermostat_zones:
            thermostat_zones.append(zone_name)

        # EnergyPlus desa fins a cinc controls numerats per termòstat; cal revisar-los tots
        # perquè un mateix objecte pot combinar consignes Dual i Single.
        for control_index in range(1, 6):
            object_type = thermostat_spec.get(f"control_{control_index}_object_type")
            object_name = thermostat_spec.get(f"control_{control_index}_name")
            if not object_type or not object_name:
                continue

            if object_type == "ThermostatSetpoint:DualSetpoint":
                dual_spec = dual_setpoints.get(object_name, {})
                _register_thermostat_candidate(
                    detected,
                    schedule_types,
                    dual_spec.get("heating_setpoint_temperature_schedule_name", ""),
                    "heating_setpoint",
                    zone_name,
                )
                _register_thermostat_candidate(
                    detected,
                    schedule_types,
                    dual_spec.get("cooling_setpoint_temperature_schedule_name", ""),
                    "cooling_setpoint",
                    zone_name,
                )
            elif object_type == "ThermostatSetpoint:SingleHeating":
                single_spec = single_heating.get(object_name, {})
                _register_thermostat_candidate(
                    detected,
                    schedule_types,
                    single_spec.get("setpoint_temperature_schedule_name", ""),
                    "heating_setpoint",
                    zone_name,
                )
            elif object_type == "ThermostatSetpoint:SingleCooling":
                single_spec = single_cooling.get(object_name, {})
                _register_thermostat_candidate(
                    detected,
                    schedule_types,
                    single_spec.get("setpoint_temperature_schedule_name", ""),
                    "cooling_setpoint",
                    zone_name,
                )

    ordered_candidates: List[Dict[str, Any]] = []
    # Ordenem els candidats perquè la UI i els YAML generats siguin estables entre execucions.
    for key in sorted(detected.keys(), key=lambda item: (item[0], item[1].lower())):
        candidate = detected[key]
        ordered_candidates.append(
            {
                "schedule_name": candidate["schedule_name"],
                "category": candidate["category"],
                "zones": tuple(sorted(candidate["zones"])),
                "references": tuple(sorted(candidate["references"])),
                "source_kind": candidate["source_kind"],
                "element_type": candidate["element_type"],
            }
        )

    return thermostat_zones, ordered_candidates


def _collect_schedule_references_from_value(
    value: Any,
    schedule_types: Dict[str, str],
) -> List[str]:
    """Recolliu referències de calendari a partir del valor."""
    references: List[str] = []
    if isinstance(value, str):
        resolved_schedule_name = resolve_schedule_name(schedule_types, value)
        if resolved_schedule_name is not None:
            references.append(resolved_schedule_name)
    elif isinstance(value, dict):
        for nested_value in value.values():
            references.extend(_collect_schedule_references_from_value(nested_value, schedule_types))
    elif isinstance(value, list):
        for nested_value in value:
            references.extend(_collect_schedule_references_from_value(nested_value, schedule_types))
    return references


def extract_schedule_numeric_values(
    building_data: Dict[str, Any],
    schedule_types: Dict[str, str],
    schedule_name: str,
    visited: Optional[set[str]] = None,
) -> List[float]:
    """Extreu valors numèrics de planificació."""
    if visited is None:
        visited = set()

    # Les planificacions poden referenciar altres planificacions; aquest conjunt evita bucles
    # quan EnergyPlus encadena objectes Schedule:* de manera recursiva.
    resolved_schedule_name = resolve_schedule_name(schedule_types, schedule_name)
    if resolved_schedule_name is None:
        return []
    if resolved_schedule_name in visited:
        return []

    object_type = schedule_types.get(resolved_schedule_name)
    if object_type is None:
        return []

    schedule_objects = building_data.get(object_type, {})
    schedule_spec = schedule_objects.get(resolved_schedule_name)
    if not isinstance(schedule_spec, dict):
        return []

    visited.add(resolved_schedule_name)

    if object_type == "Schedule:Compact":
        values: List[float] = []
        for entry in schedule_spec.get("data", []):
            if not isinstance(entry, dict):
                continue
            field_value = entry.get("field")
            if isinstance(field_value, (int, float)) and not isinstance(field_value, bool):
                values.append(float(field_value))
        return values

    if object_type == "Schedule:Constant":
        values = []
        for key, value in schedule_spec.items():
            if key == "schedule_type_limits_name":
                continue
            if isinstance(value, (int, float)) and not isinstance(value, bool):
                values.append(float(value))
        return values

    if object_type == "Schedule:Day:Hourly":
        values = []
        for key, value in schedule_spec.items():
            if key.lower().startswith("hour_") and isinstance(value, (int, float)) and not isinstance(value, bool):
                values.append(float(value))
        return values

    if object_type in {"Schedule:Year", "Schedule:Week:Daily", "Schedule:Week:Compact"}:
        values = []
        for nested_schedule_name in _collect_schedule_references_from_value(schedule_spec, schedule_types):
            values.extend(
                extract_schedule_numeric_values(
                    building_data,
                    schedule_types,
                    nested_schedule_name,
                    visited,
                )
            )
        return values

    return []


def extract_schedule_value_range(
    building_data: Dict[str, Any],
    schedule_types: Dict[str, str],
    schedule_name: str,
) -> Tuple[Optional[float], Optional[float]]:
    """Extreu l'interval de valors de la planificació."""
    values = extract_schedule_numeric_values(building_data, schedule_types, schedule_name)
    if not values:
        return None, None
    return min(values), max(values)


def summarize_references(prefix: str, references: Sequence[str]) -> str:
    """Summarize references."""
    cleaned = [reference for reference in references if reference]
    if not cleaned:
        return prefix
    if len(cleaned) == 1:
        return f"{prefix}: {cleaned[0]}"
    if len(cleaned) == 2:
        return f"{prefix}: {cleaned[0]}, {cleaned[1]}"
    return f"{prefix}: {cleaned[0]}, {cleaned[1]} +{len(cleaned) - 2}"


def extract_scheduled_hvac_controller_candidates(
    building_data: Dict[str, Any],
    schedule_types: Dict[str, str],
    excluded_schedule_names: Sequence[str],
) -> List[Dict[str, Any]]:
    """Identifica els punts de consigna programats i els gestors de disponibilitat a partir del model d'edifici que
    encara no es gestionen com a punts de consigna explícits del termòstat.

    Paràmetres:
        building_data (Dict[str, Any]): les dades epJSON carregades.
        schedule_types (Dict[str, str]): una assignació de noms de planificació a tipus de planificació.
        excluded_schedule_names (Sequence[str]): Noms de les planificacions per ignorar
        (normalment ja està assignat als termòstats).

    Retorna:
        List[Dict[str, Any]]: una llista de diccionaris genèrics HVAC candidats.
    """
    excluded = {normalize_token(name) for name in excluded_schedule_names}
    detected: Dict[Tuple[str, str], Dict[str, Any]] = {}

    scheduled_setpoints = building_data.get("SetpointManager:Scheduled", {})
    for _, manager_spec in scheduled_setpoints.items():
        if not isinstance(manager_spec, dict):
            continue

        schedule_name = resolve_schedule_name(schedule_types, manager_spec.get("schedule_name", ""))
        if schedule_name is None:
            continue
        if normalize_token(schedule_name) in excluded:
            continue
        if normalize_token(str(manager_spec.get("control_variable", ""))) != "temperature":
            continue

        # Aquests setpoints no pengen directament del termòstat, però també poden ser
        # bons actuadors RL si EnergyPlus els controla amb una schedule editable.
        key = ("scheduled_temperature_setpoint", schedule_name)
        if key not in detected:
            detected[key] = {
                "category": "scheduled_temperature_setpoint",
                "schedule_name": schedule_name,
                "zones": set(),
                "references": set(),
                "source_kind": "setpoint_manager",
                "element_type": schedule_types[schedule_name],
            }

        setpoint_target = manager_spec.get("setpoint_node_or_nodelist_name")
        if isinstance(setpoint_target, str) and setpoint_target:
            detected[key]["references"].add(setpoint_target)

    scheduled_availability = building_data.get("AvailabilityManager:Scheduled", {})
    for manager_name, manager_spec in scheduled_availability.items():
        if not isinstance(manager_spec, dict):
            continue

        schedule_name = resolve_schedule_name(schedule_types, manager_spec.get("schedule_name", ""))
        if schedule_name is None:
            continue
        if normalize_token(schedule_name) in excluded:
            continue

        key = ("availability", schedule_name)
        if key not in detected:
            detected[key] = {
                "category": "availability",
                "schedule_name": schedule_name,
                "zones": set(),
                "references": set(),
                "source_kind": "availability_manager",
                "element_type": schedule_types[schedule_name],
            }
        if manager_name:
            detected[key]["references"].add(manager_name)

    ordered_candidates: List[Dict[str, Any]] = []
    for key in sorted(detected.keys(), key=lambda item: (item[0], item[1].lower())):
        candidate = detected[key]
        ordered_candidates.append(
            {
                "schedule_name": candidate["schedule_name"],
                "category": candidate["category"],
                "zones": tuple(sorted(candidate["zones"])),
                "references": tuple(sorted(candidate["references"])),
                "source_kind": candidate["source_kind"],
                "element_type": candidate["element_type"],
            }
        )

    return ordered_candidates


def extract_storage_controller_candidates(
    building_data: Dict[str, Any],
    schedule_types: Dict[str, str],
    excluded_schedule_names: Sequence[str],
) -> List[Dict[str, Any]]:
    """Identifica els horaris de càrrega/descàrrega de la bateria que es poden exposar com a accions."""

    excluded = {normalize_token(name) for name in excluded_schedule_names}
    detected: Dict[Tuple[str, str], Dict[str, Any]] = {}

    controller_fields = (
        (
            "storage_charge_power_fraction_schedule_name",
            "battery_charge",
        ),
        (
            "storage_discharge_power_fraction_schedule_name",
            "battery_discharge",
        ),
    )

    for distribution_name, distribution_spec in building_data.get("ElectricLoadCenter:Distribution", {}).items():
        if not isinstance(distribution_spec, dict):
            continue

        # La distribució pot apuntar a un objecte Storage; guardem tots dos noms perquè
        # després la UI pugui explicar d'on surt cada actuador de bateria.
        storage_name = distribution_spec.get("electrical_storage_object_name")
        reference = distribution_name
        if storage_name:
            reference = f"{distribution_name} -> {storage_name}"

        for field_name, category in controller_fields:
            # EnergyPlus separa càrrega i descàrrega en camps diferents; els exposem com
            # actuadors diferents per no barrejar dues decisions físiques oposades.
            schedule_name = resolve_schedule_name(schedule_types, distribution_spec.get(field_name, ""))
            if schedule_name is None:
                continue
            if normalize_token(schedule_name) in excluded:
                continue

            key = (category, schedule_name)
            if key not in detected:
                detected[key] = {
                    "category": category,
                    "schedule_name": schedule_name,
                    "zones": set(),
                    "references": set(),
                    "source_kind": "electric_storage",
                    "element_type": schedule_types[schedule_name],
                }
            detected[key]["references"].add(reference)

    ordered_candidates: List[Dict[str, Any]] = []
    for key in sorted(detected.keys(), key=lambda item: (item[0], item[1].lower())):
        candidate = detected[key]
        ordered_candidates.append(
            {
                "schedule_name": candidate["schedule_name"],
                "category": candidate["category"],
                "zones": tuple(sorted(candidate["zones"])),
                "references": tuple(sorted(candidate["references"])),
                "source_kind": candidate["source_kind"],
                "element_type": candidate["element_type"],
            }
        )

    return ordered_candidates


def parse_available_handlers(available_data: str) -> Tuple[List[Tuple[str, str, str]], List[Tuple[str, str]], List[str]]:
    """Analitzar els controladors disponibles."""
    actuators: List[Tuple[str, str, str]] = []
    variables: List[Tuple[str, str]] = []
    meters: List[str] = []

    reader = csv.reader(io.StringIO(available_data or ""))
    for row in reader:
        cleaned = [cell.strip().strip(";") for cell in row if cell is not None]
        if not cleaned:
            continue
        entry_type = cleaned[0]
        if entry_type == "Actuator" and len(cleaned) >= 4:
            actuators.append((cleaned[1], cleaned[2], cleaned[3]))
        elif entry_type == "OutputVariable" and len(cleaned) >= 3:
            variables.append((cleaned[1], cleaned[2]))
        elif entry_type == "OutputMeter" and len(cleaned) >= 2:
            meters.append(cleaned[1])

    return actuators, variables, meters


def probe_available_handlers(building_path: Path, weather_path: Path) -> Tuple[bool, str, Optional[str]]:
    """Inicia un bucle mínim Sinergym per extreure exactament quins controladors EnergyPlus estan disponibles.

    Paràmetres:
        building_path (Path): camí cap al model d'edifici IDF o epJSON.
        weather_path (Path): camí cap al fitxer meteorològic EPW.

    Retorna:
        Tuple[bool, str, Optional[str]]: una tupla que conté un booleà d'èxit, un text csv de controladors sense processar i una cadena d'error si escau.
    """
    env: Optional[EplusEnv] = None
    try:
        env = EplusEnv(
            building_file=str(building_path),
            weather_files=str(weather_path),
            action_space=gym.spaces.Box(low=0, high=0, shape=(0,), dtype=np.float32),
            time_variables=DEFAULT_TIME_VARIABLES,
            variables={},
            meters={},
            actuators={},
            context={},
            reward=ProbeReward,
            reward_kwargs={},
            env_name=f"probe-{sanitize_identifier(building_path.stem, 'env')}",
            building_config={
                "runperiod": (1, 1, 1991, 1, 1, 1991),
                "timesteps_per_hour": 1,
            },
        )
        env.reset()
        return True, env.available_handlers or "", None
    except Exception as error:
        return False, "", str(error)
    finally:
        if env is not None:
            try:
                env.close()
            except Exception:
                pass


def find_available_output_variable(
    available_variables: Sequence[Tuple[str, str]],
    variable_name: str,
    variable_key: str,
) -> Optional[Tuple[str, str]]:
    """Troba la variable de sortida disponible."""
    target_name = normalize_token(variable_name)
    target_key = normalize_token(variable_key)
    for actual_name, actual_key in available_variables:
        if normalize_token(actual_name) == target_name and normalize_token(actual_key) == target_key:
            return actual_name, actual_key
    return None


def find_available_meter(
    available_meters: Sequence[str],
    meter_name: str,
) -> Optional[str]:
    """Troba el comptador disponible."""
    target_name = normalize_token(meter_name)
    for actual_name in available_meters:
        if normalize_token(actual_name) == target_name:
            return actual_name
    return None


def default_bounds_for_category(
    category: str,
    actuator_key: str,
    current_low: Optional[float],
    current_high: Optional[float],
) -> Tuple[float, float]:
    """Límits per defecte per a la categoria."""
    if category == "heating_setpoint":
        return 15.0, 22.5
    if category == "cooling_setpoint":
        return 22.5, 30.0
    if category == "availability":
        return 0.0, 1.0
    if category in {"battery_charge", "battery_discharge"}:
        return 0.0, 1.0
    if category == "scheduled_temperature_setpoint":
        if current_low is not None and current_high is not None:
            if abs(current_high - current_low) < 1e-6:
                margin = max(1.0, abs(current_low) * 0.05)
                return float(current_low - margin), float(current_high + margin)
            margin = max(0.5, (current_high - current_low) * 0.1)
            return float(current_low - margin), float(current_high + margin)
        return 15.0, 80.0
    return 0.0, 1.0


def build_actuator_label(category: str) -> str:
    """Crea una etiqueta llegible per a una categoria d'actuador."""
    if category == "heating_setpoint":
        return "Setpoint de calefacció"
    if category == "cooling_setpoint":
        return "Setpoint de refrigeració"
    if category == "availability":
        return "Disponibilitat HVAC"
    if category == "scheduled_temperature_setpoint":
        return "Setpoint de sistema HVAC"
    if category == "battery_charge":
        return "Càrrega de bateria"
    if category == "battery_discharge":
        return "Descàrrega de bateria"
    return "Control HVAC"


def _build_actuator_option(
    candidate: Dict[str, Any],
    building_data: Dict[str, Any],
    schedule_types: Dict[str, str],
) -> ActuatorOption:
    """Construeix un ActuatorOption escrit a partir d'un diccionari candidat en brut.

    Extreu l'interval de valors de la planificació de les dades de l'edifici i calcula sensiblement
    límits d'acció per defecte basats en la categoria de l'actuador.

    Paràmetres:
        candidate (Dict[str, Any]): Dicte de candidat en brut que conté schedule_name, categoria,
            zones, referències, camps element_type i source_kind.
        building_data (Dict[str, Any]): dades completes de l'edifici epJSON per a l'extracció de l'interval de planificació.
        schedule_types (Dict[str, str]): Assignació de noms de planificació a tipus de blocs de planificació.

    Retorna:
        ActuatorOption: una opció d'actuador estructurada i immutable preparada per a la representació de UI i
            configuració de l'entorn.
    """
    schedule_name = candidate["schedule_name"]
    category = candidate["category"]
    zones = tuple(candidate["zones"])
    references = tuple(candidate["references"])
    current_low, current_high = extract_schedule_value_range(
        building_data,
        schedule_types,
        schedule_name,
    )
    default_low, default_high = default_bounds_for_category(
        category,
        schedule_name,
        current_low,
        current_high,
    )
    if category in {"heating_setpoint", "cooling_setpoint"}:
        reference = summarize_references("Zones", zones)
    elif category == "scheduled_temperature_setpoint":
        reference = summarize_references("Nodes", references)
    elif category in {"battery_charge", "battery_discharge"}:
        reference = summarize_references("Bateria", references)
    else:
        reference = summarize_references("Sistemes", references)

    return ActuatorOption(
        option_id=f'{candidate["element_type"]}|Schedule Value|{schedule_name}',
        actuator_key=schedule_name,
        element_type=candidate["element_type"],
        value_type="Schedule Value",
        label=build_actuator_label(category),
        reference=reference,
        category=category,
        zones=zones,
        current_low=current_low,
        current_high=current_high,
        default_low=default_low,
        default_high=default_high,
    )


def build_training_analysis(
    building_path: Path,
    weather_path: Path,
    probe_handlers: bool = False,
) -> BuildingTrainingAnalysis:
    """Realitza una anàlisi completa d'un model d'edifici per extreure controladors HVAC entrenables.

    Paràmetres:
        building_path (Path): camí cap al model d'edifici.
        weather_path (Path): Camí al fitxer meteorològic per resoldre les restriccions de l'entorn.
        probe_handlers (bool, optional): Si True, un entorn de sonda Sinergym s'executa exactament
            determinar els components de sortida. El valor predeterminat és False.

    Retorna:
        BuildingTrainingAnalysis: una classe de dades estructurada que conté resultats de detecció per
        termòstats, variables, comptadors i actuadors.
    """
    building_data = load_building_json(building_path)
    schedule_types = extract_schedule_type_index(building_data)
    thermostat_zones, thermostat_candidates = extract_thermostat_schedule_candidates(
        building_data,
        schedule_types,
    )
    thermostat_schedule_names = [
        candidate["schedule_name"]
        for candidate in thermostat_candidates
    ]
    hvac_schedule_candidates = extract_scheduled_hvac_controller_candidates(
        building_data,
        schedule_types,
        thermostat_schedule_names,
    )
    storage_schedule_candidates = extract_storage_controller_candidates(
        building_data,
        schedule_types,
        thermostat_schedule_names,
    )

    if probe_handlers:
        probe_success, available_raw, probe_error = probe_available_handlers(building_path, weather_path)
    else:
        probe_success, available_raw, probe_error = False, "", None
    _, available_output_variables, available_meters = parse_available_handlers(available_raw)

    thermostat_actuators: List[ActuatorOption] = []
    used_schedule_keys: set[str] = set()
    for candidate in thermostat_candidates:
        option = _build_actuator_option(candidate, building_data, schedule_types)
        thermostat_actuators.append(option)
        used_schedule_keys.add(normalize_token(option.actuator_key))

    hvac_actuators: List[ActuatorOption] = []
    for candidate in [*hvac_schedule_candidates, *storage_schedule_candidates]:
        normalized_key = normalize_token(candidate["schedule_name"])
        if normalized_key in used_schedule_keys:
            continue
        hvac_actuators.append(_build_actuator_option(candidate, building_data, schedule_types))
        used_schedule_keys.add(normalized_key)

    return BuildingTrainingAnalysis(
        thermostat_zones=tuple(thermostat_zones),
        thermostat_actuators=tuple(thermostat_actuators),
        hvac_actuators=tuple(hvac_actuators),
        available_output_variables=tuple(available_output_variables),
        available_meters=tuple(available_meters),
        probe_success=probe_success,
        probe_error=probe_error,
    )
\end{lstlisting}

\Needspace{8\baselineskip}

\subsection{\texorpdfstring{\texttt{backend/afegir\_entorn\_backend.py}}{backend/afegir\_entorn\_backend.py}}

\begin{lstlisting}[style=studioPython,caption={backend/afegir\_entorn\_backend.py},label={lst:backend-afegir-entorn-backend-py}]
"""Façana de compatibilitat per al backend afegir entorns.

La implementació es divideix en mòduls centrats, però aquest fitxer manté el fitxer
públic API utilitzat per la pàgina estable Streamlit.
"""

from __future__ import annotations

from backend.afegir_entorn_analysis import (
    ProbeReward,
    _build_actuator_option,
    _collect_schedule_references_from_value,
    _register_thermostat_candidate,
    build_actuator_label,
    build_training_analysis,
    default_bounds_for_category,
    extract_schedule_numeric_values,
    extract_schedule_type_index,
    extract_schedule_value_range,
    extract_scheduled_hvac_controller_candidates,
    extract_storage_controller_candidates,
    extract_thermostat_schedule_candidates,
    find_available_meter,
    find_available_output_variable,
    load_building_json,
    parse_available_handlers,
    probe_available_handlers,
    resolve_schedule_name,
    summarize_references,
)
from backend.afegir_entorn_config import (
    add_variable_spec,
    build_action_space_expression,
    build_environment_config,
    build_registered_env_id,
    build_training_observation_config,
    build_variable_name_for_actuator,
    build_variable_output_names,
    cleanup_failed_environment,
    register_environment_from_yaml,
    validate_registered_environment,
    write_yaml_config,
)
from backend.afegir_entorn_common import (
    DEFAULT_TIME_VARIABLES,
    DEFAULT_WEATHER_VARIABILITY,
    DETAILED_HVAC_METER_CANDIDATES,
    ENERGY_METER_CANDIDATES,
    ENERGY_VARIABLE_CANDIDATES,
    ENERGYPLUS_SUBPROCESS_TIMEOUT_SECONDS,
    ENV_VALIDATION_TIMEOUT_SECONDS,
    FALLBACK_UPDATER_DIR,
    STANDARD_WEATHER_VARIABLES,
    STANDARD_ZONE_VARIABLES,
    TARGET_EPLUS_VERSION,
    TRANSITION_DOWNLOAD_URL,
)
from backend.afegir_entorn_conversion import (
    convert_idf_to_epjson,
    detect_idf_version,
    download_transition_tools,
    get_transition_updater_dir,
    needs_idf_upgrade,
    upgrade_idf_version,
)
from backend.afegir_entorn_common import (
    ActuatorOption,
    BuildingTrainingAnalysis,
    EnvironmentValidationResult,
    UpgradeResult,
)
from backend.afegir_entorn_common import (
    normalize_token,
    sanitize_identifier,
    save_uploaded_bytes,
    unique_identifier,
)
from backend.common import BUILDINGS_DIR, CFG_DIR, PKG_DIR, WEATHERS_DIR
from backend.weather_profiles import (
    WeatherProfileSuggestion,
    build_weather_profile_suggestions,
    build_weather_temperature_preview,
    build_weather_temperature_preview_figure,
    load_epw_dry_bulb_series,
    summarize_weather_file,
)

CFG_DIR.mkdir(parents=True, exist_ok=True)

__all__ = [
    "ActuatorOption",
    "BUILDINGS_DIR",
    "BuildingTrainingAnalysis",
    "CFG_DIR",
    "DEFAULT_TIME_VARIABLES",
    "DEFAULT_WEATHER_VARIABILITY",
    "DETAILED_HVAC_METER_CANDIDATES",
    "ENERGYPLUS_SUBPROCESS_TIMEOUT_SECONDS",
    "ENERGY_METER_CANDIDATES",
    "ENERGY_VARIABLE_CANDIDATES",
    "ENV_VALIDATION_TIMEOUT_SECONDS",
    "EnvironmentValidationResult",
    "FALLBACK_UPDATER_DIR",
    "PKG_DIR",
    "ProbeReward",
    "STANDARD_WEATHER_VARIABLES",
    "STANDARD_ZONE_VARIABLES",
    "TARGET_EPLUS_VERSION",
    "TRANSITION_DOWNLOAD_URL",
    "UpgradeResult",
    "WEATHERS_DIR",
    "WeatherProfileSuggestion",
    "add_variable_spec",
    "build_action_space_expression",
    "build_actuator_label",
    "build_environment_config",
    "build_registered_env_id",
    "build_training_analysis",
    "build_training_observation_config",
    "build_variable_name_for_actuator",
    "build_variable_output_names",
    "build_weather_profile_suggestions",
    "build_weather_temperature_preview",
    "build_weather_temperature_preview_figure",
    "cleanup_failed_environment",
    "convert_idf_to_epjson",
    "default_bounds_for_category",
    "detect_idf_version",
    "download_transition_tools",
    "extract_schedule_numeric_values",
    "extract_schedule_type_index",
    "extract_schedule_value_range",
    "extract_scheduled_hvac_controller_candidates",
    "extract_storage_controller_candidates",
    "extract_thermostat_schedule_candidates",
    "find_available_meter",
    "find_available_output_variable",
    "get_transition_updater_dir",
    "load_building_json",
    "load_epw_dry_bulb_series",
    "needs_idf_upgrade",
    "normalize_token",
    "parse_available_handlers",
    "probe_available_handlers",
    "register_environment_from_yaml",
    "resolve_schedule_name",
    "sanitize_identifier",
    "save_uploaded_bytes",
    "summarize_references",
    "summarize_weather_file",
    "unique_identifier",
    "upgrade_idf_version",
    "validate_registered_environment",
    "write_yaml_config",
]
\end{lstlisting}

\Needspace{8\baselineskip}

\subsection{\texorpdfstring{\texttt{backend/afegir\_entorn\_common.py}}{backend/afegir\_entorn\_common.py}}

\begin{lstlisting}[style=studioPython,caption={backend/afegir\_entorn\_common.py},label={lst:backend-afegir-entorn-common-py}]
"""Constants, tipus i helpers compartits pel flux d'afegir entorns."""

from __future__ import annotations

from dataclasses import dataclass
from pathlib import Path
import re
from typing import Optional, Tuple


FALLBACK_UPDATER_DIR = Path("/workspaces/sinergym/IDFVersionUpdater")
TRANSITION_DOWNLOAD_URL = (
    "https://github.com/NREL/EnergyPlus/releases/download/v24.1.0/"
    "EnergyPlus-24.1.0-9d7789a3ac-Linux-Ubuntu22.04-x86_64.tar.gz"
)
TARGET_EPLUS_VERSION = (24, 1)
ENERGYPLUS_SUBPROCESS_TIMEOUT_SECONDS = 300
ENV_VALIDATION_TIMEOUT_SECONDS = 180

DEFAULT_TIME_VARIABLES = ["month", "day_of_month", "hour"]
DEFAULT_WEATHER_VARIABILITY = {
    "Dry Bulb Temperature": [1.0, 0.0, 24.0],
}
STANDARD_WEATHER_VARIABLES = (
    ("Site Outdoor Air DryBulb Temperature", "outdoor_temperature", "Environment"),
    ("Site Outdoor Air Relative Humidity", "outdoor_humidity", "Environment"),
    ("Site Wind Speed", "wind_speed", "Environment"),
    ("Site Wind Direction", "wind_direction", "Environment"),
    ("Site Diffuse Solar Radiation Rate per Area", "diffuse_solar_radiation", "Environment"),
    ("Site Direct Solar Radiation Rate per Area", "direct_solar_radiation", "Environment"),
)
STANDARD_ZONE_VARIABLES = (
    ("Zone Thermostat Heating Setpoint Temperature", "htg_setpoint"),
    ("Zone Thermostat Cooling Setpoint Temperature", "clg_setpoint"),
    ("Zone Air Temperature", "air_temperature"),
    ("Zone Air Relative Humidity", "air_humidity"),
    ("Zone People Occupant Count", "people_occupant"),
)
ENERGY_VARIABLE_CANDIDATES = (
    ("Facility Total HVAC Electricity Demand Rate", "HVAC_electricity_demand_rate", "Whole Building"),
    ("Facility Total Electricity Demand Rate", "facility_electricity_demand_rate", "Whole Building"),
)
ENERGY_METER_CANDIDATES = (
    ("Electricity:HVAC", "total_electricity_hvac"),
    ("Electricity:Facility", "total_electricity_facility"),
    ("Electricity:Building", "total_electricity_building"),
)
DETAILED_HVAC_METER_CANDIDATES = (
    ("NaturalGas:HVAC", "natural_gas_hvac"),
    ("Heating:Electricity", "heating_electricity"),
    ("Cooling:Electricity", "cooling_electricity"),
    ("Fans:Electricity", "fans_electricity"),
    ("Pumps:Electricity", "pumps_electricity"),
    ("HeatRejection:Electricity", "heat_rejection_electricity"),
    ("Humidifier:Electricity", "humidifier_electricity"),
    ("HeatRecovery:Electricity", "heat_recovery_electricity"),
)


@dataclass(frozen=True)
class UpgradeResult:
    """Resultat d'un intent d'actualitzar una IDF a la versio suportada."""
    upgraded: bool
    final_version: Optional[Tuple[int, int]]
    failed_transition_version: Optional[Tuple[int, int]] = None


@dataclass(frozen=True)
class EnvironmentValidationResult:
    """Espais i observacio inicial retornats en validar un entorn registrat."""
    observation_space: object
    action_space: object
    initial_observation: object


@dataclass(frozen=True)
class ActuatorOption:
    """Opcio d'actuador candidata detectada en un model EnergyPlus."""
    option_id: str
    actuator_key: str
    element_type: str
    value_type: str
    label: str
    reference: str
    category: str
    zones: Tuple[str, ...]
    current_low: Optional[float]
    current_high: Optional[float]
    default_low: float
    default_high: float


@dataclass(frozen=True)
class BuildingTrainingAnalysis:
    """Resultat agregat de l'analisi d'un edifici per entrenament RL."""
    thermostat_zones: Tuple[str, ...]
    thermostat_actuators: Tuple[ActuatorOption, ...]
    hvac_actuators: Tuple[ActuatorOption, ...]
    available_output_variables: Tuple[Tuple[str, str], ...]
    available_meters: Tuple[str, ...]
    probe_success: bool
    probe_error: Optional[str] = None


def save_uploaded_bytes(target_path: Path, content: bytes) -> Path:
    """Desa bytes pujats per Streamlit a disc i retorna la ruta."""
    target_path.parent.mkdir(parents=True, exist_ok=True)
    with open(target_path, "wb") as file_handle:
        file_handle.write(content)
    return target_path


def sanitize_identifier(raw_value: str, fallback: str) -> str:
    """Neteja una cadena per utilitzar-la com a identificador segur."""
    sanitized = re.sub(r"[^a-zA-Z0-9]+", "_", (raw_value or "").strip().lower()).strip("_")
    return sanitized or fallback


def normalize_token(value: str) -> str:
    """Normalitza una cadena per fer comparacions flexibles."""
    normalized = (value or "").lower()
    normalized = normalized.replace("drybulb", "dry bulb")
    normalized = re.sub(r"[^a-z0-9]+", " ", normalized)
    return re.sub(r"\s+", " ", normalized).strip()


def unique_identifier(base: str, used: set[str]) -> str:
    """Genera un identificador únic afegint un sufix incremental si cal."""
    candidate = base
    suffix = 2
    while candidate in used:
        candidate = f"{base}_{suffix}"
        suffix += 1
    used.add(candidate)
    return candidate
\end{lstlisting}

\Needspace{8\baselineskip}

\subsection{\texorpdfstring{\texttt{backend/afegir\_entorn\_config.py}}{backend/afegir\_entorn\_config.py}}

\begin{lstlisting}[style=studioPython,caption={backend/afegir\_entorn\_config.py},label={lst:backend-afegir-entorn-config-py}]
"""Crea i valideu les configuracions de Sinergym des de l'entrada del formulari Studio.

Aquest mòdul és l'últim pas del flux Afegeix un entorn. Rep el
anàlisi d'edificis, variables seleccionades, comptadors, actuadors i paràmetres de recompensa,
després escriu una configuració YAML que es pot registrar com a Gymnasium
entorn.
"""

from __future__ import annotations

import json
import re
import subprocess
import sys
from pathlib import Path
from typing import Any, Dict, List, Optional, Sequence, Tuple

import gymnasium as gym
import sinergym
import yaml
from sinergym.utils.common import convert_conf_to_env_parameters

from backend.afegir_entorn_analysis import find_available_meter, find_available_output_variable
from backend.afegir_entorn_common import (
    DEFAULT_TIME_VARIABLES,
    DETAILED_HVAC_METER_CANDIDATES,
    ENERGY_METER_CANDIDATES,
    ENERGY_VARIABLE_CANDIDATES,
    ENV_VALIDATION_TIMEOUT_SECONDS,
    STANDARD_WEATHER_VARIABLES,
    STANDARD_ZONE_VARIABLES,
)
from backend.afegir_entorn_common import (
    ActuatorOption,
    BuildingTrainingAnalysis,
    EnvironmentValidationResult,
)
from backend.afegir_entorn_common import sanitize_identifier, unique_identifier

def build_variable_output_names(base_name: str, keys: Sequence[str] | str) -> List[str]:
    """Retorna àlies de variable de sortida estables per a una o més claus EnergyPlus."""
    if isinstance(keys, str):
        return [base_name]
    return [f"{key.lower().replace(' ', '_')}_{base_name}" for key in keys]


def add_variable_spec(
    variables: Dict[str, Dict[str, Any]],
    available_output_variables: Sequence[Tuple[str, str]],
    canonical_variable_name: str,
    alias_name: str,
    keys: Sequence[str] | str,
) -> List[str]:
    """Afegeix una especificació variable."""
    requested_keys = [keys] if isinstance(keys, str) else list(keys)
    matched_keys: List[str] = []
    actual_variable_name = canonical_variable_name

    for requested_key in requested_keys:
        matched = find_available_output_variable(
            available_output_variables,
            canonical_variable_name,
            requested_key,
        )
        if matched is None:
            continue
        actual_variable_name, actual_key = matched
        matched_keys.append(actual_key)

    if not matched_keys:
        if isinstance(keys, str):
            variables[canonical_variable_name] = {
                "variable_names": alias_name,
                "keys": keys,
            }
            return [alias_name]

        variables[canonical_variable_name] = {
            "variable_names": alias_name,
            "keys": list(keys),
        }
        return build_variable_output_names(alias_name, list(keys))

    output_keys: Sequence[str] | str
    if len(matched_keys) == 1:
        output_keys = matched_keys[0]
    else:
        output_keys = matched_keys

    variables[actual_variable_name] = {
        "variable_names": alias_name,
        "keys": output_keys,
    }
    return build_variable_output_names(alias_name, output_keys)


def build_training_observation_config(
    analysis: BuildingTrainingAnalysis,
) -> Tuple[Dict[str, Dict[str, Any]], Dict[str, str], List[str], List[str]]:
    """Munta la configuració d'observacions d'entrenament."""
    variables: Dict[str, Dict[str, Any]] = {}
    meters: Dict[str, str] = {}

    # Primer afegim variables meteorològiques i de zona "normals"; són les que després
    # espera la UI encara que el fitxer EnergyPlus tingui noms llargs o lleugerament diferents.
    for variable_name, alias_name, variable_key in STANDARD_WEATHER_VARIABLES:
        add_variable_spec(
            variables,
            analysis.available_output_variables,
            variable_name,
            alias_name,
            variable_key,
        )

    temperature_variable_names: List[str] = []
    for variable_name, alias_name in STANDARD_ZONE_VARIABLES:
        generated_names = add_variable_spec(
            variables,
            analysis.available_output_variables,
            variable_name,
            alias_name,
            analysis.thermostat_zones,
        )
        if alias_name == "air_temperature":
            temperature_variable_names = generated_names

    energy_variable_names: List[str] = []
    for variable_name, alias_name, variable_key in ENERGY_VARIABLE_CANDIDATES:
        matched = find_available_output_variable(
            analysis.available_output_variables,
            variable_name,
            variable_key,
        )
        if matched is None and analysis.probe_success:
            continue

        add_variable_spec(
            variables,
            analysis.available_output_variables,
            variable_name,
            alias_name,
            variable_key,
        )
        energy_variable_names = [alias_name]
        break

    if not energy_variable_names:
        # Alguns models no publiquen energia HVAC com a output variable, però sí com a meter.
        # En aquest cas fem servir el meter equivalent i mantenim el mateix alias intern.
        for meter_name, alias_name in ENERGY_METER_CANDIDATES:
            matched_meter = find_available_meter(analysis.available_meters, meter_name)
            if matched_meter is None and analysis.probe_success:
                continue
            meters[matched_meter or meter_name] = alias_name
            energy_variable_names = [alias_name]
            break

    for meter_name, alias_name in DETAILED_HVAC_METER_CANDIDATES:
        # Els meters detallats no són obligatoris per entrenar, però si hi són donen
        # gràfics molt més útils al dashboard.
        matched_meter = find_available_meter(analysis.available_meters, meter_name)
        if matched_meter is None:
            continue
        meters.setdefault(matched_meter, alias_name)

    if not temperature_variable_names:
        raise ValueError(
            "No s'han pogut construir variables de temperatura de zona a partir dels termòstats detectats."
        )
    if not energy_variable_names:
        raise ValueError(
            "No s'ha detectat cap variable o meter energètic vàlid per construir la reward."
        )

    return variables, meters, temperature_variable_names, energy_variable_names


def build_variable_name_for_actuator(
    option: ActuatorOption,
    selected_options: Sequence[ActuatorOption],
    used_names: set[str],
) -> str:
    """Genera el nom de variable associat a un actuador."""
    heating_count = sum(item.category == "heating_setpoint" for item in selected_options)
    cooling_count = sum(item.category == "cooling_setpoint" for item in selected_options)

    # Si només hi ha una consigna de fred/calor mantenim noms simples. Amb més zones,
    # el nom ha de portar zona o schedule per no crear variables duplicades.
    if option.category == "heating_setpoint":
        if heating_count == 1:
            base_name = "Heating_Setpoint_RL"
        elif len(option.zones) == 1:
            base_name = f"heating_setpoint_{sanitize_identifier(option.zones[0], 'zone')}"
        else:
            base_name = f"heating_setpoint_{sanitize_identifier(option.actuator_key, 'schedule')}"
    elif option.category == "cooling_setpoint":
        if cooling_count == 1:
            base_name = "Cooling_Setpoint_RL"
        elif len(option.zones) == 1:
            base_name = f"cooling_setpoint_{sanitize_identifier(option.zones[0], 'zone')}"
        else:
            base_name = f"cooling_setpoint_{sanitize_identifier(option.actuator_key, 'schedule')}"
    elif option.category == "scheduled_temperature_setpoint":
        base_name = f"hvac_setpoint_{sanitize_identifier(option.actuator_key, 'schedule')}"
    elif option.category == "availability":
        base_name = f"hvac_availability_{sanitize_identifier(option.actuator_key, 'schedule')}"
    elif option.category == "battery_charge":
        base_name = "Battery_Charge_Rate_RL"
    elif option.category == "battery_discharge":
        base_name = "Battery_Discharge_Rate_RL"
    else:
        base_name = f"hvac_control_{sanitize_identifier(option.actuator_key, 'actuator')}"

    return unique_identifier(base_name, used_names)


def build_action_space_expression(bounds: Sequence[Tuple[float, float]]) -> str:
    """Construeix una expressió d'espai d'acció."""
    lows = ", ".join(f"{low:.4f}" for low, _ in bounds)
    highs = ", ".join(f"{high:.4f}" for _, high in bounds)
    return (
        "gym.spaces.Box("
        f"low=np.array([{lows}], dtype=np.float32), "
        f"high=np.array([{highs}], dtype=np.float32), "
        f"shape=({len(bounds)},), dtype=np.float32)"
    )


def build_environment_config(
    id_base: str,
    building_file_name: str,
    weather_profiles: Sequence[Dict[str, str]],
    analysis: BuildingTrainingAnalysis,
    selected_actuators: Sequence[ActuatorOption],
    actuator_bounds: Dict[str, Tuple[float, float]],
    weather_variability: Optional[Dict[str, Sequence[float]]] = None,
) -> Dict[str, Any]:
    """Construeix el diccionari de configuració de l'entorn arrel Sinergym YAML.

    Paràmetres:
        id_base (str): nom base per al registre de l'entorn.
        building_file_name (str): el nom del fitxer d'estructura.
        weather_profiles (Sequence[Dict[str, str]]): Llista de mapes meteorològics (fitxa i claus).
        analysis (BuildingTrainingAnalysis): propietats detectades.
        selected_actuators (Sequence[ActuatorOption]): actuadors seleccionats del UI.
        actuator_bounds (Dict[str, Tuple[float, float]]): els límits continus de cada actuador.
        weather_variability (Optional[Dict[str, Sequence[float]]]): Diccionari de variabilitat opcional.

    Retorna:
        Dict[str, Any]: La configuració totalment estructurada està preparada per escriure com a YAML.
    """
    variables, meters, temperature_variable_names, energy_variable_names = (
        build_training_observation_config(analysis)
    )

    # Ordenem actuadors perquè el Box tingui un ordre estable: calor, fred, bateria i resta.
    # Sense això, dos registres del mateix edifici podrien generar espais d'acció diferents.
    ordered_actuators = sorted(
        selected_actuators,
        key=lambda option: (
            {
                "heating_setpoint": 0,
                "cooling_setpoint": 1,
                "battery_charge": 2,
                "battery_discharge": 3,
            }.get(option.category, 4),
            option.label.lower(),
        ),
    )

    used_variable_names: set[str] = set()
    actuators: Dict[str, Dict[str, str]] = {}
    bounds: List[Tuple[float, float]] = []
    for option in ordered_actuators:
        # actuator_key apunta a l'objecte EnergyPlus real; variable_name és el nom curt que
        # veurà Sinergym i després el model RL.
        low, high = actuator_bounds[option.option_id]
        variable_name = build_variable_name_for_actuator(option, ordered_actuators, used_variable_names)
        actuators[option.actuator_key] = {
            "variable_name": variable_name,
            "element_type": option.element_type,
            "value_type": option.value_type,
        }
        bounds.append((low, high))

    env_cfg: Dict[str, Any] = {
        "id_base": id_base,
        "building_file": building_file_name,
        "weather_specification": {
            "weather_files": [profile["weather_file"] for profile in weather_profiles],
            "keys": [profile["key"] for profile in weather_profiles],
        },
        "building_config": None,
        "max_ep_store": 3,
        "time_variables": list(DEFAULT_TIME_VARIABLES),
        "variables": variables,
        "meters": meters,
        "actuators": actuators,
        "context": {},
        "action_space": build_action_space_expression(bounds),
        "reward": "sinergym.utils.rewards:LinearReward",
        "reward_kwargs": {
            "temperature_variables": temperature_variable_names,
            "energy_variables": energy_variable_names,
            "range_comfort_winter": [20.0, 23.5],
            "range_comfort_summer": [23.0, 26.0],
            "summer_start": [6, 1],
            "summer_final": [9, 30],
            "energy_weight": 0.5,
            "lambda_energy": 1e-4,
            "lambda_temperature": 1.0,
        },
    }

    if weather_variability:
        # Només escrivim weather_variability quan l'usuari l'ha configurat; deixar-ho buit
        # pot activar comportaments diferents en versions de Sinergym.
        env_cfg["weather_variability"] = weather_variability

    return env_cfg


def write_yaml_config(cfg_path: Path, env_cfg: Dict[str, Any]) -> Path:
    """Escriu el diccionari de configuració de l'entorn generat en un fitxer YAML.

    Paràmetres:
        cfg_path (Path): Camí al fitxer de sortida YAML.
        env_cfg (Dict[str, Any]): dades de configuració del diccionari.

    Retorna:
        Path: la ruta del fitxer de sortida confirmada.
    """
    with open(cfg_path, "w", encoding="utf-8") as file_handle:
        yaml.dump(env_cfg, file_handle, sort_keys=False, allow_unicode=True)
    return cfg_path


def register_environment_from_yaml(cfg_path: Path) -> None:
    """Registra o torna a registrar els entorns des d'una configuració Sinergym YAML mitjançant Gym.

    Paràmetres:
        cfg_path (Path): camí cap a la configuració vàlida Sinergym YAML.
    """
    with open(cfg_path, "r", encoding="utf-8") as yaml_conf:
        conf = yaml.safe_load(yaml_conf)

    configurations = convert_conf_to_env_parameters(conf)
    registry = gym.envs.registration.registry
    for env_id in configurations.keys():
        registry.pop(env_id, None)
        registry.pop(env_id.replace("continuous", "discrete"), None)

    sinergym.register_envs_from_yaml(str(cfg_path))


def build_registered_env_id(id_base: str, weather_key: str, stochastic: bool = False) -> str:
    """Crea l'ID de l'entorn del Gym canònic basat en les convencions de denominació Sinergym.

    Paràmetres:
        id_base (str): Nom base donat a l'element d'edifici.
        weather_key (str): descriptor breu de la cadena meteorològica.
        stochastic (bool, optional): si la variant pretén incloure variabilitat.

    Retorna:
        str: identificador d'entorn OpenAI Gym vàlid.
    """
    suffix = "-continuous-stochastic-v1" if stochastic else "-continuous-v1"
    return f"Eplus-{id_base}-{weather_key}{suffix}"


def validate_registered_environment(env_id: str, cfg_path: Optional[Path] = None) -> EnvironmentValidationResult:
    """Executa un subprocés aïllat per activar i validar que l'entorn es registra i passa correctament.

    Paràmetres:
        env_id (str): cadena d'ID d'entorn Gym.
        cfg_path (Optional[Path], optional): camí cap a yaml si és necessari per al registre diferit.

    Retorna:
        EnvironmentValidationResult: un dict de verificació estructural que indica la matriu d'observació i els espais.
    """
    validation_script = """
import json
import sys

import gymnasium as gym
import sinergym

env_id = sys.argv[1]
cfg_path = sys.argv[2] if len(sys.argv) > 2 else None
if cfg_path:
    sinergym.register_envs_from_yaml(cfg_path)

env = gym.make(env_id)
try:
    observation = env.reset()[0]
    if hasattr(observation, "tolist"):
        observation = observation.tolist()
    payload = {
        "observation_space": repr(env.observation_space),
        "action_space": repr(env.action_space),
        "initial_observation": observation,
    }
    print(f"---JSON_START---\\n{json.dumps(payload)}\\n---JSON_END---")
finally:
    env.close()
"""

    try:
        result = subprocess.run(
            [sys.executable, "-", env_id, *( [str(cfg_path)] if cfg_path else [] )],
            input=validation_script,
            check=True,
            capture_output=True,
            text=True,
            timeout=ENV_VALIDATION_TIMEOUT_SECONDS,
        )
        match = re.search(r"---JSON_START---\n(.*?)\n---JSON_END---", result.stdout, re.DOTALL)
        if not match:
            raise RuntimeError(f"No s'ha trobat l'estructura de validació a l'output. Output obtingut:\n{result.stdout}\nErrors:\n{result.stderr}")

        payload = json.loads(match.group(1))
        return EnvironmentValidationResult(
            observation_space=payload["observation_space"],
            action_space=payload["action_space"],
            initial_observation=payload["initial_observation"],
        )
    except subprocess.TimeoutExpired as error:
        raise TimeoutError(
            "La validació de l'entorn ha superat el temps límit. "
            "Revisa la configuració del model o prova un edifici més lleuger."
        ) from error
    except subprocess.CalledProcessError as error:
        stderr = (error.stderr or "").strip()
        stdout = (error.stdout or "").strip()
        raise RuntimeError(stderr or stdout or "La validació ha fallat en arrencar l'entorn.") from error


def cleanup_failed_environment(cfg_path: Path, tmp_path: Path) -> List[Tuple[str, str]]:
    """Intenta eliminar les estructures de configuració no enllaçades o incorrectes com a resultat d'un registre fallit.

    Paràmetres:
        cfg_path (Path): Configuració interrompuda de YAML.
        tmp_path (Path): fitxer epJSON intermedi generat.

    Retorna:
        List[Tuple[str, str]]: Tuples de missatgeria de nivell de diagnòstic que mapegen quins recursos es van revertir.
    """
    messages: List[Tuple[str, str]] = []
    try:
        if cfg_path.exists():
            cfg_path.unlink()
            messages.append(("warning", "Fitxer de registre YAML revertit."))
        if tmp_path.suffix.lower() == ".idf":
            epjson_path = tmp_path.with_suffix(".epJSON")
            if epjson_path.exists():
                epjson_path.unlink()
                messages.append(("warning", "Fitxer estructural epJSON revertit."))
    except Exception as cleanup_error:
        messages.append(("error", f"Error eliminant fitxers: {cleanup_error}"))
    return messages
\end{lstlisting}

\Needspace{8\baselineskip}

\subsection{\texorpdfstring{\texttt{backend/afegir\_entorn\_conversion.py}}{backend/afegir\_entorn\_conversion.py}}

\begin{lstlisting}[style=studioPython,caption={backend/afegir\_entorn\_conversion.py},label={lst:backend-afegir-entorn-conversion-py}]
"""Conversió i actualització d'IDF/epJSON dins del flux per afegir entorns."""

from __future__ import annotations

import os
import re
import shutil
import ssl
import subprocess
import tarfile
import urllib.request
from pathlib import Path
from typing import Optional, Tuple

from backend.afegir_entorn_common import (
    ENERGYPLUS_SUBPROCESS_TIMEOUT_SECONDS,
    FALLBACK_UPDATER_DIR,
    TARGET_EPLUS_VERSION,
    TRANSITION_DOWNLOAD_URL,
)
from backend.afegir_entorn_common import UpgradeResult

def detect_idf_version(idf_path: Path) -> Optional[Tuple[int, int]]:
    """Llegeix un fitxer IDF per extreure la seva versió de destinació EnergyPlus.

    Paràmetres:
        idf_path (Path): camí cap al fitxer IDF.

    Retorna:
        Optional[Tuple[int, int]]: una tupla de nombres enters majors i menors si es troba, sinó None.
    """
    with open(idf_path, "r", encoding="utf-8", errors="ignore") as file_handle:
        content = file_handle.read()

    match = re.search(r"Version,\s*([\d\.]+)\s*;", content, re.IGNORECASE)
    if not match:
        return None

    parts = match.group(1).strip().split(".")
    if len(parts) < 2:
        return None

    return int(parts[0]), int(parts[1])


def needs_idf_upgrade(version: Tuple[int, int]) -> bool:
    """Avalua si una versió IDF requereix l'actualització a la versió del simulador de destinació.

    Paràmetres:
        version (Tuple[int, int]): versió major i menor de IDF.

    Retorna:
        bool: True si la versió és estrictament més antiga que la `TARGET_EPLUS_VERSION`.
    """
    major, minor = version
    target_major, target_minor = TARGET_EPLUS_VERSION
    return not (major > target_major or (major == target_major and minor >= target_minor))


def get_transition_updater_dir() -> Optional[Path]:
    """Recupera el directori local que conté EnergyPlus IDFVersionUpdater binaris.

    Retorna:
        Optional[Path]: camí cap a les eines d'actualització, o None si no es troba localment.
    """
    ep_base_dir = Path(os.environ.get("EPLUS_PATH", "/usr/local/EnergyPlus-24-1-0"))
    updater_dir = ep_base_dir / "PreProcess" / "IDFVersionUpdater"
    if updater_dir.exists():
        return updater_dir
    if FALLBACK_UPDATER_DIR.exists():
        return FALLBACK_UPDATER_DIR
    return None


def download_transition_tools() -> Path:
    """Baixa i extreu EnergyPlus Eines de transició al directori alternatiu.

    En lloc de descarregar tota la suite EnergyPlus, només extreu l'actualitzador PreProcess.

    Retorna:
        Path: el directori local que conté les eines de transició descarregades.
    """
    ssl._create_default_https_context = ssl._create_unverified_context
    tar_path, _ = urllib.request.urlretrieve(TRANSITION_DOWNLOAD_URL, "/tmp/eplus.tar.gz")
    FALLBACK_UPDATER_DIR.mkdir(parents=True, exist_ok=True)
    with tarfile.open(tar_path, "r:gz") as tar:
        members = []
        for member in tar.getmembers():
            if "PreProcess/IDFVersionUpdater/" in member.name:
                member.name = member.name.split("PreProcess/IDFVersionUpdater/")[1]
                if member.name:
                    members.append(member)
        tar.extractall(path=FALLBACK_UPDATER_DIR, members=members)
    Path("/tmp/eplus.tar.gz").unlink(missing_ok=True)
    return FALLBACK_UPDATER_DIR


def upgrade_idf_version(idf_path: Path, updater_dir: Path) -> UpgradeResult:
    """Actualitza seqüencialment un fitxer IDF mitjançant cadenes de transició EnergyPlus.

    Paràmetres:
        idf_path (Path): camí cap al fitxer IDF de destinació.
        updater_dir (Path): Directori que conté les eines de transició.

    Retorna:
        UpgradeResult: Estructura resumida de l'operació d'actualització.
    """
    initial_version = detect_idf_version(idf_path)
    if initial_version is None or not needs_idf_upgrade(initial_version):
        return UpgradeResult(upgraded=False, final_version=initial_version)

    current_major, current_minor = initial_version
    working_dir = Path("/tmp/eplus_upgrade")
    working_dir.mkdir(parents=True, exist_ok=True)
    working_file = working_dir / idf_path.name
    shutil.copy(idf_path, working_file)

    loops = 0
    failed_transition_version = None
    while loops < 20:
        prefix = f"Transition-V{current_major}-{current_minor}-0-to-V"
        tool = next(
            (
                file_path
                for file_path in updater_dir.glob("Transition-*")
                if file_path.name.startswith(prefix)
            ),
            None,
        )

        if not tool:
            break

        try:
            subprocess.run(
                [str(tool), str(working_file)],
                check=True,
                cwd=str(updater_dir),
                capture_output=True,
                timeout=ENERGYPLUS_SUBPROCESS_TIMEOUT_SECONDS,
            )
        except subprocess.CalledProcessError:
            failed_transition_version = (current_major, current_minor)
            break

        next_version = detect_idf_version(working_file)
        if next_version is None:
            break

        current_major, current_minor = next_version
        loops += 1

    shutil.copy(working_file, idf_path)
    final_version = detect_idf_version(idf_path)

    upgraded = False
    if initial_version is not None and final_version is not None:
        upgraded = final_version[0] > initial_version[0] or final_version[1] > initial_version[1]

    return UpgradeResult(
        upgraded=upgraded,
        final_version=final_version,
        failed_transition_version=failed_transition_version,
    )


def convert_idf_to_epjson(idf_path: Path) -> Path:
    """Converteix un fitxer EnergyPlus IDF al format epJSON mitjançant les eines CLI.

    Paràmetres:
        idf_path (Path): camí cap al fitxer estructural IDF.

    Retorna:
        Path: el camí del fitxer epJSON convertit resultant.
    """
    epjson_path = idf_path.with_suffix(".epJSON")
    subprocess.run(
        [
            "energyplus",
            "--convert-only",
            "--output-directory",
            str(idf_path.parent),
            str(idf_path),
        ],
        check=True,
        timeout=ENERGYPLUS_SUBPROCESS_TIMEOUT_SECONDS,
    )
    return epjson_path
\end{lstlisting}

\Needspace{8\baselineskip}

\subsection{\texorpdfstring{\texttt{backend/avaluar\_agent\_backend.py}}{backend/avaluar\_agent\_backend.py}}

\begin{lstlisting}[style=studioPython,caption={backend/avaluar\_agent\_backend.py},label={lst:backend-avaluar-agent-backend-py}]
"""Temps d'execució d'avaluació asíncrona per a agents Studio entrenats.

L'avaluació es pot executar per a molts passos de simulació, de manera que aquest backend aïlla el
bucle d'execució, cua de progrés i estat de cancel·lació des de la pàgina Streamlit. Això
també centralitza la càrrega de models, la creació d'entorns i la validació de l'espai d'acció
de manera que el control i l'avaluació en directe segueixen el mateix contracte.
"""

from __future__ import annotations

import queue
import threading
import time
import traceback
from dataclasses import dataclass
from functools import partial
from pathlib import Path
from typing import Any, Callable, Dict, List, Optional, Tuple

import gymnasium as gym
import numpy as np
import streamlit as st
from stable_baselines3.common.vec_env import VecEnvWrapper

import sinergym  # noqa: F401
from sinergym.utils.wrappers import CSVLogger, LoggerWrapper, NormalizeAction

from backend.model_metadata import (
    build_gym_kwargs_from_metadata,
    load_model_metadata,
    validate_action_spaces,
)
from backend.common import ONE_YEAR_STEPS
from backend.sb3_utils import (
    candidate_vecnorm,
    build_monitored_vec_env,
    env_id_from_meta_or_name,
    format_action_for_env,
    load_sb3_model_bytes,
    load_vecnormalize,
    scan_model_zips,
    should_apply_normalize_action,
)


EVALUATION_RUNTIME_STATE_KEY = "evaluation_runtime_state"


@dataclass(frozen=True)
class EvaluationResult:
    """Resum de resultats d'una execució d'un treballador d'avaluació."""
    elapsed_seconds: float
    cancelled: bool
    warnings: Tuple[str, ...] = ()


def push_evaluation_progress(
    event_queue: "queue.Queue[Dict[str, Any]]",
    job_config: Dict[str, Any],
    step_number: int,
    total_steps: int,
) -> None:
    """Envieu un esdeveniment de progrés d'avaluació a la cua d'execució."""
    push_evaluation_event(
        event_queue,
        {
            "type": "progress",
            "step_number": int(step_number or 0),
            "total_steps": int(total_steps or job_config["steps_target"]),
            "status": f"Passos: {int(step_number or 0)}/{int(total_steps or job_config['steps_target'])}",
        },
    )


def build_evaluation_env(
    env_id: str,
    gym_kwargs: Dict[str, Any],
    model_action_space: Any,
    wrapper_configs: Optional[List[Dict[str, Any]]],
) -> Any:
    """Crea un entorn embolicat per a una execució d'avaluació d'agent."""
    env = make_env(env_id, seed=0, gym_kwargs=gym_kwargs)
    env = apply_evaluation_action_contract(env, model_action_space, wrapper_configs)
    env = LoggerWrapper(env)
    env = CSVLogger(env)
    return env


def empty_evaluation_runtime() -> Dict[str, Any]:
    """Temps d'execució d'avaluació buit."""
    return {
        "running": False,
        "completed": False,
        "cancel_requested": False,
        "needs_full_rerun": False,
        "thread": None,
        "queue": None,
        "stop_event": None,
        "result": None,
        "error": None,
        "traceback": None,
        "model_path": None,
        "env_id": None,
        "steps_target": 1,
        "latest_step": 0,
        "status": "",
        "started_at": None,
    }


def ensure_evaluation_runtime() -> Dict[str, Any]:
    """Assegura el temps d'execució de l'avaluació."""
    if EVALUATION_RUNTIME_STATE_KEY not in st.session_state:
        st.session_state[EVALUATION_RUNTIME_STATE_KEY] = empty_evaluation_runtime()
    return st.session_state[EVALUATION_RUNTIME_STATE_KEY]


def reset_evaluation_runtime() -> Dict[str, Any]:
    """Restableix el temps d'execució de l'avaluació."""
    runtime = ensure_evaluation_runtime()
    previous_thread = runtime.get("thread")
    if previous_thread is not None and getattr(previous_thread, "is_alive", lambda: False)():
        return runtime

    st.session_state[EVALUATION_RUNTIME_STATE_KEY] = empty_evaluation_runtime()
    return st.session_state[EVALUATION_RUNTIME_STATE_KEY]


def push_evaluation_event(event_queue: "queue.Queue[Dict[str, Any]]", payload: Dict[str, Any]) -> None:
    """Esdeveniment d'avaluació push."""
    event_queue.put(dict(payload))


def evaluation_worker(
    job_config: Dict[str, Any],
    event_queue: "queue.Queue[Dict[str, Any]]",
    stop_event: threading.Event,
) -> None:
    """Executa l'avaluació en segon pla."""
    try:
        result = run_agent_evaluation(
            model_path=job_config["model_path"],
            env_id=job_config["env_id"],
            use_vecnorm=bool(job_config["use_vecnorm"]),
            vecnorm_path=job_config["vecnorm_path"],
            steps_target=int(job_config["steps_target"]),
            deterministic=bool(job_config["deterministic"]),
            should_stop=stop_event.is_set,
            on_progress=partial(push_evaluation_progress, event_queue, job_config),
            override_wrapper_configs=job_config.get("wrapper_configs") or None,
        )
        push_evaluation_event(event_queue, {"type": "result", "result": result})
    except Exception as exc:
        push_evaluation_event(
            event_queue,
            {
                "type": "error",
                "error": str(exc),
                "traceback": traceback.format_exc(),
            },
        )
    finally:
        push_evaluation_event(event_queue, {"type": "finished"})


def drain_evaluation_runtime(runtime: Dict[str, Any]) -> None:
    """Temps d'execució de l'avaluació del buidatge."""
    event_queue = runtime.get("queue")
    if event_queue is None:
        return

    # La UI consulta aquest runtime sovint; buidem la cua sencera en cada passada perquè
    # no quedin esdeveniments antics fent saltar el progrés enrere.
    while True:
        try:
            event = event_queue.get_nowait()
        except queue.Empty:
            break

        event_type = str(event.get("type") or "")
        if event_type == "progress":
            runtime["latest_step"] = max(
                int(runtime.get("latest_step") or 0),
                int(event.get("step_number") or 0),
            )
            runtime["steps_target"] = max(
                int(runtime.get("steps_target") or 1),
                int(event.get("total_steps") or 1),
            )
            runtime["status"] = str(event.get("status") or runtime.get("status") or "")
        elif event_type == "result":
            runtime["result"] = event.get("result")
            runtime["completed"] = True
            runtime["running"] = False
            runtime["needs_full_rerun"] = True
        elif event_type == "error":
            runtime["error"] = str(event.get("error") or "Error desconegut durant l'avaluació.")
            runtime["traceback"] = event.get("traceback")
            runtime["completed"] = True
            runtime["running"] = False
            runtime["needs_full_rerun"] = True
        elif event_type == "finished":
            runtime["completed"] = True
            runtime["running"] = False
            runtime["thread"] = None
            runtime["needs_full_rerun"] = True

    thread = runtime.get("thread")
    if thread is not None and not thread.is_alive() and runtime.get("running"):
        runtime["running"] = False
        runtime["completed"] = True
        runtime["thread"] = None
        runtime["needs_full_rerun"] = True


def sync_evaluation_runtime() -> Dict[str, Any]:
    """Sincronitza el temps d'execució de l'avaluació."""
    runtime = ensure_evaluation_runtime()
    drain_evaluation_runtime(runtime)
    return runtime


def consume_evaluation_runtime_rerun_flag(runtime: Optional[Dict[str, Any]] = None) -> bool:
    """Indicador de reexecució del clima d'execució de l'avaluació."""
    active_runtime = runtime if runtime is not None else ensure_evaluation_runtime()
    if active_runtime.get("needs_full_rerun") and not active_runtime.get("running"):
        active_runtime["needs_full_rerun"] = False
        return True
    return False


def request_evaluation_stop(runtime: Optional[Dict[str, Any]] = None) -> bool:
    """Sol·licitar l'aturada de l'avaluació."""
    active_runtime = runtime if runtime is not None else ensure_evaluation_runtime()
    if not active_runtime.get("running") or active_runtime.get("stop_event") is None:
        return False

    active_runtime["cancel_requested"] = True
    active_runtime["status"] = "Aturant avaluació..."
    active_runtime["stop_event"].set()
    return True


def start_evaluation_run(request: Dict[str, Any]) -> Dict[str, Any]:
    """Iniciar l'execució d'avaluació."""
    runtime = reset_evaluation_runtime()
    response: Dict[str, Any] = {
        "started": False,
        "errors": [],
        "runtime": runtime,
    }

    if runtime.get("running"):
        response["errors"].append("Ja hi ha una avaluació en curs.")
        return response

    model_path = Path(request.get("model_path") or "")
    env_id = str(request.get("env_id") or "").strip()
    vecnorm_path = request.get("vecnorm_path")

    if not model_path.exists():
        response["errors"].append("No s'ha trobat el model seleccionat.")
        return response
    if not env_id:
        response["errors"].append("No s'ha pogut determinar un env_id vàlid.")
        return response

    event_queue: "queue.Queue[Dict[str, Any]]" = queue.Queue()
    stop_event = threading.Event()
    job_config = {
        "model_path": model_path,
        "env_id": env_id,
        "use_vecnorm": bool(request.get("use_vecnorm")),
        "vecnorm_path": Path(vecnorm_path) if vecnorm_path else None,
        "steps_target": int(request.get("steps_target") or ONE_YEAR_STEPS),
        "deterministic": bool(request.get("deterministic", True)),
        "wrapper_configs": list(request.get("wrapper_configs") or []),
    }
    worker_thread = threading.Thread(
        target=evaluation_worker,
        args=(job_config, event_queue, stop_event),
        daemon=True,
    )

    runtime.update(
        {
            "running": True,
            "completed": False,
            "cancel_requested": False,
            "needs_full_rerun": False,
            "thread": worker_thread,
            "queue": event_queue,
            "stop_event": stop_event,
            "result": None,
            "error": None,
            "traceback": None,
            "model_path": str(model_path),
            "env_id": env_id,
            "steps_target": int(job_config["steps_target"]),
            "latest_step": 0,
            "status": "Inicialitzant avaluació...",
            "started_at": time.time(),
        }
    )
    worker_thread.start()

    response["started"] = True
    response["runtime"] = runtime
    return response


def apply_evaluation_action_contract(
    env: gym.Env,
    model_action_space: Any,
    wrapper_configs: List[Dict[str, Any]],
) -> gym.Env:
    """Aplica contracte d'acció d'avaluació."""
    from backend.entrenar_agent_wrappers import apply_training_wrappers

    if wrapper_configs:
        return apply_training_wrappers(env, wrapper_configs)
    if should_apply_normalize_action(model_action_space, env.action_space):
        return NormalizeAction(env)
    return env


def validate_action_contract(model_action_space: Any, env_action_space: Any, metadata: Dict[str, Any]) -> None:
    """Validació del contracte d'actuació."""
    validate_action_spaces(model_action_space, env_action_space, metadata)


def make_env(env_id: str, seed: int = 0, gym_kwargs: Optional[Dict[str, Any]] = None) -> gym.Env:
    """Crea l'entorn utilitzat durant l'avaluació."""
    env = gym.make(env_id, **(gym_kwargs or {}))
    env.reset(seed=seed)
    return env


def run_agent_evaluation(
    model_path: Path,
    env_id: str,
    use_vecnorm: bool,
    vecnorm_path: Optional[Path],
    steps_target: int,
    deterministic: bool,
    should_stop: Callable[[], bool],
    on_progress: Callable[[int, int], None],
    override_wrapper_configs: Optional[List[Dict[str, Any]]] = None,
) -> EvaluationResult:
    """Executa l'avaluació i retorna el resum d'execució."""

    with open(model_path, "rb") as file_handle:
        model, _ = load_sb3_model_bytes(file_handle.read(), device="cpu")

    metadata = load_model_metadata(model_path)
    if override_wrapper_configs is not None:
        wrapper_configs = override_wrapper_configs
    else:
        wrapper_configs = metadata.get("wrapper_configs", [])
    model_action_space = getattr(model, "action_space", None)

    gym_kwargs = build_gym_kwargs_from_metadata(metadata)

    env_factory = partial(build_evaluation_env, env_id, gym_kwargs, model_action_space, wrapper_configs)

    eval_warnings: List[str] = []

    venv = None
    if use_vecnorm and vecnorm_path and vecnorm_path.exists():
        # VecNormalize s'ha d'enganxar al mateix tipus d'entorn amb què es va entrenar.
        # Si no encaixa, seguim sense normalització però deixem l'avís visible.
        try:
            venv = load_vecnormalize(str(vecnorm_path), env_factory)
        except ValueError as exc:
            eval_warnings.append(
                f"VecNormalize no aplicat: {exc} "
                "Continuant sense normalitzacio d'observacions; les prediccions poden ser menys precises."
            )
    if venv is None:
        venv = build_monitored_vec_env(env_factory)

    try:
        validate_action_contract(model_action_space, venv.action_space, metadata)
    except ValueError as exc:
        eval_warnings.append(f"Avis de contracte d'accions: {exc}")

    model_obs_shape = getattr(getattr(model, "observation_space", None), "shape", None)
    env_obs_shape = getattr(getattr(venv, "observation_space", None), "shape", None)
    if model_obs_shape and env_obs_shape and tuple(model_obs_shape) != tuple(env_obs_shape):
        # Aquest error és millor aturar-lo aquí: si les observacions no tenen la mateixa
        # forma, el model podria predir accions aparentment vàlides però sense sentit.
        model_size = int(np.prod(model_obs_shape))
        env_size = int(np.prod(env_obs_shape))
        raise ValueError(
            f"L'espai d'observacions del model ({model_size} variables, forma {model_obs_shape}) "
            f"no coincideix amb l'entorn seleccionat ({env_size} variables, forma {env_obs_shape}). "
            "Assegura't d'avaluar el model amb el mateix entorn i configuracio de wrappers amb que va ser entrenat."
        )

    cancelled = False
    start_time = time.time()
    obs = venv.reset()

    try:
        for step_number in range(steps_target):
            if should_stop():
                cancelled = True
                break

            action, _ = model.predict(obs, deterministic=deterministic)
            formatted_action = format_action_for_env(action, venv)

            next_obs, rewards, dones, infos = venv.step(formatted_action)
            obs = next_obs

            if bool(np.array(dones).reshape(-1)[0]):
                # Alguns entorns acaben episodi abans dels passos demanats; reiniciem per
                # mantenir l'avaluació viva fins al límit triat per l'usuari.
                obs = venv.reset()

            if (step_number + 1) % 200 == 0 or (step_number + 1) == steps_target:
                on_progress(step_number + 1, steps_target)
    finally:
        try:
            base = venv
            if isinstance(base, VecEnvWrapper):
                base = base.venv
            if hasattr(base, "envs") and base.envs:
                base.envs[0].close()
        except Exception:
            pass

    return EvaluationResult(
        elapsed_seconds=time.time() - start_time,
        cancelled=cancelled,
        warnings=tuple(eval_warnings),
    )
\end{lstlisting}

\Needspace{8\baselineskip}

\subsection{\texorpdfstring{\texttt{backend/common.py}}{backend/common.py}}

\begin{lstlisting}[style=studioPython,caption={backend/common.py},label={lst:backend-common-py}]
"""Utilitats compartides de projecte per al backend de BEMS-RL Studio.

Aquest mòdul agrupa peces petites que abans vivien en fitxers molt curts:
camins canònics del projecte, descoberta d'entorns Sinergym i helpers genèrics
d'arguments/format UI. Són funcions utilitzades per diversos fluxos del backend.
"""

from __future__ import annotations

import inspect
from importlib.util import find_spec
from pathlib import Path
from typing import Any, Dict, List, Tuple

import gymnasium as gym


def _sinergym_package_dir() -> Path:
    """Localitza el paquet Sinergym sense executar-ne el registre d'entorns."""
    spec = find_spec("sinergym")
    if spec and spec.origin:
        return Path(spec.origin).resolve().parent
    return Path("sinergym")


PKG_DIR = _sinergym_package_dir()
BUILDINGS_DIR = PKG_DIR / "data" / "buildings"
WEATHERS_DIR = PKG_DIR / "data" / "weather"
CFG_DIR = PKG_DIR / "data" / "default_configuration"

ONE_YEAR_STEPS = 35040


def registered_env_ids() -> tuple[str, ...]:
    """Retorna els ids d'entorn exposats per Sinergym."""
    try:
        import sinergym

        env_ids = sinergym.ids()
    except Exception:
        env_ids = [env_id for env_id in gym.envs.registry.keys() if env_id.startswith("Eplus-")]
    return tuple(sorted(env_ids))


def list_registered_env_ids(*, include_discrete: bool = True) -> List[str]:
    """Llista ids Sinergym, amb opcio d'excloure entorns discrets."""
    env_ids = list(registered_env_ids())
    if include_discrete:
        return env_ids
    return [env_id for env_id in env_ids if "discrete" not in env_id.lower()]


def is_registered_env_id(env_id: str) -> bool:
    """Indica si un id forma part del cataleg registrat de Sinergym."""
    return env_id in registered_env_ids()


def filter_supported_kwargs(callable_obj: Any, kwargs: Dict[str, Any]) -> Tuple[Dict[str, Any], List[str]]:
    """Filtra kwargs no acceptats per una funcio o classe."""
    try:
        signature_target = callable_obj.__init__ if inspect.isclass(callable_obj) else callable_obj
        signature = inspect.signature(signature_target)
    except (TypeError, ValueError):
        return dict(kwargs), []

    parameters = signature.parameters
    if any(parameter.kind == inspect.Parameter.VAR_KEYWORD for parameter in parameters.values()):
        return dict(kwargs), []

    allowed_kwargs = {
        name
        for name, parameter in parameters.items()
        if name != "self"
        and parameter.kind in (inspect.Parameter.POSITIONAL_OR_KEYWORD, inspect.Parameter.KEYWORD_ONLY)
    }
    filtered_kwargs = {key: value for key, value in kwargs.items() if key in allowed_kwargs}
    dropped_kwargs = [key for key in kwargs if key not in allowed_kwargs]
    return filtered_kwargs, dropped_kwargs


def format_ui_value(value: Any) -> str:
    """Formata valors arbitraris per mostrar-los a la UI."""
    if value is None:
        return "-"
    if isinstance(value, dict):
        return ", ".join(f"{key}: {format_ui_value(item)}" for key, item in value.items()) if value else "-"
    if isinstance(value, (list, tuple, set)):
        return ", ".join(format_ui_value(item) for item in value) if value else "-"
    return str(value)
\end{lstlisting}

\Needspace{8\baselineskip}

\subsection{\texorpdfstring{\texttt{backend/entrenar\_agent\_artifacts.py}}{backend/entrenar\_agent\_artifacts.py}}

\begin{lstlisting}[style=studioPython,caption={backend/entrenar\_agent\_artifacts.py},label={lst:backend-entrenar-agent-artifacts-py}]
"""Rutes d'artefacte d'entrenament, metadades i postprocessament CSV."""

from __future__ import annotations

import copy
import csv
import json
import os
import re
from pathlib import Path
from typing import Any, Dict, List

from backend.entrenar_agent_constants import (
    DETAILED_HVAC_METERS,
    JOULES_PER_KWH,
    TRAINING_ARTIFACTS_DIR_NAME,
    TRAINING_CONFIG_FILENAME,
)


def sanitize_training_name_component(value: str) -> str:
    """Higienitzar el component del nom de la entrenament."""
    cleaned = re.sub(r"[^A-Za-z0-9]+", "-", value).strip("-").lower()
    return cleaned or "sense-nom"


def get_training_artifacts_root(base_path: Path | None = None) -> Path:
    """Retorna l'arrel dels artefactes d'entrenament."""
    return (base_path or Path.cwd()).resolve() / TRAINING_ARTIFACTS_DIR_NAME


def build_training_artifact_paths(
    training_config: Dict[str, Any],
    base_path: Path | None = None,
) -> Dict[str, Path | str]:
    """Calcula les rutes dels artefactes d'una sessió d'entrenament."""
    artifacts_root = get_training_artifacts_root(base_path)
    env_slug = sanitize_training_name_component(training_config["env_id"])
    reward_slug = sanitize_training_name_component(training_config["reward_name"])
    family_dir = artifacts_root / env_slug / reward_slug
    prefix = f"training-{env_slug}-{reward_slug}-"
    max_index = 0

    if family_dir.exists():
        for child in family_dir.iterdir():
            if not child.is_dir():
                continue
            match = re.fullmatch(rf"{re.escape(prefix)}(\d+)", child.name)
            if match:
                max_index = max(max_index, int(match.group(1)))

    artifact_name = f"{prefix}{max_index + 1:03d}"
    artifact_dir = family_dir / artifact_name
    return {
        "artifact_name": artifact_name,
        "artifact_dir": artifact_dir,
        "model_path": artifact_dir / f"{artifact_name}.zip",
        "vecnorm_path": artifact_dir / f"{artifact_name}_vecnorm.pkl",
        "eval_script_path": artifact_dir / f"{artifact_name}_eval.py",
        "training_script_path": artifact_dir / f"{artifact_name}.py",
        "config_path": artifact_dir / TRAINING_CONFIG_FILENAME,
    }


def append_detailed_meter_kwh_columns(observations_path: str | Path) -> None:
    """Afegeix columnes detallades del comptador de kWh."""
    observations_file = Path(observations_path) if observations_path else None
    if observations_file is None or not observations_file.is_file():
        return

    with open(observations_file, "r", newline="", encoding="utf-8") as file_handle:
        reader = csv.DictReader(file_handle)
        fieldnames = list(reader.fieldnames or [])
        rows = list(reader)

    if not fieldnames or not rows:
        return

    source_columns: Dict[str, str] = {}
    for alias, meter_name in DETAILED_HVAC_METERS.items():
        target_column = f"{alias}_kwh"
        if target_column in fieldnames:
            continue
        if alias in fieldnames:
            source_columns[alias] = alias
        elif meter_name in fieldnames:
            source_columns[alias] = meter_name

    if not source_columns:
        return

    # EnergyPlus escriu molts meters en joules. Afegim columnes *_kwh perquè els gràfics
    # posteriors no hagin de repetir aquesta conversió cada vegada.
    derived_columns = {source_alias: f"{source_alias}_kwh" for source_alias in source_columns}
    for row in rows:
        for source_alias, source_column in source_columns.items():
            target_column = derived_columns[source_alias]
            try:
                row[target_column] = str(float(row.get(source_column, "")) / JOULES_PER_KWH)
            except (TypeError, ValueError):
                row[target_column] = ""

    output_fields = fieldnames + list(derived_columns.values())
    with open(observations_file, "w", newline="", encoding="utf-8") as file_handle:
        writer = csv.DictWriter(file_handle, fieldnames=output_fields)
        writer.writeheader()
        writer.writerows(rows)


def append_detailed_meter_kwh_columns_in_workspace(workspace_path: str | Path | None) -> None:
    """Afegeix columnes detallades del comptador de kWh a l'espai de treball."""
    if not workspace_path:
        return

    workspace = Path(workspace_path)
    if not workspace.is_dir():
        return

    for observations_file in workspace.rglob("observations.csv"):
        try:
            append_detailed_meter_kwh_columns(observations_file)
        except Exception:
            continue


def get_runtime_workspace_path(runtime: Dict[str, Any]) -> str | None:
    """Retorna la ruta de l'espai de treball en temps d'execució."""
    if not runtime:
        return None

    try:
        paths = runtime["vec"].get_attr("workspace_path")
        if paths and paths[0]:
            return str(paths[0])
    except Exception:
        pass

    # Segon intent: alguns models guarden el VecEnv dins model.env, depenent de com s'ha
    # creat la run i de la versió d'SB3.
    try:
        paths = runtime["model"].env.get_attr("workspace_path")
        if paths and paths[0]:
            return str(paths[0])
    except Exception:
        pass

    try:
        # Últim recurs per a runs antigues: busquem la carpeta més recent que coincideixi
        # amb l'env_id i que contingui els CSV esperats.
        env_id = runtime["config"]["env_id"]
        candidates = sorted(
            [
                directory
                for directory in Path(os.getcwd()).iterdir()
                if directory.is_dir() and directory.name.startswith(f"{env_id}-res")
            ],
            key=lambda directory: directory.stat().st_mtime,
            reverse=True,
        )
        if candidates:
            return str(candidates[0])
    except Exception:
        pass

    return None


def build_training_ui_state(options: Dict[str, Any]) -> Dict[str, Any]:
    """Prepara l'estat inicial de la interfície d'entrenament."""
    excluded_keys = {"spec", "variables"}
    return {
        key: copy.deepcopy(value)
        for key, value in options.items()
        if key not in excluded_keys
    }


def list_saved_training_runs(base_path: Path | None = None) -> List[Dict[str, Any]]:
    """Llista les sessions d'entrenament desades."""
    artifacts_root = get_training_artifacts_root(base_path)
    if not artifacts_root.exists():
        return []

    runs: List[Dict[str, Any]] = []
    for config_path in artifacts_root.rglob(TRAINING_CONFIG_FILENAME):
        try:
            with open(config_path, "r", encoding="utf-8") as handle:
                payload = json.load(handle)
        except Exception:
            continue

        artifact_name = payload.get("artifact_name") or config_path.parent.name
        artifact_dir = Path(payload.get("artifact_dir") or config_path.parent)
        files = payload.get("files") or {}
        training_script_name = files.get("training_script") or f"{artifact_name}.py"
        training_script_path = artifact_dir / training_script_name

        payload["artifact_name"] = artifact_name
        payload["artifact_dir"] = str(artifact_dir)
        payload["config_path"] = str(config_path)
        payload["training_script_path"] = str(training_script_path)
        payload["training_script_exists"] = training_script_path.exists()
        runs.append(payload)

    runs.sort(
        key=lambda item: (
            item.get("created_at") or "",
            item.get("artifact_name") or "",
        ),
        reverse=True,
    )
    return runs
\end{lstlisting}

\Needspace{8\baselineskip}

\subsection{\texorpdfstring{\texttt{backend/entrenar\_agent\_charts.py}}{backend/entrenar\_agent\_charts.py}}

\begin{lstlisting}[style=studioPython,caption={backend/entrenar\_agent\_charts.py},label={lst:backend-entrenar-agent-charts-py}]
"""Gràfics d'entrenament en directe per a la pàgina d'entrenament de l'agent."""

from __future__ import annotations

import os
from pathlib import Path

import pandas as pd
import streamlit as st

from backend.entrenar_agent_session import TRAINING_WORKSPACE_KEY
from backend.grafics.data_loader import _repair_power_metrics_from_progress
from backend.grafics.episode import make_episode_metrics


def _get_workspace_from_runtime(runtime: dict) -> str | None:
    """Extreu el workspace_path del VecEnv actiu.

    Intenta tres estratègies per ordre de preferència:
    1. ``VecEnv.get_attr`` (travessa la cadena de wrappers Gymnasium).
    2. ``model.env.get_attr`` (accés directe a l'entorn intern del model).
    3. Exploració de CWD per trobar directoris que coincideixin amb ``{env_id}-res*`` (retrocés).

    Retorna:
        Camí absolut a l'espai de treball com a cadena, o ``None`` si no
        l'estratègia té èxit.
    """
    # Estratègia 1: via VecEnv.get_attr
    try:
        paths = runtime["vec"].get_attr("workspace_path")
        if paths and paths[0]:
            return str(paths[0])
    except Exception:
        pass

    # Estratègia 2: via model.env
    try:
        paths = runtime["model"].env.get_attr("workspace_path")
        if paths and paths[0]:
            return str(paths[0])
    except Exception:
        pass

    # Estratègia 3: directori més recent a CWD que coincideix amb env_id
    try:
        env_id = runtime["config"]["env_id"]
        cwd = Path(os.getcwd())
        candidates = sorted(
            [d for d in cwd.iterdir() if d.is_dir() and d.name.startswith(f"{env_id}-res")],
            key=lambda d: d.stat().st_mtime,
            reverse=True,
        )
        if candidates:
            return str(candidates[0])
    except Exception:
        pass

    return None


def render_live_training_charts() -> None:
    """Mostra els gràfics d'evolució per episodi llegint ``progress.csv``.

    Llegeix ``progress.csv`` de l'espai de treball actiu i en mostra tres
    subgràfics (recompensa, infracció de temperatura i demanda d'energia). Si no
    l'episodi encara s'ha completat, es mostra un missatge d'espera.
    """
    workspace_str = st.session_state.get(TRAINING_WORKSPACE_KEY)

    if not workspace_str:
        st.caption("Esperant la inicialitzacio del workspace...")
        return

    progress_path = Path(workspace_str) / "progress.csv"
    if not progress_path.exists():
        st.markdown(
            """
            <div class="training-live-wait">
                Els grafics apareixeran en completar-se el primer episodi.
            </div>
            """,
            unsafe_allow_html=True,
        )
        return

    try:
        progress = pd.read_csv(progress_path)
    except Exception:
        return

    if progress.empty:
        return

    try:
        progress = _repair_power_metrics_from_progress(progress)
    except Exception:
        pass

    n_episodes = len(progress)
    st.markdown(
        f"""
        <div class="training-live-status">
            <span class="training-live-badge">
                <span class="training-live-dot"></span>
                EN VIU
            </span>
            <span class="training-live-count">
                {n_episodes} episodi{"s" if n_episodes != 1 else ""} completat{"s" if n_episodes != 1 else ""}
            </span>
        </div>
        """,
        unsafe_allow_html=True,
    )

    fig = make_episode_metrics(progress)

    # Paleta de colors per a cada subgràfic
    _TRACE_COLORS = [
        "#6366f1",  # recompensa -> indigo
        "#f59e0b",  # violation  -> amber
        "#10b981",  # power      -> green
    ]
    _FILL_COLORS = [
        "rgba(99,102,241,0.08)",
        "rgba(245,158,11,0.08)",
        "rgba(16,185,129,0.08)",
    ]

    # Estil cada traça: color, amplada, marcadors i àrea de farciment
    for i, trace in enumerate(fig.data):
        trace.update(
            line=dict(color=_TRACE_COLORS[i], width=2.5),
            marker=dict(
                color="#ffffff",
                size=8,
                line=dict(color=_TRACE_COLORS[i], width=2),
            ),
            fill="tozeroy",
            fillcolor=_FILL_COLORS[i],
        )

    # Distribució manual de dominis amb prou espai per als títols de subgràfics
    row_domains = [
        [0.69, 1.00],
        [0.34, 0.60],
        [0.00, 0.25],
    ]

    # Propietats d'eix comuns
    _axis_style = dict(
        showgrid=True,
        gridcolor="rgba(0,0,0,0.05)",
        gridwidth=1,
        zeroline=False,
        linecolor="rgba(0,0,0,0.12)",
        linewidth=1,
        tickfont=dict(size=11, color="#64748b"),
        title_font=dict(size=11, color="#94a3b8"),
    )

    fig.update_layout(
        title_text="",
        height=680,
        paper_bgcolor="rgba(0,0,0,0)",
        plot_bgcolor="rgba(248,250,255,0.6)",
        margin=dict(t=30, b=20, l=50, r=10),
        font=dict(size=12, color="#334155"),
        yaxis=dict(domain=row_domains[0], **_axis_style),
        yaxis2=dict(domain=row_domains[1], **_axis_style),
        yaxis3=dict(domain=row_domains[2], **_axis_style),
        xaxis=dict(title_text="", **_axis_style),
        xaxis2=dict(title_text="", **_axis_style),
        xaxis3=dict(title_text="Episodi", **_axis_style),
    )

    # Encoixinat de l'eix X perquè els marcadors dels extrems no es retallin
    x_pad = 0.5
    x_min = float(progress["episode_num"].min()) if "episode_num" in progress.columns else 1.0
    x_max = float(progress["episode_num"].max()) if "episode_num" in progress.columns else float(len(progress))
    fig.update_xaxes(range=[x_min - x_pad, x_max + x_pad])

    # Reposicionar i estilitzar les anotacions del títol del subgràfic
    for i, annotation in enumerate(fig.layout.annotations):
        annotation.update(
            y=row_domains[i][1],
            yanchor="bottom",
            font=dict(size=12, color="#475569", family="sans-serif"),
        )

    st.plotly_chart(fig, width="stretch", key="live_training_chart")
\end{lstlisting}

\Needspace{8\baselineskip}

\subsection{\texorpdfstring{\texttt{backend/entrenar\_agent\_constants.py}}{backend/entrenar\_agent\_constants.py}}

\begin{lstlisting}[style=studioPython,caption={backend/entrenar\_agent\_constants.py},label={lst:backend-entrenar-agent-constants-py}]
"""Constants i registres per al backend de l'agent de entrenament."""

from __future__ import annotations

from collections.abc import Iterator, Mapping
from typing import Any

from gymnasium.spaces import Discrete, MultiBinary, MultiDiscrete

from backend.common import ONE_YEAR_STEPS, PKG_DIR


POLICIES = {
    "PPO": ["MlpPolicy"],
    "A2C": ["MlpPolicy"],
    "DQN": ["MlpPolicy"],
    "SAC": ["MlpPolicy"],
    "TD3": ["MlpPolicy"],
    "DDPG": ["MlpPolicy"],
}
REWARD_CLASS_NAMES = (
    "LinearReward",
    "ExpReward",
    "HourlyLinearReward",
    "EnergyCostLinearReward",
    "EnergyCostHourlyReward",
    "NormalizedLinearReward",
    "BatteryHVACReward",
    "MultiZoneReward",
    "OccupancyMultiZoneReward",
    "MultiZoneHourlyReward",
    "MultiZoneEnergyCostReward",
    "MultiZoneEnergyCostHourlyReward",
    "MultiZoneBatteryHVACReward",
    "MultizoneEnergyCostBatteryHVACReward",
    "OccupiedHoursDiscomfortReward",
)

class RewardClassRegistry(Mapping[str, Any]):
    """Registre lazy de classes de recompensa exposades per Sinergym."""

    def __iter__(self) -> Iterator[str]:
        """Itera pels noms de recompenses disponibles sense carregar les classes."""
        return iter(REWARD_CLASS_NAMES)

    def __len__(self) -> int:
        """Retorna el nombre de recompenses registrades."""
        return len(REWARD_CLASS_NAMES)

    def __getitem__(self, reward_name: str) -> Any:
        """Resol una classe de recompensa pel seu nom registrat."""
        if reward_name not in REWARD_CLASS_NAMES:
            raise KeyError(reward_name)
        from backend.model_metadata import resolve_reward_class

        reward_cls = resolve_reward_class(reward_name)
        if reward_cls is None:
            raise KeyError(reward_name)
        return reward_cls


REWARD_CLASSES = RewardClassRegistry()

COST_REWARDS = {
    "EnergyCostLinearReward",
    "EnergyCostHourlyReward",
    "MultiZoneEnergyCostReward",
    "MultiZoneEnergyCostHourlyReward",
}
MULTIZONE_REWARDS = {
    "MultiZoneReward",
    "OccupancyMultiZoneReward",
    "MultiZoneHourlyReward",
    "MultiZoneEnergyCostReward",
    "MultiZoneEnergyCostHourlyReward",
    "MultiZoneBatteryHVACReward",
    "MultizoneEnergyCostBatteryHVACReward",
}
HOURLY_REWARDS = {
    "HourlyLinearReward",
    "EnergyCostHourlyReward",
    "MultiZoneHourlyReward",
    "MultiZoneEnergyCostHourlyReward",
    "OccupiedHoursDiscomfortReward",
}
BATTERY_REWARDS = {
    "BatteryHVACReward",
    "MultiZoneBatteryHVACReward",
    "MultizoneEnergyCostBatteryHVACReward",
}

DEFAULT_MINUTES_PER_STEP = int((365 * 24 * 60) / ONE_YEAR_STEPS)

TRAINING_RUNTIME_KEY = "training_runtime"
TRAINING_RESULT_KEY = "training_result"
TRAINING_ARTIFACTS_DIR_NAME = "trainings"
TRAINING_CONFIG_FILENAME = "training_config.json"
TRAINING_REPRO_SCRIPT_FILENAME = "train_reproduce.py"
DETAILED_HVAC_METERS = {
    "natural_gas_hvac": "NaturalGas:HVAC",
    "heating_electricity": "Heating:Electricity",
    "cooling_electricity": "Cooling:Electricity",
    "fans_electricity": "Fans:Electricity",
    "pumps_electricity": "Pumps:Electricity",
    "heat_rejection_electricity": "HeatRejection:Electricity",
    "humidifier_electricity": "Humidifier:Electricity",
    "heat_recovery_electricity": "HeatRecovery:Electricity",
}
JOULES_PER_KWH = 3_600_000.0

FIELD_HELP = {
    "env_id": "Entorn de Sinergym que vols entrenar. Defineix edifici, clima, observacións i controls.",
    "algo_name": "Algorisme de Stable Baselines3 que farà l'entrenament del model.",
    "policy_name": "Política o arquitectura base que farà servir l'algorisme.",
    "reward_name": "Funcio de recompensa que transforma l'estat de l'entorn en reward a cada pas.",
    "range_winter": "Rang de confort d'hivern en C. Sortir d'aquest rang penalitza la reward.",
    "range_summer": "Rang de confort d'estiu en C. Sortir d'aquest rang penalitza la reward.",
    "energy_weight": "Pes relatiu de l'energia dins de la reward. Més alt = més pes energètic. A OccupiedHoursDiscomfortReward aquest pes s'aplica a les hores ocupades.",
    "temperature_weight": "Pes del terme de confort en rewards amb cost. El pes de cost es calcula com 1 - energia - confort.",
    "lambda_energy": "Factor base d'escala de la penalització energètica.",
    "lambda_temperature": "Factor d'escala de la penalització de temperatura.",
    "grid_energy_weight": "Pes de la potencia importada de xarxa dins de les rewards amb bateria.",
    "battery_cycle_weight": "Pes del cicle de bateria: suma de potencia de carrega i descarrega.",
    "battery_loss_weight": "Pes de les perdues termiques de la bateria si l'entorn les exposa.",
    "simultaneous_battery_weight": "Pes extra per penalitzar carrega i descarrega simultanies.",
    "lambda_grid": "Factor d'escala de la penalitzacio per importacio de xarxa.",
    "lambda_battery": "Factor d'escala de les penalitzacions de bateria.",
    "summer_start_m": "Mes d'inici del periode d'estiu usat per la reward.",
    "summer_start_d": "Dia d'inici del periode d'estiu usat per la reward.",
    "summer_final_m": "Mes final del periode d'estiu usat per la reward.",
    "summer_final_d": "Dia final del periode d'estiu usat per la reward.",
    "lambda_energy_cost": "Factor d'escala del cost economic de l'energia dins de la reward.",
    "comfort_threshold": "Marge de desviacio acceptable respecte al setpoint en rewards multizona.",
    "range_comfort_hours": "Franja horaria en que el confort te més pes dins de la reward.",
    "occupied_hours": "Franja horaria considerada d'ocupacio. Dins d'aquesta franja es considera confort; fora d'ella la reward prioritza reduir consum.",
    "occupied_discomfort_multiplier": "Multiplicador aplicat a la penalització de confort durant les hores d'ocupacio.",
    "off_hours_energy_multiplier": "Multiplicador extra del terme energètic fora d'ocupacio. Més alt = menys consum nocturn i menys preconditioning.",
    "use_energy_cost": "Activa el càlcul del cost energètic amb el fitxer de preus seleccionat.",
    "use_file_logger": "Desa un CSV extra amb el cost estimat a cada pas.",
    "energy_cost_path": "Ruta del fitxer CSV amb els preus de l'energia que farà servir el wrapper de cost.",
    "file_logger_name": "Nom del fitxer CSV on es desa el cost pas a pas.",
    "datetime_wrapper": "Afegeix variables temporals com el mes, el dia i l'hora a l'observació.",
    "previous_wrapper": "Afegeix a l'observació els valors del pas anterior per a les variables seleccionades.",
    "previous_variables": "Variables de l'observació actual que es copiaran des del pas anterior.",
    "multi_obs_wrapper": "Construeix una observació amb diversos passos consecutius per donar context temporal a l'agent.",
    "multi_obs_n": "Nombre de finestres temporals consecutives que es conservaran a l'observació.",
    "multi_obs_flatten": "Si esta activat, la pila temporal es converteix en un vector pla.",
    "normalize_obs_wrapper": "Normalitza l'observació per estabilitzar l'entrenament de l'agent.",
    "normalize_obs_auto": "Actualitza automàticament la mitjana i la variància durant l'entrenament.",
    "normalize_obs_epsilon": "Petit valor de seguretat per evitar divisions per zero en la normalitzacio.",
    "normalize_obs_mean": "Fitxer opcional amb la mitjana precomputada de les observacións.",
    "normalize_obs_var": "Fitxer opcional amb la variància precomputada de les observacións.",
    "weather_forecast_wrapper": "Afegeix a l'observació una previsió meteorològica construida a partir del fitxer climàtic de l'entorn.",
    "weather_forecast_n": "Nombre de mostres futures de previsió que s'afegiran a cada observació.",
    "weather_forecast_delta": "Separacio en passos entre mostres consecutives de previsió.",
    "weather_forecast_columns": "Variables meteorològiques que es faràn servir per construir la previsió.",
    "delta_temp_wrapper": "Calcula diferencies entre temperatures mesurades i setpoints per simplificar el control.",
    "delta_temp_variables": "Variables de temperatura reals que es compararan amb els setpoints.",
    "delta_setpoint_variables": "Variables de consigna que es faràn servir de referència per calcular el delta.",
    "reduce_obs_wrapper": "Elimina variables de l'observació per reduir-ne la dimensio.",
    "reduced_observations": "Variables que s'exclouran de l'observació final.",
    "incremental_wrapper": "Transforma accións continues en increments limitats per variable.",
    "incremental_variables": "Variables d'acció que es controlaran de forma incremental.",
    "incremental_initial_value": "Valor inicial de la variable abans d'aplicar increments.",
    "incremental_delta": "Canvi maxim admissible respecte al valor actual de la variable.",
    "incremental_step": "Pas discret intern amb que es representen els increments.",
    "discrete_incremental_wrapper": "Transforma l'espai d'accións en decisións discretes d'augment, manténiment o reduccio.",
    "discrete_incremental_initial": "Valor inicial per a cada variable controlada en mode discret incremental.",
    "discrete_incremental_delta": "Canvi maxim que es pot acumular en mode discret incremental.",
    "discrete_incremental_step": "Mida del pas aplicat a cada decisió discreta.",
    "normalize_action": "Reescala l'espai d'accións continu a un rang fix per facilitar l'entrenament.",
    "learning_rate": "Velocitat d'aprenentatge del model. Massa alt pot fer l'entrenament inestable.",
    "n_steps": "Nombre de passos recollits abans de cada actualitzacio en PPO i A2C.",
    "timesteps_per_year": f"Nombre total de passos d'entrenament. Amb els valors per defecte, 1 step = {DEFAULT_MINUTES_PER_STEP} minuts i {ONE_YEAR_STEPS} steps = 1 any.",
    "start_training": "Inicia l'entrenament amb la configuració actual.",
    "stop_training": "Atura l'entrenament al final del bloc actual i desa l'estat parcial.",
}

DEFAULT_FORECAST_COLUMNS = [
    "Dry Bulb Temperature",
    "Relative Humidity",
    "Wind Direction",
    "Wind Speed",
    "Direct Normal Radiation",
    "Diffuse Horizontal Radiation",
]

PACKAGE_DATA_DIR = PKG_DIR / "data"
ENERGY_COST_DIR = PACKAGE_DATA_DIR / "energy_cost"

FIXED_WRAPPER_ROWS = [
    {"Wrapper": "Monitor", "Actiu": "Si", "Detall": "Sempre actiu"},
    {"Wrapper": "LoggerWrapper", "Actiu": "Si", "Detall": "Sempre actiu"},
    {"Wrapper": "CSVLogger", "Actiu": "Si", "Detall": "Sempre actiu"},
]

DISCRETE_ACTION_SPACES = (Discrete, MultiDiscrete, MultiBinary)
\end{lstlisting}

\Needspace{8\baselineskip}

\subsection{\texorpdfstring{\texttt{backend/entrenar\_agent\_env.py}}{backend/entrenar\_agent\_env.py}}

\begin{lstlisting}[style=studioPython,caption={backend/entrenar\_agent\_env.py},label={lst:backend-entrenar-agent-env-py}]
"""Descobriment d'entorns i metadades per configurar entrenaments."""

from __future__ import annotations

import os
from pathlib import Path
from typing import Any, Dict, List, Sequence

import gymnasium as gym
import yaml
from gymnasium.spaces import Box

from backend.entrenar_agent_constants import DISCRETE_ACTION_SPACES
from backend.common import CFG_DIR, list_registered_env_ids


def list_registered_envs() -> List[str]:
    """Llista d'envs registrats."""
    return list_registered_env_ids()


def list_files(directory: Path, suffixes: Sequence[str]) -> List[str]:
    """Llista de fitxers."""
    if not directory.exists():
        return []
    return sorted(
        str(path)
        for path in directory.iterdir()
        if path.is_file() and path.suffix.lower() in suffixes
    )


def with_detailed_hvac_meters(meters: Dict[str, str] | None) -> Dict[str, str]:
    """Retorna només comptadors configurats.

    Inventar els metres HVAC predeterminats fa que EnergyPlus retorni els identificadors API no vàlids
    per a edificis que no els exposen, omplint eplusout.err amb errors `Severe getMeterValue`
    errors getMeterValue. Els comptadors detallats encara s'utilitzen quan es generen
    la configuració de l'entorn els va detectar explícitament.
    """
    return dict(meters or {})


def get_training_meters_for_env(env_id: str) -> Dict[str, str]:
    """Retorna els mesuradors d'entrenament per a env."""
    try:
        spec = gym.spec(env_id)
        env_kwargs = spec.kwargs or {}
        meters = env_kwargs.get("meters", {}) or {}
    except Exception:
        meters = {}
    return with_detailed_hvac_meters(meters)


def _effective_action_space_from_spec(spec: gym.envs.registration.EnvSpec, fallback: Any) -> Any:
    """Espai d'acció efectiu des de les especificacions."""
    for wrapper_spec in getattr(spec, "additional_wrappers", ()) or ():
        if getattr(wrapper_spec, "name", "") != "DiscretizeEnv":
            continue
        wrapper_kwargs = getattr(wrapper_spec, "kwargs", {}) or {}
        discrete_space = wrapper_kwargs.get("discrete_space")
        if discrete_space is not None:
            return discrete_space
    return fallback


def get_env_metadata(env_id: str) -> Dict[str, Any]:
    """Retorna les metadades de l'env."""
    spec = gym.spec(env_id)
    env_kwargs = dict(spec.kwargs or {})
    variables = env_kwargs.get("variables", {}) or {}
    meters = with_detailed_hvac_meters(env_kwargs.get("meters", {}) or {})
    env_kwargs["meters"] = meters
    actuators = env_kwargs.get("actuators", {}) or {}
    context = env_kwargs.get("context", {}) or {}
    time_variables = env_kwargs.get("time_variables", []) or []
    observation_variables = time_variables + list(variables.keys()) + list(meters.keys())
    action_variables = list(actuators.keys())
    context_variables = list(context.keys())
    base_action_space = env_kwargs.get("action_space")
    action_space = _effective_action_space_from_spec(spec, base_action_space)
    building_file = env_kwargs.get("building_file")
    building_name = os.path.basename(building_file) if building_file else ""
    return {
        "spec": spec,
        "env_kwargs": env_kwargs,
        "variables": variables,
        "meters": meters,
        "actuators": actuators,
        "context": context,
        "time_variables": time_variables,
        "observation_variables": observation_variables,
        "action_variables": action_variables,
        "context_variables": context_variables,
        "action_space": action_space,
        "base_action_space": base_action_space,
        "is_continuous": isinstance(action_space, Box),
        "is_discrete": isinstance(action_space, DISCRETE_ACTION_SPACES),
        "building_file": building_file,
        "building_name": building_name,
        "candidate_temp_vars": default_comfort_temperature_variables(observation_variables, variables),
        "candidate_energy_vars": [
            key for key in observation_variables if is_energy_variable(key, variables)
        ],
        "candidate_grid_energy_vars": default_grid_energy_variables(observation_variables),
        "candidate_battery_charge_vars": default_battery_variables(observation_variables, "charge"),
        "candidate_battery_discharge_vars": default_battery_variables(observation_variables, "discharge"),
        "candidate_battery_loss_vars": default_battery_variables(observation_variables, "loss"),
        "candidate_cost_vars": [key for key in observation_variables if ("cost" in key.lower() or "billing" in key.lower())],
        "candidate_setpoint_vars": [key for key in variables if "setpoint" in key.lower()],
    }


def is_comfort_temperature_variable(variable_name: str, variables: Dict[str, Any]) -> bool:
    """És variable la temperatura de confort."""
    variable_name_lower = str(variable_name).lower()
    if "temperature" not in variable_name_lower:
        return False
    if "outdoor" in variable_name_lower or "setpoint" in variable_name_lower:
        return False

    variable_meta = variables.get(variable_name)
    if isinstance(variable_meta, (list, tuple)):
        output_name = str(variable_meta[0]).lower() if len(variable_meta) > 0 else ""
        output_key = str(variable_meta[1]).lower() if len(variable_meta) > 1 else ""
        if "outdoor" in output_name:
            return False
        if output_key in {"environment", "site", "whole building"}:
            return False
        if "zone air temperature" in output_name:
            return True

    return "air_temperature" in variable_name_lower


def default_comfort_temperature_variables(
    variable_names: Sequence[str], variables: Dict[str, Any]
) -> List[str]:
    """Variables de temperatura de confort per defecte."""
    return [
        variable_name
        for variable_name in variable_names
        if variable_name in variables and is_comfort_temperature_variable(variable_name, variables)
    ]


def is_energy_variable(variable_name: str, variables: Dict[str, Any]) -> bool:
    """És l'energia variable."""
    variable_name_lower = str(variable_name).lower()
    if any(
        token in variable_name_lower
        for token in ("energy", "electricity", "elèctricity", "power", "hvac", "demand")
    ):
        return True

    variable_meta = variables.get(variable_name)
    if isinstance(variable_meta, (list, tuple)):
        output_name = str(variable_meta[0]).lower() if len(variable_meta) > 0 else ""
        output_key = str(variable_meta[1]).lower() if len(variable_meta) > 1 else ""
        if any(
            token in output_name
            for token in ("electricity", "power", "demand rate", "hvac")
        ):
            return True
        if any(token in output_key for token in ("whole building", "hvac")):
            return True

    return False


def default_grid_energy_variables(variable_names: Sequence[str]) -> List[str]:
    """Variables d'energia de xarxa per defecte."""
    available = list(variable_names)
    preferred = [
        "facility_net_purchased_electricity_rate",
        "facility_total_purchased_electricity_rate",
        "facility_total_electricity_demand_rate",
        "facility_total_building_electricity_demand_rate",
    ]
    for variable_name in preferred:
        if variable_name in available:
            return [variable_name]
    return [
        variable_name
        for variable_name in available
        if "purchased" in variable_name.lower() or "grid" in variable_name.lower()
    ][:1]


def default_battery_variables(variable_names: Sequence[str], kind: str) -> List[str]:
    """Variables per defecte de la bateria."""
    available = list(variable_names)
    kind_tokens = {
        "charge": ("storage_charge_power", "battery_charge"),
        "discharge": ("storage_discharge_power", "battery_discharge"),
        "loss": ("storage_thermal_loss_rate", "battery_loss"),
    }
    tokens = kind_tokens.get(kind, ())
    matches = [
        variable_name
        for variable_name in available
        if any(token in variable_name.lower() for token in tokens)
    ]
    if matches:
        return matches[:1]

    fallback_tokens = {
        "charge": ("charge",),
        "discharge": ("discharge",),
        "loss": ("loss", "thermal"),
    }.get(kind, ())
    return [
        variable_name
        for variable_name in available
        if "storage" in variable_name.lower()
        and any(token in variable_name.lower() for token in fallback_tokens)
    ][:1]


def load_default_reward_variables(
    spec: Any, env_id: str, variables: Dict[str, Any]
) -> Tuple[Dict[str, Any], List[str], List[str], List[str]]:
    """Carrega les variables de recompensa predeterminades."""
    resolved_variables = variables
    try:
        env_kwargs = dict(getattr(spec, "kwargs", {}) or {})
        meters = env_kwargs.get("meters", {}) or {}
        available_observation_names = list(dict.fromkeys([*variables.keys(), *meters.keys()]))

        spec_reward_cfg = (
            env_kwargs.get("reward_kwargs")
            if isinstance(env_kwargs.get("reward_kwargs"), dict)
            else {}
        )
        temp_vars: List[str] = list(spec_reward_cfg.get("temperature_variables") or [])
        energy_vars: List[str] = list(spec_reward_cfg.get("energy_variables") or [])
        cost_vars: List[str] = list(spec_reward_cfg.get("energy_cost_variables") or [])

        building_file = env_kwargs.get("building_file")
        base_name = os.path.splitext(os.path.basename(building_file))[0] if building_file else env_id
        cfg_file = CFG_DIR / f"{base_name}.yaml"

        if os.path.isfile(cfg_file) and not (temp_vars and energy_vars):
            # Preferim la configuració oficial de Sinergym quan existeix; porta els noms
            # de variables que l'entorn espera per defecte.
            with open(cfg_file, "r", encoding="utf-8") as cfg_handle:
                cfg = yaml.safe_load(cfg_handle) or {}
            reward_cfg = cfg.get("reward_kwargs", {}) if isinstance(cfg.get("reward_kwargs"), dict) else {}
            temp_vars = temp_vars or (
                reward_cfg.get("temperature_variables")
                or cfg.get("temperature_variables")
                or []
            )
            energy_vars = energy_vars or (
                reward_cfg.get("energy_variables")
                or cfg.get("energy_variables")
                or []
            )
            cost_vars = cost_vars or (
                reward_cfg.get("energy_cost_variables")
                or cfg.get("energy_cost_variables")
                or []
            )

        if not temp_vars:
            # Si el YAML no ajuda, fem una inferència prudent a partir dels noms de variables.
            temp_vars = default_comfort_temperature_variables(variables.keys(), variables)
        if not energy_vars:
            energy_vars = [
                key
                for key in available_observation_names
                if is_energy_variable(key, variables)
            ]
        if not cost_vars:
            cost_vars = [
                key
                for key in available_observation_names
                if "cost" in key.lower() or "billing" in key.lower()
            ]

        temp_vars = [value for value in temp_vars if value in variables]
        available_observation_set = set(available_observation_names)
        energy_vars = [value for value in energy_vars if value in available_observation_set]
        cost_vars = [value for value in cost_vars if value in available_observation_set]
        # Tornem a sanejar temperatures per evitar setpoints o variables globals que hagin
        # entrat per heurística però no siguin temperatura interior real.
        sanitized_temp_vars = default_comfort_temperature_variables(temp_vars, variables)
        if sanitized_temp_vars:
            temp_vars = sanitized_temp_vars
        return resolved_variables, temp_vars, energy_vars, cost_vars
    except Exception:
        return {}, [], [], []
\end{lstlisting}

\Needspace{8\baselineskip}

\subsection{\texorpdfstring{\texttt{backend/entrenar\_agent\_rewards.py}}{backend/entrenar\_agent\_rewards.py}}

\begin{lstlisting}[style=studioPython,caption={backend/entrenar\_agent\_rewards.py},label={lst:backend-entrenar-agent-rewards-py}]
"""Recompenses, algorismes i creadors de càrrega útil de entrenament completa."""

from __future__ import annotations

from typing import Any, Dict, List, Sequence, Tuple

from backend.entrenar_agent_artifacts import build_training_ui_state
from backend.entrenar_agent_constants import (
    BATTERY_REWARDS,
    COST_REWARDS,
    DEFAULT_MINUTES_PER_STEP,
    MULTIZONE_REWARDS,
    REWARD_CLASSES,
)
from backend.entrenar_agent_env import (
    default_battery_variables,
    default_grid_energy_variables,
    get_training_meters_for_env,
    load_default_reward_variables,
)
from backend.common import filter_supported_kwargs, format_ui_value
from backend.entrenar_agent_wrappers import (
    build_energy_cost_wrapper_reward_kwargs,
    build_wrapper_configs,
    build_wrapper_summary_rows,
)
from backend.sb3_utils import ALGOS


def build_multizone_temp_mapping(
    variables: Dict[str, Any], temp_vars: Sequence[str], pair_setpoints: bool = False
) -> Dict[str, Any]:
    """Prepara un mapa de temperatures i setpoints per zones."""
    # Les variables de Sinergym guarden el domini/zona a la segona posició. L'aprofitem
    # per emparellar cada temperatura amb els seus setpoints sense dependre només del nom.
    setpoint_vars = [key for key in variables if "setpoint" in key.lower()]
    heating_setpoints = [
        key for key in setpoint_vars if "htg" in key.lower() or "heating" in key.lower()
    ]
    cooling_setpoints = [
        key for key in setpoint_vars if "clg" in key.lower() or "cooling" in key.lower()
    ]
    temp_mapping: Dict[str, Any] = {}
    for temp_var in temp_vars:
        domain = variables[temp_var][1] if temp_var in variables else None
        matches = [
            setpoint_var
            for setpoint_var in setpoint_vars
            if domain and variables.get(setpoint_var, [None, None])[1] == domain
        ]
        if pair_setpoints:
            # Algunes rewards volen parella [calor, fred]. Si no trobem coincidència de zona,
            # fem fallback a la primera parella global disponible.
            heating_matches = [match for match in matches if match in heating_setpoints]
            cooling_matches = [match for match in matches if match in cooling_setpoints]
            if heating_matches and cooling_matches:
                temp_mapping[temp_var] = [heating_matches[0], cooling_matches[0]]
            elif matches:
                temp_mapping[temp_var] = matches[0]
            elif heating_setpoints and cooling_setpoints:
                temp_mapping[temp_var] = [heating_setpoints[0], cooling_setpoints[0]]
            elif setpoint_vars:
                temp_mapping[temp_var] = setpoint_vars[0]
        elif matches:
            temp_mapping[temp_var] = matches[0]
        elif len(setpoint_vars) == 1:
            temp_mapping[temp_var] = setpoint_vars[0]
    return temp_mapping


def build_multizone_occupancy_mapping(
    variables: Dict[str, Any], temp_vars: Sequence[str]
) -> Dict[str, str]:
    """Prepara un mapa d'ocupació per zones."""
    occupancy_vars = [
        key
        for key in variables
        if "occupant" in key.lower() or "occupancy" in key.lower()
    ]
    occupancy_mapping: Dict[str, str] = {}
    for temp_var in temp_vars:
        temp_meta = variables.get(temp_var)
        temp_domain = temp_meta[1] if isinstance(temp_meta, (list, tuple)) and len(temp_meta) > 1 else None
        matches = [
            occupancy_var
            for occupancy_var in occupancy_vars
            if temp_domain
            and isinstance(variables.get(occupancy_var), (list, tuple))
            and len(variables[occupancy_var]) > 1
            and variables[occupancy_var][1] == temp_domain
        ]
        if not matches:
            temp_prefix = temp_var.lower().replace("_air_temperature", "")
            matches = [
                occupancy_var
                for occupancy_var in occupancy_vars
                if occupancy_var.lower().startswith(temp_prefix)
            ]
        if matches:
            occupancy_mapping[temp_var] = matches[0]
    return occupancy_mapping


def _build_common_reward_kwargs(
    options: Dict[str, Any],
    temp_vars: Sequence[str],
    energy_vars: Sequence[str],
) -> Dict[str, Any]:
    """Crea els paràmetres compartits per les rewards lineals i de bateria."""
    return {
        "temperature_variables": list(temp_vars),
        "energy_variables": list(energy_vars),
        "range_comfort_winter": tuple(options["range_winter"]),
        "range_comfort_summer": tuple(options["range_summer"]),
        "summer_start": (options["summer_start_m"], options["summer_start_d"]),
        "summer_final": (options["summer_final_m"], options["summer_final_d"]),
        "energy_weight": options["energy_weight"],
        "lambda_energy": options["lambda_energy"],
        "lambda_temperature": options["lambda_temperature"],
    }


def _build_battery_reward_kwargs(
    options: Dict[str, Any],
    variables: Dict[str, Any],
) -> Dict[str, Any]:
    """Crea la part comuna de les rewards que penalitzen xarxa i bateria."""
    available_variables = variables.keys()
    grid_energy_variables = (
        list(options["grid_energy_variables"])
        if "grid_energy_variables" in options
        else default_grid_energy_variables(available_variables)
    )
    battery_charge_variables = (
        list(options["battery_charge_variables"])
        if "battery_charge_variables" in options
        else default_battery_variables(available_variables, "charge")
    )
    battery_discharge_variables = (
        list(options["battery_discharge_variables"])
        if "battery_discharge_variables" in options
        else default_battery_variables(available_variables, "discharge")
    )
    battery_loss_variables = (
        list(options["battery_loss_variables"])
        if "battery_loss_variables" in options
        else default_battery_variables(available_variables, "loss")
    )

    return {
        "grid_energy_variables": grid_energy_variables,
        "battery_charge_variables": battery_charge_variables,
        "battery_discharge_variables": battery_discharge_variables,
        "battery_loss_variables": battery_loss_variables,
        "grid_energy_weight": options.get("grid_energy_weight", 0.2),
        "battery_cycle_weight": options.get("battery_cycle_weight", 0.04),
        "battery_loss_weight": options.get("battery_loss_weight", 0.005),
        "simultaneous_battery_weight": options.get("simultaneous_battery_weight", 0.005),
        "lambda_grid": options.get("lambda_grid", options["lambda_energy"]),
        "lambda_battery": options.get("lambda_battery", options["lambda_energy"]),
    }


def _build_multizone_base_kwargs(
    options: Dict[str, Any],
    energy_vars: Sequence[str],
    temp_mapping: Dict[str, Any],
) -> Dict[str, Any]:
    """Crea la base de kwargs que comparteixen les rewards multizona."""
    return {
        "energy_variables": list(energy_vars),
        "temperature_and_setpoints_conf": temp_mapping,
        "comfort_threshold": options["comfort_threshold"],
        "energy_weight": options["energy_weight"],
        "lambda_energy": options["lambda_energy"],
        "lambda_temperature": options["lambda_temperature"],
    }


def _add_energy_cost_file_kwargs(
    reward_kwargs: Dict[str, Any],
    options: Dict[str, Any],
    cost_vars: Sequence[str],
) -> None:
    """Afegeix cost horari i fitxer de preus quan la reward ho necessita."""
    reward_kwargs.update(
        {
            "lambda_energy_cost": options["lambda_energy_cost"],
            "energy_cost_variables": list(cost_vars),
            "energy_cost_path": options["energy_cost_path"],
            "seconds_per_step": DEFAULT_MINUTES_PER_STEP * 60,
        }
    )


def _build_single_zone_reward_kwargs(
    reward_name: str,
    options: Dict[str, Any],
    reward_common: Dict[str, Any],
    battery_kwargs: Dict[str, Any],
    temp_vars: Sequence[str],
    energy_vars: Sequence[str],
    cost_vars: Sequence[str],
) -> Dict[str, Any] | None:
    """Crea kwargs per rewards que treballen amb variables globals o monozona."""
    if reward_name == "EnergyCostLinearReward":
        reward_kwargs = dict(reward_common)
        reward_kwargs["energy_cost_variables"] = list(cost_vars)
        reward_kwargs["temperature_weight"] = options["temperature_weight"]
        reward_kwargs["lambda_energy_cost"] = options["lambda_energy_cost"]
        return reward_kwargs

    if reward_name == "BatteryHVACReward":
        reward_kwargs = dict(reward_common)
        reward_kwargs["temperature_weight"] = options["temperature_weight"]
        reward_kwargs.update(battery_kwargs)
        return reward_kwargs

    if reward_name == "HourlyLinearReward":
        return {
            "temperature_variables": list(temp_vars),
            "energy_variables": list(energy_vars),
            "range_comfort_winter": tuple(options["range_winter"]),
            "range_comfort_summer": tuple(options["range_summer"]),
            "summer_start": (options["summer_start_m"], options["summer_start_d"]),
            "summer_final": (options["summer_final_m"], options["summer_final_d"]),
            "lambda_energy": options["lambda_energy"],
            "lambda_temperature": options["lambda_temperature"],
            "range_comfort_hours": tuple(options["range_comfort_hours"]),
        }

    if reward_name == "EnergyCostHourlyReward":
        reward_kwargs = dict(reward_common)
        reward_kwargs["temperature_weight"] = options["temperature_weight"]
        reward_kwargs["range_comfort_hours"] = tuple(options["range_comfort_hours"])
        _add_energy_cost_file_kwargs(reward_kwargs, options, cost_vars)
        return reward_kwargs

    if reward_name == "OccupiedHoursDiscomfortReward":
        reward_kwargs = dict(reward_common)
        reward_kwargs["occupied_hours"] = tuple(options["occupied_hours"])
        reward_kwargs["occupied_discomfort_multiplier"] = options["occupied_discomfort_multiplier"]
        reward_kwargs["off_hours_energy_multiplier"] = options["off_hours_energy_multiplier"]
        return reward_kwargs

    return None


def _build_multizone_reward_kwargs(
    reward_name: str,
    options: Dict[str, Any],
    energy_vars: Sequence[str],
    cost_vars: Sequence[str],
    temp_mapping: Dict[str, Any],
    occupancy_mapping: Dict[str, str],
    battery_kwargs: Dict[str, Any],
) -> Dict[str, Any] | None:
    """Crea kwargs per rewards que separen temperatura i consignes per zona."""
    if reward_name == "MultiZoneReward":
        return _build_multizone_base_kwargs(options, energy_vars, temp_mapping)

    if reward_name == "OccupancyMultiZoneReward":
        reward_kwargs = _build_multizone_base_kwargs(options, energy_vars, temp_mapping)
        reward_kwargs["occupancy_variables_conf"] = occupancy_mapping
        reward_kwargs["occupancy_threshold"] = 0.0
        return reward_kwargs

    if reward_name == "MultiZoneHourlyReward":
        reward_kwargs = _build_multizone_base_kwargs(options, energy_vars, temp_mapping)
        reward_kwargs["default_energy_weight"] = reward_kwargs.pop("energy_weight")
        reward_kwargs["range_comfort_hours"] = tuple(options["range_comfort_hours"])
        return reward_kwargs

    if reward_name == "MultiZoneEnergyCostReward":
        reward_kwargs = _build_multizone_base_kwargs(options, energy_vars, temp_mapping)
        reward_kwargs["temperature_weight"] = options["temperature_weight"]
        _add_energy_cost_file_kwargs(reward_kwargs, options, cost_vars)
        return reward_kwargs

    if reward_name == "MultiZoneEnergyCostHourlyReward":
        reward_kwargs = _build_multizone_base_kwargs(options, energy_vars, temp_mapping)
        reward_kwargs["temperature_weight"] = options["temperature_weight"]
        reward_kwargs["range_comfort_hours"] = tuple(options["range_comfort_hours"])
        _add_energy_cost_file_kwargs(reward_kwargs, options, cost_vars)
        return reward_kwargs

    if reward_name == "MultiZoneBatteryHVACReward":
        reward_kwargs = _build_multizone_base_kwargs(options, energy_vars, temp_mapping)
        reward_kwargs["temperature_weight"] = options["temperature_weight"]
        reward_kwargs.update(battery_kwargs)
        return reward_kwargs

    if reward_name == "MultizoneEnergyCostBatteryHVACReward":
        reward_kwargs = _build_multizone_base_kwargs(options, energy_vars, temp_mapping)
        reward_kwargs["temperature_weight"] = options["temperature_weight"]
        reward_kwargs["lambda_energy_cost"] = options["lambda_energy_cost"]
        reward_kwargs["energy_cost_weight"] = options.get("energy_cost_weight", 0.20)
        reward_kwargs["energy_cost_variables"] = list(cost_vars)
        reward_kwargs["occupied_hours"] = tuple(options.get("occupied_hours", (8, 18)))
        reward_kwargs["occupied_comfort_multiplier"] = options.get("occupied_comfort_multiplier", 2.0)
        reward_kwargs["unoccupied_comfort_multiplier"] = options.get("unoccupied_comfort_multiplier", 0.3)
        reward_kwargs.update(battery_kwargs)
        return reward_kwargs

    return None


def build_reward_kwargs(
    options: Dict[str, Any],
    variables: Dict[str, Any],
    temp_vars: Sequence[str],
    energy_vars: Sequence[str],
    cost_vars: Sequence[str],
) -> Tuple[Dict[str, Any], Dict[str, str], List[str]]:
    """Munta els kwargs de recompensa a partir de la configuració."""
    reward_name = options["reward_name"]
    reward_common = _build_common_reward_kwargs(options, temp_vars, energy_vars)

    # Les rewards multizona volen un diccionari zona -> consignes. El construïm abans
    # de seleccionar la reward perquè també serveix per validar la configuració.
    temp_mapping = build_multizone_temp_mapping(
        variables,
        temp_vars,
        pair_setpoints=reward_name in ("MultiZoneBatteryHVACReward", "MultizoneEnergyCostBatteryHVACReward"),
    )
    occupancy_mapping = build_multizone_occupancy_mapping(variables, temp_vars)
    battery_kwargs = _build_battery_reward_kwargs(options, variables)

    reward_kwargs = _build_single_zone_reward_kwargs(
        reward_name=reward_name,
        options=options,
        reward_common=reward_common,
        battery_kwargs=battery_kwargs,
        temp_vars=temp_vars,
        energy_vars=energy_vars,
        cost_vars=cost_vars,
    )
    if reward_kwargs is None:
        reward_kwargs = _build_multizone_reward_kwargs(
            reward_name=reward_name,
            options=options,
            energy_vars=energy_vars,
            cost_vars=cost_vars,
            temp_mapping=temp_mapping,
            occupancy_mapping=occupancy_mapping,
            battery_kwargs=battery_kwargs,
        )
    if reward_kwargs is None:
        reward_kwargs = dict(reward_common)

    reward_cls = REWARD_CLASSES[reward_name]
    reward_kwargs, dropped_reward_kwargs = filter_supported_kwargs(reward_cls, reward_kwargs)
    return reward_kwargs, temp_mapping, dropped_reward_kwargs


def validate_training_configuration(
    options: Dict[str, Any], temp_mapping: Dict[str, str]
) -> List[str]:
    """Valida la recompensa, l'algorisme, l'espai d'acció i la compatibilitat dels wrappers."""
    validation_errors: List[str] = []
    reward_name = options["reward_name"]
    algo_name = options["algo_name"]
    action_space_type = options.get("action_space_type")

    # Primer descartem combinacions que SB3 no accepta, abans de crear entorns o models.
    if options.get("is_discrete") and algo_name in {"SAC", "TD3", "DDPG"}:
        validation_errors.append(
            f"{algo_name} nomes es pot entrenar amb espais d'accio continus. Fes servir PPO o A2C en aquest entorn."
        )
    if action_space_type in {"MultiDiscrete", "MultiBinary"} and algo_name == "DQN":
        validation_errors.append(
            "DQN nomes admet un Discrete simple; aquest entorn fa servir un espai d'accio multidiscret."
        )

    if reward_name in COST_REWARDS and (options["energy_weight"] + options["temperature_weight"] > 1.0):
        validation_errors.append("En rewards amb cost cal complir energy_weight + temperature_weight <= 1.")
    if reward_name in MULTIZONE_REWARDS and not temp_mapping:
        validation_errors.append(
            "No s'han pogut inferir setpoints compatibles per construir la reward multizona."
        )
    if reward_name in BATTERY_REWARDS:
        # Les rewards amb bateria necessiten com a mínim saber què entra de la xarxa
        # i alguna pista de càrrega o descàrrega; si no, la penalització queda buida.
        available_variables = list((options.get("variables") or {}).keys())
        grid_vars = (
            options["grid_energy_variables"]
            if "grid_energy_variables" in options
            else default_grid_energy_variables(available_variables)
        )
        charge_vars = (
            options["battery_charge_variables"]
            if "battery_charge_variables" in options
            else default_battery_variables(available_variables, "charge")
        )
        discharge_vars = (
            options["battery_discharge_variables"]
            if "battery_discharge_variables" in options
            else default_battery_variables(available_variables, "discharge")
        )
        if not grid_vars:
            validation_errors.append(
                "No s'han detectat variables de xarxa per a la reward amb bateria."
            )
        if not charge_vars and not discharge_vars:
            validation_errors.append(
                "No s'han detectat variables de càrrega o descàrrega de bateria."
            )
    # A partir d'aquí validem dependències entre wrappers. Són regles de UI, no errors
    # de Python: millor avisar amb missatges clars que deixar que falli al mig del training.
    if options["enable_previous_wrapper"] and not options["previous_variables"]:
        validation_errors.append("PreviousObservationWrapper requereix almenys una variable.")
    if options["enable_weather_forecast_wrapper"] and not options["weather_forecast_columns"]:
        validation_errors.append("WeatherForecastingWrapper requereix almenys una variable de previsió.")
    if options["enable_delta_temp_wrapper"] and (
        not options["delta_temp_variables"] or not options["delta_setpoint_variables"]
    ):
        validation_errors.append("DeltaTempWrapper requereix variables de temperatura i almenys un setpoint.")
    if options["enable_reduce_obs_wrapper"] and options["enable_datetime_wrapper"]:
        reserved_datetime_vars = {"month", "day_of_month", "hour"}
        invalid_reduction = reserved_datetime_vars.intersection(options["reduced_observations"])
        if invalid_reduction:
            validation_errors.append(
                "No pots reduir month/day_of_month/hour mentre DatetimeWrapper estigui actiu."
            )
    if options["enable_incremental_wrapper"] and not options["incremental_variables"]:
        validation_errors.append("IncrementalWrapper requereix almenys una variable d'acció.")
    if options["enable_incremental_wrapper"] and options["enable_discrete_incremental_wrapper"]:
        validation_errors.append("IncrementalWrapper i DiscreteIncrementalWrapper no es poden activar alhora.")
    if (
        options["enable_incremental_wrapper"] or options["enable_discrete_incremental_wrapper"]
    ) and options["enable_normalize_action"]:
        validation_errors.append("NormalizeAction no es pot combinar amb els wrappers incrementals en aquesta UI.")
    if options["enable_normalize_action"] and not options.get("is_continuous", True):
        validation_errors.append("NormalizeAction nomes es pot aplicar a entorns amb action_space continu.")
    if options.get("enable_variability_context_wrapper"):
        low_values = list(options.get("variability_context_low") or [])
        high_values = list(options.get("variability_context_high") or [])
        if not low_values or len(low_values) != len(high_values):
            validation_errors.append("VariabilityContextWrapper requereix limits min/max coherents.")
        if any(high <= low for low, high in zip(low_values, high_values)):
            validation_errors.append("A VariabilityContextWrapper cada maxim ha de ser superior al minim.")
        if options.get("variability_step_frequency_max", 0) <= options.get("variability_step_frequency_min", 0):
            validation_errors.append("La frequencia maxima del context ha de ser superior a la minima.")
    if (
        options.get("enable_storage_smoothing_wrapper")
        and options.get("building_name") != "OfficeGridStorageSmoothing.epJSON"
    ):
        validation_errors.append(
            "OfficeGridStorageSmoothingActionConstraintsWrapper nomes es valid per OfficeGridStorageSmoothing.epJSON."
        )

    return validation_errors


# Àlies compatibles enrere per a l'identificador públic amb accent anterior.
globals()["vàlidate_training_configuration"] = validate_training_configuration
globals()["vālidate_training_configuration"] = validate_training_configuration


def build_algo_kwargs(options: Dict[str, Any]) -> Tuple[Dict[str, Any], List[str]]:
    """Munta els kwargs de l'algorisme a partir de la configuració."""
    candidate_kwargs = {
        "learning_rate": options["learning_rate"],
        "n_steps": options.get("n_steps"),
        "batch_size": options.get("batch_size"),
        "n_epochs": options.get("n_epochs"),
        "gamma": options.get("gamma"),
        "gae_lambda": options.get("gae_lambda"),
        "clip_range": options.get("clip_range"),
        "ent_coef": options.get("ent_coef"),
        "vf_coef": options.get("vf_coef"),
        "max_grad_norm": options.get("max_grad_norm"),
    }
    candidate_kwargs = {
        key: value
        for key, value in candidate_kwargs.items()
        if value is not None
    }
    return filter_supported_kwargs(ALGOS[options["algo_name"]], candidate_kwargs)


def assemble_training_payload(options: Dict[str, Any]) -> Dict[str, Any]:
    """Munta tota la configuració d'entrenament a partir de la UI."""
    meters = get_training_meters_for_env(options["env_id"])
    variables, temp_vars, energy_vars, cost_vars = load_default_reward_variables(
        spec=options["spec"],
        env_id=options["env_id"],
        variables=options["variables"],
    )
    reward_kwargs, temp_mapping, dropped_reward_kwargs = build_reward_kwargs(
        options=options,
        variables=variables,
        temp_vars=temp_vars,
        energy_vars=energy_vars,
        cost_vars=cost_vars,
    )
    validation_errors = validate_training_configuration(options, temp_mapping)
    if not energy_vars:
        validation_errors.append(
            "No s'han pogut inferir variables d'energia per a la reward. Revisa la configuració de l'entorn."
        )
    if options["reward_name"] == "OccupancyMultiZoneReward":
        # Aquesta reward necessita ocupació per cada temperatura/zona. Si falta algun
        # mapa, millor fallar abans que entrenar amb una penalització incompleta.
        occupancy_mapping = reward_kwargs.get("occupancy_variables_conf") or {}
        missing_occupancy_mapping = [
            temp_var
            for temp_var in temp_mapping
            if temp_var not in occupancy_mapping
        ]
        if missing_occupancy_mapping:
            validation_errors.append(
                "No s'han pogut inferir variables d'ocupacio per a totes les zones: "
                + ", ".join(missing_occupancy_mapping)
            )
    energy_cost_wrapper_reward_kwargs = build_energy_cost_wrapper_reward_kwargs(
        options,
        temp_vars,
        energy_vars,
    )
    recalculate_energy_cost_reward = options["reward_name"] in {
        "LinearReward",
        "EnergyCostLinearReward",
    }
    # El wrapper de cost pot recalcular reward només en rewards lineals; en les altres,
    # el cost ja entra dins la pròpia classe de reward.
    wrapper_configs = build_wrapper_configs(
        options,
        energy_cost_reward_kwargs=energy_cost_wrapper_reward_kwargs,
        recalculate_energy_cost_reward=recalculate_energy_cost_reward,
    )
    wrapper_rows = build_wrapper_summary_rows(wrapper_configs)
    algo_kwargs, dropped_algo_kwargs = build_algo_kwargs(options)

    # Aquest dict és el que es desa amb l'artefacte. Ha de ser suficient per repetir,
    # avaluar o inspeccionar l'entrenament sense dependre de l'estat viu de Streamlit.
    training_config = {
        "env_id": options["env_id"],
        "algo_name": options["algo_name"],
        "policy_name": options["policy_name"],
        "reward_name": options["reward_name"],
        "reward_kwargs": reward_kwargs,
        "variables": variables,
        "meters": meters,
        "temp_vars": temp_vars,
        "energy_vars": energy_vars,
        "cost_vars": cost_vars,
        "wrapper_configs": wrapper_configs,
        "wrapper_rows": wrapper_rows,
        "learning_rate": options["learning_rate"],
        "n_steps": options["n_steps"],
        "batch_size": options.get("batch_size"),
        "n_epochs": options.get("n_epochs"),
        "gamma": options.get("gamma"),
        "gae_lambda": options.get("gae_lambda"),
        "clip_range": options.get("clip_range"),
        "ent_coef": options.get("ent_coef"),
        "vf_coef": options.get("vf_coef"),
        "max_grad_norm": options.get("max_grad_norm"),
        "algo_kwargs": algo_kwargs,
        "timesteps_per_year": options["timesteps_per_year"],
    }

    return {
        "dropped_reward_kwargs": dropped_reward_kwargs,
        "dropped_algo_kwargs": dropped_algo_kwargs,
        "ui_state": build_training_ui_state(options),
        "validation_errors": validation_errors,
        "training_config": training_config,
    }
\end{lstlisting}

\Needspace{8\baselineskip}

\subsection{\texorpdfstring{\texttt{backend/entrenar\_agent\_runtime.py}}{backend/entrenar\_agent\_runtime.py}}

\begin{lstlisting}[style=studioPython,caption={backend/entrenar\_agent\_runtime.py},label={lst:backend-entrenar-agent-runtime-py}]
"""Cicle de vida d'execució per a sessions d'entrenament en directe."""

from __future__ import annotations

import copy
import json
from datetime import datetime
from functools import partial
from pathlib import Path
from typing import Any, Dict, Sequence

import gymnasium as gym
import streamlit as st
from sinergym.utils.wrappers import CSVLogger, LoggerWrapper

from backend.entrenar_agent_artifacts import (
    append_detailed_meter_kwh_columns_in_workspace,
    build_training_artifact_paths,
    get_runtime_workspace_path,
)
from backend.entrenar_agent_constants import REWARD_CLASSES, TRAINING_RUNTIME_KEY
from backend.entrenar_agent_env import get_training_meters_for_env, with_detailed_hvac_meters
from backend.entrenar_agent_wrappers import apply_training_wrappers
from backend.sb3_utils import ALGOS, build_vecnormalize_env
from backend.training_scripts import build_training_eval_script, build_training_repro_script


def clear_training_runtime() -> None:
    """Neteja el temps d'execució de l'entrenament."""
    runtime = st.session_state.pop(TRAINING_RUNTIME_KEY, None)
    if not runtime:
        return

    workspace_path = get_runtime_workspace_path(runtime)
    vec = runtime.get("vec")
    if vec is not None:
        try:
            vec.envs[0].close()
        except Exception:
            try:
                vec.close()
            except Exception:
                pass

    append_detailed_meter_kwh_columns_in_workspace(workspace_path)


def learn_training_chunk(model: Any, timesteps: int) -> None:
    """Executa un bloc d'entrenament amb la inicialització incremental pròpia de SB3."""
    model.learn(total_timesteps=int(timesteps), reset_num_timesteps=False)


def save_training_artifacts(runtime: Dict[str, Any]) -> Dict[str, Any]:
    """Deseu els artefactes d'entrenament."""
    config = runtime["config"]
    artifact_dir = Path(runtime["artifact_dir"])
    artifact_dir.mkdir(parents=True, exist_ok=False)

    runtime["model"].save(str(runtime["save_path"]))
    runtime["vec"].save(str(runtime["vecnorm_path"]))

    result = {
        "artifact_name": runtime["artifact_name"],
        "created_at": datetime.now().isoformat(timespec="seconds"),
        "artifact_dir": str(artifact_dir),
        "env_id": config["env_id"],
        "reward_name": config["reward_name"],
        "algo_name": config["algo_name"],
        "policy_name": config["policy_name"],
        "files": {
            "training_script": Path(runtime["training_script_path"]).name,
            "eval_script": Path(runtime["eval_script_path"]).name,
            "model": Path(runtime["save_path"]).name,
            "vecnorm": Path(runtime["vecnorm_path"]).name,
            "config": Path(runtime["config_path"]).name,
        },
        "ui_state": runtime.get("ui_state", {}),
        "training_config": config,
    }

    with open(runtime["config_path"], "w", encoding="utf-8") as handle:
        json.dump(result, handle, indent=2, ensure_ascii=False)

    with open(runtime["eval_script_path"], "w", encoding="utf-8") as handle:
        handle.write(build_training_eval_script(config))

    with open(runtime["training_script_path"], "w", encoding="utf-8") as handle:
        handle.write(build_training_repro_script(runtime["artifact_name"], config))

    return result


def build_training_env_instance(
    env_id: str,
    reward_name: str,
    reward_kwargs: Dict[str, Any],
    meters: Dict[str, str],
    wrapper_configs: Sequence[Dict[str, Any]],
) -> Any:
    """Crea una instància d'entorn embolicat per a l'entrenament en directe."""
    reward_cls = REWARD_CLASSES[reward_name]
    env = gym.make(env_id, meters=meters, reward=reward_cls, reward_kwargs=reward_kwargs)
    env = apply_training_wrappers(env, wrapper_configs)
    env = LoggerWrapper(env)
    env = CSVLogger(env)
    return env


def create_training_runtime(
    training_config: Dict[str, Any],
    ui_state: Dict[str, Any] | None = None,
) -> Dict[str, Any]:
    """Crea temps d'execució de la entrenament."""
    env_id = training_config["env_id"]
    reward_name = training_config["reward_name"]
    reward_kwargs = training_config["reward_kwargs"]
    meters = with_detailed_hvac_meters(training_config.get("meters") or get_training_meters_for_env(env_id))
    wrapper_configs = training_config["wrapper_configs"]
    algo_name = training_config["algo_name"]
    policy_name = training_config["policy_name"]
    learning_rate = training_config["learning_rate"]
    n_steps = training_config["n_steps"]
    timesteps_per_year = training_config["timesteps_per_year"]

    env_factory = partial(
        build_training_env_instance,
        env_id,
        reward_name,
        reward_kwargs,
        meters,
        wrapper_configs,
    )
    vec = build_vecnormalize_env(env_factory)

    model_cls = ALGOS[algo_name]
    algo_kwargs = dict(training_config.get("algo_kwargs") or {})
    algo_kwargs.setdefault("learning_rate", learning_rate)
    if algo_name in ["PPO", "A2C"] and n_steps is not None:
        algo_kwargs.setdefault("n_steps", n_steps)
    algo_kwargs["verbose"] = 0

    model = model_cls(policy_name, vec, **algo_kwargs)
    artifact_paths = build_training_artifact_paths(training_config)

    return {
        "model": model,
        "vec": vec,
        "config": training_config,
        "ui_state": copy.deepcopy(ui_state or {}),
        "artifact_name": artifact_paths["artifact_name"],
        "artifact_dir": artifact_paths["artifact_dir"],
        "save_path": artifact_paths["model_path"],
        "vecnorm_path": artifact_paths["vecnorm_path"],
        "eval_script_path": artifact_paths["eval_script_path"],
        "training_script_path": artifact_paths["training_script_path"],
        "config_path": artifact_paths["config_path"],
        "chunk_timesteps": int(algo_kwargs.get("n_steps") or n_steps) if algo_name in ["PPO", "A2C"] else min(512, timesteps_per_year),
    }
\end{lstlisting}

\Needspace{8\baselineskip}

\subsection{\texorpdfstring{\texttt{backend/entrenar\_agent\_session.py}}{backend/entrenar\_agent\_session.py}}

\begin{lstlisting}[style=studioPython,caption={backend/entrenar\_agent\_session.py},label={lst:backend-entrenar-agent-session-py}]
"""Gestió de l'estat de sessió per a la pàgina d'entrenament de l'agent."""

from __future__ import annotations

import streamlit as st

from backend.entrenar_agent_constants import TRAINING_RUNTIME_KEY
from backend.entrenar_agent_runtime import clear_training_runtime

BATTERY_REWARDS = {
    "BatteryHVACReward",
    "MultiZoneBatteryHVACReward",
    "MultizoneEnergyCostBatteryHVACReward",
}
TRAINING_WORKSPACE_KEY = "training_live_workspace_path"
TRAINING_FORM_PREFIX = "training_form_"
TRAINING_SAVED_RUN_KEY = "training_saved_run"
TRAINING_LOADED_ARTIFACT_KEY = "training_loaded_artifact"
TRAINING_RANGE_FIELDS = {
    "range_winter",
    "range_summer",
    "range_comfort_hours",
    "occupied_hours",
}


def training_form_key(field: str) -> str:
    """Retorna la clau d'estat de sessió per a un camp de formulari d'entrenament."""
    return f"{TRAINING_FORM_PREFIX}{field}"


def reset_widget_state_if_disabled(field: str, disabled: bool, value=False) -> None:
    """Restableix l'estat de sessió d'un widget quan està desactivat."""
    if disabled:
        st.session_state[training_form_key(field)] = value


def ensure_select_state(field: str, options: list[str], fallback: str | None = None) -> str:
    """Assegureu-vos que un camp de casella de selecció tingui un valor vàlid en estat de sessió.

    Si el valor actual no es troba a les opcions, se substitueix per una alternativa (si és vàlid)
    o per la primera opció.
    """
    key = training_form_key(field)
    if not options:
        st.session_state[key] = ""
        return ""

    current_value = st.session_state.get(key)
    if current_value not in options:
        st.session_state[key] = fallback if fallback in options else options[0]
    return st.session_state[key]


def selectbox_state_kwargs(field: str, options: list, fallback=None) -> dict:
    """Retorna kwargs per a selectbox sense duplicar valor per defecte i session_state."""
    key = training_form_key(field)
    current_value = st.session_state.get(key)
    if current_value in options:
        return {"key": key}

    if key in st.session_state:
        st.session_state.pop(key, None)

    if fallback in options:
        index = options.index(fallback)
    else:
        index = 0
    return {"index": index, "key": key}


def sanitize_multiselect_state(
    field: str,
    options: list[str],
    fallback: list[str] | None = None,
) -> None:
    """Elimina valors obsolets d'un camp de selecció múltiple en estat de sessió.

    Els valors que ja no estan presents a les opcions s'eliminen. Si el resultat és
    buit i es proporciona una alternativa, s'utilitzen elements de reserva vàlids.
    """
    key = training_form_key(field)
    current_value = st.session_state.get(key)
    if current_value is None:
        return

    if not isinstance(current_value, (list, tuple, set)):
        current_value = fallback or []

    sanitized = [item for item in current_value if item in options]
    if not sanitized and fallback is not None:
        sanitized = [item for item in fallback if item in options]
    st.session_state[key] = sanitized


def normalize_training_range_state() -> None:
    """Converteix els camps d'interval codificats per llista en tuples en estat de sessió.

    Streamlit serialitza les tuples lliscants com a llistes quan es restaura l'estat de la sessió
    d'una carrera guardada; aquesta funció els torna a convertir en tuples de manera que
    ``st.slider`` no genera cap desajust de tipus.
    """
    for field in TRAINING_RANGE_FIELDS:
        key = training_form_key(field)
        value = st.session_state.get(key)
        if isinstance(value, list) and len(value) == 2:
            st.session_state[key] = tuple(value)


def clear_incremental_widget_state() -> None:
    """Elimina totes les claus del widget incremental dinàmic de l'estat de sessió.

    Les claus s'identifiquen pel seu patró de prefix, que és generat per
    ``training_form_key`` per a cada widget incremental per variable.
    """
    dynamic_prefixes = (
        training_form_key("incr_init_"),
        training_form_key("incr_delta_"),
        training_form_key("incr_step_"),
        training_form_key("disc_incr_init_"),
    )
    for session_key in list(st.session_state.keys()):
        if session_key.startswith(dynamic_prefixes):
            st.session_state.pop(session_key, None)


def apply_saved_training_ui_state(ui_state: dict, envs: list[str]) -> None:
    """Restaura un estat d'entrenament desat anteriorment UI a l'estat de sessió.

    Primer esborra totes les claus del widget incrementals i després escriu cada camp des de
    ``ui_state``, converteix els camps d'interval en tuples i restableix l'entrenament
    temps d'execució perquè la configuració carregada tingui efecte a la següent execució.
    """
    clear_incremental_widget_state()
    for field, value in (ui_state or {}).items():
        normalized_value = tuple(value) if field in TRAINING_RANGE_FIELDS and isinstance(value, list) else value
        st.session_state[training_form_key(field)] = normalized_value

    ensure_select_state("env_id", envs)
    clear_training_runtime()
    st.session_state.is_training = False
    st.session_state.stop_training = False
    st.session_state.pop(TRAINING_RUNTIME_KEY, None)
    st.session_state.pop(TRAINING_WORKSPACE_KEY, None)


def seed_incremental_widget_defaults(selected_variables: list[str]) -> None:
    """Emplena els widgets incrementals per variable des de l'estat desat.

    Només escriu a claus d'estat de sessió que encara no existeixen, de manera que
    Els valors dels widgets en directe mai es sobreescriuen en tornar a renderitzar.
    """
    stored_variables = st.session_state.get(training_form_key("incremental_variables"), []) or []
    stored_initial_values = st.session_state.get(training_form_key("incremental_initial_values"), []) or []
    stored_definition = st.session_state.get(training_form_key("incremental_definition"), {}) or {}
    initial_values_map = {
        variable_name: stored_initial_values[index]
        for index, variable_name in enumerate(stored_variables)
        if index < len(stored_initial_values)
    }

    for variable_name in selected_variables:
        init_key = training_form_key(f"incr_init_{variable_name}")
        delta_key = training_form_key(f"incr_delta_{variable_name}")
        step_key = training_form_key(f"incr_step_{variable_name}")
        if init_key not in st.session_state and variable_name in initial_values_map:
            st.session_state[init_key] = initial_values_map[variable_name]

        definition_values = stored_definition.get(variable_name)
        if isinstance(definition_values, (list, tuple)) and len(definition_values) == 2:
            if delta_key not in st.session_state:
                st.session_state[delta_key] = definition_values[0]
            if step_key not in st.session_state:
                st.session_state[step_key] = definition_values[1]


def seed_discrete_incremental_widget_defaults(action_variables: list[str]) -> None:
    """Emplena els widgets incrementals discrets per variable des de l'estat desat.

    Només escriu a claus d'estat de sessió que encara no existeixen.
    """
    stored_initial_values = st.session_state.get(
        training_form_key("discrete_incremental_initial_values"),
        [],
    ) or []
    for action_index, variable_name in enumerate(action_variables):
        widget_key = training_form_key(f"disc_incr_init_{variable_name}")
        if widget_key not in st.session_state and action_index < len(stored_initial_values):
            st.session_state[widget_key] = stored_initial_values[action_index]
\end{lstlisting}

\Needspace{8\baselineskip}

\subsection{\texorpdfstring{\texttt{backend/entrenar\_agent\_wrappers.py}}{backend/entrenar\_agent\_wrappers.py}}

\begin{lstlisting}[style=studioPython,caption={backend/entrenar\_agent\_wrappers.py},label={lst:backend-entrenar-agent-wrappers-py}]
"""Configuració i aplicació de wrapper per a entorns d'entrenament."""

from __future__ import annotations

import csv
from typing import Any, Dict, List, Sequence

import numpy as np
from gymnasium import Wrapper
from gymnasium.spaces import Box

from backend.entrenar_agent_constants import FIXED_WRAPPER_ROWS
from backend.common import format_ui_value


class EnergyCostFileLogger(Wrapper):
    """Logger addicional (opcional) que escriu cost per pas en un CSV independent."""

    def __init__(self, env, filename="energy_cost_log.csv"):
        """Inicialitza la instància."""
        super().__init__(env)
        self.csvfile = open(filename, mode="w", newline="")
        self.writer = csv.writer(self.csvfile)
        self.writer.writerow(["step", "elèctricity_cost_eur_per_kwh", "energy_cost_eur"])
        self.step_count = 0

    def step(self, action):
        """Step."""
        obs, reward, terminated, truncated, info = self.env.step(action)
        self.step_count += 1
        cost_per_kwh = info.get("elèctricity_cost", 0.0)
        energy_rate = info.get("HVAC_elèctricity_demand_rate", 0.0)
        cost_total = cost_per_kwh * energy_rate
        self.writer.writerow([self.step_count, cost_per_kwh, cost_total])
        return obs, reward, terminated, truncated, info

    def reset(self, **kwargs):
        """Restableix."""
        self.step_count = 0
        return self.env.reset(**kwargs)

    def close(self):
        """Close."""
        self.csvfile.close()
        if hasattr(self.env, "close"):
            self.env.close()


def wrapper_details(params: Dict[str, Any]) -> str:
    """Detalls del wrapper."""
    if not params:
        return "-"
    return "; ".join(f"{key}={format_ui_value(value)}" for key, value in params.items())


def build_wrapper_summary_rows(wrapper_configs: Sequence[Dict[str, Any]]) -> List[Dict[str, str]]:
    """Prepara les files de resum dels wrappers seleccionats."""
    rows = [dict(row) for row in FIXED_WRAPPER_ROWS]
    for config in wrapper_configs:
        rows.append(
            {
                "Wrapper": config["name"],
                "Actiu": "Si",
                "Detall": wrapper_details(config.get("params", {})),
            }
        )
    return rows


def apply_training_wrappers(env: Any, wrapper_configs: Sequence[Dict[str, Any]]) -> Any:
    """Aplica wrappers d'entrenament."""
    # L'ordre ve guardat a la configuració i s'ha de respectar igual en training,
    # avaluació i interacció. Si es canvia, canvia també l'espai d'observació/acció.
    for wrapper_config in wrapper_configs:
        name = wrapper_config["name"]
        params = _normalize_wrapper_params(name, wrapper_config.get("params", {}))
        if name == "NormalizeAction":
            # NormalizeAction només té sentit en espais continus; així evitem embolicar
            # un Discrete per accident quan es reutilitza una configuració antiga.
            if isinstance(env.action_space, Box):
                from sinergym.utils.wrappers import NormalizeAction

                env = NormalizeAction(env, **params)
            continue
        from sinergym.utils.common import apply_wrappers_info

        env = apply_wrappers_info(env, {_wrapper_import_path(name): params})
    return env


def _normalize_wrapper_params(name: str, params: Dict[str, Any]) -> Dict[str, Any]:
    """Converteix paràmetres serialitzats per la UI al format que espera Sinergym."""
    normalized = dict(params or {})
    if name == "VariabilityContextWrapper" and "context_space" not in normalized:
        normalized["context_space"] = Box(
            low=np.array(normalized.pop("context_low"), dtype=np.float32),
            high=np.array(normalized.pop("context_high"), dtype=np.float32),
            dtype=np.float32,
        )
        normalized["step_frequency_range"] = tuple(normalized["step_frequency_range"])
    return normalized


def _wrapper_import_path(name: str) -> str:
    """Retorna la ruta importable que consumeix ``sinergym.utils.common.apply_wrappers_info``."""
    if name == "EnergyCostFileLogger":
        return "backend.entrenar_agent_wrappers:EnergyCostFileLogger"
    return f"sinergym.utils.wrappers:{name}"


def build_energy_cost_wrapper_reward_kwargs(
    options: Dict[str, Any],
    temp_vars: Sequence[str],
    energy_vars: Sequence[str],
) -> Dict[str, Any]:
    """Munta els kwargs del wrapper de costos energètics."""
    return {
        "temperature_variables": list(temp_vars),
        "energy_variables": list(energy_vars),
        "energy_cost_variables": ["energy_cost"],
        "range_comfort_winter": tuple(options["range_winter"]),
        "range_comfort_summer": tuple(options["range_summer"]),
        "energy_weight": options["energy_weight"],
        "temperature_weight": options["temperature_weight"],
        "lambda_energy": options["lambda_energy"],
        "lambda_temperature": options["lambda_temperature"],
        "lambda_energy_cost": options["lambda_energy_cost"],
    }


def build_wrapper_configs(
    options: Dict[str, Any],
    energy_cost_reward_kwargs: Dict[str, Any] | None = None,
    recalculate_energy_cost_reward: bool = True,
) -> List[Dict[str, Any]]:
    """Munta la configuració dels wrappers d'entrenament."""
    # Aquesta llista és el contracte que després es desa amb el model. Per això guardem
    # noms i paràmetres simples, fàcils de reconstruir en avaluació.
    wrapper_configs: List[Dict[str, Any]] = []
    if options["enable_previous_wrapper"]:
        wrapper_configs.append(
            {
                "name": "PreviousObservationWrapper",
                "params": {"previous_variables": options["previous_variables"]},
            }
        )
    if options["enable_datetime_wrapper"]:
        wrapper_configs.append({"name": "DatetimeWrapper", "params": {}})
    if options["enable_weather_forecast_wrapper"]:
        wrapper_configs.append(
            {
                "name": "WeatherForecastingWrapper",
                "params": {
                    "n": options["weather_forecast_n"],
                    "delta": options["weather_forecast_delta"],
                    "columns": options["weather_forecast_columns"],
                },
            }
        )
    if options["use_energy_cost"]:
        energy_cost_params = {
            "energy_cost_data_path": options["energy_cost_path"],
            "recalculate_reward": recalculate_energy_cost_reward,
        }
        if energy_cost_reward_kwargs is not None:
            energy_cost_params["reward_kwargs"] = energy_cost_reward_kwargs
        wrapper_configs.append(
            {
                "name": "EnergyCostWrapper",
                "params": energy_cost_params,
            }
        )
    if options["enable_delta_temp_wrapper"]:
        wrapper_configs.append(
            {
                "name": "DeltaTempWrapper",
                "params": {
                    "temperature_variables": options["delta_temp_variables"],
                    "setpoint_variables": options["delta_setpoint_variables"],
                },
            }
        )
    if options["enable_reduce_obs_wrapper"]:
        wrapper_configs.append(
            {
                "name": "ReduceObservationWrapper",
                "params": {"obs_reduction": options["reduced_observations"]},
            }
        )
    if options["enable_multi_obs_wrapper"]:
        wrapper_configs.append(
            {
                "name": "MultiObsWrapper",
                "params": {"n": options["multi_obs_n"], "flatten": options["multi_obs_flatten"]},
            }
        )
    if options["enable_normalize_obs_wrapper"]:
        wrapper_configs.append(
            {
                "name": "NormalizeObservation",
                "params": {
                    "automatic_update": options["normalize_obs_auto"],
                    "epsilon": options["normalize_obs_epsilon"],
                    "mean": options["normalize_obs_mean"] or None,
                    "var": options["normalize_obs_var"] or None,
                },
            }
        )
    if options.get("enable_variability_context_wrapper"):
        # Els límits del context es guarden com a llistes perquè siguin serialitzables.
        # El Box es reconstrueix quan s'apliquen els wrappers.
        wrapper_configs.append(
            {
                "name": "VariabilityContextWrapper",
                "params": {
                    "context_low": list(options.get("variability_context_low") or []),
                    "context_high": list(options.get("variability_context_high") or []),
                    "delta_value": options.get("variability_delta_value", 1.0),
                    "step_frequency_range": (
                        options.get("variability_step_frequency_min", 96),
                        options.get("variability_step_frequency_max", 96 * 7),
                    ),
                },
            }
        )
    if options["enable_incremental_wrapper"]:
        wrapper_configs.append(
            {
                "name": "IncrementalWrapper",
                "params": {
                    "incremental_variables_definition": options["incremental_definition"],
                    "initial_values": options["incremental_initial_values"],
                },
            }
        )
    if options["enable_discrete_incremental_wrapper"]:
        wrapper_configs.append(
            {
                "name": "DiscreteIncrementalWrapper",
                "params": {
                    "initial_values": options["discrete_incremental_initial_values"],
                    "delta_temp": options["discrete_incremental_delta"],
                    "step_temp": options["discrete_incremental_step"],
                },
            }
        )
    if (
        options["enable_normalize_action"]
        and options.get("is_continuous", True)
        and not options["enable_incremental_wrapper"]
        and not options["enable_discrete_incremental_wrapper"]
    ):
        # Els wrappers incrementals ja redefineixen l'escala de l'acció; combinar-los amb
        # NormalizeAction fa que l'agent aprengui en una escala difícil d'interpretar.
        wrapper_configs.append(
            {
                "name": "NormalizeAction",
                "params": {"normalize_range": (-1.0, 1.0)},
            }
        )
    if options.get("enable_storage_smoothing_wrapper"):
        wrapper_configs.append(
            {
                "name": "OfficeGridStorageSmoothingActionConstraintsWrapper",
                "params": {},
            }
        )
    if options["use_file_logger"]:
        wrapper_configs.append(
            {
                "name": "EnergyCostFileLogger",
                "params": {"filename": options["file_logger_name"]},
            }
        )
    return wrapper_configs
\end{lstlisting}

\Needspace{8\baselineskip}

\subsection{\texorpdfstring{\texttt{backend/epw\_figures.py}}{backend/epw\_figures.py}}

\begin{lstlisting}[style=studioPython,caption={backend/epw\_figures.py},label={lst:backend-epw-figures-py}]
"""Plotly creadors de figures per al visor de clima EPW."""

from __future__ import annotations

import math

import pandas as pd
import plotly.graph_objects as go
from plotly.subplots import make_subplots

from backend.visor_epw_backend import (
    DEG,
    DEG_C,
    EPW_VARIABLES,
    KWH_PER_M2,
    MID_DOT,
    MONTH_ORDER,
    WH_PER_M2,
    WIND_DIRECTION_LABELS,
    variable_label,
)

MONTH_TICK_POSITIONS = [1, 32, 60, 91, 121, 152, 182, 213, 244, 274, 305, 335]


def ui_text(value: object) -> str:
    """Retorna el text segur per a la pantalla, reparant el mojibake habitual de les metadades EPW."""

    text = "" if value is None else str(value)
    for _ in range(2):
        if not any(marker in text for marker in ("Ã", "Â", "â")):
            break
        try:
            repaired = text.encode("latin-1").decode("utf-8")
        except (UnicodeEncodeError, UnicodeDecodeError):
            break
        if repaired == text:
            break
        text = repaired
    return text


def build_active_month_axis(data_frame: pd.DataFrame) -> tuple[list[int], list[str]]:
    """Retorna les posicions de marca del mes que coincideixen amb les dades EPW visibles actualment."""

    month_frame = (
        data_frame[["month", "month_label", "day_of_year"]]
        .dropna(subset=["month", "day_of_year"])
        .sort_values(["month", "day_of_year"])
        .drop_duplicates(subset=["month"])
    )
    if month_frame.empty:
        return MONTH_TICK_POSITIONS, MONTH_ORDER

    tick_values = month_frame["day_of_year"].astype(int).tolist()
    tick_labels = month_frame["month_label"].astype(str).tolist()
    return tick_values, tick_labels


def apply_figure_style(fig: go.Figure, title: str, *, legend: bool = True, height: int = 400) -> go.Figure:
    """Aplica l'estil visual compartit a una figura Plotly."""

    top_margin = 82 if legend else 58
    fig.update_layout(
        title=dict(
            text=ui_text(title),
            x=0.02,
            xanchor="left",
            y=0.98,
            yanchor="top",
            font=dict(size=15),
        ),
        height=height,
        margin=dict(l=12, r=12, t=top_margin, b=18),
        paper_bgcolor="rgba(0,0,0,0)",
        plot_bgcolor="rgba(247,249,254,0.75)",
        hovermode="x unified",
        legend=dict(
            orientation="h",
            yanchor="bottom",
            y=1.02,
            xanchor="right",
            x=1,
            font=dict(size=12),
            bgcolor="rgba(255,255,255,0.78)",
        ),
        showlegend=legend,
        font=dict(family="Bahnschrift, Aptos, Segoe UI, sans-serif", color="#17233c"),
    )
    fig.update_xaxes(showgrid=False, zeroline=False)
    fig.update_yaxes(gridcolor="rgba(204, 214, 235, 0.55)", zeroline=False)
    return fig


def build_focus_timeseries_figure(series_frame: pd.DataFrame, variable_key: str, aggregation_label: str) -> go.Figure:
    """Crea la sèrie temporal agregada per a la variable EPW seleccionada."""

    fig = go.Figure()
    fig.add_trace(
        go.Scatter(
            x=series_frame["timestamp"],
            y=series_frame[variable_key],
            mode="lines",
            name=ui_text(variable_label(variable_key, aggregated=aggregation_label != "Hora")),
            line=dict(color=EPW_VARIABLES[variable_key]["color"], width=2.3),
            fill="tozeroy" if EPW_VARIABLES[variable_key]["aggregate"] == "sum" else None,
            fillcolor="rgba(95, 84, 249, 0.10)" if EPW_VARIABLES[variable_key]["aggregate"] == "sum" else None,
            connectgaps=False,
        )
    )
    fig.update_xaxes(title_text="Calendari EPW")
    return apply_figure_style(fig, f"S\u00e8rie temporal {MID_DOT} {aggregation_label}", height=390)


def build_monthly_climate_figure(monthly_frame: pd.DataFrame) -> go.Figure:
    """Crea la comparació mensual de temperatura i humitat."""

    fig = make_subplots(specs=[[{"secondary_y": True}]])
    fig.add_trace(
        go.Scatter(
            x=monthly_frame["month_label"],
            y=monthly_frame["dry_bulb_max"],
            mode="lines",
            line=dict(color="rgba(242, 107, 91, 0.18)", width=0.1),
            showlegend=False,
            hoverinfo="skip",
        ),
        secondary_y=False,
    )
    fig.add_trace(
        go.Scatter(
            x=monthly_frame["month_label"],
            y=monthly_frame["dry_bulb_min"],
            mode="lines",
            line=dict(color="rgba(242, 107, 91, 0.18)", width=0.1),
            fill="tonexty",
            fillcolor="rgba(242, 107, 91, 0.14)",
            name="Rang",
        ),
        secondary_y=False,
    )
    fig.add_trace(
        go.Scatter(
            x=monthly_frame["month_label"],
            y=monthly_frame["dry_bulb_mean"],
            mode="lines+markers",
            name="Temp. seca",
            line=dict(color="#f26b5b", width=2.5),
        ),
        secondary_y=False,
    )
    fig.add_trace(
        go.Bar(
            x=monthly_frame["month_label"],
            y=monthly_frame["relative_humidity_mean"],
            name="Humitat",
            marker_color="rgba(79, 142, 247, 0.65)",
        ),
        secondary_y=True,
    )
    apply_figure_style(fig, "Climatologia mensual", height=380)
    fig.update_yaxes(title_text=f"Temperatura ({DEG_C})", secondary_y=False)
    fig.update_yaxes(title_text="Humitat relativa (%)", secondary_y=True)
    return fig


def build_monthly_solar_figure(monthly_frame: pd.DataFrame) -> go.Figure:
    """Crea el resum mensual de radiació solar acumulada."""

    fig = go.Figure()
    fig.add_trace(
        go.Bar(
            x=monthly_frame["month_label"],
            y=monthly_frame["global_horizontal_radiation_kwh_m2"],
            name="GHI",
            marker_color="#f3a738",
        )
    )
    fig.add_trace(
        go.Bar(
            x=monthly_frame["month_label"],
            y=monthly_frame["direct_normal_radiation_kwh_m2"],
            name="DNI",
            marker_color="#ef7d57",
        )
    )
    fig.add_trace(
        go.Bar(
            x=monthly_frame["month_label"],
            y=monthly_frame["diffuse_horizontal_radiation_kwh_m2"],
            name="DHI",
            marker_color="#f9cb40",
        )
    )
    fig.update_layout(barmode="group")
    fig.update_yaxes(title_text=KWH_PER_M2)
    return apply_figure_style(fig, "Radiaci\u00f3 solar mensual", height=380)


def build_daily_temperature_band_figure(daily_frame: pd.DataFrame) -> go.Figure:
    """Crea la banda diària de temperatura amb mínims, mitjanes i màxims."""

    fig = go.Figure()
    fig.add_trace(
        go.Scatter(
            x=daily_frame["date"],
            y=daily_frame["dry_bulb_max"],
            mode="lines",
            line=dict(color="rgba(242, 107, 91, 0.16)", width=0.1),
            showlegend=False,
            hoverinfo="skip",
        )
    )
    fig.add_trace(
        go.Scatter(
            x=daily_frame["date"],
            y=daily_frame["dry_bulb_min"],
            mode="lines",
            line=dict(color="rgba(79, 142, 247, 0.16)", width=0.1),
            fill="tonexty",
            fillcolor="rgba(95, 84, 249, 0.16)",
            name="Rang diari",
        )
    )
    fig.add_trace(
        go.Scatter(
            x=daily_frame["date"],
            y=daily_frame["dry_bulb_mean"],
            mode="lines",
            name="Mitjana",
            line=dict(color="#5f54f9", width=2),
        )
    )
    fig.update_yaxes(title_text=DEG_C)
    return apply_figure_style(fig, "Banda di\u00e0ria de temperatura", height=360)


def build_hourly_profile_figure(hourly_frame: pd.DataFrame) -> go.Figure:
    """Crea el perfil horari mitjà de temperatura, punt de rosada i humitat."""

    fig = make_subplots(specs=[[{"secondary_y": True}]])
    fig.add_trace(
        go.Scatter(
            x=hourly_frame["hour_start"],
            y=hourly_frame["dry_bulb_mean"],
            mode="lines+markers",
            name="Temp.",
            line=dict(color="#f26b5b", width=2.4),
        ),
        secondary_y=False,
    )
    fig.add_trace(
        go.Scatter(
            x=hourly_frame["hour_start"],
            y=hourly_frame["dew_point_mean"],
            mode="lines",
            name="Rosada",
            line=dict(color="#ffb347", width=2, dash="dot"),
        ),
        secondary_y=False,
    )
    fig.add_trace(
        go.Scatter(
            x=hourly_frame["hour_start"],
            y=hourly_frame["relative_humidity_mean"],
            mode="lines",
            name="Humitat",
            line=dict(color="#4f8ef7", width=2.2),
        ),
        secondary_y=True,
    )
    apply_figure_style(fig, "Perfil horari mitj\u00e0", height=360)
    fig.update_xaxes(title_text="Hora del dia")
    fig.update_yaxes(title_text=f"Temperatura ({DEG_C})", secondary_y=False)
    fig.update_yaxes(title_text="Humitat relativa (%)", secondary_y=True)
    return fig


def build_heatmap_figure(heatmap_frame: pd.DataFrame, variable_key: str) -> go.Figure:
    """Crea un mapa de calor mes a hora per a la variable EPW seleccionada."""

    fig = go.Figure(
        data=[
            go.Heatmap(
                z=heatmap_frame.values,
                x=[int(hour) for hour in heatmap_frame.columns],
                y=list(heatmap_frame.index),
                colorscale="Blues" if variable_key == "relative_humidity_pct" else "Turbo",
                colorbar=dict(title=ui_text(variable_label(variable_key))),
                hovertemplate="Mes %{y}<br>Hora %{x}:00<br>Valor %{z:.2f}<extra></extra>",
            )
        ]
    )
    fig.update_xaxes(title_text="Hora del dia")
    fig.update_yaxes(title_text="Mes")
    return apply_figure_style(
        fig,
        f"Mapa de calor {MID_DOT} {ui_text(EPW_VARIABLES[variable_key]['label'])}",
        legend=False,
        height=390,
    )


def build_annual_heatmap_figure(
    heatmap_frame: pd.DataFrame,
    variable_key: str,
    *,
    tick_values: list[int] | None = None,
    tick_labels: list[str] | None = None,
) -> go.Figure:
    """Crea un mapa de calor anual hora a dia per a l'escaneig a escala del panell."""

    title_map = {
        "dry_bulb_c": f"Temperature ({DEG_C})",
        "direct_normal_radiation_wh_m2": f"Direct Radiation ({WH_PER_M2})",
        "wind_speed_m_s": "Wind Speed (m/s)",
        "relative_humidity_pct": "Relative Humidity (%)",
    }
    color_scale_map = {
        "dry_bulb_c": "RdBu_r",
        "direct_normal_radiation_wh_m2": "YlOrRd",
        "wind_speed_m_s": "Greens",
        "relative_humidity_pct": "Blues",
    }

    # Redueix l'escala a quantils robusts perquè els valors atípics aïllats no aplanin el mapa de calor.
    numeric_values = pd.DataFrame(heatmap_frame).to_numpy(dtype=float)
    flattened = pd.Series(numeric_values.reshape(-1)).dropna()
    if flattened.empty:
        return go.Figure()

    zmin = float(flattened.quantile(0.02))
    zmax = float(flattened.quantile(0.98))
    if variable_key == "dry_bulb_c":
        bound = max(abs(zmin), abs(zmax), 1.0)
        zmin = -bound
        zmax = bound
    elif zmin == zmax:
        zmin -= 1.0
        zmax += 1.0

    x_values = [int(column) for column in heatmap_frame.columns]
    fig = go.Figure(
        data=[
            go.Heatmap(
                z=heatmap_frame.values,
                x=x_values,
                y=list(heatmap_frame.index),
                zmin=zmin,
                zmax=zmax,
                colorscale=color_scale_map.get(variable_key, "Turbo"),
                colorbar=dict(title=ui_text(EPW_VARIABLES[variable_key]["unit"]), len=0.78, thickness=11),
                hovertemplate="Dia %{x}<br>Hora %{y}:00<br>Valor %{z:.2f}<extra></extra>",
            )
        ]
    )
    fig.update_layout(
        title=ui_text(title_map.get(variable_key, variable_label(variable_key))),
        height=220,
        margin=dict(l=10, r=10, t=40, b=10),
        paper_bgcolor="rgba(0,0,0,0)",
        plot_bgcolor="rgba(247,249,254,0.75)",
        font=dict(family="Bahnschrift, Aptos, Segoe UI, sans-serif", color="#17233c"),
    )
    fig.update_xaxes(
        tickmode="array",
        tickvals=tick_values or MONTH_TICK_POSITIONS,
        ticktext=tick_labels or MONTH_ORDER,
        showgrid=False,
        zeroline=False,
        range=[min(x_values), max(x_values)] if x_values else None,
    )
    fig.update_yaxes(
        title_text="Hora",
        tickmode="array",
        tickvals=[0, 6, 12, 18, 23],
        showgrid=False,
        zeroline=False,
        autorange="reversed",
    )
    return fig


def build_comfort_radiation_figure(radiation_frame: pd.DataFrame) -> go.Figure:
    """Crea el mapa polar de radiació amb efectes de confort per orientació."""

    if radiation_frame.empty:
        return go.Figure()

    ordered_frame = radiation_frame.sort_values(["radius_start", "theta_deg"]).reset_index(drop=True)
    band_order = list(dict.fromkeys(ordered_frame["altitude_band"].tolist()))
    # La gamma simètrica separa els efectes positius i negatius sobre el confort.
    max_abs = max(
        abs(float(ordered_frame["comfort_radiation_kwh_m2"].min())),
        abs(float(ordered_frame["comfort_radiation_kwh_m2"].max())),
        1.0,
    )

    fig = go.Figure()
    for band_index, altitude_band in enumerate(band_order):
        # Cada banda es dibuixa com un segment polar anular a l'altitud solar corresponent.
        band_frame = ordered_frame[ordered_frame["altitude_band"] == altitude_band]
        fig.add_trace(
            go.Barpolar(
                r=band_frame["radius_size"],
                theta=band_frame["theta_deg"],
                base=band_frame["radius_start"],
                width=[22.5] * len(band_frame),
                marker=dict(
                    color=band_frame["comfort_radiation_kwh_m2"],
                    colorscale=[[0.0, "#1f9d55"], [0.5, "#f4ddaf"], [1.0, "#d94d41"]],
                    cmin=-max_abs,
                    cmax=max_abs,
                    showscale=band_index == len(band_order) - 1,
                    colorbar=dict(title=KWH_PER_M2, len=0.78, thickness=11),
                ),
                customdata=list(
                    zip(
                        band_frame["altitude_band"],
                        band_frame["comfort_radiation_kwh_m2"],
                        band_frame["incident_radiation_kwh_m2"],
                    )
                ),
                hovertemplate=(
                    "%{theta}<br>"
                    "Banda solar %{customdata[0]}<br>"
                    f"Balanç de confort %{{customdata[1]:.1f}} {KWH_PER_M2}<br>"
                    f"Radiació captada %{{customdata[2]:.1f}} {KWH_PER_M2}<extra></extra>"
                ),
                opacity=0.94,
                showlegend=False,
            )
        )

    fig.update_layout(
        title=f"Benefit/Harm Radiation ({KWH_PER_M2})",
        height=390,
        margin=dict(l=10, r=10, t=46, b=10),
        paper_bgcolor="rgba(0,0,0,0)",
        font=dict(family="Bahnschrift, Aptos, Segoe UI, sans-serif", color="#17233c"),
        polar=dict(
            bgcolor="rgba(247,249,254,0.75)",
            radialaxis=dict(
                range=[0, 90],
                tickmode="array",
                tickvals=[15, 30, 45, 60, 75, 90],
                ticktext=[f"15{DEG}", f"30{DEG}", f"45{DEG}", f"60{DEG}", f"75{DEG}", f"90{DEG}"],
                gridcolor="rgba(204,214,235,0.45)",
            ),
            angularaxis=dict(
                direction="clockwise",
                rotation=90,
                tickmode="array",
                tickvals=[0, 45, 90, 135, 180, 225, 270, 315],
                ticktext=["N", "NE", "E", "SE", "S", "SW", "W", "NW"],
            ),
        ),
    )
    return fig


def build_monthly_wind_rose_grid_figure(monthly_rose_tables: dict[str, pd.DataFrame]) -> go.Figure:
    """Construeix una graella compacta de roses dels vents mensuals."""

    month_items = list(monthly_rose_tables.items())
    if not month_items:
        return go.Figure()

    max_frequency = max(
        (float(frame["frequency_pct"].max()) for _, frame in month_items if not frame.empty),
        default=1.0,
    )
    max_speed = max((float(frame["mean_speed"].max()) for _, frame in month_items if not frame.empty), default=1.0)
    if max_frequency <= 0:
        max_frequency = 1.0
    if max_speed <= 0:
        max_speed = 1.0

    # Fem una graella màxima de quatre columnes perquè els dotze mesos càpiguen sense
    # convertir cada rosa del vent en una miniatura il·legible.
    columns = min(4, max(1, len(month_items)))
    rows = max(1, math.ceil(len(month_items) / columns))
    specs = [[{"type": "polar"} for _ in range(columns)] for _ in range(rows)]
    subplot_titles = [month_label for month_label, _ in month_items]
    empty_slots = rows * columns - len(subplot_titles)
    if empty_slots > 0:
        subplot_titles.extend([""] * empty_slots)

    fig = make_subplots(
        rows=rows,
        cols=columns,
        specs=specs,
        subplot_titles=subplot_titles,
        horizontal_spacing=0.05,
        vertical_spacing=0.12,
    )

    for month_index, (month_label, rose_frame) in enumerate(month_items):
        row = month_index // columns + 1
        col = month_index % columns + 1
        ordered_frame = rose_frame.copy()
        # Ordenem les direccions manualment; si no, Plotly les podria ordenar alfabèticament.
        ordered_frame["direction"] = pd.Categorical(
            ordered_frame["direction"],
            categories=list(WIND_DIRECTION_LABELS),
            ordered=True,
        )
        ordered_frame = ordered_frame.sort_values("direction")

        fig.add_trace(
            go.Barpolar(
                r=ordered_frame["frequency_pct"],
                theta=ordered_frame["direction"].astype(str),
                marker=dict(
                    color=ordered_frame["mean_speed"],
                    colorscale="Turbo",
                    cmin=0.0,
                    cmax=max_speed,
                    showscale=month_index == len(month_items) - 1,
                    colorbar=dict(title="m/s", len=0.44, thickness=10, y=0.18),
                ),
                opacity=0.9,
                hovertemplate=(
                    f"{month_label}<br>"
                    "%{theta}<br>"
                    "Freq. %{r:.1f}%<br>"
                    "Vent mitjà %{marker.color:.2f} m/s<extra></extra>"
                ),
                showlegend=False,
            ),
            row=row,
            col=col,
        )

    fig.update_layout(
        title="Wind Rose",
        height=245 * rows + 105,
        margin=dict(l=10, r=10, t=70, b=12),
        paper_bgcolor="rgba(0,0,0,0)",
        font=dict(family="Bahnschrift, Aptos, Segoe UI, sans-serif", color="#17233c"),
    )
    fig.update_annotations(
        font=dict(size=11, color="#17233c"),
        yshift=14,
    )

    for polar_index in range(1, len(month_items) + 1):
        polar_name = "polar" if polar_index == 1 else f"polar{polar_index}"
        fig.layout[polar_name].update(
            bgcolor="rgba(247,249,254,0.75)",
            radialaxis=dict(
                range=[0, max_frequency],
                showticklabels=False,
                ticks="",
                gridcolor="rgba(204,214,235,0.35)",
            ),
            angularaxis=dict(
                direction="clockwise",
                rotation=90,
                tickmode="array",
                tickvals=[0, 90, 180, 270],
                ticktext=["N", "E", "S", "W"],
                tickfont=dict(size=8),
            ),
        )

    return fig
\end{lstlisting}

\Needspace{8\baselineskip}

\subsection{\texorpdfstring{\texttt{backend/gestionar\_arxius\_backend.py}}{backend/gestionar\_arxius\_backend.py}}

\begin{lstlisting}[style=studioPython,caption={backend/gestionar\_arxius\_backend.py},label={lst:backend-gestionar-arxius-backend-py}]
"""Gestió de fitxers per a actius, models i dades meteorològiques de Studio.

Aquest mòdul fa una còpia de seguretat de la pàgina del navegador de fitxers. Formata metadades del sistema de fitxers,
resumeix els fitxers meteorològics EPW, descobreix artefactes de model entrenats i s'aplica
operacions de còpia/supressió limitades utilitzades per la UI de Streamlit.
"""

import os
import shutil
from dataclasses import dataclass
from datetime import datetime
from pathlib import Path

import pandas as pd
import streamlit as st


ROOT_PATH = Path.cwd().resolve()
DATA_BTC_PATH = ROOT_PATH / 'sinergym' / 'data' / 'buildings'
DATA_WEA_PATH = ROOT_PATH / 'sinergym' / 'data' / 'weather'
MODEL_FILE_EXTENSIONS = (".zip", ".pkl")
MODEL_ROOT_DIRECTORY_NAMES = {"models", "trainings"}


@dataclass(frozen=True)
class ExplorerItem:
    """Fila del navegador de fitxers serialitzable que es mostra a la pàgina de gestió d'actius."""
    name: str
    path: str
    size: str
    mtime: str
    is_dir: bool


def get_size_format(b, factor=1024, suffix="B"):
    """Retorna una mida compacta i llegible, com ara ``12.40KB``."""
    for unit in ["", "K", "M", "G", "T", "P", "E", "Z"]:
        if b < factor:
            return f"{b:.2f}{unit}{suffix}"
        b /= factor
    return f"{b:.2f}Y{suffix}"


@st.cache_data(show_spinner=False)
def parse_epw_overview(epw_file):
    """Llegeix les metadades d'ubicació i la temperatura mitjana del bulb sec d'un fitxer EPW."""
    with open(epw_file, 'r', encoding='utf-8', errors='ignore') as handle:
        lines = handle.readlines()

    loc = lines[0].strip().split(',')
    location = {
        'city': loc[1] if len(loc) > 1 else '-',
        'state': loc[2] if len(loc) > 2 else '-',
        'country': loc[3] if len(loc) > 3 else '-',
        'latitude': float(loc[6]) if len(loc) > 6 else 0.0,
        'longitude': float(loc[7]) if len(loc) > 7 else 0.0,
    }

    df = pd.read_csv(epw_file, skiprows=8, header=None)
    avg_temp = float(df.iloc[:, 6].mean())
    climate_zone = 'hot' if avg_temp > 22 else 'cold' if avg_temp < 12 else 'cool'
    return location, avg_temp, climate_zone


@st.cache_data(show_spinner=False)
def load_weather_map_rows(root_path_str):
    """Retorna les files de fitxers meteorològics preparats per al mapa descobertes a continuació ``root_path_str``."""
    root_path = Path(root_path_str)
    rows = []
    color_map = {
        'hot': [230, 120, 40],
        'cool': [70, 130, 180],
        'cold': [80, 190, 220],
    }

    for epw_path in sorted(root_path.rglob('*.epw')):
        try:
            location, avg_temp, climate_zone = parse_epw_overview(str(epw_path))
        except Exception:
            continue

        color = color_map.get(climate_zone, [120, 120, 120])
        rows.append({
            'file_name': epw_path.name,
            'relative_path': str(epw_path.relative_to(root_path)),
            'city': location['city'],
            'country': location['country'],
            'latitude': location['latitude'],
            'longitude': location['longitude'],
            'climate': climate_zone,
            'avg_temp': round(avg_temp, 2),
            'color_r': color[0],
            'color_g': color[1],
            'color_b': color[2],
        })

    return pd.DataFrame(rows)


def delete_item(path_str):
    """Suprimeix l'element."""
    path = Path(path_str)
    try:
        if path.is_dir():
            shutil.rmtree(path)
        else:
            if path.exists():
                path.unlink()
    except Exception as exc:
        return str(exc)
    return None


def list_explorer_items(current_path, root_path, filter_func=None):
    """Llista els elements de l'explorador."""
    try:
        with os.scandir(current_path) as iterator:
            entries = list(iterator)
    except Exception as exc:
        return (), str(exc)

    folders = []
    files = []

    for entry in entries:
        if filter_func:
            if not filter_func(entry.name, entry.is_dir(), current_path == root_path):
                continue

        try:
            stat = entry.stat()
            mtime = datetime.fromtimestamp(stat.st_mtime).strftime('%Y-%m-%d %H:%M')
            size_str = "DIR" if entry.is_dir() else get_size_format(stat.st_size)

            item = ExplorerItem(
                name=entry.name,
                path=entry.path,
                size=size_str,
                mtime=mtime,
                is_dir=entry.is_dir(),
            )

            if entry.is_dir():
                folders.append(item)
            else:
                files.append(item)
        except Exception:
            continue

    folders.sort(key=lambda item: item.name.lower())
    files.sort(key=lambda item: item.name.lower())
    return tuple(folders + files), None


def filter_trainings(name, is_dir, is_root):
    """Entrenaments de filtre."""
    if is_root:
        return is_dir and (
            name == "trainings"
            or name.startswith("Eplus-")
            or name.startswith("ppo_")
        )
    return True


def filter_models(name, is_dir, is_root):
    """Models de filtre."""
    name_lower = name.lower()
    if not is_dir:
        return name_lower.endswith(MODEL_FILE_EXTENSIONS)
    if is_root:
        return (
            name_lower in MODEL_ROOT_DIRECTORY_NAMES
            or name_lower.startswith("eplus-")
            or name_lower.startswith("ppo_")
            or name_lower.startswith("training-")
        )
    return True


def filter_weather(name, is_dir, is_root):
    """Filtre el temps."""
    if not is_dir:
        return name.endswith(".epw")
    return True


def filter_envs(name, is_dir, is_root):
    """Filtre envs."""
    if not is_dir:
        return name.endswith(".epJSON") or name.endswith(".idf")
    return True
\end{lstlisting}

\Needspace{8\baselineskip}

\subsection{\texorpdfstring{\texttt{backend/interaccionar\_agent\_backend.py}}{backend/interaccionar\_agent\_backend.py}}

\begin{lstlisting}[style=studioPython,caption={backend/interaccionar\_agent\_backend.py},label={lst:backend-interaccionar-agent-backend-py}]
"""Execució en directe per provar agents entrenats dins d'un entorn Sinergym.

La pàgina de control en directe utilitza aquest mòdul per crear entorns Sinergym amb wrappers,
carregueu les polítiques Stable-Baselines3, valideu la compatibilitat de les accions i traduïu-les
observacions/accions en estructures fàcils de mostrar. Evita intencionadament
Streamlit, de manera que les mateixes comprovacions poden donar suport a l'avaluació i la documentació.
"""

from functools import partial
from pathlib import Path
from typing import Any, Callable, Dict, List, Optional, Sequence, Tuple

import gymnasium as gym
import numpy as np
from stable_baselines3.common.base_class import BaseAlgorithm
from stable_baselines3.common.vec_env import VecNormalize
from sinergym.utils.wrappers import CSVLogger, LoggerWrapper, NormalizeAction

from backend.model_metadata import (
    build_gym_kwargs_from_metadata,
    load_model_metadata,
    validate_action_spaces,
)
from backend.sb3_utils import (
    candidate_vecnorm,
    build_monitored_vec_env,
    env_id_from_meta_or_name,
    format_action_for_env,
    load_sb3_model_bytes,
    load_vecnormalize,
    scan_model_zips,
    should_apply_normalize_action,
)
from backend.common import is_registered_env_id


def _make_runtime_env(
    selected_env_id: str,
    wrapper_configs: Optional[Sequence[Dict[str, Any]]] = None,
    *,
    trust_training_metadata: bool = False,
    model_action_space: Any = None,
    gym_kwargs: Optional[Dict[str, Any]] = None,
) -> Any:
    """Crea una instància d'entorn d'interacció en directe embolicada."""
    env = gym.make(selected_env_id, **(gym_kwargs or {}))
    if wrapper_configs:
        from backend.entrenar_agent_wrappers import apply_training_wrappers

        env = apply_training_wrappers(env, wrapper_configs)
    elif not trust_training_metadata and should_apply_normalize_action(model_action_space, env.action_space):
        env = NormalizeAction(env)
    env = LoggerWrapper(env)
    env = CSVLogger(env)
    return env


def mode_requires_model(mode_label: str) -> bool:
    """El mode requereix un model."""
    lowered = (mode_label or "").lower()
    return "interactiva" in lowered or "avaluaci" in lowered


def load_model_training_metadata(model_path: Path) -> Dict[str, Any]:
    """Carrega les metadades d'entrenament del model."""
    return load_model_metadata(model_path)


def build_env_factory(
    selected_env_id: str,
    wrapper_configs: Optional[Sequence[Dict[str, Any]]] = None,
    *,
    trust_training_metadata: bool = False,
    model_action_space: Any = None,
    gym_kwargs: Optional[Dict[str, Any]] = None,
) -> Callable[[], Any]:
    """Crea la funció que inicialitza l'entorn de la sessió d'interacció."""
    return partial(
        _make_runtime_env,
        selected_env_id,
        wrapper_configs,
        trust_training_metadata=trust_training_metadata,
        model_action_space=model_action_space,
        gym_kwargs=gym_kwargs,
    )


def validate_action_contract(
    model_obj: Optional[BaseAlgorithm],
    vec_env: Any,
    metadata: Dict[str, Any],
) -> None:
    """Validació del contracte d'actuació."""
    if model_obj is None:
        return

    model_action_space = getattr(model_obj, "action_space", None)
    env_action_space = getattr(vec_env, "action_space", None)
    validate_action_spaces(model_action_space, env_action_space, metadata)


def _space_shape(space: Any) -> Tuple[int, ...] | None:
    """Space shape."""
    shape = getattr(space, "shape", None)
    if shape is None:
        return None
    try:
        return tuple(int(value) for value in shape)
    except Exception:
        return None


def get_model_observation_size(model_obj: Optional[BaseAlgorithm]) -> Optional[int]:
    """Retorna la mida d'observació del model."""
    if model_obj is None:
        return None
    shape = _space_shape(getattr(model_obj, "observation_space", None))
    if not shape:
        return None
    return int(np.prod(shape))


def get_observation_size_from_space(space: Any) -> Optional[int]:
    """Retorna la mida d'observació des de l'espai."""
    shape = _space_shape(space)
    if not shape:
        return None
    return int(np.prod(shape))


def get_observation_size_from_array(obs: Any) -> Optional[int]:
    """Retorna la mida de l'observació de la matriu."""
    try:
        arr = np.asarray(obs)
        if arr.ndim > 1:
            arr = arr.reshape(arr.shape[0], -1)[0]
        return int(arr.size)
    except Exception:
        return None


def flatten_observation_values(values: Any) -> List[float]:
    """Aplanar els valors d'observació."""
    if values is None:
        return []
    try:
        return [float(value) for value in np.asarray(values, dtype=np.float32).reshape(-1)]
    except Exception:
        flattened: List[float] = []
        try:
            iterator = list(values)
        except Exception:
            iterator = []
        for value in iterator:
            if isinstance(value, (list, tuple, np.ndarray)):
                flattened.extend(flatten_observation_values(value))
                continue
            try:
                flattened.append(float(value))
            except Exception:
                pass
        return flattened


def get_runtime_observation_variables(
    core_env: Any,
    obs: Any = None,
    model_obj: Optional[BaseAlgorithm] = None,
    vec_env: Any = None,
) -> List[str]:
    """Retorna les variables d'observació del clima d'execució."""
    obs_vars = get_wrapper_variables(core_env, "observation_variables")
    target_size = (
        get_model_observation_size(model_obj)
        or get_observation_size_from_space(getattr(vec_env, "observation_space", None))
        or get_observation_size_from_array(obs)
        or len(obs_vars)
    )

    if not target_size:
        return obs_vars
    if len(obs_vars) == target_size:
        return obs_vars
    if obs_vars and target_size % len(obs_vars) == 0:
        repeats = target_size // len(obs_vars)
        if repeats > 1:
            return [
                f"{variable}_hist{block + 1}"
                for block in range(repeats)
                for variable in obs_vars
            ]
    return [f"obs_{index + 1:03d}" for index in range(target_size)]


def coerce_observation_values(values: Sequence[float], expected_size: Optional[int]) -> List[float]:
    """Converteix els valors d'observació."""
    coerced = flatten_observation_values(values)
    if expected_size is None or len(coerced) == expected_size:
        return coerced
    if len(coerced) > expected_size:
        return coerced[:expected_size]
    fill_value = coerced[-1] if coerced else 0.0
    return coerced + [fill_value] * (expected_size - len(coerced))


def initialize_runtime(
    selected_env_id: str,
    mode_label: str,
    model_path: Optional[Path] = None,
) -> Tuple[Any, Any, Optional[BaseAlgorithm], Any, List[str]]:
    """Inicialitza l'execució de l'entorn."""
    init_warnings: List[str] = []
    model_obj: Optional[BaseAlgorithm] = None
    vecnorm_path: Optional[Path] = None
    metadata: Dict[str, Any] = {}
    if mode_requires_model(mode_label) and model_path is not None:
        with open(model_path, "rb") as file_handle:
            model_obj, _ = load_sb3_model_bytes(file_handle.read(), device="cpu")
        vecnorm_path = candidate_vecnorm(model_path)
        metadata = load_model_training_metadata(model_path)

    metadata_env_id = metadata.get("env_id")
    if isinstance(metadata_env_id, str) and is_registered_env_id(metadata_env_id):
        # Si el model porta env_id al metadata, el fem prevaldre sobre el text escrit a mà.
        # Això evita carregar una política entrenada amb un espai d'accions diferent.
        selected_env_id = metadata_env_id

    has_wrapper_contract = "wrapper_configs" in metadata
    wrapper_configs = list(metadata.get("wrapper_configs") or [])
    gym_kwargs = build_gym_kwargs_from_metadata(metadata)
    env_factory = build_env_factory(
        selected_env_id,
        wrapper_configs=wrapper_configs,
        trust_training_metadata=has_wrapper_contract,
        model_action_space=getattr(model_obj, "action_space", None),
        gym_kwargs=gym_kwargs or None,
    )

    vec_env = None
    if mode_requires_model(mode_label) and vecnorm_path and vecnorm_path.exists():
        # VecNormalize només és segur si es reconstrueix amb la mateixa funció de creació d'entorn.
        try:
            vec_env = load_vecnormalize(str(vecnorm_path), env_factory)
        except ValueError as exc:
            init_warnings.append(
                f"VecNormalize no aplicat: {exc} "
                "Continuant sense normalitzacio d'observacions; les prediccions del model poden ser menys precises."
            )
    if vec_env is None:
        vec_env = build_monitored_vec_env(env_factory)

    try:
        validate_action_contract(model_obj, vec_env, metadata)
    except ValueError as exc:
        init_warnings.append(f"Avis de contracte d'accions: {exc}")

    model_obs_shape = _space_shape(getattr(model_obj, "observation_space", None)) if model_obj else None
    env_obs_shape = _space_shape(getattr(vec_env, "observation_space", None))
    if model_obs_shape and env_obs_shape and model_obs_shape != env_obs_shape:
        # En mode interactiu no aturem directament: avisem i deixem que l'usuari decideixi,
        # perquè pot estar fent una prova ràpida de compatibilitat.
        model_size = int(np.prod(model_obs_shape))
        env_size = int(np.prod(env_obs_shape))
        init_warnings.append(
            f"Avis de contracte d'observacions: el model espera {model_size} variables {model_obs_shape} "
            f"pero l'entorn en dona {env_size} {env_obs_shape}. "
            "Les prediccions de l'agent no seran fiables."
        )

    core_env = vec_env.envs[0] if hasattr(vec_env, "envs") else vec_env
    obs = vec_env.reset()
    return vec_env, core_env, model_obj, obs, init_warnings


def get_wrapper_variables(core_env: Any, attr_name: str) -> List[Any]:
    """Retorna les variables del wrapper."""
    if hasattr(core_env, "get_wrapper_attr"):
        values = core_env.get_wrapper_attr(attr_name)
        if values is None:
            return []
        return list(values)
    return []


def extract_policy_test_defaults(
    obs: Any,
    info_dict: Dict[str, Any],
    obs_vars: Sequence[str],
    vec_env: Any,
    core_env: Any,
) -> List[float]:
    """Extreu els valors predeterminats de la prova de política."""
    raw_obs_values = obs[0] if getattr(obs, "ndim", 0) > 1 else obs
    has_raw = False
    if isinstance(info_dict, dict) and "observation" in info_dict:
        # Si l'entorn ens dona l'observació crua a info, és la millor base per al formulari.
        if len(info_dict["observation"]) == len(obs_vars):
            raw_obs_values = info_dict["observation"]
            has_raw = True

    if not has_raw:
        if isinstance(vec_env, VecNormalize):
            # En avaluació estàtica volem valors humans; desfem VecNormalize quan és possible.
            try:
                raw_obs_values = vec_env.unnormalize_obs(np.array([raw_obs_values]))[0]
            except Exception:
                pass

        temp = core_env
        while hasattr(temp, "env"):
            if type(temp).__name__ == "NormalizeObservation":
                # Alguns wrappers normalitzen sense VecNormalize; busquem aquest cas caminant
                # la cadena de wrappers.
                try:
                    std = np.sqrt(temp.obs_rms.var + temp.epsilon)
                    raw_obs_values = (raw_obs_values * std) + temp.obs_rms.mean
                except Exception:
                    pass
                break
            temp = temp.env

    return flatten_observation_values(raw_obs_values)


def randomize_observation_values(obs_vars: Sequence[str]) -> List[float]:
    """Aleatoritzar els valors d'observació."""
    import random

    new_vals: List[float] = []
    for var_name in obs_vars:
        lowered = var_name.lower()
        if "temperature" in lowered or "temp" in lowered:
            new_vals.append(random.uniform(5.0, 35.0))
        elif "power" in lowered or "rate" in lowered:
            new_vals.append(random.uniform(0.0, 5000.0))
        elif "month" in lowered:
            new_vals.append(float(random.randint(1, 12)))
        elif "day" in lowered:
            new_vals.append(float(random.randint(1, 31)))
        elif "hour" in lowered:
            new_vals.append(float(random.randint(0, 23)))
        else:
            new_vals.append(random.uniform(0.0, 1.0))
    return new_vals


def infer_unit(var_name: str) -> str:
    """Infer unit."""
    lowered = var_name.lower()
    if "temperature" in lowered or "temp" in lowered or "setpoint" in lowered:
        return "°C"
    if "power" in lowered or "rate" in lowered or "demand" in lowered:
        return "W"
    if "humidity" in lowered:
        return "%"
    if "speed" in lowered:
        return "m/s"
    if "direction" in lowered or "angle" in lowered:
        return "°"
    if "radiation" in lowered:
        return "W/m²"
    if "month" in lowered or "day" in lowered or "hour" in lowered:
        return ""
    return ""


def normalize_policy_observation(
    custom_obs: Sequence[float],
    core_env: Any,
    vec_env: Any,
    expected_size: Optional[int] = None,
) -> np.ndarray:
    """Normalitza l'observació de les polítiques."""
    expected_shape = _space_shape(getattr(vec_env, "observation_space", None))
    if expected_size is None and expected_shape:
        expected_size = int(np.prod(expected_shape))

    normalized_values = coerce_observation_values(custom_obs, expected_size)
    if expected_shape and int(np.prod(expected_shape)) == len(normalized_values):
        normalized = np.array(normalized_values, dtype=np.float32).reshape((1, *expected_shape))
    else:
        normalized = np.array([normalized_values], dtype=np.float32)

    temp = core_env
    while hasattr(temp, "env"):
        if type(temp).__name__ == "NormalizeObservation":
            try:
                std = np.sqrt(temp.obs_rms.var + temp.epsilon)
                normalized[0] = (normalized[0] - temp.obs_rms.mean) / std
            except Exception:
                pass
            break
        temp = temp.env

    if isinstance(vec_env, VecNormalize):
        return vec_env.normalize_obs(normalized)
    return normalized


def unnormalize_sinergym_action(action_array: Any, env_obj: Any) -> Any:
    """Unnormalize sinergym action."""
    temp = env_obj
    while hasattr(temp, "env"):
        if type(temp).__name__ == "NormalizeAction":
            return temp.action(action_array)
        temp = temp.env
    return action_array


def prepare_action_display(action: Any, vec_env: Any, core_env: Any) -> Any:
    """Prepara la visualització de l'acció."""
    action_formatted = format_action_for_env(action, vec_env)
    if isinstance(action_formatted, np.ndarray) and len(action_formatted.shape) > 1:
        action_display = action_formatted[0]
    else:
        action_display = action_formatted
    return unnormalize_sinergym_action(action_display, core_env)


def extract_display_observation(
    obs: Any,
    info_dict: Dict[str, Any],
    obs_vars: Sequence[str],
    vec_env: Any,
) -> Tuple[Any, Any]:
    """Extreu l'observació de visualització."""
    actual_obs = obs[0] if isinstance(obs, np.ndarray) and len(obs.shape) > 1 else obs
    raw_obs_values: Any = []
    has_raw = False

    # Preferim l'observació crua que torna Sinergym a info; és la que té unitats humanes.
    if isinstance(info_dict, dict) and "observation" in info_dict:
        last_raw_obs = info_dict["observation"]
        if len(last_raw_obs) == len(obs_vars):
            raw_obs_values = last_raw_obs
            has_raw = True

    if not has_raw and isinstance(vec_env, VecNormalize):
        # Si la run passa per VecNormalize, el model veu valors normalitzats. Per mostrar-los
        # a pantalla els tornem a l'escala original quan es pot.
        try:
            raw_obs_values = vec_env.unnormalize_obs(np.array([actual_obs]))[0]
            has_raw = True
        except Exception:
            pass

    display_obs = flatten_observation_values(raw_obs_values if has_raw else actual_obs)
    return actual_obs, display_obs


def extract_info_dict(infos: Any) -> Dict[str, Any]:
    """Extreu el diccionari info d'un step de VecEnv o Gym."""
    if isinstance(infos, (tuple, list, np.ndarray)):
        return infos[0] if len(infos) > 0 else {}
    if isinstance(infos, dict):
        return infos
    return {}


def build_history_entry(
    step_number: int,
    reward: float,
    info_dict: Dict[str, Any],
    next_obs: Any,
    obs_vars: Sequence[str],
    vec_env: Any,
) -> Dict[str, Any]:
    """Prepara una entrada d'historial per a una acció de l'agent."""
    hist_entry: Dict[str, Any] = {"pas": step_number, "reward": reward}

    raw_vals: Any = []
    has_raw = False
    if isinstance(info_dict, dict) and "observation" in info_dict:
        # Guardem historial en escala real perquè els KPI del panell siguin interpretables.
        last_raw = info_dict["observation"]
        if len(last_raw) == len(obs_vars):
            raw_vals = last_raw
            has_raw = True

    if not has_raw and isinstance(vec_env, VecNormalize):
        try:
            obs_2d = np.array([next_obs[0] if len(next_obs.shape) > 1 else next_obs])
            raw_vals = vec_env.unnormalize_obs(obs_2d)[0]
            has_raw = True
        except Exception:
            pass

    vals_to_log = flatten_observation_values(
        raw_vals if has_raw else (next_obs[0] if len(next_obs.shape) > 1 else next_obs)
    )
    if len(vals_to_log) == len(obs_vars):
        for var_name, value in zip(obs_vars, vals_to_log):
            hist_entry[var_name] = value

    return hist_entry


def run_environment_steps(
    vec_env: Any,
    core_env: Any,
    action_to_take: Any,
    current_step: int,
    repeat_n: int = 1,
    obs_vars: Optional[Sequence[str]] = None,
) -> Dict[str, Any]:
    """Executa els passos de l'entorn."""
    if hasattr(core_env, "action_space"):
        action_formatted = format_action_for_env(action_to_take, vec_env)
    else:
        action_formatted = action_to_take

    runtime_obs_vars = (
        list(obs_vars)
        if obs_vars is not None
        else get_wrapper_variables(core_env, "observation_variables")
    )
    history_entries: List[Dict[str, Any]] = []
    next_obs = None
    reward = 0.0
    info_dict: Dict[str, Any] = {}
    episode_finished = False
    reset_obs = None

    for offset in range(repeat_n):
        next_obs, rewards, dones, infos = vec_env.step(action_formatted)
        reward = rewards[0] if isinstance(rewards, (list, np.ndarray)) else rewards
        done = dones[0] if isinstance(dones, (list, np.ndarray)) else dones
        info_dict = extract_info_dict(infos)

        history_entries.append(
            build_history_entry(
                step_number=current_step + offset + 1,
                reward=reward,
                info_dict=info_dict,
                next_obs=next_obs,
                obs_vars=runtime_obs_vars,
                vec_env=vec_env,
            )
        )

        if done:
            episode_finished = True
            reset_obs = vec_env.reset()
            next_obs = reset_obs
            break

    return {
        "next_obs": next_obs,
        "reward": reward,
        "info_dict": info_dict,
        "history_entries": history_entries,
        "episode_finished": episode_finished,
        "reset_obs": reset_obs,
    }


def find_matching_column(row: Dict[str, Any], keywords: Sequence[str]) -> Optional[str]:
    """Troba una columna compatible."""
    for col_name in row.keys():
        if any(keyword.lower() in col_name.lower() for keyword in keywords):
            return col_name
    return None
\end{lstlisting}

\Needspace{8\baselineskip}

\subsection{\texorpdfstring{\texttt{backend/mapa\_backend.py}}{backend/mapa\_backend.py}}

\begin{lstlisting}[style=studioPython,caption={backend/mapa\_backend.py},label={lst:backend-mapa-backend-py}]
"""Utilitats de representació de mapes compartides construïdes a pydeck i Streamlit."""

from __future__ import annotations

import pandas as pd
import pydeck as pdk
import streamlit as st

from page_styles.theme import build_pydeck_tooltip


def _chart_width(use_container_width: bool) -> str:
    """Tradueix el flag antic de Streamlit al valor de width actual."""
    return "stretch" if use_container_width else "content"


def render_location_map(
    latitude: float,
    longitude: float,
    *,
    zoom: float = 4.5,
    color: tuple[int, int, int, int] = (95, 84, 249, 220),
    radius: int = 22000,
    tooltip_html: str | None = None,
    use_container_width: bool = True,
) -> None:
    """Dibuixa un mapa de dispersió d'un sol punt centrat en les coordenades donades.

    tooltip_html accepta HTML preformatats (no calen referències de columna pydeck).
    """

    point_frame = pd.DataFrame([{"latitude": latitude, "longitude": longitude}])
    tooltip = build_pydeck_tooltip(tooltip_html, variant="location") if tooltip_html else None
    st.pydeck_chart(
        pdk.Deck(
            initial_view_state=pdk.ViewState(
                latitude=latitude,
                longitude=longitude,
                zoom=zoom,
                pitch=0,
            ),
            layers=[
                pdk.Layer(
                    "ScatterplotLayer",
                    data=point_frame,
                    get_position="[longitude, latitude]",
                    get_fill_color=list(color),
                    get_radius=radius,
                    pickable=tooltip_html is not None,
                )
            ],
            tooltip=tooltip,
        ),
        width=_chart_width(use_container_width),
    )


def render_multipoint_map(
    data_frame: pd.DataFrame,
    *,
    center_latitude: float | None = None,
    center_longitude: float | None = None,
    zoom: float = 2,
    color: tuple[int, int, int, int] = (255, 0, 0, 200),
    radius: int = 55000,
    tooltip_html: str | None = None,
    use_container_width: bool = True,
) -> None:
    """Dibuixa un mapa de dispersió de diversos punts des d'un DataFrame amb latitud/longitud.

    La vista es centra a les coordenades mitjanes tret que center_latitude/center_longitude
    es proporcionen de manera explícita. tooltip_html admet marcadors de posició de nom de columna pydeck
    (e.g. '<b>{file_name}</b><br/>{ciutat}').
    """

    lat_center = center_latitude if center_latitude is not None else float(data_frame["latitude"].mean())
    lon_center = center_longitude if center_longitude is not None else float(data_frame["longitude"].mean())
    tooltip = build_pydeck_tooltip(tooltip_html, variant="multipoint") if tooltip_html else None
    st.pydeck_chart(
        pdk.Deck(
            initial_view_state=pdk.ViewState(
                latitude=lat_center,
                longitude=lon_center,
                zoom=zoom,
                pitch=0,
            ),
            layers=[
                pdk.Layer(
                    "ScatterplotLayer",
                    data=data_frame,
                    get_position="[longitude, latitude]",
                    get_fill_color=list(color),
                    get_radius=radius,
                    pickable=tooltip_html is not None,
                )
            ],
            tooltip=tooltip,
        ),
        width=_chart_width(use_container_width),
    )
\end{lstlisting}

\Needspace{8\baselineskip}

\subsection{\texorpdfstring{\texttt{backend/model\_metadata.py}}{backend/model\_metadata.py}}

\begin{lstlisting}[style=studioPython,caption={backend/model\_metadata.py},label={lst:backend-model-metadata-py}]
"""Utilitats compartides per reconstruir models entrenats de l'Studio."""

from __future__ import annotations

import json
from pathlib import Path
from typing import Any, Dict

from backend.sb3_utils import action_spaces_compatible, describe_action_space


def load_model_metadata(model_path: Path) -> Dict[str, Any]:
    """Carrega les metadades d'entrenament associades a un model SB3."""
    metadata_path = model_path.with_name(model_path.stem + "_metadata.json")
    if metadata_path.exists():
        with open(metadata_path, "r", encoding="utf-8") as handle:
            metadata = json.load(handle)
        training_config = metadata.get("training_config") if isinstance(metadata.get("training_config"), dict) else {}
        if not metadata.get("env_id"):
            metadata["env_id"] = training_config.get("env_id")
        if not metadata.get("reward_name"):
            metadata["reward_name"] = training_config.get("reward_name")
        if "wrapper_configs" not in metadata and "wrapper_configs" in training_config:
            metadata["wrapper_configs"] = training_config.get("wrapper_configs", [])
        if "meters" not in metadata and "meters" in training_config:
            metadata["meters"] = training_config.get("meters")
        if "reward_kwargs" not in metadata and "reward_kwargs" in training_config:
            metadata["reward_kwargs"] = training_config.get("reward_kwargs")
        metadata["metadata_source"] = str(metadata_path)
        return metadata

    training_config_path = model_path.parent / "training_config.json"
    if training_config_path.exists():
        with open(training_config_path, "r", encoding="utf-8") as handle:
            payload = json.load(handle)
        training_config = dict(payload.get("training_config") or {})
        return {
            "env_id": payload.get("env_id") or training_config.get("env_id"),
            "reward_name": payload.get("reward_name") or training_config.get("reward_name"),
            "wrapper_configs": training_config.get("wrapper_configs", []),
            "meters": training_config.get("meters"),
            "reward_kwargs": training_config.get("reward_kwargs"),
            "training_config": training_config,
            "artifact_name": payload.get("artifact_name"),
            "metadata_source": str(training_config_path),
        }
    return {}


def build_gym_kwargs_from_metadata(metadata: Dict[str, Any]) -> Dict[str, Any]:
    """Reconstrueix els kwargs de Gym a partir de les metadades d'entrenament."""
    training_config = metadata.get("training_config") or {}
    meters = metadata.get("meters") or training_config.get("meters")
    reward_name = metadata.get("reward_name") or training_config.get("reward_name")
    reward_kwargs = metadata.get("reward_kwargs") or training_config.get("reward_kwargs") or {}

    kwargs: Dict[str, Any] = {}
    if not reward_name and not meters:
        return kwargs

    try:
        from backend.entrenar_agent_env import with_detailed_hvac_meters

        kwargs["meters"] = with_detailed_hvac_meters(meters or {})
        reward_cls = resolve_reward_class(reward_name)
        if reward_cls is not None:
            kwargs["reward"] = reward_cls
            if reward_kwargs:
                kwargs["reward_kwargs"] = reward_kwargs
    except Exception:
        pass
    return kwargs


def resolve_reward_class(reward_name: str | None) -> Any:
    """Resol una reward usant les classes exposades per Sinergym."""
    if not reward_name:
        return None
    if ":" in reward_name:
        from sinergym.utils.common import import_from_path

        return import_from_path(reward_name)
    import sinergym.utils.rewards as reward_module

    return getattr(reward_module, reward_name, None)


def validate_action_spaces(
    model_action_space: Any,
    env_action_space: Any,
    metadata: Dict[str, Any],
) -> None:
    """Valida que el model i l'entorn comparteixin el mateix contracte d'accions."""
    if action_spaces_compatible(model_action_space, env_action_space):
        return

    wrappers = metadata.get("wrapper_configs") or []
    wrapper_names = ", ".join(str(item.get("name")) for item in wrappers if isinstance(item, dict)) or "cap detectat"
    source = metadata.get("metadata_source") or "metadades no trobades"
    raise ValueError(
        "L'espai d'accions de l'entorn no coincideix amb el que espera el model. "
        f"Model: {describe_action_space(model_action_space)}; "
        f"entorn: {describe_action_space(env_action_space)}. "
        f"Wrappers detectats: {wrapper_names}. Font: {source}."
    )
\end{lstlisting}

\Needspace{8\baselineskip}

\subsection{\texorpdfstring{\texttt{backend/mostrar\_entorn\_backend.py}}{backend/mostrar\_entorn\_backend.py}}

\begin{lstlisting}[style=studioPython,caption={backend/mostrar\_entorn\_backend.py},label={lst:backend-mostrar-entorn-backend-py}]
"""Backend de metadades de l'entorn: carrega actius, PV, emmagatzematge i dades per pintar la UI."""

from __future__ import annotations

import json
from typing import Any, Dict, List, Optional, Tuple

import pandas as pd
import streamlit as st

from backend.viewer_3d_backend import build_surface_records


@st.cache_data(show_spinner=False)
def load_epjson(path: str) -> dict:
    """Llegeix i retorna un fitxer epJSON del disc (emmagatzemat a la memòria cau)."""
    with open(path, 'r') as fh:
        return json.load(fh)


def _safe_float(value: Any) -> Optional[float]:
    """Retorna float(valor) o None si la conversió falla."""
    try:
        return None if value is None else float(value)
    except Exception:
        return None


def _getf(data: dict, *keys: str) -> Optional[float]:
    """Retorna el primer valor numèric trobat entre les claus donades, o None."""
    for key in keys:
        value = data.get(key)
        try:
            if value is None:
                continue
            return float(value)
        except Exception:
            continue
    return None


def _derive_kibam_metrics(
    obj: dict,
) -> Tuple[Optional[float], Optional[float], Optional[float]]:
    """Deriva (energy_kwh, p_charge_kw, p_discharge_kw) d'un objecte de bateria KiBaM.

    Processa ElectricLoadCenter:Emmagatzematge:camps de bateria per estimar el nivell de paquet
    capacitat energètica i límits de potència de càrrega/descàrrega.
    """
    cap_ah = _getf(obj, "maximum_module_capacity")
    voltage_full = _getf(obj, "fully_charged_module_open_circuit_voltage")
    voltage_cut = _getf(obj, "module_cut_off_voltage")
    c_rate = _getf(obj, "module_charge_rate_limit")
    current_discharge = _getf(obj, "maximum_module_discharging_current")

    n_series = int(_getf(obj, "number_of_battery_modules_in_series") or 1)
    n_parallel = int(_getf(obj, "number_of_battery_modules_in_parallel") or 1)

    if cap_ah is None:
        return None, None, None

    if voltage_full and voltage_cut:
        voltage_nominal = 0.5 * (voltage_full + voltage_cut)
    elif voltage_full:
        voltage_nominal = 0.95 * voltage_full
    elif voltage_cut:
        voltage_nominal = 1.10 * voltage_cut
    else:
        voltage_nominal = 12.0

    voltage_pack = voltage_nominal * n_series
    cap_ah_pack = cap_ah * n_parallel
    energy_kwh = (cap_ah_pack * voltage_pack) / 1000.0

    p_charge_kw = None
    if c_rate is not None:
        p_charge_kw = (c_rate * cap_ah * n_parallel * voltage_pack) / 1000.0

    p_discharge_kw = None
    if current_discharge is not None:
        p_discharge_kw = (current_discharge * n_parallel * voltage_pack) / 1000.0

    return energy_kwh, p_charge_kw, p_discharge_kw


def parse_pv_and_storage(epjson: dict) -> Dict[str, Any]:
    """Analitza generadors de PV i objectes d'emmagatzematge de la bateria des d'un epJSON dict.

    Admet generador: fotovoltaic, generador: PVWatts,
    ElectricLoadCenter:Emmagatzematge:Simple, :Bateria i :LiIonNMCBattery.
    Retorna un dict amb les claus 'pv' i 'emmagatzematge', cadascuna una llista de dicts normalitzats.
    """
    output: Dict[str, List[dict]] = {'pv': [], 'storage': []}

    # EnergyPlus té dues famílies habituals de generadors solars. Les normalitzem a la
    # mateixa estructura perquè el visor no hagi de conèixer cada tipus.
    for key in ("Generator:Photovoltaic", "Generator:PVWatts"):
        for name, obj in epjson.get(key, {}).items():
            output['pv'].append({
                'name': name,
                'surface_name': obj.get("surface_name"),
                'type': key,
                'capacity_dc': obj.get("dc_system_capacity"),
                'num_series': obj.get("number_of_modules_in_series"),
                'num_strings': obj.get("number_of_series_strings"),
                'performance_obj': obj.get("module_performance_name"),
                'array_geometry_type': obj.get("array_geometry_type"),
            })

    for name, obj in epjson.get("ElectricLoadCenter:Storage:Simple", {}).items():
        # En Storage:Simple l'eficiència de volta completa es pot estimar multiplicant
        # càrrega i descàrrega, si totes dues existeixen.
        eff_charge = _safe_float(obj.get("nominal_energètic_efficiency_for_charging"))
        eff_discharge = _safe_float(obj.get("nominal_discharging_energètic_efficiency"))
        eff_roundtrip = None
        if eff_charge is not None and eff_discharge is not None:
            try:
                eff_roundtrip = float(eff_charge) * float(eff_discharge)
            except Exception:
                pass
        output['storage'].append({
            'name': name,
            'type': "ElectricLoadCenter:Storage:Simple",
            'energy_kwh': _safe_float(obj.get("maximum_storage_capacity")),
            'p_charge_kw': _safe_float(obj.get("maximum_power_for_charging")),
            'p_discharge_kw': _safe_float(obj.get("maximum_power_for_discharging")),
            'eff_roundtrip': eff_roundtrip,
            'zone_name': obj.get("zone_name"),
        })

    for name, obj in epjson.get("ElectricLoadCenter:Storage:Battery", {}).items():
        energy_kwh, p_charge_kw, p_discharge_kw = _derive_kibam_metrics(obj)
        output['storage'].append({
            'name': name,
            'type': "ElectricLoadCenter:Storage:Battery",
            'energy_kwh': None if energy_kwh is None else round(energy_kwh, 2),
            'p_charge_kw': None if p_charge_kw is None else round(p_charge_kw, 2),
            'p_discharge_kw': None if p_discharge_kw is None else round(p_discharge_kw, 2),
            'eff_roundtrip': None,
            'zone_name': obj.get("zone_name"),
        })

    for name, obj in epjson.get("ElectricLoadCenter:Storage:LiIonNMCBattery", {}).items():
        output['storage'].append({
            'name': name,
            'type': "ElectricLoadCenter:Storage:LiIonNMCBattery",
            'energy_kwh': None,
            'p_charge_kw': None,
            'p_discharge_kw': None,
            'eff_roundtrip': None,
            'zone_name': obj.get("zone_name"),
        })

    return output


@st.cache_data(show_spinner=False)
def load_environment_assets(epjson_path: str) -> Dict[str, Any]:
    """Carrega i prepareu tots els actius principals per a un entorn Sinergym (emmagatzemats a la memòria cau).

    Llegeix el fitxer epJSON, crea registres de superfície enriquits i analitza
    PV/objectes d'emmagatzematge. Retorna un dict unificat preparat per a UI.
    """
    data = load_epjson(epjson_path)
    enriched = build_surface_records(data)
    energy = parse_pv_and_storage(data)
    return {
        'records': enriched['records'],
        'zones': enriched['zones'],
        'types': enriched['types'],
        'center': enriched['center'],
        'z_clip_range': enriched['z_clip_range'],
        'name_index': enriched['name_index'],
        'pv': energy['pv'],
        'storage': energy['storage'],
    }


def summarize_pv(
    records: List[dict],
    name_index: Dict[str, int],
    pv_list: List[dict],
) -> Dict[str, Any]:
    """Calcula mètriques agregades dels panells PV associats a superfícies de l'edifici.

    Recompte de retorns, llista de noms de superfície actives, àrea total, inclinació mitjana i
    azimut mitjà. L'àrea i els angles són None quan no es troben superfícies coincidents.
    """
    used_surfaces = []
    total_area = 0.0
    tilts: List[float] = []
    azimuths: List[float] = []
    for pv in pv_list:
        surface_name = pv.get('surface_name')
        if surface_name in name_index:
            record = records[name_index[surface_name]]
            used_surfaces.append(surface_name)
            total_area += record['area']
            tilts.append(record['tilt'])
            azimuths.append(record['azimuth'])
    return {
        'count': len(pv_list),
        'surfaces_with_pv': sorted(set(used_surfaces)),
        'area_m2': total_area,
        'avg_tilt': sum(tilts) / len(tilts) if tilts else None,
        'avg_azimuth': sum(azimuths) / len(azimuths) if azimuths else None,
    }


@st.cache_data(show_spinner=False)
def parse_epw_metadata_and_temperature(epw_file: str):
    """Llegeix les metadades d'ubicació i la temperatura anual mitjana d'un fitxer EPW (emmagatzemat a la memòria cau).

    Retorna (location_dict, avg_temp_celsius, climate_zone_label).
    """
    with open(epw_file, 'r') as fh:
        lines = fh.readlines()

    loc = lines[0].strip().split(',')
    loc_data = {
        'city': loc[1], 'state': loc[2], 'country': loc[3],
        'latitude': float(loc[6]), 'longitude': float(loc[7]),
        'timezone': float(loc[8]), 'elevation': float(loc[9]),
    }
    df = pd.read_csv(epw_file, skiprows=8, header=None)
    avg_temp = df.iloc[:, 6].mean()
    climate_zone = "hot" if avg_temp > 22 else "cold" if avg_temp < 12 else "cool"
    return loc_data, avg_temp, climate_zone


def _format_ui_value(value: Any) -> str:
    """Formata qualsevol valor per mostrar-lo, retornant '-' per als valors buits/None."""
    if value is None:
        return '-'
    if isinstance(value, (list, tuple, set, dict)) and len(value) == 0:
        return '-'
    return str(value)


def _tuple_mapping_to_df(mapping: Dict[str, Any], columns: List[str]) -> pd.DataFrame:
    """Converteix un dict d'entrades amb valors de tupla a una etiqueta DataFrame."""
    rows = []
    for alias, values in (mapping or {}).items():
        flat = list(values) if isinstance(values, (list, tuple)) else [values]
        padded = flat + [''] * max(0, len(columns) - len(flat))
        row: Dict[str, Any] = {'Alias': alias}
        for index, column in enumerate(columns):
            row[column] = _format_ui_value(padded[index])
        rows.append(row)
    return pd.DataFrame(rows)


def _sequence_to_df(values: List[Any], column_name: str) -> pd.DataFrame:
    """Converteix una llista plana de valors en una sola columna DataFrame."""
    return pd.DataFrame([{column_name: _format_ui_value(v)} for v in (values or [])])


def _reward_summary_frames(
    reward_kwargs: Dict[str, Any],
) -> Tuple[pd.DataFrame, pd.DataFrame]:
    """Divideix els kwargs de recompensa en escalar i agrupa (list-valued) DataFrames."""
    scalar_rows: List[dict] = []
    grouped_rows: List[dict] = []
    for key, value in (reward_kwargs or {}).items():
        if isinstance(value, (list, tuple)):
            grouped_rows.append({'Parametre': key, 'Valor': ', '.join(map(_format_ui_value, value))})
        else:
            scalar_rows.append({'Parametre': key, 'Valor': _format_ui_value(value)})
    return pd.DataFrame(scalar_rows), pd.DataFrame(grouped_rows)
\end{lstlisting}

\Needspace{8\baselineskip}

\subsection{\texorpdfstring{\texttt{backend/resultats\_backend.py}}{backend/resultats\_backend.py}}

\begin{lstlisting}[style=studioPython,caption={backend/resultats\_backend.py},label={lst:backend-resultats-backend-py}]
"""Càrrega de dades, preparació de KPI i descobriment d'execució per a la pàgina de resultats.

El backend de resultats converteix els artefactes de entrenament i avaluació persistents en el
Estructura ``DashboardData`` consumida per Streamlit i pel generador d'informes.
Manté deliberadament el descobriment del sistema de fitxers, la normalització CSV i les dades d'acció
alineació fora del fitxer de pàgina, de manera que el mateix comportament es pot reutilitzar de manera automatitzada
exportacions.
"""

from __future__ import annotations

import copy
import glob
import json
import os
import re
from dataclasses import dataclass
from pathlib import Path
from typing import Any

import pandas as pd
import plotly.graph_objects as go

from backend.grafics.time_utils import _ensure_datetime_index


@dataclass(frozen=True)
class RunOption:
    """Directori d'execució seleccionable que es mostra a la pàgina de resultats."""

    path: str
    name: str

    def __str__(self) -> str:
        """Str."""
        return self.name


@dataclass(frozen=True)
class ResultsPageState:
    """S'han resolt les entrades necessàries abans que es pugui representar la pàgina de resultats."""

    default_base_dir: str
    run_dirs: tuple[RunOption, ...]
    warning_message: str | None = None


@dataclass(frozen=True)
class RunArtifacts:
    """Fitxers d'artefactes derivats d'una execució seleccionada."""

    progress_path: str
    observations_path: str
    error_message: str | None = None


@dataclass
class DashboardData:
    """Es requereix un paquet de dades carregat per representar el panell de resultats en línia."""

    progress: pd.DataFrame
    obs: pd.DataFrame
    actions: pd.DataFrame
    metrics_dict: dict
    zone_opts: list
    has_zones: bool
    all_zone_options: list
    has_occupied_scope: bool
    comfort_scope_options: list
    default_comfort_scope: str
    env_name: str
    run_mode: str
    timesteps_num: int
    model_name: str
    progress_path: str = ""
    observations_path: str = ""
    actions_path: str = ""
    yaml_cfg: Any = None


_IGNORED_RUN_PARTS = {"__pycache__", ".git", ".venv", "venv", "node_modules"}
_ACTION_CONTEXT_COLUMNS = (
    "datetime",
    "timestamp",
    "month",
    "day_of_month",
    "hour",
    "time_elapsed(hours)",
)


def _is_valid_run_directory(run_directory: Path) -> bool:
    """Retorna True si el directori conté tant progress.csv com almenys un observations.csv."""
    if any(part in _IGNORED_RUN_PARTS for part in run_directory.parts):
        return False
    if not (run_directory / "progress.csv").exists():
        return False
    observation_files = glob.glob(
        os.path.join(str(run_directory), "**", "observations.csv"),
        recursive=True,
    )
    return bool(observation_files)


def _build_run_display_name(run_directory: Path, base_directory: Path) -> str:
    """Retorna una etiqueta llegible per a un directori d'execució relativa al directori base."""
    try:
        relative_path = run_directory.relative_to(base_directory)
        return str(relative_path)
    except ValueError:
        return run_directory.name


def build_results_page_state(base_dir: str | None = None) -> ResultsPageState:
    """Escaneja el directori base i construeix l'estat inicial de resultats."""

    resolved_base_dir = Path.cwd() if base_dir is None else Path(base_dir)
    progress_files = glob.glob(
        os.path.join(str(resolved_base_dir), "**", "progress.csv"),
        recursive=True,
    )
    run_directories = sorted(
        {
            Path(path).parent
            for path in progress_files
            if _is_valid_run_directory(Path(path).parent)
        },
        key=os.path.getmtime,
        reverse=True,
    )

    if not run_directories:
        return ResultsPageState(
            default_base_dir=str(Path.cwd()),
            run_dirs=(),
            warning_message="No s'ha trobat cap execució amb progress.csv i observations.csv.",
        )

    return ResultsPageState(
        default_base_dir=str(Path.cwd()),
        run_dirs=tuple(
            RunOption(
                path=str(run_directory),
                name=_build_run_display_name(run_directory, resolved_base_dir),
            )
            for run_directory in run_directories
        ),
    )


def get_run_artifacts(selected_run: str) -> RunArtifacts:
    """Resol els fitxers d'artefactes principals associats a una execució seleccionada."""

    latest_progress = os.path.join(selected_run, "progress.csv")
    observation_files = glob.glob(
        os.path.join(selected_run, "**", "observations.csv"),
        recursive=True,
    )

    if not (os.path.exists(latest_progress) and observation_files):
        return RunArtifacts(
            progress_path=latest_progress,
            observations_path="",
            error_message="Falten progress.csv o observations.csv en l'execució seleccionada.",
        )

    sorted_observation_files = sorted(
        observation_files,
        key=os.path.getmtime,
        reverse=True,
    )
    selected_observation = (
        sorted_observation_files[1]
        if len(sorted_observation_files) > 1
        else sorted_observation_files[0]
    )

    return RunArtifacts(
        progress_path=latest_progress,
        observations_path=selected_observation,
    )


def load_action_data(observations_path: str) -> tuple[pd.DataFrame, str]:
    """Carrega la primera acció admesa CSV que es troba al costat del fitxer d'observacions."""

    monitor_dir = Path(observations_path).parent
    for filename in ("simulated_actions.csv", "agent_actions.csv"):
        candidate = monitor_dir / filename
        if not candidate.is_file():
            continue
        try:
            return pd.read_csv(candidate), str(candidate)
        except Exception:
            continue
    return pd.DataFrame(), ""


def align_actions_with_observations(
    actions: pd.DataFrame,
    obs: pd.DataFrame,
) -> pd.DataFrame:
    """Alineeu les files d'acció amb les marques de temps d'observació i columnes temporals útils."""

    if actions.empty or obs.empty:
        return actions

    aligned = actions.copy()
    obs_index = obs.index
    action_count = len(aligned)
    if len(obs_index) == action_count + 1:
        context = obs.iloc[1:action_count + 1]
    else:
        context = obs.iloc[:action_count]

    aligned = aligned.iloc[:len(context)].copy()
    aligned.index = context.index
    for column in _ACTION_CONTEXT_COLUMNS:
        if column in context.columns and column not in aligned.columns:
            aligned[column] = context[column].to_numpy()
    return aligned


def select_actions_for_obs(actions: pd.DataFrame, obs: pd.DataFrame) -> pd.DataFrame:
    """Retorna les files d'acció que corresponen al sector d'observació proporcionat."""

    if actions.empty or obs.empty:
        return pd.DataFrame()
    if isinstance(actions.index, pd.DatetimeIndex) and isinstance(obs.index, pd.DatetimeIndex):
        selected = actions.loc[actions.index.isin(obs.index)].copy()
        if not selected.empty:
            return selected.sort_index(kind="mergesort")
    return actions.iloc[:len(obs)].copy().sort_index(kind="mergesort")


def is_radiant_run(data: DashboardData) -> bool:
    """Retorna si l'execució seleccionada pertany a una configuració d'edifici radiant."""

    run_text = f"{data.env_name} {data.model_name} {data.yaml_cfg}".lower()
    return "radiant" in run_text


def select_radiant_action_data(actions: pd.DataFrame, data: DashboardData) -> pd.DataFrame:
    """Retorna les dades d'acció només quan l'execució seleccionada utilitza el control del sòl radiant."""

    if actions.empty or not is_radiant_run(data):
        return pd.DataFrame(index=actions.index)
    return actions


def load_dashboard_data(progress_path: str, observations_path: str) -> DashboardData:
    """Carrega i retorna el paquet de dades complet utilitzat pel panell de resultats."""

    import yaml

    from backend.grafics.comfort_scope import has_occupied_data
    from backend.grafics.data_loader import load_data
    from backend.grafics.figures_zones import get_zone_options
    from backend.grafics.metrics import compute_metrics

    yaml_cfg = None
    try:
        yaml_candidates = list(Path(progress_path).parent.glob("*.yaml"))
        if yaml_candidates:
            with open(yaml_candidates[0], "r", encoding="utf-8") as config_file:
                yaml_cfg = yaml.safe_load(config_file)
    except Exception:
        pass

    progress, obs = load_data(progress_path, observations_path)
    obs = _ensure_observation_time_columns(obs)

    actions, actions_path = load_action_data(observations_path)
    actions = align_actions_with_observations(actions, obs)

    metrics_dict = compute_metrics(progress, obs, yaml_cfg=yaml_cfg)
    zone_opts = get_zone_options(obs, yaml_cfg)
    has_zones = len(zone_opts) > 1
    all_zone_options = [{"label": "Global (Totes)", "value": "ALL"}] + zone_opts
    has_occupied_scope = has_occupied_data(obs)
    comfort_scope_options = [{"label": "Totes les hores", "value": "all"}]
    if has_occupied_scope:
        comfort_scope_options.append(
            {"label": "Nomes hores ocupades", "value": "occupied"}
        )
    default_comfort_scope = "occupied" if has_occupied_scope else "all"

    run_folder = Path(progress_path).parent
    run_folder_name = run_folder.name
    match_env = re.match(r"(.*?)(?:-res\d+)?$", run_folder_name)
    env_name = match_env.group(1) if match_env else run_folder_name
    is_eval = len(progress) < 10
    run_mode = "Avaluacio" if is_eval else "Entrenament"

    if "time/total_timesteps" in progress.columns:
        timesteps_num = int(progress["time/total_timesteps"].max())
    elif "length(timesteps)" in progress.columns:
        timesteps_num = int(progress["length(timesteps)"].sum())
    else:
        timesteps_num = len(obs)

    model_name = resolve_model_name_for_run(progress_path, env_name)

    return DashboardData(
        progress=progress,
        obs=obs,
        actions=actions,
        metrics_dict=metrics_dict,
        zone_opts=zone_opts,
        has_zones=has_zones,
        all_zone_options=all_zone_options,
        has_occupied_scope=has_occupied_scope,
        comfort_scope_options=comfort_scope_options,
        default_comfort_scope=default_comfort_scope,
        env_name=env_name,
        run_mode=run_mode,
        timesteps_num=timesteps_num,
        model_name=model_name,
        progress_path=progress_path,
        observations_path=observations_path,
        actions_path=actions_path,
        yaml_cfg=yaml_cfg,
    )


def resolve_model_name_for_run(progress_path: str, env_name: str) -> str:
    """Retorna el nom del model desat més probable per a una execució de resultats."""

    run_folder = Path(progress_path).parent
    env_slug = _slugify_model_lookup(env_name)
    fallback = "Desconegut / No Especificat"

    metadata_candidate = _model_name_from_training_metadata(run_folder, env_name, env_slug)
    if metadata_candidate:
        return metadata_candidate

    zip_candidate = _model_name_from_zip_search(run_folder, env_name, env_slug)
    if zip_candidate:
        return zip_candidate

    return fallback


def _slugify_model_lookup(value: str) -> str:
    """Normalitza els noms igual que els directoris d'artefactes d'entrenament."""

    cleaned = re.sub(r"[^A-Za-z0-9]+", "-", str(value)).strip("-").lower()
    return cleaned or "sense-nom"


def _candidate_model_roots(run_folder: Path) -> list[Path]:
    """Retorna les arrels probables que continguin artefactes de model desats."""

    cwd = Path.cwd()
    roots = [
        run_folder,
        run_folder.parent,
        cwd / "trainings",
        cwd / "models",
        cwd,
    ]
    unique_roots: list[Path] = []
    seen: set[Path] = set()
    for root in roots:
        try:
            resolved = root.resolve()
        except Exception:
            resolved = root
        if resolved in seen or not root.exists():
            continue
        seen.add(resolved)
        unique_roots.append(root)
    return unique_roots


def _model_name_from_training_metadata(
    run_folder: Path,
    env_name: str,
    env_slug: str,
) -> str | None:
    """Troba un training_config.json que coincideixi i retorna la base del fitxer del model."""

    config_candidates: list[tuple[int, float, str]] = []
    progress_mtime = _safe_mtime(run_folder / "progress.csv")

    # Busquem training_config.json propers a la run. Si hi ha diversos candidats, guanya
    # el que coincideix millor amb l'env_id i amb una data de modificació més propera.
    for root in _candidate_model_roots(run_folder):
        if root.name in _IGNORED_RUN_PARTS:
            continue
        try:
            config_paths = (
                root.rglob("training_config.json")
                if root.name in {"trainings", "models"}
                else root.glob("training_config.json")
            )
            for config_path in config_paths:
                try:
                    with open(config_path, "r", encoding="utf-8") as handle:
                        payload = json.load(handle)
                except Exception:
                    continue

                payload_env = str(payload.get("env_id") or "")
                payload_slug = _slugify_model_lookup(payload_env)
                if payload_env != env_name and payload_slug != env_slug:
                    continue

                artifact_name = str(payload.get("artifact_name") or config_path.parent.name)
                files = payload.get("files") if isinstance(payload.get("files"), dict) else {}
                model_file = str(files.get("model") or f"{artifact_name}.zip")
                model_stem = Path(model_file).stem or artifact_name
                model_path = config_path.parent / model_file
                distance = abs(_safe_mtime(model_path if model_path.exists() else config_path) - progress_mtime)
                score = 2 if payload_env == env_name else 1
                config_candidates.append((score, -distance, model_stem))
        except Exception:
            continue

    if not config_candidates:
        return None

    config_candidates.sort(reverse=True)
    return config_candidates[0][2]


def _model_name_from_zip_search(
    run_folder: Path,
    env_name: str,
    env_slug: str,
) -> str | None:
    """Troba un model zip probable quan les metadades no estiguin disponibles."""

    env_text = env_name.lower()
    zip_candidates: list[tuple[int, float, str]] = []
    progress_mtime = _safe_mtime(run_folder / "progress.csv")

    for root in _candidate_model_roots(run_folder):
        try:
            zip_paths = (
                root.rglob("*.zip")
                if root.name in {"trainings", "models"}
                else root.glob("*.zip")
            )
            for zip_path in zip_paths:
                stem_lower = zip_path.stem.lower()
                stem_slug = _slugify_model_lookup(zip_path.stem)
                score = 0
                if env_slug and env_slug in stem_slug:
                    score = 2
                elif env_text and env_text in stem_lower:
                    score = 1
                if score <= 0:
                    continue
                distance = abs(_safe_mtime(zip_path) - progress_mtime)
                zip_candidates.append((score, -distance, zip_path.stem))
        except Exception:
            continue

    if not zip_candidates:
        return None

    zip_candidates.sort(reverse=True)
    return zip_candidates[0][2]


def _safe_mtime(path: Path) -> float:
    """Retorna un fitxer mtime, o 0 quan no estigui disponible."""

    try:
        return path.stat().st_mtime
    except OSError:
        return 0.0


def load_comparison_data(run_path: str) -> DashboardData | None:
    """Carrega les dades del panell per a una execució de comparació, retornant None quan no és vàlid."""

    artifacts = get_run_artifacts(run_path)
    if artifacts.error_message:
        return None
    try:
        return load_dashboard_data(artifacts.progress_path, artifacts.observations_path)
    except Exception:
        return None


def prepare_obs_time_index(obs: pd.DataFrame) -> pd.DataFrame:
    """Retorna les observacions amb un DatetimeIndex coherent per a visualitzacions de sèries temporals reals."""

    return _ensure_datetime_index(obs)


def build_real_period_options(obs: pd.DataFrame, period_kind: str) -> list[tuple[str, str]]:
    """Prepara opcions de dia, setmana o mes per a les sèries temporals en brut."""

    df = prepare_obs_time_index(obs)
    if df.empty or not isinstance(df.index, pd.DatetimeIndex):
        return []

    if period_kind == "day":
        days = pd.Index(df.index.floor("D").unique()).sort_values()
        return [
            (day.strftime("%Y-%m-%d"), day.strftime("%d/%m/%Y"))
            for day in days
        ]

    if period_kind == "week":
        normalized = df.index.normalize()
        week_starts = pd.Index(
            normalized - pd.to_timedelta(normalized.dayofweek, unit="D")
        ).unique().sort_values()
        return [
            (
                start.strftime("%Y-%m-%d"),
                f"{start.strftime('%d/%m/%Y')} - {(start + pd.Timedelta(days=6)).strftime('%d/%m/%Y')}",
            )
            for start in week_starts
        ]

    months = pd.Index(df.index.to_period("M").unique()).sort_values()
    return [
        (month.strftime("%Y-%m"), month.strftime("%m/%Y"))
        for month in months
    ]


def slice_obs_for_real_period(
    obs: pd.DataFrame,
    period_kind: str,
    selected_value: str,
    *,
    allow_fallback: bool = False,
) -> tuple[pd.DataFrame, str | None]:
    """Filtreu les observacions al període de dia, setmana o mes seleccionat."""

    df = prepare_obs_time_index(obs)
    if df.empty or not selected_value:
        return df.iloc[0:0].copy(), None

    if period_kind == "day":
        target = pd.Timestamp(selected_value).normalize()
        mask = df.index.floor("D") == target
        if not mask.any() and allow_fallback:
            # Les runs de comparació poden tenir un any diferent; en aquest cas comparem
            # el mateix dia i mes, que és el que l'usuari espera visualment.
            mask = (df.index.month == target.month) & (df.index.day == target.day)
        result = df.loc[mask].copy().sort_index(kind="mergesort")
        if result.empty:
            return result, target.strftime("%d/%m/%Y")
        actual_day = result.index[0].normalize()
        return result, actual_day.strftime("%d/%m/%Y")

    if period_kind == "week":
        target = pd.Timestamp(selected_value).normalize()
        week_start = df.index.normalize() - pd.to_timedelta(df.index.dayofweek, unit="D")
        mask = week_start == target
        if not mask.any() and allow_fallback:
            iso_target = target.isocalendar()
            iso_current = df.index.isocalendar()
            mask = iso_current.week.astype(int) == int(iso_target.week)
        result = df.loc[mask].copy().sort_index(kind="mergesort")
        if result.empty:
            end = target + pd.Timedelta(days=6)
            return result, f"{target.strftime('%d/%m/%Y')} - {end.strftime('%d/%m/%Y')}"
        actual_start = result.index.min().normalize()
        actual_start = actual_start - pd.Timedelta(days=int(actual_start.dayofweek))
        actual_end = actual_start + pd.Timedelta(days=6)
        return result, f"{actual_start.strftime('%d/%m/%Y')} - {actual_end.strftime('%d/%m/%Y')}"

    target = pd.Period(selected_value, freq="M")
    mask = df.index.to_period("M") == target
    if not mask.any() and allow_fallback:
        mask = df.index.month == int(target.month)
    result = df.loc[mask].copy().sort_index(kind="mergesort")
    if result.empty:
        return result, target.strftime("%m/%Y")
    actual_month = result.index[0].to_period("M")
    return result, actual_month.strftime("%m/%Y")


def apply_real_period_axis(
    fig: go.Figure,
    obs: pd.DataFrame,
    period_kind: str,
) -> go.Figure:
    """Mapeja els rastres de figures de sèries temporals en brut en un eix de temps relatiu compacte."""

    if not has_figure_data(fig):
        return fig

    df = prepare_obs_time_index(obs)
    if df.empty or not isinstance(df.index, pd.DatetimeIndex):
        return fig

    start = df.index.min()
    elapsed = pd.Series(
        (df.index - start).total_seconds(),
        index=df.index,
        dtype=float,
    )

    if period_kind == "day":
        relative_x = elapsed / 3600.0
        markers = pd.Series(df.index.floor("h"), index=df.index)
        start_flags = markers.ne(markers.shift())
        tickvals = relative_x[start_flags].tolist()
        ticktext = [stamp.strftime("%H:%M") for stamp in markers[start_flags].tolist()]
        axis_title = "Hores del dia"
    elif period_kind == "week":
        relative_x = elapsed / 86400.0
        markers = pd.Series(df.index.floor("D"), index=df.index)
        start_flags = markers.ne(markers.shift())
        tickvals = relative_x[start_flags].tolist()
        ticktext = [f"Dia {idx}" for idx in range(1, len(tickvals) + 1)]
        axis_title = "Dies del període"
    else:
        relative_x = elapsed / 86400.0
        markers = pd.Series(df.index.floor("D"), index=df.index)
        start_flags = markers.ne(markers.shift())
        tickvals = relative_x[start_flags].tolist()
        ticktext = [str(stamp.day) for stamp in markers[start_flags].tolist()]
        axis_title = "Dies del mes"

    if not tickvals:
        return fig

    relative_x_values = relative_x.tolist()
    period_start_x_values = relative_x[start_flags].tolist()
    hour_markers = pd.Series(df.index.floor("h"), index=df.index)
    hour_start_flags = hour_markers.ne(hour_markers.shift())
    hour_start_x_values = relative_x[hour_start_flags].tolist()
    for trace in fig.data:
        try:
            # Cada figura pot venir amb granularitats diferents: punts raw, agregats per hora
            # o agregats pel període. Triem l'eix relatiu mirant quants punts té la traça.
            x_values = getattr(trace, "x", None)
            if x_values is None:
                continue
            if len(x_values) == len(relative_x_values):
                trace.x = relative_x_values
            elif len(x_values) == len(hour_start_x_values):
                trace.x = hour_start_x_values
            elif len(x_values) == len(period_start_x_values):
                trace.x = period_start_x_values
        except Exception:
            continue

    fig.update_xaxes(
        title=axis_title,
        tickmode="array",
        tickvals=tickvals,
        ticktext=ticktext,
        range=[float(relative_x.iloc[0]), float(relative_x.iloc[-1])],
    )
    return fig


def flatten_trace_values(values: Any) -> list[Any]:
    """Retorna valors escalars de les matrius de traça Plotly possiblement imbricades."""

    # Plotly pot guardar dades com list, Series, ndarray o fins i tot DataFrame.
    # Ho aplanem tot per poder decidir si una figura té valors numèrics reals.
    if values is None:
        return []

    if isinstance(values, (str, bytes)):
        return [values]

    if isinstance(values, pd.DataFrame):
        return values.to_numpy().ravel().tolist()

    if isinstance(values, pd.Series):
        return values.tolist()

    if hasattr(values, "ravel"):
        try:
            return values.ravel().tolist()
        except Exception:
            pass

    if hasattr(values, "to_numpy"):
        try:
            return values.to_numpy().ravel().tolist()
        except Exception:
            pass

    try:
        iterator = list(values)
    except TypeError:
        return [values]

    flattened: list[Any] = []
    for item in iterator:
        # Les traces de heatmap solen venir imbricades; fem recursió només amb llistes
        # i tuples per no tractar strings com a seqüències de caràcters.
        if isinstance(item, (str, bytes)):
            flattened.append(item)
            continue
        if isinstance(item, pd.Series):
            flattened.extend(item.tolist())
            continue
        if isinstance(item, pd.DataFrame):
            flattened.extend(item.to_numpy().ravel().tolist())
            continue
        if hasattr(item, "ravel"):
            try:
                flattened.extend(item.ravel().tolist())
                continue
            except Exception:
                pass
        if isinstance(item, (list, tuple)):
            flattened.extend(flatten_trace_values(item))
            continue
        flattened.append(item)

    return flattened


def trace_has_numeric_values(trace: go.BaseTraceType) -> bool:
    """Retorna si una traça Plotly conté almenys un valor numèric."""

    for attribute in ("y", "z", "r", "values"):
        values = flatten_trace_values(getattr(trace, attribute, None))
        if not values:
            continue
        numeric_values = pd.to_numeric(pd.Series(values), errors="coerce").dropna()
        if not numeric_values.empty:
            return True
    return False


def has_figure_data(fig: go.Figure) -> bool:
    """Retorna si una figura Plotly té almenys un rastre amb valors traçables."""

    if fig is None:
        return False

    traces = getattr(fig, "data", [])
    if not traces:
        return False

    return any(trace_has_numeric_values(trace) for trace in traces)


def has_action_figure_data(fig: go.Figure) -> bool:
    """Retorna si una figura d'accions d'agent conté una acció real diferent de zero."""

    if not has_figure_data(fig):
        return False

    for trace in getattr(fig, "data", []):
        values = flatten_trace_values(getattr(trace, "y", None))
        if not values:
            continue
        numeric_values = pd.to_numeric(pd.Series(values), errors="coerce").dropna()
        if not numeric_values.empty and bool((numeric_values.abs() > 0).any()):
            return True

    return False


def add_comfort_percentage_kpi(
    kpi_dict: dict,
    df: pd.DataFrame,
    comfort_config: Any = None,
) -> None:
    """Afegeix un percentatge d'hores de confort KPI al diccionari KPI proporcionat al seu lloc."""

    violation_column = next(
        (column for column in df.columns if re.search(r"(?i)temp[_]?violation", column)),
        None,
    )
    if violation_column is not None:
        series = pd.to_numeric(df[violation_column], errors="coerce").dropna()
        if not series.empty:
            kpi_dict["% Hores en Confort"] = round(float((series <= 1e-9).mean() * 100), 1)
            return

    try:
        from backend.grafics.comfort import _ensure_temp_violation_column

        with_violation = _ensure_temp_violation_column(df, comfort_config)
        violation = pd.to_numeric(with_violation["temp_violation"], errors="coerce").dropna()
        if not violation.empty:
            kpi_dict["% Hores en Confort"] = round(float((violation <= 1e-9).mean() * 100), 1)
            return
    except Exception:
        pass

    average_violation = kpi_dict.get("Avg Temp Violation (All Hours, C)")
    if average_violation is not None:
        kpi_dict["% Hores en Confort"] = (
            100.0 if float(average_violation) == 0.0 else None
        )


def overlay_comparison_traces(
    main_fig: go.Figure,
    comp_fig: go.Figure,
    label_suffix: str = " (ref)",
) -> go.Figure:
    """Superposeu traces de referència en una figura principal quan el tipus de gràfic ho admet."""

    if not has_figure_data(comp_fig):
        return main_fig

    main_barmode = (main_fig.layout.barmode or "").lower()
    if main_barmode in ("stack", "relative"):
        return main_fig

    if any(getattr(trace, "type", "") == "pie" for trace in main_fig.data) or any(
        getattr(trace, "type", "") == "pie" for trace in comp_fig.data
    ):
        return main_fig
    if any(getattr(trace, "type", "") == "heatmap" for trace in main_fig.data) or any(
        getattr(trace, "type", "") == "heatmap" for trace in comp_fig.data
    ):
        return main_fig

    result = go.Figure(main_fig)
    main_bar_index = 0
    for trace in result.data:
        if isinstance(trace, go.Bar):
            # Les barres de la run principal i la referència han d'anar en grups separats,
            # no apilades, perquè el delta sigui llegible d'un cop d'ull.
            trace.update(
                offsetgroup=f"main-{main_bar_index}",
                alignmentgroup="comparison-bars",
            )
            main_bar_index += 1
        elif isinstance(trace, (go.Scatter, go.Scattergl)) and getattr(trace, "stackgroup", None):
            trace.update(stackgroup=f"{trace.stackgroup}-main")

    has_bar_traces = any(isinstance(trace, go.Bar) for trace in comp_fig.data)
    comp_bar_index = 0

    for trace in comp_fig.data:
        trace_copy = copy.deepcopy(trace)
        original_name = getattr(trace_copy, "name", None) or ""
        trace_copy.name = original_name + label_suffix

        if isinstance(trace_copy, (go.Scatter, go.Scattergl)):
            # Les línies de referència van puntejades per no confondre-les amb la run actual.
            if getattr(trace_copy, "stackgroup", None):
                trace_copy.update(stackgroup=f"{trace_copy.stackgroup}-ref")
            trace_copy.update(line=dict(dash="dot", width=3.0), opacity=0.92)
        elif isinstance(trace_copy, go.Bar):
            trace_copy.update(
                marker_line_width=1,
                offsetgroup=f"ref-{comp_bar_index}",
                alignmentgroup="comparison-bars",
                opacity=0.58,
            )
            comp_bar_index += 1
        else:
            trace_copy.update(opacity=0.55)

        result.add_trace(trace_copy)

    if has_bar_traces:
        result.update_layout(barmode="group", bargap=0.18, bargroupgap=0.08)

    return result


def _ensure_observation_time_columns(obs: pd.DataFrame) -> pd.DataFrame:
    """Assegureu-vos que les observacions incloguin camps de mes i hora quan hi hagi marques de temps."""

    if "month" in obs.columns and "hour" in obs.columns:
        return obs

    if "timestamp" in obs.columns:
        timestamps = pd.to_datetime(obs["timestamp"], errors="coerce")
    elif "datetime" in obs.columns:
        timestamps = pd.to_datetime(obs["datetime"], errors="coerce")
    else:
        return obs

    result = obs.copy()
    result["month"] = timestamps.dt.month
    result["hour"] = timestamps.dt.hour
    result.index = timestamps
    return result
\end{lstlisting}

\Needspace{8\baselineskip}

\subsection{\texorpdfstring{\texttt{backend/resultats\_report\_backend.py}}{backend/resultats\_report\_backend.py}}

\begin{lstlisting}[style=studioPython,caption={backend/resultats\_report\_backend.py},label={lst:backend-resultats-report-backend-py}]
"""Generació d'informes HTML i PDF per a les execucions de Studio completades.

El backend de l'informe comparteix els carregadors del panell de resultats i la figura Plotly
builders, després organitza els gràfics disponibles en un document estructurat. Figures
sense dades utilitzables es salten, de manera que els informes exportats reflecteixen els artefactes
produït realment per una tirada.
"""

from __future__ import annotations

import base64
import math
from dataclasses import dataclass
from datetime import datetime
from html import escape
from numbers import Number

import plotly.graph_objects as go

from backend.grafics.comfort_scope import comfort_scope_label
from backend.grafics.battery import (
    make_battery_charge_with_price_plot,
    make_battery_power_plot,
    make_battery_soc_plot,
    make_battery_vs_grid_plot,
    make_energy_price_plot,
)
from backend.grafics.comfort import (
    _fix_colors_for_visibility,
    make_comfort_compliance,
    make_violation_bars,
)
from backend.grafics.control import make_agent_actions_plot, make_radiant_control_plot
from backend.grafics.episode import make_episode_metrics
from backend.grafics.heatmaps import (
    make_heatmap,
    make_violation_heatmap_percent,
    make_zone_temperature_heatmaps,
)
from backend.grafics.hvac import make_hvac_consumption_plot, make_hvac_meter_breakdown_plot
from backend.grafics.style import style_figure_semantics
from backend.grafics.thermal import (
    make_indoor_humidity_plot,
    make_indoor_temperature_plot,
    make_setpoints_plot,
    make_setpoints_vs_indoor_plot,
)
from backend.grafics.figures_zones import filter_obs_by_zone
from backend.grafics.kpis import compute_kpis
from backend.resultats_backend import (
    DashboardData,
    add_comfort_percentage_kpi,
    has_action_figure_data,
    has_figure_data,
    load_dashboard_data,
    select_actions_for_obs,
    select_radiant_action_data,
)
from page_styles.reports import REPORT_HTML_CSS


@dataclass(frozen=True)
class ReportFigure:
    """Plotly gràfic més la secció d'informe on hauria d'aparèixer."""

    section: str
    title: str
    figure: go.Figure


def generate_report_bytes(
    progress_path: str,
    observations_path: str,
    agg_mode: str,
    season: str,
    comfort_scope: str,
    zone_str: str,
) -> tuple[bytes, str]:
    """Crea un informe que es pot descarregar per a una execució completada.

    La funció retorna PDF bytes quan ``pdfkit`` i ``wkhtmltopdf`` són
    disponible. Si la representació estàtica PDF no està disponible, retorna un informe HTML
    amb el mateix contingut perquè UI encara pugui oferir una exportació completa.
    """

    data = load_dashboard_data(progress_path, observations_path)
    selected_zone = None if zone_str in ("ALL", "") else zone_str
    zone_obs = filter_obs_by_zone(data.obs, selected_zone, data.yaml_cfg)
    action_data = select_actions_for_obs(data.actions, zone_obs)
    action_data = select_radiant_action_data(action_data, data)

    kpis = compute_kpis(
        obs=zone_obs,
        cost_hourly=data.metrics_dict.get("cost_hourly"),
        selected_zone=selected_zone,
        comfort_scope=comfort_scope,
        comfort_config=data.yaml_cfg,
    )
    kpis.pop("Avg Energy Cost (EUR/h)", None)
    add_comfort_percentage_kpi(kpis, zone_obs, data.yaml_cfg)

    figures = _build_report_figures(
        data=data,
        zone_obs=zone_obs,
        action_data=action_data,
        agg_mode=agg_mode,
        season=season,
        comfort_scope=comfort_scope,
    )
    return _render_report_document(
        data=data,
        kpis=kpis,
        figures=figures,
        agg_mode=agg_mode,
        season=season,
        selected_zone=selected_zone,
        comfort_scope=comfort_scope,
    )


def _build_report_figures(
    data: DashboardData,
    zone_obs,
    action_data,
    agg_mode: str,
    season: str,
    comfort_scope: str,
) -> tuple[ReportFigure, ...]:
    """Prepara les figures incloses en un informe de resultats exportat."""

    # L'informe reutilitza les mateixes factories del dashboard per evitar que PDF i UI
    # expliquin històries diferents.
    comfort_mode = "raw" if agg_mode == "day" else agg_mode
    pivot = data.metrics_dict.get("pivot_consumption")
    action_fig = make_agent_actions_plot(action_data, agg_mode, season)
    if not has_action_figure_data(action_fig):
        action_fig = go.Figure()

    return (
        ReportFigure(
            "Confort i clima interior",
            "Confort (%)",
            make_comfort_compliance(
                zone_obs,
                comfort_mode,
                season,
                comfort_scope=comfort_scope,
                comfort_config=data.yaml_cfg,
            ),
        ),
        ReportFigure(
            "Confort i clima interior",
            "Temperatura interior vs exterior",
            make_indoor_temperature_plot(zone_obs, agg_mode, season),
        ),
        ReportFigure(
            "Confort i clima interior",
            "Humitat interior vs exterior",
            make_indoor_humidity_plot(zone_obs, agg_mode, season),
        ),
        ReportFigure(
            "Confort i clima interior",
            "Infraccions de confort",
            _fix_colors_for_visibility(
                make_violation_bars(
                    zone_obs,
                    agg_mode,
                    season,
                    comfort_scope=comfort_scope,
                    comfort_config=data.yaml_cfg,
                ),
                kind="violation",
                mode=agg_mode,
            ),
        ),
        ReportFigure(
            "Confort i clima interior",
            "Mapa de calor d'infraccions de confort (%)",
            make_violation_heatmap_percent(
                zone_obs,
                season=season,
                comfort_config=data.yaml_cfg,
            ),
        ),
        ReportFigure(
            "Confort i clima interior",
            "Temperatura per zona (mes x hora)",
            make_zone_temperature_heatmaps(
                zone_obs,
                season=season,
                agg="mean",
                max_cols=3,
            ),
        ),
        ReportFigure(
            "Consum i energia",
            "Consum HVAC (kWh)",
            _fix_colors_for_visibility(
                make_hvac_consumption_plot(zone_obs, agg_mode, season),
                kind="hvac",
                mode=agg_mode,
            ),
        ),
        ReportFigure(
            "Consum i energia",
            "Desglossament HVAC per meter (kWh)",
            make_hvac_meter_breakdown_plot(zone_obs, agg_mode, season),
        ),
        ReportFigure(
            "Consum i energia",
            "Preu energia (EUR/kWh)",
            make_energy_price_plot(zone_obs, agg_mode, season),
        ),
        ReportFigure(
            "Consum i energia",
            "Mapa de calor global (consum)",
            (
                make_heatmap(pivot, "Total HVAC Consumption (kWh)")
                if pivot is not None and not pivot.empty
                else go.Figure()
            ),
        ),
        ReportFigure(
            "Control HVAC",
            "Actuació tèrmica (setpoints)",
            make_setpoints_plot(zone_obs, agg_mode, season),
        ),
        ReportFigure(
            "Control HVAC",
            "Setpoints vs temperatura interior",
            make_setpoints_vs_indoor_plot(zone_obs, agg_mode, season),
        ),
        ReportFigure(
            "Control HVAC",
            "Control radiant",
            make_radiant_control_plot(zone_obs, agg_mode, season),
        ),
        ReportFigure("Control HVAC", "Accions de l'agent", action_fig),
        ReportFigure(
            "Bateria i xarxa",
            "Càrrega/descàrrega bateria",
            make_battery_power_plot(zone_obs, agg_mode, season),
        ),
        ReportFigure(
            "Bateria i xarxa",
            "SOC bateria",
            make_battery_soc_plot(zone_obs, agg_mode, season),
        ),
        ReportFigure(
            "Bateria i xarxa",
            "Energia bateria vs xarxa",
            make_battery_vs_grid_plot(zone_obs, agg_mode, season),
        ),
        ReportFigure(
            "Bateria i xarxa",
            "Bateria vs preu energia",
            make_battery_charge_with_price_plot(zone_obs, agg_mode, season),
        ),
        ReportFigure(
            "Episodi",
            "Mètriques d'episodi",
            make_episode_metrics(data.progress),
        ),
    )


def _render_report_document(
    data: DashboardData,
    kpis: dict,
    figures: tuple[ReportFigure, ...],
    agg_mode: str,
    season: str,
    selected_zone: str | None,
    comfort_scope: str,
) -> tuple[bytes, str]:
    """Genera l'informe HTML i el converteix a PDF quan pdfkit està disponible."""

    # Muntem primer HTML complet: és fàcil de revisar al navegador i després pdfkit el pot
    # convertir sense haver de mantenir una plantilla paral·lela.
    zone_label = selected_zone if selected_zone else "Totes"
    scope_label = comfort_scope_label(comfort_scope)
    rendered_figures, chart_count, static_charts = _render_report_figures(figures)
    generated_at = datetime.now().strftime("%d/%m/%Y %H:%M")
    timesteps_label = f"{data.timesteps_num:,}".replace(",", ".")

    html_parts: list[str] = [f"""<!DOCTYPE html>
<html>
<head>
  <meta charset="utf-8">
  <title>Sinergym Dashboard Report</title>
  {REPORT_HTML_CSS}
</head>
<body>
  <main class="report-shell">
  <section class="cover">
    <p class="eyebrow">BEMS-RL Studio · Informe de resultats</p>
    <h1>Resum {escape(data.run_mode)} - Sinergym</h1>
    <p class="subtitle">
      Informe generat automàticament amb els indicadors principals i tots els gràfics
      disponibles per a l'execució seleccionada.
    </p>
    <div class="summary-grid">
      <div class="summary-item">
        <div class="summary-label">Entorn</div>
        <div class="summary-value">{escape(data.env_name)}</div>
      </div>
      <div class="summary-item">
        <div class="summary-label">Timesteps</div>
        <div class="summary-value">{timesteps_label}</div>
      </div>
      <div class="summary-item">
        <div class="summary-label">Model</div>
        <div class="summary-value">{escape(data.model_name)}</div>
      </div>
      <div class="summary-item">
        <div class="summary-label">Gràfics inclosos</div>
        <div class="summary-value">{chart_count}</div>
      </div>
    </div>
  </section>

  <section class="filter-strip">
    <div class="filter-item">
      <div class="filter-label">Generat</div>
      <div class="filter-value">{generated_at}</div>
    </div>
    <div class="filter-item">
      <div class="filter-label">Agregació</div>
      <div class="filter-value">{escape(agg_mode.upper())}</div>
    </div>
    <div class="filter-item">
      <div class="filter-label">Estació</div>
      <div class="filter-value">{escape(season)}</div>
    </div>
    <div class="filter-item">
      <div class="filter-label">Zona / confort</div>
      <div class="filter-value">{escape(zone_label)} · {escape(scope_label)}</div>
    </div>
  </section>

  <h2 class="block-title">Indicadors clau</h2>
  <div class="kpi-row">"""]

    for key, value in kpis.items():
        value_string = _format_kpi_value(key, value)
        html_parts.append(
            f'<div class="kpi"><div class="kpi-val">{escape(value_string)}</div>'
            f'<div class="kpi-lbl">{escape(key)}</div></div>'
        )

    html_parts.append("</div>")
    if rendered_figures:
        html_parts.extend(rendered_figures)
    else:
        html_parts.append(
            '<div class="empty-note">No hi ha gràfics amb dades disponibles per als filtres seleccionats.</div>'
        )
    html_parts.append("</main></body></html>")
    html = "\n".join(html_parts)

    if static_charts:
        try:
            import pdfkit

            options = {
                "page-size": "A4",
                "margin-top": "0.45in",
                "margin-right": "0.45in",
                "margin-bottom": "0.45in",
                "margin-left": "0.45in",
                "encoding": "UTF-8",
                "enable-local-file-access": None,
                "disable-smart-shrinking": None,
                "quiet": "",
            }
            return pdfkit.from_string(html, False, options=options), "pdf"
        except Exception:
            pass

    return html.encode("utf-8"), "html"


def _format_kpi_value(key: str, value) -> str:
    """Formata els valors KPI de manera coherent per a l'informe exportat."""

    if value is None:
        return "N/D"
    if isinstance(value, Number):
        numeric = float(value)
        if not math.isfinite(numeric):
            return "N/D"
        if key.strip().startswith("%"):
            return f"{numeric:.1f}%"
        return f"{numeric:.2f}"
    return str(value)


def _render_report_figures(
    figures: tuple[ReportFigure, ...],
) -> tuple[list[str], int, bool]:
    """Retorna fragments seccionats de HTML per a cada figura d'informe amb dades visibles."""

    rendered: list[str] = []
    current_section: str | None = None
    chart_count = 0
    static_charts = True
    plotlyjs_included = False

    for item in figures:
        if not has_figure_data(item.figure):
            continue

        chart_html, is_static, plotlyjs_used = _render_single_report_figure(
            item,
            include_plotlyjs=not plotlyjs_included,
        )
        if chart_html is None:
            continue

        if item.section != current_section:
            if current_section is not None:
                rendered.append("</section>")
            current_section = item.section
            rendered.append(
                f'<section class="report-section"><h2 class="section-heading">'
                f"{escape(item.section)}</h2>"
            )

        rendered.append(chart_html)
        chart_count += 1
        static_charts = static_charts and is_static
        plotlyjs_included = plotlyjs_included or plotlyjs_used

    if current_section is not None:
        rendered.append("</section>")

    return rendered, chart_count, static_charts


def _render_single_report_figure(
    item: ReportFigure,
    *,
    include_plotlyjs: bool,
) -> tuple[str | None, bool, bool]:
    """Converteix una figura a imatge estàtica o, si cal, a HTML interactiu."""

    export_fig = _prepare_export_figure(item.figure)
    try:
        image_bytes = export_fig.to_image(format="png", engine="kaleido", scale=2)
        image_b64 = base64.b64encode(image_bytes).decode()
        return (
            f'<div class="chart"><h3>{escape(item.title)}</h3>'
            f'<img alt="{escape(item.title)}" src="data:image/png;base64,{image_b64}"></div>',
            True,
            False,
        )
    except Exception:
        try:
            chart_html = export_fig.to_html(
                full_html=False,
                include_plotlyjs=True if include_plotlyjs else False,
                config={"displayModeBar": False, "responsive": True},
            )
            return (
                f'<div class="chart"><h3>{escape(item.title)}</h3>{chart_html}</div>',
                False,
                include_plotlyjs,
            )
        except Exception:
            return None, False, False


def _prepare_export_figure(fig: go.Figure) -> go.Figure:
    """Copia i normalitzeu una figura Plotly per a la representació d'informes."""

    export_fig = style_figure_semantics(go.Figure(fig))
    current_height = export_fig.layout.height
    try:
        height = int(current_height) if current_height else 520
    except (TypeError, ValueError):
        height = 520
    height = max(460, min(height, 980))
    export_fig.update_layout(
        width=1100,
        height=height,
        margin=dict(l=66, r=36, t=78, b=58),
        font=dict(size=13, color="#17233c"),
        title=dict(
            text=export_fig.layout.title.text,
            x=0.01,
            xanchor="left",
            font=dict(size=18, color="#17233c"),
        ),
        legend=dict(
            orientation="h",
            yanchor="bottom",
            y=1.02,
            xanchor="right",
            x=1,
            bgcolor="rgba(255,255,255,0.92)",
            bordercolor="#d7deeb",
            borderwidth=1,
        ),
    )
    return export_fig
\end{lstlisting}

\Needspace{8\baselineskip}

\subsection{\texorpdfstring{\texttt{backend/sb3\_utils.py}}{backend/sb3\_utils.py}}

\begin{lstlisting}[style=studioPython,caption={backend/sb3\_utils.py},label={lst:backend-sb3-utils-py}]
"""Descoberta, càrrega i format d'accions per a models Stable-Baselines3.

Els fluxos de treball d'avaluació i control en directe utilitzen aquest mòdul per localitzar polítiques desades,
carregar fitxers zip SB3, adjuntar estadístiques VecNormalize, deduir Sinergym compatibles
entorns i convertir les accions en representacions amigables amb UI.
"""

from __future__ import annotations

import io
import re
from pathlib import Path
from typing import Any, Callable, Dict, List, Optional, Tuple

import numpy as np
import pandas as pd
from gymnasium.spaces import Box, Discrete, MultiBinary, MultiDiscrete
from stable_baselines3 import A2C, DDPG, DQN, PPO, SAC, TD3
from stable_baselines3.common.base_class import BaseAlgorithm
from stable_baselines3.common.env_util import make_vec_env
from stable_baselines3.common.save_util import load_from_zip_file
from stable_baselines3.common.vec_env import VecEnv, VecNormalize

from backend.common import list_registered_env_ids


ALGOS: Dict[str, type] = {"PPO": PPO, "A2C": A2C, "DQN": DQN, "SAC": SAC, "TD3": TD3, "DDPG": DDPG}
VECNORM_SUFFIXES: List[str] = ["vecnormalize.pkl", "vecnorm.pkl", "vec_norm.pkl", "normalize.pkl"]


def scan_model_zips(dir_list_text: str) -> pd.DataFrame:
    """Escaneja directoris separats per comes per trobar models Stable-Baselines3 ``.zip``."""
    rows: List[Dict[str, Any]] = []
    for chunk in dir_list_text.split(","):
        root = Path(chunk.strip())
        if not root.exists():
            continue
        for path in root.rglob("*.zip"):
            rows.append(
                {
                    "name": path.name,
                    "stem": path.stem,
                    "path": str(path.resolve()),
                    "dir": str(path.parent.resolve()),
                }
            )
    if not rows:
        return pd.DataFrame(columns=["name", "stem", "path", "dir"])
    return pd.DataFrame(rows).sort_values(by=["stem", "name"]).reset_index(drop=True)


def candidate_vecnorm(model_path: Path) -> Optional[Path]:
    """Retorna el fitxer d'estadístiques VecNormalize més probable al costat d'un zip de model."""
    folder = model_path.parent
    stem = model_path.stem.lower()
    for suffix in VECNORM_SUFFIXES:
        for candidate in [folder / f"{model_path.stem}_{suffix}", folder / f"{stem}_{suffix}"]:
            if candidate.exists():
                return candidate
    for file_path in folder.glob("*.pkl"):
        if any(file_path.name.lower().endswith(suffix) for suffix in VECNORM_SUFFIXES):
            return file_path
    return None


def _algo_cls_from_meta(meta: Dict[str, Any]) -> Optional[type]:
    """Resol la classe d'algorisme SB3 declarada a les metadades del model."""
    algo = (meta or {}).get("algo", None)
    if isinstance(algo, str) and algo.upper() in ALGOS:
        return ALGOS[algo.upper()]
    return None


def load_sb3_model_bytes(model_bytes: bytes, device: str = "cpu") -> Tuple[BaseAlgorithm, Dict[str, Any]]:
    """Carrega un model Stable-Baselines3 des de bytes ZIP i retorna les metadades del model."""
    buffer = io.BytesIO(model_bytes)
    meta: Dict[str, Any] = {}
    try:
        meta, _, _ = load_from_zip_file(buffer, device=device, print_system_info=False)
    except Exception:
        meta = {}

    algo_cls = _algo_cls_from_meta(meta)
    buffer.seek(0)
    model: Optional[BaseAlgorithm] = None

    if algo_cls is not None:
        model = algo_cls.load(buffer, device=device)
    else:
        for algo_type in ALGOS.values():
            try:
                buffer.seek(0)
                model = algo_type.load(buffer, device=device)
                break
            except Exception:
                continue

    if model is None:
        raise ValueError("No s'ha pogut carregar el model SB3.")

    return model, meta


def env_id_from_meta_or_name(meta: Dict[str, Any], zip_stem: str) -> Optional[str]:
    """Dedueix un identificador d'entorn Sinergym registrat a partir de les metadades del model o el nom del fitxer."""
    candidates = list_registered_env_ids()
    for key in ("env_id", "env", "gym_id", "environment"):
        value = meta.get(key)
        if isinstance(value, str) and value in candidates:
            return value

    match = re.search(r"(Eplus-[A-Za-z0-9_\-]+-v\d+)", zip_stem)
    if match:
        return match.group(1)

    tokens = [token for token in re.split(r"[_\-]+", zip_stem.lower()) if len(token) >= 3]
    best, best_score = None, 0
    for env_id in candidates:
        score = sum(token in env_id.lower() for token in tokens)
        if score > best_score:
            best, best_score = env_id, score
    return best


def build_monitored_vec_env(env_factory: Callable[[], Any], *, seed: int | None = None) -> VecEnv:
    """Crea un ``VecEnv`` d'un sol entorn delegant ``Monitor`` i seeds a SB3."""
    return make_vec_env(env_factory, n_envs=1, seed=seed)


def build_vecnormalize_env(
    env_factory: Callable[[], Any],
    *,
    seed: int | None = None,
    norm_obs: bool = True,
    norm_reward: bool = True,
    training: bool = True,
) -> VecNormalize:
    """Crea un ``VecNormalize`` sobre el ``VecEnv`` monitoritzat de SB3."""
    venv = VecNormalize(
        build_monitored_vec_env(env_factory, seed=seed),
        norm_obs=norm_obs,
        norm_reward=norm_reward,
    )
    venv.training = training
    if not training:
        venv.norm_reward = False
    return venv


def load_vecnormalize(vec_path: str, env_factory: Callable[[], Any]) -> VecNormalize:
    """Carrega les estadístiques de VecNormalize sobre un entorn vectoritzat acabat de crear."""
    venv = build_monitored_vec_env(env_factory)
    try:
        venv = VecNormalize.load(vec_path, venv)
    except AssertionError as exc:
        msg = str(exc)
        if "spaces must have the same shape" in msg:
            import re as _re
            m = _re.search(r"\((\d+),\).*?\((\d+),\)", msg)
            detail = f" ({m.group(1)} vs {m.group(2)})" if m else ""
            raise ValueError(
                f"L'espai d'observacio del VecNormalize no coincideix amb l'entorn seleccionat{detail}. "
                "Assegura't d'avaluar el model amb el mateix entorn amb que va ser entrenat."
            ) from exc
        raise
    venv.training = False
    venv.norm_reward = False
    return venv


def _unwrap_action_space(vecenv: VecEnv):
    """Unwrap action space."""
    return vecenv.action_space


def _space_shape_tuple(space: Any) -> Optional[Tuple[int, ...]]:
    """Space shape tuple."""
    shape = getattr(space, "shape", None)
    if shape is None:
        return None
    try:
        return tuple(int(value) for value in shape)
    except Exception:
        return None


def is_normalized_box_space(space: Any) -> bool:
    """Is normalized box space."""
    if not isinstance(space, Box):
        return False
    low = np.asarray(space.low, dtype=np.float32)
    high = np.asarray(space.high, dtype=np.float32)
    return bool(np.allclose(low, -1.0) and np.allclose(high, 1.0))


def action_spaces_compatible(expected_space: Any, actual_space: Any) -> bool:
    """Espais d'acció compatibles."""
    if expected_space is None or actual_space is None:
        return True
    # Comparar només el tipus no és suficient: en Box importen forma i límits; en
    # MultiDiscrete importen els nvec. Això evita acceptar models quasi compatibles.
    if isinstance(expected_space, Box) and isinstance(actual_space, Box):
        if _space_shape_tuple(expected_space) != _space_shape_tuple(actual_space):
            return False
        return bool(
            np.allclose(expected_space.low, actual_space.low)
            and np.allclose(expected_space.high, actual_space.high)
        )
    if isinstance(expected_space, Discrete) and isinstance(actual_space, Discrete):
        return int(expected_space.n) == int(actual_space.n)
    if isinstance(expected_space, MultiDiscrete) and isinstance(actual_space, MultiDiscrete):
        return bool(np.array_equal(expected_space.nvec, actual_space.nvec))
    if isinstance(expected_space, MultiBinary) and isinstance(actual_space, MultiBinary):
        return _space_shape_tuple(expected_space) == _space_shape_tuple(actual_space)
    return type(expected_space) is type(actual_space)


def should_apply_normalize_action(model_action_space: Any, env_action_space: Any) -> bool:
    """S'ha d'aplicar l'acció de normalització."""
    if not isinstance(model_action_space, Box) or not isinstance(env_action_space, Box):
        return False
    if _space_shape_tuple(model_action_space) != _space_shape_tuple(env_action_space):
        return False
    return is_normalized_box_space(model_action_space) and not is_normalized_box_space(env_action_space)


def describe_action_space(space: Any) -> str:
    """Describe action space."""
    if isinstance(space, Box):
        return f"Box(shape={space.shape}, low={np.asarray(space.low).tolist()}, high={np.asarray(space.high).tolist()})"
    if isinstance(space, Discrete):
        return f"Discrete(n={space.n})"
    if isinstance(space, MultiDiscrete):
        return f"MultiDiscrete(nvec={np.asarray(space.nvec).tolist()})"
    if isinstance(space, MultiBinary):
        return f"MultiBinary(shape={space.shape})"
    return type(space).__name__


def _as_action_array(action: Any) -> np.ndarray:
    """Converteix una acció SB3 en ndarray sense assumir encara la forma final."""
    if isinstance(action, np.ndarray):
        return action
    if isinstance(action, (list, tuple)):
        return np.asarray(action)
    return np.asarray([action])


def _reshape_action_batch(
    action_values: np.ndarray,
    action_shape: Tuple[int, ...],
    n_envs: int,
    dtype: Any,
) -> np.ndarray:
    """Porta una acció plana a la forma ``(n_envs, *action_shape)`` esperada per VecEnv."""
    action_dimension = int(np.prod(action_shape))
    flat_action = action_values.astype(dtype).reshape(-1)

    # SB3 pot retornar una sola acció, una acció per entorn o una forma intermèdia.
    # Aquí fem que totes acabin amb la mateixa estructura abans d'aplicar límits.
    if flat_action.size == 1 and action_dimension > 1:
        flat_action = np.repeat(flat_action.item(), action_dimension)
    elif flat_action.size > action_dimension and flat_action.size % action_dimension == 0:
        pass
    elif flat_action.size > action_dimension:
        flat_action = flat_action[:action_dimension]
    elif flat_action.size < action_dimension:
        fill_value = flat_action[-1] if flat_action.size else 0
        flat_action = np.concatenate(
            [flat_action, np.repeat(fill_value, action_dimension - flat_action.size)]
        )

    action_batch = flat_action.reshape(-1, *action_shape)
    if action_batch.shape[0] == 1 and n_envs > 1:
        action_batch = np.repeat(action_batch, n_envs, axis=0)
    elif action_batch.shape[0] > n_envs:
        action_batch = action_batch[:n_envs]
    elif action_batch.shape[0] < n_envs:
        action_batch = np.concatenate(
            [action_batch] + [action_batch[-1:]] * (n_envs - action_batch.shape[0]),
            axis=0,
        )
    return action_batch.astype(dtype)


def format_action_for_env(action: Any, vecenv: VecEnv) -> np.ndarray:
    """
    Converteix l'acció que retorna el model al format esperat pel VecEnv:
      - Discrete: array shape (n_envs,), dtype=int64
      - Box: array shape (n_envs, act_dim), dtype=float32 (clip a [low, high])
    """
    action_space = _unwrap_action_space(vecenv)
    n_envs = getattr(vecenv, "num_envs", 1)

    # SB3 pot retornar escalar, llista o ndarray segons l'algorisme; primer ho normalitzem
    # per no tenir quatre camins diferents fent el mateix.
    normalized_action = _as_action_array(action)

    if isinstance(action_space, Discrete):
        # En Discrete cada entorn paral·lel espera un sol enter; si falta algun valor,
        # repetim l'últim perquè VecEnv no falli per forma incorrecta.
        normalized_action = normalized_action.reshape(-1)
        if normalized_action.size == 1 and n_envs > 1:
            normalized_action = np.repeat(int(normalized_action.item()), n_envs)
        if normalized_action.size > n_envs:
            normalized_action = normalized_action[:n_envs]
        if normalized_action.size < n_envs:
            normalized_action = np.pad(normalized_action, (0, n_envs - normalized_action.size), mode="edge")
        normalized_action = normalized_action.astype(np.int64)
        normalized_action = np.clip(normalized_action, 0, action_space.n - 1)
        return normalized_action

    if isinstance(action_space, MultiDiscrete):
        # MultiDiscrete és el cas més fàcil de trencar: l'acció pot venir plana o ja
        # agrupada per entorns. La portem a la forma (n_envs, *shape) abans de retallar.
        nvec = np.asarray(action_space.nvec, dtype=np.int64)
        action_shape = tuple(int(value) for value in nvec.shape)
        normalized_action = _reshape_action_batch(normalized_action, action_shape, n_envs, np.int64)
        return np.clip(normalized_action, 0, nvec - 1).astype(np.int64)

    if isinstance(action_space, MultiBinary):
        # MultiBinary comparteix la mateixa idea que MultiDiscrete, però només accepta 0/1.
        # El clip final evita que un model carregat amb petites diferències retorni valors fora de rang.
        action_shape = tuple(int(value) for value in action_space.shape)
        normalized_action = _reshape_action_batch(normalized_action, action_shape, n_envs, np.int64)
        return np.clip(normalized_action, 0, 1).astype(np.int8)

    if isinstance(action_space, Box):
        # En espais continus mantenim float32 i fem clip als límits reals de l'entorn.
        # Això és especialment important quan hi ha wrappers que canvien l'escala de l'acció.
        action_shape = tuple(int(value) for value in action_space.shape)
        normalized_action = _reshape_action_batch(normalized_action, action_shape, n_envs, np.float32)
        low = np.array(action_space.low, dtype=np.float32).reshape(1, *action_shape)
        high = np.array(action_space.high, dtype=np.float32).reshape(1, *action_shape)
        normalized_action = np.clip(normalized_action, low, high).astype(np.float32)
        return normalized_action

    # Fallback conservador per a espais poc habituals: es prioritza que el VecEnv rebi
    # una mida coherent, encara que després el mateix entorn acabi rebutjant l'acció.
    normalized_action = normalized_action.reshape(-1)
    if normalized_action.size == 1 and n_envs > 1:
        normalized_action = np.repeat(normalized_action.item(), n_envs)
    elif normalized_action.size > n_envs:
        normalized_action = normalized_action[:n_envs]
    elif normalized_action.size < n_envs:
        normalized_action = np.pad(normalized_action, (0, n_envs - normalized_action.size), mode="edge")
    return normalized_action.astype(np.int64)
\end{lstlisting}

\Needspace{8\baselineskip}

\subsection{\texorpdfstring{\texttt{backend/simular\_entorn\_backend.py}}{backend/simular\_entorn\_backend.py}}

\begin{lstlisting}[style=studioPython,caption={backend/simular\_entorn\_backend.py},label={lst:backend-simular-entorn-backend-py}]
"""Simulacions de referència per a execucions no controlades o basades en regles.

La pàgina de simulació utilitza aquest backend per inspeccionar entorns registrats, executar
episodis de referència, capturen el progrés i persisteixen la mateixa forma d'artefacte esperada
pel panell de resultats. Això ofereix als usuaris una prova de referència abans de comparar
agents formats.
"""

from __future__ import annotations

import csv
import glob
import os
import time
from dataclasses import dataclass
from pathlib import Path
from typing import Callable

import gymnasium as gym
import numpy as np
from stable_baselines3.common.monitor import Monitor

from sinergym.utils.rewards import BaseReward
from sinergym.utils.wrappers import CSVLogger, LoggerWrapper

from backend.entrenar_agent_constants import DETAILED_HVAC_METERS, JOULES_PER_KWH
from backend.entrenar_agent_env import with_detailed_hvac_meters
from backend.common import list_registered_env_ids

NO_CONTROL_ACTION = np.empty((0,), dtype=np.float32)
BASELINE_TIMESTEPS_PER_HOUR = 4
BASELINE_STEP_MINUTES = 60.0 / BASELINE_TIMESTEPS_PER_HOUR


@dataclass(frozen=True)
class EnvironmentSummary:
    """Descripció compacta d'un entorn Sinergym registrat."""
    env_id: str
    env_name: str
    building_file: str
    weather_count: int
    variables_count: int
    meters_count: int
    actuators_count: int
    reward_name: str
    step_minutes: float | None


@dataclass(frozen=True)
class BaselineSimulationResult:
    """Sortides persistents i metadades d'estat per a una simulació de referència."""
    env_id: str
    run_path: str
    progress_path: str
    observations_path: str
    elapsed_seconds: float
    completed_steps: int
    cancelled: bool


class ScheduleBaselineReward(BaseReward):
    """Reward zero que conserva les mètriques físiques per comparar contra agents entrenats."""

    def __init__(
        self,
        temperature_variables: list[str] | None = None,
        energy_variables: list[str] | None = None,
        range_comfort_winter: tuple[float, float] = (20.0, 23.5),
        range_comfort_summer: tuple[float, float] = (23.0, 26.0),
        summer_start: tuple[int, int] = (6, 1),
        summer_final: tuple[int, int] = (9, 30),
    ) -> None:
        """Inicialitza la instància."""
        super().__init__()
        self.temperature_variables = list(temperature_variables or [])
        self.energy_variables = list(energy_variables or [])
        self.range_comfort_winter = tuple(range_comfort_winter)
        self.range_comfort_summer = tuple(range_comfort_summer)
        self.summer_start = tuple(summer_start)
        self.summer_final = tuple(summer_final)

    def __call__(self, obs_dict: dict[str, float]) -> tuple[float, dict[str, float]]:
        """Executa l'objecte cridable."""
        total_power_demand = self._compute_power_demand(obs_dict)
        total_temperature_violation = self._compute_temperature_violation(obs_dict)
        reward_terms = {
            "energy_term": 0.0,
            "comfort_term": 0.0,
            "energy_penalty": -total_power_demand,
            "comfort_penalty": -total_temperature_violation,
            "total_power_demand": total_power_demand,
            "total_temperature_violation": total_temperature_violation,
            "reward_weight": 0.0,
        }
        return 0.0, reward_terms

    def _compute_power_demand(self, obs_dict: dict[str, float]) -> float:
        """Calcula la demanda de potència."""
        total = 0.0
        for variable in self.energy_variables:
            value = obs_dict.get(variable)
            if value is None:
                continue
            try:
                total += float(value)
            except (TypeError, ValueError):
                continue
        return total

    def _compute_temperature_violation(self, obs_dict: dict[str, float]) -> float:
        """Calcula la violació de temperatura."""
        month = _safe_int(obs_dict.get("month"), default=1)
        day = _safe_int(obs_dict.get("day_of_month"), default=1)
        comfort_range = (
            self.range_comfort_summer
            if _is_within_summer((month, day), self.summer_start, self.summer_final)
            else self.range_comfort_winter
        )

        violation = 0.0
        for variable in self.temperature_variables:
            value = obs_dict.get(variable)
            if value is None:
                continue
            try:
                temperature = float(value)
            except (TypeError, ValueError):
                continue
            lower, upper = comfort_range
            violation += max(lower - temperature, 0.0) + max(temperature - upper, 0.0)
        return violation


def list_environment_ids() -> list[str]:
    """Retorna els entorns Sinergym que tenen sentit per a una línia de base de planificació."""

    return list_registered_env_ids(include_discrete=False)


def describe_environment(env_id: str) -> EnvironmentSummary:
    """Crea un resum compacte utilitzat per la interfície."""

    spec = gym.spec(env_id)
    kwargs = spec.kwargs or {}
    meters = with_detailed_hvac_meters(kwargs.get("meters") or {})
    weather_files = kwargs.get("weather_files") or []
    if isinstance(weather_files, str):
        weather_count = 1
    else:
        weather_count = len(weather_files)

    reward_cls = kwargs.get("reward")
    reward_name = getattr(reward_cls, "__name__", str(reward_cls))
    return EnvironmentSummary(
        env_id=env_id,
        env_name=str(kwargs.get("env_name") or env_id),
        building_file=str(kwargs.get("building_file") or ""),
        weather_count=weather_count,
        variables_count=len(kwargs.get("variables") or {}),
        meters_count=len(meters),
        actuators_count=len(kwargs.get("actuators") or {}),
        reward_name=reward_name,
        step_minutes=BASELINE_STEP_MINUTES,
    )


def run_baseline_simulation(
    env_id: str,
    steps_target: int,
    seed: int | None,
    should_stop: Callable[[], bool],
    on_progress: Callable[[int, int], None],
) -> BaselineSimulationResult:
    """Executa una línia de base sense control utilitzant els planificacions per defecte d'EnergyPlus."""

    env = _make_baseline_env(env_id)
    workspace_path = ""
    completed_steps = 0
    cancelled = False
    start_time = time.time()

    try:
        env.reset(seed=seed)
        workspace_path = str(env.get_wrapper_attr("workspace_path"))

        for step_number in range(steps_target):
            if should_stop():
                cancelled = True
                break

            _, _, terminated, truncated, _ = env.step(NO_CONTROL_ACTION)
            completed_steps = step_number + 1

            if terminated or truncated:
                env.reset()

            if completed_steps % 200 == 0 or completed_steps == steps_target:
                on_progress(completed_steps, steps_target)
    finally:
        try:
            env.close()
        except Exception:
            pass

    progress_path, observations_path = _resolve_run_artifacts(workspace_path)
    _append_detailed_meter_kwh_columns(observations_path)
    return BaselineSimulationResult(
        env_id=env_id,
        run_path=workspace_path,
        progress_path=progress_path,
        observations_path=observations_path,
        elapsed_seconds=time.time() - start_time,
        completed_steps=completed_steps,
        cancelled=cancelled,
    )


def _make_baseline_env(env_id: str):
    """Crea l'entorn de baseline."""
    env_summary = describe_environment(env_id)
    spec = gym.spec(env_id)
    spec_kwargs = dict(spec.kwargs or {})
    baseline_reward_kwargs = _build_baseline_reward_kwargs(spec_kwargs)
    baseline_name = f"{env_summary.env_name}-baseline"
    building_config = dict(spec_kwargs.get("building_config") or {})
    building_config["timesteps_per_hour"] = BASELINE_TIMESTEPS_PER_HOUR
    meters = with_detailed_hvac_meters(spec_kwargs.get("meters") or {})
    env = gym.make(
        env_id,
        action_space=gym.spaces.Box(low=0.0, high=0.0, shape=(0,), dtype=np.float32),
        actuators={},
        building_config=building_config,
        env_name=baseline_name,
        meters=meters,
        reward=ScheduleBaselineReward,
        reward_kwargs=baseline_reward_kwargs,
    )
    env = LoggerWrapper(env)
    env = CSVLogger(env)
    env = Monitor(env)
    return env


def _resolve_run_artifacts(run_path: str) -> tuple[str, str]:
    """Retorna les sortides principals CSV de l'execució que acaba d'acabar."""

    progress_path = str(Path(run_path) / "progress.csv")
    observation_files = sorted(
        glob.glob(os.path.join(run_path, "**", "observations.csv"), recursive=True),
        key=os.path.getmtime,
        reverse=True,
    )
    observations_path = observation_files[0] if observation_files else ""
    return progress_path, observations_path


def _append_detailed_meter_kwh_columns(observations_path: str) -> None:
    """Afegeix columnes detallades del comptador de kWh."""
    if not observations_path or not os.path.exists(observations_path):
        return

    with open(observations_path, "r", newline="", encoding="utf-8") as file_handle:
        reader = csv.DictReader(file_handle)
        fieldnames = list(reader.fieldnames or [])
        rows = list(reader)

    if not fieldnames or not rows:
        return

    # Els meters detallats de Sinergym poden quedar en joules al CSV; afegim *_kwh una
    # sola vegada perquè el dashboard posterior no faci la conversió repetida.
    derived_columns = {
        alias: f"{alias}_kwh"
        for alias in DETAILED_HVAC_METERS
        if alias in fieldnames and f"{alias}_kwh" not in fieldnames
    }
    if not derived_columns:
        return

    for row in rows:
        for source_column, target_column in derived_columns.items():
            try:
                row[target_column] = str(float(row.get(source_column, "")) / JOULES_PER_KWH)
            except (TypeError, ValueError):
                row[target_column] = ""

    output_fields = fieldnames + list(derived_columns.values())
    with open(observations_path, "w", newline="", encoding="utf-8") as file_handle:
        writer = csv.DictWriter(file_handle, fieldnames=output_fields)
        writer.writeheader()
        writer.writerows(rows)


def _build_baseline_reward_kwargs(env_kwargs: dict) -> dict[str, object]:
    """Crea kwargs de recompensa de referència."""
    reward_kwargs = dict(env_kwargs.get("reward_kwargs") or {})
    variables = env_kwargs.get("variables") or {}
    meters = env_kwargs.get("meters") or {}

    # La baseline no optimitza res, però igualment necessita variables de confort i energia
    # perquè els gràfics comparatius surtin amb les mateixes mètriques que un agent.
    temperature_variables = reward_kwargs.get("temperature_variables")
    if not temperature_variables:
        temperature_variables = [
            name for name in variables.keys()
            if "air_temperature" in name.lower()
        ]

    energy_variables = reward_kwargs.get("energy_variables")
    if not energy_variables:
        energy_candidates = list(variables.keys()) + list(meters.keys())
        energy_variables = [
            name for name in energy_candidates
            if any(token in name.lower() for token in ("hvac", "electric", "power", "demand", "energy"))
        ]

    return {
        "temperature_variables": list(temperature_variables or []),
        "energy_variables": list(energy_variables or []),
        "range_comfort_winter": tuple(reward_kwargs.get("range_comfort_winter") or (20.0, 23.5)),
        "range_comfort_summer": tuple(reward_kwargs.get("range_comfort_summer") or (23.0, 26.0)),
        "summer_start": tuple(reward_kwargs.get("summer_start") or (6, 1)),
        "summer_final": tuple(reward_kwargs.get("summer_final") or (9, 30)),
    }


def _safe_int(value: object, *, default: int) -> int:
    """Safe int."""
    try:
        return int(float(value))
    except (TypeError, ValueError):
        return default


def _is_within_summer(current: tuple[int, int], start: tuple[int, int], end: tuple[int, int]) -> bool:
    """Is within summer."""
    return start <= current <= end
\end{lstlisting}

\Needspace{8\baselineskip}

\subsection{\texorpdfstring{\texttt{backend/training\_scripts.py}}{backend/training\_scripts.py}}

\begin{lstlisting}[style=studioPython,caption={backend/training\_scripts.py},label={lst:backend-training-scripts-py}]
"""Scripts d'ajuda reproduïbles per a execucions BEMS-RL desades."""

from __future__ import annotations

from pprint import pformat
from typing import Any, Dict


def build_training_eval_script(training_config: Dict[str, Any]) -> str:
    """Crea un script breu d'avaluació al costat de la sessió d'entrenament."""
    lines = [
        "# Script d'avaluacio autogenerat (no s'executa automaticament)",
        "# Entorn i recompensa amb la mateixa configuracio que l'entrenament.",
        "",
        f"env_id = {training_config['env_id']!r}",
        f"meters = {pformat(training_config.get('meters', {}), sort_dicts=False)}",
        f"reward_name = {training_config['reward_name']!r}",
        f"reward_kwargs = {pformat(training_config['reward_kwargs'], sort_dicts=False)}",
        f"wrapper_configs = {pformat(training_config.get('wrapper_configs', []), sort_dicts=False)}",
        "",
    ]
    return "\n".join(lines)


def build_training_repro_script(
    artifact_name: str,
    training_config: Dict[str, Any],
) -> str:
    """Crea un script autònom que pugui reproduir una configuració d'entrenament desada."""
    config_literal = pformat(training_config, sort_dicts=False, width=100)
    # Retornem el fitxer complet com a string perquè l'artefacte guardat sigui executable
    # fora de Streamlit, amb la mateixa configuració que s'ha fet servir a la UI.
    return f'''"""Script d'entrenament reproduible autogenerat per BEMS-RL-STUDIO."""

import csv
from pathlib import Path

import gymnasium as gym
import numpy as np
from gymnasium import Wrapper
from gymnasium.spaces import Box
from stable_baselines3 import A2C, DDPG, DQN, PPO, SAC, TD3
from stable_baselines3.common.env_util import make_vec_env
from stable_baselines3.common.vec_env import VecNormalize
import sinergym.utils.rewards as reward_module
import sinergym.utils.wrappers as wrapper_module


TRAINING_NAME = {artifact_name!r}
TRAINING_CONFIG = {config_literal}
ALGOS = {{"PPO": PPO, "A2C": A2C, "DQN": DQN, "SAC": SAC, "TD3": TD3, "DDPG": DDPG}}
DETAILED_HVAC_METERS = {{
    "natural_gas_hvac": "NaturalGas:HVAC",
    "heating_electricity": "Heating:Electricity",
    "cooling_electricity": "Cooling:Electricity",
    "fans_electricity": "Fans:Electricity",
    "pumps_electricity": "Pumps:Electricity",
    "heat_rejection_electricity": "HeatRejection:Electricity",
    "humidifier_electricity": "Humidifier:Electricity",
    "heat_recovery_electricity": "HeatRecovery:Electricity",
}}
JOULES_PER_KWH = 3_600_000.0


class EnergyCostFileLogger(Wrapper):
    def __init__(self, env, filename="energy_cost_log.csv"):
        super().__init__(env)
        self.csvfile = open(filename, mode="w", newline="")
        self.writer = csv.writer(self.csvfile)
        self.writer.writerow(["step", "electricity_cost_eur_per_kwh", "energy_cost_eur"])
        self.step_count = 0

    def step(self, action):
        obs, reward, terminated, truncated, info = self.env.step(action)
        self.step_count += 1
        cost_per_kwh = info.get("electricity_cost", 0.0)
        energy_rate = info.get("HVAC_electricity_demand_rate", 0.0)
        self.writer.writerow([self.step_count, cost_per_kwh, cost_per_kwh * energy_rate])
        return obs, reward, terminated, truncated, info

    def reset(self, **kwargs):
        self.step_count = 0
        return self.env.reset(**kwargs)

    def close(self):
        self.csvfile.close()
        if hasattr(self.env, "close"):
            self.env.close()


def get_reward_class(name):
    reward_cls = getattr(reward_module, name, None)
    if reward_cls is None:
        raise ImportError(
            f"La reward '{{name}}' no existeix a sinergym.utils.rewards en aquesta instal-lacio."
        )
    return reward_cls


def get_wrapper_class(name):
    if name == "EnergyCostFileLogger":
        return EnergyCostFileLogger
    wrapper_cls = getattr(wrapper_module, name, None)
    if wrapper_cls is None:
        raise ImportError(
            f"El wrapper '{{name}}' no existeix a sinergym.utils.wrappers en aquesta instal-lacio."
        )
    return wrapper_cls


def apply_training_wrappers(env, wrapper_configs):
    for wrapper_config in wrapper_configs:
        name = wrapper_config["name"]
        params = wrapper_config.get("params", {{}})
        wrapper_cls = get_wrapper_class(name)
        if name == "PreviousObservationWrapper":
            env = wrapper_cls(env, **params)
        elif name == "DatetimeWrapper":
            env = wrapper_cls(env)
        elif name == "WeatherForecastingWrapper":
            env = wrapper_cls(env, **params)
        elif name == "EnergyCostWrapper":
            env = wrapper_cls(env, **params)
        elif name == "DeltaTempWrapper":
            env = wrapper_cls(env, **params)
        elif name == "ReduceObservationWrapper":
            env = wrapper_cls(env, **params)
        elif name == "MultiObsWrapper":
            env = wrapper_cls(env, **params)
        elif name == "NormalizeObservation":
            env = wrapper_cls(env, **params)
        elif name == "VariabilityContextWrapper":
            context_space = Box(
                low=np.array(params["context_low"], dtype=np.float32),
                high=np.array(params["context_high"], dtype=np.float32),
                dtype=np.float32,
            )
            env = wrapper_cls(
                env,
                context_space=context_space,
                delta_value=params["delta_value"],
                step_frequency_range=tuple(params["step_frequency_range"]),
            )
        elif name == "OfficeGridStorageSmoothingActionConstraintsWrapper":
            env = wrapper_cls(env)
        elif name == "IncrementalWrapper":
            env = wrapper_cls(env, **params)
        elif name == "DiscreteIncrementalWrapper":
            env = wrapper_cls(env, **params)
        elif name == "NormalizeAction" and isinstance(env.action_space, Box):
            env = wrapper_cls(env, **params)
        elif name == "EnergyCostFileLogger":
            env = wrapper_cls(env, **params)
    return env


def with_detailed_hvac_meters(meters):
    return dict(meters or {{}})


def get_training_meters():
    configured_meters = TRAINING_CONFIG.get("meters")
    if configured_meters is not None:
        return with_detailed_hvac_meters(configured_meters)

    try:
        spec = gym.spec(TRAINING_CONFIG["env_id"])
        meters = (spec.kwargs or {{}}).get("meters", {{}}) or {{}}
    except Exception:
        meters = {{}}
    return with_detailed_hvac_meters(meters)


def make_env():
    reward_cls = get_reward_class(TRAINING_CONFIG["reward_name"])
    meters = get_training_meters()
    env = gym.make(
        TRAINING_CONFIG["env_id"],
        meters=meters,
        reward=reward_cls,
        reward_kwargs=TRAINING_CONFIG["reward_kwargs"],
    )
    env = apply_training_wrappers(env, TRAINING_CONFIG.get("wrapper_configs", []))
    env = get_wrapper_class("LoggerWrapper")(env)
    env = get_wrapper_class("CSVLogger")(env)
    return env


def get_workspace_path(vec):
    try:
        paths = vec.get_attr("workspace_path")
        if paths and paths[0]:
            return str(paths[0])
    except Exception:
        pass
    return None


def append_detailed_meter_kwh_columns(observations_path):
    observations_file = Path(observations_path) if observations_path else None
    if observations_file is None or not observations_file.is_file():
        return

    with open(observations_file, "r", newline="", encoding="utf-8") as file_handle:
        reader = csv.DictReader(file_handle)
        fieldnames = list(reader.fieldnames or [])
        rows = list(reader)

    if not fieldnames or not rows:
        return

    source_columns = {{}}
    for alias, meter_name in DETAILED_HVAC_METERS.items():
        target_column = alias + "_kwh"
        if target_column in fieldnames:
            continue
        if alias in fieldnames:
            source_columns[alias] = alias
        elif meter_name in fieldnames:
            source_columns[alias] = meter_name

    if not source_columns:
        return

    derived_columns = {{source_alias: source_alias + "_kwh" for source_alias in source_columns}}
    for row in rows:
        for source_alias, source_column in source_columns.items():
            target_column = derived_columns[source_alias]
            try:
                row[target_column] = str(float(row.get(source_column, "")) / JOULES_PER_KWH)
            except (TypeError, ValueError):
                row[target_column] = ""

    output_fields = fieldnames + list(derived_columns.values())
    with open(observations_file, "w", newline="", encoding="utf-8") as file_handle:
        writer = csv.DictWriter(file_handle, fieldnames=output_fields)
        writer.writeheader()
        writer.writerows(rows)


def append_detailed_meter_kwh_columns_in_workspace(workspace_path):
    if not workspace_path:
        return

    workspace = Path(workspace_path)
    if not workspace.is_dir():
        return

    for observations_file in workspace.rglob("observations.csv"):
        try:
            append_detailed_meter_kwh_columns(observations_file)
        except Exception:
            continue


def main() -> None:
    artifact_dir = Path(__file__).resolve().parent
    vec = make_vec_env(make_env, n_envs=1)
    vec = VecNormalize(vec, norm_obs=True, norm_reward=True)
    workspace_path = None

    try:
        algo_cls = ALGOS[TRAINING_CONFIG["algo_name"]]
        algo_kwargs = dict(TRAINING_CONFIG.get("algo_kwargs") or {{}})
        algo_kwargs.setdefault("learning_rate", TRAINING_CONFIG["learning_rate"])
        if TRAINING_CONFIG["algo_name"] in {{"PPO", "A2C"}} and TRAINING_CONFIG.get("n_steps") is not None:
            algo_kwargs.setdefault("n_steps", TRAINING_CONFIG["n_steps"])
        algo_kwargs["verbose"] = 1

        model = algo_cls(TRAINING_CONFIG["policy_name"], vec, **algo_kwargs)
        model.learn(total_timesteps=TRAINING_CONFIG["timesteps_per_year"])
        model.save(str(artifact_dir / f"{{TRAINING_NAME}}.zip"))
        vec.save(str(artifact_dir / f"{{TRAINING_NAME}}_vecnorm.pkl"))
        workspace_path = get_workspace_path(vec)
    finally:
        if workspace_path is None:
            workspace_path = get_workspace_path(vec)
        vec.close()
        append_detailed_meter_kwh_columns_in_workspace(workspace_path)


if __name__ == "__main__":
    main()
'''
\end{lstlisting}

\Needspace{8\baselineskip}

\subsection{\texorpdfstring{\texttt{backend/viewer\_3d\_backend.py}}{backend/viewer\_3d\_backend.py}}

\begin{lstlisting}[style=studioPython,caption={backend/viewer\_3d\_backend.py},label={lst:backend-viewer-3d-backend-py}]
"""Backend del visor d'edificis 3D: processament de geometria i generació de figures Plotly."""

from __future__ import annotations

import math
from typing import Any, Dict, List, Tuple

import numpy as np
import plotly.graph_objects as go
import streamlit as st
import trimesh

SURFACE_STYLES: Dict[str, Dict[str, Any]] = {
    'Wall':      {'color': 'orange',   'opacity': 0.45},
    'Roof':      {'color': 'green',    'opacity': 0.45},
    'Floor':     {'color': 'blue',     'opacity': 0.35},
    'Ceiling':   {'color': 'purple',   'opacity': 0.35},
    'Window':    {'color': 'skyblue',  'opacity': 0.60},
    'Door':      {'color': '#8B4513', 'opacity': 0.90},
    'Glassdoor': {'color': 'cyan',     'opacity': 0.60},
    'Unknown':   {'color': 'darkgray', 'opacity': 0.50},
}
PV_STYLE_COLOR = "#222222"
PV_STYLE_OPACITY = 0.90


def extract_vertices_general(surface: dict) -> List[Tuple[float, float, float]]:
    """Extreu vèrtexs d'un dict de superfície.

    Admet tant una llista de "vèrtexs" com camps de coordenades indexats separats
    (vertex_N_x_coordinate / vertex_N_y_coordinate / vertex_N_z_coordinate).
    """
    vertices = []
    if "vertices" in surface:
        for vertex in surface["vertices"]:
            x = vertex.get("vertex_x_coordinate")
            y = vertex.get("vertex_y_coordinate")
            z = vertex.get("vertex_z_coordinate")
            if None not in (x, y, z):
                vertices.append((x, y, z))
    else:
        index = 1
        while f"vertex_{index}_x_coordinate" in surface:
            x = surface.get(f"vertex_{index}_x_coordinate")
            y = surface.get(f"vertex_{index}_y_coordinate")
            z = surface.get(f"vertex_{index}_z_coordinate")
            if None not in (x, y, z):
                vertices.append((x, y, z))
            index += 1
    return vertices


def _classify_type(surf_type_raw: str, kind: str) -> str:
    """Normalise a raw surface type string to a canonical type label."""
    surface_type = (surf_type_raw or "Unknown").strip().lower()
    if kind == "sub":
        if "door" in surface_type and "glass" in surface_type:
            return "Glassdoor"
        if "door" in surface_type:
            return "Door"
        return "Window"
    if "wall" in surface_type:
        return "Wall"
    if "roof" in surface_type:
        return "Roof"
    if "ceiling" in surface_type:
        return "Ceiling"
    if "floor" in surface_type:
        return "Floor"
    return "Unknown"


def _triangulate_fan(n: int) -> List[Tuple[int, int, int]]:
    """Retorna l'índex de triangulació en ventall triplicat per a un polígon convex amb n vèrtexs."""
    return [(0, i, i + 1) for i in range(1, n - 1)] if n >= 3 else []


def _tilt_azimuth_from_normal(
    nx: float, ny: float, nz: float
) -> Tuple[float, float, str]:
    """Calcula inclinació, azimut i orientació cardinal a partir d'una normal de superfície.

    Retorna:
        inclinació: Angle des de l'horitzontal en graus (0 = pla).
        azimut: en sentit horari des del nord en graus (0=N, 90=E, 180=S, 270=W).
        orientació: etiqueta cardinal ('N', 'NE', ..., 'H' per a gairebé horitzontal).
    """
    norm = float(np.linalg.norm([nx, ny, nz])) or 1e-9
    tilt = math.degrees(math.acos(abs(nz) / norm))
    azimuth = (math.degrees(math.atan2(nx, ny)) + 360.0) % 360.0
    directions = ['N', 'NE', 'E', 'SE', 'S', 'SO', 'O', 'NO']
    index = int(((azimuth + 22.5) % 360) // 45)
    orientation = directions[index]
    if tilt < 15:
        orientation = 'H'
    return tilt, azimuth, orientation


def _polygon_mesh(vertices: List[Tuple[float, float, float]]) -> trimesh.Trimesh:
    """Crea una malla triangular mínima a partir d'una superfície EnergyPlus."""
    return trimesh.Trimesh(
        vertices=np.asarray(vertices, dtype=float),
        faces=np.asarray(_triangulate_fan(len(vertices)), dtype=np.int64),
        process=False,
    )


def _polygon_geometry(vertices: List[Tuple[float, float, float]]) -> Tuple[float, Tuple[float, float, float]]:
    """Retorna àrea i normal mitjana ponderada per àrea d'una superfície."""
    mesh = _polygon_mesh(vertices)
    if len(mesh.faces) == 0:
        return 0.0, (0.0, 0.0, 1.0)
    normal = np.average(mesh.face_normals, axis=0, weights=mesh.area_faces)
    return float(mesh.area), tuple(float(value) for value in trimesh.util.unitize(normal))


def _surface_hover_text(record: Dict[str, Any]) -> str:
    """Crea una etiqueta de flotació llegible pels humans per a un registre de superfície."""
    geometry_line = f"Area: {record['area']:.2f} m2"
    if record['orientation'] == 'H':
        orientation_line = f"Tilt: {record['tilt']:.1f} deg | Orientation: horizontal"
    else:
        orientation_line = (
            f"Tilt: {record['tilt']:.1f} deg | "
            f"Azimuth: {record['azimuth']:.1f} deg ({record['orientation']})"
        )
    return (
        f"Surface: {record['name']}<br>"
        f"Zone: {record['zone']}<br>"
        f"Type: {record['type']} ({record['kind']})<br>"
        f"{geometry_line}<br>"
        f"{orientation_line}"
    )


def _scaled_offset_polygon(
    vertices: List[Tuple[float, float, float]],
    offset: float = 0.05,
    shrink: float = 0.97,
) -> List[Tuple[float, float, float]]:
    """Retorna una còpia reduïda i amb desplaçament normal d'un polígon.

    Paràmetres:
        vèrtexs: vèrtexs de polígons originals.
        desplaçament: desplaçament al llarg de la unitat normal per evitar la lluita en z.
        contracció: factor d'escala [0..1] aplicat en relació al baricentre.
    """
    if not vertices:
        return vertices
    points = np.asarray(vertices, dtype=float)
    centroid = points.mean(axis=0)
    _area, normal = _polygon_geometry(vertices)
    output = []
    for point in points:
        shifted = centroid + shrink * (point - centroid) + np.asarray(normal) * offset
        output.append(tuple(float(value) for value in shifted))
    return output


def _collect_surfaces_with_names(epjson: dict) -> List[Dict[str, Any]]:
    """Recull totes les superfícies d'un epJSON dict.

    Retorna una llista de dicts amb claus: nom, zona, type_raw, tipus ('principal'/'sub'), vèrtexs.
    """
    output = []
    for name, surface in epjson.get("BuildingSurface:Detailed", {}).items():
        vertices = extract_vertices_general(surface)
        if len(vertices) >= 3:
            output.append({
                'name': name,
                'zone': surface.get("zone_name", "Unknown"),
                'type_raw': surface.get("surface_type", "Unknown"),
                'kind': 'main',
                'vertices': vertices,
            })
    for name, surface in epjson.get("FenestrationSurface:Detailed", {}).items():
        vertices = extract_vertices_general(surface)
        if len(vertices) >= 3:
            output.append({
                'name': name,
                'zone': surface.get("building_surface_name", "Unknown"),
                'type_raw': surface.get("surface_type", "Unknown"),
                'kind': 'sub',
                'vertices': vertices,
            })
    return output


def _bounds(
    vertices_list: List[List[Tuple[float, float, float]]],
) -> Tuple[float, float, float, float, float, float]:
    """Retorna (xmin, xmax, ymin, ymax, zmin, zmax) el quadre delimitador de tots els vèrtexs proporcionats."""
    points = [point for vertices in vertices_list for point in vertices]
    if not points:
        return (0.0, 0.0, 0.0, 0.0, 0.0, 0.0)
    bounds = trimesh.points.PointCloud(np.asarray(points, dtype=float)).bounds
    return (
        float(bounds[0][0]),
        float(bounds[1][0]),
        float(bounds[0][1]),
        float(bounds[1][1]),
        float(bounds[0][2]),
        float(bounds[1][2]),
    )


def _color_for(record: Dict[str, Any]) -> str:
    """Retorna el color de visualització d'un registre de superfície en funció del seu tipus canònic."""
    return SURFACE_STYLES.get(record['type'], SURFACE_STYLES['Unknown'])['color']


@st.cache_data(show_spinner=False)
def build_surface_records(epjson: dict) -> Dict[str, Any]:
    """Crea registres de superfície enriquits a partir d'un dic epJSON (emmagatzemat en memòria cau).

    Cada registre conté: id, nom, zona, tipus (normalitzat), tipus (principal/sub),
    vèrtexs originals, vèrtexs centrats, àrea, inclinació, azimut, orientació,
    i la mitjana z en coordenades centrades.

    Retorna un dict amb claus: registres, zones, tipus, centre, z_clip_range, name_index.
    El name_index mapes el nom de la superfície per registrar l'índex només per a les superfícies principals.
    """
    items = _collect_surfaces_with_names(epjson)
    vertices_list = [item['vertices'] for item in items]
    xmn, xmx, ymn, ymx, zmn, zmx = _bounds(vertices_list)
    cx = (xmn + xmx) / 2.0
    cy = (ymn + ymx) / 2.0
    cz = (zmn + zmx) / 2.0

    records: List[Dict[str, Any]] = []
    zones: set = set()
    types: set = set()
    name_index: Dict[str, int] = {}

    for index, item in enumerate(items):
        normalized_type = _classify_type(item['type_raw'], item['kind'])
        zones.add(item['zone'])
        types.add(normalized_type)

        area, (nx, ny, nz) = _polygon_geometry(item['vertices'])
        tilt, azimuth, orientation = _tilt_azimuth_from_normal(nx, ny, nz)
        vertices_c = [(x - cx, y - cy, z - cz) for (x, y, z) in item['vertices']]

        record: Dict[str, Any] = {
            'id': index,
            'name': item['name'],
            'zone': item['zone'],
            'type': normalized_type,
            'kind': item['kind'],
            'vertices': item['vertices'],
            'vertices_c': vertices_c,
            'area': area,
            'tilt': tilt,
            'azimuth': azimuth,
            'orientation': orientation,
            'z_mean_c': sum(v[2] for v in vertices_c) / len(vertices_c),
        }
        records.append(record)
        if item['kind'] == 'main':
            name_index[item['name']] = index

    z_values = [r['z_mean_c'] for r in records] or [0.0]
    return {
        'records': records,
        'zones': sorted(zones),
        'types': sorted(types),
        'center': (cx, cy, cz),
        'z_clip_range': (float(min(z_values)), float(max(z_values))),
        'name_index': name_index,
    }


def plot_3d_model(
    records: List[Dict[str, Any]],
    visible_types: Dict[str, bool],
    visible_zones: Dict[str, bool],
    z_clip: Tuple[float, float],
    pv_list: List[Dict[str, Any]],
    name_index: Dict[str, int],
) -> go.Figure:
    """Construeix una figura 3D Plotly de la geometria de l'edifici.

    Les superfícies es filtren per tipus, zona i rang de clip d'alçada.
    Els panells PV es representen com una capa addicional superposada a les seves superfícies host.
    """
    fig = go.Figure()
    zmin, zmax = z_clip

    for record in records:
        # El filtre es fa abans de crear traces perquè Plotly es torna lent si li passem
        # superfícies ocultes i després només les amaguem a la llegenda.
        if not visible_types.get(record['type'], False):
            continue
        if not visible_zones.get(record['zone'], False):
            continue
        if not (zmin <= record['z_mean_c'] <= zmax):
            continue

        vertices = record['vertices_c']
        x, y, z = zip(*vertices)
        # EnergyPlus dona polígons; Mesh3d necessita triangles. El fan és suficient per a
        # superfícies habituals d'edifici, que normalment són convexes o gairebé planes.
        triangles = _triangulate_fan(len(vertices))
        if not triangles:
            continue
        i_idx, j_idx, k_idx = zip(*triangles)
        color = _color_for(record)
        opacity = SURFACE_STYLES.get(record['type'], SURFACE_STYLES['Unknown'])['opacity']

        fig.add_trace(go.Mesh3d(
            x=x, y=y, z=z,
            i=i_idx, j=j_idx, k=k_idx,
            color=color, opacity=opacity,
            name=f"{record['zone']} - {record['type']}",
            hovertext=_surface_hover_text(record),
            hoverinfo='text',
            showscale=False,
        ))

        if record['kind'] == 'sub':
            fig.add_trace(go.Scatter3d(
                x=list(x) + [x[0]],
                y=list(y) + [y[0]],
                z=list(z) + [z[0]],
                mode='lines',
                line=dict(color='black', width=2),
                name="Subsurface edge",
                showlegend=False,
                hoverinfo='skip',
            ))

    if pv_list and visible_types.get('PV', False):
        for pv in pv_list:
            surface_name = pv.get('surface_name')
            if surface_name not in name_index:
                continue
            record = records[name_index[surface_name]]
            if not (zmin <= record['z_mean_c'] <= zmax):
                continue
            # El panell PV es dibuixa una mica separat i encongit perquè no faci z-fighting
            # amb la superfície de la coberta o façana que el suporta.
            pv_polygon = _scaled_offset_polygon(record['vertices_c'], offset=0.03, shrink=0.96)
            px, py, pz = zip(*pv_polygon)
            triangles = _triangulate_fan(len(pv_polygon))
            if not triangles:
                continue
            i_idx, j_idx, k_idx = zip(*triangles)
            fig.add_trace(go.Mesh3d(
                x=px, y=py, z=pz,
                i=i_idx, j=j_idx, k=k_idx,
                color=PV_STYLE_COLOR, opacity=PV_STYLE_OPACITY,
                name=f"PV - {pv.get('name')}",
                hovertext=(
                    f"Panel: {pv.get('name')}<br>"
                    f"Surface: {surface_name}<br>"
                    f"DC capacity: {pv.get('capacity_dc', '-')}"
                ),
                hoverinfo='text',
                showscale=False,
            ))
            fig.add_trace(go.Scatter3d(
                x=list(px) + [px[0]],
                y=list(py) + [py[0]],
                z=list(pz) + [pz[0]],
                mode='lines',
                line=dict(color='black', width=2),
                name="PV edge",
                showlegend=False,
                hoverinfo='skip',
            ))

    fig.update_layout(
        scene=dict(xaxis_title="X", yaxis_title="Y", zaxis_title="Z", aspectmode='data'),
        margin=dict(l=0, r=0, b=0, t=0),
        height=550,
        showlegend=True,
    )
    return fig
\end{lstlisting}

\Needspace{8\baselineskip}

\subsection{\texorpdfstring{\texttt{backend/visor\_epw\_backend.py}}{backend/visor\_epw\_backend.py}}

\begin{lstlisting}[style=studioPython,caption={backend/visor\_epw\_backend.py},label={lst:backend-visor-epw-backend-py}]
"""Utilitats de preparació de dades per al visor climàtic EPW.

El backend del visor EPW localitza fitxers meteorològics d'EnergyPlus,
normalitza els camps EPW en brut, afegeix atributs de calendari i produeix
taules de resum per a la UI de Streamlit. Les sortides es mantenen com a
DataFrames perquè els gràfics i les descàrregues es puguin construir sense
tornar a llegir el fitxer meteorològic.
"""

from __future__ import annotations

from pathlib import Path

import numpy as np
import pandas as pd
import streamlit as st


ROOT_PATH = Path.cwd().resolve()
WEATHERS_DIR = ROOT_PATH / "sinergym" / "data" / "weather"
WEATHER_ROOT = WEATHERS_DIR if WEATHERS_DIR.exists() else ROOT_PATH
DEG = "\N{DEGREE SIGN}"
DEG_C = f"{DEG}C"
SQUARE_2 = "\N{SUPERSCRIPT TWO}"
WH_PER_M2 = f"Wh/m{SQUARE_2}"
KWH_PER_M2 = f"kWh/m{SQUARE_2}"
MID_DOT = "\N{MIDDLE DOT}"

MONTH_LABELS = {
    1: "Gen",
    2: "Feb",
    3: "Mar",
    4: "Abr",
    5: "Mai",
    6: "Jun",
    7: "Jul",
    8: "Ago",
    9: "Set",
    10: "Oct",
    11: "Nov",
    12: "Des",
}
MONTH_ORDER = list(MONTH_LABELS.values())
WIND_DIRECTION_LABELS = (
    "N",
    "NNE",
    "NE",
    "ENE",
    "E",
    "ESE",
    "SE",
    "SSE",
    "S",
    "SSW",
    "SW",
    "WSW",
    "W",
    "WNW",
    "NW",
    "NNW",
)
TIME_AGGREGATION_RULES = {
    "Hora": None,
    "Dia": "D",
    "Setmana": "W-MON",
    "Mes": "M",
}
CATALOG_SEPARATOR = f" {MID_DOT} "

EPW_COLUMNS = [
    "year",
    "month",
    "day",
    "hour",
    "minute",
    "data_source",
    "dry_bulb_c",
    "dew_point_c",
    "relative_humidity_pct",
    "atmospheric_pressure_pa",
    "extraterrestrial_horizontal_radiation_wh_m2",
    "extraterrestrial_direct_normal_radiation_wh_m2",
    "horizontal_infrared_radiation_wh_m2",
    "global_horizontal_radiation_wh_m2",
    "direct_normal_radiation_wh_m2",
    "diffuse_horizontal_radiation_wh_m2",
    "global_horizontal_illuminance_lux",
    "direct_normal_illuminance_lux",
    "diffuse_horizontal_illuminance_lux",
    "zenith_luminance_cd_m2",
    "wind_direction_deg",
    "wind_speed_m_s",
    "total_sky_cover_tenths",
    "opaque_sky_cover_tenths",
    "visibility_km",
    "ceiling_height_m",
    "present_weather_observation",
    "present_weather_codes",
    "precipitable_water_mm",
    "aerosol_optical_depth_thousandths",
    "snow_depth_cm",
    "days_since_last_snowfall",
    "albedo",
    "liquid_precipitation_depth_mm",
    "liquid_precipitation_quantity_hr",
]
NUMERIC_EPW_COLUMNS = tuple(column for column in EPW_COLUMNS if column != "data_source")

EPW_VARIABLES = {
    "dry_bulb_c": {
        "label": "Temperatura seca",
        "unit": DEG_C,
        "aggregate": "mean",
        "color": "#f26b5b",
    },
    "relative_humidity_pct": {
        "label": "Humitat relativa",
        "unit": "%",
        "aggregate": "mean",
        "color": "#4f8ef7",
    },
    "wind_speed_m_s": {
        "label": "Velocitat del vent",
        "unit": "m/s",
        "aggregate": "mean",
        "color": "#2ca58d",
    },
    "direct_normal_radiation_wh_m2": {
        "label": "Radiació directa normal",
        "unit": WH_PER_M2,
        "aggregate": "sum",
        "color": "#ef7d57",
        "scale_to_kilo": True,
        "display_unit_aggregated": KWH_PER_M2,
    },
}

MISSING_SENTINELS = {
    "dry_bulb_c": {99.9},
    "dew_point_c": {99.9},
    "relative_humidity_pct": {999.0},
    "atmospheric_pressure_pa": {999999.0},
    "extraterrestrial_horizontal_radiation_wh_m2": {9999.0},
    "extraterrestrial_direct_normal_radiation_wh_m2": {9999.0},
    "horizontal_infrared_radiation_wh_m2": {9999.0},
    "global_horizontal_radiation_wh_m2": {9999.0},
    "direct_normal_radiation_wh_m2": {9999.0},
    "diffuse_horizontal_radiation_wh_m2": {9999.0},
    "global_horizontal_illuminance_lux": {999999.0},
    "direct_normal_illuminance_lux": {999999.0},
    "diffuse_horizontal_illuminance_lux": {999999.0},
    "zenith_luminance_cd_m2": {999999.0, 9999.0},
    "wind_direction_deg": {999.0},
    "wind_speed_m_s": {999.0},
    "total_sky_cover_tenths": {99.0},
    "opaque_sky_cover_tenths": {99.0},
    "visibility_km": {9999.0},
    "ceiling_height_m": {99999.0},
    "precipitable_water_mm": {999.0},
    "snow_depth_cm": {999.0},
    "days_since_last_snowfall": {99.0},
    "albedo": {999.0},
    "liquid_precipitation_depth_mm": {999.0},
    "liquid_precipitation_quantity_hr": {99.0},
}

COMFORT_RADIATION_SOURCE_COLUMNS = (
    "dry_bulb_c",
    "direct_normal_radiation_wh_m2",
    "diffuse_horizontal_radiation_wh_m2",
    "day_of_year",
    "hour_start",
)
COMFORT_RADIATION_COLUMNS = (
    "direction",
    "theta_deg",
    "altitude_band",
    "radius_start",
    "radius_size",
    "comfort_radiation_kwh_m2",
    "incident_radiation_kwh_m2",
)


def _parse_optional_float(value: str | None) -> float:
    """Retorna un camp numèric EPW com a flotant, o NaN quan el valor està en blanc o no és vàlid."""

    if value is None or value == "":
        return np.nan
    try:
        return float(value)
    except ValueError:
        return np.nan


def _header_payload(header_lines: list[str], line_index: int) -> str:
    """Retorna la càrrega útil després de la primera coma en una línia de capçalera EPW."""

    if len(header_lines) <= line_index or "," not in header_lines[line_index]:
        return "-"
    return header_lines[line_index].split(",", 1)[1].strip() or "-"


def _classify_climate(annual_mean_temp: float) -> str:
    """Classifica la família climàtica a partir de la temperatura mitjana anual de bulb sec."""

    if annual_mean_temp > 22:
        return "Càlid"
    if annual_mean_temp < 12:
        return "Fred"
    return "Mixt"


def _metadata_coordinates(metadata: dict[str, object]) -> tuple[float, float, float] | None:
    """Retorna valors vàlids de latitud, longitud i zona horària de les metadades EPW."""

    try:
        coordinates = (
            float(metadata.get("latitude")),
            float(metadata.get("longitude")),
            float(metadata.get("timezone")),
        )
    except (TypeError, ValueError):
        return None
    if any(pd.isna(value) for value in coordinates):
        return None
    return coordinates


def variable_label(variable_key: str, *, aggregated: bool = False) -> str:
    """Retorna l'etiqueta UI i la unitat de visualització per a una clau variable EPW."""

    spec = EPW_VARIABLES[variable_key]
    if aggregated and spec.get("display_unit_aggregated"):
        unit = spec["display_unit_aggregated"]
    else:
        unit = spec["unit"]
    return f"{spec['label']} ({unit})" if unit else spec["label"]


def _parse_location_line(line: str) -> dict[str, object]:
    """Analitza la línia de capçalera EPW LOCATION als camps de metadades del lloc."""

    parts = [part.strip() for part in line.strip().split(",")]
    return {
        "city": parts[1] if len(parts) > 1 else "-",
        "state": parts[2] if len(parts) > 2 else "-",
        "country": parts[3] if len(parts) > 3 else "-",
        "source": parts[4] if len(parts) > 4 else "-",
        "wmo": parts[5] if len(parts) > 5 else "-",
        "latitude": _parse_optional_float(parts[6] if len(parts) > 6 else None),
        "longitude": _parse_optional_float(parts[7] if len(parts) > 7 else None),
        "timezone": _parse_optional_float(parts[8] if len(parts) > 8 else None),
        "elevation_m": _parse_optional_float(parts[9] if len(parts) > 9 else None),
    }


def _read_epw_header(epw_path: Path) -> list[str]:
    """Llegeix i torna les vuit línies de capçalera estàndard EPW."""

    with epw_path.open("r", encoding="utf-8", errors="ignore") as file_handle:
        return [file_handle.readline().rstrip("\n") for _ in range(8)]


def _build_epw_metadata(epw_path: Path, header_lines: list[str]) -> dict[str, object]:
    """Crea metadades de fitxer i ubicació a partir de les línies de capçalera EPW."""

    metadata = _parse_location_line(header_lines[0] if header_lines else "")
    metadata.update(
        {
            "file_name": epw_path.name,
            "stem": epw_path.stem,
            "comments_1": _header_payload(header_lines, 5),
            "comments_2": _header_payload(header_lines, 6),
            "data_periods": header_lines[7] if len(header_lines) > 7 else "-",
            "header_lines": header_lines,
        }
    )
    return metadata


def _read_epw_data(epw_path: Path) -> pd.DataFrame:
    """Llegeix les files de dades per hora EPW en un DataFrame sense processar."""

    raw_data_frame = pd.read_csv(
        epw_path,
        skiprows=8,
        header=None,
        names=EPW_COLUMNS,
        encoding="utf-8",
        encoding_errors="ignore",
    )
    return raw_data_frame


def _clean_epw_data(data_frame: pd.DataFrame) -> pd.DataFrame:
    """Converteix EPW columnes numèriques i substituir EPW els sentinelles de valors que falten per NaN."""

    for column in NUMERIC_EPW_COLUMNS:
        data_frame[column] = pd.to_numeric(data_frame[column], errors="coerce")

    for column, sentinels in MISSING_SENTINELS.items():
        if column in data_frame.columns:
            data_frame[column] = data_frame[column].replace(list(sentinels), np.nan)

    return data_frame


def _add_time_features(data_frame: pd.DataFrame) -> pd.DataFrame:
    """Afegeix segells de temps normalitzats i atributs de calendari utilitzades pels gràfics de visualització."""

    data_frame["hour_start"] = data_frame["hour"].fillna(1).clip(1, 24).astype(int) - 1
    data_frame["source_year"] = data_frame["year"].fillna(2001).astype("Int64")
    data_frame["timestamp"] = pd.to_datetime(
        {
            "year": 2001,
            "month": data_frame["month"].fillna(1).astype(int),
            "day": data_frame["day"].fillna(1).astype(int),
            "hour": data_frame["hour_start"],
            "minute": 0,
        },
        errors="coerce",
    )
    data_frame["date"] = data_frame["timestamp"].dt.normalize()
    data_frame["month_label"] = pd.Categorical(
        data_frame["month"].map(MONTH_LABELS),
        categories=MONTH_ORDER,
        ordered=True,
    )
    data_frame["day_of_year"] = data_frame["timestamp"].dt.dayofyear
    return data_frame


def _add_climate_metadata(metadata: dict[str, object], data_frame: pd.DataFrame) -> dict[str, object]:
    """Afegeix mètriques climàtiques derivades a les metadades EPW carregades."""

    annual_mean_temp = float(data_frame["dry_bulb_c"].mean())
    metadata["annual_mean_temp"] = annual_mean_temp
    metadata["climate_family"] = _classify_climate(annual_mean_temp)
    return metadata


@st.cache_data(show_spinner=False)
def list_epw_catalog(root_path_str: str) -> tuple[dict[str, str], ...]:
    """Escaneja un directori arrel i retorna entrades de catàleg cercables per a tots els fitxers EPW."""

    root_path = Path(root_path_str)
    rows: list[dict[str, str]] = []

    for epw_path in sorted(root_path.rglob("*.epw")):
        city = "-"
        country = "-"
        try:
            with epw_path.open("r", encoding="utf-8", errors="ignore") as file_handle:
                first_line = file_handle.readline()
            location = _parse_location_line(first_line)
            city = str(location["city"])
            country = str(location["country"])
        except OSError:
            pass

        relative_path = epw_path.relative_to(root_path).as_posix()
        display_name = f"{city}, {country}{CATALOG_SEPARATOR}{relative_path}" if city != "-" else relative_path
        search_blob = " ".join([display_name, epw_path.stem, relative_path]).lower()
        rows.append(
            {
                "path": str(epw_path),
                "relative_path": relative_path,
                "display_name": display_name,
                "search_blob": search_blob,
            }
        )

    return tuple(rows)


@st.cache_data(show_spinner=False)
def load_epw_bundle(epw_path_str: str) -> dict[str, object]:
    """Carrega un fitxer EPW i retorna'n les metadades més un clima enriquit DataFrame."""

    epw_path = Path(epw_path_str)
    header_lines = _read_epw_header(epw_path)
    data_frame = _read_epw_data(epw_path)
    data_frame = _clean_epw_data(data_frame)
    data_frame = _add_time_features(data_frame)
    metadata = _build_epw_metadata(epw_path, header_lines)
    metadata = _add_climate_metadata(metadata, data_frame)

    return {
        "metadata": metadata,
        "data": data_frame,
    }


def summarize_epw_metrics(data_frame: pd.DataFrame) -> dict[str, float]:
    """Retorna les mètriques del clima, el vent, la humitat, la pressió i la radiació."""

    has_rows = not data_frame.empty
    ghi_total_kwh_m2 = np.nan
    dni_total_kwh_m2 = np.nan
    dhi_total_kwh_m2 = np.nan
    if has_rows:
        ghi_total_kwh_m2 = float(data_frame["global_horizontal_radiation_wh_m2"].sum() / 1000.0)
        dni_total_kwh_m2 = float(data_frame["direct_normal_radiation_wh_m2"].sum() / 1000.0)
        dhi_total_kwh_m2 = float(data_frame["diffuse_horizontal_radiation_wh_m2"].sum() / 1000.0)

    return {
        "records": float(len(data_frame)),
        "mean_temp": float(data_frame["dry_bulb_c"].mean()) if has_rows else np.nan,
        "min_temp": float(data_frame["dry_bulb_c"].min()) if has_rows else np.nan,
        "max_temp": float(data_frame["dry_bulb_c"].max()) if has_rows else np.nan,
        "mean_rh": float(data_frame["relative_humidity_pct"].mean()) if has_rows else np.nan,
        "mean_wind": float(data_frame["wind_speed_m_s"].mean()) if has_rows else np.nan,
        "max_wind": float(data_frame["wind_speed_m_s"].max()) if has_rows else np.nan,
        "ghi_total_kwh_m2": ghi_total_kwh_m2,
        "dni_total_kwh_m2": dni_total_kwh_m2,
        "dhi_total_kwh_m2": dhi_total_kwh_m2,
    }


def build_monthly_summary(data_frame: pd.DataFrame) -> pd.DataFrame:
    """Agrupa les principals variables tèrmiques, d'humitat, de vent i solars per mes."""

    monthly = (
        data_frame.groupby(["month", "month_label"], observed=True)
        .agg(
            dry_bulb_mean=("dry_bulb_c", "mean"),
            dry_bulb_min=("dry_bulb_c", "min"),
            dry_bulb_max=("dry_bulb_c", "max"),
            relative_humidity_mean=("relative_humidity_pct", "mean"),
            global_horizontal_radiation_kwh_m2=("global_horizontal_radiation_wh_m2", "sum"),
            direct_normal_radiation_kwh_m2=("direct_normal_radiation_wh_m2", "sum"),
            diffuse_horizontal_radiation_kwh_m2=("diffuse_horizontal_radiation_wh_m2", "sum"),
        )
        .reset_index()
        .sort_values("month")
    )
    radiation_columns = [
        "global_horizontal_radiation_kwh_m2",
        "direct_normal_radiation_kwh_m2",
        "diffuse_horizontal_radiation_kwh_m2",
    ]
    monthly[radiation_columns] = monthly[radiation_columns] / 1000.0
    return monthly


def aggregate_epw_timeseries(
    data_frame: pd.DataFrame,
    variable_key: str,
    aggregation_label: str,
) -> pd.DataFrame:
    """Agrega una variable EPW a la sèrie horària, diària, setmanal o mensual sol·licitada."""

    series_frame = data_frame[["timestamp", variable_key]].dropna().sort_values("timestamp")
    if series_frame.empty:
        return series_frame

    rule = TIME_AGGREGATION_RULES.get(aggregation_label)
    if not rule:
        return series_frame.reset_index(drop=True)

    variable_spec = EPW_VARIABLES[variable_key]
    aggregate_mode = variable_spec["aggregate"]
    if aggregate_mode == "sum":
        aggregated = series_frame.set_index("timestamp").resample(rule).sum(numeric_only=True).reset_index()
    else:
        aggregated = series_frame.set_index("timestamp").resample(rule).mean(numeric_only=True).reset_index()

    if variable_spec.get("scale_to_kilo"):
        aggregated[variable_key] = aggregated[variable_key] / 1000.0

    return aggregated.reset_index(drop=True)


def build_daily_temperature_profile(data_frame: pd.DataFrame) -> pd.DataFrame:
    """Retorna els valors diaris mínims, mitjans i màxims de temperatura de bulb sec."""

    daily = (
        data_frame.dropna(subset=["date"])
        .groupby("date", observed=False)
        .agg(
            dry_bulb_min=("dry_bulb_c", "min"),
            dry_bulb_mean=("dry_bulb_c", "mean"),
            dry_bulb_max=("dry_bulb_c", "max"),
        )
        .reset_index()
    )
    return daily


def build_hourly_profile(data_frame: pd.DataFrame) -> pd.DataFrame:
    """Retorna perfils horaris mitjans per a les principals variables meteorològiques."""

    hourly = (
        data_frame.groupby("hour_start", observed=False)
        .agg(
            dry_bulb_mean=("dry_bulb_c", "mean"),
            dew_point_mean=("dew_point_c", "mean"),
            relative_humidity_mean=("relative_humidity_pct", "mean"),
        )
        .reset_index()
    )
    return hourly


def build_heatmap_table(data_frame: pd.DataFrame, variable_key: str) -> pd.DataFrame:
    """Crea una taula dinàmica mes a hora per a la variable EPW seleccionada."""

    heatmap = (
        data_frame.pivot_table(
            index="month_label",
            columns="hour_start",
            values=variable_key,
            aggfunc="mean",
            observed=False,
        )
        .reindex(MONTH_ORDER)
    )
    return heatmap


def build_annual_hourly_heatmap_table(data_frame: pd.DataFrame, variable_key: str) -> pd.DataFrame:
    """Crea una taula dinàmica hora a dia de l'any per a la variable EPW seleccionada."""

    heatmap = (
        data_frame.pivot_table(
            index="hour_start",
            columns="day_of_year",
            values=variable_key,
            aggfunc="mean",
            observed=False,
        )
        .sort_index()
        .reindex(index=list(range(24)))
    )
    return heatmap


def _solar_position_frame(
    data_frame: pd.DataFrame,
    latitude_deg: float,
    longitude_deg: float,
    timezone_hours: float,
) -> pd.DataFrame:
    """Estima l'altitud i l'azimut solar per a cada fila horària EPW."""

    solar_frame = data_frame[["day_of_year", "hour_start"]].copy()
    solar_frame = solar_frame.dropna()
    if solar_frame.empty:
        return pd.DataFrame(columns=["altitude_deg", "azimuth_deg"])

    day_of_year = solar_frame["day_of_year"].to_numpy(dtype=float)
    local_hour = solar_frame["hour_start"].to_numpy(dtype=float) + 0.5

    gamma = 2.0 * np.pi / 365.0 * (day_of_year - 1.0 + (local_hour - 12.0) / 24.0)
    equation_of_time = 229.18 * (
        0.000075
        + 0.001868 * np.cos(gamma)
        - 0.032077 * np.sin(gamma)
        - 0.014615 * np.cos(2.0 * gamma)
        - 0.040849 * np.sin(2.0 * gamma)
    )
    declination = (
        0.006918
        - 0.399912 * np.cos(gamma)
        + 0.070257 * np.sin(gamma)
        - 0.006758 * np.cos(2.0 * gamma)
        + 0.000907 * np.sin(2.0 * gamma)
        - 0.002697 * np.cos(3.0 * gamma)
        + 0.00148 * np.sin(3.0 * gamma)
    )

    true_solar_minutes = (local_hour * 60.0 + equation_of_time + 4.0 * longitude_deg - 60.0 * timezone_hours) % 1440.0
    hour_angle_deg = true_solar_minutes / 4.0 - 180.0
    hour_angle_rad = np.deg2rad(hour_angle_deg)
    latitude_rad = np.deg2rad(latitude_deg)

    cos_zenith = (
        np.sin(latitude_rad) * np.sin(declination)
        + np.cos(latitude_rad) * np.cos(declination) * np.cos(hour_angle_rad)
    )
    cos_zenith = np.clip(cos_zenith, -1.0, 1.0)
    zenith_rad = np.arccos(cos_zenith)
    altitude_deg = np.rad2deg(np.pi / 2.0 - zenith_rad)

    azimuth_deg = (
        np.rad2deg(
            np.arctan2(
                np.sin(hour_angle_rad),
                np.cos(hour_angle_rad) * np.sin(latitude_rad) - np.tan(declination) * np.cos(latitude_rad),
            )
        )
        + 180.0
    ) % 360.0

    return pd.DataFrame(
        {
            "altitude_deg": altitude_deg,
            "azimuth_deg": azimuth_deg,
        },
        index=solar_frame.index,
    )


def build_comfort_radiation_table(data_frame: pd.DataFrame, metadata: dict[str, object]) -> pd.DataFrame:
    """Estima la radiació incident per orientació, ponderada pel benefici de calefacció o refrigeració."""

    empty_result = pd.DataFrame(columns=COMFORT_RADIATION_COLUMNS)
    coordinates = _metadata_coordinates(metadata)
    if coordinates is None:
        return empty_result.copy()
    latitude, longitude, timezone = coordinates

    radiation_frame = data_frame[list(COMFORT_RADIATION_SOURCE_COLUMNS)].dropna().copy()
    if radiation_frame.empty:
        return empty_result.copy()

    solar_frame = _solar_position_frame(radiation_frame, latitude, longitude, timezone)
    if solar_frame.empty:
        return empty_result.copy()

    radiation_frame = radiation_frame.join(solar_frame)
    radiation_frame = radiation_frame[radiation_frame["altitude_deg"] > 0.0].copy()
    if radiation_frame.empty:
        return empty_result.copy()

    altitude_edges = np.array([0.0, 15.0, 30.0, 45.0, 60.0, 75.0, 90.0])
    altitude_labels = tuple(
        f"{int(start)}-{int(end)}{DEG}" for start, end in zip(altitude_edges[:-1], altitude_edges[1:])
    )
    altitude_indices = np.digitize(radiation_frame["altitude_deg"].to_numpy(), altitude_edges, right=False) - 1
    altitude_indices = np.clip(altitude_indices, 0, len(altitude_labels) - 1)

    temperature = radiation_frame["dry_bulb_c"].to_numpy(dtype=float)
    direct_normal = radiation_frame["direct_normal_radiation_wh_m2"].to_numpy(dtype=float)
    diffuse_horizontal = radiation_frame["diffuse_horizontal_radiation_wh_m2"].to_numpy(dtype=float)
    altitude_rad = np.deg2rad(radiation_frame["altitude_deg"].to_numpy(dtype=float))
    azimuth_deg = radiation_frame["azimuth_deg"].to_numpy(dtype=float)

    benefit_weight = np.clip((18.0 - temperature) / 8.0, 0.0, 1.0)
    harm_weight = np.clip((temperature - 24.0) / 8.0, 0.0, 1.0)
    comfort_weight = benefit_weight - harm_weight

    rows: list[dict[str, float | str]] = []
    orientation_centers = np.arange(0.0, 360.0, 22.5)

    for direction_label, theta_deg in zip(WIND_DIRECTION_LABELS, orientation_centers):
        azimuth_delta_rad = np.deg2rad(((azimuth_deg - theta_deg + 180.0) % 360.0) - 180.0)
        incident_direct = direct_normal * np.maximum(0.0, np.cos(altitude_rad) * np.cos(azimuth_delta_rad))
        incident_vertical = incident_direct + diffuse_horizontal * 0.5
        signed_incident = incident_vertical * comfort_weight

        orientation_frame = pd.DataFrame(
            {
                "altitude_band": [altitude_labels[int(index)] for index in altitude_indices],
                "comfort_radiation_kwh_m2": signed_incident / 1000.0,
                "incident_radiation_kwh_m2": incident_vertical / 1000.0,
            }
        )
        summary = (
            orientation_frame.groupby("altitude_band", observed=False)
            .agg(
                comfort_radiation_kwh_m2=("comfort_radiation_kwh_m2", "sum"),
                incident_radiation_kwh_m2=("incident_radiation_kwh_m2", "sum"),
            )
            .reindex(altitude_labels, fill_value=0.0)
            .reset_index()
        )

        for band_index, band_row in summary.iterrows():
            rows.append(
                {
                    "direction": direction_label,
                    "theta_deg": float(theta_deg),
                    "altitude_band": str(band_row["altitude_band"]),
                    "radius_start": float(altitude_edges[band_index]),
                    "radius_size": float(altitude_edges[band_index + 1] - altitude_edges[band_index]),
                    "comfort_radiation_kwh_m2": float(band_row["comfort_radiation_kwh_m2"]),
                    "incident_radiation_kwh_m2": float(band_row["incident_radiation_kwh_m2"]),
                }
            )

    return pd.DataFrame(rows)


def _build_wind_rose_table(data_frame: pd.DataFrame) -> pd.DataFrame:
    """Freqüència del vent de retorn i velocitat mitjana agrupades en sectors de brúixola."""

    wind_frame = data_frame[["wind_direction_deg", "wind_speed_m_s"]].dropna()
    if wind_frame.empty:
        return pd.DataFrame(columns=["direction", "frequency_pct", "mean_speed"])

    adjusted = (wind_frame["wind_direction_deg"] + 11.25) % 360
    sector_indices = ((adjusted / 22.5).astype(int) % len(WIND_DIRECTION_LABELS)).to_numpy()
    sectors = [WIND_DIRECTION_LABELS[int(index)] for index in sector_indices]
    rose = (
        pd.DataFrame(
            {
                "direction": sectors,
                "wind_speed_m_s": wind_frame["wind_speed_m_s"].to_numpy(),
            }
        )
        .groupby("direction", as_index=False)
        .agg(
            count=("wind_speed_m_s", "size"),
            mean_speed=("wind_speed_m_s", "mean"),
        )
        .set_index("direction")
        .reindex(WIND_DIRECTION_LABELS, fill_value=0.0)
        .reset_index()
    )
    total = rose["count"].sum()
    rose["frequency_pct"] = np.where(total > 0, rose["count"] / total * 100.0, 0.0)
    return rose


def build_monthly_wind_rose_tables(data_frame: pd.DataFrame) -> dict[str, pd.DataFrame]:
    """Construeix una taula resum de la rosa dels vents per a cada mes natural."""

    month_tables: dict[str, pd.DataFrame] = {}
    for month_number, month_label in MONTH_LABELS.items():
        month_frame = data_frame[data_frame["month"] == month_number]
        rose = _build_wind_rose_table(month_frame)
        if not rose.empty:
            month_tables[month_label] = rose
    return month_tables


def build_download_frame(data_frame: pd.DataFrame) -> pd.DataFrame:
    """Retorna les columnes EPW netes que s'exporten com a CSV des de UI."""

    export_columns = [
        "timestamp",
        "source_year",
        "month",
        "day",
        "hour",
        "dry_bulb_c",
        "dew_point_c",
        "relative_humidity_pct",
        "wind_speed_m_s",
        "wind_direction_deg",
        "atmospheric_pressure_pa",
        "global_horizontal_radiation_wh_m2",
        "direct_normal_radiation_wh_m2",
        "diffuse_horizontal_radiation_wh_m2",
        "total_sky_cover_tenths",
        "opaque_sky_cover_tenths",
    ]
    export_frame = data_frame[export_columns].copy()
    export_frame["timestamp"] = export_frame["timestamp"].dt.strftime("%Y-%m-%d %H:%M")
    return export_frame
\end{lstlisting}

\Needspace{8\baselineskip}

\subsection{\texorpdfstring{\texttt{backend/weather\_profiles.py}}{backend/weather\_profiles.py}}

\begin{lstlisting}[style=studioPython,caption={backend/weather\_profiles.py},label={lst:backend-weather-profiles-py}]
"""Anàlisi de perfils meteorològics i vista prèvia per crear entorns."""

from __future__ import annotations

import csv
import re
from dataclasses import dataclass
from pathlib import Path
from typing import Dict, List, Optional, Sequence, Tuple

import numpy as np
import plotly.graph_objects as go


EPW_HEADER_ROW_COUNT = 8
EPW_DRY_BULB_COLUMN_INDEX = 6
WEATHER_PREVIEW_RANDOM_SEED = 42
MONTH_START_DAYS = (0, 31, 59, 90, 120, 151, 181, 212, 243, 273, 304, 334)


def _sanitize_weather_identifier(raw_value: str, fallback: str) -> str:
    """Retorna un identificador estable en minúscules per a les claus del perfil meteorològic."""

    slug = re.sub(r"[^a-zA-Z0-9]+", "_", raw_value.strip().lower()).strip("_")
    return slug or fallback


def _unique_weather_identifier(base: str, used: set[str]) -> str:
    """Retorna un identificador meteorològic únic dins del conjunt proporcionat."""

    candidate = base
    suffix = 2
    while candidate in used:
        candidate = f"{base}_{suffix}"
        suffix += 1
    used.add(candidate)
    return candidate


@dataclass(frozen=True)
class WeatherProfileSuggestion:
    """Descriu un suggeriment de fitxer meteorològic per a la creació de l'entorn UI."""

    file_name: str
    path: str
    suggested_key: str
    climate_label: str
    mean_dry_bulb: Optional[float]


def summarize_weather_file(weather_path: Path) -> Tuple[Optional[float], str]:
    """Analitza un fitxer EnergyPlus Weather (EPW) per suggerir una classificació climàtica.

    Paràmetres:
        weather_path (Path): camí cap al fitxer EPW.

    Retorna:
        Tuple[Optional[float], str]: una tupla que conté la temperatura mitjana de bulb sec (si s'analitza correctament),
        i una cadena de category ('calent', 'cool', 'mixt' o 'personalitzat').
    """
    try:
        with open(weather_path, "r", encoding="utf-8", errors="ignore") as file_handle:
            lines = file_handle.readlines()[8:]
    except OSError:
        return None, "custom"

    temperatures: List[float] = []
    for line in lines:
        parts = line.strip().split(",")
        if len(parts) <= 6:
            continue
        try:
            temperatures.append(float(parts[6]))
        except ValueError:
            continue

    if not temperatures:
        return None, "custom"

    mean_dry_bulb = sum(temperatures) / len(temperatures)
    if mean_dry_bulb >= 18.0:
        return mean_dry_bulb, "hot"
    if mean_dry_bulb <= 12.0:
        return mean_dry_bulb, "cool"
    return mean_dry_bulb, "mixed"


def _parse_epw_int(value: str) -> Optional[int]:
    """Analitza un nombre enter d'una cel·la EPW tot tolerant les cadenes decimals."""

    try:
        return int(float(value))
    except (TypeError, ValueError):
        return None


def _parse_epw_float(value: str) -> Optional[float]:
    """Analitza un nombre flotant d'una cel·la EPW i retorna None quan la cel·la no sigui vàlida."""

    try:
        return float(value)
    except (TypeError, ValueError):
        return None


def _calculate_epw_day_of_year(row: Sequence[str], fallback_index: int) -> int:
    """Retorna el dia de l'any basat en 1 representat per una fila de dades EPW."""

    fallback_day = fallback_index // 24 + 1
    if len(row) < 3:
        return fallback_day

    month = _parse_epw_int(row[1])
    day = _parse_epw_int(row[2])
    if month is None or day is None or month < 1 or month > 12 or day < 1:
        return fallback_day

    return MONTH_START_DAYS[month - 1] + day


def load_epw_dry_bulb_series(weather_path: Path) -> List[Dict[str, float]]:
    """Carrega els registres de temperatura de bulb sec cada hora des d'un fitxer meteorològic EPW.

    Paràmetres:
        weather_path (Path): camí cap al fitxer meteorològic EPW.

    Retorna:
        List[Dict[str, float]]: registres horàries amb dia de l'any, índex d'hores,
        i temperatura de bulb sec en graus centígrads. S'ometen files no vàlides.
    """

    records: List[Dict[str, float]] = []
    try:
        with open(weather_path, "r", encoding="utf-8", errors="ignore", newline="") as file_handle:
            reader = csv.reader(file_handle)
            for _ in range(EPW_HEADER_ROW_COUNT):
                next(reader, None)

            for row in reader:
                if len(row) <= EPW_DRY_BULB_COLUMN_INDEX:
                    continue
                temperature = _parse_epw_float(row[EPW_DRY_BULB_COLUMN_INDEX])
                if temperature is None:
                    continue

                record_index = len(records)
                records.append(
                    {
                        "annual_hour": float(record_index + 1),
                        "day_of_year": float(_calculate_epw_day_of_year(row, record_index)),
                        "dry_bulb_temperature_c": temperature,
                    }
                )
    except OSError:
        return []

    return records


def _build_ornstein_uhlenbeck_noise(
    row_count: int,
    sigma: float,
    mu: float,
    tau: float,
    *,
    seed: int = WEATHER_PREVIEW_RANDOM_SEED,
) -> np.ndarray:
    """Construeix un soroll determinista Ornstein-Uhlenbeck per al gràfic de vista prèvia del clima."""

    safe_tau = max(float(tau), 1e-9)
    sigma_bis = float(sigma) * np.sqrt(2.0 / safe_tau)
    rng = np.random.default_rng(seed)
    noise = np.zeros(row_count)

    for index in range(row_count - 1):
        noise[index + 1] = (
            noise[index]
            + (-(noise[index] - float(mu)) / safe_tau)
            + sigma_bis * rng.normal()
        )

    return noise


def build_weather_temperature_preview(
    weather_path: Path,
    sigma: float,
    mu: float,
    tau: float,
) -> List[Dict[str, float]]:
    """Crea dades de vista prèvia de la temperatura diària per a la variació estocàstica del clima.

    La vista prèvia reflecteix la pertorbació meteorològica d'Ornstein-Uhlenbeck amb Sinergym
    una llavor aleatòria fixa, de manera que UI s'actualitza de manera previsible quan l'usuari canvia el
    controls sigma, mu o tau.

    Paràmetres:
        weather_path (Path): fitxer EPW utilitzat com a font climàtica base.
        sigma (float): paràmetre de desviació estàndard per al procés d'OU.
        mu (float): paràmetre mitjà per al procés d'OU.
        tau (float): paràmetre constant de temps per al procés d'OU, en hores.

    Retorna:
        List[Dict[str, float]]: temperatures mitjanes diàries per al EPW original,
        la vista prèvia modificada determinista i una banda sigma visual.
    """

    hourly_records = load_epw_dry_bulb_series(weather_path)
    if not hourly_records:
        return []

    noise = _build_ornstein_uhlenbeck_noise(
        len(hourly_records),
        sigma=float(sigma),
        mu=float(mu),
        tau=float(tau),
    )
    daily_accumulators: Dict[int, Dict[str, float]] = {}

    for index, record in enumerate(hourly_records):
        day = int(record["day_of_year"])
        base_temperature = record["dry_bulb_temperature_c"]
        expected_temperature = base_temperature + float(mu)
        modified_temperature = base_temperature + float(noise[index])
        accumulator = daily_accumulators.setdefault(
            day,
            {
                "count": 0.0,
                "base_temperature_c": 0.0,
                "expected_temperature_c": 0.0,
                "modified_temperature_c": 0.0,
                "lower_temperature_c": 0.0,
                "upper_temperature_c": 0.0,
            },
        )
        accumulator["count"] += 1.0
        accumulator["base_temperature_c"] += base_temperature
        accumulator["expected_temperature_c"] += expected_temperature
        accumulator["modified_temperature_c"] += modified_temperature
        accumulator["lower_temperature_c"] += expected_temperature - float(sigma)
        accumulator["upper_temperature_c"] += expected_temperature + float(sigma)

    preview_records: List[Dict[str, float]] = []
    for day, accumulator in sorted(daily_accumulators.items()):
        count = max(accumulator.pop("count"), 1.0)
        preview_records.append(
            {
                "day_of_year": float(day),
                "base_temperature_c": accumulator["base_temperature_c"] / count,
                "expected_temperature_c": accumulator["expected_temperature_c"] / count,
                "modified_temperature_c": accumulator["modified_temperature_c"] / count,
                "lower_temperature_c": accumulator["lower_temperature_c"] / count,
                "upper_temperature_c": accumulator["upper_temperature_c"] / count,
            }
        )

    return preview_records


def build_weather_temperature_preview_figure(preview_records: Sequence[Dict[str, float]]) -> go.Figure:
    """Crea la figura de previsualització anual de la temperatura per a la pàgina Afegeix un entorn.

    Paràmetres:
        preview_records (Sequence[Dict[str, float]]): dades de vista prèvia de la temperatura diària.

    Retorna:
        go.Figure: Plotly figura que compara la temperatura de bulb sec original i modificada.
    """

    day_values = [record["day_of_year"] for record in preview_records]
    max_day = max(day_values) if day_values else 365.0
    fig = go.Figure()
    fig.add_trace(
        go.Scatter(
            x=day_values,
            y=[record["upper_temperature_c"] for record in preview_records],
            mode="lines",
            line=dict(width=0),
            showlegend=False,
            hoverinfo="skip",
        )
    )
    fig.add_trace(
        go.Scatter(
            x=day_values,
            y=[record["lower_temperature_c"] for record in preview_records],
            mode="lines",
            fill="tonexty",
            fillcolor="rgba(95, 84, 249, 0.13)",
            line=dict(width=0),
            name="Rang sigma",
            hoverinfo="skip",
        )
    )
    fig.add_trace(
        go.Scatter(
            x=day_values,
            y=[record["base_temperature_c"] for record in preview_records],
            mode="lines",
            name="EPW original",
            line=dict(color="#3b82f6", width=2),
            hovertemplate="Dia %{x:.0f}<br>Original %{y:.1f} °C<extra></extra>",
        )
    )
    fig.add_trace(
        go.Scatter(
            x=day_values,
            y=[record["modified_temperature_c"] for record in preview_records],
            mode="lines",
            name="Variant simulada",
            line=dict(color="#ef4444", width=2.2),
            hovertemplate="Dia %{x:.0f}<br>Simulada %{y:.1f} °C<extra></extra>",
        )
    )
    fig.update_layout(
        height=420,
        margin=dict(l=42, r=24, t=34, b=46),
        template="plotly_white",
        paper_bgcolor="#ffffff",
        plot_bgcolor="#f7f9fe",
        font=dict(color="#17233c", family="Bahnschrift, Aptos, Segoe UI"),
        legend=dict(orientation="h", yanchor="bottom", y=1.02, xanchor="right", x=1),
        xaxis=dict(
            title="Dia de l'any",
            range=[1, max(365.0, float(max_day))],
            gridcolor="#dce3f2",
            zeroline=False,
        ),
        yaxis=dict(
            title="Temperatura seca (°C)",
            gridcolor="#dce3f2",
            zeroline=False,
        ),
    )
    return fig


def build_weather_profile_suggestions(weather_paths: Sequence[Path]) -> List[WeatherProfileSuggestion]:
    """Crea configuracions de perfils per a fitxers meteorològics amb noms i categorització fàcils d'utilitzar.

    Paràmetres:
        weather_paths (Sequence[Path]): una col·lecció de camins que apunten a fitxers EPW.

    Retorna:
        List[WeatherProfileSuggestion]: Configuracions recomanades i resums climàtics per als fitxers.
    """
    used_keys: set[str] = set()
    suggestions: List[WeatherProfileSuggestion] = []

    for path in weather_paths:
        mean_dry_bulb, suggested_family = summarize_weather_file(path)
        if mean_dry_bulb is None:
            climate_label = "Sense resum climàtic automàtic"
            base_key = _sanitize_weather_identifier(path.stem, "weather")
        else:
            if suggested_family == "hot":
                climate_label = f"Clima càlid ({mean_dry_bulb:.1f} °C de bulb sec mitjà)"
            elif suggested_family == "cool":
                climate_label = f"Clima fred ({mean_dry_bulb:.1f} °C de bulb sec mitjà)"
            else:
                climate_label = f"Clima mixt ({mean_dry_bulb:.1f} °C de bulb sec mitjà)"
            base_key = suggested_family

        unique_key = _unique_weather_identifier(base_key, used_keys)
        suggestions.append(
            WeatherProfileSuggestion(
                file_name=path.name,
                path=str(path),
                suggested_key=unique_key,
                climate_label=climate_label,
                mean_dry_bulb=mean_dry_bulb,
            )
        )

    return suggestions
\end{lstlisting}

\section{Codi: Grafics i metriques}

\Needspace{8\baselineskip}

\subsection{\texorpdfstring{\texttt{backend/grafics/\_\_init\_\_.py}}{backend/grafics/\_\_init\_\_.py}}

\begin{lstlisting}[style=studioPython,caption={backend/grafics/\_\_init\_\_.py},label={lst:backend-grafics---init---py}]
"""Gràfics Plotly i mètriques compartides per a resultats i informes.

Els mòduls d'aquest paquet converteixen les dades d'execució normalitzades en xifres reutilitzables:
compliment de comoditat, consum HVAC, preus de l'energia, comportament de la bateria, acció
senyals i mapes de calor temporals. Per tant, haurien de romandre independents de Streamlit
el panell i el generador d'informes poden compartir la mateixa lògica visual.
"""
\end{lstlisting}

\Needspace{8\baselineskip}

\subsection{\texorpdfstring{\texttt{backend/grafics/battery.py}}{backend/grafics/battery.py}}

\begin{lstlisting}[style=studioPython,caption={backend/grafics/battery.py},label={lst:backend-grafics-battery-py}]
"""Dades del quadre de comandament de la bateria, la xarxa i les tarifes."""

from __future__ import annotations

import pandas as pd
import plotly.graph_objects as go

from backend.grafics.column_utils import (
    _BATTERY_CHARGE_POWER_COLUMNS,
    _BATTERY_DISCHARGE_ENERGY_COLUMNS,
    _BATTERY_DISCHARGE_POWER_COLUMNS,
    _GRID_ENERGY_COLUMNS,
    _GRID_KWH_COLUMNS,
    _GRID_POWER_COLUMNS,
    _find_battery_state_column,
    _find_energy_price_column,
    _find_first_present_column,
)
from backend.grafics.energy_units import (
    _battery_state_display_series,
    _energy_price_series_eur_per_kwh,
    _energy_to_kwh,
)
from backend.grafics.plot_helpers import (
    add_energy_bar_trace,
    add_mode_scatter_trace,
    build_mode_scatter_trace,
    energy_axis_title,
)
from backend.grafics.style import STUDIO_CHART_COLORS, _apply_figure_theme
from backend.grafics.time_utils import (
    _apply_mode_xaxis,
    _auto_scale_series,
    _ensure_datetime_index,
    _infer_timestep_hours,
)


def _add_battery_power_trace(
    fig: go.Figure,
    base: pd.DataFrame,
    series: str | None,
    label: str,
    sign: int,
    mode: str,
    season: str,
    units: str,
) -> None:
    """Afegeix un rastre de càrrega o descàrrega a la figura de la bateria."""
    color = (
        STUDIO_CHART_COLORS["battery_charge"]
        if sign > 0
        else STUDIO_CHART_COLORS["battery_discharge"]
    )
    add_mode_scatter_trace(
        fig,
        base,
        series,
        mode,
        season,
        name=label,
        color=color,
        units=units,
        sign=sign,
        hovertemplate="%{hovertext}<br>%{y:.2f} " + units + "<extra></extra>",
    )


def _add_battery_grid_bar_trace(
    fig: go.Figure,
    base: pd.DataFrame,
    column: str,
    label: str,
    color: str,
    mode: str,
    season: str,
) -> None:
    """Afegeix una traça de la barra de bateria contra la xarxa per al mode d'agregació seleccionat."""
    add_energy_bar_trace(
        fig,
        base,
        column,
        label=label,
        color=color,
        mode=mode,
        season=season,
    )


def make_battery_power_plot(obs: pd.DataFrame, mode: str, season: str) -> go.Figure:
    """Crea una figura Plotly per a la bateria."""
    base = _ensure_datetime_index(obs).copy()
    fig = go.Figure()

    charge_power_col = _find_first_present_column(base.columns, _BATTERY_CHARGE_POWER_COLUMNS)
    discharge_power_col = _find_first_present_column(base.columns, _BATTERY_DISCHARGE_POWER_COLUMNS)
    state_col = _find_battery_state_column(base.columns)
    has_charge_p = charge_power_col is not None
    has_discharge_p = discharge_power_col is not None
    has_state = state_col is not None

    if not (has_charge_p or has_discharge_p or has_state):
        return go.Figure()

    scale_suffix = ""
    # Reserva: només tenim battery_charge_state → fem servir la diferència per pas
    if not (has_charge_p or has_discharge_p):
        s = pd.to_numeric(base[state_col], errors="coerce")
        d = s.diff().fillna(0.0)                # +: carrega, −: descarrega
        d, scale_suffix = _auto_scale_series(d) # millora llegibilitat
        base["__charge__"]    = d.clip(lower=0.0)        # ≥ 0
        base["__discharge__"] = -d.clip(upper=0.0)       # ≥ 0
        charge_col, discharge_col = "__charge__", "__discharge__"
        units = f"Δstate/step {scale_suffix}" if scale_suffix else "Δstate/step"
    else:
        charge_col = charge_power_col
        discharge_col = discharge_power_col
        units = "W"

    # Convenció del gràfic: càrrega positiva, descàrrega negativa
    _add_battery_power_trace(fig, base, charge_col, "Charge (+)", +1, mode, season, units)
    _add_battery_power_trace(fig, base, discharge_col, "Discharge (−)", -1, mode, season, units)

    # Eixos i línia y=0
    _apply_mode_xaxis(fig, mode)
    fig.update_yaxes(title=f"Battery Power ({units})")
    fig.add_hline(y=0, line_dash="dash", line_color="black")
    fig.update_layout(title=f"Battery Charge/Discharge – {mode.capitalize()} ({season})")
    return fig



def make_battery_soc_plot(obs: pd.DataFrame, mode: str, season: str) -> go.Figure:
    """Crea una figura Plotly per a la bateria."""
    base = _ensure_datetime_index(obs).copy()
    fig = go.Figure()

    state_col = _find_battery_state_column(base.columns)
    if state_col is None:
        return go.Figure()

    s = pd.to_numeric(base[state_col], errors="coerce")
    if s.dropna().empty:
        return go.Figure()

    s, units, state_label = _battery_state_display_series(base, state_col)
    base["__soc__"] = s

    trace = build_mode_scatter_trace(
        base,
        "__soc__",
        mode,
        season,
        name=state_label,
        color=STUDIO_CHART_COLORS["battery"],
        units=units,
        line_width=2.0 if mode == "raw" else 2.4,
        hovertemplate="%{hovertext}<br>%{y:.4f} " + units + "<extra></extra>",
        meta="keep_color",
    )
    if trace is None:
        return go.Figure()
    fig.add_trace(trace)
    _apply_mode_xaxis(fig, mode)

    fig.update_yaxes(title=f"{state_label} ({units})")
    fig.update_layout(title=f"{state_label} – {mode.capitalize()} ({'All' if mode=='month' else season})")
    return fig

def _grid_source_unit_kind(column_name: str) -> str:
    """Tipus d'unitat font de la xarxa."""
    label = column_name.lower()
    if "kwh" in label:
        return "kwh"
    if "energy" in label or "joule" in label:
        return "joule"
    if "rate" in label or "power" in label:
        return "watt"
    return "watt"


def make_energy_price_plot(obs: pd.DataFrame, mode: str, season: str) -> go.Figure:
    """Crea una figura de Plotly per al preu de l'energia."""
    from backend.grafics.observation_columns import normalize_observation_columns

    base = _ensure_datetime_index(normalize_observation_columns(obs)).copy()
    price_col = _find_energy_price_column(base.columns)
    if price_col is None:
        return go.Figure()

    base["__energy_price_eur_per_kwh__"] = _energy_price_series_eur_per_kwh(base, price_col)
    price_col = "__energy_price_eur_per_kwh__"
    if base[price_col].dropna().empty:
        return go.Figure()

    trace = build_mode_scatter_trace(
        base,
        price_col,
        mode,
        season,
        name="Energy Price",
        color=STUDIO_CHART_COLORS["price"],
        units="EUR/kWh",
        line_width=2.4 if mode == "raw" else 2.8,
        hovertemplate="%{hovertext}<br>%{y:.4f} EUR/kWh<extra></extra>",
    )
    if trace is None:
        return go.Figure()

    fig = go.Figure(trace)
    _apply_mode_xaxis(fig, mode)
    fig.update_yaxes(title="Price (EUR/kWh)")
    fig.update_layout(title=f"Energy Price - {mode.capitalize()} ({season})")
    return fig


def make_battery_vs_grid_plot(obs: pd.DataFrame, mode: str, season: str) -> go.Figure:
    """Crea una figura Plotly per a la bateria i la xarxa."""
    base = _ensure_datetime_index(obs).copy()

    discharge_power_col = _find_first_present_column(base.columns, _BATTERY_DISCHARGE_POWER_COLUMNS)
    discharge_energy_col = _find_first_present_column(base.columns, _BATTERY_DISCHARGE_ENERGY_COLUMNS)
    grid_col = _find_first_present_column(
        base.columns,
        _GRID_POWER_COLUMNS + _GRID_ENERGY_COLUMNS + _GRID_KWH_COLUMNS,
    )

    if discharge_power_col is None and discharge_energy_col is None and grid_col is None:
        return go.Figure()

    ts_hours = _infer_timestep_hours(base)

    if discharge_power_col is not None:
        bat_w = pd.to_numeric(base[discharge_power_col], errors="coerce").clip(lower=0.0)
        base["__bat_kwh__"] = _energy_to_kwh(bat_w, "watt", ts_hours)
    elif discharge_energy_col is not None:
        base["__bat_kwh__"] = _energy_to_kwh(
            pd.to_numeric(base[discharge_energy_col], errors="coerce").clip(lower=0.0),
            "joule",
        )
    if grid_col is not None:
        grid_kind = _grid_source_unit_kind(grid_col)
        grid_values = pd.to_numeric(base[grid_col], errors="coerce").clip(lower=0.0)
        base["__grid_kwh__"] = _energy_to_kwh(grid_values, grid_kind, ts_hours)

    bat_color = STUDIO_CHART_COLORS["battery_charge"]
    grid_color = "#eab308"
    fig = go.Figure()

    _add_battery_grid_bar_trace(fig, base, "__bat_kwh__", "Energia Bateria", bat_color, mode, season)
    _add_battery_grid_bar_trace(fig, base, "__grid_kwh__", "Compra Xarxa", grid_color, mode, season)

    _apply_mode_xaxis(fig, mode)
    y_title = energy_axis_title(mode)

    fig.update_layout(barmode="group")

    fig.update_yaxes(title=y_title)
    fig.update_layout(title=f"Energia Bateria vs Xarxa - {mode.capitalize()} ({season})")
    return fig


# Gràfic de càrrega/descarrega de bateria amb preu d'energia.

def make_battery_charge_with_price_plot(obs: pd.DataFrame, mode: str, season: str) -> go.Figure:
    """Crea una figura Plotly per a la càrrega de la bateria amb el preu."""
    from plotly.subplots import make_subplots as _make_subplots
    from backend.grafics.observation_columns import normalize_observation_columns

    base = _ensure_datetime_index(normalize_observation_columns(obs)).copy()

    # Detectem les columnes de bateria que existeixen en aquest CSV.
    charge_col = _find_first_present_column(base.columns, _BATTERY_CHARGE_POWER_COLUMNS)
    discharge_col = _find_first_present_column(base.columns, _BATTERY_DISCHARGE_POWER_COLUMNS)
    state_col = _find_battery_state_column(base.columns)
    has_bat = charge_col is not None or discharge_col is not None or state_col is not None

    # El preu pot venir amb noms diferents segons el wrapper utilitzat.
    price_col_raw = _find_energy_price_column(base.columns)
    has_price = price_col_raw is not None

    if not has_bat and not has_price:
        return go.Figure()

    # Si només tenim l'estat de càrrega, estimem càrrega/descarrega amb la diferència.
    units = "W"
    if has_bat:
        if charge_col is None and discharge_col is None:
            s = pd.to_numeric(base[state_col], errors="coerce")
            d = s.diff().fillna(0.0)
            d, _sfx = _auto_scale_series(d)
            base["__charge__"] = d.clip(lower=0.0)
            base["__discharge__"] = -d.clip(upper=0.0)
            charge_col, discharge_col = "__charge__", "__discharge__"
            units = f"Δstate/step {_sfx}".strip() if _sfx else "Δstate/step"

    # Normalitzem el preu a EUR/kWh perquè l'eix sigui comparable entre runs.
    price_col = "__price__"
    if has_price:
        base[price_col] = _energy_price_series_eur_per_kwh(base, price_col_raw)
        if base[price_col].dropna().empty:
            has_price = False

    # Posem la bateria a l'eix principal i el preu a l'eix secundari.
    fig = _make_subplots(specs=[[{"secondary_y": True}]])

    bat_charge_color = STUDIO_CHART_COLORS["battery_charge"]
    bat_discharge_color = STUDIO_CHART_COLORS["battery_discharge"]
    price_color = STUDIO_CHART_COLORS["price"]

    # Carrega (eix primari)
    if has_bat:
        charge_trace = build_mode_scatter_trace(
            base,
            charge_col,
            mode,
            season,
            name=f"Carrega (+) [{units}]",
            color=bat_charge_color,
            units=units,
            sign=+1,
            line_width=2.6,
            hovertemplate="%{hovertext}<br>Carrega: %{y:.2f} " + units + "<extra></extra>",
            plain_hovertemplate="Carrega: %{y:.2f} " + units + "<extra></extra>",
        )
        if charge_trace is not None:
            fig.add_trace(
                charge_trace,
                secondary_y=False,
            )

        # Descarrega (eix primari, valors negatius)
        discharge_trace = build_mode_scatter_trace(
            base,
            discharge_col,
            mode,
            season,
            name=f"Descarrega (-) [{units}]",
            color=bat_discharge_color,
            units=units,
            sign=-1,
            line_width=2.6,
            fill="tozeroy",
            fillcolor="rgba(234,88,12,0.10)",  # battery_discharge amb alpha
            hovertemplate="%{hovertext}<br>Descarrega: %{y:.2f} " + units + "<extra></extra>",
            plain_hovertemplate="Descarrega: %{y:.2f} " + units + "<extra></extra>",
        )
        if discharge_trace is not None:
            fig.add_trace(
                discharge_trace,
                secondary_y=False,
            )

        fig.add_hline(y=0, line_dash="dash", line_color=STUDIO_CHART_COLORS["neutral_soft"],
                      line_width=1.2)

    # Preu (eix secundari)
    price_trace = (
        build_mode_scatter_trace(
            base,
            price_col,
            mode,
            season,
            name="Preu Energia (EUR/kWh)",
            color=price_color,
            units="EUR/kWh",
            line_width=2.2,
            marker_size=5,
            hovertemplate="%{hovertext}<br>Preu: %{y:.4f} EUR/kWh<extra></extra>",
            plain_hovertemplate="Preu: %{y:.4f} EUR/kWh<extra></extra>",
        )
        if has_price
        else None
    )
    if price_trace is not None:
        fig.add_trace(
            price_trace,
            secondary_y=True,
        )

    # Ajustem els eixos segons l'agregació escollida.
    x_titles = {"hour": "Hour of Day", "day": "Day (1..365)", "month": "Month", "raw": "Time (step)"}
    x_dticks = {"hour": 1, "month": 1}
    fig.update_xaxes(title_text=x_titles.get(mode, ""), tickmode="linear" if mode in x_dticks else None,
                     dtick=x_dticks.get(mode))
    fig.update_yaxes(title_text=f"Potencia Bateria ({units})", secondary_y=False,
                     zeroline=True, zerolinecolor=STUDIO_CHART_COLORS["neutral_soft"])
    fig.update_yaxes(title_text="Preu Energia (EUR/kWh)", secondary_y=True,
                     showgrid=False)

    fig.update_layout(
        title=f"Carrega/Descarrega Bateria vs Preu Energia – {mode.capitalize()} ({season})",
        legend=dict(orientation="h", yanchor="bottom", y=1.02, xanchor="right", x=1),
    )
    return _apply_figure_theme(fig)
\end{lstlisting}

\Needspace{8\baselineskip}

\subsection{\texorpdfstring{\texttt{backend/grafics/column\_utils.py}}{backend/grafics/column\_utils.py}}

\begin{lstlisting}[style=studioPython,caption={backend/grafics/column\_utils.py},label={lst:backend-grafics-column-utils-py}]
"""Detecció de columnes disponibles per construir les figures del panell."""

from __future__ import annotations

import pandas as pd

def _ensure_air_temperature(df: pd.DataFrame) -> pd.DataFrame:
    """Assegura 'air_temperature'; si no hi és, fa la mitjana de columnes *_air_temperature."""
    dfc = df.copy()
    if "air_temperature" not in dfc.columns:
        zone_cols = [c for c in dfc.columns if c.endswith("_air_temperature")]
        if zone_cols:
            dfc["air_temperature"] = dfc[zone_cols].mean(axis=1)
    return dfc


def _find_first_present_column(columns, candidates) -> str | None:
    """Troba la primera columna disponible."""
    available = set(columns)
    for candidate in candidates:
        if candidate in available:
            return candidate

    lower_lookup = {str(column).lower(): column for column in columns}
    for candidate in candidates:
        actual = lower_lookup.get(str(candidate).lower())
        if actual is not None:
            return actual
    return None


_BATTERY_STATE_COLUMNS = (
    "battery_charge_state",
    "storage_charge_state",
    "electric_storage_charge_state",
    "storage_battery_charge_state",
    "storage_battery_charge_fraction",
    "battery_charge_fraction",
    "battery_state_of_charge",
    "storage_battery_soc",
    "battery_soc",
    "state_of_charge",
)

_BATTERY_CHARGE_POWER_COLUMNS = (
    "storage_charge_power",
    "battery_charge_power",
    "battery_charge_rate",
)

_BATTERY_DISCHARGE_POWER_COLUMNS = (
    "storage_discharge_power",
    "battery_discharge_power",
    "battery_discharge_rate",
)

_BATTERY_DISCHARGE_ENERGY_COLUMNS = (
    "storage_discharge_energy",
    "battery_discharge_energy",
)

_GRID_POWER_COLUMNS = (
    "facility_net_purchased_electricity_rate",
    "facility_total_purchased_electricity_rate",
    "net_electricity_purchased",
    "grid_electricity_rate",
)

_GRID_ENERGY_COLUMNS = (
    "facility_net_purchased_electricity_energy",
    "facility_total_purchased_electricity_energy",
    "net_purchased_electricity_energy",
    "total_purchased_electricity_energy",
)

_GRID_KWH_COLUMNS = (
    "facility_net_purchased_electricity_kwh",
    "facility_total_purchased_electricity_kwh",
    "net_purchased_electricity_kwh",
    "total_purchased_electricity_kwh",
)

_ENERGY_PRICE_COLUMNS = (
    "energy_cost",
    "electricity_cost",
    "elèctricity_cost",
    "price_eur_per_kwh",
    "energy_price_kwh",
    "energy_cost_eur_per_kwh",
    "electricity_cost_eur_per_kwh",
)


def _find_battery_state_column(columns) -> str | None:
    """Cerca la columna d'estat de la bateria."""
    return _find_first_present_column(columns, _BATTERY_STATE_COLUMNS)


def _find_energy_price_column(columns) -> str | None:
    """Busca la columna del preu de l'energia."""
    return _find_first_present_column(columns, _ENERGY_PRICE_COLUMNS)
\end{lstlisting}

\Needspace{8\baselineskip}

\subsection{\texorpdfstring{\texttt{backend/grafics/comfort.py}}{backend/grafics/comfort.py}}

\begin{lstlisting}[style=studioPython,caption={backend/grafics/comfort.py},label={lst:backend-grafics-comfort-py}]
"""Compliment de la comoditat i xifres d'incompliment de la temperatura."""

from __future__ import annotations

import numpy as np
import pandas as pd
import plotly.express as px
import plotly.graph_objects as go

from backend.grafics.column_utils import _ensure_air_temperature
from backend.grafics.comfort_scope import comfort_scope_label, filter_by_comfort_scope
from backend.grafics.style import STUDIO_CHART_COLORS, _apply_figure_theme
from backend.grafics.time_utils import (
    _apply_mode_xaxis,
    _ensure_datetime_index,
    _month_series,
    _raw_axis_data,
    _seasonal_marker_colors,
    _series_for_mode,
    filter_by_season,
)

DEFAULT_WINTER_COMFORT_RANGE = (20.0, 23.5)
DEFAULT_SUMMER_COMFORT_RANGE = (23.0, 26.0)
DEFAULT_SUMMER_START = (6, 1)
DEFAULT_SUMMER_FINAL = (9, 30)


def _violation_marker(color=None) -> dict:
    """Retorna el marcador comú de barres de violació tèrmica."""
    return dict(
        color=color or STUDIO_CHART_COLORS["violation"],
        line=dict(color=STUDIO_CHART_COLORS["violation_edge"], width=0.8),
    )


def _add_violation_bar(fig: go.Figure, *, x, y, name: str, color=None) -> None:
    """Afegeix una barra de violació amb el mateix estil visual."""
    fig.add_bar(
        x=x,
        y=y,
        name=name,
        marker=_violation_marker(color),
        meta="keep_color",
    )


def _extract_reward_kwargs(config) -> dict:
    """Extreu kwargs de recompensa."""
    if not isinstance(config, dict):
        return {}
    if any(
        key in config
        for key in (
            "range_comfort_winter",
            "range_comfort_summer",
            "summer_start",
            "summer_final",
        )
    ):
        return config
    reward_kwargs = config.get("reward_kwargs")
    if isinstance(reward_kwargs, dict):
        return reward_kwargs
    training_config = config.get("training_config")
    if isinstance(training_config, dict):
        reward_kwargs = training_config.get("reward_kwargs")
        if isinstance(reward_kwargs, dict):
            return reward_kwargs
    for value in config.values():
        if isinstance(value, dict):
            nested = _extract_reward_kwargs(value)
            if nested:
                return nested
    return {}


def _coerce_pair(value, default_pair):
    """Converteix una parella de valors."""
    if isinstance(value, (list, tuple)) and len(value) >= 2:
        try:
            return float(value[0]), float(value[1])
        except (TypeError, ValueError):
            pass
    return default_pair


def _coerce_month_day(value, default_pair):
    """Converteix una parella de mes i dia."""
    if isinstance(value, (list, tuple)) and len(value) >= 2:
        try:
            return int(value[0]), int(value[1])
        except (TypeError, ValueError):
            pass
    return default_pair


def _comfort_bounds_from_reward_kwargs(df: pd.DataFrame, comfort_config=None):
    """Retorna els límits de comoditat de reward_kwargs. Els punts de consigna del termòstat són controls,
    no són criteris de confort, de manera que aquí no s'utilitzen intencionadament.
    """
    reward_kwargs = _extract_reward_kwargs(comfort_config)
    winter_lower, winter_upper = _coerce_pair(
        reward_kwargs.get("range_comfort_winter"),
        DEFAULT_WINTER_COMFORT_RANGE,
    )
    summer_lower, summer_upper = _coerce_pair(
        reward_kwargs.get("range_comfort_summer"),
        DEFAULT_SUMMER_COMFORT_RANGE,
    )
    summer_start = _coerce_month_day(
        reward_kwargs.get("summer_start"),
        DEFAULT_SUMMER_START,
    )
    summer_final = _coerce_month_day(
        reward_kwargs.get("summer_final"),
        DEFAULT_SUMMER_FINAL,
    )

    months = _month_series(df).astype(int)
    if "day_of_month" in df.columns:
        days = pd.to_numeric(df["day_of_month"], errors="coerce").fillna(1).astype(int)
    elif isinstance(df.index, pd.DatetimeIndex):
        days = pd.Series(df.index.day, index=df.index).astype(int)
    else:
        days = pd.Series(1, index=df.index)

    current = months * 100 + days
    start_code = summer_start[0] * 100 + summer_start[1]
    final_code = summer_final[0] * 100 + summer_final[1]
    if start_code <= final_code:
        is_summer = (current >= start_code) & (current <= final_code)
    else:
        is_summer = (current >= start_code) | (current <= final_code)

    lower = np.where(is_summer, summer_lower, winter_lower)
    upper = np.where(is_summer, summer_upper, winter_upper)
    return pd.Series(lower, index=df.index), pd.Series(upper, index=df.index)


def _categorize_comfort(df: pd.DataFrame, comfort_config=None) -> pd.DataFrame:
    """
    Afegeix:
      - air_temperature (si cal)
      - month/hour (si cal)
      - comfort_cat: 'Massa fred' | 'En confort' | 'Massa calor'
    """
    base = _ensure_datetime_index(df).copy()
    if "month" not in base.columns:
        base["month"] = base.index.month
    if "hour" not in base.columns:
        base["hour"] = base.index.hour

    base = _ensure_air_temperature(base)  # crea 'air_temperature' si no existeix
    lower, upper = _comfort_bounds_from_reward_kwargs(base, comfort_config)

    if "air_temperature" not in base.columns:
        if "temp_violation" in base.columns:
            violation = pd.to_numeric(base["temp_violation"], errors="coerce")
            valid = violation.notna()
            in_comfort = valid & (violation <= 1e-9)
            cat = pd.Series(pd.NA, index=base.index, dtype="object")
            cat.loc[in_comfort] = "En confort"
            cat.loc[valid & ~in_comfort] = "Massa calor"
            base["comfort_cat"] = pd.Categorical(
                cat,
                categories=["Massa fred", "En confort", "Massa calor"],
                ordered=True,
            )
            return base
        base["comfort_cat"] = pd.Categorical(
            [np.nan] * len(base),
            categories=["Massa fred", "En confort", "Massa calor"],
            ordered=True,
        )
        return base

    t_in = pd.to_numeric(base["air_temperature"], errors="coerce")
    lower = pd.to_numeric(lower, errors="coerce")
    upper = pd.to_numeric(upper, errors="coerce")
    if "temp_violation" in base.columns:
        violation = pd.to_numeric(base["temp_violation"], errors="coerce")
    else:
        violation = np.maximum(lower - t_in, 0) + np.maximum(t_in - upper, 0)
        base["temp_violation"] = violation

    valid = violation.notna()
    in_comfort = valid & (violation <= 1e-9)
    cold = valid & (~in_comfort) & (t_in < lower)
    cat = pd.Series(pd.NA, index=base.index, dtype="object")
    cat.loc[in_comfort] = "En confort"
    cat.loc[cold] = "Massa fred"
    cat.loc[valid & ~in_comfort & ~cold] = "Massa calor"
    base["comfort_cat"] = pd.Categorical(
        cat, categories=["Massa fred", "En confort", "Massa calor"], ordered=True
    )
    return base


def _ensure_temp_violation_column(df: pd.DataFrame, comfort_config=None) -> pd.DataFrame:
    """Assegura una columna `temp_violation` per pas de temps.

    Utilitza el valor del monitor quan existeix. En cas contrari, deriva la violació de
    reward_kwargs rangs de comoditat, tornant als rangs predeterminats de Sinergym.
    Els punts de consigna del termòstat són controls i no s'utilitzen intencionadament aquí.
    """
    base = _ensure_datetime_index(df).copy()
    if "month" not in base.columns:
        base["month"] = base.index.month
    if "hour" not in base.columns:
        base["hour"] = base.index.hour
    base = _ensure_air_temperature(base)

    if "temp_violation" in base.columns:
        base["temp_violation"] = pd.to_numeric(base["temp_violation"], errors="coerce")
        return base

    if "air_temperature" not in base.columns:
        base["temp_violation"] = np.nan
        return base

    lower, upper = _comfort_bounds_from_reward_kwargs(base, comfort_config)
    t_in = pd.to_numeric(base["air_temperature"], errors="coerce")
    lower = pd.to_numeric(lower, errors="coerce")
    upper = pd.to_numeric(upper, errors="coerce")
    base["temp_violation"] = np.maximum(lower - t_in, 0) + np.maximum(t_in - upper, 0)
    return base


# Gràfic del percentatge d'hores en confort.
def make_comfort_compliance(
    obs: pd.DataFrame,
    mode: str,
    season: str,
    comfort_scope: str = "all",
    comfort_config=None,
) -> go.Figure:
    """Crea el gràfic de percentatge d'hores dins el rang de confort."""
    scoped_obs = filter_by_comfort_scope(obs, comfort_scope)
    scope_label = comfort_scope_label(comfort_scope)
    df = _categorize_comfort(scoped_obs, comfort_config)

    # Només apliquem filtre d’estació quan té sentit
    if mode in ("raw", "hour", "day"):
        df = filter_by_season(df, season)

    categories = ["Massa fred", "En confort", "Massa calor"]
    colors_map = {
        "Massa fred": STUDIO_CHART_COLORS["comfort_cold"],
        "En confort": STUDIO_CHART_COLORS["comfort_ok"],
        "Massa calor": STUDIO_CHART_COLORS["comfort_hot"],
    }

    if df.empty:
        return go.Figure()

    if mode == "raw":
        counts = df["comfort_cat"].value_counts().reindex(categories).fillna(0)
        fig = px.pie(values=counts.values, names=counts.index, hole=0.45,
                     title=f"% d’hores en confort – Donut ({season} | {scope_label})",
                     color=counts.index, color_discrete_map=colors_map)
        fig.update_traces(textinfo="percent+label", marker=dict(line=dict(color="#ffffff", width=2)))
        fig.update_layout(legend_title_text="")
        return _apply_figure_theme(fig)

    # Agreguem en barres apilades per veure el repartiment de confort.
    if mode == "hour":
        idx = df["hour"].astype(int); x_title="Hora del dia"; x_order=list(range(24))
    elif mode == "day":
        df = df.copy(); df["__day__"] = df.index.floor("D").astype(str)
        idx = df["__day__"]; x_title="Dia"; x_order=None
    elif mode == "month":
        # IMPORTANT: en month NO fem filtre per estació
        df = _categorize_comfort(scoped_obs, comfort_config).copy()
        df["__month__"] = df.index.month
        idx = df["__month__"].astype(int); x_title="Mes"; x_order=list(range(1,13))
    else:
        return make_comfort_compliance(
            scoped_obs,
            "raw",
            season,
            comfort_scope=comfort_scope,
            comfort_config=comfort_config,
        )

    tbl = pd.crosstab(idx, df["comfort_cat"])
    tbl = (tbl.div(tbl.sum(axis=1).replace(0, np.nan), axis=0) * 100).fillna(0.0)
    tbl = tbl.reindex(columns=categories, fill_value=0.0)
    if x_order is not None:
        tbl = tbl.reindex(index=x_order, fill_value=0.0)

    fig = go.Figure()
    for cat in categories:
        fig.add_bar(x=tbl.index, y=tbl[cat].values, name=cat, marker_color=colors_map[cat])

    fig.update_layout(
        barmode="stack",
        title=f"% d’hores en confort – {mode.capitalize()} ({'All' if mode=='month' else season} | {scope_label})",
        legend_title_text="",
    )
    fig.update_yaxes(title="% d’hores", range=[0, 100])
    fig.update_xaxes(title=x_title, tickmode="linear", dtick=1 if mode in ("hour","month") else None)
    return _apply_figure_theme(fig)


# Gràfic de violació de confort en barres.
def make_violation_bars(
    obs: pd.DataFrame,
    mode: str,
    season: str,
    comfort_scope: str = "all",
    comfort_config=None,
) -> go.Figure:
    """Crea el gràfic de barres de desviació de confort."""
    scoped_obs = filter_by_comfort_scope(obs, comfort_scope)
    scope_label = comfort_scope_label(comfort_scope)
    base = _ensure_temp_violation_column(scoped_obs, comfort_config)

    fig = go.Figure()

    if mode == "hour":
        df = filter_by_season(base, season)
        g = df.groupby(df.index.hour)["temp_violation"].mean()
        x_vals = list(range(24))
        y_vals = [g.get(h, float("nan")) for h in x_vals]
        _add_violation_bar(fig, x=x_vals, y=y_vals, name="Mitjana horària")
        fig.update_xaxes(title="Hora del dia", tickmode="linear", dtick=1)

    elif mode == "day":
        df = filter_by_season(base, season)
        g = df.groupby(df.index.floor("D"))["temp_violation"].mean()
        fig.add_scatter(
            x=g.index.astype(str), y=g.values, name="Mitjana diària",
            mode="lines",
            line=dict(color=STUDIO_CHART_COLORS["violation"], width=2),
            meta="keep_color"
        )
        fig.update_xaxes(title="Dia")

    elif mode == "month":
        df = base.copy()  # ressaltem amb estació via colors, no filtrem
        g = df.groupby(df.index.month)["temp_violation"].mean()
        x_vals = list(range(1, 12 + 1))
        y_vals = [g.get(m, float("nan")) for m in x_vals]
        c_viol = STUDIO_CHART_COLORS["violation"]
        marker_colors = _seasonal_marker_colors(x_vals, season, active_color=c_viol)
        _add_violation_bar(fig, x=x_vals, y=y_vals, name="Mitjana mensual", color=marker_colors)
        fig.update_xaxes(title="Mes", tickmode="linear", dtick=1)

    elif mode == "raw":
        df = filter_by_season(base, season)
        _add_violation_bar(fig, x=np.arange(len(df)), y=df["temp_violation"].values, name="Registres")
        fig.update_xaxes(title="Time (step)")

    fig.update_yaxes(title="Temperature Violation (°C)")
    fig.update_layout(
        title=f"Temperature Violation (Bars) – {mode.capitalize()} ({season} | {scope_label})"
    )
    return fig


def _fix_colors_for_visibility(fig, *, kind: str, mode: str):
    """
    Aplica colors/contorns més marcats a barres en 'daily' i 'raw'.
    kind: 'violation' | 'hvac' | qualsevol altre (ignora).
    """
    if mode not in ("day", "raw"):
        return fig

    try:
        for tr in getattr(fig, "data", []):
            if getattr(tr, "type", "") != "bar":
                continue
            default_color = "#1f2937" if kind == "violation" else "#264653"
            color = getattr(getattr(tr, "marker", None), "color", None)
            if color is None or (isinstance(color, str) and color.lower() in {
                "lightgray", "#d3d3d3", "rgb(211,211,211)", "rgba(211,211,211,1.0)"
            }):
                tr.marker.color = default_color
            tr.marker.line = dict(color="rgba(0,0,0,0.25)", width=0.4)
            tr.opacity = 0.95
    except Exception:
        pass
    return fig


# Línia única de violació de temperatura amb els filtres actius.
def make_violation_single_line(
    obs: pd.DataFrame,
    mode: str,
    season: str,
    comfort_config=None,
) -> go.Figure:
    """Crea la línia de desviació de temperatura."""
    base = _ensure_temp_violation_column(obs, comfort_config)

    fig = go.Figure()
    violation_color = STUDIO_CHART_COLORS["violation"]

    if mode in ("hour", "day", "month"):
        data = _series_for_mode(base, "temp_violation", mode, season)
        if data is not None:
            x_vals, y_vals, hover, trace_mode = data
            marker_colors = (
                _seasonal_marker_colors(x_vals, season, active_color=violation_color, muted_color="#7f7f7f")
                if mode == "month"
                else violation_color
            )
            fig.add_trace(go.Scatter(
                x=x_vals,
                y=y_vals,
                mode="lines+markers" if mode in ("hour", "day", "month") else trace_mode,
                name={"hour": "Hourly Mean", "day": "Daily Mean", "month": "Monthly Mean"}[mode],
                line=dict(color=violation_color, width=2.6),
                marker=dict(color=marker_colors, size=6),
                meta="keep_color",
                hovertext=hover,
                hovertemplate=(
                    "%{hovertext}<br>%{y:.2f}°C<extra></extra>"
                    if hover is not None
                    else None
                ),
            ))
        _apply_mode_xaxis(fig, mode)

    elif mode == "raw":
        # En raw dibuixem primer tot l'any i, si hi ha filtre de temporada, el ressaltem
        # a sobre per no perdre el context.
        df_all = filter_by_season(base, "All")
        x_all, y_all, hover_all = _raw_axis_data(df_all, "temp_violation")
        fig.add_trace(go.Scatter(x=x_all, y=y_all, mode="lines", name="All Records",
                                 line=dict(color=violation_color, width=1.8), meta="keep_color", opacity=0.92,
                                 hovertext=hover_all, hovertemplate="%{hovertext}<br>%{y:.2f}°C<extra></extra>"))
        if season != "All":
            df_season = filter_by_season(base, season)
            if not df_season.empty:
                x_s, y_s, hover_s = _raw_axis_data(df_season, "temp_violation")
                fig.add_trace(go.Scatter(x=x_s, y=y_s, mode="lines",
                                         name=f"{season} Records", line=dict(color=violation_color, width=2.4), meta="keep_color",
                                         hovertext=hover_s, hovertemplate="%{hovertext}<br>%{y:.2f}°C<extra></extra>"))
        fig.update_xaxes(title="Time (step)")

    fig.update_yaxes(title="Temperature Violation (°C)")
    fig.update_layout(title=f"Temperature Violation – {mode.capitalize()} ({season})")
    return fig
\end{lstlisting}

\Needspace{8\baselineskip}

\subsection{\texorpdfstring{\texttt{backend/grafics/comfort\_scope.py}}{backend/grafics/comfort\_scope.py}}

\begin{lstlisting}[style=studioPython,caption={backend/grafics/comfort\_scope.py},label={lst:backend-grafics-comfort-scope-py}]
"""Filtres d'abast de confort compartits per KPI, gràfics i informes.

Aquest fitxer decideix si una mètrica de confort s'ha de calcular sobre totes les hores o només sobre
les hores amb ocupació, i manté el mateix criteri al dashboard i als informes descarregables.
"""

from __future__ import annotations

from typing import Optional

import pandas as pd


def derive_occupied_mask(obs: pd.DataFrame) -> Optional[pd.Series]:
    """Inferir una màscara booleana per als registres ocupats quan les dades ho proporcionen."""
    if obs is None or obs.empty:
        return None

    explicit_cols = (
        "occupied_hour",
        "people_occupant",
        "people_occupants",
        "occupancy",
        "occupants",
        "is_occupied",
    )
    for col in explicit_cols:
        if col in obs.columns:
            values = pd.to_numeric(obs[col], errors="coerce").fillna(0.0)
            return pd.Series(values > 0.0, index=obs.index)

    fuzzy_cols = [
        c for c in obs.columns
        if any(token in c.lower() for token in ("occupant", "occupancy", "people"))
    ]
    if fuzzy_cols:
        values = (
            obs[fuzzy_cols]
            .apply(pd.to_numeric, errors="coerce")
            .fillna(0.0)
            .sum(axis=1)
        )
        return pd.Series(values > 0.0, index=obs.index)

    return None


def has_occupied_data(obs: pd.DataFrame) -> bool:
    """Retorna si les observacions contenen prou dades d'ocupació per filtrar per àmbit."""
    return derive_occupied_mask(obs) is not None


def filter_by_comfort_scope(obs: pd.DataFrame, comfort_scope: str) -> pd.DataFrame:
    """Retorna totes les files o només les files ocupades en funció de l'àmbit seleccionat."""
    if comfort_scope != "occupied":
        return obs.copy()

    mask = derive_occupied_mask(obs)
    if mask is None:
        return obs.copy()
    return obs.loc[mask].copy()


def comfort_scope_label(comfort_scope: str) -> str:
    """Retorna l'etiqueta de visualització associada a un identificador d'abast de confort."""
    return "Occupied Hours" if comfort_scope == "occupied" else "All Hours"
\end{lstlisting}

\Needspace{8\baselineskip}

\subsection{\texorpdfstring{\texttt{backend/grafics/control.py}}{backend/grafics/control.py}}

\begin{lstlisting}[style=studioPython,caption={backend/grafics/control.py},label={lst:backend-grafics-control-py}]
"""Figures de control i acció per a HVAC i terra radiant."""

from __future__ import annotations

import pandas as pd
import plotly.graph_objects as go
from plotly.subplots import make_subplots

from backend.grafics.style import pick_semantic_trace_color
from backend.grafics.time_utils import (
    _apply_mode_xaxis,
    _ensure_datetime_index,
    _mode_axis_config,
    _series_for_mode,
)


def _mean_series(base: pd.DataFrame, columns: list[str]) -> pd.Series:
    """Retorna la mitjana numèrica per files de les columnes seleccionades."""
    numeric = base[columns].apply(pd.to_numeric, errors="coerce")
    return numeric.mean(axis=1)


def _add_radiant_control_trace(
    fig: go.Figure,
    base: pd.DataFrame,
    series_name: str,
    label: str,
    secondary_y: bool,
    suffix: str,
    mode: str,
    season: str,
) -> None:
    """Afegeix un traça de control radiant a la figura."""
    data = _series_for_mode(base, series_name, mode, season)
    if data is None:
        return
    x_values, y_values, hover, trace_mode = data
    color = pick_semantic_trace_color(label)
    scatter_kwargs = dict(
        x=x_values,
        y=y_values,
        mode=trace_mode,
        name=label,
        line=dict(color=color, width=2.8, shape="hv" if mode in ("hour", "raw") else "linear"),
    )
    if "markers" in trace_mode:
        scatter_kwargs["marker"] = dict(color=color, size=6)
    if hover is not None:
        scatter_kwargs["hovertext"] = hover
        scatter_kwargs["hovertemplate"] = f"%{{hovertext}}<br>%{{y:.2f}} {suffix}<extra></extra>"
    else:
        scatter_kwargs["hovertemplate"] = f"%{{x}}<br>%{{y:.2f}} {suffix}<extra></extra>"
    fig.add_trace(go.Scatter(**scatter_kwargs), secondary_y=secondary_y)


def _pretty_action_name(column_name: str) -> str:
    """Retorna una etiqueta de visualització per a una columna d'acció."""
    return column_name.replace("_", " ").replace("radiant ", "radiant: ").title()


def _add_agent_action_trace(
    fig: go.Figure,
    base: pd.DataFrame,
    column: str,
    row: int,
    color: str,
    temp_columns: list[str],
    mode: str,
    season: str,
) -> None:
    """Afegeix una traça d'acció a la figura de diverses files."""
    data = _series_for_mode(base, column, mode, season)
    if data is None:
        return
    x_values, y_values, hover, trace_mode = data
    is_temp = column in temp_columns
    scatter_kwargs = dict(
        x=x_values,
        y=y_values,
        mode=trace_mode,
        name=_pretty_action_name(column),
        legendgroup="agent_actions",
        showlegend=False,
        line=dict(color=color, width=2.5, shape="hv" if mode in ("hour", "raw") else "linear"),
    )
    if "markers" in trace_mode:
        scatter_kwargs["marker"] = dict(color=color, size=6)
    suffix = " C" if is_temp else ""
    if hover is not None:
        scatter_kwargs["hovertext"] = hover
        scatter_kwargs["hovertemplate"] = f"%{{hovertext}}<br>%{{y:.2f}}{suffix}<extra></extra>"
    else:
        scatter_kwargs["hovertemplate"] = f"%{{x}}<br>%{{y:.2f}}{suffix}<extra></extra>"
    fig.add_trace(go.Scatter(**scatter_kwargs), row=row, col=1)


def make_radiant_control_plot(obs: pd.DataFrame, mode: str, season: str) -> go.Figure:
    """Crea una figura Plotly per al control radiant."""
    base = _ensure_datetime_index(obs).copy()

    # Els noms dels sensors radiants varien molt entre models. Fem una detecció flexible
    # i després els reduïm a tres sèries llegibles: disponibilitat, entrada i sortida.
    lower_cols = {col: col.lower() for col in base.columns}
    availability_cols = [
        col
        for col, label in lower_cols.items()
        if "radiant" in label
        and ("availavility" in label or "availability" in label or "operation_mode" in label or "operation mode" in label)
    ]
    inlet_cols = [
        col
        for col, label in lower_cols.items()
        if "radiant" in label and "inlet" in label and "temperature" in label
    ]
    outlet_cols = [
        col
        for col, label in lower_cols.items()
        if "radiant" in label and "outlet" in label and "temperature" in label
    ]

    if not availability_cols and not inlet_cols and not outlet_cols:
        return go.Figure()

    if availability_cols:
        # Si hi ha més d'una zona radiant, la mitjana dona una lectura global del mode actiu.
        base["__radiant_availability_pct__"] = _mean_series(base, availability_cols) * 100.0
    if inlet_cols:
        base["__radiant_inlet_temp__"] = _mean_series(base, inlet_cols)
    if outlet_cols:
        base["__radiant_outlet_temp__"] = _mean_series(base, outlet_cols)

    fig = make_subplots(specs=[[{"secondary_y": True}]])

    _add_radiant_control_trace(fig, base, "__radiant_outlet_temp__", "Radiant Outlet Temp", False, "C", mode, season)
    _add_radiant_control_trace(fig, base, "__radiant_inlet_temp__", "Radiant Inlet Temp", False, "C", mode, season)
    _add_radiant_control_trace(fig, base, "__radiant_availability_pct__", "Radiant Operation Mode", True, "%", mode, season)

    _apply_mode_xaxis(fig, mode)

    fig.update_yaxes(title="Temperature (C)", secondary_y=False)
    fig.update_yaxes(title="Operation Mode (%)", range=[-5, 105], secondary_y=True)
    fig.update_layout(title=f"Radiant Control - {mode.capitalize()} ({season})")
    return fig


def make_agent_actions_plot(actions: pd.DataFrame, mode: str, season: str) -> go.Figure:
    """Crea una figura Plotly per a les accions de l'agent."""
    if actions is None or actions.empty:
        return go.Figure()

    time_columns = {"datetime", "timestamp", "month", "day_of_month", "hour", "time_elapsed(hours)"}
    base = _ensure_datetime_index(actions).copy().sort_index(kind="mergesort")
    action_columns = [col for col in base.columns if col not in time_columns]
    if not action_columns:
        return go.Figure()

    for col in action_columns:
        base[col] = pd.to_numeric(base[col], errors="coerce")

    temp_cols = [
        col
        for col in action_columns
        if "temperature" in col.lower() or "temp" in col.lower()
    ]
    binary_cols = [col for col in action_columns if col not in temp_cols]
    ordered_cols = temp_cols + binary_cols
    if not ordered_cols:
        return go.Figure()

    for col in temp_cols:
        valid = base[col].dropna()
        if not valid.empty and valid.min() >= 0 and valid.max() <= 20:
            # Alguns wrappers guarden deltes o consignes normalitzades; si totes cauen en
            # 0..20 les desplacem per representar-les com a consignes de temperatura llegibles.
            base[col] = base[col] + 25.0

    rows = len(ordered_cols)
    # Les accions de temperatura necessiten més alçada que les binàries, on només volem
    # veure si l'actuador està actiu o no.
    row_heights = [1.7 if col in temp_cols else 0.55 for col in ordered_cols]
    fig = make_subplots(
        rows=rows,
        cols=1,
        shared_xaxes=True,
        vertical_spacing=0.018 if rows > 1 else 0.0,
        row_heights=row_heights,
        subplot_titles=[_pretty_action_name(col) for col in ordered_cols],
    )

    palette = ["#0f766e", "#7c3aed", "#2563eb", "#ea580c", "#16a34a", "#dc2626", "#0891b2", "#ca8a04"]
    for row, col in enumerate(ordered_cols, start=1):
        _add_agent_action_trace(fig, base, col, row, palette[(row - 1) % len(palette)], temp_cols, mode, season)

    x_axis_title, x_axis_kwargs = _mode_axis_config(mode)

    for row, col in enumerate(ordered_cols, start=1):
        fig.update_xaxes(**x_axis_kwargs, row=row, col=1)
        if col in temp_cols:
            fig.update_yaxes(title="C", row=row, col=1)
        else:
            if mode == "raw":
                fig.update_yaxes(
                    title="ON/OFF",
                    range=[-0.08, 1.08],
                    tickmode="array",
                    tickvals=[0, 1],
                    ticktext=["OFF", "ON"],
                    row=row,
                    col=1,
                )
            else:
                fig.update_yaxes(
                    title="ON fraction",
                    range=[-0.05, 1.05],
                    tickmode="array",
                    tickvals=[0, 0.5, 1],
                    row=row,
                    col=1,
                )
    fig.update_xaxes(title=x_axis_title, row=rows, col=1)
    figure_height = 260 + 120 * len(temp_cols) + 58 * len(binary_cols)
    fig.update_layout(
        title=f"Agent Actions - {mode.capitalize()} ({season})",
        height=max(360, figure_height),
        showlegend=False,
    )
    return fig
\end{lstlisting}

\Needspace{8\baselineskip}

\subsection{\texorpdfstring{\texttt{backend/grafics/data\_loader.py}}{backend/grafics/data\_loader.py}}

\begin{lstlisting}[style=studioPython,caption={backend/grafics/data\_loader.py},label={lst:backend-grafics-data-loader-py}]
"""Carrega i normalitza els fitxers CSV de resultats per a panells i informes.

La pàgina de resultats i el backend de l'informe passen per aquí per llegir el progrés de Sinergym,
arreglar buits de mètriques habituals, estandarditzar noms de columnes i retornar DataFrames preparats
per generar KPIs i figures Plotly.
"""

from functools import lru_cache
from pathlib import Path

import pandas as pd

from backend.grafics.observation_columns import (
    add_meter_kwh_columns,
    normalize_observation_columns,
)
from backend.grafics.time_utils import sanitize_observation_time_columns

_DEFAULT_PROGRESS_HEADERS = [
    "episode_num",
    "mean_reward",
    "std_reward",
    "mean_reward_comfort_term",
    "std_reward_comfort_term",
    "mean_reward_energy_term",
    "std_reward_energy_term",
    "mean_comfort_penalty",
    "std_comfort_penalty",
    "mean_energy_penalty",
    "std_energy_penalty",
    "mean_temperature_violation",
    "std_temperature_violation",
    "mean_power_demand",
    "std_power_demand",
    "cumulative_power_demand",
    "comfort_violation_time(%)",
    "length(timesteps)",
    "time_elapsed(hours)",
    "terminated",
    "truncated",
]

_NUMERIC_PROGRESS_COLUMNS = tuple(
    column
    for column in _DEFAULT_PROGRESS_HEADERS
    if column not in {"terminated", "truncated"}
)


def _csv_signature(path: str | Path) -> tuple[str, int, int]:
    """Retorna una clau de memòria cau que canvia quan canvia el fitxer CSV."""

    csv_path = Path(path)
    stat = csv_path.stat()
    return str(csv_path), int(stat.st_mtime_ns), int(stat.st_size)


@lru_cache(maxsize=32)
def _read_csv_cached(
    path: str,
    mtime_ns: int,
    size: int,
    mode: str,
    names: tuple[str, ...],
) -> pd.DataFrame:
    """Llegeix un fitxer CSV mitjançant una memòria cau en procés conscient de mtime."""

    _ = (mtime_ns, size)
    if mode == "progress_no_header":
        return pd.read_csv(path, header=None, names=list(names))
    if mode == "low_memory_false":
        return pd.read_csv(path, low_memory=False)
    return pd.read_csv(path)


def _read_csv(path: str | Path, *, mode: str = "default", names: tuple[str, ...] = ()) -> pd.DataFrame:
    """Llegeix un fitxer CSV i en retorna una còpia segura per mutar."""

    csv_path, mtime_ns, size = _csv_signature(path)
    return _read_csv_cached(csv_path, mtime_ns, size, mode, names).copy()


def _fill_missing_or_zero(target: pd.Series, replacement: pd.Series) -> pd.Series:
    """Ompliu els valors objectiu que falten o zero amb una sèrie de substitució numèrica."""

    target = pd.to_numeric(target, errors="coerce")
    replacement = pd.to_numeric(replacement, errors="coerce")
    replace_mask = target.isna() | (
        (target == 0) & replacement.notna() & (replacement != 0)
    )
    result = target.copy()
    result.loc[replace_mask] = replacement.loc[replace_mask]
    return result


def _repair_power_metrics_from_progress(progress: pd.DataFrame) -> pd.DataFrame:
    """Repareu les mètriques de potència de progrés quan només hi hagi valors acumulatius o mitjans."""

    repaired = progress.copy()

    if (
        "cumulative_power_demand" in repaired.columns
        and "length(timesteps)" in repaired.columns
    ):
        derived_mean_from_cumulative = (
            pd.to_numeric(repaired["cumulative_power_demand"], errors="coerce")
            / pd.to_numeric(repaired["length(timesteps)"], errors="coerce").replace(0, pd.NA)
        )
        if "mean_power_demand" in repaired.columns:
            repaired["mean_power_demand"] = _fill_missing_or_zero(
                repaired["mean_power_demand"], derived_mean_from_cumulative
            )
        else:
            repaired["mean_power_demand"] = derived_mean_from_cumulative

    if (
        "mean_power_demand" in repaired.columns
        and "length(timesteps)" in repaired.columns
    ):
        derived_cumulative = (
            pd.to_numeric(repaired["mean_power_demand"], errors="coerce")
            * pd.to_numeric(repaired["length(timesteps)"], errors="coerce")
        )
        if "cumulative_power_demand" in repaired.columns:
            repaired["cumulative_power_demand"] = _fill_missing_or_zero(
                repaired["cumulative_power_demand"], derived_cumulative
            )
        else:
            repaired["cumulative_power_demand"] = derived_cumulative

    return repaired


def _drop_repeated_header_rows(data: pd.DataFrame) -> pd.DataFrame:
    """Elimina capçaleres CSV repetides accidentalment dins dels fitxers de registre."""

    if data.empty:
        return data

    header_mask = pd.Series(False, index=data.index)
    for column in data.columns:
        header_mask |= data[column].astype(str).str.strip().eq(str(column))

    if not header_mask.any():
        return data

    return data.loc[~header_mask].reset_index(drop=True)


def _coerce_progress_numeric_columns(progress: pd.DataFrame) -> pd.DataFrame:
    """Converteix les mètriques de progrés conegudes en valors numèrics."""

    coerced = progress.copy()
    for column in _NUMERIC_PROGRESS_COLUMNS:
        if column in coerced.columns:
            coerced[column] = pd.to_numeric(coerced[column], errors="coerce")
    return coerced


def load_data(progress_path, obs_path):
    """Carrega els fitxers de progrés i observació CSV utilitzats pel panell de resultats."""

    progress = _read_csv(progress_path)

    # Comprova si el progress.csv té capçaleres que falten (comú quan s'afegeixen registres)
    if "mean_reward" not in progress.columns and "episode_num" not in progress.columns:
        if len(progress.columns) == len(_DEFAULT_PROGRESS_HEADERS):
            progress = _read_csv(
                progress_path,
                mode="progress_no_header",
                names=tuple(_DEFAULT_PROGRESS_HEADERS),
            )

    progress = _drop_repeated_header_rows(progress)
    progress = _coerce_progress_numeric_columns(progress)
    progress = _repair_power_metrics_from_progress(progress)

    obs = _read_csv(obs_path)

    infos_path = Path(obs_path).with_name("infos.csv")
    if infos_path.exists():
        try:
            infos = _read_csv(infos_path, mode="low_memory_false")
            if len(infos) == len(obs):
                for col in infos.columns:
                    if col not in obs.columns:
                        obs[col] = infos[col]
        except Exception:
            pass

    obs = sanitize_observation_time_columns(obs)
    obs = normalize_observation_columns(obs)
    obs = add_meter_kwh_columns(obs)

    return progress, obs
\end{lstlisting}

\Needspace{8\baselineskip}

\subsection{\texorpdfstring{\texttt{backend/grafics/energy\_units.py}}{backend/grafics/energy\_units.py}}

\begin{lstlisting}[style=studioPython,caption={backend/grafics/energy\_units.py},label={lst:backend-grafics-energy-units-py}]
"""Conversions d'energia, preu i unitats de bateria per als gràfics."""

from __future__ import annotations

import pandas as pd

from backend.grafics.time_utils import _auto_scale_series

def _energy_price_series_eur_per_kwh(df: pd.DataFrame, col: str) -> pd.Series:
    """Preu de l'energia sèrie eur per kwh."""
    price = pd.to_numeric(df[col], errors="coerce")
    label = col.lower()
    if "eur_per_kwh" in label or "price_eur_per_kwh" in label or "energy_price_kwh" in label:
        return price

    median_abs = price.abs().dropna().median()
    if pd.notna(median_abs) and median_abs > 5:
        return price / 1000.0
    return price


def _energy_to_kwh(values: pd.Series, unit_kind: str, timestep_hours: float | None = None) -> pd.Series:
    """Energy to kwh."""
    numeric = pd.to_numeric(values, errors="coerce")
    if unit_kind == "kwh":
        return numeric
    if unit_kind == "joule":
        return numeric / 3_600_000.0
    if unit_kind == "watt":
        hours = 0.0 if timestep_hours is None else float(timestep_hours)
        return (numeric / 1000.0) * hours
    return numeric


def _battery_state_display_series(
    base: pd.DataFrame,
    state_col: str,
) -> tuple[pd.Series, str, str]:
    """Retorna la sèrie de visualització, la unitat i l'etiqueta per a les sortides de l'estat de la bateria."""
    raw = pd.to_numeric(base[state_col], errors="coerce")
    label = state_col.lower()

    if label in {"storage_battery_charge_state"}:
        # EnergyPlus ElectricLoadCenter: Emmagatzematge: Informes de l'estat de càrrega de la bateria en Ah.
        return raw, "Ah", "Battery Charge State"

    if (
        "fraction" in label
        or "soc" in label
        or label in {"battery_state_of_charge", "state_of_charge"}
    ):
        finite = raw.dropna()
        if not finite.empty and finite.abs().max() <= 1.5:
            return raw * 100.0, "%", "Battery State of Charge"
        return raw, "%", "Battery State of Charge"

    if label in {"battery_charge_state", "storage_charge_state", "electric_storage_charge_state"}:
        # EnergyPlus ElectricLoadCenter: Emmagatzematge: estat de càrrega dels informes simples en J.
        return raw / 3_600_000.0, "kWh", "Battery Stored Energy"

    finite_abs = raw.abs().dropna()
    if not finite_abs.empty and finite_abs.median() > 1e5:
        return raw / 3_600_000.0, "kWh", "Battery Stored Energy"

    scaled, suffix = _auto_scale_series(raw)
    unit = f"raw {suffix}".strip() if suffix else "raw"
    return scaled, unit, "Battery State"
\end{lstlisting}

\Needspace{8\baselineskip}

\subsection{\texorpdfstring{\texttt{backend/grafics/episode.py}}{backend/grafics/episode.py}}

\begin{lstlisting}[style=studioPython,caption={backend/grafics/episode.py},label={lst:backend-grafics-episode-py}]
"""Xifres del quadre de comandament a nivell d'episodi."""

from __future__ import annotations

import numpy as np
import pandas as pd
import plotly.graph_objects as go
from plotly.subplots import make_subplots

from backend.grafics.style import STUDIO_CHART_COLORS, _apply_figure_theme

def _resolve_episode_mean_power(progress: pd.DataFrame) -> pd.Series | None:
    """Resol la potència mitjana de l'episodi."""
    mean_power = None
    if "mean_power_demand" in progress.columns:
        mean_power = pd.to_numeric(progress["mean_power_demand"], errors="coerce")

    if mean_power is None:
        mean_power = pd.Series(np.nan, index=progress.index, dtype=float)

    derived_mean = None
    if (
        "cumulative_power_demand" in progress.columns
        and "length(timesteps)" in progress.columns
    ):
        derived_mean = (
            pd.to_numeric(progress["cumulative_power_demand"], errors="coerce")
            / pd.to_numeric(progress["length(timesteps)"], errors="coerce").replace(0, np.nan)
        )

    if derived_mean is not None:
        replace_mask = mean_power.isna() | (
            (mean_power == 0) & derived_mean.notna() & (derived_mean != 0)
        )
        mean_power = mean_power.copy()
        mean_power.loc[replace_mask] = derived_mean.loc[replace_mask]

    if mean_power.isna().all():
        return None

    return mean_power

# Mètriques d'episodi.
def make_episode_metrics(progress: pd.DataFrame) -> go.Figure:
    """Mètriques d'episodi amb 3 subtrames:
      1) Recompensa
      2) Violació de la temperatura
      3) Demanda mitjana de potència
    """
    x = progress["episode_num"] if "episode_num" in progress.columns else (progress.index + 1)
    episode_vals = x.tolist()
    episode_ticks = [str(v) for v in episode_vals]

    fig = make_subplots(
        rows=3, cols=1,
        subplot_titles=[
            "Reward per Episode",
            "Average Temperature Violation (°C)",
            "Mean Power Demand per Episode (W)",
        ],
    )

    if "mean_reward" in progress.columns:
        fig.add_trace(
            go.Scatter(
                x=x,
                y=progress["mean_reward"],
                mode="lines+markers",
                name="Reward",
                line=dict(color=STUDIO_CHART_COLORS["reward"], width=3),
                marker=dict(color=STUDIO_CHART_COLORS["reward"], size=7),
            ),
            row=1,
            col=1,
        )

    if "mean_temperature_violation" in progress.columns:
        fig.add_trace(
            go.Scatter(
                x=x,
                y=progress["mean_temperature_violation"],
                mode="lines+markers",
                name="Violation",
                line=dict(color=STUDIO_CHART_COLORS["violation"], width=3),
                marker=dict(color=STUDIO_CHART_COLORS["violation"], size=7),
            ),
            row=2, col=1,
        )

    power_series = _resolve_episode_mean_power(progress)
    if power_series is not None:
        fig.add_trace(
            go.Scatter(
                x=x,
                y=power_series,
                mode="lines+markers",
                name="Mean Power",
                line=dict(color=STUDIO_CHART_COLORS["consumption"], width=3),
                marker=dict(color=STUDIO_CHART_COLORS["consumption"], size=7),
            ),
            row=3, col=1,
        )

    for r in range(1, 4):
        fig.update_xaxes(
            row=r, col=1,
            title_text="Episode",
            tickmode="array",
            tickvals=episode_vals,
            ticktext=episode_ticks,
        )

    fig.update_layout(height=850, title_text="Episode Metrics", showlegend=False)
    return _apply_figure_theme(fig)
\end{lstlisting}

\Needspace{8\baselineskip}

\subsection{\texorpdfstring{\texttt{backend/grafics/figures\_zones.py}}{backend/grafics/figures\_zones.py}}

\begin{lstlisting}[style=studioPython,caption={backend/grafics/figures\_zones.py},label={lst:backend-grafics-figures-zones-py}]
# -*- coding: utf-8 -*-
"""Detecció i filtratge de zones per als panells de resultats.

Aquest mòdul detecta zones tèrmiques des de la configuració de YAML o de la columna d'observació
patrons, i mapes les columnes de la zona seleccionada als camps genèrics que fan servir
les figures comunes.
"""

from __future__ import annotations

from typing import Dict, List, Optional
import re

import pandas as pd

# Patrons flexibles per detectar variables de zona sense dependre d'un únic prefix.
VAR_PATTERNS: Dict[str, re.Pattern] = {
    # temperatura interior de la zona
    "air_temperature": re.compile(r"(?i)(?:^|_)air[_]?temperature(?:$|_)"),
    # consignes
    "htg_setpoint":   re.compile(r"(?i)(?:^|_)htg[_]?setpoint(?:$|_)|(?:^|_)heating[_]?setpoint(?:$|_)"),
    "clg_setpoint":   re.compile(r"(?i)(?:^|_)clg[_]?setpoint(?:$|_)|(?:^|_)cooling[_]?setpoint(?:$|_)"),
    # violació (si la registres per zona)
    "temp_violation": re.compile(r"(?i)(?:^|_)temp[_]?violation(?:$|_)"),
    # opcional (no es fan servir a les figures, però ajuda a detectar zones)
    "air_humidity":   re.compile(r"(?i)(?:^|_)air[_]?humidity(?:$|_)"),
    "people":         re.compile(r"(?i)(?:^|_)people(?:$|_)|(?:^|_)occupant(?:$|_)"),
}

# Quines variables mapejarem a genèriques per a les figures
GENERIC_MAP_KEYS = ("air_temperature", "htg_setpoint", "clg_setpoint", "temp_violation")


# Detecció de zones a partir del YAML o dels noms de columna.
def _zones_from_yaml(yaml_cfg: Optional[dict]) -> List[str]:
    """Zones definides a YAML sota variables→air_temperature→keys."""
    if not yaml_cfg or not isinstance(yaml_cfg.get("variables"), dict):
        return []
    for var_name, var_def in yaml_cfg["variables"].items():
        if str(var_def.get("variable_names")).lower() == "air_temperature":
            keys = var_def.get("keys")
            if isinstance(keys, list) and keys:
                return [str(k).strip() for k in keys if str(k).strip()]
            if isinstance(keys, str) and keys.strip():
                return [keys.strip()]
    return []


def _try_extract_zone_name(col: str, pattern: re.Pattern) -> Optional[str]:
    """
    Donada una columna i un patró de variable (p.ex. air_temperature),
    retorna el 'nom de zona' eliminant aquesta part.
    Exemple: 'zone1_office_air_temperature' -> 'zone1_office'
    """
    m = pattern.search(col)
    if not m:
        return None
    start, end = m.span()
    # Retirem el separador '_' veí si cal
    prefix = (col[:start].rstrip("_"))
    suffix = (col[end:].lstrip("_"))
    # Preferim el prefix (estil 'zone1_office_air_temperature').
    zone = prefix if prefix else suffix
    zone = zone.strip("_")
    return zone or None


def _zones_from_columns(obs: pd.DataFrame) -> List[str]:
    """Infereix zones a partir de columnes que contenen variables de zona conegudes."""
    candidate_sets: List[set] = []
    for var, pat in VAR_PATTERNS.items():
        cols = [c for c in obs.columns if pat.search(c)]
        zones = set()
        for c in cols:
            z = _try_extract_zone_name(c, pat)
            if z:
                zones.add(z)
        if len(zones) >= 2:
            candidate_sets.append(zones)

    if not candidate_sets:
        # Últim recurs: busca *_air_temperature estrictament al final
        zones = set()
        for c in obs.columns:
            if c.endswith("_air_temperature"):
                zones.add(c[: -len("_air_temperature")])
        return sorted(zones)

    # Tria el conjunt amb més cobertura (més zones detectades)
    zones = max(candidate_sets, key=len)
    return sorted(zones)


def detect_zones(obs: pd.DataFrame, yaml_cfg: Optional[dict] = None) -> List[str]:
    """API pública de detecció de zones."""
    zones = _zones_from_yaml(yaml_cfg)
    if zones:
        return zones
    return _zones_from_columns(obs)


# Mapatge de columnes originals cap a noms genèrics per zona.
def _zone_columns(obs: pd.DataFrame, zone: str) -> Dict[str, str]:
    """
    Per una zona concreta, retorna un mapatge {var_key -> nom_columna} segons VAR_PATTERNS.
    Exemple: {'air_temperature': 'zone1_office_air_temperature', ...}
    """
    mapping: Dict[str, str] = {}
    for key, pat in VAR_PATTERNS.items():
        for c in obs.columns:
            if pat.search(c):
                z = _try_extract_zone_name(c, pat)
                if z and z.lower() == zone.lower():
                    # si hi ha duplicats, manté el primer (ordre del CSV)
                    mapping.setdefault(key, c)
    return mapping


def filter_obs_by_zone(obs: pd.DataFrame, zone: Optional[str], yaml_cfg: Optional[dict] = None) -> pd.DataFrame:
    """Retorna les observacions assignades a la zona seleccionada.

    Es conserven les columnes globals, s'eliminen les columnes que pertanyen a altres zones,
    i els valors de la zona seleccionada s'exposen mitjançant els noms genèrics utilitzats per
    Constructors de figures comuns: ``air_temperature``, ``htg_setpoint``,
    ``clg_setpoint`` i ``temp_violation``.
    """
    zones = detect_zones(obs, yaml_cfg)
    if not zone or len(zones) <= 1 or zone not in zones:
        return obs

    # Construïm llistat de columnes de zona vs globals. No podem fiar-nos només del prefix
    # perquè alguns CSV barregen noms de zona, alias i variables globals al mateix fitxer.
    zone_cols_by_zone: Dict[str, set] = {z: set() for z in zones}
    for key, pat in VAR_PATTERNS.items():
        for c in obs.columns:
            if pat.search(c):
                z = _try_extract_zone_name(c, pat)
                if z in zone_cols_by_zone:
                    zone_cols_by_zone[z].add(c)

    keep_cols: List[str] = []
    for c in obs.columns:
        # Si pertany a alguna zona…
        matched_zone = None
        for z, cols in zone_cols_by_zone.items():
            if c in cols:
                matched_zone = z
                break
        if matched_zone is None:
            # global
            keep_cols.append(c)
        elif matched_zone == zone:
            keep_cols.append(c)
        # si és d'una altra zona, s'omet

    out = obs[keep_cols].copy()

    # Mapem la zona triada a noms genèrics perquè les figures comunes no hagin de saber
    # com es deia originalment cada columna al CSV.
    mapping = _zone_columns(out, zone)
    for key in GENERIC_MAP_KEYS:
        col = mapping.get(key)
        if col and col in out.columns:
            out[key] = out[col]

    # Evita que _ensure_air_temperature faci mitjanes entre zones
    # Si tenim la genèrica 'air_temperature', podem eliminar altres *_air_temperature
    if "air_temperature" in out.columns:
        candidates = [c for c in out.columns if VAR_PATTERNS["air_temperature"].search(c)]
        for c in candidates:
            if c != mapping.get("air_temperature"):
                # Drop amb errors='ignore' per si ja no hi és
                out.drop(columns=[c], inplace=True, errors="ignore")

    return out


# Opcions de zona per als selectors de la interfície.
def get_zone_options(obs: pd.DataFrame, yaml_cfg: Optional[dict] = None) -> List[Dict[str, str]]:
    """Retorna les opcions de zona."""
    zones = detect_zones(obs, yaml_cfg)
    return [{"label": z, "value": z} for z in zones]
\end{lstlisting}

\Needspace{8\baselineskip}

\subsection{\texorpdfstring{\texttt{backend/grafics/heatmaps.py}}{backend/grafics/heatmaps.py}}

\begin{lstlisting}[style=studioPython,caption={backend/grafics/heatmaps.py},label={lst:backend-grafics-heatmaps-py}]
"""Xifres de mapes de calor per als resums del panell global i de zona."""

from __future__ import annotations

import numpy as np
import pandas as pd
import plotly.express as px
import plotly.graph_objects as go
from plotly.subplots import make_subplots

from backend.grafics.column_utils import _ensure_air_temperature
from backend.grafics.comfort import _comfort_bounds_from_reward_kwargs
from backend.grafics.time_utils import _ensure_datetime_index, filter_by_season

def make_heatmap(pivot_df: pd.DataFrame, label: str) -> go.Figure:
    """Mapa de calor genèric donat un dataframe pivotat (Mes × Hora)."""
    if pivot_df is None or pivot_df.empty:
        return go.Figure()

    fig = px.imshow(
        pivot_df, origin="lower",
        labels={"x": "Hora", "y": "Mes", "color": label},
        aspect="auto", color_continuous_scale="Viridis",
    )
    fig.update_layout(title=f"Mapa de calor {label} (Mes x Hora)")
    return fig


def make_violation_heatmap_percent(
    obs: pd.DataFrame,
    season: str = "All",
    comfort_config=None,
) -> go.Figure:
    """
    Mapa de calor (Mes x Hora) amb el % de timesteps en violacio de confort.
      - Usa temp_violation si existeix; si no, reward_kwargs o defaults.
      - Respecta filtre d'estació via months (Winter/Spring/Summer/Autumn/All).
      - Escala de color 0..100 (%).
    """
    # Preparem les columnes temporals que necessita el mapa.
    base = _ensure_datetime_index(obs).copy()
    if "month" not in base.columns:
        base["month"] = base.index.month
    if "hour" not in base.columns:
        base["hour"] = base.index.hour
    base = _ensure_air_temperature(base)

    # Marquem cada timestep com a violació o confort.
    if "temp_violation" in base.columns:
        viol_raw = pd.to_numeric(base["temp_violation"], errors="coerce")
        viol = (viol_raw > 0.0).where(viol_raw.notna(), np.nan)
    else:
        lower, upper = _comfort_bounds_from_reward_kwargs(base, comfort_config)
        tin = pd.to_numeric(base["air_temperature"], errors="coerce")
        viol = ((tin < lower) | (tin > upper)).where(tin.notna(), np.nan)

    base = base.assign(__viol__=viol.astype(float))

    # El filtre d'estació s'aplica abans de construir la matriu hora-mes.
    base = filter_by_season(base, season)

    # Cada cel·la guarda el percentatge de violació per mes i hora.
    if base.empty:
        pivot = pd.DataFrame(index=range(1, 13), columns=range(24), dtype=float)
    else:
        pivot = (
            base.groupby(["month", "hour"])["__viol__"]
                .mean()
                .mul(100.0)
                .unstack()                             # columnes = hora
                .reindex(index=range(1, 13))          # mesos 1..12
                .reindex(columns=range(24))           # hores 0..23
        )

    # Pintem el percentatge resultant en una escala comuna de 0 a 100.
    fig = px.imshow(
        pivot,
        origin="lower",
        labels={"x": "Hora", "y": "Mes", "color": "Violacions (%)"},
        aspect="auto",
        color_continuous_scale="Viridis",
    )
    fig.update_coloraxes(cmin=0, cmax=100)
    fig.update_layout(title=f"Violacions de confort (%) – Hora × Mes ({season})")
    return fig

# Heatmaps de temperatura per zona.


def _pretty_zone_temperature_name(name: str) -> str:
    """Retorna una etiqueta de zona llegible des d'un nom de columna de temperatura de l'aire."""
    base = name[: -len("_air_temperature")]
    base = base.replace("_", " ").strip()
    return base if base else name


def _zone_temperature_columns(obs: pd.DataFrame) -> list[tuple[str, str]]:
    """
    Detecta columnes de temperatura d'interior per zona.
    Retorna llista [(nom_zona,"colname"), …].
    Si no hi ha columnes *_air_temperature però existeix 'air_temperature',
    retorna [('All','air_temperature')].
    """
    cols = [c for c in obs.columns if c.endswith("_air_temperature")]
    if cols:
        return [(_pretty_zone_temperature_name(c), c) for c in cols]
    if "air_temperature" in obs.columns:
        return [("All", "air_temperature")]
    return []


def _ensure_month_hour(df: pd.DataFrame) -> pd.DataFrame:
    """Assegura columnes 'month' i 'hour' a partir de l'índex temporal."""
    base = _ensure_datetime_index(df).copy()
    if "month" not in base.columns:
        base["month"] = base.index.month
    if "hour" not in base.columns:
        base["hour"] = base.index.hour
    return base


def _zone_heatmap_layout(nrows: int, ncols: int) -> tuple[float, float, int]:
    """Retorna l'espai i l'alçada segurs de la subparcel·la per a edificis grans de diverses zones."""
    horizontal_spacing = 0.06 if ncols > 1 else 0.0
    if nrows <= 1:
        vertical_spacing = 0.0
    else:
        vertical_spacing = min(0.045, 0.9 / max(nrows - 1, 1))
    height = max(560, 250 * nrows + 120)
    return horizontal_spacing, vertical_spacing, height


def compute_zone_temperature_pivots(
    obs: pd.DataFrame, season: str = "All", agg: str = "mean"
) -> dict[str, pd.DataFrame]:
    """
    Construeix, per a cada zona, un pivot Mes×Hora amb la temperatura d'interior.
    - Respecta el filtre d'estació (via filter_by_season).
    - Reindexa sempre a mesos 1..12 i hores 0..23 (cel·les inexistents -> NaN).
    - 'agg' pot ser 'mean' o 'median'.

    Retorna: {nom_zona: DataFrame (index=1..12 mesos, columns=0..23 hores)}
    """
    base = _ensure_month_hour(_ensure_air_temperature(obs))
    base = filter_by_season(base, season)

    zones = _zone_temperature_columns(base)
    if not zones:
        return {}

    aggfun = "median" if str(agg).lower() == "median" else "mean"

    pivots: dict[str, pd.DataFrame] = {}
    for zname, col in zones:
        s = pd.to_numeric(base[col], errors="coerce")
        g = (
            s.groupby([base["month"].astype(int), base["hour"].astype(int)])
             .agg(aggfun)
             .unstack()  # columnes = hour
        )
        # assegura el grid complet Mes(1..12) × Hora(0..23)
        g = g.reindex(index=range(1, 13)).reindex(columns=range(24))
        pivots[zname] = g

    return pivots


def make_zone_temperature_heatmaps(
    obs: pd.DataFrame,
    season: str = "All",
    agg: str = "mean",
    max_cols: int = 3,
) -> go.Figure:
    """Crea mapes de calor de temperatura interior mensuals per hores per a cada zona.

    La figura generada utilitza una escala de colors compartida entre zones, canònica
    eixos mes/hora, el filtre de temporada seleccionat i la mitjana o la mediana
    agregació. El creador de pivot omet les dades de zones buides o incompletes.

    Paràmetres:
        obs: DataFrame d'observació.
        season: filtre de temporada: ``All``, ``Winter``, ``Spring``, ``Summer`` o
            ``Autumn``.
        agg: mètode d'agregació, ja sigui ``mean`` o ``median``.
        max_cols: nombre màxim de columnes de subgràfic.

    Retorna:
        Una figura Plotly amb un mapa de calor per zona, o un únic mapa de calor quan el
        run només conté una zona.
    """
    pivots = compute_zone_temperature_pivots(obs, season=season, agg=agg)
    if not pivots:
        return go.Figure()

    # Escala comuna: min/max sobre totes les zones (ignorant NaN)
    all_vals = np.concatenate([p.values.flatten() for p in pivots.values()])
    finite_vals = all_vals[np.isfinite(all_vals)]
    if finite_vals.size:
        cmin = float(np.nanmin(finite_vals))
        cmax = float(np.nanmax(finite_vals))
    else:
        cmin, cmax = np.nan, np.nan
    if not np.isfinite(cmin) or not np.isfinite(cmax) or cmin == cmax:
        # reserva segura
        cmin, cmax = 15.0, 30.0

    znames = list(pivots.keys())
    n = len(znames)
    ncols = min(max_cols, n)
    nrows = (n + ncols - 1) // ncols
    horizontal_spacing, vertical_spacing, height = _zone_heatmap_layout(nrows, ncols)

    fig = make_subplots(
        rows=nrows,
        cols=ncols,
        subplot_titles=znames,
        horizontal_spacing=horizontal_spacing,
        vertical_spacing=vertical_spacing,
    )

    # Afegeix un heatmap per zona, compartint escala; una sola colorbar (últim subplot)
    for i, zname in enumerate(znames):
        r = i // ncols + 1
        c = i % ncols + 1
        pivot = pivots[zname]
        showscale = (i == n - 1)  # només a l'últim
        heat = go.Heatmap(
            z=pivot.values,
            x=list(pivot.columns),    # 0..23
            y=list(pivot.index),      # 1..12
            colorscale="Viridis",
            zmin=cmin,
            zmax=cmax,
            colorbar=dict(title="°C"),
            showscale=showscale,
        )
        fig.add_trace(heat, row=r, col=c)
        fig.update_xaxes(title_text="Hora", row=r, col=c, tickmode="linear", dtick=2)
        fig.update_yaxes(title_text="Mes",  row=r, col=c, autorange="reversed")  # 1..12 de dalt a baix

    fig.update_layout(
        title=f"Temperatura interior per zona – Mes × Hora ({season}, {agg})",
        margin=dict(l=60, r=20, t=60, b=40),
        height=height,
    )
    return fig
\end{lstlisting}

\Needspace{8\baselineskip}

\subsection{\texorpdfstring{\texttt{backend/grafics/hvac.py}}{backend/grafics/hvac.py}}

\begin{lstlisting}[style=studioPython,caption={backend/grafics/hvac.py},label={lst:backend-grafics-hvac-py}]
"""HVAC xifres de consum i avaria del comptador."""

from __future__ import annotations

import pandas as pd
import plotly.graph_objects as go

from backend.grafics.plot_helpers import (
    add_energy_bar_trace,
    energy_axis_title,
    energy_series_for_mode,
    raw_hourly_total_series,
)
from backend.grafics.style import STUDIO_CHART_COLORS, _alpha_color, _apply_figure_theme
from backend.grafics.time_utils import (
    _canon_day_total_series,
    _ensure_datetime_index,
    _infer_timestep_hours,
    _raw_axis_data,
    filter_by_season,
)


def _with_hvac_consumption_kwh(df: pd.DataFrame, hvac_col: str, unit_kind: str) -> pd.DataFrame:
    """Afegeix una columna temporal amb el consum HVAC per timestep en kWh."""
    from backend.grafics.observation_columns import convert_hvac_source_to_kwh

    base = df.copy()
    base["__hvac_consumption_kwh__"] = convert_hvac_source_to_kwh(
        base[hvac_col],
        unit_kind,
        _infer_timestep_hours(base),
    )
    return base


def make_hvac_consumption_plot(
    obs: pd.DataFrame,
    mode: str,
    season: str,
    *,
    raw_as_rate: bool = False,
    raw_as_line: bool = False,
) -> go.Figure:
    """Crea una figura Plotly per al consum de climatització."""
    from backend.grafics.observation_columns import (
        HVAC_POWER_COLUMN,
        find_hvac_consumption_source,
        normalize_observation_columns,
    )

    base = _ensure_datetime_index(normalize_observation_columns(obs))

    fig = go.Figure()
    hvac_source = find_hvac_consumption_source(base.columns)
    if hvac_source is None:
        return go.Figure()
    hvac_col, hvac_unit = hvac_source

    if raw_as_rate and mode == "raw" and hvac_unit == "watt":
        # En mode interactiu a vegades volem veure potència instantània, no energia per pas.
        # Si la font és W, podem mostrar-la directament com kW.
        df = filter_by_season(base, season)
        if df.empty:
            return go.Figure()
        x_vals, y_vals, hover = _raw_axis_data(df, HVAC_POWER_COLUMN)
        y_kw = (pd.Series(y_vals, dtype=float) / 1000.0).tolist()
        fig.add_trace(
            go.Scatter(
                x=x_vals,
                y=y_kw,
                mode="lines",
                name="HVAC Demand Rate",
                line=dict(color=STUDIO_CHART_COLORS["consumption"], width=2.0),
                hovertext=hover,
                hovertemplate="%{hovertext}<br>%{y:.2f} kW<extra></extra>",
            )
        )
        fig.update_xaxes(title="Time (step)")
        fig.update_yaxes(title="Demand Rate (kW)")
        fig.update_layout(title=f"HVAC Demand Rate - {mode.capitalize()} ({season})")
        return fig

    base = _with_hvac_consumption_kwh(base, hvac_col, hvac_unit)
    energy_col = "__hvac_consumption_kwh__"

    if mode == "hour":
        add_energy_bar_trace(
            fig,
            base,
            energy_col,
            label="HVAC Consumption",
            color=STUDIO_CHART_COLORS["consumption"],
            mode=mode,
            season=season,
        )
        fig.update_xaxes(title="Hour of Day", tickmode="linear", dtick=1)
        y_axis_title = energy_axis_title(mode)

    elif mode == "day":
        series_data = energy_series_for_mode(base, energy_col, mode, season)
        if series_data is None:
            return go.Figure()
        x_vals, y_vals, hover = series_data
        fig.add_trace(
            go.Scatter(
                x=x_vals,
                y=y_vals,
                mode="lines",
                name="HVAC Consumption",
                line=dict(color=STUDIO_CHART_COLORS["consumption"], width=2.4),
                hovertext=hover,
                hovertemplate="%{hovertext}<br>%{y:.2f} kWh<extra></extra>",
            )
        )
        fig.update_xaxes(title="Day (1..365)", range=[1, 365], tickmode="linear", tick0=1, dtick=10)
        y_axis_title = energy_axis_title(mode)

    elif mode == "month":
        add_energy_bar_trace(
            fig,
            base,
            energy_col,
            label="HVAC Consumption",
            color=STUDIO_CHART_COLORS["consumption"],
            mode=mode,
            season=season,
        )
        fig.update_xaxes(title="Month", tickmode="linear", dtick=1)
        y_axis_title = energy_axis_title(mode)

    elif mode == "raw":
        # La vista raw es reagrupa a hores perquè els CSV de simulació poden tenir pas de
        # 15 minuts i una barra per timestep seria massa sorollosa.
        df = filter_by_season(base, season)
        if df.empty:
            return go.Figure()
        x_vals, y_vals, hover = raw_hourly_total_series(df, energy_col)
        if not y_vals:
            return go.Figure()

        if raw_as_line:
            fig.add_trace(
                go.Scatter(
                    x=x_vals,
                    y=y_vals,
                    mode="lines",
                    name="HVAC Consumption",
                    line=dict(color=STUDIO_CHART_COLORS["consumption"], width=2.2),
                    meta="keep_color",
                    hovertext=hover,
                    hovertemplate="%{hovertext}<br>%{y:.2f} kWh<extra></extra>",
                )
            )
        else:
            fig.add_trace(go.Bar(
                x=x_vals,
                y=y_vals,
                name="HVAC Consumption",
                marker_color=STUDIO_CHART_COLORS["consumption"],
                hovertext=hover,
                hovertemplate="%{hovertext}<br>%{y:.2f} kWh<extra></extra>",
            ))
        fig.update_xaxes(title="Time (hour)")
        y_axis_title = energy_axis_title(mode, raw_title="Hourly Consumption (kWh)")
    else:
        y_axis_title = "Consumption (kWh)"

    fig.update_yaxes(title=y_axis_title)
    fig.update_layout(title=f"HVAC Consumption - {mode.capitalize()} ({season})")
    return fig


# Desglossament del consum HVAC per meters disponibles.

def make_hvac_meter_breakdown_plot(
    obs: pd.DataFrame,
    mode: str,
    season: str,
) -> go.Figure:
    """Crea una figura Plotly per a l'avaria del comptador de climatització."""
    from backend.grafics.observation_columns import (
        DETAILED_HVAC_METER_KWH_COLUMNS,
        DETAILED_HVAC_METER_LABELS,
        ELECTRIC_HVAC_METER_ALIASES,
        HVAC_ENERGY_METER_COLUMNS,
        HVAC_POWER_COLUMN,
        add_meter_kwh_columns,
        convert_hvac_source_to_kwh,
        find_hvac_consumption_source,
        normalize_observation_columns,
    )

    base = _ensure_datetime_index(add_meter_kwh_columns(normalize_observation_columns(obs)))
    meter_series = []
    # Alguns IDF només tenen el total HVAC i altres exposen meters separats per fans,
    # pumps, heating, cooling, etc. Afegim només les sèries que realment tenen energia.
    for meter_alias, kwh_column in DETAILED_HVAC_METER_KWH_COLUMNS.items():
        if kwh_column not in base.columns:
            continue
        numeric_values = pd.to_numeric(base[kwh_column], errors="coerce")
        if not numeric_values.dropna().abs().gt(1e-12).any():
            continue
        meter_series.append(
            (
                kwh_column,
                DETAILED_HVAC_METER_LABELS.get(meter_alias, meter_alias),
            )
        )
    fig = go.Figure()
    colors = {
        "Natural Gas HVAC": "#92400e",
        "Heating Electricity": STUDIO_CHART_COLORS["heating"],
        "Cooling Electricity": STUDIO_CHART_COLORS["cooling"],
        "Fans Electricity": STUDIO_CHART_COLORS["neutral"],
        "Pumps Electricity": "#0891b2",
        "Heat Rejection Electricity": "#f97316",
        "Humidifier Electricity": "#7c3aed",
        "Heat Recovery Electricity": "#16a34a",
        "Other / Unmetered HVAC Electricity": STUDIO_CHART_COLORS["consumption_muted"],
    }
    electric_detail_columns = [
        DETAILED_HVAC_METER_KWH_COLUMNS[meter_alias]
        for meter_alias in ELECTRIC_HVAC_METER_ALIASES
        if DETAILED_HVAC_METER_KWH_COLUMNS.get(meter_alias) in base.columns
    ]

    hvac_source = find_hvac_consumption_source(base.columns)
    if hvac_source is not None:
        hvac_col, hvac_unit = hvac_source
        hvac_col_is_total_electric = (
            hvac_col in HVAC_ENERGY_METER_COLUMNS
            or hvac_col.lower() in {column.lower() for column in HVAC_ENERGY_METER_COLUMNS}
            or hvac_col == HVAC_POWER_COLUMN
        )
        if hvac_col_is_total_electric:
            # Quan tenim total elèctric i també submeters, el "Other" és el residu.
            # Això ajuda a detectar consum HVAC que EnergyPlus no ha separat per categoria.
            base["__hvac_total_electricity_kwh__"] = convert_hvac_source_to_kwh(
                base[hvac_col],
                hvac_unit,
                _infer_timestep_hours(base),
            )
            if electric_detail_columns:
                electric_detail_sum = pd.DataFrame(
                    {
                        column: pd.to_numeric(base[column], errors="coerce")
                        for column in electric_detail_columns
                    },
                    index=base.index,
                ).sum(axis=1, min_count=1).fillna(0.0)
            else:
                electric_detail_sum = pd.Series(0.0, index=base.index)
            base["__other_hvac_electricity_kwh__"] = (
                pd.to_numeric(base["__hvac_total_electricity_kwh__"], errors="coerce")
                - electric_detail_sum
            ).clip(lower=0.0)
            if (
                base["__other_hvac_electricity_kwh__"]
                .dropna()
                .abs()
                .gt(1e-9)
                .any()
            ):
                meter_series.append(
                    (
                        "__other_hvac_electricity_kwh__",
                        "Other / Unmetered HVAC Electricity",
                    )
                )

    if not meter_series:
        return go.Figure()

    if mode == "hour":
        # Per hora fem mitjana entre dies: és millor per veure patrons típics que no pas
        # sumar tot l'any en una sola barra enorme.
        for kwh_column, label in meter_series:
            add_energy_bar_trace(
                fig,
                base,
                kwh_column,
                label=label,
                color=colors.get(label, STUDIO_CHART_COLORS["neutral"]),
                mode=mode,
                season=season,
            )
        fig.update_xaxes(title="Hour of Day", tickmode="linear", dtick=1)
        y_axis_title = energy_axis_title(mode)

    elif mode == "day":
        # En vista diària interessa més la proporció entre meters que la sèrie completa,
        # per això fem un donut amb el consum diari mitjà.
        df = filter_by_season(base, season)
        if df.empty:
            return go.Figure()
        labels = []
        values = []
        marker_colors = []
        for kwh_column, label in meter_series:
            _, y_vals, _ = _canon_day_total_series(df, kwh_column)
            daily_mean = pd.to_numeric(pd.Series(y_vals), errors="coerce").dropna().mean()
            if pd.isna(daily_mean) or daily_mean <= 1e-12:
                continue
            labels.append(label)
            values.append(float(daily_mean))
            marker_colors.append(colors.get(label, STUDIO_CHART_COLORS["neutral"]))

        if not values:
            return go.Figure()

        fig.add_trace(
            go.Pie(
                labels=labels,
                values=values,
                hole=0.42,
                marker=dict(colors=marker_colors, line=dict(color="#ffffff", width=2)),
                textinfo="percent+label",
                hovertemplate="%{label}<br>Mean daily consumption %{value:.2f} kWh<extra></extra>",
            )
        )
        fig.update_layout(
            title=f"HVAC Meter Breakdown - Mean Daily Share ({season})",
            showlegend=True,
        )
        return _apply_figure_theme(fig)

    elif mode == "month":
        # En mesos agrupem per any i mes abans de fer la mitjana, així dos anys simulats
        # no queden barrejats com si fossin un únic mes gegant.
        for kwh_column, label in meter_series:
            add_energy_bar_trace(
                fig,
                base,
                kwh_column,
                label=label,
                color=colors.get(label, STUDIO_CHART_COLORS["neutral"]),
                mode=mode,
                season=season,
            )
        fig.update_xaxes(title="Month", tickmode="linear", dtick=1)
        y_axis_title = energy_axis_title(mode)

    elif mode == "raw":
        df = filter_by_season(base, season)
        if df.empty:
            return go.Figure()
        for kwh_column, label in meter_series:
            x_vals, y_vals, hover = _raw_axis_data(df, kwh_column)
            color = colors.get(label, STUDIO_CHART_COLORS["neutral"])
            fig.add_trace(
                go.Scatter(
                    x=x_vals,
                    y=y_vals,
                    mode="lines",
                    name=label,
                    line=dict(color=color, width=1.8),
                    fillcolor=_alpha_color(color, 0.42),
                    stackgroup="hvac_meter_breakdown",
                    hovertext=hover,
                    hovertemplate="%{hovertext}<br>%{y:.4f} kWh<extra></extra>",
                )
            )
        fig.update_xaxes(title="Time (step)")
        y_axis_title = "Consumption per Timestep (kWh)"
    else:
        y_axis_title = "Consumption (kWh)"

    fig.update_yaxes(title=y_axis_title)
    fig.update_layout(
        title=f"HVAC Meter Breakdown - {mode.capitalize()} ({season})",
        barmode="group",
        bargap=0.18,
        bargroupgap=0.08,
    )
    return _apply_figure_theme(fig)
\end{lstlisting}

\Needspace{8\baselineskip}

\subsection{\texorpdfstring{\texttt{backend/grafics/kpis.py}}{backend/grafics/kpis.py}}

\begin{lstlisting}[style=studioPython,caption={backend/grafics/kpis.py},label={lst:backend-grafics-kpis-py}]
"""Càlculs de KPI per als resums de resultats de BEMS-RL Studio.

Aquest mòdul deriva mètriques operatives compactes a partir de les observacions
normalitzades: desviacions de confort, consum HVAC, cost energètic i resums de
temperatura de zona.
"""

import re

import pandas as pd

from backend.grafics.comfort_scope import comfort_scope_label, derive_occupied_mask
from backend.grafics.observation_columns import (
    HVAC_POWER_COLUMN,
    normalize_observation_columns,
)


def _normalize_name(s: str) -> str:
    """Normalitza els noms de zona perquè coincideixin amb les columnes de manera més tolerant."""
    s = (s or "").lower()
    s = s.replace(" ", "_").replace("-", "_")
    return re.sub(r"[^a-z0-9_]+", "", s)


def _pick_zone_temperature_series(
    obs: pd.DataFrame,
    selected_zone: str | None,
) -> pd.Series | None:
    """Retorna la sèrie de temperatura interior utilitzada per les targetes KPI."""
    zone_temp_cols = [c for c in obs.columns if c.endswith("_air_temperature")]
    if zone_temp_cols:
        if not selected_zone or selected_zone == "All":
            return obs[zone_temp_cols].mean(axis=1)

        target = _normalize_name(selected_zone)
        cand_map = {
            _normalize_name(c[:-len("_air_temperature")]): c for c in zone_temp_cols
        }
        col = cand_map.get(target)
        if col is not None:
            return pd.to_numeric(obs[col], errors="coerce")
        for key, col in cand_map.items():
            if target in key or key in target:
                return pd.to_numeric(obs[col], errors="coerce")
        return obs[zone_temp_cols].mean(axis=1)

    if "air_temperature" in obs.columns:
        return pd.to_numeric(obs["air_temperature"], errors="coerce")

    return None


def _temperature_violation_frame(obs: pd.DataFrame, comfort_config=None) -> pd.DataFrame | None:
    """Retorna les observacions amb una infracció de confort alineada per registre."""
    if "temp_violation" in obs.columns:
        result = obs.copy()
        result["temp_violation"] = pd.to_numeric(result["temp_violation"], errors="coerce")
        return result

    try:
        from backend.grafics.comfort import _ensure_temp_violation_column

        with_violation = _ensure_temp_violation_column(obs, comfort_config)
    except Exception:
        return None

    if "temp_violation" not in with_violation.columns:
        return None
    with_violation = with_violation.copy()
    with_violation["temp_violation"] = pd.to_numeric(
        with_violation["temp_violation"],
        errors="coerce",
    )
    return with_violation


def compute_kpis(
    obs: pd.DataFrame,
    cost_hourly: pd.DataFrame | None,
    selected_zone: str | None = None,
    comfort_scope: str = "all",
    comfort_config=None,
):
    """Calcula les targetes KPI mostrades al panell i exportades a PDF.

    Els KPI de temperatura respecten la zona seleccionada. HVAC i els KPI de costos segueixen sent globals.
    La desviació de confort es mostra amb una etiqueta d'abast explícita per a hores ocupades
    les recompenses no es comparen amb una mètrica de totes les hores per accident.
    """
    obs = normalize_observation_columns(obs)
    kpis: dict[str, float | None] = {}

    tin_series = _pick_zone_temperature_series(obs, selected_zone)
    if tin_series is not None:
        # Si l'usuari ha triat una zona, els KPI de temperatura han de parlar només
        # d'aquella zona, no de la mitjana global de l'edifici.
        kpis["Avg Temp Interior (C)"] = round(tin_series.mean(), 2)
        kpis["Max Temp Interior (C)"] = round(tin_series.max(), 2)
        kpis["Min Temp Interior (C)"] = round(tin_series.min(), 2)
    else:
        if "air_temperature" not in obs.columns:
            zone_temp_cols = [c for c in obs.columns if c.endswith("_air_temperature")]
            tmp = obs[zone_temp_cols].mean(axis=1) if zone_temp_cols else None
        else:
            tmp = obs["air_temperature"]
        if tmp is not None:
            tmp = pd.to_numeric(tmp, errors="coerce")
            kpis["Avg Temp Interior (C)"] = round(tmp.mean(), 2)
            kpis["Max Temp Interior (C)"] = round(tmp.max(), 2)
            kpis["Min Temp Interior (C)"] = round(tmp.min(), 2)
        else:
            kpis["Avg Temp Interior (C)"] = None
            kpis["Max Temp Interior (C)"] = None
            kpis["Min Temp Interior (C)"] = None

    if HVAC_POWER_COLUMN in obs.columns:
        hvac = pd.to_numeric(obs[HVAC_POWER_COLUMN], errors="coerce")
        kpis["Avg HVAC Power (kW)"] = round(hvac.mean() / 1000, 2)
    else:
        kpis["Avg HVAC Power (kW)"] = None

    if cost_hourly is not None and "energy_cost_eur" in cost_hourly.columns:
        energy_cost = pd.to_numeric(cost_hourly["energy_cost_eur"], errors="coerce")
        kpis["Avg Energy Cost (EUR/h)"] = round(energy_cost.mean(), 2)
    else:
        kpis["Avg Energy Cost (EUR/h)"] = None

    violation_frame = _temperature_violation_frame(obs, comfort_config)
    primary_label = comfort_scope_label(comfort_scope)

    if violation_frame is None:
        kpis[f"Avg Temp Violation ({primary_label}, C)"] = None
        return kpis

    violation = violation_frame["temp_violation"]
    occupied_mask = derive_occupied_mask(violation_frame)

    # En rewards amb ocupació, comparar totes les hores amb hores ocupades pot portar a
    # conclusions falses. Per això calculem l'abast principal i deixem l'altre com a context.
    scoped_values = violation
    if comfort_scope == "occupied" and occupied_mask is not None:
        scoped_values = violation.loc[occupied_mask]
    scoped_values = scoped_values.dropna()
    scoped_mean = scoped_values.mean() if not scoped_values.empty else None
    kpis[f"Avg Temp Violation ({primary_label}, C)"] = (
        round(scoped_mean, 2) if scoped_mean is not None else None
    )

    if occupied_mask is not None:
        all_values = violation.dropna()
        occupied_values = violation.loc[occupied_mask].dropna()
        all_mean = all_values.mean() if not all_values.empty else None
        occupied_mean = occupied_values.mean() if not occupied_values.empty else None
        kpis["Occupied Samples (%)"] = round(float(occupied_mask.mean()) * 100.0, 2)

        if comfort_scope == "occupied":
            kpis["Avg Temp Violation (All Hours, C)"] = (
                round(all_mean, 2) if all_mean is not None else None
            )
        else:
            kpis["Avg Temp Violation (Occupied Hours, C)"] = (
                round(occupied_mean, 2) if occupied_mean is not None else None
            )

    return kpis
\end{lstlisting}

\Needspace{8\baselineskip}

\subsection{\texorpdfstring{\texttt{backend/grafics/metrics.py}}{backend/grafics/metrics.py}}

\begin{lstlisting}[style=studioPython,caption={backend/grafics/metrics.py},label={lst:backend-grafics-metrics-py}]
"""Agrega les dades d'observació i progrés en mètriques preparades per a gràfics.

Les funcions aquí converteixen la sortida en brut Sinergym en mensual, per hora i
DataFrames orientats a la comoditat consumits pel panell de resultats, Plotly builders
i informes exportats.
"""

import numpy as np
import pandas as pd
from typing import Optional

from backend.grafics.observation_columns import (
    ENERGY_COST_COLUMN,
    HVAC_POWER_COLUMN,
    convert_hvac_source_to_kwh,
    find_hvac_consumption_source,
    normalize_observation_columns,
)
from backend.grafics.time_utils import _infer_timestep_hours, sanitize_observation_time_columns


def _price_to_eur_per_kwh(values: pd.Series) -> pd.Series:
    """Normalitza els valors dels preus de l'electricitat a EUR/kWh quan semblen ser EUR/MWh."""
    price = pd.to_numeric(values, errors="coerce")
    median_abs = price.abs().dropna().median()
    if pd.notna(median_abs) and median_abs > 5:
        return price / 1000.0
    return price


def compute_metrics(progress: pd.DataFrame, obs: pd.DataFrame, yaml_cfg: Optional[dict] = None) -> dict:
    """
    Calcula tots els DataFrames agregats per a figures,
    validant contra les dades disponibles i opcionalment el fitxer YAML.
    """
    obs = sanitize_observation_time_columns(normalize_observation_columns(obs.copy()))

    # Els CSV poden venir amb timestamp/datetime o amb columnes ja separades. Normalitzem
    # month/hour aquí perquè totes les figures de resultats puguin reutilitzar el mateix output.
    if 'month' not in obs.columns:
        if 'timestamp' in obs.columns:
            obs['timestamp'] = pd.to_datetime(obs['timestamp'])
            obs['month'] = obs['timestamp'].dt.month
        elif 'datetime' in obs.columns:
            obs['datetime'] = pd.to_datetime(obs['datetime'])
            obs['month'] = obs['datetime'].dt.month
        else:
            raise ValueError("La columna 'month' no existeix ni es pot derivar.")
    if 'hour' not in obs.columns:
        if 'timestamp' in obs.columns:
            obs['timestamp'] = pd.to_datetime(obs['timestamp'])
            obs['hour'] = obs['timestamp'].dt.hour
        elif 'datetime' in obs.columns:
            obs['datetime'] = pd.to_datetime(obs['datetime'])
            obs['hour'] = obs['datetime'].dt.hour
        else:
            raise ValueError("La columna 'hour' no existeix ni es pot derivar.")

    available = set(obs.columns)

    # Si el CSV és multizona i no porta air_temperature global, fem una mitjana simple
    # només com a fallback per a gràfics globals.
    if 'air_temperature' not in available:
        zone_temp_cols = [c for c in available if c.endswith('_air_temperature')]
        if zone_temp_cols:
            obs['air_temperature'] = obs[zone_temp_cols].mean(axis=1)
            available.add('air_temperature')

    if HVAC_POWER_COLUMN not in available:
        for cand in ('hvac_power', 'HVAC_power', 'hvac_load'):
            if cand in available:
                obs[HVAC_POWER_COLUMN] = obs[cand]
                available.add(HVAC_POWER_COLUMN)
                break

    if 'temp_violation' not in available and 'air_temperature' in available:
        from backend.grafics.comfort import _ensure_temp_violation_column

        obs = _ensure_temp_violation_column(obs, yaml_cfg)
        available.add('temp_violation')

    bins = 12
    if len(progress) > 0 and 'mean_reward' in progress.columns:
        # progress.csv no sempre té calendari real; repartim els punts en 12 blocs per
        # tenir una lectura mensual aproximada sense inventar dates.
        progress = progress.copy()
        progress['mean_reward'] = pd.to_numeric(progress['mean_reward'], errors='coerce')
        progress['month_bin'] = (np.arange(len(progress)) * bins // len(progress)) + 1
        monthly_reward = (
            progress.groupby('month_bin')['mean_reward']
                    .mean()
                    .reset_index(name='mean_reward')
        )
        monthly_reward['month'] = monthly_reward['month_bin'].astype(int)
    else:
        monthly_reward = pd.DataFrame({'month': range(1, bins + 1), 'mean_reward': [0] * bins})

    hourly_all = obs.groupby('hour').mean(numeric_only=True).reset_index()
    winter = obs[obs['month'].isin([12, 1, 2])].groupby('hour').mean(numeric_only=True).reset_index()
    summer = obs[obs['month'].isin([6, 7, 8])].groupby('hour').mean(numeric_only=True).reset_index()
    monthly_obs = (
        obs.groupby('month').mean(numeric_only=True)
           .reindex(index=range(1, 13), fill_value=0)
           .reset_index()
    )

    step_hours = _infer_timestep_hours(obs)
    hvac_source = find_hvac_consumption_source(obs.columns)
    if hvac_source is not None:
        # Aquí unifiquem W, J o meters ja en kWh. La resta del dashboard treballa sempre
        # amb aquesta columna normalitzada.
        hvac_col, hvac_unit = hvac_source
        obs["hvac_consumption_kwh"] = convert_hvac_source_to_kwh(
            obs[hvac_col],
            hvac_unit,
            step_hours,
        )
        available.add("hvac_consumption_kwh")

    cost_hourly = cost_monthly = None
    if ENERGY_COST_COLUMN in available and 'hvac_consumption_kwh' in available:
        obs['energy_price_kwh'] = _price_to_eur_per_kwh(obs[ENERGY_COST_COLUMN])
        obs['energy_cost_eur'] = obs['hvac_consumption_kwh'] * obs['energy_price_kwh']
        cost_hourly = obs.groupby('hour')['energy_cost_eur'].sum().reset_index()
        cost_monthly = (
            obs.groupby('month')['energy_cost_eur']
               .sum()
               .reindex(index=range(1, 13), fill_value=0)
               .reset_index()
        )

    violation_hourly = eps_winter = eps_summer = season_line = None
    if 'temp_violation' in available:
        # Separem hivern/estiu perquè les consignes de confort poden ser diferents i la
        # mitjana anual amagaria bastant el problema.
        violation_hourly = obs.groupby('hour')['temp_violation'].mean().reset_index(name='mean_violation')
        eps_winter = obs[obs['month'].isin([12, 1, 2])].groupby('hour')['temp_violation'].mean().reset_index(name='winter_violation')
        eps_summer = obs[obs['month'].isin([6, 7, 8])].groupby('hour')['temp_violation'].mean().reset_index(name='summer_violation')
        season_line = pd.merge(eps_winter, eps_summer, on='hour', how='outer').fillna(0)

    pivot_violation = None
    if 'temp_violation' in available:
        pivot_violation = (
            obs.groupby(['month', 'hour'])['temp_violation']
               .mean()
               .unstack(fill_value=0)
               .reindex(index=range(1, 13), fill_value=0)
        )

    pivot_consumption = None
    if 'hvac_consumption_kwh' in available:
        pivot_consumption = (
            obs.groupby(['month', 'hour'])['hvac_consumption_kwh']
               .sum()
               .unstack(fill_value=0)
               .reindex(index=range(1, 13), fill_value=0)
        )

    return {
        'monthly_reward': monthly_reward,
        'hourly_all': hourly_all,
        'winter': winter,
        'summer': summer,
        'monthly_obs': monthly_obs,
        'cost_hourly': cost_hourly,
        'cost_monthly': cost_monthly,
        'violation_hourly': violation_hourly,
        'eps_winter': eps_winter,
        'eps_summer': eps_summer,
        'season_line': season_line,
        'pivot_violation': pivot_violation,
        'pivot_consumption': pivot_consumption
    }
\end{lstlisting}

\Needspace{8\baselineskip}

\subsection{\texorpdfstring{\texttt{backend/grafics/observation\_columns.py}}{backend/grafics/observation\_columns.py}}

\begin{lstlisting}[style=studioPython,caption={backend/grafics/observation\_columns.py},label={lst:backend-grafics-observation-columns-py}]
"""Utilitats de normalització de columnes per a fitxers d'observació Sinergym.

Els entorns Sinergym i les configuracions personalitzades poden emetre comptadors equivalents
sota diferents noms o unitats. Aquest mòdul centralitza els àlies, HVAC energia
conversió i creació de columnes normalitzades perquè tots els gràfics llegeixin igual
camps canònics.
"""

from __future__ import annotations

from typing import Iterable

import pandas as pd

HVAC_POWER_COLUMN = "HVAC_electricity_demand_rate"
ENERGY_COST_COLUMN = "energy_cost"
JOULES_PER_KWH = 3_600_000.0

HVAC_ENERGY_METER_COLUMNS = (
    "total_electricity_hvac",
    "total_electricity_HVAC",
    "Electricity:HVAC",
    "electricity_hvac",
)

DETAILED_HVAC_METER_COLUMNS = {
    "natural_gas_hvac": (
        "NaturalGas:HVAC",
        "natural_gas_hvac",
        "naturalgas_hvac",
    ),
    "heating_electricity": (
        "Heating:Electricity",
        "heating_electricity",
    ),
    "cooling_electricity": (
        "Cooling:Electricity",
        "cooling_electricity",
    ),
    "fans_electricity": (
        "Fans:Electricity",
        "fans_electricity",
    ),
    "pumps_electricity": (
        "Pumps:Electricity",
        "pumps_electricity",
    ),
    "heat_rejection_electricity": (
        "HeatRejection:Electricity",
        "heat_rejection_electricity",
    ),
    "humidifier_electricity": (
        "Humidifier:Electricity",
        "Humidification:Electricity",
        "humidifier_electricity",
        "humidification_electricity",
    ),
    "heat_recovery_electricity": (
        "HeatRecovery:Electricity",
        "heat_recovery_electricity",
    ),
}

DETAILED_HVAC_METER_LABELS = {
    "natural_gas_hvac": "Natural Gas HVAC",
    "heating_electricity": "Heating Electricity",
    "cooling_electricity": "Cooling Electricity",
    "fans_electricity": "Fans Electricity",
    "pumps_electricity": "Pumps Electricity",
    "heat_rejection_electricity": "Heat Rejection Electricity",
    "humidifier_electricity": "Humidifier Electricity",
    "heat_recovery_electricity": "Heat Recovery Electricity",
}

ELECTRIC_HVAC_METER_ALIASES = (
    "heating_electricity",
    "cooling_electricity",
    "fans_electricity",
    "pumps_electricity",
    "heat_rejection_electricity",
    "humidifier_electricity",
    "heat_recovery_electricity",
)

DETAILED_HVAC_METER_KWH_COLUMNS = {
    alias: f"{alias}_kwh"
    for alias in DETAILED_HVAC_METER_COLUMNS
}

FALLBACK_ELECTRICITY_METER_COLUMNS = (
    "total_electricity_facility",
    "total_electricity_building",
)

_COLUMN_ALIASES = {
    HVAC_POWER_COLUMN: (
        HVAC_POWER_COLUMN,
        "HVAC_elèctricity_demand_rate",
        "hvac_power",
        "HVAC_power",
        "hvac_load",
    ),
    ENERGY_COST_COLUMN: (
        ENERGY_COST_COLUMN,
        "electricity_cost",
        "elèctricity_cost",
    ),
}

_COLUMN_ALIASES.update(
    {
        canonical: (canonical, *aliases)
        for canonical, aliases in DETAILED_HVAC_METER_COLUMNS.items()
    }
)


def find_first_available_column(
    columns: Iterable[str], candidates: Iterable[str]
) -> str | None:
    """Retorna la primera columna disponible dins la llista de candidats."""
    available = set(columns)
    for candidate in candidates:
        if candidate in available:
            return candidate
    return None


def normalize_observation_columns(obs: pd.DataFrame) -> pd.DataFrame:
    """Normalitza les columnes d'observació."""
    normalized = obs.copy()
    for canonical, aliases in _COLUMN_ALIASES.items():
        if canonical in normalized.columns:
            continue
        alias = find_first_available_column(normalized.columns, aliases)
        if alias and alias != canonical:
            normalized[canonical] = normalized[alias]
    return normalized


def find_hvac_consumption_source(columns: Iterable[str]) -> tuple[str, str] | None:
    """Retorna la millor font de consum HVAC i el seu tipus d'unitat.

    EnergyPlus els comptadors informen l'energia per pas de temps en Joules, mentre que el
    La variable de sortida de la taxa de demanda informa de la potència en watts.
    """
    available = set(columns)
    lower_lookup = {column.lower(): column for column in columns}
    for column in HVAC_ENERGY_METER_COLUMNS:
        if column in available:
            return column, "joule"
        actual_column = lower_lookup.get(column.lower())
        if actual_column is not None:
            return actual_column, "joule"
    if HVAC_POWER_COLUMN in available:
        return HVAC_POWER_COLUMN, "watt"
    actual_power_column = lower_lookup.get(HVAC_POWER_COLUMN.lower())
    if actual_power_column is not None:
        return actual_power_column, "watt"
    for column in FALLBACK_ELECTRICITY_METER_COLUMNS:
        if column in available:
            return column, "joule"
        actual_column = lower_lookup.get(column.lower())
        if actual_column is not None:
            return actual_column, "joule"
    return None


def convert_hvac_source_to_kwh(
    values: pd.Series,
    unit_kind: str,
    timestep_hours: float | None = None,
) -> pd.Series:
    """Converteix la font de consum HVAC a kWh."""
    numeric = pd.to_numeric(values, errors="coerce")
    if unit_kind == "joule":
        return numeric / JOULES_PER_KWH
    if unit_kind == "watt":
        hours = 0.0 if timestep_hours is None else float(timestep_hours)
        return (numeric / 1000.0) * hours
    return numeric


def add_meter_kwh_columns(obs: pd.DataFrame) -> pd.DataFrame:
    """Afegeix columnes derivades de kWh per a les sortides conegudes del comptador EnergyPlus en Joules."""
    enriched = obs.copy()
    for meter_alias, kwh_column in DETAILED_HVAC_METER_KWH_COLUMNS.items():
        if meter_alias not in enriched.columns or kwh_column in enriched.columns:
            continue
        enriched[kwh_column] = convert_hvac_source_to_kwh(
            enriched[meter_alias],
            "joule",
        )
    return enriched
\end{lstlisting}

\Needspace{8\baselineskip}

\subsection{\texorpdfstring{\texttt{backend/grafics/plot\_helpers.py}}{backend/grafics/plot\_helpers.py}}

\begin{lstlisting}[style=studioPython,caption={backend/grafics/plot\_helpers.py},label={lst:backend-grafics-plot-helpers-py}]
"""Helpers Plotly comuns per als gràfics agregats del dashboard."""

from __future__ import annotations

import numpy as np
import pandas as pd
import plotly.graph_objects as go

from backend.grafics.time_utils import (
    _canon_day_total_series,
    _hour_series,
    _infer_timestep_hours,
    _month_series,
    _raw_axis_data,
    _seasonal_marker_colors,
    _series_for_mode,
    filter_by_season,
)


def build_mode_scatter_trace(
    base: pd.DataFrame,
    column: str | None,
    mode: str,
    season: str,
    *,
    name: str,
    color: str,
    units: str = "",
    sign: float = 1.0,
    line_width: float = 2.4,
    marker_size: int = 6,
    trace_mode: str | None = None,
    hovertemplate: str | None = None,
    plain_hovertemplate: str | None = None,
    line_shape: str | None = None,
    fill: str | None = None,
    fillcolor: str | None = None,
    meta: str | None = None,
) -> go.Scatter | None:
    """Crea una traça scatter usant la mateixa agregació temporal que la resta de gràfics."""
    if column is None:
        return None
    data = _series_for_mode(base, column, mode, season, sign=sign)
    if data is None:
        return None

    x_values, y_values, hover, inferred_mode = data
    selected_mode = trace_mode or inferred_mode
    marker_color = (
        _seasonal_marker_colors(x_values, season, active_color=color)
        if mode == "month"
        else color
    )
    line = dict(color=color, width=line_width)
    if line_shape is not None:
        line["shape"] = line_shape

    trace = go.Scatter(
        x=x_values,
        y=y_values,
        mode=selected_mode,
        name=name,
        line=line,
        marker=dict(color=marker_color, size=marker_size),
        fill=fill,
        fillcolor=fillcolor,
        meta=meta,
    )
    if hover is not None:
        trace.hovertext = hover
        trace.hovertemplate = hovertemplate or _default_hovertemplate(units)
    elif plain_hovertemplate is not None or hovertemplate is not None:
        trace.hovertemplate = plain_hovertemplate or hovertemplate
    return trace


def add_mode_scatter_trace(
    fig: go.Figure,
    base: pd.DataFrame,
    column: str | None,
    mode: str,
    season: str,
    **kwargs,
) -> go.Scatter | None:
    """Afegeix una traça scatter agregada a una figura si la columna existeix."""
    trace = build_mode_scatter_trace(base, column, mode, season, **kwargs)
    if trace is not None:
        fig.add_trace(trace)
    return trace


def energy_series_for_mode(
    base: pd.DataFrame,
    column: str,
    mode: str,
    season: str,
    *,
    raw_hourly: bool = False,
) -> tuple[list, list, list | None] | None:
    """Retorna sèries d'energia kWh amb sumes coherents per mode temporal."""
    if column not in base.columns:
        return None

    if mode == "hour":
        df = filter_by_season(base, season)
        if df.empty:
            return None
        day_keys = pd.Index(df.index.floor("D"), name="day")
        hour_keys = pd.Index(_hour_series(df), name="hour")
        slots = df.groupby([day_keys, hour_keys])[column].sum()
        grouped = slots.groupby(level="hour").mean()
        x_values = list(range(24))
        return x_values, [grouped.get(hour, float("nan")) for hour in x_values], None

    if mode == "day":
        df = filter_by_season(base, season)
        if df.empty:
            return None
        return _canon_day_total_series(df, column)

    if mode == "month":
        df = base if season == "All" else filter_by_season(base, season)
        if df.empty:
            return None
        if isinstance(df.index, pd.DatetimeIndex):
            monthly = df.groupby([df.index.year, df.index.month])[column].sum()
            grouped = monthly.groupby(level=1).mean()
        else:
            grouped = df.groupby(_month_series(df))[column].sum()
        x_values = list(range(1, 13))
        return x_values, [grouped.get(month, float("nan")) for month in x_values], None

    if mode == "raw":
        df = filter_by_season(base, season)
        if df.empty:
            return None
        return raw_hourly_total_series(df, column) if raw_hourly else _raw_axis_data(df, column)

    return None


def add_energy_bar_trace(
    fig: go.Figure,
    base: pd.DataFrame,
    column: str,
    *,
    label: str,
    color: str,
    mode: str,
    season: str,
    hover_units: str = "kWh",
    raw_hourly: bool = False,
) -> go.Bar | None:
    """Afegeix barres d'energia agregades amb hover homogeni."""
    data = energy_series_for_mode(base, column, mode, season, raw_hourly=raw_hourly)
    if data is None:
        return None

    x_values, y_values, hover = data
    marker_color = (
        _seasonal_marker_colors(x_values, season, active_color=color)
        if mode == "month"
        else color
    )
    trace = go.Bar(
        x=x_values,
        y=y_values,
        name=label,
        marker=dict(color=marker_color, line=dict(color=color, width=0.5)),
    )
    if hover is not None:
        trace.hovertext = hover
        decimals = 4 if mode == "raw" and not raw_hourly else 2
        trace.hovertemplate = (
            f"%{{hovertext}}<br>%{{y:.{decimals}f}} {hover_units}<extra></extra>"
        )
    fig.add_trace(trace)
    return trace


def raw_hourly_total_series(df: pd.DataFrame, column: str) -> tuple[list, list, list]:
    """Agrupa una sèrie raw en totals horaris per evitar gràfics de barres massa densos."""
    values = pd.to_numeric(df[column], errors="coerce")

    if isinstance(df.index, pd.DatetimeIndex):
        hourly = values.resample("h").sum(min_count=1).dropna()
        x_values = list(range(1, len(hourly) + 1))
        hover = [stamp.strftime("%m-%d %H:00") for stamp in hourly.index]
        return x_values, hourly.to_list(), hover

    timestep_hours = _infer_timestep_hours(df)
    steps_per_hour = max(1, int(round(1.0 / timestep_hours))) if timestep_hours > 0 else 4
    groups = np.arange(len(values)) // steps_per_hour
    hourly = values.groupby(groups).sum(min_count=1).dropna()
    x_values = (hourly.index.to_numpy(dtype=int) + 1).tolist()
    hover = [f"Hour {hour}" for hour in x_values]
    return x_values, hourly.to_list(), hover


def energy_axis_title(mode: str, *, raw_title: str = "Consumption per Timestep (kWh)") -> str:
    """Retorna el títol Y habitual per a gràfics de consum energètic."""
    return {
        "hour": "Mean Hourly Consumption (kWh)",
        "day": "Daily Consumption (kWh)",
        "month": "Monthly Consumption (kWh)",
        "raw": raw_title,
    }.get(mode, "Consumption (kWh)")


def _default_hovertemplate(units: str) -> str:
    """Construeix el tooltip Plotly per defecte amb unitats opcionals."""
    if not units:
        return "%{hovertext}<br>%{y:.2f}<extra></extra>"
    return f"%{{hovertext}}<br>%{{y:.2f}} {units}<extra></extra>"
\end{lstlisting}

\Needspace{8\baselineskip}

\subsection{\texorpdfstring{\texttt{backend/grafics/style.py}}{backend/grafics/style.py}}

\begin{lstlisting}[style=studioPython,caption={backend/grafics/style.py},label={lst:backend-grafics-style-py}]
"""Estil Plotly compartit per a les figures del panell BEMS-RL Studio."""

from __future__ import annotations

import numpy as np
import plotly.graph_objects as go


# Paleta i format comú de les figures.
STUDIO_CHART_COLORS = {
    "heating": "#c81e1e",
    "heating_ref": "#ef4444",
    "cooling": "#1d4ed8",
    "cooling_ref": "#38bdf8",
    "consumption": "#f59e0b",
    "consumption_muted": "#fbbf24",
    "consumption_edge": "#b45309",
    "violation": "#dc2626",
    "violation_muted": "#f87171",
    "violation_edge": "#991b1b",
    "indoor": "#0f766e",
    "outdoor": "#64748b",
    "comfort_ok": "#16a34a",
    "comfort_cold": "#2563eb",
    "comfort_hot": "#dc2626",
    "battery": "#7c3aed",
    "battery_charge": "#15803d",
    "battery_discharge": "#ea580c",
    "grid_purchase": "#ca8a04",
    "price": "#0891b2",
    "reward": "#6d28d9",
    "neutral": "#475569",
    "neutral_soft": "#cbd5e1",
    "grid": "#d7deeb",
    "paper": "#ffffff",
    "plot": "#f8fbff",
    "text": "#17233c",
}

STUDIO_HEATMAP_SCALES = {
    "consumption": "YlOrBr",
    "violation": "Reds",
    "temperature": "RdYlBu_r",
}


def pick_semantic_trace_color(
    trace_name: str | None,
    *,
    reference: bool = False,
    fallback: str | None = None,
) -> str:
    """Retorna un color semàntic segons el nom de la traça."""
    label = (trace_name or "").lower()

    # Les figures arriben de mòduls diferents i no sempre comparteixen palette. Fem servir
    # paraules del nom de la traça per mantenir colors coherents sense acoblar els gràfics.
    if "heating" in label or "htg" in label:
        return STUDIO_CHART_COLORS["heating_ref" if reference else "heating"]
    if "cooling" in label or "clg" in label:
        return STUDIO_CHART_COLORS["cooling_ref" if reference else "cooling"]
    if "viol" in label:
        return STUDIO_CHART_COLORS["violation_muted" if reference else "violation"]
    if "massa fred" in label or "cold" in label:
        return STUDIO_CHART_COLORS["comfort_cold"]
    if "massa calor" in label or "hot" in label:
        return STUDIO_CHART_COLORS["comfort_hot"]
    if "confort" in label or "comfort" in label:
        return STUDIO_CHART_COLORS["comfort_ok"]
    if "hvac" in label or "consum" in label or "demand rate" in label or "mean power" in label:
        return STUDIO_CHART_COLORS["consumption_muted" if reference else "consumption"]
    if "price" in label or "preu" in label or "eur/kwh" in label:
        return STUDIO_CHART_COLORS["price"]
    if "reward" in label:
        return STUDIO_CHART_COLORS["reward"]
    if "outdoor" in label:
        return STUDIO_CHART_COLORS["outdoor"]
    if "radiant operation" in label or "radiant availability" in label:
        return STUDIO_CHART_COLORS["battery"]
    if "radiant inlet" in label:
        return STUDIO_CHART_COLORS["price"]
    if "radiant outlet" in label or "radiant water" in label:
        return STUDIO_CHART_COLORS["battery_discharge"]
    if "discharge" in label:
        return STUDIO_CHART_COLORS["battery_discharge"]
    if "charge" in label:
        return STUDIO_CHART_COLORS["battery_charge"]
    if "battery" in label or "stored energy" in label or "soc" in label or "bateria" in label:
        return STUDIO_CHART_COLORS["battery_charge"]
    if "xarxa" in label or "compra" in label or "grid purchase" in label:
        return STUDIO_CHART_COLORS["grid_purchase"]
    if "indoor" in label:
        return STUDIO_CHART_COLORS["indoor"]

    return fallback or STUDIO_CHART_COLORS["neutral"]


def _alpha_color(hex_color: str, alpha: float = 0.30) -> str:
    """Converteix un color hex a rgba amb transparència."""
    h = hex_color.lstrip("#")
    r, g, b = int(h[0:2], 16), int(h[2:4], 16), int(h[4:6], 16)
    return f"rgba({r},{g},{b},{alpha})"

def _apply_figure_theme(fig: go.Figure, *, heatmap: bool = False) -> go.Figure:
    """Aplica un estil més contrastat i homogeni a les figures."""
    fig.update_layout(
        template="plotly_white",
        paper_bgcolor=STUDIO_CHART_COLORS["paper"],
        plot_bgcolor=STUDIO_CHART_COLORS["plot"],
        font=dict(color=STUDIO_CHART_COLORS["text"], size=14),
        title=dict(
            x=0.02,
            xanchor="left",
            font=dict(size=20, color=STUDIO_CHART_COLORS["text"]),
        ),
        legend=dict(
            bgcolor="rgba(255,255,255,0.9)",
            bordercolor=STUDIO_CHART_COLORS["grid"],
            borderwidth=1,
        ),
        hoverlabel=dict(
            bgcolor="#ffffff",
            font_color=STUDIO_CHART_COLORS["text"],
        ),
        margin=dict(l=60, r=28, t=72, b=56),
    )
    fig.update_xaxes(
        showgrid=not heatmap,
        gridcolor=STUDIO_CHART_COLORS["grid"],
        linecolor=STUDIO_CHART_COLORS["grid"],
        ticks="outside",
        zeroline=False,
    )
    fig.update_yaxes(
        showgrid=not heatmap,
        gridcolor=STUDIO_CHART_COLORS["grid"],
        linecolor=STUDIO_CHART_COLORS["grid"],
        ticks="outside",
        zeroline=False,
    )
    return fig


def style_figure_semantics(fig: go.Figure) -> go.Figure:
    """Recolora traces segons la seva semàntica i aplica el tema comú."""
    heatmap_like = False

    for trace in getattr(fig, "data", []):
        trace_type = getattr(trace, "type", "")
        name = getattr(trace, "name", None) or ""
        base_name = name.replace(" (ref)", "")
        is_reference = "(ref)" in name.lower()
        if getattr(trace, "legendgroup", None) == "agent_actions":
            continue
        if getattr(trace, "meta", None) == "keep_color":
            continue
        # Algunes traces ja porten colors especials (per exemple heatmaps o accions).
        # Les altres es recoloren aquí perquè el dashboard sembli una sola peça.
        color = pick_semantic_trace_color(base_name, reference=is_reference)

        if trace_type in ("scatter", "scattergl"):
            line = getattr(trace, "line", None)
            width = max(float(getattr(line, "width", 0) or 0), 2.8 if is_reference else 2.6)
            marker = getattr(trace, "marker", None)
            marker_color = getattr(marker, "color", None)
            trace.line = dict(
                color=color,
                width=width,
            )
            if "markers" in str(getattr(trace, "mode", "")):
                trace.marker = dict(
                    color=marker_color if isinstance(marker_color, (list, tuple, np.ndarray)) and not is_reference else color,
                    size=max(int(getattr(marker, "size", 0) or 0), 6),
                )
            if is_reference:
                trace.opacity = max(float(getattr(trace, "opacity", 1.0) or 1.0), 0.92)

        elif trace_type == "bar":
            marker = getattr(trace, "marker", None)
            marker_color = getattr(marker, "color", None)
            edge_color = (
                STUDIO_CHART_COLORS["violation_edge"]
                if color in (STUDIO_CHART_COLORS["violation"], STUDIO_CHART_COLORS["violation_muted"])
                else STUDIO_CHART_COLORS["consumption_edge"]
            )
            trace.marker = dict(
                color=marker_color if isinstance(marker_color, (list, tuple, np.ndarray)) and not is_reference else color,
                line=dict(color=edge_color, width=0.8),
            )
            if is_reference:
                trace.opacity = min(float(getattr(trace, "opacity", 1.0) or 1.0), 0.58)

        elif trace_type == "heatmap":
            heatmap_like = True
            title_text = str(getattr(getattr(fig.layout, "title", None), "text", "") or "").lower()
            if "viol" in title_text:
                trace.colorscale = STUDIO_HEATMAP_SCALES["violation"]
            elif "consum" in title_text or "hvac" in title_text:
                trace.colorscale = STUDIO_HEATMAP_SCALES["consumption"]
            else:
                trace.colorscale = STUDIO_HEATMAP_SCALES["temperature"]

    return _apply_figure_theme(fig, heatmap=heatmap_like)
\end{lstlisting}

\Needspace{8\baselineskip}

\subsection{\texorpdfstring{\texttt{backend/grafics/thermal.py}}{backend/grafics/thermal.py}}

\begin{lstlisting}[style=studioPython,caption={backend/grafics/thermal.py},label={lst:backend-grafics-thermal-py}]
"""Dades d'entrada de confort tèrmic: temperatura, humitat i consignes."""

from __future__ import annotations

import pandas as pd
import plotly.graph_objects as go

from backend.grafics.column_utils import _ensure_air_temperature
from backend.grafics.plot_helpers import build_mode_scatter_trace
from backend.grafics.style import pick_semantic_trace_color
from backend.grafics.time_utils import (
    _apply_mode_xaxis,
    _canon_day_series,
    _ensure_datetime_index,
    _get_outdoor_temperature_series,
    _hour_series,
    _month_series,
    _raw_axis_data,
    filter_by_season,
)


def make_indoor_temperature_plot(obs: pd.DataFrame, mode: str, season: str) -> go.Figure:
    """Dibuixa la temperatura interior com a barres i l'exterior com una línia amb marcadors."""
    base = _ensure_datetime_index(_ensure_air_temperature(obs))
    out_series = _get_outdoor_temperature_series(base)

    _C_BAR = "rgba(79, 142, 247, 0.65)"   # blau EPW (igual que Humitat)
    _C_OUT = "#f26b5b"                     # salmó EPW (igual que Temp. seca)
    _C_OUT_EDGE = "#f26b5b"

    fig = go.Figure()
    if "air_temperature" not in base.columns:
        return go.Figure()

    if mode in ("hour", "month"):
        # Hora i mes són agregats; per això fem barres de mitjana interior i línia exterior
        # amb el mateix eix temporal.
        if mode == "hour":
            df = filter_by_season(base, season)
            grp = _hour_series(df)
            x_vals = list(range(24))
            x_title, tick_cfg = "Hour of Day", dict(tickmode="linear", dtick=1)
        else:
            df = base if season == "All" else filter_by_season(base, season)
            grp = _month_series(df)
            x_vals = list(range(1, 13))
            x_title, tick_cfg = "Month", dict(tickmode="linear", dtick=1)

        g_mean = df.groupby(grp)["air_temperature"].mean()
        y_in   = [g_mean.get(v, float("nan")) for v in x_vals]

        fig.add_trace(go.Bar(
            x=x_vals, y=y_in, name="Indoor Temp",
            marker_color=_C_BAR,
            meta="keep_color",
            hovertemplate="%{x}<br>Mitjana: %{y:.1f}°C<extra></extra>",
        ))
        if out_series is not None:
            g_out = df.groupby(grp)[out_series.name].mean()
            y_out = [g_out.get(v, float("nan")) for v in x_vals]
            fig.add_trace(go.Scatter(
                x=x_vals, y=y_out, mode="lines+markers", name="Outdoor Temp",
                line=dict(color=_C_OUT, width=2.5),
                marker=dict(color=_C_OUT, size=7, symbol="circle",
                            line=dict(color="white", width=1.5)),
                meta="keep_color",
                hovertemplate="%{x}<br>%{y:.1f}°C<extra></extra>",
            ))
        fig.update_xaxes(title=x_title, **tick_cfg)
        fig.update_layout(bargap=0.22)

    elif mode == "day":
        # La vista diària conserva la forma anual però redueix el soroll de cada timestep.
        df = filter_by_season(base, season)
        x_in, y_in, hover_in = _canon_day_series(df, "air_temperature")
        fig.add_trace(go.Scatter(
            x=x_in, y=y_in, mode="lines", name="Indoor Temp",
            line=dict(color="rgba(79, 142, 247, 0.9)", width=2),
            meta="keep_color",
            hovertext=hover_in,
            hovertemplate="%{hovertext}<br>%{y:.1f}°C<extra></extra>",
        ))
        if out_series is not None:
            x_out, y_out, hover_out = _canon_day_series(
                df.rename(columns={out_series.name: "_out"}), "_out"
            )
            fig.add_trace(go.Scatter(
                x=x_out, y=y_out, mode="lines", name="Outdoor Temp",
                line=dict(color=_C_OUT, width=2),
                meta="keep_color",
                hovertext=hover_out,
                hovertemplate="%{hovertext}<br>%{y:.1f}°C<extra></extra>",
            ))
        fig.update_xaxes(title="Day (1..365)", tickmode="linear", tick0=1, dtick=10)

    elif mode == "raw":
        # Raw manté el pas original de simulació, útil per revisar pics puntuals.
        df = filter_by_season(base, season)
        x_vals, y_in, hover_in = _raw_axis_data(df, "air_temperature")
        fig.add_trace(go.Scatter(
            x=x_vals, y=y_in, mode="lines", name="Indoor Temp",
            line=dict(color="rgba(79, 142, 247, 0.9)", width=2.2, shape="spline"),
            meta="keep_color",
            hovertext=hover_in,
            hovertemplate="%{hovertext}<br>%{y:.1f}°C<extra></extra>",
        ))
        if out_series is not None:
            x_v2, y_out, hover_out = _raw_axis_data(
                df.rename(columns={out_series.name: "_out"}), "_out"
            )
            fig.add_trace(go.Scatter(
                x=x_v2, y=y_out, mode="lines", name="Outdoor Temp",
                line=dict(color=_C_OUT, width=1.8, shape="spline"),
                meta="keep_color",
                hovertext=hover_out,
                hovertemplate="%{hovertext}<br>%{y:.1f}°C<extra></extra>",
            ))
        fig.update_xaxes(title="Time (step)")

    fig.update_yaxes(title="Temperature (°C)")
    fig.update_layout(title=f"Indoor vs Outdoor Temperature – {mode.capitalize()} ({season})")
    return fig


# Humitat interior amb humitat exterior de referència.
def make_indoor_humidity_plot(obs: pd.DataFrame, mode: str, season: str) -> go.Figure:
    """Dibuixa la humitat interior com a barres o línies i l'exterior amb marcadors."""
    base = _ensure_datetime_index(obs.copy())
    out_series = base["outdoor_humidity"] if "outdoor_humidity" in base.columns else None

    _C_IN_LINE = "rgba(0, 200, 255, 1.0)"
    _C_OUT_BAR = "rgba(15, 45, 120, 0.75)"

    fig = go.Figure()
    if "air_humidity" not in base.columns:
        return go.Figure()

    if mode in ("hour", "month"):
        # Igual que amb temperatura: agregats per hora/mes, amb humitat interior com a
        # línia principal i exterior com a referència.
        if mode == "hour":
            df = filter_by_season(base, season)
            grp = _hour_series(df)
            x_vals = list(range(24))
            x_title, tick_cfg = "Hour of Day", dict(tickmode="linear", dtick=1)
        else:
            df = base if season == "All" else filter_by_season(base, season)
            grp = _month_series(df)
            x_vals = list(range(1, 13))
            x_title, tick_cfg = "Month", dict(tickmode="linear", dtick=1)

        g_mean = df.groupby(grp)["air_humidity"].mean()
        y_in   = [g_mean.get(v, float("nan")) for v in x_vals]

        fig.add_trace(go.Scatter(
            x=x_vals, y=y_in, mode="lines+markers", name="Indoor Humidity",
            line=dict(color=_C_IN_LINE, width=2.5),
            marker=dict(color=_C_IN_LINE, size=7, symbol="circle",
                        line=dict(color="white", width=1.5)),
            meta="keep_color",
            hovertemplate="%{x}<br>Mitjana: %{y:.1f}%<extra></extra>",
        ))
        if out_series is not None:
            g_out = df.groupby(grp)[out_series.name].mean()
            y_out = [g_out.get(v, float("nan")) for v in x_vals]
            fig.add_trace(go.Bar(
                x=x_vals, y=y_out, name="Outdoor Humidity",
                marker_color=_C_OUT_BAR,
                meta="keep_color",
                hovertemplate="%{x}<br>%{y:.1f}%<extra></extra>",
            ))
        fig.update_xaxes(title=x_title, **tick_cfg)
        fig.update_layout(bargap=0.22)

    elif mode == "day":
        df = filter_by_season(base, season)
        x_in, y_in, hover_in = _canon_day_series(df, "air_humidity")
        fig.add_trace(go.Scatter(
            x=x_in, y=y_in, mode="lines", name="Indoor Humidity",
            line=dict(color=_C_IN_LINE, width=2),
            meta="keep_color",
            hovertext=hover_in,
            hovertemplate="%{hovertext}<br>%{y:.1f}%<extra></extra>",
        ))
        if out_series is not None:
            x_out, y_out, hover_out = _canon_day_series(
                df.rename(columns={out_series.name: "_out"}), "_out"
            )
            fig.add_trace(go.Scatter(
                x=x_out, y=y_out, mode="lines", name="Outdoor Humidity",
                line=dict(color=_C_OUT_BAR, width=2),
                meta="keep_color",
                hovertext=hover_out,
                hovertemplate="%{hovertext}<br>%{y:.1f}%<extra></extra>",
            ))
        fig.update_xaxes(title="Day (1..365)", tickmode="linear", tick0=1, dtick=10)

    elif mode == "raw":
        # Raw és útil per detectar pics d'humitat que quedarien amagats en la mitjana.
        df = filter_by_season(base, season)
        x_vals, y_in, hover_in = _raw_axis_data(df, "air_humidity")
        fig.add_trace(go.Scatter(
            x=x_vals, y=y_in, mode="lines", name="Indoor Humidity",
            line=dict(color=_C_IN_LINE, width=2.2, shape="spline"),
            meta="keep_color",
            hovertext=hover_in,
            hovertemplate="%{hovertext}<br>%{y:.1f}%<extra></extra>",
        ))
        if out_series is not None:
            x_v2, y_out, hover_out = _raw_axis_data(
                df.rename(columns={out_series.name: "_out"}), "_out"
            )
            fig.add_trace(go.Scatter(
                x=x_v2, y=y_out, mode="lines", name="Outdoor Humidity",
                line=dict(color=_C_OUT_BAR, width=1.8, shape="spline"),
                meta="keep_color",
                hovertext=hover_out,
                hovertemplate="%{hovertext}<br>%{y:.1f}%<extra></extra>",
            ))
        fig.update_xaxes(title="Time (step)")

    fig.update_yaxes(title="Humidity (%)")
    fig.update_layout(title=f"Indoor vs Outdoor Humidity – {mode.capitalize()} ({season})")
    return fig


# Temperatures de consigna de calefacció i refrigeració.
def _with_global_setpoint_columns(base: pd.DataFrame) -> tuple[pd.DataFrame, bool]:
    """Afegeix consignes globals mitjanes quan les dades venen per zona."""
    htg_cols = [column for column in base.columns if "htg_setpoint" in column.lower()]
    clg_cols = [column for column in base.columns if "clg_setpoint" in column.lower()]

    if htg_cols and "htg_setpoint" not in base.columns:
        base["htg_setpoint"] = base[htg_cols].mean(axis=1)
    if clg_cols and "clg_setpoint" not in base.columns:
        base["clg_setpoint"] = base[clg_cols].mean(axis=1)
    return base, bool(htg_cols or clg_cols)


def _build_setpoint_trace(
    base: pd.DataFrame,
    series_name: str,
    label: str,
    line_style: dict,
    mode: str,
    season: str,
) -> go.Scatter | None:
    """Crea una traça relacionada amb el punt de consigna per al mode d'agregació seleccionat."""
    trace = build_mode_scatter_trace(
        base,
        series_name,
        mode,
        season,
        name=label,
        color=line_style["color"],
        units="°C",
        line_width=float(line_style.get("width", 2.4)),
        line_shape=line_style.get("shape"),
        hovertemplate="%{hovertext}<br>%{y:.2f}°C<extra></extra>",
    )
    return trace


def make_setpoints_plot(obs: pd.DataFrame, mode: str, season: str) -> go.Figure:
    """Crea una figura Plotly per als punts de consigna."""
    base = _ensure_datetime_index(obs).copy()

    base, has_setpoints = _with_global_setpoint_columns(base)
    if not has_setpoints:
        return go.Figure()

    fig = go.Figure()

    htg_color = pick_semantic_trace_color("Heating Setpoint")
    clg_color = pick_semantic_trace_color("Cooling Setpoint")
    htg_style = dict(color=htg_color, width=3.2, shape="hv" if mode in ("hour", "raw") else "linear")
    clg_style = dict(color=clg_color, width=3.2, shape="hv" if mode in ("hour", "raw") else "linear")
    tr1 = _build_setpoint_trace(base, "htg_setpoint", "Heating Setpoint", htg_style, mode, season)
    tr2 = _build_setpoint_trace(base, "clg_setpoint", "Cooling Setpoint", clg_style, mode, season)
    if tr1: fig.add_trace(tr1)
    if tr2: fig.add_trace(tr2)

    _apply_mode_xaxis(fig, mode)

    fig.update_yaxes(title="Temperature (°C)")
    fig.update_layout(title=f"Setpoints – {mode.capitalize()} ({season})")
    return fig


# Comparació entre consignes i temperatura interior.
def make_setpoints_vs_indoor_plot(obs: pd.DataFrame, mode: str, season: str) -> go.Figure:
    """Crea una figura Plotly per a punts de consigna vs interior."""
    base = _ensure_datetime_index(obs).copy()

    base, has_setpoints = _with_global_setpoint_columns(base)
    if not has_setpoints and "air_temperature" not in base.columns:
        return go.Figure()

    fig = go.Figure()

    htg_color = pick_semantic_trace_color("Heating Setpoint")
    clg_color = pick_semantic_trace_color("Cooling Setpoint")
    ind_color = pick_semantic_trace_color("indoor")

    htg_style = dict(color=htg_color, width=3.2, shape="hv" if mode in ("hour", "raw") else "linear")
    clg_style = dict(color=clg_color, width=3.2, shape="hv" if mode in ("hour", "raw") else "linear")
    ind_style = dict(color=ind_color, width=2.0, shape="spline" if mode == "raw" else "linear")

    tr1 = _build_setpoint_trace(base, "htg_setpoint", "Heating Setpoint", htg_style, mode, season)
    tr2 = _build_setpoint_trace(base, "clg_setpoint", "Cooling Setpoint", clg_style, mode, season)
    tr3 = _build_setpoint_trace(base, "air_temperature", "Indoor Temperature", ind_style, mode, season)

    if tr1: fig.add_trace(tr1)
    if tr2: fig.add_trace(tr2)
    if tr3: fig.add_trace(tr3)

    _apply_mode_xaxis(fig, mode)

    fig.update_yaxes(title="Temperature (°C)")
    fig.update_layout(title=f"Setpoints vs Indoor Temperature – {mode.capitalize()} ({season})")
    return fig
\end{lstlisting}

\Needspace{8\baselineskip}

\subsection{\texorpdfstring{\texttt{backend/grafics/time\_utils.py}}{backend/grafics/time\_utils.py}}

\begin{lstlisting}[style=studioPython,caption={backend/grafics/time\_utils.py},label={lst:backend-grafics-time-utils-py}]
"""Filtrat temporal i ajuda d'eix per a figures del panell."""

from __future__ import annotations

import numpy as np
import pandas as pd

from backend.grafics.style import _alpha_color


def _bounded_int_series(
    df: pd.DataFrame,
    column: str,
    lower: int,
    upper: int,
) -> pd.Series:
    """Retorna una sèrie entera anul·lable amb valors temporals impossibles eliminats."""

    values = pd.to_numeric(df[column], errors="coerce").round()
    values = values.where(values.between(lower, upper))
    return values.astype("Int64")


def _coerce_temporal_parts(df: pd.DataFrame) -> tuple[pd.Series, pd.Series, pd.Series]:
    """Retorna les columnes mes/dia/hora netejades per a la reconstrucció de data i hora."""

    month = _bounded_int_series(df, "month", 1, 12)
    day = _bounded_int_series(df, "day_of_month", 1, 31)
    hour = _bounded_int_series(df, "hour", 0, 23)
    return month, day, hour


def sanitize_observation_time_columns(df: pd.DataFrame) -> pd.DataFrame:
    """Converteix les columnes d'observació temporal i emmascarar valors físicament impossibles."""

    result = df.copy()
    bounds = {
        "month": (1, 12),
        "day_of_month": (1, 31),
        "hour": (0, 23),
    }
    for column, (lower, upper) in bounds.items():
        if column in result.columns:
            result[column] = _bounded_int_series(result, column, lower, upper)
    return result


def _infer_timestep_hours(df: pd.DataFrame) -> float:
    """Infereix la durada del timestep en hores a partir del temps disponible."""
    if "time_elapsed(hours)" in df.columns:
        elapsed = pd.to_numeric(df["time_elapsed(hours)"], errors="coerce")
        elapsed_diffs = elapsed.diff().dropna()
        elapsed_diffs = elapsed_diffs[(elapsed_diffs > 0) & np.isfinite(elapsed_diffs)]
        if not elapsed_diffs.empty:
            return float(elapsed_diffs.median())

    time_col = None
    if "timestamp" in df.columns:
        time_col = pd.to_datetime(df["timestamp"], errors="coerce")
    elif "datetime" in df.columns:
        time_col = pd.to_datetime(df["datetime"], errors="coerce")

    if time_col is not None:
        time_diffs = time_col.diff().dropna().dt.total_seconds() / 3600.0
        time_diffs = time_diffs[(time_diffs > 0) & np.isfinite(time_diffs)]
        if not time_diffs.empty:
            return float(time_diffs.median())

    if isinstance(df.index, pd.DatetimeIndex) and len(df.index) > 1:
        diffs = pd.Series(df.index, index=df.index).diff().dropna()
        diffs_h = diffs.dt.total_seconds() / 3600.0
        diffs_h = diffs_h[(diffs_h > 0) & np.isfinite(diffs_h)]
        if not diffs_h.empty:
            return float(diffs_h.median())

    if all(col in df.columns for col in ("month", "day_of_month", "hour")):
        mdh = df[["month", "day_of_month", "hour"]].apply(pd.to_numeric, errors="coerce")
        valid = mdh.dropna()
        if not valid.empty:
            hour_blocks = valid.ne(valid.shift()).any(axis=1).cumsum()
            block_sizes = hour_blocks.value_counts().sort_index()
            block_sizes = block_sizes[block_sizes > 0]
            if not block_sizes.empty:
                return 1.0 / float(block_sizes.median())

    return 0.25


def _seasonal_marker_colors(
    x_vals: list[int],
    season: str,
    *,
    active_color: str,
    muted_color: str | None = None,
) -> list[str]:
    """Seasonal marker colors."""
    months_sel = _season_months(season)
    if not months_sel:
        return [active_color] * len(x_vals)
    dim = muted_color if muted_color is not None else _alpha_color(active_color, 0.30)
    return [active_color if month in months_sel else dim for month in x_vals]

def _season_months(name: str):
    """Season months."""
    seasons = {
        "Winter": [12, 1, 2],
        "Spring": [3, 4, 5],
        "Summer": [6, 7, 8],
        "Autumn": [9, 10, 11],
    }
    return seasons.get(name, None)  # None -> Tot


def _ensure_datetime_index(df: pd.DataFrame) -> pd.DataFrame:
    """
    Garanteix un DatetimeIndex per reagrupar per hores/dies/mesos.
    Si no hi ha 'datetime' ni 'timestamp', reconstrueix un índex sintètic
    coherent amb els timesteps de 15 minuts del projecte.
    """
    dfc = df.copy()
    if isinstance(dfc.index, pd.DatetimeIndex):
        dfc = dfc[~dfc.index.isna()]
        return dfc.sort_index(kind="mergesort")
    if "datetime" in dfc.columns:
        dfc["datetime"] = pd.to_datetime(dfc["datetime"], errors="coerce")
        dfc = dfc.set_index("datetime").sort_index()
        dfc = dfc[~dfc.index.isna()]
        return dfc
    if "timestamp" in dfc.columns:
        dfc["datetime"] = pd.to_datetime(dfc["timestamp"], errors="coerce")
        dfc = dfc.set_index("datetime").sort_index()
        dfc = dfc[~dfc.index.isna()]
        return dfc

    if all(col in dfc.columns for col in ("month", "day_of_month", "hour")):
        month, day, hour = _coerce_temporal_parts(dfc)

        # Ancorem la sèrie en el primer registre complet: així respectem simulacions que no
        # comencen exactament a les 00:00.
        if "time_elapsed(hours)" in dfc.columns:
            elapsed = pd.to_numeric(dfc["time_elapsed(hours)"], errors="coerce")
            valid = elapsed.notna() & month.notna() & day.notna() & hour.notna()
            if valid.any():
                first_idx = valid[valid].index[0]
                first_pos = dfc.index.get_loc(first_idx)
                first_elapsed = float(elapsed.iloc[first_pos])
                base_minutes = int(round((first_elapsed % 1.0) * 60.0)) % 60
                base_ts = pd.Timestamp(
                    year=2001,
                    month=int(month.iloc[first_pos]),
                    day=int(day.iloc[first_pos]),
                    hour=int(hour.iloc[first_pos]),
                    minute=base_minutes,
                )
                offset_minutes = ((elapsed - first_elapsed) * 60.0).round()
                try:
                    dt = base_ts + pd.to_timedelta(offset_minutes, unit="m")
                    dfc.index = pd.DatetimeIndex(dt)
                    dfc = dfc[~dfc.index.isna()]
                    return dfc.sort_index(kind="mergesort")
                except (OverflowError, ValueError):
                    pass

        hour_blocks = (
            pd.DataFrame({"month": month, "day": day, "hour": hour})
            .ne(pd.DataFrame({"month": month, "day": day, "hour": hour}).shift())
            .any(axis=1)
            .cumsum()
        )
        # Quan no hi ha temps acumulat, repartim cada bloc d'una mateixa hora en quarts
        # d'hora consecutius, que és el pas habitual dels CSV de Sinergym.
        minute_slot = hour_blocks.groupby(hour_blocks).cumcount() * 15
        dt = pd.to_datetime(
            {
                "year": 2001,
                "month": month.astype("float64"),
                "day": day.astype("float64"),
                "hour": hour.astype("float64"),
                "minute": minute_slot.astype("float64"),
            },
            errors="coerce",
        )
        dfc.index = pd.DatetimeIndex(dt)
        dfc = dfc[~dfc.index.isna()]
        return dfc.sort_index(kind="mergesort")

    # Reserva defensiva: mantenim 15 minuts per pas.
    dfc.index = pd.date_range("2001-01-01", periods=len(dfc), freq="15min")
    return dfc.sort_index(kind="mergesort")

def _month_series(df: pd.DataFrame) -> pd.Series:
    """
    Sèrie 'mes' canònica 1..12:
      - Prioritza columna 'month' si existeix (normalitza possibles 0/13).
      - Si no, usa el DatetimeIndex.
    """
    if "month" in df.columns:
        s = _bounded_int_series(df, "month", 1, 12)
        if s.notna().any():
            return pd.Series(s, index=df.index)
    if isinstance(df.index, pd.DatetimeIndex):
        return pd.Series(df.index.month, index=df.index)
    # amb índex sintètic ja és DateTimeIndex
    return pd.Series(pd.to_datetime(df.index).month, index=df.index)


def _hour_series(df: pd.DataFrame) -> pd.Series:
    """
    Sèrie 'hora' canònica 0..23:
      - Prioritza columna 'hour' si existeix (normalitza 24→0).
      - Si no, usa el DatetimeIndex.
    """
    if "hour" in df.columns:
        h = _bounded_int_series(df, "hour", 0, 23)
        if h.notna().any():
            return pd.Series(h, index=df.index)
    if isinstance(df.index, pd.DatetimeIndex):
        return pd.Series(df.index.hour, index=df.index)
    return pd.Series(pd.to_datetime(df.index).hour, index=df.index)


def filter_by_season(df: pd.DataFrame, season: str) -> pd.DataFrame:
    """Filtra per estació amb la sèrie de mesos canònica."""
    months = _season_months(season)
    if not months:
        return df
    mon = _month_series(df)
    return df[mon.isin(months)]


def _get_outdoor_temperature_series(df: pd.DataFrame) -> pd.Series | None:
    """Retorna la columna de temperatura exterior utilitzada per les observacions generades."""
    for col in ("outdoor_temperature", "outdoor_air_temperature"):
        if col in df.columns:
            return pd.to_numeric(df[col], errors="coerce")
    return None


def _has_mdh(df: pd.DataFrame) -> bool:
    """Té columnes month/day_of_month/hour del CSV?"""
    return all(c in df.columns for c in ("month", "day_of_month", "hour"))


def _canon_day_series(df: pd.DataFrame, col: str):
    """
    Retorna (x, y, hover) per a una sèrie agregada DIÀRIA sense anys:
      - Si hi ha columnes 'month' i 'day_of_month', fa mean per (MM,DD).
      - Si només hi ha DatetimeIndex, resample('D') i mitjana per (MM,DD) (uneix anys).
    Sempre torna UNA observació per MM-DD i X=1..365.
    """
    if _has_mdh(df):
        month, day, _ = _coerce_temporal_parts(df)
        working = df.assign(month=month, day_of_month=day)
        g = working.groupby(["month", "day_of_month"])[col].mean().reset_index()
        g = g.dropna(subset=["month", "day_of_month"])
        g["order"] = pd.to_datetime(
            {"year": 2001, "month": g["month"].astype(int), "day": g["day_of_month"].astype(int)},
            errors="coerce",
        )
        g = g.sort_values("order")
        x = g["order"].dt.dayofyear.tolist()
        y = g[col].to_list()
        hover = [f"{int(m):02d}-{int(d):02d}" for m, d in zip(g["month"], g["day_of_month"])]
        return x, y, hover

    if isinstance(df.index, pd.DatetimeIndex):
        s = df[col].resample("D").mean()
        g = s.groupby([s.index.month, s.index.day]).mean()
        months = np.array([i[0] for i in g.index], dtype=int)
        days = np.array([i[1] for i in g.index], dtype=int)
        order = pd.to_datetime({"year": 2001, "month": months, "day": days}, errors="coerce")
        ord_idx = np.argsort(order.values)
        x = order.dt.dayofyear.to_numpy()[ord_idx].tolist()
        y = g.to_numpy()[ord_idx].tolist()
        hover = [f"{int(m):02d}-{int(d):02d}" for m, d in zip(months[ord_idx], days[ord_idx])]
        return x, y, hover

    # Reserva defensiva (24 registres ~ 1 dia)
    ids = np.arange(len(df)) // 24 + 1
    g = pd.Series(df[col].to_numpy()).groupby(ids).mean()
    x = g.index.tolist()
    y = g.values.tolist()
    hover = [f"Day {i}" for i in x]
    return x, y, hover


def _canon_day_total_series(df: pd.DataFrame, col: str):
    """Retorna (x, y, hover) per a totals diaris en un eix canònic MM-DD."""
    if isinstance(df.index, pd.DatetimeIndex):
        daily = pd.to_numeric(df[col], errors="coerce").resample("D").sum(min_count=1)
        g = daily.groupby([daily.index.month, daily.index.day]).mean()
        months = np.array([idx[0] for idx in g.index], dtype=int)
        days = np.array([idx[1] for idx in g.index], dtype=int)
        order = pd.to_datetime({"year": 2001, "month": months, "day": days}, errors="coerce")
        ord_idx = np.argsort(order.values)
        x = order.dt.dayofyear.to_numpy()[ord_idx].tolist()
        y = g.to_numpy()[ord_idx].tolist()
        hover = [
            f"{int(month):02d}-{int(day):02d}"
            for month, day in zip(months[ord_idx], days[ord_idx])
        ]
        return x, y, hover

    daily_ids = np.arange(len(df)) // 24 + 1
    g = pd.to_numeric(pd.Series(df[col].to_numpy()), errors="coerce").groupby(daily_ids).sum()
    x = g.index.tolist()
    y = g.values.tolist()
    hover = [f"Day {day}" for day in x]
    return x, y, hover


def _raw_axis_data(df: pd.DataFrame, col: str):
    """X=1..N (pas), hover='MM-DD HHh' si es pot."""
    x = list(range(1, len(df) + 1))
    if _has_mdh(df):
        hover = [
            (
                f"{int(m):02d}-{int(d):02d} {int(h):02d}h"
                if pd.notna(m) and pd.notna(d) and pd.notna(h)
                else f"Step {idx}"
            )
            for idx, (m, d, h) in enumerate(
                zip(df["month"], df["day_of_month"], df["hour"]),
                start=1,
            )
        ]
    elif isinstance(df.index, pd.DatetimeIndex):
        hover = [f"{ix.month:02d}-{ix.day:02d} {ix.hour:02d}h" for ix in df.index]
    else:
        hover = [f"Step {i}" for i in x]
    y = df[col].to_list()
    return x, y, hover


def _mode_axis_config(mode: str) -> tuple[str, dict]:
    """Retorna el títol i la configuració d'eix X per al mode temporal."""
    if mode == "hour":
        return "Hour of Day", {"tickmode": "linear", "dtick": 1}
    if mode == "day":
        return "Day (1..365)", {"tickmode": "linear", "tick0": 1, "dtick": 10}
    if mode == "month":
        return "Month", {"tickmode": "linear", "dtick": 1}
    return "Time (step)", {}


def _apply_mode_xaxis(fig, mode: str, **kwargs):
    """Aplica a una figura Plotly l'eix X habitual del mode temporal."""
    title, axis_kwargs = _mode_axis_config(mode)
    axis_kwargs.update(kwargs)
    return fig.update_xaxes(title=title, **axis_kwargs)


def _series_for_mode(
    base: pd.DataFrame,
    column: str,
    mode: str,
    season: str,
    *,
    sign: float = 1.0,
) -> tuple[list, list, list | None, str] | None:
    """Retorna dades x/y/hover/trace-mode per a una sèrie mitjana segons el mode."""
    if column not in base.columns:
        return None

    if mode == "hour":
        df = filter_by_season(base, season)
        if df.empty:
            return None
        grouped = df.groupby(_hour_series(df))[column].mean() * sign
        x_values = list(range(24))
        return x_values, [grouped.get(hour, float("nan")) for hour in x_values], None, "lines+markers"

    if mode == "day":
        df = filter_by_season(base, season)
        if df.empty:
            return None
        x_values, y_values, hover = _canon_day_series(df, column)
        return x_values, [value * sign for value in y_values], hover, "lines"

    if mode == "month":
        df = base if season == "All" else filter_by_season(base, season)
        if df.empty:
            return None
        grouped = df.groupby(_month_series(df))[column].mean() * sign
        x_values = list(range(1, 13))
        return x_values, [grouped.get(month, float("nan")) for month in x_values], None, "lines+markers"

    if mode == "raw":
        df = filter_by_season(base, season)
        if df.empty:
            return None
        x_values, y_values, hover = _raw_axis_data(df, column)
        return x_values, [value * sign for value in y_values], hover, "lines"

    return None


def _auto_scale_series(dfcol: pd.Series):
    """Retorna sèrie escalada i sufix ('', '×1e3', '×1e6', '×1e9')."""
    absmax = float(np.nanmax(np.abs(dfcol.to_numpy()))) if len(dfcol) else 0.0
    if absmax >= 1e9: return dfcol / 1e9, "×1e9"
    if absmax >= 1e6: return dfcol / 1e6, "×1e6"
    if absmax >= 1e3: return dfcol / 1e3, "×1e3"
    return dfcol, ""
\end{lstlisting}

\end{document}
